\documentclass[12pt, openany, twoside]{book}

% --- LANGUAGE & FONT SETUP ---
\usepackage{fontspec}
\usepackage{polyglossia}
\setdefaultlanguage{sanskrit}
\usepackage{etoolbox}

% --- COLOR SETUP ---
\usepackage[dvipsnames]{xcolor} % <--- ADDED for \textcolor support

% --- CUSTOM COLORS (Matching HTML color scheme conditionally) ---
\definecolor{AccentBlue}{HTML}{000000}      
\definecolor{AccentPurple}{HTML}{000000}    
\definecolor{AccentGreen}{HTML}{000000}     
\definecolor{AccentBrown}{HTML}{000000}     
\definecolor{SwaraRed}{HTML}{000000}        
\definecolor{ModifierDarkBlue}{HTML}{000000}
\definecolor{ModifierGreen}{HTML}{000000}
\definecolor{ModifierSkyBlue}{HTML}{000000}
\definecolor{vedicmaroon}{RGB}{0, 0, 0}

% Default: AdishilaVedic
\setmainfont[Scale=1.35, Script=Devanagari, Renderer=Harfbuzz, Ligatures=TeX, Path=C:/Users/sekha/OneDrive/Documents/GitHub/jaimineeyasamavedam/fonts/, AutoFakeBold=1.18, BoldFont=AdishilaVedicBold.ttf]{AdishilaVedic.ttf}
\newfontfamily\devafont[Scale=1.35, Script=Devanagari, Renderer=Harfbuzz, Ligatures=TeX, Path=C:/Users/sekha/OneDrive/Documents/GitHub/jaimineeyasamavedam/fonts/, AutoFakeBold=1.18, BoldFont=AdishilaVedicBold.ttf]{AdishilaVedic.ttf}
\newfontfamily\headerfont[Scale=1.45, Script=Devanagari, Renderer=Harfbuzz, Ligatures=TeX, Path=C:/Users/sekha/OneDrive/Documents/GitHub/jaimineeyasamavedam/fonts/]{AdishilaVedicBold.ttf}
\newfontfamily\smallredfont[Scale=0.95, Script=Devanagari, Renderer=Harfbuzz, Ligatures=TeX, Color=SwaraRed, Path=C:/Users/sekha/OneDrive/Documents/GitHub/jaimineeyasamavedam/fonts/]{AdishilaVedic.ttf}
\newfontfamily\swarafont[Scale=0.88, Script=Grantha, Ligatures=TeX, AutoFakeBold=3.5, AutoFakeSlant, Path=C:/Users/sekha/OneDrive/Documents/GitHub/jaimineeyasamavedam/fonts/]{JaimineeyaSwara.ttf}
\newfontfamily\arcfontDeva[Scale=1.25, Script=Grantha, Ligatures=TeX, AutoFakeBold=3.5, Path=C:/Users/sekha/OneDrive/Documents/GitHub/jaimineeyasamavedam/fonts/]{JaimineeyaSwara.ttf}

\setcounter{tocdepth}{1} % Printed TOC: Level 0 (Chapter) and Level 1 (Section) only

% --- PACKAGES ---
\usepackage{ragged2e} 
\usepackage{stackengine} 
\usepackage{fancyhdr}    
\usepackage[perpage,norule]{footmisc} % Resets footnote numbering on every page - MUST BE BEFORE hyperref
\usepackage{hyperref}    
\usepackage{bookmark}
\bookmarksetup{depth=2, numbered=false} % PDF Bookmarks Navigation: Includes individual Arsheyams (Subsections)
\usepackage{imakeidx}    
\usepackage{needspace}
\usepackage{longtable}
\usepackage{titlesec}
\titleformat{\chapter}[block]{\normalfont\huge\bfseries\centering}{}{0pt}{\headerfont}
\titlespacing*{\chapter}{0pt}{-27.5pt}{-4.7pt}
\titleformat{\section}[block]{\normalfont\Large\bfseries\centering}{}{0pt}{\headerfont}
\titlespacing*{\section}{0pt}{0.0pt}{4pt}
\usepackage{setspace}
\setstretch{1.90}
% --- CRITICAL: MICROTYPE FOR SPACING ---
% This allows LaTeX to stretch letters by 20% to avoid gaps between words 
\usepackage[stretch=20, shrink=20]{microtype}  % <--- THE MAGIC FIX for spacing

% --- GEOMETRY ---
\usepackage[
  paperwidth=180mm,
  paperheight=228mm,
  twoside,
  inner=18mm,       % Inside / Spine margin (8mm inside border)
  outer=15mm,       % Outside margin (5mm inside border)
  top=19mm,         % Top margin including header (9mm inside border)
  bottom=19mm,      % Bottom margin including footer (7mm inside border)
  headheight=22pt,
  headsep=15.0pt,
  includehead,
  includefoot
]{geometry} 

% --- FOOTNOTE LINE SETUP ---
\newcommand{\drawstarline}{%
  \begin{center}
    \vspace{-0.2em}
    \textcolor{SwaraRed}{* * *}
    \vspace{-0.5em}
  \end{center}
}

% --- SPACING & BREAKING RULES ---
% User Rule 1: No word breaking across lines (strictly enforced globally and per-language)
\hyphenpenalty=10000     
\exhyphenpenalty=10000
\appto\captionssanskrit{\hyphenpenalty=10000\relax\exhyphenpenalty=10000\relax}
\AtBeginDocument{\hyphenpenalty=10000\relax\exhyphenpenalty=10000\relax}

% --- HYPHENATION & SPACING RULES (CRITICAL) ---
\lefthyphenmin=3      % Never split first 2 letters (e.g., prevents "ha-")
\righthyphenmin=3     % Never split last 2 letters (e.g., prevents "-ih")
\tolerance=500       % Lower tolerance forces better breaks (standard is 200)
\pretolerance=100     % Tries to break without hyphens first
\setlength{\emergencystretch}{3em} % Allows line breaks without spilling into margin
\setlength{\hfuzz}{2pt}
\setlength{\parindent}{0pt} 

% --- STACKING COMMANDS (Strictly No Phantom Spaces) ---

\newlength{\swarabasew}
\newlength{\swarasw}
% 1. STANDARD: Swara markings BELOW the mantra text
\newcommand{\swarastack}[2]{%
  \settowidth{\swarabasew}{#1}%
  \settowidth{\swarasw}{#2}%
  \ifdim\swarasw>\swarabasew
    \mbox{\hspace{0.5\swarasw-0.5\swarabasew}#1\hspace{-0.5\swarabasew-0.5\swarasw}\makebox[\swarasw][c]{\raisebox{-1.05ex}[0pt][0pt]{#2}}}%
  \else
    \mbox{#1\hspace{-\swarabasew}\makebox[\swarabasew][c]{\raisebox{-1.05ex}[0pt][0pt]{#2}}}%
  \fi
}
\newcommand{\stackcenter}[2]{\swarastack{#1}{#2}}
\newcommand{\stackleft}[2]{\swarastack{#1}{#2}}

% --- DANDA WITH ARC MACRO (MOD-A1: Centered Over Danda) ---
% AdishilaVedic: Stem is at 0.802*dimen0. Balance whitespace around danda:
\newcommand{\dandaWithArc}[2][\char"E00D]{%
  \begingroup%
  \settowidth{\dimen0}{#2}%
  \hspace{-0.30\dimen0}%
  \mbox{\hspace{0.802\dimen0}\makebox[0pt][c]{{\arcfontDeva \textcolor{ModifierSkyBlue}{\raisebox{0.92ex}{#1}}}}\hspace{-0.802\dimen0}#2}%
  \hspace{0.30\dimen0}%
  \endgroup%
}

% --- OVERHEAD CONJUNCT ARC (MOD-A2: Centered Over Syllable) ---
\newcommand{\arcOverSyllable}[2][\char"E02E]{%
  \begingroup%
  \settowidth{\dimen0}{#2}%
  \mbox{\hspace{0.5\dimen0}\makebox[0pt][c]{{\Large\swarafont \textcolor{ModifierSkyBlue}{\textbf{\raisebox{1.18ex}{#1}}}}}\hspace{-0.5\dimen0}#2}%
  \endgroup%
}

% --- LOWER UNDER-SLASH (MOD-G: Positioned Below-Right of Syllable) ---
\newcommand{\modGUnder}[1]{%
  \begingroup%
  \settowidth{\dimen0}{#1}%
  \mbox{#1\hspace{-0.5\dimen0}\makebox[0pt][c]{\swarafont\textcolor{ModifierSkyBlue}{\textbf{\raisebox{-0.30ex}[0pt][0pt]{\hspace{0.45em}\char"E003}}}}\hspace{0.5\dimen0}}%
  \endgroup%
}

% --- UNDERBAR MODIFIER (MOD-UNDERBAR: In-line Under-dash to the Right) ---
\newcommand{\underbarMark}{\hspace{0.05em}{\textcolor{ModifierSkyBlue}{\raisebox{-0.22ex}{\rule{0.35em}{0.28ex}}}}\hspace{0.06em}}

% --- PEAK ELEVATION CARET (MOD-B: Spans 2 Syllables with Swara on Apex) ---
\newcommand{\caretWithSwara}[1]{%
  \rlap{%
    \swarafont \textcolor{ModifierSkyBlue}{\raisebox{1.18ex}{\char"E005}}%
    \settowidth{\dimen0}{\swarafont\char"E005}%
    \hspace{-0.5\dimen0}%
    \makebox[0pt][c]{{\devafont\raisebox{2.35ex}[0pt][0pt]{{\smallredfont \textcolor{SwaraRed}{#1}}}}}%
    \hspace{0.5\dimen0}%
  }%
}

% --- INDEX INITIALIZATION ---
\indexsetup{level=\chapter*, toclevel=chapter, othercode={\footnotesize\singlespacing}}
\makeindex[columns=2, title=वर्णानुक्रमणिका, intoc]

% --- HYPERREF ---
\hypersetup{
    colorlinks=true,
    linkcolor=black,     
    filecolor=magenta,      
    urlcolor=cyan,
    pdftitle={Vedic Text},
    pdfpagemode=UseOutlines, 
}



% --- STACK CONFIG ---
\def\stackalignment{c} 
\def\stacktype{L}     
\setstackgap{L}{0.9\baselineskip} 

% --- HEADER CONFIG ---
\setlength{\headheight}{22pt} 
\addtolength{\topmargin}{-10pt}
\pagestyle{fancy}
\fancyhf{}
\fancyhead[LE,RO]{\thepage}
\fancyhead[LO]{\rightmark}   
\fancyhead[RE]{\leftmark}    
\renewcommand{\headrulewidth}{0.4pt}

% --- FOOTER CONFIG ---
\fancyfoot[C]{\footnotesize jaimineeyasamavedam.org (v1.03)}
\renewcommand{\footrulewidth}{0.4pt}
% Jinja helpers
\newcommand{\VAR}[1]{}
\newcommand{\BLOCK}[1]{}

% Command to format accent marks with larger size and bold
\newcommand{\accentmark}[2]{%
    {\fontsize{#1pt}{#1pt}\selectfont\bfseries\addfontfeature{FakeBold=3}#2}%
}

% --- ACCENT OVERLAP ADJUSTMENT ---
% Defines a small negative space (kerning) to pull the next character closer
% Used when two "danger" accents (Anudatta/Kampa) appear consecutively.
\newcommand{\accentadj}{\kern0.15ex}

% newunicodechar overrides removed to ensure strict zero-space ||N|| and tight danda binding
% --- CUSTOM FOOTNOTE STYLE: Devanagari Numerals & Higher Position ---
\makeatletter
\def\devanagarinumeral#1{\expandafter\dodevanagari\number#1\relax}
\def\dodevanagari#1{\ifx#1\relax\else\ifcase#1०\or१\or२\or३\or४\or५\or६\or७\or८\or९\fi\expandafter\dodevanagari\fi}
\renewcommand{\thefootnote}{\devanagarinumeral{\value{footnote}}}
\renewcommand\@makefnmark{\hbox{\textsuperscript{\raisebox{1.2ex}{\normalfont\@thefnmark}}}}
\makeatother
% ------------------------

% --- DECORATIVE DAISY PAGE BORDER ---
% Seamless floral daisy chain along all four edges with zero corner whitespace gaps.
\usepackage{tikz}
\usepackage{eso-pic}
\usepackage{bbding}
\newcommand{\jsvdaisy}{{\color{vedicmaroon}\SixFlowerRemovedOpenPetal}}

\newsavebox{\jsvpageborderbox}
\AtBeginDocument{%
  \sbox{\jsvpageborderbox}{%
    \begin{tikzpicture}
      \useasboundingbox (0,0) rectangle (\paperwidth, \paperheight);
      
      \coordinate (BL) at (10mm, 12mm);
      \coordinate (BR) at (\dimexpr\paperwidth-10mm\relax, 12mm);
      \coordinate (TL) at (10mm, \dimexpr\paperheight-10mm\relax);
      \coordinate (TR) at (\dimexpr\paperwidth-10mm\relax, \dimexpr\paperheight-10mm\relax);

      \pgfmathsetmacro{\targetpitch}{4.0}
      \pgfmathsetmacro{\totalW}{\paperwidth/1mm - 20}
      \pgfmathsetmacro{\totalH}{\paperheight/1mm - 22}
      
      \pgfmathtruncatemacro{\N}{round(\totalW / \targetpitch)}
      \pgfmathtruncatemacro{\M}{round(\totalH / \targetpitch)}
      
      % Horizontal chains (bottom and top) from k=0 to N
      \foreach \k in {0,...,\N} {
        \path (BL) -- (BR) coordinate[pos=\k/\N] (B\k);
        \node[anchor=center, inner sep=0, outer sep=0] at (B\k) {\jsvdaisy};
        
        \path (TL) -- (TR) coordinate[pos=\k/\N] (T\k);
        \node[anchor=center, inner sep=0, outer sep=0] at (T\k) {\jsvdaisy};
      }
      
      % Vertical chains (left and right) from j=1 to M-1
      \pgfmathtruncatemacro{\MminusOne}{\M - 1}
      \foreach \j in {1,...,\MminusOne} {
        \path (BL) -- (TL) coordinate[pos=\j/\M] (L\j);
        \node[anchor=center, inner sep=0, outer sep=0] at (L\j) {\jsvdaisy};
        
        \path (BR) -- (TR) coordinate[pos=\j/\M] (R\j);
        \node[anchor=center, inner sep=0, outer sep=0] at (R\j) {\jsvdaisy};
      }
    \end{tikzpicture}%
  }%
}
\AddToShipoutPictureBG{%
  \AtPageLowerLeft{\usebox{\jsvpageborderbox}}%
}
% ------------------------------------

\begin{document}

\renewcommand{\contentsname}{॥ अनुक्रमणिका ॥}

\begin{titlepage}
    \thispagestyle{empty}
    \begin{singlespace}
    \begin{center}
        % SS1: Exact text from Cover page.docx, single line spacing, Harfbuzz
        \vspace*{-0.3cm}
        {\fontsize{8.5pt}{11pt}\selectfont
        ॥ श्री: ॥\\[0.15em]
        ॥ श्रीमते जैमिन्याचार्याय नमः ॥\ ॥ श्री तलवकारगुरवे नमः ॥\\[0.15em]
        ॥ श्री गुरुभ्यो नमः ॥\par}
        
        \vfill
        
        % SS2: Title & Subtitle from Cover page docx
        {\fontsize{24pt}{30pt}\selectfont \textbf{जैमिनीयसामवेद-संहिता}\par}
        \vskip 0.35em
        {\fontsize{15pt}{19pt}\selectfont \textbf{तलवकारशाखा-प्रकृति भाग:}\par}
        
        \vfill
        \vfill
        
        % SS3: Editors (Dual) from Cover page docx
        {\fontsize{15pt}{19pt}\selectfont
        \textbf{सम्पादकौ}\\[0.30em]
        {\fontsize{13pt}{17pt}\selectfont वेदमूर्तिः मकरभूषणम् अय्यङ्गार्}\\[0.18em]
        {\fontsize{13pt}{17pt}\selectfont ब्रह्मश्रीः शैलेष देशपाण्डे}\par}
        
        \vfill
        \vfill
        
        % SS4: Publisher & Date from Cover page docx
        {\fontsize{12pt}{15pt}\selectfont
        \textbf{प्रकाशकः}\\[0.25em]
        {\fontsize{11pt}{14pt}\selectfont आर्ष विद्या वेदपाठशाला, कोडुन्तिरपुळ्ळि}\\[0.18em]
        {\fontsize{11pt}{14pt}\selectfont पालक्काड्, केरळम्}\\[0.18em]
        {\fontsize{11pt}{14pt}\selectfont २०२६}\par}
        
        \vspace*{-0.2cm}
    \end{center}
    \end{singlespace}
\end{titlepage}

% --- TOC ---
{
  \RaggedRight 
  \tableofcontents
}
\newpage

% ==========================================
% JINJA2 LOOP STRUCTURE
% ==========================================


    % --- SUPERSECTION ---
    \clearpage           
    \phantomsection              
    \chapter*{\centering आग्नेयपाठः}
    \addcontentsline{toc}{chapter}{आग्नेयपाठः}
    \markboth{आग्नेयपाठः}{आग्नेयपाठः}
    
    % GLUE FIX: Prevent page break between Chapter and First Section
    \nopagebreak

        

            % --- SECTION ---
            
            \phantomsection  

            \section*{\centering प्रथम खण्डः}
            \addcontentsline{toc}{section}{प्रथम खण्डः} 
            \markright{प्रथम खण्डः}
            \nopagebreak % Keep section title with first mantra
            
            % --- SECTION-LEVEL PROCEDURE LINK ---

                % --- SUBSECTIONS (Mantras) ---
                
                % Conditional rendering based on output_mode
                \par\filbreak

\phantomsection

\addcontentsline{toc}{subsection}{ऋक् १}

{\centering \footnotesize \textcolor{AccentPurple}{॥ भारद्वाजो बार्हस्पत्यः अग्निः गायत्री ॥} \par}

\vspace{0.12em}

{\centering \textbf{\textcolor{AccentBlue}{अग्ना\raisebox{0.25ex}{\accentmark{12}{\char"0951}}आ\raisebox{0.25ex}{\accentmark{15}{\char"1CD2}}याहीवी\raisebox{0.25ex}{\accentmark{15}{\char"1CD2}}ता\raisebox{0.25ex}{\accentmark{15}{\char"1CD2}}ये\kern 0.05em ।\hskip 0.12em plus 0.04em minus 0.02em गृ\raisebox{0.25ex}{\accentmark{15}{\char"1CD2}}णा\raisebox{0.25ex}{\accentmark{15}{\char"1CD2}}नोहव्यादा\raisebox{0.25ex}{\accentmark{12}{\char"0951}}तये\kern 0.05em ।\hskip 0.12em plus 0.04em minus 0.02em नीहो\raisebox{0.25ex}{\accentmark{15}{\char"1CD2}}ता\raisebox{0.25ex}{\accentmark{12}{\char"0951}}सत्सीबर्हीषि\raisebox{0.25ex}{\accentmark{12}{\char"0951}} \nolinebreak\hskip 0.08em \mbox{॥१॥}}} \par}

\vspace{0.8em}
                % --- SUBSECTIONS (Mantras) ---
                
                % Conditional rendering based on output_mode
                \par\filbreak

\phantomsection

\addcontentsline{toc}{subsection}{ऋक् २}

{\centering \footnotesize \textcolor{AccentPurple}{॥ भारद्वाजो बार्हस्पत्यः अग्निः गायत्री ॥} \par}

\vspace{0.12em}

{\centering \textbf{\textcolor{AccentBlue}{त्वा\raisebox{0.25ex}{\accentmark{12}{\char"0951}}माग्नेयज्ञा\raisebox{0.25ex}{\accentmark{15}{\char"1CD2}}नाम्\kern 0.05em ।\hskip 0.12em plus 0.04em minus 0.02em हो\raisebox{0.25ex}{\accentmark{12}{\char"0951}}ताविश्वे\raisebox{0.25ex}{\accentmark{12}{\char"0951}}षांहितः\kern 0.05em ।\hskip 0.12em plus 0.04em minus 0.02em देवेभि\raisebox{0.25ex}{\accentmark{12}{\char"0951}}र्मानूषे\raisebox{0.25ex}{\accentmark{12}{\char"0951}}जने \nolinebreak\hskip 0.08em \mbox{॥२॥}}} \par}

\vspace{0.8em}
                % --- SUBSECTIONS (Mantras) ---
                
                % Conditional rendering based on output_mode
                \par\filbreak

\phantomsection

\addcontentsline{toc}{subsection}{ऋक् ३}

{\centering \footnotesize \textcolor{AccentPurple}{॥ मेधातिथिः काण्वः अग्निः गायत्री ॥} \par}

\vspace{0.12em}

{\centering \textbf{\textcolor{AccentBlue}{अग्निंदूतं\raisebox{0.25ex}{\accentmark{15}{\char"1CD2}}वृ\raisebox{0.25ex}{\accentmark{12}{\char"0951}}णीमहे\kern 0.05em ।\hskip 0.12em plus 0.04em minus 0.02em हो\raisebox{0.25ex}{\accentmark{15}{\char"1CD2}}ता\raisebox{0.25ex}{\accentmark{12}{\char"0951}}रंविश्वा\raisebox{0.25ex}{\accentmark{15}{\char"1CD2}}वे\raisebox{0.25ex}{\accentmark{12}{\char"0951}}दसम्\kern 0.05em ।\hskip 0.12em plus 0.04em minus 0.02em अस्या\raisebox{0.15ex}{\accentmark{12}{\char"1CF9}}यज्ञास्या\raisebox{0.25ex}{\accentmark{12}{\char"0951}}सुक्रा\raisebox{0.25ex}{\accentmark{12}{\char"0951}}तुम् \nolinebreak\hskip 0.08em \mbox{॥३॥}}} \par}

\vspace{0.8em}
                % --- SUBSECTIONS (Mantras) ---
                
                % Conditional rendering based on output_mode
                \par\filbreak

\phantomsection

\addcontentsline{toc}{subsection}{ऋक् ४}

{\centering \footnotesize \textcolor{AccentPurple}{॥ भारद्वाजो बार्हस्पत्यः अग्निः गायत्री ॥} \par}

\vspace{0.12em}

{\centering \textbf{\textcolor{AccentBlue}{अग्निर्वृ\raisebox{0.25ex}{\accentmark{12}{\char"0951}}णि\raisebox{0.25ex}{\accentmark{12}{\char"0951}}जङ्घनात्\kern 0.05em ।\hskip 0.12em plus 0.04em minus 0.02em द्रविणास्यु\raisebox{0.15ex}{\accentmark{12}{\char"1CF9}}र्वीपन्या\raisebox{0.25ex}{\accentmark{15}{\char"1CD2}}या\kern 0.05em ।\hskip 0.12em plus 0.04em minus 0.02em सामी\raisebox{0.25ex}{\accentmark{15}{\char"1CD2}}द्धश्शुक्रा\raisebox{0.25ex}{\accentmark{15}{\char"1CD2}}आहूतः \nolinebreak\hskip 0.08em \mbox{॥४॥}}} \par}

\vspace{0.8em}
                % --- SUBSECTIONS (Mantras) ---
                
                % Conditional rendering based on output_mode
                \par\filbreak

\phantomsection

\addcontentsline{toc}{subsection}{ऋक् ५}

{\centering \footnotesize \textcolor{AccentPurple}{॥ उशनाः अग्निः गायत्री ॥} \par}

\vspace{0.12em}

{\centering \textbf{\textcolor{AccentBlue}{प्रेष्ठं\raisebox{0.15ex}{\accentmark{12}{\char"1CF8}}वोआतीथिम्\kern 0.05em ।\hskip 0.12em plus 0.04em minus 0.02em स्तू\raisebox{0.25ex}{\accentmark{15}{\char"1CD2}}षे\raisebox{0.25ex}{\accentmark{15}{\char"1CD2}}मित्रामीवाप्रीयम्\kern 0.05em ।\hskip 0.12em plus 0.04em minus 0.02em अ\raisebox{0.25ex}{\accentmark{15}{\char"1CD2}}ग्नेराथन्ना\raisebox{0.25ex}{\accentmark{15}{\char"1CD2}}वेद्यम् \nolinebreak\hskip 0.08em \mbox{॥५॥}}} \par}

\vspace{0.8em}
                % --- SUBSECTIONS (Mantras) ---
                
                % Conditional rendering based on output_mode
                \par\filbreak

\phantomsection

\addcontentsline{toc}{subsection}{ऋक् ६}

{\centering \footnotesize \textcolor{AccentPurple}{॥ सुदीतिपुरुमीढा वांगिरसौ तयोर्वाऽन्यतरः अग्निः गायत्री ॥} \par}

\vspace{0.12em}

{\centering \textbf{\textcolor{AccentBlue}{त्वन्नोआग्नेमा\raisebox{0.25ex}{\accentmark{15}{\char"1CD2}}हो\raisebox{0.25ex}{\accentmark{12}{\char"0951}}भिः\kern 0.05em ।\hskip 0.12em plus 0.04em minus 0.02em पाहीविश्वा\raisebox{0.25ex}{\accentmark{12}{\char"0951}}स्याआ\raisebox{0.25ex}{\accentmark{15}{\char"1CD2}}रा\raisebox{0.25ex}{\accentmark{12}{\char"0951}}तेः\kern 0.05em ।\hskip 0.12em plus 0.04em minus 0.02em ऊत\raisebox{0.25ex}{\accentmark{15}{\char"1CD2}}द्वीषोमर्क्त्या\raisebox{0.25ex}{\accentmark{12}{\char"0951}}स्या \nolinebreak\hskip 0.08em \mbox{॥६॥}}} \par}

\vspace{0.8em}
                % --- SUBSECTIONS (Mantras) ---
                
                % Conditional rendering based on output_mode
                \par\filbreak

\phantomsection

\addcontentsline{toc}{subsection}{ऋक् ७}

{\centering \footnotesize \textcolor{AccentPurple}{॥ भारद्वाजो बार्हस्पत्यः अग्निः गायत्री ॥} \par}

\vspace{0.12em}

{\centering \textbf{\textcolor{AccentBlue}{एह्यू\raisebox{0.25ex}{\accentmark{15}{\char"1CD2}}षुब्रा\raisebox{0.25ex}{\accentmark{15}{\char"1CD2}}वा\raisebox{0.25ex}{\accentmark{12}{\char"0951}}णिते\kern 0.05em ।\hskip 0.12em plus 0.04em minus 0.02em अग्नाइत्थे\raisebox{0.15ex}{\accentmark{12}{\char"1CF9}}ता\raisebox{0.25ex}{\accentmark{12}{\char"0951}}रा\raisebox{0.25ex}{\accentmark{12}{\char"0951}}गीराः\kern 0.05em ।\hskip 0.12em plus 0.04em minus 0.02em ए\raisebox{0.25ex}{\accentmark{15}{\char"1CD2}}भिर्वार्धा\raisebox{0.25ex}{\accentmark{12}{\char"0951}}साइ\raisebox{0.25ex}{\accentmark{15}{\char"1CD2}}न्दू\raisebox{0.25ex}{\accentmark{12}{\char"0951}}भिः \nolinebreak\hskip 0.08em \mbox{॥७॥}}} \par}

\vspace{0.8em}
                % --- SUBSECTIONS (Mantras) ---
                
                % Conditional rendering based on output_mode
                \par\filbreak

\phantomsection

\addcontentsline{toc}{subsection}{ऋक् ८}

{\centering \footnotesize \textcolor{AccentPurple}{॥ वत्सः काण्वः अग्निः गायत्री ॥} \par}

\vspace{0.12em}

{\centering \textbf{\textcolor{AccentBlue}{आ\raisebox{0.15ex}{\accentmark{12}{\char"1CF8}}तेवात्सो\raisebox{0.25ex}{\accentmark{15}{\char"1CD2}}मानोयमात्\kern 0.05em ।\hskip 0.12em plus 0.04em minus 0.02em परामा\raisebox{0.15ex}{\accentmark{12}{\char"1CF9}}ञ्चित्साधास्थात्\kern 0.05em ।\hskip 0.12em plus 0.04em minus 0.02em अ\raisebox{0.25ex}{\accentmark{15}{\char"1CD2}}ग्ने\raisebox{0.25ex}{\accentmark{15}{\char"1CD2}}त्वाङ्कामयोगी\raisebox{0.25ex}{\accentmark{15}{\char"1CD2}}रा\raisebox{0.25ex}{\accentmark{12}{\char"0951}} \nolinebreak\hskip 0.08em \mbox{॥८॥}}} \par}

\vspace{0.8em}
                % --- SUBSECTIONS (Mantras) ---
                
                % Conditional rendering based on output_mode
                \par\filbreak

\phantomsection

\addcontentsline{toc}{subsection}{ऋक् ९}

{\centering \footnotesize \textcolor{AccentPurple}{॥ भारद्वाजो बार्हस्पत्यः अग्निः गायत्री ॥} \par}

\vspace{0.12em}

{\centering \textbf{\textcolor{AccentBlue}{त्वामाग्नेपु\raisebox{0.15ex}{\accentmark{12}{\char"1CF9}}ष्काराद\raisebox{0.25ex}{\accentmark{15}{\char"1CD2}}धी\raisebox{0.25ex}{\accentmark{12}{\char"0951}}\kern 0.05em ।\hskip 0.12em plus 0.04em minus 0.02em आ\raisebox{0.15ex}{\accentmark{12}{\char"1CF8}}थ\raisebox{0.25ex}{\accentmark{12}{\char"0951}}र्वाणीरा\raisebox{0.25ex}{\accentmark{12}{\char"0951}}मन्धत\kern 0.05em ।\hskip 0.12em plus 0.04em minus 0.02em मूर्ध्नो\raisebox{0.25ex}{\accentmark{15}{\char"1CD2}}विश्वास्यावाघा\raisebox{0.25ex}{\accentmark{15}{\char"1CD2}}ता\raisebox{0.25ex}{\accentmark{12}{\char"0951}}ः \nolinebreak\hskip 0.08em \mbox{॥९॥}}} \par}

\vspace{0.8em}
                % --- SUBSECTIONS (Mantras) ---
                
                % Conditional rendering based on output_mode
                \par\filbreak

\phantomsection

\addcontentsline{toc}{subsection}{ऋक् १०}

{\centering \footnotesize \textcolor{AccentPurple}{॥ वामदेवः अग्निः गायत्री ॥} \par}

\vspace{0.12em}

{\centering \textbf{\textcolor{AccentBlue}{अग्नेवी\raisebox{0.15ex}{\accentmark{12}{\char"1CF8}}वा\raisebox{0.25ex}{\accentmark{15}{\char"1CD2}}स्वादा\raisebox{0.25ex}{\accentmark{15}{\char"1CD2}}भारा\kern 0.05em ।\hskip 0.12em plus 0.04em minus 0.02em अस्मभ्यामू\raisebox{0.25ex}{\accentmark{15}{\char"1CD2}}ता\raisebox{0.15ex}{\accentmark{12}{\char"1CF9}}ये\raisebox{0.25ex}{\accentmark{12}{\char"0951}}माहे\raisebox{0.25ex}{\accentmark{15}{\char"1CD2}}\kern 0.05em ।\hskip 0.12em plus 0.04em minus 0.02em देवोह्यासी\raisebox{0.25ex}{\accentmark{12}{\char"0951}}नोदृशे \nolinebreak\hskip 0.08em \mbox{॥१०॥}}} \par}

\vspace{0.8em}

            % --- End of Kandah Closing Floral Line & Counts ---
            \par\nopagebreak\vspace{0.9em plus 0.3em minus 0.1em}
            \centerline{\textcolor{black!35}{\rule{1.8cm}{0.35pt}}\;\raisebox{-0.22ex}{\textcolor{vedicmaroon}{\large\SixFlowerPetalDotted}}\;\textcolor{black!35}{\rule{1.8cm}{0.35pt}}}
            \nopagebreak\vspace{0.45em}
            \centerline{\devafont\textbf{\textcolor{AccentBlue}{॥ ऋक् सङ्ख्या : १० ॥}}}
        
            \clearpage

            % --- SECTION ---
            
            \phantomsection  

            \section*{\centering द्वितीय खण्डः}
            \addcontentsline{toc}{section}{द्वितीय खण्डः} 
            \markright{द्वितीय खण्डः}
            \nopagebreak % Keep section title with first mantra
            
            % --- SECTION-LEVEL PROCEDURE LINK ---

                % --- SUBSECTIONS (Mantras) ---
                
                % Conditional rendering based on output_mode
                \par\filbreak

\phantomsection

\addcontentsline{toc}{subsection}{ऋक् १}

{\centering \footnotesize \textcolor{AccentPurple}{॥ आयुन्ङ्क्श्वाहिः (ऋ विरूपांगिरसः) अग्निः गायत्री ॥} \par}

\vspace{0.12em}

{\centering \textbf{\textcolor{AccentBlue}{नामास्ते\raisebox{0.25ex}{\accentmark{15}{\char"1CD2}}अग्ना\raisebox{0.25ex}{\accentmark{15}{\char"1CD2}}ओज\raisebox{0.25ex}{\accentmark{12}{\char"0951}}से\kern 0.05em ।\hskip 0.12em plus 0.04em minus 0.02em गृणन्तीदे\raisebox{0.25ex}{\accentmark{15}{\char"1CD2}}वा\raisebox{0.25ex}{\accentmark{12}{\char"0951}}कृष्टायाः\kern 0.05em ।\hskip 0.12em plus 0.04em minus 0.02em आमै\raisebox{0.25ex}{\accentmark{15}{\char"1CD2}}रामित्रा\raisebox{0.25ex}{\accentmark{15}{\char"1CD2}}मर्दया \nolinebreak\hskip 0.08em \mbox{॥१॥}}} \par}

\vspace{0.8em}
                % --- SUBSECTIONS (Mantras) ---
                
                % Conditional rendering based on output_mode
                \par\filbreak

\phantomsection

\addcontentsline{toc}{subsection}{ऋक् २}

{\centering \footnotesize \textcolor{AccentPurple}{॥ वामदेवो गौतमः अग्निः गायत्री ॥} \par}

\vspace{0.12em}

{\centering \textbf{\textcolor{AccentBlue}{दूतंवो\raisebox{0.25ex}{\accentmark{12}{\char"0951}}विश्वावे\raisebox{0.25ex}{\accentmark{12}{\char"0951}}दसम्\kern 0.05em ।\hskip 0.12em plus 0.04em minus 0.02em हव्यावाहामामा\raisebox{0.25ex}{\accentmark{15}{\char"1CD2}}र्क्त्यम्\raisebox{0.25ex}{\accentmark{12}{\char"0951}}\kern 0.05em ।\hskip 0.12em plus 0.04em minus 0.02em याजीष्ठमृन्जसे\raisebox{0.25ex}{\accentmark{15}{\char"1CD2}}गीरा \nolinebreak\hskip 0.08em \mbox{॥२॥}}} \par}

\vspace{0.8em}
                % --- SUBSECTIONS (Mantras) ---
                
                % Conditional rendering based on output_mode
                \par\filbreak

\phantomsection

\addcontentsline{toc}{subsection}{ऋक् ३}

{\centering \footnotesize \textcolor{AccentPurple}{॥ प्रयोगो भागर्वः अग्निः गायत्री ॥} \par}

\vspace{0.12em}

{\centering \textbf{\textcolor{AccentBlue}{ऊपात्वाजामायो\raisebox{0.25ex}{\accentmark{15}{\char"1CD2}}गीरा\raisebox{0.25ex}{\accentmark{12}{\char"0951}}ः\kern 0.05em ।\hskip 0.12em plus 0.04em minus 0.02em देदीशतिर्हा\raisebox{0.25ex}{\accentmark{12}{\char"0951}}विष्कृताः\kern 0.05em ।\hskip 0.12em plus 0.04em minus 0.02em वायोरा\raisebox{0.25ex}{\accentmark{12}{\char"0951}}नीकेअस्थिरन् \nolinebreak\hskip 0.08em \mbox{॥३॥}}} \par}

\vspace{0.8em}
                % --- SUBSECTIONS (Mantras) ---
                
                % Conditional rendering based on output_mode
                \par\filbreak

\phantomsection

\addcontentsline{toc}{subsection}{ऋक् ४}

{\centering \footnotesize \textcolor{AccentPurple}{॥ मधुच्छन्दा वैश्वामित्रः अग्निः गायत्री ॥} \par}

\vspace{0.12em}

{\centering \textbf{\textcolor{AccentBlue}{ऊ\raisebox{0.25ex}{\accentmark{12}{\char"0951}}पात्वाग्नेदीवे\raisebox{0.25ex}{\accentmark{15}{\char"1CD2}}दीवे\kern 0.05em ।\hskip 0.12em plus 0.04em minus 0.02em दो\raisebox{0.25ex}{\accentmark{15}{\char"1CD2}}षा\raisebox{0.25ex}{\accentmark{12}{\char"0951}}वस्तार्धीयावायम्\raisebox{0.25ex}{\accentmark{15}{\char"1CD2}}\kern 0.05em ।\hskip 0.12em plus 0.04em minus 0.02em ना\raisebox{0.25ex}{\accentmark{15}{\char"1CD2}}मोभा\raisebox{0.15ex}{\accentmark{12}{\char"1CF9}}र\raisebox{0.25ex}{\accentmark{12}{\char"0951}}न्ताएमा\raisebox{0.25ex}{\accentmark{12}{\char"0951}}सि \nolinebreak\hskip 0.08em \mbox{॥४॥}}} \par}

\vspace{0.8em}
                % --- SUBSECTIONS (Mantras) ---
                
                % Conditional rendering based on output_mode
                \par\filbreak

\phantomsection

\addcontentsline{toc}{subsection}{ऋक् ५}

{\centering \footnotesize \textcolor{AccentPurple}{॥ शुनः शेप आजीर्गितः अग्निः गायत्री ॥} \par}

\vspace{0.12em}

{\centering \textbf{\textcolor{AccentBlue}{जरा\raisebox{0.25ex}{\accentmark{12}{\char"0951}}बोधाताव्दीविड्ढी\kern 0.05em ।\hskip 0.12em plus 0.04em minus 0.02em वीशेवी\raisebox{0.25ex}{\accentmark{15}{\char"1CD2}}शेय\raisebox{0.25ex}{\accentmark{12}{\char"0951}}ज्ञीयाया\kern 0.05em ।\hskip 0.12em plus 0.04em minus 0.02em स्तो\raisebox{0.15ex}{\accentmark{12}{\char"1CF8}}मंरुद्रायादृशीक\raisebox{0.25ex}{\accentmark{15}{\char"1CD2}}म् \nolinebreak\hskip 0.08em \mbox{॥५॥}}} \par}

\vspace{0.8em}
                % --- SUBSECTIONS (Mantras) ---
                
                % Conditional rendering based on output_mode
                \par\filbreak

\phantomsection

\addcontentsline{toc}{subsection}{ऋक् ६}

{\centering \footnotesize \textcolor{AccentPurple}{॥ मेधातिथिः अग्निः गायत्री ॥} \par}

\vspace{0.12em}

{\centering \textbf{\textcolor{AccentBlue}{प्रातित्यञ्चा\raisebox{0.25ex}{\accentmark{15}{\char"1CD2}}रुमाध्वारम्\raisebox{0.25ex}{\accentmark{15}{\char"1CD2}}\kern 0.05em ।\hskip 0.12em plus 0.04em minus 0.02em गो\raisebox{0.15ex}{\accentmark{12}{\char"1CF8}}पी\raisebox{0.25ex}{\accentmark{12}{\char"0951}}था\raisebox{0.25ex}{\accentmark{15}{\char"1CD2}}यप्रा\raisebox{0.25ex}{\accentmark{15}{\char"1CD2}}हू\raisebox{0.25ex}{\accentmark{12}{\char"0951}}यसे\kern 0.05em ।\hskip 0.12em plus 0.04em minus 0.02em मा\raisebox{0.25ex}{\accentmark{15}{\char"1CD2}}रुद्भीरग्नाआग\raisebox{0.25ex}{\accentmark{12}{\char"0951}}हि \nolinebreak\hskip 0.08em \mbox{॥६॥}}} \par}

\vspace{0.8em}
                % --- SUBSECTIONS (Mantras) ---
                
                % Conditional rendering based on output_mode
                \par\filbreak

\phantomsection

\addcontentsline{toc}{subsection}{ऋक् ७}

{\centering \footnotesize \textcolor{AccentPurple}{॥ शुनः शेप आजीर्गितः अग्निः गायत्री ॥} \par}

\vspace{0.12em}

{\centering \textbf{\textcolor{AccentBlue}{अ\raisebox{0.25ex}{\accentmark{12}{\char"0951}}श्वन्न\raisebox{0.25ex}{\accentmark{15}{\char"1CD2}}त्वावारा\raisebox{0.25ex}{\accentmark{12}{\char"0951}}व\raisebox{0.25ex}{\accentmark{15}{\char"1CD2}}न्तम्\kern 0.05em ।\hskip 0.12em plus 0.04em minus 0.02em व\raisebox{0.15ex}{\accentmark{12}{\char"1CF8}}न्द\raisebox{0.15ex}{\accentmark{12}{\char"1CF9}}ध्याआग्नि\raisebox{0.25ex}{\accentmark{15}{\char"1CD2}}न्नामोभिः\kern 0.05em ।\hskip 0.12em plus 0.04em minus 0.02em स\raisebox{0.25ex}{\accentmark{15}{\char"1CD2}}म्म्रा\raisebox{0.25ex}{\accentmark{15}{\char"1CD2}}जन्तमाध्वाराणा\raisebox{0.25ex}{\accentmark{12}{\char"0951}}म् \nolinebreak\hskip 0.08em \mbox{॥७॥}}} \par}

\vspace{0.8em}
                % --- SUBSECTIONS (Mantras) ---
                
                % Conditional rendering based on output_mode
                \par\filbreak

\phantomsection

\addcontentsline{toc}{subsection}{ऋक् ८}

{\centering \footnotesize \textcolor{AccentPurple}{॥ प्रयोगो भागर्वः अग्निः गायत्री ॥} \par}

\vspace{0.12em}

{\centering \textbf{\textcolor{AccentBlue}{और्वा\raisebox{0.25ex}{\accentmark{15}{\char"1CD2}}भृगूवच्छूचिम्\kern 0.05em ।\hskip 0.12em plus 0.04em minus 0.02em आप्नावा\raisebox{0.25ex}{\accentmark{15}{\char"1CD2}}नावादाहु\raisebox{0.25ex}{\accentmark{12}{\char"0951}}वे\kern 0.05em ।\hskip 0.12em plus 0.04em minus 0.02em अग्निं\raisebox{0.15ex}{\accentmark{12}{\char"1CF9}}सामुद्रावाससम् \nolinebreak\hskip 0.08em \mbox{॥८॥}}} \par}

\vspace{0.8em}
                % --- SUBSECTIONS (Mantras) ---
                
                % Conditional rendering based on output_mode
                \par\filbreak

\phantomsection

\addcontentsline{toc}{subsection}{ऋक् ९}

{\centering \footnotesize \textcolor{AccentPurple}{॥ प्रयोगो भागर्वः अग्निः गायत्री ॥} \par}

\vspace{0.12em}

{\centering \textbf{\textcolor{AccentBlue}{अ\raisebox{0.25ex}{\accentmark{15}{\char"1CD2}}ग्नी\raisebox{0.25ex}{\accentmark{15}{\char"1CD2}}मिन्धानो\raisebox{0.25ex}{\accentmark{12}{\char"0951}}माना\raisebox{0.25ex}{\accentmark{12}{\char"0951}}सा\kern 0.05em ।\hskip 0.12em plus 0.04em minus 0.02em धी\raisebox{0.25ex}{\accentmark{12}{\char"0951}}यंसचेतामा\raisebox{0.25ex}{\accentmark{15}{\char"1CD2}}र्त्या\raisebox{0.25ex}{\accentmark{12}{\char"0951}}ः\kern 0.05em ।\hskip 0.12em plus 0.04em minus 0.02em अ\raisebox{0.25ex}{\accentmark{15}{\char"1CD2}}ग्नी\raisebox{0.25ex}{\accentmark{15}{\char"1CD2}}मिन्धेवीवा\raisebox{0.25ex}{\accentmark{15}{\char"1CD2}}स्वा\raisebox{0.25ex}{\accentmark{12}{\char"0951}}भिः \nolinebreak\hskip 0.08em \mbox{॥९॥}}} \par}

\vspace{0.8em}
                % --- SUBSECTIONS (Mantras) ---
                
                % Conditional rendering based on output_mode
                \par\filbreak

\phantomsection

\addcontentsline{toc}{subsection}{ऋक् १०}

{\centering \footnotesize \textcolor{AccentPurple}{॥ वत्सः काण्वः अग्निः गायत्री ॥} \par}

\vspace{0.12em}

{\centering \textbf{\textcolor{AccentBlue}{आदी\raisebox{0.25ex}{\accentmark{15}{\char"1CD2}}त्प्रत्ना\raisebox{0.25ex}{\accentmark{15}{\char"1CD2}}स्यारे\raisebox{0.25ex}{\accentmark{15}{\char"1CD2}}ता\raisebox{0.25ex}{\accentmark{12}{\char"0951}}सः\kern 0.05em ।\hskip 0.12em plus 0.04em minus 0.02em ज्योतिःपश्यन्तीवासा\raisebox{0.25ex}{\accentmark{15}{\char"1CD2}}रम्\kern 0.05em ।\hskip 0.12em plus 0.04em minus 0.02em पा\raisebox{0.25ex}{\accentmark{15}{\char"1CD2}}रोयातिध्या\raisebox{0.15ex}{\accentmark{12}{\char"1CF9}}तेदीवि\raisebox{0.25ex}{\accentmark{15}{\char"1CD2}} \nolinebreak\hskip 0.08em \mbox{॥१०॥}}} \par}

\vspace{0.8em}

            % --- End of Kandah Closing Floral Line & Counts ---
            \par\nopagebreak\vspace{0.9em plus 0.3em minus 0.1em}
            \centerline{\textcolor{black!35}{\rule{1.8cm}{0.35pt}}\;\raisebox{-0.22ex}{\textcolor{vedicmaroon}{\large\SixFlowerPetalDotted}}\;\textcolor{black!35}{\rule{1.8cm}{0.35pt}}}
            \nopagebreak\vspace{0.45em}
            \centerline{\devafont\textbf{\textcolor{AccentBlue}{॥ ऋक् सङ्ख्या : १० ॥}}}
        
            \clearpage

            % --- SECTION ---
            
            \phantomsection  

            \section*{\centering तृतीय खण्डः}
            \addcontentsline{toc}{section}{तृतीय खण्डः} 
            \markright{तृतीय खण्डः}
            \nopagebreak % Keep section title with first mantra
            
            % --- SECTION-LEVEL PROCEDURE LINK ---

                % --- SUBSECTIONS (Mantras) ---
                
                % Conditional rendering based on output_mode
                \par\filbreak

\phantomsection

\addcontentsline{toc}{subsection}{ऋक् १}

{\centering \footnotesize \textcolor{AccentPurple}{॥ प्रयोगो भागर्वः अग्निः गायत्री ॥} \par}

\vspace{0.12em}

{\centering \textbf{\textcolor{AccentBlue}{अग्निंवोवृधन्तम्\kern 0.05em ।\hskip 0.12em plus 0.04em minus 0.02em अध्वाराणांपुरूतामम्\kern 0.05em ।\hskip 0.12em plus 0.04em minus 0.02em अच्छानाप्नेसाहा\raisebox{0.25ex}{\accentmark{15}{\char"1CD2}}स्वते \nolinebreak\hskip 0.08em \mbox{॥१॥}}} \par}

\vspace{0.8em}
                % --- SUBSECTIONS (Mantras) ---
                
                % Conditional rendering based on output_mode
                \par\filbreak

\phantomsection

\addcontentsline{toc}{subsection}{ऋक् २}

{\centering \footnotesize \textcolor{AccentPurple}{॥ भाद्वाजो बार्हस्पत्यः अग्निः गायत्री ॥} \par}

\vspace{0.12em}

{\centering \textbf{\textcolor{AccentBlue}{अग्निस्तिग्मे\raisebox{0.15ex}{\accentmark{12}{\char"1CF9}}ना\raisebox{0.25ex}{\accentmark{12}{\char"0951}}शोचीषा\raisebox{0.25ex}{\accentmark{12}{\char"0951}}\kern 0.05em ।\hskip 0.12em plus 0.04em minus 0.02em यंसद्वि\raisebox{0.25ex}{\accentmark{15}{\char"1CD2}}श्वंन्याटट्यत्रीणम्\kern 0.05em ।\hskip 0.12em plus 0.04em minus 0.02em अ\raisebox{0.25ex}{\accentmark{15}{\char"1CD2}}ग्नि\raisebox{0.25ex}{\accentmark{15}{\char"1CD2}}र्नो\raisebox{0.25ex}{\accentmark{12}{\char"0951}}वंसतेरायिम्\raisebox{0.15ex}{\accentmark{12}{\char"1CF8}} \nolinebreak\hskip 0.08em \mbox{॥२॥}}} \par}

\vspace{0.8em}
                % --- SUBSECTIONS (Mantras) ---
                
                % Conditional rendering based on output_mode
                \par\filbreak

\phantomsection

\addcontentsline{toc}{subsection}{ऋक् ३}

{\centering \footnotesize \textcolor{AccentPurple}{॥ वामदेवो गौतमः अग्निः गायत्री ॥} \par}

\vspace{0.12em}

{\centering \textbf{\textcolor{AccentBlue}{अग्नेमृ\raisebox{0.15ex}{\accentmark{12}{\char"1CF9}}ळामाहंआ\raisebox{0.25ex}{\accentmark{12}{\char"0951}}सि\raisebox{0.25ex}{\accentmark{12}{\char"0951}}\kern 0.05em ।\hskip 0.12em plus 0.04em minus 0.02em आ\raisebox{0.25ex}{\accentmark{15}{\char"1CD2}}याआ\raisebox{0.15ex}{\accentmark{12}{\char"1CF9}}दे\raisebox{0.25ex}{\accentmark{12}{\char"0951}}वायुन्ज\raisebox{0.25ex}{\accentmark{12}{\char"0951}}नम्\kern 0.05em ।\hskip 0.12em plus 0.04em minus 0.02em ईये\raisebox{0.15ex}{\accentmark{12}{\char"1CF8}}थाबर्ही\raisebox{0.25ex}{\accentmark{12}{\char"0951}}रासादम् \nolinebreak\hskip 0.08em \mbox{॥३॥}}} \par}

\vspace{0.8em}
                % --- SUBSECTIONS (Mantras) ---
                
                % Conditional rendering based on output_mode
                \par\filbreak

\phantomsection

\addcontentsline{toc}{subsection}{ऋक् ४}

{\centering \textbf{\textcolor{AccentBlue}{॥ वसिष्ठोमैत्रावरुणिःअग्निःगायत्री ॥ ॥ भाद्वाजोबार्हस्पत्यःअग्निःगायत्री ॥ अ\raisebox{0.25ex}{\accentmark{15}{\char"1CD2}}ग्नेर\raisebox{0.15ex}{\accentmark{12}{\char"1CF9}}क्षाणोअं\raisebox{0.25ex}{\accentmark{15}{\char"1CD2}}हा\raisebox{0.25ex}{\accentmark{12}{\char"0951}}सः\kern 0.05em ।\hskip 0.12em plus 0.04em minus 0.02em प्रा\raisebox{0.25ex}{\accentmark{15}{\char"1CD2}}तीस्मदेवरिषाताः\kern 0.05em ।\hskip 0.12em plus 0.04em minus 0.02em ता\raisebox{0.25ex}{\accentmark{15}{\char"1CD2}}पीष्ठैराज\raisebox{0.25ex}{\accentmark{15}{\char"1CD2}}रोदहा \nolinebreak\hskip 0.08em \mbox{॥४॥}}} \par}

\vspace{0.8em}
                % --- SUBSECTIONS (Mantras) ---
                
                % Conditional rendering based on output_mode
                \par\filbreak

\phantomsection

\addcontentsline{toc}{subsection}{ऋक् ५}

{\centering \textbf{\textcolor{AccentBlue}{अ\raisebox{0.15ex}{\accentmark{12}{\char"1CF8}}ग्नेयुंक्ष्वाहीयेता\raisebox{0.25ex}{\accentmark{15}{\char"1CD2}}वा\raisebox{0.25ex}{\accentmark{12}{\char"0951}}\kern 0.05em ।\hskip 0.12em plus 0.04em minus 0.02em अश्वा\raisebox{0.25ex}{\accentmark{12}{\char"0951}}सोदेवासाधावाः\kern 0.05em ।\hskip 0.12em plus 0.04em minus 0.02em आ\raisebox{0.25ex}{\accentmark{15}{\char"1CD2}}रं\raisebox{0.25ex}{\accentmark{12}{\char"0951}}वा\raisebox{0.15ex}{\accentmark{12}{\char"1CF9}}ह\raisebox{0.25ex}{\accentmark{12}{\char"0951}}न्त्याशा\raisebox{0.25ex}{\accentmark{15}{\char"1CD2}}वा\raisebox{0.25ex}{\accentmark{12}{\char"0951}}ः \nolinebreak\hskip 0.08em \mbox{॥५॥}}} \par}

\vspace{0.8em}
                % --- SUBSECTIONS (Mantras) ---
                
                % Conditional rendering based on output_mode
                \par\filbreak

\phantomsection

\addcontentsline{toc}{subsection}{ऋक् ६}

{\centering \footnotesize \textcolor{AccentPurple}{॥ वसिष्ठो मैत्रावरुणिः अग्निः गायत्री ॥} \par}

\vspace{0.12em}

{\centering \textbf{\textcolor{AccentBlue}{नित्वा\raisebox{0.25ex}{\accentmark{12}{\char"0951}}नक्ष्यविस्पते\kern 0.05em ।\hskip 0.12em plus 0.04em minus 0.02em धू\raisebox{0.25ex}{\accentmark{15}{\char"1CD2}}म\raisebox{0.25ex}{\accentmark{12}{\char"0951}}न्त\raisebox{0.25ex}{\accentmark{15}{\char"1CD2}}न्धीमहेवायम्\kern 0.05em ।\hskip 0.12em plus 0.04em minus 0.02em सु\raisebox{0.25ex}{\accentmark{15}{\char"1CD2}}वी\raisebox{0.25ex}{\accentmark{15}{\char"1CD2}}रामग्नआ\raisebox{0.15ex}{\accentmark{12}{\char"1CF8}}हुतः \nolinebreak\hskip 0.08em \mbox{॥६॥}}} \par}

\vspace{0.8em}
                % --- SUBSECTIONS (Mantras) ---
                
                % Conditional rendering based on output_mode
                \par\filbreak

\phantomsection

\addcontentsline{toc}{subsection}{ऋक् ७}

{\centering \footnotesize \textcolor{AccentPurple}{॥ विरूप आंगिरसः अग्निः गायत्री ॥} \par}

\vspace{0.12em}

{\centering \textbf{\textcolor{AccentBlue}{अ\raisebox{0.25ex}{\accentmark{15}{\char"1CD2}}ग्नि\raisebox{0.25ex}{\accentmark{15}{\char"1CD2}}र्मूर्द्धादीवाःकाकू\raisebox{0.25ex}{\accentmark{15}{\char"1CD2}}त्\kern 0.05em ।\hskip 0.12em plus 0.04em minus 0.02em पा\raisebox{0.25ex}{\accentmark{15}{\char"1CD2}}तीःपृथिव्या\raisebox{0.25ex}{\accentmark{15}{\char"1CD2}}आ\raisebox{0.25ex}{\accentmark{15}{\char"1CD2}}यम्\kern 0.05em ।\hskip 0.12em plus 0.04em minus 0.02em आपांरेतांसिजिन्वति \nolinebreak\hskip 0.08em \mbox{॥७॥}}} \par}

\vspace{0.8em}
                % --- SUBSECTIONS (Mantras) ---
                
                % Conditional rendering based on output_mode
                \par\filbreak

\phantomsection

\addcontentsline{toc}{subsection}{ऋक् ८}

{\centering \footnotesize \textcolor{AccentPurple}{॥ शुनः शेप आजीर्गितः अग्निः गायत्री ॥} \par}

\vspace{0.12em}

{\centering \textbf{\textcolor{AccentBlue}{ई\raisebox{0.25ex}{\accentmark{15}{\char"1CD2}}मामू\raisebox{0.25ex}{\accentmark{12}{\char"0951}}षुत्वाम\raisebox{0.25ex}{\accentmark{15}{\char"1CD2}}स्मा\raisebox{0.25ex}{\accentmark{12}{\char"0951}}कम्\kern 0.05em ।\hskip 0.12em plus 0.04em minus 0.02em सा\raisebox{0.25ex}{\accentmark{15}{\char"1CD2}}नि\raisebox{0.15ex}{\accentmark{12}{\char"1CF9}}ङ्गा\raisebox{0.25ex}{\accentmark{12}{\char"0951}}यात्र\raisebox{0.25ex}{\accentmark{15}{\char"1CD2}}न्नव्या\raisebox{0.25ex}{\accentmark{12}{\char"0951}}सम्\kern 0.05em ।\hskip 0.12em plus 0.04em minus 0.02em अग्ने\raisebox{0.25ex}{\accentmark{12}{\char"0951}}दे\raisebox{0.25ex}{\accentmark{15}{\char"1CD2}}वे\raisebox{0.25ex}{\accentmark{15}{\char"1CD2}}षुप्रा\raisebox{0.25ex}{\accentmark{15}{\char"1CD2}}वोचः \nolinebreak\hskip 0.08em \mbox{॥८॥}}} \par}

\vspace{0.8em}
                % --- SUBSECTIONS (Mantras) ---
                
                % Conditional rendering based on output_mode
                \par\filbreak

\phantomsection

\addcontentsline{toc}{subsection}{ऋक् ९}

{\centering \footnotesize \textcolor{AccentPurple}{॥ गोपवन अत्रयेः अग्निः गायत्री ॥} \par}

\vspace{0.12em}

{\centering \textbf{\textcolor{AccentBlue}{तन्त्वा\raisebox{0.25ex}{\accentmark{12}{\char"0951}}गोपा\raisebox{0.25ex}{\accentmark{15}{\char"1CD2}}वा\raisebox{0.25ex}{\accentmark{12}{\char"0951}}नोगिरा\kern 0.05em ।\hskip 0.12em plus 0.04em minus 0.02em जा\raisebox{0.25ex}{\accentmark{15}{\char"1CD2}}नि\raisebox{0.25ex}{\accentmark{12}{\char"0951}}ष्ठदग्नेअङ्गिरः\kern 0.05em ।\hskip 0.12em plus 0.04em minus 0.02em सा\raisebox{0.25ex}{\accentmark{15}{\char"1CD2}}पा\raisebox{0.25ex}{\accentmark{12}{\char"0951}}वकाश्रू\raisebox{0.25ex}{\accentmark{12}{\char"0951}}धीहा\raisebox{0.25ex}{\accentmark{15}{\char"1CD2}}व\raisebox{0.25ex}{\accentmark{12}{\char"0951}}म् \nolinebreak\hskip 0.08em \mbox{॥९॥}}} \par}

\vspace{0.8em}
                % --- SUBSECTIONS (Mantras) ---
                
                % Conditional rendering based on output_mode
                \par\filbreak

\phantomsection

\addcontentsline{toc}{subsection}{ऋक् १०}

{\centering \textbf{\textcolor{AccentBlue}{॥ वामदेवोगौतमःअग्निःगायत्री ॥ ॥ प्रस्कण्वःअग्निःगायत्री ॥ पा\raisebox{0.25ex}{\accentmark{15}{\char"1CD2}}रिवा\raisebox{0.25ex}{\accentmark{15}{\char"1CD2}}जपतीःकावि\raisebox{0.25ex}{\accentmark{15}{\char"1CD2}}ः\kern 0.05em ।\hskip 0.12em plus 0.04em minus 0.02em अ\raisebox{0.25ex}{\accentmark{15}{\char"1CD2}}ग्नि\raisebox{0.25ex}{\accentmark{15}{\char"1CD2}}ऋ\raisebox{0.25ex}{\accentmark{15}{\char"1CD2}}हव्या\raisebox{0.25ex}{\accentmark{15}{\char"1CD2}}न्या\raisebox{0.25ex}{\accentmark{12}{\char"0951}}क्रमीत्\kern 0.05em ।\hskip 0.12em plus 0.04em minus 0.02em दध\raisebox{0.25ex}{\accentmark{15}{\char"1CD2}}द्रत्ना\raisebox{0.25ex}{\accentmark{15}{\char"1CD2}}नीदाशू\raisebox{0.25ex}{\accentmark{15}{\char"1CD2}}षे\raisebox{0.25ex}{\accentmark{15}{\char"1CD2}} \nolinebreak\hskip 0.08em \mbox{॥१०॥}}} \par}

\vspace{0.8em}
                % --- SUBSECTIONS (Mantras) ---
                
                % Conditional rendering based on output_mode
                \par\filbreak

\phantomsection

\addcontentsline{toc}{subsection}{ऋक् ११}

{\centering \textbf{\textcolor{AccentBlue}{ऊदूत्यं\raisebox{0.25ex}{\accentmark{15}{\char"1CD2}}जाता\raisebox{0.25ex}{\accentmark{15}{\char"1CD2}}वे\raisebox{0.25ex}{\accentmark{12}{\char"0951}}दसम्\kern 0.05em ।\hskip 0.12em plus 0.04em minus 0.02em दे\raisebox{0.25ex}{\accentmark{15}{\char"1CD2}}वं\raisebox{0.25ex}{\accentmark{15}{\char"1CD2}}वाहन्तींकेता\raisebox{0.25ex}{\accentmark{15}{\char"1CD2}}वा\raisebox{0.25ex}{\accentmark{12}{\char"0951}}ः\kern 0.05em ।\hskip 0.12em plus 0.04em minus 0.02em दृशे\raisebox{0.15ex}{\accentmark{12}{\char"1CF9}}वि\raisebox{0.15ex}{\accentmark{12}{\char"1CF9}}श्वा\raisebox{0.25ex}{\accentmark{12}{\char"0951}}यासूर्यम् \nolinebreak\hskip 0.08em \mbox{॥११॥}}} \par}

\vspace{0.8em}
                % --- SUBSECTIONS (Mantras) ---
                
                % Conditional rendering based on output_mode
                \par\filbreak

\phantomsection

\addcontentsline{toc}{subsection}{ऋक् १२}

{\centering \footnotesize \textcolor{AccentPurple}{॥ मेधातिथिः काण्वः अग्निः गायत्री ॥} \par}

\vspace{0.12em}

{\centering \textbf{\textcolor{AccentBlue}{का\raisebox{0.25ex}{\accentmark{15}{\char"1CD2}}वी\raisebox{0.25ex}{\accentmark{15}{\char"1CD2}}माग्नी\raisebox{0.25ex}{\accentmark{15}{\char"1CD2}}मू\raisebox{0.25ex}{\accentmark{15}{\char"1CD2}}पा\raisebox{0.25ex}{\accentmark{12}{\char"0951}}स्तुहि\kern 0.05em ।\hskip 0.12em plus 0.04em minus 0.02em स\raisebox{0.25ex}{\accentmark{15}{\char"1CD2}}त्या\raisebox{0.25ex}{\accentmark{15}{\char"1CD2}}धर्मा\raisebox{0.25ex}{\accentmark{12}{\char"0951}}णमाध्वारे\kern 0.05em ।\hskip 0.12em plus 0.04em minus 0.02em दे\raisebox{0.25ex}{\accentmark{15}{\char"1CD2}}वा\raisebox{0.25ex}{\accentmark{15}{\char"1CD2}}मा\raisebox{0.25ex}{\accentmark{12}{\char"0951}}मीवाचा\raisebox{0.25ex}{\accentmark{15}{\char"1CD2}}ता\raisebox{0.25ex}{\accentmark{12}{\char"0951}}नम् \nolinebreak\hskip 0.08em \mbox{॥१२॥}}} \par}

\vspace{0.8em}
                % --- SUBSECTIONS (Mantras) ---
                
                % Conditional rendering based on output_mode
                \par\filbreak

\phantomsection

\addcontentsline{toc}{subsection}{ऋक् १३}

{\centering \footnotesize \textcolor{AccentPurple}{॥ सिन्धुद्वीप आम्बरीषः त्रित आत्यो वा अग्निः गायत्री ॥} \par}

\vspace{0.12em}

{\centering \textbf{\textcolor{AccentBlue}{श\raisebox{0.15ex}{\accentmark{12}{\char"1CF8}}न्नो\raisebox{0.25ex}{\accentmark{12}{\char"0951}}देवी\raisebox{0.25ex}{\accentmark{15}{\char"1CD2}}राभीष्टा\raisebox{0.25ex}{\accentmark{12}{\char"0951}}ये\kern 0.05em ।\hskip 0.12em plus 0.04em minus 0.02em श\raisebox{0.25ex}{\accentmark{15}{\char"1CD2}}न्नो\raisebox{0.25ex}{\accentmark{12}{\char"0951}}भूवन्तूपीता\raisebox{0.25ex}{\accentmark{15}{\char"1CD2}}ये\raisebox{0.25ex}{\accentmark{12}{\char"0951}}\kern 0.05em ।\hskip 0.12em plus 0.04em minus 0.02em श\raisebox{0.25ex}{\accentmark{15}{\char"1CD2}}य्योराभि\raisebox{0.25ex}{\accentmark{15}{\char"1CD2}}स्रावन्तूनः \nolinebreak\hskip 0.08em \mbox{॥१३॥}}} \par}

\vspace{0.8em}
                % --- SUBSECTIONS (Mantras) ---
                
                % Conditional rendering based on output_mode
                \par\filbreak

\phantomsection

\addcontentsline{toc}{subsection}{ऋक् १४}

{\centering \footnotesize \textcolor{AccentPurple}{॥ उशना काव्यः अग्निः गायत्री ॥} \par}

\vspace{0.12em}

{\centering \textbf{\textcolor{AccentBlue}{कस्यानू\raisebox{0.25ex}{\accentmark{15}{\char"1CD2}}नम्पारी\raisebox{0.25ex}{\accentmark{12}{\char"0951}}णसि\kern 0.05em ।\hskip 0.12em plus 0.04em minus 0.02em धीयो\raisebox{0.25ex}{\accentmark{12}{\char"0951}}जिन्वसिसत्पते\kern 0.05em ।\hskip 0.12em plus 0.04em minus 0.02em गो\raisebox{0.25ex}{\accentmark{15}{\char"1CD2}}षा\raisebox{0.25ex}{\accentmark{15}{\char"1CD2}}तायस्यातागी\raisebox{0.25ex}{\accentmark{15}{\char"1CD2}}रा\raisebox{0.25ex}{\accentmark{15}{\char"1CD2}}ः \nolinebreak\hskip 0.08em \mbox{॥१४॥}}} \par}

\vspace{0.8em}

            % --- End of Kandah Closing Floral Line & Counts ---
            \par\nopagebreak\vspace{0.9em plus 0.3em minus 0.1em}
            \centerline{\textcolor{black!35}{\rule{1.8cm}{0.35pt}}\;\raisebox{-0.22ex}{\textcolor{vedicmaroon}{\large\SixFlowerPetalDotted}}\;\textcolor{black!35}{\rule{1.8cm}{0.35pt}}}
            \nopagebreak\vspace{0.45em}
            \centerline{\devafont\textbf{\textcolor{AccentBlue}{॥ ऋक् सङ्ख्या : १४ ॥}}}
        
            \clearpage

            % --- SECTION ---
            
            \phantomsection  

            \section*{\centering चतुर्थ खण्डः}
            \addcontentsline{toc}{section}{चतुर्थ खण्डः} 
            \markright{चतुर्थ खण्डः}
            \nopagebreak % Keep section title with first mantra
            
            % --- SECTION-LEVEL PROCEDURE LINK ---

                % --- SUBSECTIONS (Mantras) ---
                
                % Conditional rendering based on output_mode
                \par\filbreak

\phantomsection

\addcontentsline{toc}{subsection}{ऋक् १}

{\centering \footnotesize \textcolor{AccentPurple}{॥ शय्युबार्हस्पत्यः अग्निः बृहती ॥} \par}

\vspace{0.12em}

{\centering \textbf{\textcolor{AccentBlue}{यज्ञा\raisebox{0.25ex}{\accentmark{15}{\char"1CD2}}याज्ञा\raisebox{0.25ex}{\accentmark{15}{\char"1CD2}}वोअग्ना\raisebox{0.25ex}{\accentmark{15}{\char"1CD2}}ये\raisebox{0.25ex}{\accentmark{12}{\char"0951}}\kern 0.05em ।\hskip 0.12em plus 0.04em minus 0.02em गी\raisebox{0.25ex}{\accentmark{15}{\char"1CD2}}रा\raisebox{0.25ex}{\accentmark{15}{\char"1CD2}}गी\raisebox{0.25ex}{\accentmark{12}{\char"0951}}राचाद\raisebox{0.25ex}{\accentmark{15}{\char"1CD2}}क्षासे\kern 0.05em ।\hskip 0.12em plus 0.04em minus 0.02em प्र\raisebox{0.15ex}{\accentmark{12}{\char"1CF8}}प्रा\raisebox{0.25ex}{\accentmark{15}{\char"1CD2}}वाया\raisebox{0.25ex}{\accentmark{15}{\char"1CD2}}मामृ\raisebox{0.15ex}{\accentmark{12}{\char"1CF9}}त\raisebox{0.25ex}{\accentmark{15}{\char"1CD2}}ञ्जा\raisebox{0.25ex}{\accentmark{12}{\char"0951}}तावेदसम्\kern 0.05em ।\hskip 0.12em plus 0.04em minus 0.02em प्रीयम्मित्र\raisebox{0.25ex}{\accentmark{15}{\char"1CD2}}न्ना\raisebox{0.25ex}{\accentmark{15}{\char"1CD2}}शंसिषम् \nolinebreak\hskip 0.08em \mbox{॥१॥}}} \par}

\vspace{0.8em}
                % --- SUBSECTIONS (Mantras) ---
                
                % Conditional rendering based on output_mode
                \par\filbreak

\phantomsection

\addcontentsline{toc}{subsection}{ऋक् २}

{\centering \footnotesize \textcolor{AccentPurple}{॥ भगः प्रागाथः अग्निः बृहती ॥} \par}

\vspace{0.12em}

{\centering \textbf{\textcolor{AccentBlue}{पाहिनोअग्नाएका\raisebox{0.25ex}{\accentmark{12}{\char"0951}}या\kern 0.05em ।\hskip 0.12em plus 0.04em minus 0.02em पा\raisebox{0.25ex}{\accentmark{15}{\char"1CD2}}ह्यूताद्वि\raisebox{0.25ex}{\accentmark{15}{\char"1CD2}}तीया\raisebox{0.25ex}{\accentmark{15}{\char"1CD2}}या\kern 0.05em ।\hskip 0.12em plus 0.04em minus 0.02em पा\raisebox{0.25ex}{\accentmark{15}{\char"1CD2}}हीगीर्भीस्तीसृ\raisebox{0.25ex}{\accentmark{15}{\char"1CD2}}भीरूर्जां\raisebox{0.25ex}{\accentmark{12}{\char"0951}}पते\kern 0.05em ।\hskip 0.12em plus 0.04em minus 0.02em पाही\raisebox{0.15ex}{\accentmark{12}{\char"1CF9}}चातासृभी\raisebox{0.25ex}{\accentmark{12}{\char"0951}}र्वसो \nolinebreak\hskip 0.08em \mbox{॥२॥}}} \par}

\vspace{0.8em}
                % --- SUBSECTIONS (Mantras) ---
                
                % Conditional rendering based on output_mode
                \par\filbreak

\phantomsection

\addcontentsline{toc}{subsection}{ऋक् ३}

{\centering \footnotesize \textcolor{AccentPurple}{॥ शय्युबार्हस्पत्यः अग्निः बृहती ॥} \par}

\vspace{0.12em}

{\centering \textbf{\textcolor{AccentBlue}{बृ\raisebox{0.25ex}{\accentmark{15}{\char"1CD2}}हा\raisebox{0.25ex}{\accentmark{15}{\char"1CD2}}त्भी\raisebox{0.25ex}{\accentmark{12}{\char"0951}}रग्नेरर्च्ची\raisebox{0.25ex}{\accentmark{12}{\char"0951}}भिः\kern 0.05em ।\hskip 0.12em plus 0.04em minus 0.02em शुक्रेणा\raisebox{0.25ex}{\accentmark{12}{\char"0951}}देवाशोची\raisebox{0.25ex}{\accentmark{15}{\char"1CD2}}षा\raisebox{0.25ex}{\accentmark{12}{\char"0951}}\kern 0.05em ।\hskip 0.12em plus 0.04em minus 0.02em भारा\raisebox{0.25ex}{\accentmark{15}{\char"1CD2}}द्वा\raisebox{0.25ex}{\accentmark{12}{\char"0951}}जेसमीधा\kern 0.05em ।\hskip 0.12em plus 0.04em minus 0.02em नो\raisebox{0.25ex}{\accentmark{15}{\char"1CD2}}या\raisebox{0.25ex}{\accentmark{12}{\char"0951}}विष्यारेवा\raisebox{0.25ex}{\accentmark{15}{\char"1CD2}}त्पा\raisebox{0.25ex}{\accentmark{12}{\char"0951}}वकदीदिहि \nolinebreak\hskip 0.08em \mbox{॥३॥}}} \par}

\vspace{0.8em}
                % --- SUBSECTIONS (Mantras) ---
                
                % Conditional rendering based on output_mode
                \par\filbreak

\phantomsection

\addcontentsline{toc}{subsection}{ऋक् ४}

{\centering \footnotesize \textcolor{AccentPurple}{॥ वसिष्ठो मैत्रावरुणिः अग्निः बृहती ॥} \par}

\vspace{0.12em}

{\centering \textbf{\textcolor{AccentBlue}{त्वय्या\raisebox{0.25ex}{\accentmark{12}{\char"0951}}ग्नेस्वा\raisebox{0.25ex}{\accentmark{15}{\char"1CD2}}हुतः\kern 0.05em ।\hskip 0.12em plus 0.04em minus 0.02em प्रीया\raisebox{0.25ex}{\accentmark{15}{\char"1CD2}}सा\raisebox{0.25ex}{\accentmark{15}{\char"1CD2}}सन्तूसूराया\raisebox{0.25ex}{\accentmark{12}{\char"0951}}ः\kern 0.05em ।\hskip 0.12em plus 0.04em minus 0.02em यन्ता\raisebox{0.25ex}{\accentmark{15}{\char"1CD2}}रोयेमाघा\raisebox{0.15ex}{\accentmark{12}{\char"1CF9}}वा\raisebox{0.25ex}{\accentmark{12}{\char"0951}}नोज\raisebox{0.25ex}{\accentmark{15}{\char"1CD2}}ना\raisebox{0.25ex}{\accentmark{12}{\char"0951}}नाम्\kern 0.05em ।\hskip 0.12em plus 0.04em minus 0.02em ऊर्वन्दय\raisebox{0.15ex}{\accentmark{12}{\char"1CF8}}न्तागो\raisebox{0.25ex}{\accentmark{15}{\char"1CD2}}नाम् \nolinebreak\hskip 0.08em \mbox{॥४॥}}} \par}

\vspace{0.8em}
                % --- SUBSECTIONS (Mantras) ---
                
                % Conditional rendering based on output_mode
                \par\filbreak

\phantomsection

\addcontentsline{toc}{subsection}{ऋक् ५}

{\centering \footnotesize \textcolor{AccentPurple}{॥ भगः प्रागाथः अग्निः बृहती ॥} \par}

\vspace{0.12em}

{\centering \textbf{\textcolor{AccentBlue}{अग्नेजरी\raisebox{0.25ex}{\accentmark{15}{\char"1CD2}}तार्वीस्पातिः\kern 0.05em ।\hskip 0.12em plus 0.04em minus 0.02em तपानो\raisebox{0.25ex}{\accentmark{15}{\char"1CD2}}देवारक्षा\raisebox{0.25ex}{\accentmark{15}{\char"1CD2}}साः\kern 0.05em ।\hskip 0.12em plus 0.04em minus 0.02em अ\raisebox{0.25ex}{\accentmark{15}{\char"1CD2}}प्रो\raisebox{0.25ex}{\accentmark{12}{\char"0951}}षिवान्ग्रहापा\raisebox{0.25ex}{\accentmark{12}{\char"0951}}तिर्माहंआ\raisebox{0.25ex}{\accentmark{12}{\char"0951}}सि\kern 0.05em ।\hskip 0.12em plus 0.04em minus 0.02em दी\raisebox{0.25ex}{\accentmark{12}{\char"0951}}वा\raisebox{0.25ex}{\accentmark{15}{\char"1CD2}}स्पायू\raisebox{0.25ex}{\accentmark{15}{\char"1CD2}}र्दुरोणायुः \nolinebreak\hskip 0.08em \mbox{॥५॥}}} \par}

\vspace{0.8em}
                % --- SUBSECTIONS (Mantras) ---
                
                % Conditional rendering based on output_mode
                \par\filbreak

\phantomsection

\addcontentsline{toc}{subsection}{ऋक् ६}

{\centering \footnotesize \textcolor{AccentPurple}{॥ प्रस्कण्वः काण्वः अग्निः बृहती ॥} \par}

\vspace{0.12em}

{\centering \textbf{\textcolor{AccentBlue}{अग्नेवी\raisebox{0.25ex}{\accentmark{15}{\char"1CD2}}वा\raisebox{0.25ex}{\accentmark{12}{\char"0951}}स्वादूषा\raisebox{0.25ex}{\accentmark{15}{\char"1CD2}}साः\kern 0.05em ।\hskip 0.12em plus 0.04em minus 0.02em चि\raisebox{0.25ex}{\accentmark{15}{\char"1CD2}}त्रं\raisebox{0.25ex}{\accentmark{15}{\char"1CD2}}रा\raisebox{0.25ex}{\accentmark{15}{\char"1CD2}}धोआ\raisebox{0.25ex}{\accentmark{15}{\char"1CD2}}मा\raisebox{0.25ex}{\accentmark{12}{\char"0951}}र्क्त्याम्\raisebox{0.25ex}{\accentmark{12}{\char"0951}}\kern 0.05em ।\hskip 0.12em plus 0.04em minus 0.02em आदाशू\raisebox{0.25ex}{\accentmark{15}{\char"1CD2}}षेजातवेदोवा\raisebox{0.25ex}{\accentmark{15}{\char"1CD2}}हात्वम्\raisebox{0.25ex}{\accentmark{15}{\char"1CD2}}\kern 0.05em ।\hskip 0.12em plus 0.04em minus 0.02em अद्या\raisebox{0.25ex}{\accentmark{15}{\char"1CD2}}देवं\raisebox{0.15ex}{\accentmark{12}{\char"1CF9}}ऊ\raisebox{0.25ex}{\accentmark{12}{\char"0951}}ष\raisebox{0.15ex}{\accentmark{12}{\char"1CF8}}र्बुधाः \nolinebreak\hskip 0.08em \mbox{॥६॥}}} \par}

\vspace{0.8em}
                % --- SUBSECTIONS (Mantras) ---
                
                % Conditional rendering based on output_mode
                \par\filbreak

\phantomsection

\addcontentsline{toc}{subsection}{ऋक् ७}

{\centering \footnotesize \textcolor{AccentPurple}{॥ शय्युबार्हस्पत्यः (तृणपाणिः) अग्निः बृहती ॥} \par}

\vspace{0.12em}

{\centering \textbf{\textcolor{AccentBlue}{त्व\raisebox{0.15ex}{\accentmark{12}{\char"1CF8}}न्नाश्चित्रावूत्य\kern 0.05em ।\hskip 0.12em plus 0.04em minus 0.02em वा\raisebox{0.25ex}{\accentmark{15}{\char"1CD2}}सो\raisebox{0.25ex}{\accentmark{15}{\char"1CD2}}राधांसिचो\raisebox{0.25ex}{\accentmark{15}{\char"1CD2}}दया\kern 0.05em ।\hskip 0.12em plus 0.04em minus 0.02em अ\raisebox{0.25ex}{\accentmark{15}{\char"1CD2}}स्या\raisebox{0.25ex}{\accentmark{15}{\char"1CD2}}रायेत्वा\raisebox{0.25ex}{\accentmark{15}{\char"1CD2}}मा\raisebox{0.25ex}{\accentmark{12}{\char"0951}}ग्नेरा\raisebox{0.25ex}{\accentmark{12}{\char"0951}}थीरासी\raisebox{0.25ex}{\accentmark{15}{\char"1CD2}}\kern 0.05em ।\hskip 0.12em plus 0.04em minus 0.02em वीदागा\raisebox{0.25ex}{\accentmark{15}{\char"1CD2}}धन्तूचेतूनाः \nolinebreak\hskip 0.08em \mbox{॥७॥}}} \par}

\vspace{0.8em}
                % --- SUBSECTIONS (Mantras) ---
                
                % Conditional rendering based on output_mode
                \par\filbreak

\phantomsection

\addcontentsline{toc}{subsection}{ऋक् ८}

{\centering \footnotesize \textcolor{AccentPurple}{॥ भगः प्रागाथः अग्निः बृहती ॥} \par}

\vspace{0.12em}

{\centering \textbf{\textcolor{AccentBlue}{त्वा\raisebox{0.25ex}{\accentmark{15}{\char"1CD2}}मीत्सप्रा\raisebox{0.25ex}{\accentmark{15}{\char"1CD2}}था\raisebox{0.25ex}{\accentmark{12}{\char"0951}}\raisebox{0.25ex}{\accentmark{15}{\char"1CD2}}असि\kern 0.05em ।\hskip 0.12em plus 0.04em minus 0.02em अग्नेत्राता\raisebox{0.25ex}{\accentmark{15}{\char"1CD2}}ऋताःकाविः\kern 0.05em ।\hskip 0.12em plus 0.04em minus 0.02em त्वांविप्रा\raisebox{0.25ex}{\accentmark{12}{\char"0951}}सस्समीधानदीदिवः\kern 0.05em ।\hskip 0.12em plus 0.04em minus 0.02em आ\raisebox{0.25ex}{\accentmark{15}{\char"1CD2}}वीवासन्तीवेधासाः \nolinebreak\hskip 0.08em \mbox{॥८॥}}} \par}

\vspace{0.8em}
                % --- SUBSECTIONS (Mantras) ---
                
                % Conditional rendering based on output_mode
                \par\filbreak

\phantomsection

\addcontentsline{toc}{subsection}{ऋक् ९}

{\centering \footnotesize \textcolor{AccentPurple}{॥ भगः प्रागाथः अग्निः बृहती ॥} \par}

\vspace{0.12em}

{\centering \textbf{\textcolor{AccentBlue}{आनो\raisebox{0.25ex}{\accentmark{12}{\char"0951}}आग्नेवायो\raisebox{0.25ex}{\accentmark{15}{\char"1CD2}}वृ\raisebox{0.25ex}{\accentmark{15}{\char"1CD2}}धम्\raisebox{0.25ex}{\accentmark{12}{\char"0951}}\kern 0.05em ।\hskip 0.12em plus 0.04em minus 0.02em रायिंपा\raisebox{0.25ex}{\accentmark{12}{\char"0951}}वाका\raisebox{0.25ex}{\accentmark{12}{\char"0951}}शंस्यम्\raisebox{0.25ex}{\accentmark{12}{\char"0951}}\kern 0.05em ।\hskip 0.12em plus 0.04em minus 0.02em रास्वा\raisebox{0.25ex}{\accentmark{12}{\char"0951}}चन\raisebox{0.25ex}{\accentmark{15}{\char"1CD2}}उपमातेपूरूस्पृहम्\raisebox{0.25ex}{\accentmark{12}{\char"0951}}\kern 0.05em ।\hskip 0.12em plus 0.04em minus 0.02em सू\raisebox{0.25ex}{\accentmark{12}{\char"0951}}नीं\raisebox{0.25ex}{\accentmark{15}{\char"1CD2}}तीसू\raisebox{0.25ex}{\accentmark{15}{\char"1CD2}}याशस्तरम् \nolinebreak\hskip 0.08em \mbox{॥९॥}}} \par}

\vspace{0.8em}
                % --- SUBSECTIONS (Mantras) ---
                
                % Conditional rendering based on output_mode
                \par\filbreak

\phantomsection

\addcontentsline{toc}{subsection}{ऋक् १०}

{\centering \footnotesize \textcolor{AccentPurple}{॥ सोभिरः काण्वः अग्निः बृहती ॥} \par}

\vspace{0.12em}

{\centering \textbf{\textcolor{AccentBlue}{यो\raisebox{0.25ex}{\accentmark{15}{\char"1CD2}}विश्वाद\raisebox{0.15ex}{\accentmark{12}{\char"1CF9}}यातेवासू\raisebox{0.25ex}{\accentmark{12}{\char"0951}}\kern 0.05em ।\hskip 0.12em plus 0.04em minus 0.02em हो\raisebox{0.15ex}{\accentmark{12}{\char"1CF8}}ता\raisebox{0.25ex}{\accentmark{12}{\char"0951}}मन्द्रोजना\raisebox{0.25ex}{\accentmark{12}{\char"0951}}नाम्\kern 0.05em ।\hskip 0.12em plus 0.04em minus 0.02em मा\raisebox{0.25ex}{\accentmark{15}{\char"1CD2}}धोर्नपात्राप्राथा\raisebox{0.25ex}{\accentmark{15}{\char"1CD2}}मान्यस्मै\kern 0.05em ।\hskip 0.12em plus 0.04em minus 0.02em प्रस्तो\raisebox{0.25ex}{\accentmark{15}{\char"1CD2}}मायन्त्वग्ना\raisebox{0.25ex}{\accentmark{15}{\char"1CD2}}ये \nolinebreak\hskip 0.08em \mbox{॥१०॥}}} \par}

\vspace{0.8em}

            % --- End of Kandah Closing Floral Line & Counts ---
            \par\nopagebreak\vspace{0.9em plus 0.3em minus 0.1em}
            \centerline{\textcolor{black!35}{\rule{1.8cm}{0.35pt}}\;\raisebox{-0.22ex}{\textcolor{vedicmaroon}{\large\SixFlowerPetalDotted}}\;\textcolor{black!35}{\rule{1.8cm}{0.35pt}}}
            \nopagebreak\vspace{0.45em}
            \centerline{\devafont\textbf{\textcolor{AccentBlue}{॥ ऋक् सङ्ख्या : १० ॥}}}
        
            \clearpage

            % --- SECTION ---
            
            \phantomsection  

            \section*{\centering पञ्चम खण्डः}
            \addcontentsline{toc}{section}{पञ्चम खण्डः} 
            \markright{पञ्चम खण्डः}
            \nopagebreak % Keep section title with first mantra
            
            % --- SECTION-LEVEL PROCEDURE LINK ---

                % --- SUBSECTIONS (Mantras) ---
                
                % Conditional rendering based on output_mode
                \par\filbreak

\phantomsection

\addcontentsline{toc}{subsection}{ऋक् १}

{\centering \footnotesize \textcolor{AccentPurple}{॥ वसिष्ठो मैत्रावरुणिः अग्निः बृहती ॥} \par}

\vspace{0.12em}

{\centering \textbf{\textcolor{AccentBlue}{ए\raisebox{0.25ex}{\accentmark{15}{\char"1CD2}}नावो\raisebox{0.25ex}{\accentmark{15}{\char"1CD2}}अग्नि\raisebox{0.25ex}{\accentmark{15}{\char"1CD2}}न्ना\raisebox{0.25ex}{\accentmark{15}{\char"1CD2}}मा\raisebox{0.25ex}{\accentmark{12}{\char"0951}}सा\kern 0.05em ।\hskip 0.12em plus 0.04em minus 0.02em ऊर्जो\raisebox{0.15ex}{\accentmark{12}{\char"1CF9}}नापातामाहू\raisebox{0.25ex}{\accentmark{12}{\char"0951}}वे\kern 0.05em ।\hskip 0.12em plus 0.04em minus 0.02em प्री\raisebox{0.25ex}{\accentmark{15}{\char"1CD2}}यञ्चेतिष्ठामारातिं\raisebox{0.15ex}{\accentmark{12}{\char"1CF9}}स्वाध्वारम्\raisebox{0.25ex}{\accentmark{15}{\char"1CD2}}\kern 0.05em ।\hskip 0.12em plus 0.04em minus 0.02em विश्वा\raisebox{0.25ex}{\accentmark{12}{\char"0951}}स्यादूता\raisebox{0.15ex}{\accentmark{12}{\char"1CF8}}मामृत\raisebox{0.25ex}{\accentmark{12}{\char"0951}}म् \nolinebreak\hskip 0.08em \mbox{॥१॥}}} \par}

\vspace{0.8em}
                % --- SUBSECTIONS (Mantras) ---
                
                % Conditional rendering based on output_mode
                \par\filbreak

\phantomsection

\addcontentsline{toc}{subsection}{ऋक् २}

{\centering \footnotesize \textcolor{AccentPurple}{॥ भगः प्रागाथः अग्निः बृहती ॥} \par}

\vspace{0.12em}

{\centering \textbf{\textcolor{AccentBlue}{शे\raisebox{0.25ex}{\accentmark{15}{\char"1CD2}}षेवा\raisebox{0.25ex}{\accentmark{15}{\char"1CD2}}ने\raisebox{0.25ex}{\accentmark{12}{\char"0951}}षूमातृ\raisebox{0.25ex}{\accentmark{12}{\char"0951}}षु\kern 0.05em ।\hskip 0.12em plus 0.04em minus 0.02em स\raisebox{0.25ex}{\accentmark{15}{\char"1CD2}}न्त्वामर्क्ता\raisebox{0.25ex}{\accentmark{12}{\char"0951}}सइन्धते\kern 0.05em ।\hskip 0.12em plus 0.04em minus 0.02em आत\raisebox{0.25ex}{\accentmark{12}{\char"0951}}न्द्रोहव्यं\raisebox{0.25ex}{\accentmark{15}{\char"1CD2}}वा\raisebox{0.25ex}{\accentmark{12}{\char"0951}}हसिहाविष्कृता\raisebox{0.25ex}{\accentmark{12}{\char"0951}}ः\kern 0.05em ।\hskip 0.12em plus 0.04em minus 0.02em आ\raisebox{0.25ex}{\accentmark{15}{\char"1CD2}}दि\raisebox{0.25ex}{\accentmark{15}{\char"1CD2}}द्देवेषु\raisebox{0.25ex}{\accentmark{12}{\char"0951}}राजसि \nolinebreak\hskip 0.08em \mbox{॥२॥}}} \par}

\vspace{0.8em}
                % --- SUBSECTIONS (Mantras) ---
                
                % Conditional rendering based on output_mode
                \par\filbreak

\phantomsection

\addcontentsline{toc}{subsection}{ऋक् ३}

{\centering \footnotesize \textcolor{AccentPurple}{॥ सोभिरः काण्वः अग्निः बृहती ॥} \par}

\vspace{0.12em}

{\centering \textbf{\textcolor{AccentBlue}{अ\raisebox{0.25ex}{\accentmark{15}{\char"1CD2}}द\raisebox{0.25ex}{\accentmark{12}{\char"0951}}र्शीगातूवि\raisebox{0.25ex}{\accentmark{15}{\char"1CD2}}क्ता\raisebox{0.25ex}{\accentmark{12}{\char"0951}}मः\kern 0.05em ।\hskip 0.12em plus 0.04em minus 0.02em यस्मि\raisebox{0.25ex}{\accentmark{12}{\char"0951}}न्व्राता\raisebox{0.15ex}{\accentmark{12}{\char"1CF9}}न्याद\raisebox{0.25ex}{\accentmark{15}{\char"1CD2}}धु\raisebox{0.25ex}{\accentmark{15}{\char"1CD2}}ः\kern 0.05em ।\hskip 0.12em plus 0.04em minus 0.02em ऊपोषूजाता\raisebox{0.15ex}{\accentmark{12}{\char"1CF9}}मा\raisebox{0.15ex}{\accentmark{12}{\char"1CF9}}र्यास्यावार्थानम्\kern 0.05em ।\hskip 0.12em plus 0.04em minus 0.02em अ\raisebox{0.25ex}{\accentmark{15}{\char"1CD2}}ग्नि\raisebox{0.25ex}{\accentmark{15}{\char"1CD2}}न्ना\raisebox{0.25ex}{\accentmark{12}{\char"0951}}क्षन्तूनोगीरा\raisebox{0.25ex}{\accentmark{12}{\char"0951}}ः \nolinebreak\hskip 0.08em \mbox{॥३॥}}} \par}

\vspace{0.8em}
                % --- SUBSECTIONS (Mantras) ---
                
                % Conditional rendering based on output_mode
                \par\filbreak

\phantomsection

\addcontentsline{toc}{subsection}{ऋक् ४}

{\centering \footnotesize \textcolor{AccentPurple}{॥ मनवुर्वैस्वस्तः अग्निः बृहती ॥} \par}

\vspace{0.12em}

{\centering \textbf{\textcolor{AccentBlue}{अग्नी\raisebox{0.25ex}{\accentmark{15}{\char"1CD2}}रूत्थेपूरो\raisebox{0.25ex}{\accentmark{15}{\char"1CD2}}हित\raisebox{0.25ex}{\accentmark{15}{\char"1CD2}}ः\kern 0.05em ।\hskip 0.12em plus 0.04em minus 0.02em ग्रावा\raisebox{0.25ex}{\accentmark{12}{\char"0951}}णोबर्हीरा\raisebox{0.25ex}{\accentmark{12}{\char"0951}}ध्वारे\raisebox{0.25ex}{\accentmark{15}{\char"1CD2}}\kern 0.05em ।\hskip 0.12em plus 0.04em minus 0.02em ऋचा\raisebox{0.25ex}{\accentmark{15}{\char"1CD2}}या\raisebox{0.25ex}{\accentmark{12}{\char"0951}}मिमरुतोब्रह्मणस्पते\kern 0.05em ।\hskip 0.12em plus 0.04em minus 0.02em देवंआ\raisebox{0.25ex}{\accentmark{15}{\char"1CD2}}वोवा\raisebox{0.25ex}{\accentmark{15}{\char"1CD2}}रेण्यम् \nolinebreak\hskip 0.08em \mbox{॥४॥}}} \par}

\vspace{0.8em}
                % --- SUBSECTIONS (Mantras) ---
                
                % Conditional rendering based on output_mode
                \par\filbreak

\phantomsection

\addcontentsline{toc}{subsection}{ऋक् ५}

{\centering \footnotesize \textcolor{AccentPurple}{॥ सुदीतिपुरुमी{-}ळावांगिरसौ अग्निः बृहती ॥} \par}

\vspace{0.12em}

{\centering \textbf{\textcolor{AccentBlue}{अग्नी\raisebox{0.15ex}{\accentmark{12}{\char"1CF9}}मीळिष्वा\raisebox{0.25ex}{\accentmark{15}{\char"1CD2}}वा\raisebox{0.25ex}{\accentmark{12}{\char"0951}}से\kern 0.05em ।\hskip 0.12em plus 0.04em minus 0.02em गा\raisebox{0.25ex}{\accentmark{15}{\char"1CD2}}था\raisebox{0.25ex}{\accentmark{12}{\char"0951}}भीश्शीरा\raisebox{0.25ex}{\accentmark{15}{\char"1CD2}}शो\raisebox{0.25ex}{\accentmark{12}{\char"0951}}चिषम्\kern 0.05em ।\hskip 0.12em plus 0.04em minus 0.02em अ\raisebox{0.25ex}{\accentmark{15}{\char"1CD2}}ग्निं\raisebox{0.25ex}{\accentmark{15}{\char"1CD2}}राये\raisebox{0.25ex}{\accentmark{15}{\char"1CD2}}पू\raisebox{0.25ex}{\accentmark{12}{\char"0951}}रूमिढाश्रूतन्ना\raisebox{0.25ex}{\accentmark{15}{\char"1CD2}}रा\raisebox{0.25ex}{\accentmark{12}{\char"0951}}ः\kern 0.05em ।\hskip 0.12em plus 0.04em minus 0.02em अग्निस्सू\raisebox{0.25ex}{\accentmark{12}{\char"0951}}दीताये\raisebox{0.15ex}{\accentmark{12}{\char"1CF9}}च्छार्दिः \nolinebreak\hskip 0.08em \mbox{॥५॥}}} \par}

\vspace{0.8em}
                % --- SUBSECTIONS (Mantras) ---
                
                % Conditional rendering based on output_mode
                \par\filbreak

\phantomsection

\addcontentsline{toc}{subsection}{ऋक् ६}

{\centering \footnotesize \textcolor{AccentPurple}{॥ प्रस्कण्वः काण्वः अग्निः बृहती ॥} \par}

\vspace{0.12em}

{\centering \textbf{\textcolor{AccentBlue}{श्रूधि\raisebox{0.25ex}{\accentmark{15}{\char"1CD2}}श्रुत्कर्णावन्ही\raisebox{0.25ex}{\accentmark{12}{\char"0951}}भिः\kern 0.05em ।\hskip 0.12em plus 0.04em minus 0.02em देवैरा\raisebox{0.25ex}{\accentmark{12}{\char"0951}}ग्ने\raisebox{0.25ex}{\accentmark{15}{\char"1CD2}}सायावा\raisebox{0.25ex}{\accentmark{12}{\char"0951}}भिः\kern 0.05em ।\hskip 0.12em plus 0.04em minus 0.02em आसी\raisebox{0.25ex}{\accentmark{12}{\char"0951}}दतुबर्ही\raisebox{0.25ex}{\accentmark{15}{\char"1CD2}}षी\raisebox{0.25ex}{\accentmark{12}{\char"0951}}मित्रो\raisebox{0.15ex}{\accentmark{12}{\char"1CF9}}अर्यामा\kern 0.05em ।\hskip 0.12em plus 0.04em minus 0.02em प्रात\raisebox{0.25ex}{\accentmark{15}{\char"1CD2}}र्या\raisebox{0.25ex}{\accentmark{15}{\char"1CD2}}वा\raisebox{0.25ex}{\accentmark{12}{\char"0951}}भीराध्वारे\raisebox{0.25ex}{\accentmark{15}{\char"1CD2}} \nolinebreak\hskip 0.08em \mbox{॥६॥}}} \par}

\vspace{0.8em}
                % --- SUBSECTIONS (Mantras) ---
                
                % Conditional rendering based on output_mode
                \par\filbreak

\phantomsection

\addcontentsline{toc}{subsection}{ऋक् ७}

{\centering \footnotesize \textcolor{AccentPurple}{॥ सोभिरः काण्वः अग्निः बृहती ॥} \par}

\vspace{0.12em}

{\centering \textbf{\textcolor{AccentBlue}{प्रदै\raisebox{0.25ex}{\accentmark{15}{\char"1CD2}}वो\raisebox{0.25ex}{\accentmark{12}{\char"0951}}दासोअग्निः\kern 0.05em ।\hskip 0.12em plus 0.04em minus 0.02em दे\raisebox{0.25ex}{\accentmark{15}{\char"1CD2}}वइन्द्रोनाम\raisebox{0.25ex}{\accentmark{15}{\char"1CD2}}ज्मा\raisebox{0.25ex}{\accentmark{15}{\char"1CD2}}ना\kern 0.05em ।\hskip 0.12em plus 0.04em minus 0.02em आ\raisebox{0.15ex}{\accentmark{12}{\char"1CF8}}नू\raisebox{0.25ex}{\accentmark{15}{\char"1CD2}}मातारं\raisebox{0.25ex}{\accentmark{12}{\char"0951}}पृथिवींवीवा\raisebox{0.25ex}{\accentmark{12}{\char"0951}}वृते\kern 0.05em ।\hskip 0.12em plus 0.04em minus 0.02em तस्थौ\raisebox{0.15ex}{\accentmark{12}{\char"1CF9}}ना\raisebox{0.15ex}{\accentmark{12}{\char"1CF9}}का\raisebox{0.25ex}{\accentmark{12}{\char"0951}}स्याशर्माणि \nolinebreak\hskip 0.08em \mbox{॥७॥}}} \par}

\vspace{0.8em}
                % --- SUBSECTIONS (Mantras) ---
                
                % Conditional rendering based on output_mode
                \par\filbreak

\phantomsection

\addcontentsline{toc}{subsection}{ऋक् ८}

{\centering \footnotesize \textcolor{AccentPurple}{॥ मेधातिमध्येआतिथी काण्वौ इन्द्रः बृहती ॥} \par}

\vspace{0.12em}

{\centering \textbf{\textcolor{AccentBlue}{आधा\raisebox{0.25ex}{\accentmark{12}{\char"0951}}ज्मो\raisebox{0.25ex}{\accentmark{15}{\char"1CD2}}आधा\raisebox{0.25ex}{\accentmark{12}{\char"0951}}वादी\raisebox{0.25ex}{\accentmark{12}{\char"0951}}वः\kern 0.05em ।\hskip 0.12em plus 0.04em minus 0.02em बृ\raisebox{0.25ex}{\accentmark{15}{\char"1CD2}}हातो\raisebox{0.15ex}{\accentmark{12}{\char"1CF9}}रो\raisebox{0.25ex}{\accentmark{12}{\char"0951}}चानाद\raisebox{0.25ex}{\accentmark{15}{\char"1CD2}}धी\raisebox{0.25ex}{\accentmark{12}{\char"0951}}\kern 0.05em ।\hskip 0.12em plus 0.04em minus 0.02em आया\raisebox{0.25ex}{\accentmark{15}{\char"1CD2}}वा\raisebox{0.15ex}{\accentmark{12}{\char"1CF8}}र्धस्वा\raisebox{0.25ex}{\accentmark{15}{\char"1CD2}}तन्वा\raisebox{0.25ex}{\accentmark{15}{\char"1CD2}}गीरा\raisebox{0.25ex}{\accentmark{15}{\char"1CD2}}मामा\kern 0.05em ।\hskip 0.12em plus 0.04em minus 0.02em आजाता\raisebox{0.25ex}{\accentmark{15}{\char"1CD2}}सूकृतॊ\raisebox{0.25ex}{\accentmark{12}{\char"0951}}पृणा \nolinebreak\hskip 0.08em \mbox{॥८॥}}} \par}

\vspace{0.8em}
                % --- SUBSECTIONS (Mantras) ---
                
                % Conditional rendering based on output_mode
                \par\filbreak

\phantomsection

\addcontentsline{toc}{subsection}{ऋक् ९}

{\centering \footnotesize \textcolor{AccentPurple}{॥ विश्वामित्रो गाथिनः अग्निः बृहती ॥} \par}

\vspace{0.12em}

{\centering \textbf{\textcolor{AccentBlue}{काया\raisebox{0.25ex}{\accentmark{12}{\char"0951}}मानोवा\raisebox{0.25ex}{\accentmark{15}{\char"1CD2}}नात्वम्\raisebox{0.25ex}{\accentmark{15}{\char"1CD2}}\kern 0.05em ।\hskip 0.12em plus 0.04em minus 0.02em य\raisebox{0.15ex}{\accentmark{12}{\char"1CF8}}न्मा\raisebox{0.25ex}{\accentmark{12}{\char"0951}}त्ररा\raisebox{0.25ex}{\accentmark{15}{\char"1CD2}}ज\raisebox{0.25ex}{\accentmark{12}{\char"0951}}गन्नाप\raisebox{0.25ex}{\accentmark{15}{\char"1CD2}}ः\kern 0.05em ।\hskip 0.12em plus 0.04em minus 0.02em ना\raisebox{0.25ex}{\accentmark{15}{\char"1CD2}}ताक्ते\raisebox{0.25ex}{\accentmark{12}{\char"0951}}अग्नेप्रामृ\raisebox{0.15ex}{\accentmark{12}{\char"1CF9}}षेनीवार्त्ता\raisebox{0.25ex}{\accentmark{12}{\char"0951}}नम्\kern 0.05em ।\hskip 0.12em plus 0.04em minus 0.02em य\raisebox{0.25ex}{\accentmark{15}{\char"1CD2}}द्दूरेस\raisebox{0.25ex}{\accentmark{15}{\char"1CD2}}न्नी\raisebox{0.25ex}{\accentmark{15}{\char"1CD2}}हाभूवः \nolinebreak\hskip 0.08em \mbox{॥९॥}}} \par}

\vspace{0.8em}
                % --- SUBSECTIONS (Mantras) ---
                
                % Conditional rendering based on output_mode
                \par\filbreak

\phantomsection

\addcontentsline{toc}{subsection}{ऋक् १०}

{\centering \footnotesize \textcolor{AccentPurple}{॥ कण्वो घौरः अग्निः बृहती ॥} \par}

\vspace{0.12em}

{\centering \textbf{\textcolor{AccentBlue}{नित्वा\raisebox{0.25ex}{\accentmark{12}{\char"0951}}माग्ने\raisebox{0.25ex}{\accentmark{15}{\char"1CD2}}मानू\raisebox{0.25ex}{\accentmark{12}{\char"0951}}र्दधे\kern 0.05em ।\hskip 0.12em plus 0.04em minus 0.02em ज्यो\raisebox{0.25ex}{\accentmark{15}{\char"1CD2}}तीर्ज\raisebox{0.25ex}{\accentmark{12}{\char"0951}}ना\raisebox{0.25ex}{\accentmark{12}{\char"0951}}याशा\raisebox{0.25ex}{\accentmark{15}{\char"1CD2}}श्वा\raisebox{0.25ex}{\accentmark{12}{\char"0951}}ते\kern 0.05em ।\hskip 0.12em plus 0.04em minus 0.02em दीदेथाक\raisebox{0.15ex}{\accentmark{12}{\char"1CF9}}ण्वऋताजातउक्षीतः\kern 0.05em ।\hskip 0.12em plus 0.04em minus 0.02em यन्ना\raisebox{0.25ex}{\accentmark{12}{\char"0951}}मास्य\raisebox{0.15ex}{\accentmark{12}{\char"1CF9}}न्तीकृष्टायाः \nolinebreak\hskip 0.08em \mbox{॥१०॥}}} \par}

\vspace{0.8em}

            % --- End of Kandah Closing Floral Line & Counts ---
            \par\nopagebreak\vspace{0.9em plus 0.3em minus 0.1em}
            \centerline{\textcolor{black!35}{\rule{1.8cm}{0.35pt}}\;\raisebox{-0.22ex}{\textcolor{vedicmaroon}{\large\SixFlowerPetalDotted}}\;\textcolor{black!35}{\rule{1.8cm}{0.35pt}}}
            \nopagebreak\vspace{0.45em}
            \centerline{\devafont\textbf{\textcolor{AccentBlue}{॥ ऋक् सङ्ख्या : १० ॥}}}
        
            \clearpage

            % --- SECTION ---
            
            \phantomsection  

            \section*{\centering षष्ठ खण्डः}
            \addcontentsline{toc}{section}{षष्ठ खण्डः} 
            \markright{षष्ठ खण्डः}
            \nopagebreak % Keep section title with first mantra
            
            % --- SECTION-LEVEL PROCEDURE LINK ---

                % --- SUBSECTIONS (Mantras) ---
                
                % Conditional rendering based on output_mode
                \par\filbreak

\phantomsection

\addcontentsline{toc}{subsection}{ऋक् १}

{\centering \footnotesize \textcolor{AccentPurple}{॥ वसिष्ठो मैत्रावरुणिः अग्निः बृहती ॥} \par}

\vspace{0.12em}

{\centering \textbf{\textcolor{AccentBlue}{दे\raisebox{0.25ex}{\accentmark{15}{\char"1CD2}}वो\raisebox{0.25ex}{\accentmark{15}{\char"1CD2}}वो\raisebox{0.25ex}{\accentmark{12}{\char"0951}}द्रविणोदाः\kern 0.05em ।\hskip 0.12em plus 0.04em minus 0.02em पूर्णां\raisebox{0.25ex}{\accentmark{12}{\char"0951}}वी\raisebox{0.25ex}{\accentmark{12}{\char"0951}}वाष्ट्वासी\raisebox{0.25ex}{\accentmark{15}{\char"1CD2}}च\raisebox{0.25ex}{\accentmark{12}{\char"0951}}म्\kern 0.05em ।\hskip 0.12em plus 0.04em minus 0.02em उद्वा\raisebox{0.25ex}{\accentmark{12}{\char"0951}}सिन्चा\raisebox{0.25ex}{\accentmark{15}{\char"1CD2}}ध्वामू\raisebox{0.25ex}{\accentmark{15}{\char"1CD2}}पा\raisebox{0.25ex}{\accentmark{12}{\char"0951}}वापृणध्वं\kern 0.05em ।\hskip 0.12em plus 0.04em minus 0.02em आदीद्वो\raisebox{0.25ex}{\accentmark{12}{\char"0951}}देवा\raisebox{0.25ex}{\accentmark{15}{\char"1CD2}}ओहते \nolinebreak\hskip 0.08em \mbox{॥१॥}}} \par}

\vspace{0.8em}
                % --- SUBSECTIONS (Mantras) ---
                
                % Conditional rendering based on output_mode
                \par\filbreak

\phantomsection

\addcontentsline{toc}{subsection}{ऋक् २}

{\centering \footnotesize \textcolor{AccentPurple}{॥ कण्वो घौरः ब्रह्मणस्पतिः बृहती ॥} \par}

\vspace{0.12em}

{\centering \textbf{\textcolor{AccentBlue}{प्रैतू\raisebox{0.15ex}{\accentmark{12}{\char"1CF9}}ब्रह्माण\raisebox{0.25ex}{\accentmark{15}{\char"1CD2}}स्पा\raisebox{0.25ex}{\accentmark{15}{\char"1CD2}}तीः\kern 0.05em ।\hskip 0.12em plus 0.04em minus 0.02em प्रा\raisebox{0.25ex}{\accentmark{15}{\char"1CD2}}देव्येतूसूनृ\raisebox{0.25ex}{\accentmark{15}{\char"1CD2}}ता\kern 0.05em ।\hskip 0.12em plus 0.04em minus 0.02em अच्छावीरं\raisebox{0.15ex}{\accentmark{12}{\char"1CF9}}नर्यं\raisebox{0.25ex}{\accentmark{12}{\char"0951}}पङ्तीरा\raisebox{0.25ex}{\accentmark{12}{\char"0951}}धसम्\kern 0.05em ।\hskip 0.12em plus 0.04em minus 0.02em दे\raisebox{0.25ex}{\accentmark{15}{\char"1CD2}}वा\raisebox{0.25ex}{\accentmark{15}{\char"1CD2}}यज्ञ\raisebox{0.25ex}{\accentmark{15}{\char"1CD2}}न्ना\raisebox{0.25ex}{\accentmark{15}{\char"1CD2}}यन्तुनः \nolinebreak\hskip 0.08em \mbox{॥२॥}}} \par}

\vspace{0.8em}
                % --- SUBSECTIONS (Mantras) ---
                
                % Conditional rendering based on output_mode
                \par\filbreak

\phantomsection

\addcontentsline{toc}{subsection}{ऋक् ३}

{\centering \footnotesize \textcolor{AccentPurple}{॥ कण्वो घौरः यूपः बृहती ॥} \par}

\vspace{0.12em}

{\centering \textbf{\textcolor{AccentBlue}{ऊर्ध्वाऊषू\raisebox{0.15ex}{\accentmark{12}{\char"1CF9}}णा\raisebox{0.25ex}{\accentmark{12}{\char"0951}}ऊता\raisebox{0.25ex}{\accentmark{15}{\char"1CD2}}ये\raisebox{0.25ex}{\accentmark{12}{\char"0951}}\kern 0.05em ।\hskip 0.12em plus 0.04em minus 0.02em तिष्ठा\raisebox{0.25ex}{\accentmark{12}{\char"0951}}देवो\raisebox{0.15ex}{\accentmark{12}{\char"1CF9}}ना\raisebox{0.15ex}{\accentmark{12}{\char"1CF9}}सा\raisebox{0.25ex}{\accentmark{12}{\char"0951}}वीता\kern 0.05em ।\hskip 0.12em plus 0.04em minus 0.02em ऊर्ध्वोवा\raisebox{0.25ex}{\accentmark{15}{\char"1CD2}}ज\raisebox{0.25ex}{\accentmark{12}{\char"0951}}स्यासानीता\raisebox{0.15ex}{\accentmark{12}{\char"1CF9}}यादन्जी\raisebox{0.25ex}{\accentmark{15}{\char"1CD2}}भि\raisebox{0.25ex}{\accentmark{12}{\char"0951}}ः\kern 0.05em ।\hskip 0.12em plus 0.04em minus 0.02em वा\raisebox{0.25ex}{\accentmark{15}{\char"1CD2}}घा\raisebox{0.15ex}{\accentmark{12}{\char"1CF9}}त्भी\raisebox{0.25ex}{\accentmark{12}{\char"0951}}र्वीभ्वाया\raisebox{0.25ex}{\accentmark{12}{\char"0951}}महे \nolinebreak\hskip 0.08em \mbox{॥३॥}}} \par}

\vspace{0.8em}
                % --- SUBSECTIONS (Mantras) ---
                
                % Conditional rendering based on output_mode
                \par\filbreak

\phantomsection

\addcontentsline{toc}{subsection}{ऋक् ४}

{\centering \footnotesize \textcolor{AccentPurple}{॥ सोभिरः काण्वः अग्निः बृहती ॥} \par}

\vspace{0.12em}

{\centering \textbf{\textcolor{AccentBlue}{प्र\raisebox{0.15ex}{\accentmark{12}{\char"1CF8}}योरा\raisebox{0.25ex}{\accentmark{12}{\char"0951}}येनीनीषते\kern 0.05em ।\hskip 0.12em plus 0.04em minus 0.02em म\raisebox{0.25ex}{\accentmark{15}{\char"1CD2}}र्तोयस्तेवासोदा\raisebox{0.25ex}{\accentmark{12}{\char"0951}}शा\raisebox{0.25ex}{\accentmark{12}{\char"0951}}त्\kern 0.05em ।\hskip 0.12em plus 0.04em minus 0.02em सावीरन्धाक्तेअ\raisebox{0.25ex}{\accentmark{15}{\char"1CD2}}ग्निउत्थाशंसी\raisebox{0.25ex}{\accentmark{12}{\char"0951}}नम्\kern 0.05em ।\hskip 0.12em plus 0.04em minus 0.02em इत्माना\raisebox{0.25ex}{\accentmark{12}{\char"0951}}सहस्रापोषी\raisebox{0.25ex}{\accentmark{15}{\char"1CD2}}णम् \nolinebreak\hskip 0.08em \mbox{॥४॥}}} \par}

\vspace{0.8em}
                % --- SUBSECTIONS (Mantras) ---
                
                % Conditional rendering based on output_mode
                \par\filbreak

\phantomsection

\addcontentsline{toc}{subsection}{ऋक् ५}

{\centering \footnotesize \textcolor{AccentPurple}{॥ कण्वो घौरः अग्निः बृहती ॥} \par}

\vspace{0.12em}

{\centering \textbf{\textcolor{AccentBlue}{प्रावो\raisebox{0.25ex}{\accentmark{12}{\char"0951}}यंभंपू\raisebox{0.25ex}{\accentmark{12}{\char"0951}}रूणाम्\raisebox{0.25ex}{\accentmark{15}{\char"1CD2}}\kern 0.05em ।\hskip 0.12em plus 0.04em minus 0.02em वीशा\raisebox{0.25ex}{\accentmark{15}{\char"1CD2}}न्दे\raisebox{0.25ex}{\accentmark{12}{\char"0951}}वायातीनाम्\kern 0.05em ।\hskip 0.12em plus 0.04em minus 0.02em अग्निंसूक्तेभिर्वाचो\raisebox{0.25ex}{\accentmark{12}{\char"0951}}भिर्वृणीमहे\kern 0.05em ।\hskip 0.12em plus 0.04em minus 0.02em यंसमीदन्याइ\raisebox{0.25ex}{\accentmark{15}{\char"1CD2}}न्धा\raisebox{0.25ex}{\accentmark{15}{\char"1CD2}}ते \nolinebreak\hskip 0.08em \mbox{॥५॥}}} \par}

\vspace{0.8em}
                % --- SUBSECTIONS (Mantras) ---
                
                % Conditional rendering based on output_mode
                \par\filbreak

\phantomsection

\addcontentsline{toc}{subsection}{ऋक् ६}

{\centering \footnotesize \textcolor{AccentPurple}{॥ उत्कीलः कात्यः अग्निः बृहती ॥} \par}

\vspace{0.12em}

{\centering \textbf{\textcolor{AccentBlue}{आया\raisebox{0.25ex}{\accentmark{15}{\char"1CD2}}मग्नी\raisebox{0.25ex}{\accentmark{15}{\char"1CD2}}सूवी\raisebox{0.25ex}{\accentmark{15}{\char"1CD2}}र्या\raisebox{0.25ex}{\accentmark{12}{\char"0951}}स्या\kern 0.05em ।\hskip 0.12em plus 0.04em minus 0.02em ई\raisebox{0.25ex}{\accentmark{15}{\char"1CD2}}शेहिसौ\raisebox{0.25ex}{\accentmark{15}{\char"1CD2}}भा\raisebox{0.25ex}{\accentmark{15}{\char"1CD2}}गस्य\kern 0.05em ।\hskip 0.12em plus 0.04em minus 0.02em रायाई\raisebox{0.25ex}{\accentmark{15}{\char"1CD2}}शेस्वापत्यास्यागो\raisebox{0.25ex}{\accentmark{15}{\char"1CD2}}मा\raisebox{0.25ex}{\accentmark{12}{\char"0951}}तः\kern 0.05em ।\hskip 0.12em plus 0.04em minus 0.02em ईशे\raisebox{0.25ex}{\accentmark{12}{\char"0951}}हिवृ\raisebox{0.25ex}{\accentmark{15}{\char"1CD2}}त्राहा\raisebox{0.25ex}{\accentmark{15}{\char"1CD2}}थानाम् \nolinebreak\hskip 0.08em \mbox{॥६॥}}} \par}

\vspace{0.8em}
                % --- SUBSECTIONS (Mantras) ---
                
                % Conditional rendering based on output_mode
                \par\filbreak

\phantomsection

\addcontentsline{toc}{subsection}{ऋक् ७}

{\centering \footnotesize \textcolor{AccentPurple}{॥ वसिष्ठो मैत्रावरुणिः अग्निः बृहती ॥} \par}

\vspace{0.12em}

{\centering \textbf{\textcolor{AccentBlue}{त्वा\raisebox{0.15ex}{\accentmark{12}{\char"1CF8}}मा\raisebox{0.25ex}{\accentmark{12}{\char"0951}}ग्नेगृहा\raisebox{0.25ex}{\accentmark{15}{\char"1CD2}}पा\raisebox{0.25ex}{\accentmark{12}{\char"0951}}तिः\kern 0.05em ।\hskip 0.12em plus 0.04em minus 0.02em त्वं\raisebox{0.25ex}{\accentmark{15}{\char"1CD2}}होता\raisebox{0.25ex}{\accentmark{12}{\char"0951}}नोआ\raisebox{0.15ex}{\accentmark{12}{\char"1CF8}}ध्वारे\kern 0.05em ।\hskip 0.12em plus 0.04em minus 0.02em त्वा\raisebox{0.25ex}{\accentmark{15}{\char"1CD2}}म्पो\raisebox{0.25ex}{\accentmark{15}{\char"1CD2}}ता\raisebox{0.25ex}{\accentmark{12}{\char"0951}}विश्वावारप्रा\raisebox{0.25ex}{\accentmark{15}{\char"1CD2}}चे\raisebox{0.25ex}{\accentmark{12}{\char"0951}}तः\kern 0.05em ।\hskip 0.12em plus 0.04em minus 0.02em य\raisebox{0.25ex}{\accentmark{15}{\char"1CD2}}क्षीया\raisebox{0.15ex}{\accentmark{12}{\char"1CF9}}सी\raisebox{0.25ex}{\accentmark{12}{\char"0951}}चावार्यम् \nolinebreak\hskip 0.08em \mbox{॥७॥}}} \par}

\vspace{0.8em}
                % --- SUBSECTIONS (Mantras) ---
                
                % Conditional rendering based on output_mode
                \par\filbreak

\phantomsection

\addcontentsline{toc}{subsection}{ऋक् ८}

{\centering \footnotesize \textcolor{AccentPurple}{॥ विश्वामित्रो गाथिनः अग्निः बृहती ॥} \par}

\vspace{0.12em}

{\centering \textbf{\textcolor{AccentBlue}{साखायस्त्वाववृमहे\kern 0.05em ।\hskip 0.12em plus 0.04em minus 0.02em दे\raisebox{0.25ex}{\accentmark{15}{\char"1CD2}}वम्मर्ता\raisebox{0.25ex}{\accentmark{15}{\char"1CD2}}साऊता\raisebox{0.25ex}{\accentmark{15}{\char"1CD2}}ये\kern 0.05em ।\hskip 0.12em plus 0.04em minus 0.02em आ\raisebox{0.25ex}{\accentmark{15}{\char"1CD2}}पान्ना\raisebox{0.25ex}{\accentmark{12}{\char"0951}}पातंसूभा\raisebox{0.15ex}{\accentmark{12}{\char"1CF9}}गं\raisebox{0.25ex}{\accentmark{12}{\char"0951}}सूदंसा\raisebox{0.25ex}{\accentmark{12}{\char"0951}}सम्\kern 0.05em ।\hskip 0.12em plus 0.04em minus 0.02em सु\raisebox{0.25ex}{\accentmark{15}{\char"1CD2}}प्रा\raisebox{0.25ex}{\accentmark{15}{\char"1CD2}}तू\raisebox{0.25ex}{\accentmark{12}{\char"0951}}र्तीमानेहा\raisebox{0.25ex}{\accentmark{15}{\char"1CD2}}सम् \nolinebreak\hskip 0.08em \mbox{॥८॥}}} \par}

\vspace{0.8em}

            % --- End of Kandah Closing Floral Line & Counts ---
            \par\nopagebreak\vspace{0.9em plus 0.3em minus 0.1em}
            \centerline{\textcolor{black!35}{\rule{1.8cm}{0.35pt}}\;\raisebox{-0.22ex}{\textcolor{vedicmaroon}{\large\SixFlowerPetalDotted}}\;\textcolor{black!35}{\rule{1.8cm}{0.35pt}}}
            \nopagebreak\vspace{0.45em}
            \centerline{\devafont\textbf{\textcolor{AccentBlue}{॥ ऋक् सङ्ख्या : ८ ॥}}}
        
            \clearpage

            % --- SECTION ---
            
            \phantomsection  

            \section*{\centering सप्तम खण्डः}
            \addcontentsline{toc}{section}{सप्तम खण्डः} 
            \markright{सप्तम खण्डः}
            \nopagebreak % Keep section title with first mantra
            
            % --- SECTION-LEVEL PROCEDURE LINK ---

                % --- SUBSECTIONS (Mantras) ---
                
                % Conditional rendering based on output_mode
                \par\filbreak

\phantomsection

\addcontentsline{toc}{subsection}{ऋक् १}

{\centering \footnotesize \textcolor{AccentPurple}{॥ श्यावाश्वो वामदेवो वा अग्निः त्रिष्टुप् ॥} \par}

\vspace{0.12em}

{\centering \textbf{\textcolor{AccentBlue}{आजू\raisebox{0.25ex}{\accentmark{12}{\char"0951}}होताहविषा\raisebox{0.25ex}{\accentmark{12}{\char"0951}}मर्जयध्वम्\kern 0.05em ।\hskip 0.12em plus 0.04em minus 0.02em नीहो\raisebox{0.25ex}{\accentmark{15}{\char"1CD2}}ता\raisebox{0.25ex}{\accentmark{12}{\char"0951}}रंगृहा\raisebox{0.25ex}{\accentmark{15}{\char"1CD2}}पा\raisebox{0.25ex}{\accentmark{12}{\char"0951}}तिंदधिध्वम्\kern 0.05em ।\hskip 0.12em plus 0.04em minus 0.02em ईळा\raisebox{0.25ex}{\accentmark{12}{\char"0951}}स्पातेनामा\raisebox{0.25ex}{\accentmark{12}{\char"0951}}सारा\raisebox{0.25ex}{\accentmark{15}{\char"1CD2}}ताहा\raisebox{0.25ex}{\accentmark{12}{\char"0951}}व्यम्\kern 0.05em ।\hskip 0.12em plus 0.04em minus 0.02em सपर्यातायाजतम्पस्त्याठ\raisebox{0.25ex}{\accentmark{12}{\char"0951}}नाम् \nolinebreak\hskip 0.08em \mbox{॥१॥}}} \par}

\vspace{0.8em}
                % --- SUBSECTIONS (Mantras) ---
                
                % Conditional rendering based on output_mode
                \par\filbreak

\phantomsection

\addcontentsline{toc}{subsection}{ऋक् २}

{\centering \footnotesize \textcolor{AccentPurple}{॥ उपस्तुतो वाहिष्ठव्यः अग्निः जगती ॥} \par}

\vspace{0.12em}

{\centering \textbf{\textcolor{AccentBlue}{चित्र\raisebox{0.25ex}{\accentmark{15}{\char"1CD2}}इच्छीशो\raisebox{0.25ex}{\accentmark{15}{\char"1CD2}}स्ता\raisebox{0.25ex}{\accentmark{15}{\char"1CD2}}रुणास्यावक्षा\raisebox{0.25ex}{\accentmark{15}{\char"1CD2}}था\raisebox{0.25ex}{\accentmark{12}{\char"0951}}ः\kern 0.05em ।\hskip 0.12em plus 0.04em minus 0.02em न\raisebox{0.25ex}{\accentmark{15}{\char"1CD2}}यो\raisebox{0.25ex}{\accentmark{15}{\char"1CD2}}माता\raisebox{0.15ex}{\accentmark{12}{\char"1CF9}}रा\raisebox{0.25ex}{\accentmark{12}{\char"0951}}वन्वे\raisebox{0.25ex}{\accentmark{15}{\char"1CD2}}तीधा\raisebox{0.25ex}{\accentmark{15}{\char"1CD2}}ता\raisebox{0.25ex}{\accentmark{12}{\char"0951}}वे\kern 0.05em ।\hskip 0.12em plus 0.04em minus 0.02em अ\raisebox{0.25ex}{\accentmark{15}{\char"1CD2}}नुधा\raisebox{0.25ex}{\accentmark{15}{\char"1CD2}}याद\raisebox{0.25ex}{\accentmark{15}{\char"1CD2}}जीज\raisebox{0.25ex}{\accentmark{12}{\char"0951}}नाद\raisebox{0.25ex}{\accentmark{15}{\char"1CD2}}धा\raisebox{0.25ex}{\accentmark{12}{\char"0951}}चीदा\raisebox{0.25ex}{\accentmark{15}{\char"1CD2}}\kern 0.05em ।\hskip 0.12em plus 0.04em minus 0.02em वावक्षस्सद्योमा\raisebox{0.25ex}{\accentmark{12}{\char"0951}}हीदूत्याटट्यञ्चा\raisebox{0.25ex}{\accentmark{15}{\char"1CD2}}र॑\raisebox{0.25ex}{\accentmark{12}{\char"0951}}न् \nolinebreak\hskip 0.08em \mbox{॥२॥}}} \par}

\vspace{0.8em}
                % --- SUBSECTIONS (Mantras) ---
                
                % Conditional rendering based on output_mode
                \par\filbreak

\phantomsection

\addcontentsline{toc}{subsection}{ऋक् ३}

{\centering \footnotesize \textcolor{AccentPurple}{॥ बृहदुक्तो वामदेव्यः अग्निः त्रिष्टुप् ॥} \par}

\vspace{0.12em}

{\centering \textbf{\textcolor{AccentBlue}{ईद\raisebox{0.25ex}{\accentmark{15}{\char"1CD2}}न्ताएकं\raisebox{0.25ex}{\accentmark{12}{\char"0951}}पारा\raisebox{0.15ex}{\accentmark{12}{\char"1CF9}}ऊ\raisebox{0.25ex}{\accentmark{12}{\char"0951}}ताएकम्\raisebox{0.25ex}{\accentmark{12}{\char"0951}}\kern 0.05em ।\hskip 0.12em plus 0.04em minus 0.02em तृतीये\raisebox{0.25ex}{\accentmark{12}{\char"0951}}नज्योती\raisebox{0.25ex}{\accentmark{12}{\char"0951}}षासं\raisebox{0.25ex}{\accentmark{15}{\char"1CD2}}वीशस्वा\kern 0.05em ।\hskip 0.12em plus 0.04em minus 0.02em सं\raisebox{0.25ex}{\accentmark{15}{\char"1CD2}}वे\raisebox{0.25ex}{\accentmark{15}{\char"1CD2}}शान\raisebox{0.15ex}{\accentmark{12}{\char"1CF9}}स्तन्वेटट्यचारू\raisebox{0.25ex}{\accentmark{15}{\char"1CD2}}रेधि\kern 0.05em ।\hskip 0.12em plus 0.04em minus 0.02em प्रियो\raisebox{0.25ex}{\accentmark{12}{\char"0951}}देवा\raisebox{0.25ex}{\accentmark{15}{\char"1CD2}}ना\raisebox{0.25ex}{\accentmark{12}{\char"0951}}म्पा\raisebox{0.25ex}{\accentmark{12}{\char"0951}}रामेज\raisebox{0.25ex}{\accentmark{15}{\char"1CD2}}नित्रे \nolinebreak\hskip 0.08em \mbox{॥३॥}}} \par}

\vspace{0.8em}
                % --- SUBSECTIONS (Mantras) ---
                
                % Conditional rendering based on output_mode
                \par\filbreak

\phantomsection

\addcontentsline{toc}{subsection}{ऋक् ४}

{\centering \footnotesize \textcolor{AccentPurple}{॥ कुत्स आंगिरसः अग्निः जगती ॥} \par}

\vspace{0.12em}

{\centering \textbf{\textcolor{AccentBlue}{इमं\raisebox{0.25ex}{\accentmark{15}{\char"1CD2}}स्तोमामर्हातेजाता\raisebox{0.25ex}{\accentmark{15}{\char"1CD2}}वेदसे\kern 0.05em ।\hskip 0.12em plus 0.04em minus 0.02em राथा\raisebox{0.25ex}{\accentmark{12}{\char"0951}}मी\raisebox{0.15ex}{\accentmark{12}{\char"1CF9}}वा\raisebox{0.25ex}{\accentmark{15}{\char"1CD2}}संमा\raisebox{0.25ex}{\accentmark{12}{\char"0951}}हेमामा\raisebox{0.25ex}{\accentmark{12}{\char"0951}}नीषा\raisebox{0.25ex}{\accentmark{15}{\char"1CD2}}या\kern 0.05em ।\hskip 0.12em plus 0.04em minus 0.02em भ\raisebox{0.25ex}{\accentmark{15}{\char"1CD2}}द्रा\raisebox{0.25ex}{\accentmark{15}{\char"1CD2}}हीनःप्रामा\raisebox{0.25ex}{\accentmark{12}{\char"0951}}तीरस्यासंसादि\kern 0.05em ।\hskip 0.12em plus 0.04em minus 0.02em अग्ने\raisebox{0.25ex}{\accentmark{12}{\char"0951}}सख्येमारी\raisebox{0.25ex}{\accentmark{12}{\char"0951}}षामावायन्ता\raisebox{0.25ex}{\accentmark{15}{\char"1CD2}}वा\raisebox{0.25ex}{\accentmark{12}{\char"0951}} \nolinebreak\hskip 0.08em \mbox{॥४॥}}} \par}

\vspace{0.8em}
                % --- SUBSECTIONS (Mantras) ---
                
                % Conditional rendering based on output_mode
                \par\filbreak

\phantomsection

\addcontentsline{toc}{subsection}{ऋक् ५}

{\centering \footnotesize \textcolor{AccentPurple}{॥ भरद्वाजो बार्हस्पत्यः अग्निः त्रिष्टुप् ॥} \par}

\vspace{0.12em}

{\centering \textbf{\textcolor{AccentBlue}{मूर्धानंदीवोआ\raisebox{0.25ex}{\accentmark{12}{\char"0951}}रातिं\raisebox{0.15ex}{\accentmark{12}{\char"1CF9}}पृ\raisebox{0.25ex}{\accentmark{12}{\char"0951}}थिव्या\raisebox{0.25ex}{\accentmark{15}{\char"1CD2}}ः\kern 0.05em ।\hskip 0.12em plus 0.04em minus 0.02em वैश्वा\raisebox{0.25ex}{\accentmark{15}{\char"1CD2}}नारा\raisebox{0.25ex}{\accentmark{15}{\char"1CD2}}मृतआ\raisebox{0.25ex}{\accentmark{15}{\char"1CD2}}जाता\raisebox{0.25ex}{\accentmark{15}{\char"1CD2}}म\raisebox{0.25ex}{\accentmark{15}{\char"1CD2}}ग्निम्\kern 0.05em ।\hskip 0.12em plus 0.04em minus 0.02em काविं\raisebox{0.25ex}{\accentmark{15}{\char"1CD2}}सम्म्राजमा\raisebox{0.15ex}{\accentmark{12}{\char"1CF9}}तीथि\raisebox{0.25ex}{\accentmark{12}{\char"0951}}ञ्ज\raisebox{0.25ex}{\accentmark{12}{\char"0951}}ना\raisebox{0.25ex}{\accentmark{12}{\char"0951}}नाम्\kern 0.05em ।\hskip 0.12em plus 0.04em minus 0.02em आ\raisebox{0.25ex}{\accentmark{15}{\char"1CD2}}स\raisebox{0.25ex}{\accentmark{15}{\char"1CD2}}न्नापा\raisebox{0.25ex}{\accentmark{12}{\char"0951}}त्रञ्जनयन्तादेवाः \nolinebreak\hskip 0.08em \mbox{॥५॥}}} \par}

\vspace{0.8em}
                % --- SUBSECTIONS (Mantras) ---
                
                % Conditional rendering based on output_mode
                \par\filbreak

\phantomsection

\addcontentsline{toc}{subsection}{ऋक् ६}

{\centering \footnotesize \textcolor{AccentPurple}{॥ भरद्वाजो बार्हस्पत्यः अग्निः ॥} \par}

\vspace{0.12em}

{\centering \textbf{\textcolor{AccentBlue}{वित्वा\raisebox{0.25ex}{\accentmark{15}{\char"1CD2}}दापो\raisebox{0.25ex}{\accentmark{15}{\char"1CD2}}नापा\raisebox{0.25ex}{\accentmark{15}{\char"1CD2}}र्वातस्यापृष्ठात्\kern 0.05em ।\hskip 0.12em plus 0.04em minus 0.02em उक्थे\raisebox{0.25ex}{\accentmark{15}{\char"1CD2}}भी\raisebox{0.25ex}{\accentmark{12}{\char"0951}}रग्नेर्जनयन्तादे\raisebox{0.25ex}{\accentmark{15}{\char"1CD2}}वा\raisebox{0.25ex}{\accentmark{15}{\char"1CD2}}ः\kern 0.05em ।\hskip 0.12em plus 0.04em minus 0.02em तन्त्वा\raisebox{0.25ex}{\accentmark{15}{\char"1CD2}}गीरा\raisebox{0.25ex}{\accentmark{12}{\char"0951}}स्सू\raisebox{0.25ex}{\accentmark{12}{\char"0951}}ष्टू\raisebox{0.25ex}{\accentmark{12}{\char"0951}}तायो\raisebox{0.25ex}{\accentmark{12}{\char"0951}}र्वाजयन्ति\kern 0.05em ।\hskip 0.12em plus 0.04em minus 0.02em आ\raisebox{0.25ex}{\accentmark{15}{\char"1CD2}}जि\raisebox{0.15ex}{\accentmark{12}{\char"1CF9}}न्नागीर्वावाहो\raisebox{0.25ex}{\accentmark{12}{\char"0951}}जीग्यूरा\raisebox{0.25ex}{\accentmark{15}{\char"1CD2}}श्वाः \nolinebreak\hskip 0.08em \mbox{॥६॥}}} \par}

\vspace{0.8em}
                % --- SUBSECTIONS (Mantras) ---
                
                % Conditional rendering based on output_mode
                \par\filbreak

\phantomsection

\addcontentsline{toc}{subsection}{ऋक् ७}

{\centering \footnotesize \textcolor{AccentPurple}{॥ वामदेवो गौतमः अग्निः ॥} \par}

\vspace{0.12em}

{\centering \textbf{\textcolor{AccentBlue}{आ\raisebox{0.25ex}{\accentmark{15}{\char"1CD2}}वोरा\raisebox{0.25ex}{\accentmark{15}{\char"1CD2}}जानमध्वर\raisebox{0.15ex}{\accentmark{12}{\char"1CF9}}स्या\raisebox{0.25ex}{\accentmark{12}{\char"0951}}रूद्रं\raisebox{0.25ex}{\accentmark{15}{\char"1CD2}}\kern 0.05em ।\hskip 0.12em plus 0.04em minus 0.02em हो\raisebox{0.25ex}{\accentmark{15}{\char"1CD2}}ता\raisebox{0.25ex}{\accentmark{12}{\char"0951}}रंसत्याया\raisebox{0.25ex}{\accentmark{15}{\char"1CD2}}जंरो\raisebox{0.25ex}{\accentmark{15}{\char"1CD2}}द॑\raisebox{0.25ex}{\accentmark{12}{\char"0951}}स्योः\kern 0.05em ।\hskip 0.12em plus 0.04em minus 0.02em अग्निंपूरा\raisebox{0.25ex}{\accentmark{15}{\char"1CD2}}तानाइत्नो\raisebox{0.25ex}{\accentmark{15}{\char"1CD2}}राचिक्तात्\kern 0.05em ।\hskip 0.12em plus 0.04em minus 0.02em हिरण्यरूपामा\raisebox{0.25ex}{\accentmark{15}{\char"1CD2}}वा\raisebox{0.25ex}{\accentmark{12}{\char"0951}}सेकृणुध्वम् \nolinebreak\hskip 0.08em \mbox{॥७॥}}} \par}

\vspace{0.8em}
                % --- SUBSECTIONS (Mantras) ---
                
                % Conditional rendering based on output_mode
                \par\filbreak

\phantomsection

\addcontentsline{toc}{subsection}{ऋक् ८}

{\centering \footnotesize \textcolor{AccentPurple}{॥ वसिष्ठो मैत्रावरुणिः अग्निः ॥} \par}

\vspace{0.12em}

{\centering \textbf{\textcolor{AccentBlue}{इन्धेरा\raisebox{0.25ex}{\accentmark{15}{\char"1CD2}}जासामर्यो\raisebox{0.15ex}{\accentmark{12}{\char"1CF9}}नामो\raisebox{0.25ex}{\accentmark{12}{\char"0951}}भिः\kern 0.05em ।\hskip 0.12em plus 0.04em minus 0.02em यस्याप्रा\raisebox{0.15ex}{\accentmark{12}{\char"1CF9}}ती\raisebox{0.25ex}{\accentmark{12}{\char"0951}}कामा\raisebox{0.25ex}{\accentmark{12}{\char"0951}}हूतङ्घृतेना\kern 0.05em ।\hskip 0.12em plus 0.04em minus 0.02em नारो\raisebox{0.25ex}{\accentmark{15}{\char"1CD2}}हव्येभीरीडतेसाबाधाः\kern 0.05em ।\hskip 0.12em plus 0.04em minus 0.02em अ\raisebox{0.25ex}{\accentmark{15}{\char"1CD2}}ग्नी\raisebox{0.25ex}{\accentmark{15}{\char"1CD2}}राग्रामूषा\raisebox{0.25ex}{\accentmark{15}{\char"1CD2}}सा\raisebox{0.25ex}{\accentmark{15}{\char"1CD2}}माशोचि \nolinebreak\hskip 0.08em \mbox{॥८॥}}} \par}

\vspace{0.8em}
                % --- SUBSECTIONS (Mantras) ---
                
                % Conditional rendering based on output_mode
                \par\filbreak

\phantomsection

\addcontentsline{toc}{subsection}{ऋक् ९}

{\centering \footnotesize \textcolor{AccentPurple}{॥ त्रिशिरास्त्वाष्ट्र्ः अग्निः त्रिष्टुप् ॥} \par}

\vspace{0.12em}

{\centering \textbf{\textcolor{AccentBlue}{प्रा\raisebox{0.25ex}{\accentmark{15}{\char"1CD2}}केतुना\raisebox{0.25ex}{\accentmark{12}{\char"0951}}बृहाता\raisebox{0.15ex}{\accentmark{12}{\char"1CF9}}यात्य\raisebox{0.25ex}{\accentmark{15}{\char"1CD2}}ग्नि\raisebox{0.25ex}{\accentmark{15}{\char"1CD2}}ः\kern 0.05em ।\hskip 0.12em plus 0.04em minus 0.02em आ\raisebox{0.25ex}{\accentmark{15}{\char"1CD2}}रो\raisebox{0.25ex}{\accentmark{12}{\char"0951}}दसी\raisebox{0.25ex}{\accentmark{12}{\char"0951}}वृषाभो\raisebox{0.25ex}{\accentmark{15}{\char"1CD2}}रो\raisebox{0.25ex}{\accentmark{12}{\char"0951}}रवीति\kern 0.05em ।\hskip 0.12em plus 0.04em minus 0.02em दीवा\raisebox{0.25ex}{\accentmark{15}{\char"1CD2}}श्चीद\raisebox{0.25ex}{\accentmark{15}{\char"1CD2}}न्तादूपामा\raisebox{0.25ex}{\accentmark{15}{\char"1CD2}}मू\raisebox{0.25ex}{\accentmark{15}{\char"1CD2}}दानाट्\kern 0.05em ।\hskip 0.12em plus 0.04em minus 0.02em आ\raisebox{0.25ex}{\accentmark{15}{\char"1CD2}}पा\raisebox{0.25ex}{\accentmark{15}{\char"1CD2}}मूपास्थे\raisebox{0.25ex}{\accentmark{12}{\char"0951}}माहीषो\raisebox{0.25ex}{\accentmark{15}{\char"1CD2}}वावर्धा \nolinebreak\hskip 0.08em \mbox{॥९॥}}} \par}

\vspace{0.8em}
                % --- SUBSECTIONS (Mantras) ---
                
                % Conditional rendering based on output_mode
                \par\filbreak

\phantomsection

\addcontentsline{toc}{subsection}{ऋक् १०}

{\centering \footnotesize \textcolor{AccentPurple}{॥ वसिष्ठो मैत्रावरुणिः अग्निः त्रिपदाविराटगायत्री ॥} \par}

\vspace{0.12em}

{\centering \textbf{\textcolor{AccentBlue}{अग्निन्नारो\raisebox{0.25ex}{\accentmark{15}{\char"1CD2}}दीधितिभिरारण्यो\raisebox{0.25ex}{\accentmark{12}{\char"0951}}ः\kern 0.05em ।\hskip 0.12em plus 0.04em minus 0.02em ह\raisebox{0.25ex}{\accentmark{15}{\char"1CD2}}स्ता\raisebox{0.25ex}{\accentmark{12}{\char"0951}}च्युतञ्जनयतप्राशस्तम्\kern 0.05em ।\hskip 0.12em plus 0.04em minus 0.02em दूरेदृशङ्गृहा\raisebox{0.25ex}{\accentmark{15}{\char"1CD2}}पा\raisebox{0.25ex}{\accentmark{12}{\char"0951}}तीमाथव्युम् \nolinebreak\hskip 0.08em \mbox{॥१०॥}}} \par}

\vspace{0.8em}

            % --- End of Kandah Closing Floral Line & Counts ---
            \par\nopagebreak\vspace{0.9em plus 0.3em minus 0.1em}
            \centerline{\textcolor{black!35}{\rule{1.8cm}{0.35pt}}\;\raisebox{-0.22ex}{\textcolor{vedicmaroon}{\large\SixFlowerPetalDotted}}\;\textcolor{black!35}{\rule{1.8cm}{0.35pt}}}
            \nopagebreak\vspace{0.45em}
            \centerline{\devafont\textbf{\textcolor{AccentBlue}{॥ ऋक् सङ्ख्या : १० ॥}}}
        
            \clearpage

            % --- SECTION ---
            
            \phantomsection  

            \section*{\centering अष्टम खण्डः}
            \addcontentsline{toc}{section}{अष्टम खण्डः} 
            \markright{अष्टम खण्डः}
            \nopagebreak % Keep section title with first mantra
            
            % --- SECTION-LEVEL PROCEDURE LINK ---

                % --- SUBSECTIONS (Mantras) ---
                
                % Conditional rendering based on output_mode
                \par\filbreak

\phantomsection

\addcontentsline{toc}{subsection}{ऋक् १}

{\centering \footnotesize \textcolor{AccentPurple}{॥ बुधगविष्ठिरावात्रेयौ अग्निः त्रिष्टुप् ॥} \par}

\vspace{0.12em}

{\centering \textbf{\textcolor{AccentBlue}{आ\raisebox{0.25ex}{\accentmark{15}{\char"1CD2}}बो\raisebox{0.25ex}{\accentmark{12}{\char"0951}}ध्यग्निस्सा\raisebox{0.25ex}{\accentmark{15}{\char"1CD2}}मीधाज\raisebox{0.25ex}{\accentmark{15}{\char"1CD2}}ना\raisebox{0.25ex}{\accentmark{12}{\char"0951}}नां\kern 0.05em ।\hskip 0.12em plus 0.04em minus 0.02em प्रा\raisebox{0.15ex}{\accentmark{12}{\char"1CF8}}तीधेनू\raisebox{0.25ex}{\accentmark{15}{\char"1CD2}}मीवायातीमू\raisebox{0.25ex}{\accentmark{15}{\char"1CD2}}षा\raisebox{0.25ex}{\accentmark{15}{\char"1CD2}}सम्\kern 0.05em ।\hskip 0.12em plus 0.04em minus 0.02em यह्वा\raisebox{0.15ex}{\accentmark{12}{\char"1CF8}}इ\raisebox{0.15ex}{\accentmark{12}{\char"1CF8}}वप्रा\raisebox{0.25ex}{\accentmark{15}{\char"1CD2}}वाया\raisebox{0.25ex}{\accentmark{15}{\char"1CD2}}मूज्जीहा\raisebox{0.25ex}{\accentmark{12}{\char"0951}}नाः\kern 0.05em ।\hskip 0.12em plus 0.04em minus 0.02em प्रा\raisebox{0.15ex}{\accentmark{12}{\char"1CF8}}भान्नावासस्राते\raisebox{0.25ex}{\accentmark{15}{\char"1CD2}}नाकामाच्छा \nolinebreak\hskip 0.08em \mbox{॥१॥}}} \par}

\vspace{0.8em}
                % --- SUBSECTIONS (Mantras) ---
                
                % Conditional rendering based on output_mode
                \par\filbreak

\phantomsection

\addcontentsline{toc}{subsection}{ऋक् २}

{\centering \footnotesize \textcolor{AccentPurple}{॥ वत्सप्रिर्भालन्दनः अग्निः त्रिष्टुप् ॥} \par}

\vspace{0.12em}

{\centering \textbf{\textcolor{AccentBlue}{प्राभूर्जय\raisebox{0.25ex}{\accentmark{15}{\char"1CD2}}न्तम्माहीं\raisebox{0.15ex}{\accentmark{12}{\char"1CF9}}वीपोधाम्\kern 0.05em ।\hskip 0.12em plus 0.04em minus 0.02em मूरै\raisebox{0.25ex}{\accentmark{15}{\char"1CD2}}रा\raisebox{0.25ex}{\accentmark{15}{\char"1CD2}}मू\raisebox{0.25ex}{\accentmark{15}{\char"1CD2}}रंपूरा\raisebox{0.25ex}{\accentmark{15}{\char"1CD2}}न्द\raisebox{0.25ex}{\accentmark{12}{\char"0951}}र्माण\raisebox{0.25ex}{\accentmark{12}{\char"0951}}म्\kern 0.05em ।\hskip 0.12em plus 0.04em minus 0.02em ना\raisebox{0.25ex}{\accentmark{15}{\char"1CD2}}यान्तंगी\raisebox{0.25ex}{\accentmark{15}{\char"1CD2}}र्भिर्वाना\raisebox{0.25ex}{\accentmark{15}{\char"1CD2}}धिय\raisebox{0.25ex}{\accentmark{12}{\char"0951}}न्धाः\kern 0.05em ।\hskip 0.12em plus 0.04em minus 0.02em हारी\raisebox{0.25ex}{\accentmark{12}{\char"0951}}श्माश्रून्नावर्मणाधानार्चि\raisebox{0.25ex}{\accentmark{12}{\char"0951}}म् \nolinebreak\hskip 0.08em \mbox{॥२॥}}} \par}

\vspace{0.8em}
                % --- SUBSECTIONS (Mantras) ---
                
                % Conditional rendering based on output_mode
                \par\filbreak

\phantomsection

\addcontentsline{toc}{subsection}{ऋक् ३}

{\centering \footnotesize \textcolor{AccentPurple}{॥ भरद्वाजो बार्हस्पत्यः पूषा त्रिष्टुप् ॥} \par}

\vspace{0.12em}

{\centering \textbf{\textcolor{AccentBlue}{शु\raisebox{0.25ex}{\accentmark{15}{\char"1CD2}}क्रन्ते\raisebox{0.25ex}{\accentmark{12}{\char"0951}}अन्य\raisebox{0.25ex}{\accentmark{12}{\char"0951}}द्याजतन्ते\raisebox{0.25ex}{\accentmark{12}{\char"0951}}अन्यात्\kern 0.05em ।\hskip 0.12em plus 0.04em minus 0.02em वि\raisebox{0.25ex}{\accentmark{12}{\char"0951}}षूरूपेआ\raisebox{0.15ex}{\accentmark{12}{\char"1CF8}}हा\raisebox{0.25ex}{\accentmark{12}{\char"0951}}निद्यौरीवा\raisebox{0.25ex}{\accentmark{12}{\char"0951}}सी\kern 0.05em ।\hskip 0.12em plus 0.04em minus 0.02em विश्वा\raisebox{0.15ex}{\accentmark{12}{\char"1CF9}}हीमाया\raisebox{0.25ex}{\accentmark{15}{\char"1CD2}}आवासिस्वधावन्\kern 0.05em ।\hskip 0.12em plus 0.04em minus 0.02em भ\raisebox{0.25ex}{\accentmark{15}{\char"1CD2}}द्रा\raisebox{0.25ex}{\accentmark{15}{\char"1CD2}}तेपूषन्नीहारातीरस्तू \nolinebreak\hskip 0.08em \mbox{॥३॥}}} \par}

\vspace{0.8em}
                % --- SUBSECTIONS (Mantras) ---
                
                % Conditional rendering based on output_mode
                \par\filbreak

\phantomsection

\addcontentsline{toc}{subsection}{ऋक् ४}

{\centering \footnotesize \textcolor{AccentPurple}{॥ विश्वामित्रो गाथिनः अग्निः त्रिष्टुप् ॥} \par}

\vspace{0.12em}

{\centering \textbf{\textcolor{AccentBlue}{ई\raisebox{0.25ex}{\accentmark{12}{\char"0951}}ळामग्ने\raisebox{0.15ex}{\accentmark{12}{\char"1CF9}}पू\raisebox{0.25ex}{\accentmark{12}{\char"0951}}रूदं\raisebox{0.15ex}{\accentmark{12}{\char"1CF8}}सं\raisebox{0.15ex}{\accentmark{12}{\char"1CF8}}सा\raisebox{0.15ex}{\accentmark{12}{\char"1CF8}}निंगोः\kern 0.05em ।\hskip 0.12em plus 0.04em minus 0.02em श\raisebox{0.25ex}{\accentmark{15}{\char"1CD2}}श्वक्तामं\raisebox{0.25ex}{\accentmark{15}{\char"1CD2}}हा\raisebox{0.25ex}{\accentmark{15}{\char"1CD2}}वा\raisebox{0.25ex}{\accentmark{12}{\char"0951}}मानायासा\raisebox{0.25ex}{\accentmark{15}{\char"1CD2}}धा\raisebox{0.25ex}{\accentmark{12}{\char"0951}}ः\kern 0.05em ।\hskip 0.12em plus 0.04em minus 0.02em स्यानस्सूनुस्ताना\raisebox{0.25ex}{\accentmark{12}{\char"0951}}योर्वीजा\raisebox{0.25ex}{\accentmark{15}{\char"1CD2}}वा\kern 0.05em ।\hskip 0.12em plus 0.04em minus 0.02em अ\raisebox{0.25ex}{\accentmark{15}{\char"1CD2}}ग्नेसाते\raisebox{0.25ex}{\accentmark{12}{\char"0951}}सूमात्य\raisebox{0.15ex}{\accentmark{12}{\char"1CF9}}र्भूत्वस्मे \nolinebreak\hskip 0.08em \mbox{॥४॥}}} \par}

\vspace{0.8em}
                % --- SUBSECTIONS (Mantras) ---
                
                % Conditional rendering based on output_mode
                \par\filbreak

\phantomsection

\addcontentsline{toc}{subsection}{ऋक् ५}

{\centering \footnotesize \textcolor{AccentPurple}{॥ वत्सप्रिर्भालन्दनः अग्निः त्रिष्टुप् ॥} \par}

\vspace{0.12em}

{\centering \textbf{\textcolor{AccentBlue}{प्रहो\raisebox{0.15ex}{\accentmark{12}{\char"1CF8}}ता\raisebox{0.25ex}{\accentmark{12}{\char"0951}}जातो\raisebox{0.25ex}{\accentmark{12}{\char"0951}}माहा\raisebox{0.15ex}{\accentmark{12}{\char"1CF9}}न्ना\raisebox{0.25ex}{\accentmark{15}{\char"1CD2}}भोवीत्\kern 0.05em ।\hskip 0.12em plus 0.04em minus 0.02em नृषात्मा\raisebox{0.25ex}{\accentmark{12}{\char"0951}}सीददपां\raisebox{0.15ex}{\accentmark{12}{\char"1CF9}}वीवार्ते\raisebox{0.25ex}{\accentmark{12}{\char"0951}}\kern 0.05em ।\hskip 0.12em plus 0.04em minus 0.02em दधद्यो\raisebox{0.25ex}{\accentmark{12}{\char"0951}}धायीसू\raisebox{0.25ex}{\accentmark{12}{\char"0951}}ते\raisebox{0.25ex}{\accentmark{15}{\char"1CD2}}वा\raisebox{0.25ex}{\accentmark{15}{\char"1CD2}}यांसि\raisebox{0.25ex}{\accentmark{12}{\char"0951}}\kern 0.05em ।\hskip 0.12em plus 0.04em minus 0.02em यन्तावासूनीवीधाते\raisebox{0.15ex}{\accentmark{12}{\char"1CF9}}ता\raisebox{0.25ex}{\accentmark{12}{\char"0951}}नूपाः \nolinebreak\hskip 0.08em \mbox{॥५॥}}} \par}

\vspace{0.8em}
                % --- SUBSECTIONS (Mantras) ---
                
                % Conditional rendering based on output_mode
                \par\filbreak

\phantomsection

\addcontentsline{toc}{subsection}{ऋक् ६}

{\centering \textbf{\textcolor{AccentBlue}{॥ वसिष्ठोमैत्रावरुणिःअग्निःत्रिष्टुप् ॥ ॥ विश्वामित्रोगाथिनःअग्निःत्रिष्टुप् ॥ प्रासम्राजमा\raisebox{0.25ex}{\accentmark{15}{\char"1CD2}}सू\raisebox{0.25ex}{\accentmark{12}{\char"0951}}रस्यप्रा\raisebox{0.25ex}{\accentmark{12}{\char"0951}}शस्तम्\kern 0.05em ।\hskip 0.12em plus 0.04em minus 0.02em पुंसःकृष्टीना\raisebox{0.15ex}{\accentmark{12}{\char"1CF9}}मा\raisebox{0.25ex}{\accentmark{12}{\char"0951}}नूमाद्यास्या\raisebox{0.25ex}{\accentmark{12}{\char"0951}}\kern 0.05em ।\hskip 0.12em plus 0.04em minus 0.02em इन्द्रस्येवप्रा\raisebox{0.25ex}{\accentmark{15}{\char"1CD2}}तावा\raisebox{0.15ex}{\accentmark{12}{\char"1CF9}}सःकृता\raisebox{0.25ex}{\accentmark{15}{\char"1CD2}}नि\kern 0.05em ।\hskip 0.12em plus 0.04em minus 0.02em व\raisebox{0.25ex}{\accentmark{15}{\char"1CD2}}न्द\raisebox{0.15ex}{\accentmark{12}{\char"1CF9}}द्वीरावन्दमानाविवष्टु \nolinebreak\hskip 0.08em \mbox{॥६॥}}} \par}

\vspace{0.8em}
                % --- SUBSECTIONS (Mantras) ---
                
                % Conditional rendering based on output_mode
                \par\filbreak

\phantomsection

\addcontentsline{toc}{subsection}{ऋक् ७}

{\centering \textbf{\textcolor{AccentBlue}{आरण्यो\raisebox{0.25ex}{\accentmark{15}{\char"1CD2}}र्नीहीतोजातावे\raisebox{0.25ex}{\accentmark{12}{\char"0951}}दाः\kern 0.05em ।\hskip 0.12em plus 0.04em minus 0.02em ग\raisebox{0.15ex}{\accentmark{12}{\char"1CF8}}र्भा\raisebox{0.25ex}{\accentmark{12}{\char"0951}}यीवेत्सू\raisebox{0.15ex}{\accentmark{12}{\char"1CF9}}भृ\raisebox{0.25ex}{\accentmark{12}{\char"0951}}तोगर्भिणीभिः\kern 0.05em ।\hskip 0.12em plus 0.04em minus 0.02em दीवे\raisebox{0.15ex}{\accentmark{12}{\char"1CF9}}दीवाईश्वोजागृवा\raisebox{0.25ex}{\accentmark{15}{\char"1CD2}}त्मिः\kern 0.05em ।\hskip 0.12em plus 0.04em minus 0.02em हा\raisebox{0.25ex}{\accentmark{15}{\char"1CD2}}वि\raisebox{0.25ex}{\accentmark{15}{\char"1CD2}}ष्मा\raisebox{0.25ex}{\accentmark{12}{\char"0951}}द्धि\raisebox{0.15ex}{\accentmark{12}{\char"1CF9}}र्मानुष्ये\raisebox{0.25ex}{\accentmark{12}{\char"0951}}भीरग्नि\raisebox{0.25ex}{\accentmark{15}{\char"1CD2}}ः \nolinebreak\hskip 0.08em \mbox{॥७॥}}} \par}

\vspace{0.8em}
                % --- SUBSECTIONS (Mantras) ---
                
                % Conditional rendering based on output_mode
                \par\filbreak

\phantomsection

\addcontentsline{toc}{subsection}{ऋक् ८}

{\centering \footnotesize \textcolor{AccentPurple}{॥ पायुर्भारद्वाजः अग्निः त्रिष्टुप् ॥} \par}

\vspace{0.12em}

{\centering \textbf{\textcolor{AccentBlue}{सा\raisebox{0.25ex}{\accentmark{15}{\char"1CD2}}नाद\raisebox{0.25ex}{\accentmark{12}{\char"0951}}अग्नेमृणसि\raisebox{0.25ex}{\accentmark{12}{\char"0951}}यतूधाना\raisebox{0.25ex}{\accentmark{12}{\char"0951}}न्\kern 0.05em ।\hskip 0.12em plus 0.04em minus 0.02em न\raisebox{0.25ex}{\accentmark{15}{\char"1CD2}}त्वारक्षां\raisebox{0.15ex}{\accentmark{12}{\char"1CF9}}सीपृ\raisebox{0.25ex}{\accentmark{15}{\char"1CD2}}ता\raisebox{0.25ex}{\accentmark{12}{\char"0951}}नासुजिग्युः\kern 0.05em ।\hskip 0.12em plus 0.04em minus 0.02em आनू\raisebox{0.25ex}{\accentmark{12}{\char"0951}}दहासाहामूरान्काया\raisebox{0.25ex}{\accentmark{15}{\char"1CD2}}दाः\kern 0.05em ।\hskip 0.12em plus 0.04em minus 0.02em मा\raisebox{0.15ex}{\accentmark{12}{\char"1CF8}}ते\raisebox{0.25ex}{\accentmark{12}{\char"0951}}हेत्यामू\raisebox{0.25ex}{\accentmark{12}{\char"0951}}क्षातादै\raisebox{0.25ex}{\accentmark{15}{\char"1CD2}}व्या\raisebox{0.25ex}{\accentmark{12}{\char"0951}}याः \nolinebreak\hskip 0.08em \mbox{॥८॥}}} \par}

\vspace{0.8em}

            % --- End of Kandah Closing Floral Line & Counts ---
            \par\nopagebreak\vspace{0.9em plus 0.3em minus 0.1em}
            \centerline{\textcolor{black!35}{\rule{1.8cm}{0.35pt}}\;\raisebox{-0.22ex}{\textcolor{vedicmaroon}{\large\SixFlowerPetalDotted}}\;\textcolor{black!35}{\rule{1.8cm}{0.35pt}}}
            \nopagebreak\vspace{0.45em}
            \centerline{\devafont\textbf{\textcolor{AccentBlue}{॥ ऋक् सङ्ख्या : ८ ॥}}}
        
            \clearpage

            % --- SECTION ---
            
            \phantomsection  

            \section*{\centering नवम खण्डः}
            \addcontentsline{toc}{section}{नवम खण्डः} 
            \markright{नवम खण्डः}
            \nopagebreak % Keep section title with first mantra
            
            % --- SECTION-LEVEL PROCEDURE LINK ---

                % --- SUBSECTIONS (Mantras) ---
                
                % Conditional rendering based on output_mode
                \par\filbreak

\phantomsection

\addcontentsline{toc}{subsection}{ऋक् १}

{\centering \footnotesize \textcolor{AccentPurple}{॥ गय आत्रयेः अग्निः अनुष्टुप् ॥} \par}

\vspace{0.12em}

{\centering \textbf{\textcolor{AccentBlue}{अग्ना\raisebox{0.15ex}{\accentmark{12}{\char"1CF9}}ओजीष्ठामा\raisebox{0.25ex}{\accentmark{15}{\char"1CD2}}भा\raisebox{0.25ex}{\accentmark{12}{\char"0951}}रा\kern 0.05em ।\hskip 0.12em plus 0.04em minus 0.02em धुम्नामस्माभ्या\raisebox{0.25ex}{\accentmark{12}{\char"0951}}मधिगो\kern 0.05em ।\hskip 0.12em plus 0.04em minus 0.02em प्रा\raisebox{0.15ex}{\accentmark{12}{\char"1CF8}}नो\raisebox{0.15ex}{\accentmark{12}{\char"1CF8}}रा\raisebox{0.25ex}{\accentmark{12}{\char"0951}}येपा\raisebox{0.25ex}{\accentmark{15}{\char"1CD2}}नीयसे\kern 0.05em ।\hskip 0.12em plus 0.04em minus 0.02em रत्सीवा\raisebox{0.15ex}{\accentmark{12}{\char"1CF9}}जायापन्थाम् \nolinebreak\hskip 0.08em \mbox{॥१॥}}} \par}

\vspace{0.8em}
                % --- SUBSECTIONS (Mantras) ---
                
                % Conditional rendering based on output_mode
                \par\filbreak

\phantomsection

\addcontentsline{toc}{subsection}{ऋक् २}

{\centering \footnotesize \textcolor{AccentPurple}{॥ वामदेवः अग्निः अनुष्टुप् ॥} \par}

\vspace{0.12em}

{\centering \textbf{\textcolor{AccentBlue}{या\raisebox{0.15ex}{\accentmark{12}{\char"1CF8}}दीवीरो\raisebox{0.25ex}{\accentmark{12}{\char"0951}}आनूष्यात्\kern 0.05em ।\hskip 0.12em plus 0.04em minus 0.02em अ\raisebox{0.25ex}{\accentmark{15}{\char"1CD2}}ग्नी\raisebox{0.25ex}{\accentmark{15}{\char"1CD2}}मि\raisebox{0.25ex}{\accentmark{15}{\char"1CD2}}न्धीतामार्क्त्याः\kern 0.05em ।\hskip 0.12em plus 0.04em minus 0.02em आ\raisebox{0.25ex}{\accentmark{15}{\char"1CD2}}जुह्वद्धेर्व्यमानूषाक्\kern 0.05em ।\hskip 0.12em plus 0.04em minus 0.02em शर्माभक्षितादै\raisebox{0.25ex}{\accentmark{15}{\char"1CD2}}व्य\raisebox{0.25ex}{\accentmark{12}{\char"0951}}म् \nolinebreak\hskip 0.08em \mbox{॥२॥}}} \par}

\vspace{0.8em}
                % --- SUBSECTIONS (Mantras) ---
                
                % Conditional rendering based on output_mode
                \par\filbreak

\phantomsection

\addcontentsline{toc}{subsection}{ऋक् ३}

{\centering \footnotesize \textcolor{AccentPurple}{॥ भरद्वाजो बार्हस्पत्यः अग्निः अनुष्टुप् ॥} \par}

\vspace{0.12em}

{\centering \textbf{\textcolor{AccentBlue}{त्वे\raisebox{0.25ex}{\accentmark{15}{\char"1CD2}}ष\raisebox{0.15ex}{\accentmark{12}{\char"1CF9}}स्ते\raisebox{0.25ex}{\accentmark{12}{\char"0951}}धूमा\raisebox{0.25ex}{\accentmark{15}{\char"1CD2}}ऋण्वति\kern 0.05em ।\hskip 0.12em plus 0.04em minus 0.02em दिविसत्शुक्रआ\raisebox{0.25ex}{\accentmark{15}{\char"1CD2}}त\raisebox{0.25ex}{\accentmark{12}{\char"0951}}तः\kern 0.05em ।\hskip 0.12em plus 0.04em minus 0.02em सू\raisebox{0.25ex}{\accentmark{15}{\char"1CD2}}रोनहिद्यूतात्वं\raisebox{0.25ex}{\accentmark{15}{\char"1CD2}}\kern 0.05em ।\hskip 0.12em plus 0.04em minus 0.02em कृपापावा\raisebox{0.25ex}{\accentmark{15}{\char"1CD2}}कारो\raisebox{0.25ex}{\accentmark{15}{\char"1CD2}}चा\raisebox{0.25ex}{\accentmark{12}{\char"0951}}से \nolinebreak\hskip 0.08em \mbox{॥३॥}}} \par}

\vspace{0.8em}
                % --- SUBSECTIONS (Mantras) ---
                
                % Conditional rendering based on output_mode
                \par\filbreak

\phantomsection

\addcontentsline{toc}{subsection}{ऋक् ४}

{\centering \footnotesize \textcolor{AccentPurple}{॥ भरद्वाजो बार्हस्पत्यः अग्निः अनुष्टुप् ॥} \par}

\vspace{0.12em}

{\centering \textbf{\textcolor{AccentBlue}{त्वं\raisebox{0.25ex}{\accentmark{12}{\char"0951}}हीक्षैता\raisebox{0.25ex}{\accentmark{15}{\char"1CD2}}व\raisebox{0.25ex}{\accentmark{15}{\char"1CD2}}द्या\raisebox{0.25ex}{\accentmark{15}{\char"1CD2}}शाः\kern 0.05em ।\hskip 0.12em plus 0.04em minus 0.02em अ\raisebox{0.15ex}{\accentmark{12}{\char"1CF8}}ग्ने\raisebox{0.25ex}{\accentmark{12}{\char"0951}}मित्रो\raisebox{0.25ex}{\accentmark{15}{\char"1CD2}}नापा\raisebox{0.25ex}{\accentmark{15}{\char"1CD2}}त्या\raisebox{0.25ex}{\accentmark{12}{\char"0951}}से\kern 0.05em ।\hskip 0.12em plus 0.04em minus 0.02em त्वंवि\raisebox{0.25ex}{\accentmark{12}{\char"0951}}चर्षा\raisebox{0.25ex}{\accentmark{12}{\char"0951}}णेश्रा\raisebox{0.25ex}{\accentmark{15}{\char"1CD2}}वा\raisebox{0.25ex}{\accentmark{12}{\char"0951}}ः\kern 0.05em ।\hskip 0.12em plus 0.04em minus 0.02em वासो\raisebox{0.25ex}{\accentmark{12}{\char"0951}}पूष्टिन्ना\raisebox{0.25ex}{\accentmark{15}{\char"1CD2}}पूष्यसि \nolinebreak\hskip 0.08em \mbox{॥४॥}}} \par}

\vspace{0.8em}
                % --- SUBSECTIONS (Mantras) ---
                
                % Conditional rendering based on output_mode
                \par\filbreak

\phantomsection

\addcontentsline{toc}{subsection}{ऋक् ५}

{\centering \footnotesize \textcolor{AccentPurple}{॥ द्वितो भृगवुाहा आत्रयेः अग्निः अनुष्टुप् ॥} \par}

\vspace{0.12em}

{\centering \textbf{\textcolor{AccentBlue}{प्रा\raisebox{0.15ex}{\accentmark{12}{\char"1CF8}}ता\raisebox{0.25ex}{\accentmark{12}{\char"0951}}रग्निःपूरूप्रीय\raisebox{0.25ex}{\accentmark{15}{\char"1CD2}}ः\kern 0.05em ।\hskip 0.12em plus 0.04em minus 0.02em वीश\raisebox{0.15ex}{\accentmark{12}{\char"1CF9}}स्तावेता\raisebox{0.25ex}{\accentmark{15}{\char"1CD2}}तीथिः\kern 0.05em ।\hskip 0.12em plus 0.04em minus 0.02em विश्वेयास्मिन्ना\raisebox{0.25ex}{\accentmark{15}{\char"1CD2}}मार्ये\kern 0.05em ।\hskip 0.12em plus 0.04em minus 0.02em हव्यम्मर्ता\raisebox{0.25ex}{\accentmark{12}{\char"0951}}साइन्धाते \nolinebreak\hskip 0.08em \mbox{॥५॥}}} \par}

\vspace{0.8em}
                % --- SUBSECTIONS (Mantras) ---
                
                % Conditional rendering based on output_mode
                \par\filbreak

\phantomsection

\addcontentsline{toc}{subsection}{ऋक् ६}

{\centering \footnotesize \textcolor{AccentPurple}{॥ वस्युव आत्रयेाः अग्निः अनुष्टुप् ॥} \par}

\vspace{0.12em}

{\centering \textbf{\textcolor{AccentBlue}{य\raisebox{0.15ex}{\accentmark{12}{\char"1CF8}}द्वाही\raisebox{0.25ex}{\accentmark{12}{\char"0951}}ष्ठन्ता\raisebox{0.25ex}{\accentmark{12}{\char"0951}}दग्ना\raisebox{0.25ex}{\accentmark{15}{\char"1CD2}}ये\kern 0.05em ।\hskip 0.12em plus 0.04em minus 0.02em बृहा\raisebox{0.25ex}{\accentmark{15}{\char"1CD2}}द\raisebox{0.25ex}{\accentmark{12}{\char"0951}}र्चविभा\raisebox{0.25ex}{\accentmark{12}{\char"0951}}वसो\kern 0.05em ।\hskip 0.12em plus 0.04em minus 0.02em माही\raisebox{0.25ex}{\accentmark{12}{\char"0951}}षीवत्वा\raisebox{0.25ex}{\accentmark{15}{\char"1CD2}}द्रायिः\kern 0.05em ।\hskip 0.12em plus 0.04em minus 0.02em त्व\raisebox{0.15ex}{\accentmark{12}{\char"1CF8}}द्वा\raisebox{0.25ex}{\accentmark{12}{\char"0951}}जाऊ\raisebox{0.25ex}{\accentmark{15}{\char"1CD2}}दी\raisebox{0.25ex}{\accentmark{12}{\char"0951}}रते \nolinebreak\hskip 0.08em \mbox{॥६॥}}} \par}

\vspace{0.8em}
                % --- SUBSECTIONS (Mantras) ---
                
                % Conditional rendering based on output_mode
                \par\filbreak

\phantomsection

\addcontentsline{toc}{subsection}{ऋक् ७}

{\centering \footnotesize \textcolor{AccentPurple}{॥ गोपवन आत्रयेः अग्निः अनुष्टुप् ॥} \par}

\vspace{0.12em}

{\centering \textbf{\textcolor{AccentBlue}{वी\raisebox{0.25ex}{\accentmark{15}{\char"1CD2}}शोवी\raisebox{0.25ex}{\accentmark{15}{\char"1CD2}}शोवोआ\raisebox{0.25ex}{\accentmark{15}{\char"1CD2}}ती\raisebox{0.25ex}{\accentmark{12}{\char"0951}}थिम्\kern 0.05em ।\hskip 0.12em plus 0.04em minus 0.02em वज\raisebox{0.25ex}{\accentmark{15}{\char"1CD2}}यन्ताःपूरूप्रीयम्\kern 0.05em ।\hskip 0.12em plus 0.04em minus 0.02em अ\raisebox{0.25ex}{\accentmark{15}{\char"1CD2}}ग्निं\raisebox{0.15ex}{\accentmark{12}{\char"1CF9}}वो\raisebox{0.25ex}{\accentmark{12}{\char"0951}}दुर्यंवाचाः\kern 0.05em ।\hskip 0.12em plus 0.04em minus 0.02em स्तूषेशुषास्याम\raisebox{0.25ex}{\accentmark{15}{\char"1CD2}}न्मा\raisebox{0.25ex}{\accentmark{12}{\char"0951}}भिः \nolinebreak\hskip 0.08em \mbox{॥७॥}}} \par}

\vspace{0.8em}
                % --- SUBSECTIONS (Mantras) ---
                
                % Conditional rendering based on output_mode
                \par\filbreak

\phantomsection

\addcontentsline{toc}{subsection}{ऋक् ८}

{\centering \footnotesize \textcolor{AccentPurple}{॥ पूरुरात्रयेः अग्निः अनुष्टुप् ॥} \par}

\vspace{0.12em}

{\centering \textbf{\textcolor{AccentBlue}{बृ\raisebox{0.25ex}{\accentmark{15}{\char"1CD2}}ह\raisebox{0.15ex}{\accentmark{12}{\char"1CF9}}द्वा\raisebox{0.15ex}{\accentmark{12}{\char"1CF9}}योहीभा\raisebox{0.25ex}{\accentmark{15}{\char"1CD2}}ना\raisebox{0.25ex}{\accentmark{15}{\char"1CD2}}वे\kern 0.05em ।\hskip 0.12em plus 0.04em minus 0.02em अ\raisebox{0.15ex}{\accentmark{12}{\char"1CF8}}र्चा\raisebox{0.25ex}{\accentmark{12}{\char"0951}}देवा\raisebox{0.25ex}{\accentmark{15}{\char"1CD2}}याग्ना\raisebox{0.25ex}{\accentmark{15}{\char"1CD2}}ये\kern 0.05em ।\hskip 0.12em plus 0.04em minus 0.02em यम्मित्र\raisebox{0.25ex}{\accentmark{15}{\char"1CD2}}न्नप्रा\raisebox{0.25ex}{\accentmark{15}{\char"1CD2}}शा\raisebox{0.25ex}{\accentmark{12}{\char"0951}}स्तये\kern 0.05em ।\hskip 0.12em plus 0.04em minus 0.02em म\raisebox{0.25ex}{\accentmark{15}{\char"1CD2}}र्क्ता\raisebox{0.25ex}{\accentmark{12}{\char"0951}}सोदधी\raisebox{0.15ex}{\accentmark{12}{\char"1CF9}}रे\raisebox{0.25ex}{\accentmark{12}{\char"0951}}पूरः \nolinebreak\hskip 0.08em \mbox{॥८॥}}} \par}

\vspace{0.8em}
                % --- SUBSECTIONS (Mantras) ---
                
                % Conditional rendering based on output_mode
                \par\filbreak

\phantomsection

\addcontentsline{toc}{subsection}{ऋक् ९}

{\centering \footnotesize \textcolor{AccentPurple}{॥ गोपवन आत्रयेः अग्निः अनुष्टुप् ॥} \par}

\vspace{0.12em}

{\centering \textbf{\textcolor{AccentBlue}{आ\raisebox{0.25ex}{\accentmark{15}{\char"1CD2}}ग\raisebox{0.25ex}{\accentmark{12}{\char"0951}}न्मवृत्राह\raisebox{0.25ex}{\accentmark{15}{\char"1CD2}}न्तामम्\kern 0.05em ।\hskip 0.12em plus 0.04em minus 0.02em ज्ये\raisebox{0.15ex}{\accentmark{12}{\char"1CF8}}ष्ठा\raisebox{0.25ex}{\accentmark{12}{\char"0951}}मग्निमाना\raisebox{0.25ex}{\accentmark{15}{\char"1CD2}}वम्\kern 0.05em ।\hskip 0.12em plus 0.04em minus 0.02em यस्म\raisebox{0.25ex}{\accentmark{12}{\char"0951}}श्रूतार्वन्नाक्षे\kern 0.05em ।\hskip 0.12em plus 0.04em minus 0.02em बृहाद\raisebox{0.25ex}{\accentmark{12}{\char"0951}}नीका\raisebox{0.25ex}{\accentmark{12}{\char"0951}}इध्याते \nolinebreak\hskip 0.08em \mbox{॥९॥}}} \par}

\vspace{0.8em}
                % --- SUBSECTIONS (Mantras) ---
                
                % Conditional rendering based on output_mode
                \par\filbreak

\phantomsection

\addcontentsline{toc}{subsection}{ऋक् १०}

{\centering \footnotesize \textcolor{AccentPurple}{॥ वामदेवः, कश्यपो वा मारीचो, मनुर्वा वैवस्वत, उभौ वा अग्निः अनुष्टुप् ॥} \par}

\vspace{0.12em}

{\centering \textbf{\textcolor{AccentBlue}{जाता\raisebox{0.15ex}{\accentmark{12}{\char"1CF9}}ःपा\raisebox{0.25ex}{\accentmark{15}{\char"1CD2}}रे\raisebox{0.25ex}{\accentmark{12}{\char"0951}}णाध\raisebox{0.25ex}{\accentmark{15}{\char"1CD2}}र्माणा\kern 0.05em ।\hskip 0.12em plus 0.04em minus 0.02em यत्सा\raisebox{0.25ex}{\accentmark{15}{\char"1CD2}}वृ\raisebox{0.15ex}{\accentmark{12}{\char"1CF9}}त्भी\raisebox{0.25ex}{\accentmark{12}{\char"0951}}स्साहा\raisebox{0.25ex}{\accentmark{15}{\char"1CD2}}भू\raisebox{0.25ex}{\accentmark{12}{\char"0951}}वः\kern 0.05em ।\hskip 0.12em plus 0.04em minus 0.02em पी\raisebox{0.25ex}{\accentmark{15}{\char"1CD2}}ताय\raisebox{0.25ex}{\accentmark{15}{\char"1CD2}}त्काश्या\raisebox{0.25ex}{\accentmark{15}{\char"1CD2}}पा\raisebox{0.25ex}{\accentmark{12}{\char"0951}}स्याग्निः\kern 0.05em ।\hskip 0.12em plus 0.04em minus 0.02em श्रद्धा\raisebox{0.15ex}{\accentmark{12}{\char"1CF9}}माता\raisebox{0.25ex}{\accentmark{12}{\char"0951}}मानूःकाविः \nolinebreak\hskip 0.08em \mbox{॥१०॥}}} \par}

\vspace{0.8em}

            % --- End of Kandah Closing Floral Line & Counts ---
            \par\nopagebreak\vspace{0.9em plus 0.3em minus 0.1em}
            \centerline{\textcolor{black!35}{\rule{1.8cm}{0.35pt}}\;\raisebox{-0.22ex}{\textcolor{vedicmaroon}{\large\SixFlowerPetalDotted}}\;\textcolor{black!35}{\rule{1.8cm}{0.35pt}}}
            \nopagebreak\vspace{0.45em}
            \centerline{\devafont\textbf{\textcolor{AccentBlue}{॥ ऋक् सङ्ख्या : १० ॥}}}
        
            \clearpage

            % --- SECTION ---
            
            \phantomsection  

            \section*{\centering दशम खण्डः}
            \addcontentsline{toc}{section}{दशम खण्डः} 
            \markright{दशम खण्डः}
            \nopagebreak % Keep section title with first mantra
            
            % --- SECTION-LEVEL PROCEDURE LINK ---

                % --- SUBSECTIONS (Mantras) ---
                
                % Conditional rendering based on output_mode
                \par\filbreak

\phantomsection

\addcontentsline{toc}{subsection}{ऋक् १}

{\centering \footnotesize \textcolor{AccentPurple}{॥ अग्निस्तापसः विश्वेदेवाः अनुष्टुप् ॥} \par}

\vspace{0.12em}

{\centering \textbf{\textcolor{AccentBlue}{सोमंरा\raisebox{0.15ex}{\accentmark{12}{\char"1CF9}}जा\raisebox{0.25ex}{\accentmark{12}{\char"0951}}नंवा\raisebox{0.25ex}{\accentmark{15}{\char"1CD2}}रूणम्\kern 0.05em ।\hskip 0.12em plus 0.04em minus 0.02em अ\raisebox{0.25ex}{\accentmark{15}{\char"1CD2}}ग्नी\raisebox{0.25ex}{\accentmark{15}{\char"1CD2}}मन्वारा\raisebox{0.25ex}{\accentmark{12}{\char"0951}}भामहे\kern 0.05em ।\hskip 0.12em plus 0.04em minus 0.02em आ\raisebox{0.25ex}{\accentmark{15}{\char"1CD2}}दित्यंविष्णुंसू\raisebox{0.25ex}{\accentmark{15}{\char"1CD2}}र्य\raisebox{0.25ex}{\accentmark{12}{\char"0951}}म्\kern 0.05em ।\hskip 0.12em plus 0.04em minus 0.02em ब्र\raisebox{0.25ex}{\accentmark{15}{\char"1CD2}}ह्मा\raisebox{0.15ex}{\accentmark{12}{\char"1CF9}}ण\raisebox{0.25ex}{\accentmark{12}{\char"0951}}ञ्चाबृहस्पा\raisebox{0.25ex}{\accentmark{15}{\char"1CD2}}तिम् \nolinebreak\hskip 0.08em \mbox{॥१॥}}} \par}

\vspace{0.8em}
                % --- SUBSECTIONS (Mantras) ---
                
                % Conditional rendering based on output_mode
                \par\filbreak

\phantomsection

\addcontentsline{toc}{subsection}{ऋक् २}

{\centering \footnotesize \textcolor{AccentPurple}{॥ वामदेवो कश्यपः, असितो दवेलो वा अंगिराः अनुष्टुप् ॥} \par}

\vspace{0.12em}

{\centering \textbf{\textcolor{AccentBlue}{ईता\raisebox{0.25ex}{\accentmark{15}{\char"1CD2}}एता\raisebox{0.25ex}{\accentmark{15}{\char"1CD2}}ऊदारूहन्\kern 0.05em ।\hskip 0.12em plus 0.04em minus 0.02em दीव\raisebox{0.25ex}{\accentmark{15}{\char"1CD2}}ःपृष्ठान्यारू\raisebox{0.25ex}{\accentmark{12}{\char"0951}}हन्\kern 0.05em ।\hskip 0.12em plus 0.04em minus 0.02em प्राभूर्जयो\raisebox{0.25ex}{\accentmark{12}{\char"0951}}या\raisebox{0.15ex}{\accentmark{12}{\char"1CF9}}थापाथा\raisebox{0.25ex}{\accentmark{15}{\char"1CD2}}\kern 0.05em ।\hskip 0.12em plus 0.04em minus 0.02em उ\raisebox{0.25ex}{\accentmark{15}{\char"1CD2}}द्यामङ्गी\raisebox{0.25ex}{\accentmark{12}{\char"0951}}रसोययुः \nolinebreak\hskip 0.08em \mbox{॥२॥}}} \par}

\vspace{0.8em}
                % --- SUBSECTIONS (Mantras) ---
                
                % Conditional rendering based on output_mode
                \par\filbreak

\phantomsection

\addcontentsline{toc}{subsection}{ऋक् ३}

{\centering \footnotesize \textcolor{AccentPurple}{॥ वामदेवो कश्यपः, असितो दवेलो वा अग्निः अनुष्टुप् ॥} \par}

\vspace{0.12em}

{\centering \textbf{\textcolor{AccentBlue}{रा\raisebox{0.25ex}{\accentmark{15}{\char"1CD2}}ये\raisebox{0.25ex}{\accentmark{15}{\char"1CD2}}आग्नेमाहेत्वा\kern 0.05em ।\hskip 0.12em plus 0.04em minus 0.02em दा\raisebox{0.15ex}{\accentmark{12}{\char"1CF9}}ना\raisebox{0.25ex}{\accentmark{12}{\char"0951}}यासा\raisebox{0.25ex}{\accentmark{15}{\char"1CD2}}मी\raisebox{0.25ex}{\accentmark{12}{\char"0951}}धीमहि\kern 0.05em ।\hskip 0.12em plus 0.04em minus 0.02em ई\raisebox{0.15ex}{\accentmark{12}{\char"1CF8}}ळिष्वाहीमा\raisebox{0.25ex}{\accentmark{12}{\char"0951}}हेवृषन्\kern 0.05em ।\hskip 0.12em plus 0.04em minus 0.02em द्या\raisebox{0.15ex}{\accentmark{12}{\char"1CF8}}वा\raisebox{0.25ex}{\accentmark{12}{\char"0951}}होत्रा\raisebox{0.25ex}{\accentmark{12}{\char"0951}}या\raisebox{0.25ex}{\accentmark{12}{\char"0951}}पृथिवी\raisebox{0.25ex}{\accentmark{15}{\char"1CD2}} \nolinebreak\hskip 0.08em \mbox{॥३॥}}} \par}

\vspace{0.8em}
                % --- SUBSECTIONS (Mantras) ---
                
                % Conditional rendering based on output_mode
                \par\filbreak

\phantomsection

\addcontentsline{toc}{subsection}{ऋक् ४}

{\centering \footnotesize \textcolor{AccentPurple}{॥ सोमाहुतिभागर्वः अग्निः अनुष्टुप् ॥} \par}

\vspace{0.12em}

{\centering \textbf{\textcolor{AccentBlue}{दधन्वे\raisebox{0.25ex}{\accentmark{15}{\char"1CD2}}वाया\raisebox{0.25ex}{\accentmark{15}{\char"1CD2}}दीमानू\kern 0.05em ।\hskip 0.12em plus 0.04em minus 0.02em वो\raisebox{0.15ex}{\accentmark{12}{\char"1CF9}}चा\raisebox{0.25ex}{\accentmark{12}{\char"0951}}त्ब्रह्मे\raisebox{0.25ex}{\accentmark{15}{\char"1CD2}}तीवे\raisebox{0.25ex}{\accentmark{15}{\char"1CD2}}रूतात्\kern 0.05em ।\hskip 0.12em plus 0.04em minus 0.02em पा\raisebox{0.25ex}{\accentmark{15}{\char"1CD2}}रीवि\raisebox{0.15ex}{\accentmark{12}{\char"1CF9}}श्वा\raisebox{0.25ex}{\accentmark{12}{\char"0951}}नीका\raisebox{0.25ex}{\accentmark{15}{\char"1CD2}}व्या\raisebox{0.25ex}{\accentmark{12}{\char"0951}}\kern 0.05em ।\hskip 0.12em plus 0.04em minus 0.02em ने\raisebox{0.25ex}{\accentmark{15}{\char"1CD2}}मीश्चक्रा\raisebox{0.25ex}{\accentmark{15}{\char"1CD2}}मी\raisebox{0.25ex}{\accentmark{12}{\char"0951}}वाभुवात् \nolinebreak\hskip 0.08em \mbox{॥४॥}}} \par}

\vspace{0.8em}
                % --- SUBSECTIONS (Mantras) ---
                
                % Conditional rendering based on output_mode
                \par\filbreak

\phantomsection

\addcontentsline{toc}{subsection}{ऋक् ५}

{\centering \footnotesize \textcolor{AccentPurple}{॥ पायुभारद्वाजः अग्निः अनुष्टुप् ॥} \par}

\vspace{0.12em}

{\centering \textbf{\textcolor{AccentBlue}{त्वामा\raisebox{0.25ex}{\accentmark{12}{\char"0951}}ग्नेवा\raisebox{0.15ex}{\accentmark{12}{\char"1CF9}}सुंरिहा\kern 0.05em ।\hskip 0.12em plus 0.04em minus 0.02em रूद्रं\raisebox{0.15ex}{\accentmark{12}{\char"1CF9}}आ\raisebox{0.25ex}{\accentmark{12}{\char"0951}}दित्यंऊता\kern 0.05em ।\hskip 0.12em plus 0.04em minus 0.02em याजस्वाध्वार\raisebox{0.25ex}{\accentmark{15}{\char"1CD2}}ञ्ज\raisebox{0.25ex}{\accentmark{15}{\char"1CD2}}न\raisebox{0.25ex}{\accentmark{12}{\char"0951}}म्\kern 0.05em ।\hskip 0.12em plus 0.04em minus 0.02em आ\raisebox{0.25ex}{\accentmark{15}{\char"1CD2}}नूजातङ्घृ\raisebox{0.15ex}{\accentmark{12}{\char"1CF9}}तप्रू\raisebox{0.25ex}{\accentmark{15}{\char"1CD2}}षम् \nolinebreak\hskip 0.08em \mbox{॥५॥}}} \par}

\vspace{0.8em}
                % --- SUBSECTIONS (Mantras) ---
                
                % Conditional rendering based on output_mode
                \par\filbreak

\phantomsection

\addcontentsline{toc}{subsection}{ऋक् ६}

{\centering \footnotesize \textcolor{AccentPurple}{॥ प्रस्कण्वः काण्वः अग्निः अनुष्टुप् ॥} \par}

\vspace{0.12em}

{\centering \textbf{\textcolor{AccentBlue}{प्रत्य\raisebox{0.25ex}{\accentmark{12}{\char"0951}}ग्नेहा\raisebox{0.15ex}{\accentmark{12}{\char"1CF9}}रा\raisebox{0.25ex}{\accentmark{12}{\char"0951}}साहा\raisebox{0.25ex}{\accentmark{15}{\char"1CD2}}राः\kern 0.05em ।\hskip 0.12em plus 0.04em minus 0.02em शृणा\raisebox{0.25ex}{\accentmark{15}{\char"1CD2}}हीविश्वा\raisebox{0.25ex}{\accentmark{15}{\char"1CD2}}तप्रा\raisebox{0.25ex}{\accentmark{15}{\char"1CD2}}ती\kern 0.05em ।\hskip 0.12em plus 0.04em minus 0.02em या\raisebox{0.25ex}{\accentmark{15}{\char"1CD2}}तूधा\raisebox{0.25ex}{\accentmark{15}{\char"1CD2}}ना\raisebox{0.25ex}{\accentmark{12}{\char"0951}}स्या\raisebox{0.15ex}{\accentmark{12}{\char"1CF8}}रक्षासा\raisebox{0.25ex}{\accentmark{12}{\char"0951}}ः\kern 0.05em ।\hskip 0.12em plus 0.04em minus 0.02em बलंन्युब्जावीर्य\raisebox{0.25ex}{\accentmark{12}{\char"0951}}म् \nolinebreak\hskip 0.08em \mbox{॥६॥}}} \par}

\vspace{0.8em}

            % --- End of Kandah Closing Floral Line & Counts ---
            \par\nopagebreak\vspace{0.9em plus 0.3em minus 0.1em}
            \centerline{\textcolor{black!35}{\rule{1.8cm}{0.35pt}}\;\raisebox{-0.22ex}{\textcolor{vedicmaroon}{\large\SixFlowerPetalDotted}}\;\textcolor{black!35}{\rule{1.8cm}{0.35pt}}}
            \nopagebreak\vspace{0.45em}
            \centerline{\devafont\textbf{\textcolor{AccentBlue}{॥ ऋक् सङ्ख्या : ६ ॥}}}
        
            \clearpage

            % --- SECTION ---
            
            \phantomsection  

            \section*{\centering एकादश खण्डः}
            \addcontentsline{toc}{section}{एकादश खण्डः} 
            \markright{एकादश खण्डः}
            \nopagebreak % Keep section title with first mantra
            
            % --- SECTION-LEVEL PROCEDURE LINK ---

                % --- SUBSECTIONS (Mantras) ---
                
                % Conditional rendering based on output_mode
                \par\filbreak

\phantomsection

\addcontentsline{toc}{subsection}{ऋक् १}

{\centering \footnotesize \textcolor{AccentPurple}{॥ दीर्घतमा औचथ्यः अग्निः उष्णिक् ॥} \par}

\vspace{0.12em}

{\centering \textbf{\textcolor{AccentBlue}{पूरू\raisebox{0.25ex}{\accentmark{15}{\char"1CD2}}त्वा\raisebox{0.25ex}{\accentmark{12}{\char"0951}}दाशीवं\raisebox{0.25ex}{\accentmark{15}{\char"1CD2}}वो\raisebox{0.25ex}{\accentmark{12}{\char"0951}}चे\raisebox{0.15ex}{\accentmark{12}{\char"1CF8}}त्\kern 0.05em ।\hskip 0.12em plus 0.04em minus 0.02em आरीरग्नेता\raisebox{0.15ex}{\accentmark{12}{\char"1CF9}}वास्वीदा\raisebox{0.25ex}{\accentmark{15}{\char"1CD2}}\kern 0.05em ।\hskip 0.12em plus 0.04em minus 0.02em तोद\raisebox{0.25ex}{\accentmark{15}{\char"1CD2}}स्ये\raisebox{0.25ex}{\accentmark{15}{\char"1CD2}}वशारणआ\raisebox{0.25ex}{\accentmark{12}{\char"0951}}माहास्या\raisebox{0.25ex}{\accentmark{12}{\char"0951}} \nolinebreak\hskip 0.08em \mbox{॥१॥}}} \par}

\vspace{0.8em}
                % --- SUBSECTIONS (Mantras) ---
                
                % Conditional rendering based on output_mode
                \par\filbreak

\phantomsection

\addcontentsline{toc}{subsection}{ऋक् २}

{\centering \footnotesize \textcolor{AccentPurple}{॥ विश्वामित्रो गाथिनः अग्निः उष्णिक् ॥} \par}

\vspace{0.12em}

{\centering \textbf{\textcolor{AccentBlue}{प्र\raisebox{0.15ex}{\accentmark{12}{\char"1CF8}}होत्रे\raisebox{0.25ex}{\accentmark{12}{\char"0951}}पूर्वीयंवाचा\raisebox{0.25ex}{\accentmark{12}{\char"0951}}ः\kern 0.05em ।\hskip 0.12em plus 0.04em minus 0.02em अग्नाये\raisebox{0.25ex}{\accentmark{12}{\char"0951}}भरताबृहा\raisebox{0.25ex}{\accentmark{15}{\char"1CD2}}त्\kern 0.05em ।\hskip 0.12em plus 0.04em minus 0.02em वीपांज्योती\raisebox{0.25ex}{\accentmark{12}{\char"0951}}षिबि\raisebox{0.15ex}{\accentmark{12}{\char"1CF9}}भ्रा\raisebox{0.25ex}{\accentmark{12}{\char"0951}}तेना\raisebox{0.25ex}{\accentmark{15}{\char"1CD2}}वेधासे\raisebox{0.25ex}{\accentmark{12}{\char"0951}} \nolinebreak\hskip 0.08em \mbox{॥२॥}}} \par}

\vspace{0.8em}
                % --- SUBSECTIONS (Mantras) ---
                
                % Conditional rendering based on output_mode
                \par\filbreak

\phantomsection

\addcontentsline{toc}{subsection}{ऋक् ३}

{\centering \footnotesize \textcolor{AccentPurple}{॥ गौतमो राहूगणः अग्निः उष्णिक् ॥} \par}

\vspace{0.12em}

{\centering \textbf{\textcolor{AccentBlue}{अग्नेवाज\raisebox{0.25ex}{\accentmark{12}{\char"0951}}स्यागोमा\raisebox{0.25ex}{\accentmark{12}{\char"0951}}तः\kern 0.05em ।\hskip 0.12em plus 0.04em minus 0.02em ई\raisebox{0.25ex}{\accentmark{15}{\char"1CD2}}शा\raisebox{0.25ex}{\accentmark{12}{\char"0951}}नस्सहसोयहो\kern 0.05em ।\hskip 0.12em plus 0.04em minus 0.02em अ\raisebox{0.15ex}{\accentmark{12}{\char"1CF8}}स्मे\raisebox{0.25ex}{\accentmark{15}{\char"1CD2}}देहीजातवेदोमा\raisebox{0.25ex}{\accentmark{15}{\char"1CD2}}हिश्रावाः \nolinebreak\hskip 0.08em \mbox{॥३॥}}} \par}

\vspace{0.8em}
                % --- SUBSECTIONS (Mantras) ---
                
                % Conditional rendering based on output_mode
                \par\filbreak

\phantomsection

\addcontentsline{toc}{subsection}{ऋक् ४}

{\centering \footnotesize \textcolor{AccentPurple}{॥ विश्वामित्रो गाथिनः अग्निः उष्णिक् ॥} \par}

\vspace{0.12em}

{\centering \textbf{\textcolor{AccentBlue}{अ\raisebox{0.15ex}{\accentmark{12}{\char"1CF8}}ग्ने\raisebox{0.25ex}{\accentmark{15}{\char"1CD2}}या\raisebox{0.25ex}{\accentmark{15}{\char"1CD2}}जीष्ठोआध्वारे\kern 0.05em ।\hskip 0.12em plus 0.04em minus 0.02em देवान्दे\raisebox{0.25ex}{\accentmark{12}{\char"0951}}वायातेयाजा\kern 0.05em ।\hskip 0.12em plus 0.04em minus 0.02em होता\raisebox{0.25ex}{\accentmark{12}{\char"0951}}मन्द्रो\raisebox{0.15ex}{\accentmark{12}{\char"1CF9}}वी\raisebox{0.25ex}{\accentmark{12}{\char"0951}}राज\raisebox{0.25ex}{\accentmark{15}{\char"1CD2}}स्यातिसृधाः \nolinebreak\hskip 0.08em \mbox{॥४॥}}} \par}

\vspace{0.8em}
                % --- SUBSECTIONS (Mantras) ---
                
                % Conditional rendering based on output_mode
                \par\filbreak

\phantomsection

\addcontentsline{toc}{subsection}{ऋक् ५}

{\centering \footnotesize \textcolor{AccentPurple}{॥ त्रित आत्यः पवमानः सोमः उष्णिक् ॥} \par}

\vspace{0.12em}

{\centering \textbf{\textcolor{AccentBlue}{जज्ञान\raisebox{0.25ex}{\accentmark{15}{\char"1CD2}}स्सप्ता\raisebox{0.25ex}{\accentmark{15}{\char"1CD2}}मातृ\raisebox{0.25ex}{\accentmark{15}{\char"1CD2}}भिः\kern 0.05em ।\hskip 0.12em plus 0.04em minus 0.02em मे\raisebox{0.25ex}{\accentmark{15}{\char"1CD2}}धामाशा\raisebox{0.25ex}{\accentmark{12}{\char"0951}}सातश्रिये\raisebox{0.25ex}{\accentmark{15}{\char"1CD2}}\kern 0.05em ।\hskip 0.12em plus 0.04em minus 0.02em आ\raisebox{0.15ex}{\accentmark{12}{\char"1CF8}}यन्ध्रुवोरायीणा\raisebox{0.15ex}{\accentmark{12}{\char"1CF9}}ञ्चीके\raisebox{0.25ex}{\accentmark{12}{\char"0951}}तादा \nolinebreak\hskip 0.08em \mbox{॥५॥}}} \par}

\vspace{0.8em}
                % --- SUBSECTIONS (Mantras) ---
                
                % Conditional rendering based on output_mode
                \par\filbreak

\phantomsection

\addcontentsline{toc}{subsection}{ऋक् ६}

{\centering \footnotesize \textcolor{AccentPurple}{॥ इरिम्बिठः काण्वः अदितिः उष्णिक् ॥} \par}

\vspace{0.12em}

{\centering \textbf{\textcolor{AccentBlue}{ऊत\raisebox{0.25ex}{\accentmark{15}{\char"1CD2}}स्या\raisebox{0.25ex}{\accentmark{15}{\char"1CD2}}नोदीवामातिः\kern 0.05em ।\hskip 0.12em plus 0.04em minus 0.02em आदी\raisebox{0.25ex}{\accentmark{12}{\char"0951}}तीरूत्यगमात्\kern 0.05em ।\hskip 0.12em plus 0.04em minus 0.02em साश\raisebox{0.15ex}{\accentmark{12}{\char"1CF8}}न्तातामायास्कारादपश्री\raisebox{0.25ex}{\accentmark{15}{\char"1CD2}}धाः \nolinebreak\hskip 0.08em \mbox{॥६॥}}} \par}

\vspace{0.8em}
                % --- SUBSECTIONS (Mantras) ---
                
                % Conditional rendering based on output_mode
                \par\filbreak

\phantomsection

\addcontentsline{toc}{subsection}{ऋक् ७}

{\centering \footnotesize \textcolor{AccentPurple}{॥ विश्वमना वैयश्वः अग्निः उष्णिक् ॥} \par}

\vspace{0.12em}

{\centering \textbf{\textcolor{AccentBlue}{ईळि\raisebox{0.25ex}{\accentmark{12}{\char"0951}}ष्वाहि\raisebox{0.15ex}{\accentmark{12}{\char"1CF9}}प्रा\raisebox{0.25ex}{\accentmark{12}{\char"0951}}तिव्यम्\raisebox{0.25ex}{\accentmark{12}{\char"0951}}\kern 0.05em ।\hskip 0.12em plus 0.04em minus 0.02em या\raisebox{0.25ex}{\accentmark{15}{\char"1CD2}}ज\raisebox{0.25ex}{\accentmark{12}{\char"0951}}स्वाजातावे\raisebox{0.25ex}{\accentmark{12}{\char"0951}}दसम्\kern 0.05em ।\hskip 0.12em plus 0.04em minus 0.02em चरिष्णू\raisebox{0.25ex}{\accentmark{12}{\char"0951}}धूमामागृभीतशोचिषम् \nolinebreak\hskip 0.08em \mbox{॥७॥}}} \par}

\vspace{0.8em}
                % --- SUBSECTIONS (Mantras) ---
                
                % Conditional rendering based on output_mode
                \par\filbreak

\phantomsection

\addcontentsline{toc}{subsection}{ऋक् ८}

{\centering \footnotesize \textcolor{AccentPurple}{॥ विश्वमना वैयश्वः अग्निः उष्णिक् ॥} \par}

\vspace{0.12em}

{\centering \textbf{\textcolor{AccentBlue}{नत\raisebox{0.25ex}{\accentmark{15}{\char"1CD2}}स्यामाया\raisebox{0.15ex}{\accentmark{12}{\char"1CF9}}याचाना\raisebox{0.25ex}{\accentmark{15}{\char"1CD2}}\kern 0.05em ।\hskip 0.12em plus 0.04em minus 0.02em री\raisebox{0.25ex}{\accentmark{15}{\char"1CD2}}पू\raisebox{0.25ex}{\accentmark{15}{\char"1CD2}}रीशीतामाक्र्त्याः\kern 0.05em ।\hskip 0.12em plus 0.04em minus 0.02em यो\raisebox{0.25ex}{\accentmark{15}{\char"1CD2}}अग्ना\raisebox{0.15ex}{\accentmark{12}{\char"1CF9}}ये\raisebox{0.25ex}{\accentmark{12}{\char"0951}}ददा\raisebox{0.15ex}{\accentmark{12}{\char"1CF9}}शाहव्या\raisebox{0.25ex}{\accentmark{15}{\char"1CD2}}दा\raisebox{0.25ex}{\accentmark{12}{\char"0951}}तये \nolinebreak\hskip 0.08em \mbox{॥८॥}}} \par}

\vspace{0.8em}
                % --- SUBSECTIONS (Mantras) ---
                
                % Conditional rendering based on output_mode
                \par\filbreak

\phantomsection

\addcontentsline{toc}{subsection}{ऋक् ९}

{\centering \footnotesize \textcolor{AccentPurple}{॥ ऋजिश्वा भारद्वाजः विश्वेदेवाः उष्णिक् ॥} \par}

\vspace{0.12em}

{\centering \textbf{\textcolor{AccentBlue}{अपत्यं\raisebox{0.15ex}{\accentmark{12}{\char"1CF9}}वृ\raisebox{0.25ex}{\accentmark{12}{\char"0951}}जीनंरिपु\raisebox{0.25ex}{\accentmark{15}{\char"1CD2}}म्\kern 0.05em ।\hskip 0.12em plus 0.04em minus 0.02em स्ते\raisebox{0.25ex}{\accentmark{15}{\char"1CD2}}ना\raisebox{0.25ex}{\accentmark{15}{\char"1CD2}}मा\raisebox{0.25ex}{\accentmark{12}{\char"0951}}ग्नेदूराध्यम्\kern 0.05em ।\hskip 0.12em plus 0.04em minus 0.02em दे\raisebox{0.25ex}{\accentmark{15}{\char"1CD2}}वीष्ठमत्स्यस\raisebox{0.25ex}{\accentmark{12}{\char"0951}}त्पतेकृधीसूगम् \nolinebreak\hskip 0.08em \mbox{॥९॥}}} \par}

\vspace{0.8em}
                % --- SUBSECTIONS (Mantras) ---
                
                % Conditional rendering based on output_mode
                \par\filbreak

\phantomsection

\addcontentsline{toc}{subsection}{ऋक् १०}

{\centering \footnotesize \textcolor{AccentPurple}{॥ विश्वमना वैयश्वः अग्निः उष्णिक् ॥} \par}

\vspace{0.12em}

{\centering \textbf{\textcolor{AccentBlue}{श्रु\raisebox{0.25ex}{\accentmark{15}{\char"1CD2}}ष्ट्या\raisebox{0.25ex}{\accentmark{15}{\char"1CD2}}ग्नेना\raisebox{0.25ex}{\accentmark{15}{\char"1CD2}}वा\raisebox{0.25ex}{\accentmark{12}{\char"0951}}स्यमे\kern 0.05em ।\hskip 0.12em plus 0.04em minus 0.02em स्तोमा\raisebox{0.25ex}{\accentmark{12}{\char"0951}}स्यवीरविस्पते\kern 0.05em ।\hskip 0.12em plus 0.04em minus 0.02em नीमा\raisebox{0.25ex}{\accentmark{15}{\char"1CD2}}यिनस्तपा\raisebox{0.25ex}{\accentmark{15}{\char"1CD2}}सारक्षा\raisebox{0.25ex}{\accentmark{15}{\char"1CD2}}सो\raisebox{0.25ex}{\accentmark{12}{\char"0951}}दह \nolinebreak\hskip 0.08em \mbox{॥१०॥}}} \par}

\vspace{0.8em}

            % --- End of Kandah Closing Floral Line & Counts ---
            \par\nopagebreak\vspace{0.9em plus 0.3em minus 0.1em}
            \centerline{\textcolor{black!35}{\rule{1.8cm}{0.35pt}}\;\raisebox{-0.22ex}{\textcolor{vedicmaroon}{\large\SixFlowerPetalDotted}}\;\textcolor{black!35}{\rule{1.8cm}{0.35pt}}}
            \nopagebreak\vspace{0.45em}
            \centerline{\devafont\textbf{\textcolor{AccentBlue}{॥ ऋक् सङ्ख्या : १० ॥}}}
        
            \clearpage

            % --- SECTION ---
            
            \phantomsection  

            \section*{\centering द्वादश खण्डः}
            \addcontentsline{toc}{section}{द्वादश खण्डः} 
            \markright{द्वादश खण्डः}
            \nopagebreak % Keep section title with first mantra
            
            % --- SECTION-LEVEL PROCEDURE LINK ---

                % --- SUBSECTIONS (Mantras) ---
                
                % Conditional rendering based on output_mode
                \par\filbreak

\phantomsection

\addcontentsline{toc}{subsection}{ऋक् १}

{\centering \footnotesize \textcolor{AccentPurple}{॥ प्रयोगो भागर्वः (सोभिरः काण्वः) अग्निः उष्णिक् ॥} \par}

\vspace{0.12em}

{\centering \textbf{\textcolor{AccentBlue}{प्रमंही\raisebox{0.25ex}{\accentmark{12}{\char"0951}}ष्ठायगायता\kern 0.05em ।\hskip 0.12em plus 0.04em minus 0.02em ऋताव्नेबृहाते\raisebox{0.25ex}{\accentmark{15}{\char"1CD2}}शुक्राशो\raisebox{0.25ex}{\accentmark{12}{\char"0951}}चिषे\kern 0.05em ।\hskip 0.12em plus 0.04em minus 0.02em उ\raisebox{0.25ex}{\accentmark{15}{\char"1CD2}}प\raisebox{0.25ex}{\accentmark{15}{\char"1CD2}}स्तूता\raisebox{0.15ex}{\accentmark{12}{\char"1CF9}}सो\raisebox{0.25ex}{\accentmark{12}{\char"0951}}अग्नाये \nolinebreak\hskip 0.08em \mbox{॥१॥}}} \par}

\vspace{0.8em}
                % --- SUBSECTIONS (Mantras) ---
                
                % Conditional rendering based on output_mode
                \par\filbreak

\phantomsection

\addcontentsline{toc}{subsection}{ऋक् २}

{\centering \footnotesize \textcolor{AccentPurple}{॥ सोभिरः काण्वः अग्निः उष्णिक् ॥} \par}

\vspace{0.12em}

{\centering \textbf{\textcolor{AccentBlue}{प्र\raisebox{0.15ex}{\accentmark{12}{\char"1CF8}}सोअग्नेतावोतीभीः\kern 0.05em ।\hskip 0.12em plus 0.04em minus 0.02em सू\raisebox{0.25ex}{\accentmark{15}{\char"1CD2}}वीरा\raisebox{0.25ex}{\accentmark{12}{\char"0951}}भिस्तारातीवा\raisebox{0.25ex}{\accentmark{15}{\char"1CD2}}जकर्म\raisebox{0.25ex}{\accentmark{12}{\char"0951}}भिः\kern 0.05em ।\hskip 0.12em plus 0.04em minus 0.02em या\raisebox{0.25ex}{\accentmark{15}{\char"1CD2}}स्यत्वं\raisebox{0.25ex}{\accentmark{15}{\char"1CD2}}सख्यमा\raisebox{0.25ex}{\accentmark{15}{\char"1CD2}}वी\raisebox{0.25ex}{\accentmark{15}{\char"1CD2}}थ \nolinebreak\hskip 0.08em \mbox{॥२॥}}} \par}

\vspace{0.8em}
                % --- SUBSECTIONS (Mantras) ---
                
                % Conditional rendering based on output_mode
                \par\filbreak

\phantomsection

\addcontentsline{toc}{subsection}{ऋक् ३}

{\centering \footnotesize \textcolor{AccentPurple}{॥ सोभिरः काण्वः अग्निः उष्णिक् ॥} \par}

\vspace{0.12em}

{\centering \textbf{\textcolor{AccentBlue}{तं\raisebox{0.25ex}{\accentmark{15}{\char"1CD2}}गूर्धायास्व\raisebox{0.25ex}{\accentmark{15}{\char"1CD2}}र्णारम्\kern 0.05em ।\hskip 0.12em plus 0.04em minus 0.02em देवासो\raisebox{0.25ex}{\accentmark{12}{\char"0951}}देवा\raisebox{0.15ex}{\accentmark{12}{\char"1CF9}}मा\raisebox{0.25ex}{\accentmark{12}{\char"0951}}रातिन्द\raisebox{0.25ex}{\accentmark{12}{\char"0951}}धन्विरे\kern 0.05em ।\hskip 0.12em plus 0.04em minus 0.02em देवात्रा\raisebox{0.25ex}{\accentmark{15}{\char"1CD2}}हव्या\raisebox{0.25ex}{\accentmark{15}{\char"1CD2}}मू\raisebox{0.25ex}{\accentmark{12}{\char"0951}}हिषे \nolinebreak\hskip 0.08em \mbox{॥३॥}}} \par}

\vspace{0.8em}
                % --- SUBSECTIONS (Mantras) ---
                
                % Conditional rendering based on output_mode
                \par\filbreak

\phantomsection

\addcontentsline{toc}{subsection}{ऋक् ४}

{\centering \footnotesize \textcolor{AccentPurple}{॥ प्रयोगो भागर्वः सोभिरः काण्वो वा अग्निः उष्णिक् ॥} \par}

\vspace{0.12em}

{\centering \textbf{\textcolor{AccentBlue}{मानो\raisebox{0.25ex}{\accentmark{12}{\char"0951}}हृ\raisebox{0.25ex}{\accentmark{12}{\char"0951}}णिथाआ\raisebox{0.25ex}{\accentmark{15}{\char"1CD2}}तीथिम्\kern 0.05em ।\hskip 0.12em plus 0.04em minus 0.02em वा\raisebox{0.15ex}{\accentmark{12}{\char"1CF9}}सूरग्नि\raisebox{0.25ex}{\accentmark{15}{\char"1CD2}}ःपू\raisebox{0.25ex}{\accentmark{12}{\char"0951}}रूप्राशस्ता\raisebox{0.25ex}{\accentmark{15}{\char"1CD2}}एषः\kern 0.05em ।\hskip 0.12em plus 0.04em minus 0.02em या\raisebox{0.25ex}{\accentmark{15}{\char"1CD2}}सुहोता\raisebox{0.25ex}{\accentmark{12}{\char"0951}}सूआध्वारः \nolinebreak\hskip 0.08em \mbox{॥४॥}}} \par}

\vspace{0.8em}
                % --- SUBSECTIONS (Mantras) ---
                
                % Conditional rendering based on output_mode
                \par\filbreak

\phantomsection

\addcontentsline{toc}{subsection}{ऋक् ५}

{\centering \footnotesize \textcolor{AccentPurple}{॥ सोभिरः काण्वः अग्निः उष्णिक् ॥} \par}

\vspace{0.12em}

{\centering \textbf{\textcolor{AccentBlue}{भद्रो\raisebox{0.15ex}{\accentmark{12}{\char"1CF9}}नो\raisebox{0.25ex}{\accentmark{12}{\char"0951}}अग्नि\raisebox{0.25ex}{\accentmark{15}{\char"1CD2}}राहूतः\kern 0.05em ।\hskip 0.12em plus 0.04em minus 0.02em भ\raisebox{0.25ex}{\accentmark{15}{\char"1CD2}}द्रा\raisebox{0.25ex}{\accentmark{15}{\char"1CD2}}रातिस्सू\raisebox{0.25ex}{\accentmark{15}{\char"1CD2}}भागभद्रो\raisebox{0.15ex}{\accentmark{12}{\char"1CF9}}आ\raisebox{0.25ex}{\accentmark{12}{\char"0951}}ध्वार\raisebox{0.25ex}{\accentmark{15}{\char"1CD2}}ः\kern 0.05em ।\hskip 0.12em plus 0.04em minus 0.02em भद्रा\raisebox{0.15ex}{\accentmark{12}{\char"1CF9}}ऊतप्रा\raisebox{0.25ex}{\accentmark{15}{\char"1CD2}}शा\raisebox{0.25ex}{\accentmark{12}{\char"0951}}स्तयः \nolinebreak\hskip 0.08em \mbox{॥५॥}}} \par}

\vspace{0.8em}
                % --- SUBSECTIONS (Mantras) ---
                
                % Conditional rendering based on output_mode
                \par\filbreak

\phantomsection

\addcontentsline{toc}{subsection}{ऋक् ६}

{\centering \footnotesize \textcolor{AccentPurple}{॥ सोभिरः काण्वः अग्निः उष्णिक् ॥} \par}

\vspace{0.12em}

{\centering \textbf{\textcolor{AccentBlue}{या\raisebox{0.25ex}{\accentmark{15}{\char"1CD2}}जी\raisebox{0.25ex}{\accentmark{12}{\char"0951}}ष्ठं\raisebox{0.25ex}{\accentmark{12}{\char"0951}}त्वाववृमहे\kern 0.05em ।\hskip 0.12em plus 0.04em minus 0.02em देवन्दे\raisebox{0.25ex}{\accentmark{12}{\char"0951}}वत्रा\raisebox{0.15ex}{\accentmark{12}{\char"1CF9}}हो\raisebox{0.15ex}{\accentmark{12}{\char"1CF9}}ता\raisebox{0.25ex}{\accentmark{12}{\char"0951}}रामा\raisebox{0.25ex}{\accentmark{15}{\char"1CD2}}मा\raisebox{0.25ex}{\accentmark{12}{\char"0951}}र्त्यम्\kern 0.05em ।\hskip 0.12em plus 0.04em minus 0.02em अस्या\raisebox{0.25ex}{\accentmark{15}{\char"1CD2}}यज्ञा\raisebox{0.15ex}{\accentmark{12}{\char"1CF9}}स्या\raisebox{0.25ex}{\accentmark{12}{\char"0951}}सुक्रातु\raisebox{0.25ex}{\accentmark{12}{\char"0951}}म् \nolinebreak\hskip 0.08em \mbox{॥६॥}}} \par}

\vspace{0.8em}
                % --- SUBSECTIONS (Mantras) ---
                
                % Conditional rendering based on output_mode
                \par\filbreak

\phantomsection

\addcontentsline{toc}{subsection}{ऋक् ७}

{\centering \footnotesize \textcolor{AccentPurple}{॥ सोभिरः काण्वः अग्निः उष्णिक् ॥} \par}

\vspace{0.12em}

{\centering \textbf{\textcolor{AccentBlue}{ता\raisebox{0.25ex}{\accentmark{15}{\char"1CD2}}दग्नेद्युम्नामा\raisebox{0.25ex}{\accentmark{15}{\char"1CD2}}भारा\kern 0.05em ।\hskip 0.12em plus 0.04em minus 0.02em य\raisebox{0.25ex}{\accentmark{15}{\char"1CD2}}त्सासा\raisebox{0.25ex}{\accentmark{15}{\char"1CD2}}हासाद\raisebox{0.25ex}{\accentmark{12}{\char"0951}}नेकञ्ची\raisebox{0.25ex}{\accentmark{12}{\char"0951}}दत्री\raisebox{0.25ex}{\accentmark{15}{\char"1CD2}}णम्\kern 0.05em ।\hskip 0.12em plus 0.04em minus 0.02em मन्यु\raisebox{0.25ex}{\accentmark{15}{\char"1CD2}}ञ्जना\raisebox{0.25ex}{\accentmark{12}{\char"0951}}स्यादूढ्यम् \nolinebreak\hskip 0.08em \mbox{॥७॥}}} \par}

\vspace{0.8em}
                % --- SUBSECTIONS (Mantras) ---
                
                % Conditional rendering based on output_mode
                \par\filbreak

\phantomsection

\addcontentsline{toc}{subsection}{ऋक् ८}

{\centering \footnotesize \textcolor{AccentPurple}{॥ विश्वमना वैयश्वः अग्निः उष्णिक् ॥} \par}

\vspace{0.12em}

{\centering \textbf{\textcolor{AccentBlue}{य\raisebox{0.25ex}{\accentmark{15}{\char"1CD2}}द्वा\raisebox{0.25ex}{\accentmark{15}{\char"1CD2}}ऊविश्पा\raisebox{0.15ex}{\accentmark{12}{\char"1CF9}}तीश्शितः\kern 0.05em ।\hskip 0.12em plus 0.04em minus 0.02em सुप्रीतो\raisebox{0.25ex}{\accentmark{12}{\char"0951}}मा\raisebox{0.25ex}{\accentmark{15}{\char"1CD2}}नुषोवीशे\kern 0.05em ।\hskip 0.12em plus 0.04em minus 0.02em विश्वेदग्निप्रा\raisebox{0.25ex}{\accentmark{15}{\char"1CD2}}तिर\raisebox{0.25ex}{\accentmark{12}{\char"0951}}क्षांसिसेधति \nolinebreak\hskip 0.08em \mbox{॥८॥}}} \par}

\vspace{0.8em}
                % --- SUBSECTIONS (Mantras) ---
                
                % Conditional rendering based on output_mode
                \par\filbreak

\phantomsection

\addcontentsline{toc}{subsection}{ऋक् ९}

{\centering \textbf{\textcolor{AccentBlue}{आ\raisebox{0.25ex}{\accentmark{15}{\char"1CD2}}या\raisebox{0.25ex}{\accentmark{15}{\char"1CD2}}मग्नि\raisebox{0.25ex}{\accentmark{15}{\char"1CD2}}श्रेष्ठातमः\kern 0.05em ।\hskip 0.12em plus 0.04em minus 0.02em आ\raisebox{0.25ex}{\accentmark{15}{\char"1CD2}}यं\raisebox{0.25ex}{\accentmark{15}{\char"1CD2}}वामाधू\raisebox{0.25ex}{\accentmark{15}{\char"1CD2}}मक्तमः\kern 0.05em ।\hskip 0.12em plus 0.04em minus 0.02em आयं\raisebox{0.25ex}{\accentmark{15}{\char"1CD2}}साहस्रासा\raisebox{0.25ex}{\accentmark{15}{\char"1CD2}}तमः\kern 0.05em ।\hskip 0.12em plus 0.04em minus 0.02em अ\raisebox{0.25ex}{\accentmark{15}{\char"1CD2}}स्मि\raisebox{0.25ex}{\accentmark{15}{\char"1CD2}}न्ना\raisebox{0.25ex}{\accentmark{12}{\char"0951}}स्तुसूवीर्यम् \nolinebreak\hskip 0.08em \mbox{॥९॥}}} \par}

\vspace{0.8em}
                % --- SUBSECTIONS (Mantras) ---
                
                % Conditional rendering based on output_mode
                \par\filbreak

\phantomsection

\addcontentsline{toc}{subsection}{ऋक् १०}

{\centering \textbf{\textcolor{AccentBlue}{आ\raisebox{0.25ex}{\accentmark{15}{\char"1CD2}}दित्य\raisebox{0.25ex}{\accentmark{15}{\char"1CD2}}ःशुक्रा\raisebox{0.25ex}{\accentmark{15}{\char"1CD2}}ऊ\raisebox{0.25ex}{\accentmark{15}{\char"1CD2}}दगात्पूरास्तात्\kern 0.05em ।\hskip 0.12em plus 0.04em minus 0.02em ज्यो\raisebox{0.15ex}{\accentmark{12}{\char"1CF8}}तिःकृण्वन्वी\raisebox{0.25ex}{\accentmark{15}{\char"1CD2}}तामो\raisebox{0.25ex}{\accentmark{12}{\char"0951}}बाधामानः\kern 0.05em ।\hskip 0.12em plus 0.04em minus 0.02em आ\raisebox{0.25ex}{\accentmark{15}{\char"1CD2}}भा\raisebox{0.25ex}{\accentmark{15}{\char"1CD2}}स\raisebox{0.25ex}{\accentmark{12}{\char"0951}}मानाप्रादीशो\raisebox{0.25ex}{\accentmark{12}{\char"0951}}नूसा\raisebox{0.25ex}{\accentmark{15}{\char"1CD2}}र्वाः\kern 0.05em ।\hskip 0.12em plus 0.04em minus 0.02em भद्रस्या\raisebox{0.25ex}{\accentmark{15}{\char"1CD2}}कर्ता\raisebox{0.25ex}{\accentmark{12}{\char"0951}}रूचायन्नाआगात् \nolinebreak\hskip 0.08em \mbox{॥१०॥}}} \par}

\vspace{0.8em}

            % --- End of Kandah Closing Floral Line & Counts ---
            \par\nopagebreak\vspace{0.9em plus 0.3em minus 0.1em}
            \centerline{\textcolor{black!35}{\rule{1.8cm}{0.35pt}}\;\raisebox{-0.22ex}{\textcolor{vedicmaroon}{\large\SixFlowerPetalDotted}}\;\textcolor{black!35}{\rule{1.8cm}{0.35pt}}}
            \nopagebreak\vspace{0.45em}
            \centerline{\devafont\textbf{\textcolor{AccentBlue}{॥ ऋक् सङ्ख्या : १० ॥}}}


    % --- SUPERSECTION ---
    \clearpage           
    \phantomsection              
    \chapter*{\centering तद्वपाठः}
    \addcontentsline{toc}{chapter}{तद्वपाठः}
    \markboth{तद्वपाठः}{तद्वपाठः}
    
    % GLUE FIX: Prevent page break between Chapter and First Section
    \nopagebreak

        

            % --- SECTION ---
            
            \phantomsection  

            \section*{\centering प्रथम खण्डः}
            \addcontentsline{toc}{section}{प्रथम खण्डः} 
            \markright{प्रथम खण्डः}
            \nopagebreak % Keep section title with first mantra
            
            % --- SECTION-LEVEL PROCEDURE LINK ---

                % --- SUBSECTIONS (Mantras) ---
                
                % Conditional rendering based on output_mode
                \par\filbreak

\phantomsection

\addcontentsline{toc}{subsection}{ऋक् १}

{\centering \footnotesize \textcolor{AccentPurple}{॥ शय्युबार्हस्पत्यः इन्द्रः गायत्री ॥} \par}

\vspace{0.12em}

{\centering \textbf{\textcolor{AccentBlue}{त\raisebox{0.25ex}{\accentmark{15}{\char"1CD2}}द्वो\raisebox{0.25ex}{\accentmark{15}{\char"1CD2}}गायासू\raisebox{0.25ex}{\accentmark{15}{\char"1CD2}}तेसा\raisebox{0.25ex}{\accentmark{15}{\char"1CD2}}चा\kern 0.05em ।\hskip 0.12em plus 0.04em minus 0.02em पुरुहूता\raisebox{0.25ex}{\accentmark{15}{\char"1CD2}}यासात्वा\raisebox{0.25ex}{\accentmark{12}{\char"0951}}ने\kern 0.05em ।\hskip 0.12em plus 0.04em minus 0.02em शं\raisebox{0.25ex}{\accentmark{15}{\char"1CD2}}य\raisebox{0.25ex}{\accentmark{15}{\char"1CD2}}त्ग\raisebox{0.25ex}{\accentmark{15}{\char"1CD2}}वेनाशाकीने \nolinebreak\hskip 0.08em \mbox{॥१॥}}} \par}

\vspace{0.8em}
                % --- SUBSECTIONS (Mantras) ---
                
                % Conditional rendering based on output_mode
                \par\filbreak

\phantomsection

\addcontentsline{toc}{subsection}{ऋक् २}

{\centering \footnotesize \textcolor{AccentPurple}{॥ श्रुतकक्षः सुकक्षो वा आंगिरसः इन्द्रः गायत्री ॥} \par}

\vspace{0.12em}

{\centering \textbf{\textcolor{AccentBlue}{य\raisebox{0.15ex}{\accentmark{12}{\char"1CF8}}स्ते\raisebox{0.25ex}{\accentmark{12}{\char"0951}}नूनंशा\raisebox{0.25ex}{\accentmark{12}{\char"0951}}तकृतो\raisebox{0.25ex}{\accentmark{12}{\char"0951}}\kern 0.05em ।\hskip 0.12em plus 0.04em minus 0.02em इन्द्र\raisebox{0.25ex}{\accentmark{12}{\char"0951}}घुम्नीता\raisebox{0.25ex}{\accentmark{12}{\char"0951}}मो\raisebox{0.25ex}{\accentmark{15}{\char"1CD2}}मा\raisebox{0.25ex}{\accentmark{15}{\char"1CD2}}दा\raisebox{0.25ex}{\accentmark{12}{\char"0951}}ः\kern 0.05em ।\hskip 0.12em plus 0.04em minus 0.02em ते\raisebox{0.15ex}{\accentmark{12}{\char"1CF8}}ना\raisebox{0.25ex}{\accentmark{12}{\char"0951}}नूनम्मादे\raisebox{0.25ex}{\accentmark{12}{\char"0951}}मदेः \nolinebreak\hskip 0.08em \mbox{॥२॥}}} \par}

\vspace{0.8em}
                % --- SUBSECTIONS (Mantras) ---
                
                % Conditional rendering based on output_mode
                \par\filbreak

\phantomsection

\addcontentsline{toc}{subsection}{ऋक् ३}

{\centering \footnotesize \textcolor{AccentPurple}{॥ हयतः प्रागाथः इन्द्रः (अग्निर्हवींषि वा) गायत्री ॥} \par}

\vspace{0.12em}

{\centering \textbf{\textcolor{AccentBlue}{गा\raisebox{0.25ex}{\accentmark{15}{\char"1CD2}}वाऊ\raisebox{0.25ex}{\accentmark{15}{\char"1CD2}}पा\raisebox{0.25ex}{\accentmark{12}{\char"0951}}वादावाटे\raisebox{0.25ex}{\accentmark{15}{\char"1CD2}}\kern 0.05em ।\hskip 0.12em plus 0.04em minus 0.02em माही\raisebox{0.25ex}{\accentmark{15}{\char"1CD2}}यज्ञास्या\raisebox{0.15ex}{\accentmark{12}{\char"1CF9}}रप्सू\raisebox{0.25ex}{\accentmark{12}{\char"0951}}दाः\kern 0.05em ।\hskip 0.12em plus 0.04em minus 0.02em ऊ\raisebox{0.25ex}{\accentmark{15}{\char"1CD2}}भाक\raisebox{0.25ex}{\accentmark{15}{\char"1CD2}}र्णा\raisebox{0.25ex}{\accentmark{12}{\char"0951}}हीरण्या\raisebox{0.25ex}{\accentmark{15}{\char"1CD2}}या \nolinebreak\hskip 0.08em \mbox{॥३॥}}} \par}

\vspace{0.8em}
                % --- SUBSECTIONS (Mantras) ---
                
                % Conditional rendering based on output_mode
                \par\filbreak

\phantomsection

\addcontentsline{toc}{subsection}{ऋक् ४}

{\centering \footnotesize \textcolor{AccentPurple}{॥ श्रुतकक्षः (सुकक्षो वा) आंगिरसः इन्द्रः गायत्री ॥} \par}

\vspace{0.12em}

{\centering \textbf{\textcolor{AccentBlue}{आ\raisebox{0.25ex}{\accentmark{15}{\char"1CD2}}राम\raisebox{0.15ex}{\accentmark{12}{\char"1CF9}}श्व\raisebox{0.25ex}{\accentmark{12}{\char"0951}}यगायता\kern 0.05em ।\hskip 0.12em plus 0.04em minus 0.02em श्रू\raisebox{0.15ex}{\accentmark{12}{\char"1CF8}}ता\raisebox{0.25ex}{\accentmark{12}{\char"0951}}कक्षा\raisebox{0.25ex}{\accentmark{15}{\char"1CD2}}रङ्ग\raisebox{0.25ex}{\accentmark{15}{\char"1CD2}}वे\kern 0.05em ।\hskip 0.12em plus 0.04em minus 0.02em आ\raisebox{0.25ex}{\accentmark{15}{\char"1CD2}}रामि\raisebox{0.15ex}{\accentmark{12}{\char"1CF9}}न्द्रा\raisebox{0.25ex}{\accentmark{12}{\char"0951}}स्याधा\raisebox{0.25ex}{\accentmark{15}{\char"1CD2}}म्ने\raisebox{0.15ex}{\accentmark{12}{\char"1CF8}} \nolinebreak\hskip 0.08em \mbox{॥४॥}}} \par}

\vspace{0.8em}
                % --- SUBSECTIONS (Mantras) ---
                
                % Conditional rendering based on output_mode
                \par\filbreak

\phantomsection

\addcontentsline{toc}{subsection}{ऋक् ५}

{\centering \footnotesize \textcolor{AccentPurple}{॥ श्रुतकक्षः (सुकक्षो वा सुकक्षः) आंगिरसः इन्द्रः गायत्री ॥} \par}

\vspace{0.12em}

{\centering \textbf{\textcolor{AccentBlue}{त\raisebox{0.25ex}{\accentmark{15}{\char"1CD2}}मि\raisebox{0.25ex}{\accentmark{15}{\char"1CD2}}न्द्रं\raisebox{0.25ex}{\accentmark{12}{\char"0951}}वाजयामसि\kern 0.05em ।\hskip 0.12em plus 0.04em minus 0.02em मा\raisebox{0.15ex}{\accentmark{12}{\char"1CF8}}हे\raisebox{0.25ex}{\accentmark{12}{\char"0951}}वृत्रायाह\raisebox{0.25ex}{\accentmark{15}{\char"1CD2}}न्ता\raisebox{0.25ex}{\accentmark{12}{\char"0951}}वे\kern 0.05em ।\hskip 0.12em plus 0.04em minus 0.02em स\raisebox{0.25ex}{\accentmark{15}{\char"1CD2}}वृ\raisebox{0.25ex}{\accentmark{15}{\char"1CD2}}षा\raisebox{0.25ex}{\accentmark{12}{\char"0951}}वृषाभो\raisebox{0.25ex}{\accentmark{15}{\char"1CD2}}भू\raisebox{0.25ex}{\accentmark{12}{\char"0951}}वात् \nolinebreak\hskip 0.08em \mbox{॥५॥}}} \par}

\vspace{0.8em}
                % --- SUBSECTIONS (Mantras) ---
                
                % Conditional rendering based on output_mode
                \par\filbreak

\phantomsection

\addcontentsline{toc}{subsection}{ऋक् ६}

{\centering \footnotesize \textcolor{AccentPurple}{॥ देवजामय इन्द्रमातरः ऋषिकाः इन्द्रः गायत्री ॥} \par}

\vspace{0.12em}

{\centering \textbf{\textcolor{AccentBlue}{त्वा\raisebox{0.15ex}{\accentmark{12}{\char"1CF8}}मिन्द्राब\raisebox{0.25ex}{\accentmark{15}{\char"1CD2}}लाद\raisebox{0.25ex}{\accentmark{15}{\char"1CD2}}धी\raisebox{0.25ex}{\accentmark{12}{\char"0951}}\kern 0.05em ।\hskip 0.12em plus 0.04em minus 0.02em साहा\raisebox{0.25ex}{\accentmark{15}{\char"1CD2}}सा\raisebox{0.25ex}{\accentmark{15}{\char"1CD2}}जाता\raisebox{0.25ex}{\accentmark{15}{\char"1CD2}}ओ\raisebox{0.25ex}{\accentmark{15}{\char"1CD2}}ज\raisebox{0.25ex}{\accentmark{12}{\char"0951}}सा\kern 0.05em ।\hskip 0.12em plus 0.04em minus 0.02em त्वं\raisebox{0.15ex}{\accentmark{12}{\char"1CF8}}सन्वृष\raisebox{0.25ex}{\accentmark{12}{\char"0951}}न्वृषेद\raisebox{0.25ex}{\accentmark{12}{\char"0951}}सि \nolinebreak\hskip 0.08em \mbox{॥६॥}}} \par}

\vspace{0.8em}
                % --- SUBSECTIONS (Mantras) ---
                
                % Conditional rendering based on output_mode
                \par\filbreak

\phantomsection

\addcontentsline{toc}{subsection}{ऋक् ७}

{\centering \footnotesize \textcolor{AccentPurple}{॥ गोषूकश्वसूनूकाण्वायनाै इन्द्रः गायत्री ॥} \par}

\vspace{0.12em}

{\centering \textbf{\textcolor{AccentBlue}{यज्ञाइ\raisebox{0.25ex}{\accentmark{15}{\char"1CD2}}न्द्रा\raisebox{0.25ex}{\accentmark{12}{\char"0951}}मवर्धयात्\kern 0.05em ।\hskip 0.12em plus 0.04em minus 0.02em य\raisebox{0.25ex}{\accentmark{15}{\char"1CD2}}त्भूमिंव्या\raisebox{0.25ex}{\accentmark{15}{\char"1CD2}}वा\raisebox{0.25ex}{\accentmark{12}{\char"0951}}र्क्तयात्\kern 0.05em ।\hskip 0.12em plus 0.04em minus 0.02em चक्रा\raisebox{0.15ex}{\accentmark{12}{\char"1CF9}}णाओपाश\raisebox{0.25ex}{\accentmark{15}{\char"1CD2}}न्दीवि \nolinebreak\hskip 0.08em \mbox{॥७॥}}} \par}

\vspace{0.8em}
                % --- SUBSECTIONS (Mantras) ---
                
                % Conditional rendering based on output_mode
                \par\filbreak

\phantomsection

\addcontentsline{toc}{subsection}{ऋक् ८}

{\centering \footnotesize \textcolor{AccentPurple}{॥ गोषूकश्वसूनूकाण्वायनाै इन्द्रः गायत्री ॥} \par}

\vspace{0.12em}

{\centering \textbf{\textcolor{AccentBlue}{या\raisebox{0.15ex}{\accentmark{12}{\char"1CF8}}दिन्द्राहं\raisebox{0.15ex}{\accentmark{12}{\char"1CF8}}या\raisebox{0.25ex}{\accentmark{15}{\char"1CD2}}थात्वम्\raisebox{0.25ex}{\accentmark{15}{\char"1CD2}}\kern 0.05em ।\hskip 0.12em plus 0.04em minus 0.02em ईशी\raisebox{0.25ex}{\accentmark{12}{\char"0951}}यावा\raisebox{0.25ex}{\accentmark{15}{\char"1CD2}}स्वाए\raisebox{0.25ex}{\accentmark{15}{\char"1CD2}}काई\raisebox{0.25ex}{\accentmark{15}{\char"1CD2}}त्\kern 0.05em ।\hskip 0.12em plus 0.04em minus 0.02em स्तोता\raisebox{0.25ex}{\accentmark{15}{\char"1CD2}}मेगोसाखास्यात् \nolinebreak\hskip 0.08em \mbox{॥८॥}}} \par}

\vspace{0.8em}
                % --- SUBSECTIONS (Mantras) ---
                
                % Conditional rendering based on output_mode
                \par\filbreak

\phantomsection

\addcontentsline{toc}{subsection}{ऋक् ९}

{\centering \footnotesize \textcolor{AccentPurple}{॥ मेधातिथिः काण्वः प्रियमेधश्चांगिरसः इन्द्रः गायत्री ॥} \par}

\vspace{0.12em}

{\centering \textbf{\textcolor{AccentBlue}{प\raisebox{0.25ex}{\accentmark{15}{\char"1CD2}}न्यं\raisebox{0.25ex}{\accentmark{12}{\char"0951}}पन्यामि\raisebox{0.25ex}{\accentmark{15}{\char"1CD2}}स्तो\raisebox{0.25ex}{\accentmark{12}{\char"0951}}तरः\kern 0.05em ।\hskip 0.12em plus 0.04em minus 0.02em आ\raisebox{0.25ex}{\accentmark{15}{\char"1CD2}}धा\raisebox{0.25ex}{\accentmark{12}{\char"0951}}वातामा\raisebox{0.25ex}{\accentmark{15}{\char"1CD2}}द्या\raisebox{0.25ex}{\accentmark{12}{\char"0951}}या\kern 0.05em ।\hskip 0.12em plus 0.04em minus 0.02em सो\raisebox{0.15ex}{\accentmark{12}{\char"1CF8}}मं\raisebox{0.25ex}{\accentmark{12}{\char"0951}}वीरा\raisebox{0.25ex}{\accentmark{15}{\char"1CD2}}यासूरा\raisebox{0.25ex}{\accentmark{12}{\char"0951}}या \nolinebreak\hskip 0.08em \mbox{॥९॥}}} \par}

\vspace{0.8em}
                % --- SUBSECTIONS (Mantras) ---
                
                % Conditional rendering based on output_mode
                \par\filbreak

\phantomsection

\addcontentsline{toc}{subsection}{ऋक् १०}

{\centering \footnotesize \textcolor{AccentPurple}{॥ मेधातिथिः काण्वः प्रियमेधश्चांगिरसः इन्द्रः गायत्री ॥} \par}

\vspace{0.12em}

{\centering \textbf{\textcolor{AccentBlue}{ईदं\raisebox{0.25ex}{\accentmark{15}{\char"1CD2}}वासो\raisebox{0.25ex}{\accentmark{15}{\char"1CD2}}सूताम\raisebox{0.25ex}{\accentmark{15}{\char"1CD2}}न्धाः\kern 0.05em ।\hskip 0.12em plus 0.04em minus 0.02em पीबासूपू\raisebox{0.25ex}{\accentmark{12}{\char"0951}}र्णा\raisebox{0.25ex}{\accentmark{12}{\char"0951}}मू\raisebox{0.25ex}{\accentmark{15}{\char"1CD2}}दरम्\raisebox{0.25ex}{\accentmark{12}{\char"0951}}\kern 0.05em ।\hskip 0.12em plus 0.04em minus 0.02em आनाभइन्द्रा\raisebox{0.25ex}{\accentmark{15}{\char"1CD2}}रीमाते \nolinebreak\hskip 0.08em \mbox{॥१०॥}}} \par}

\vspace{0.8em}

            % --- End of Kandah Closing Floral Line & Counts ---
            \par\nopagebreak\vspace{0.9em plus 0.3em minus 0.1em}
            \centerline{\textcolor{black!35}{\rule{1.8cm}{0.35pt}}\;\raisebox{-0.22ex}{\textcolor{vedicmaroon}{\large\SixFlowerPetalDotted}}\;\textcolor{black!35}{\rule{1.8cm}{0.35pt}}}
            \nopagebreak\vspace{0.45em}
            \centerline{\devafont\textbf{\textcolor{AccentBlue}{॥ ऋक् सङ्ख्या : १० ॥}}}
        
            \clearpage

            % --- SECTION ---
            
            \phantomsection  

            \section*{\centering द्वितीय खण्डः}
            \addcontentsline{toc}{section}{द्वितीय खण्डः} 
            \markright{द्वितीय खण्डः}
            \nopagebreak % Keep section title with first mantra
            
            % --- SECTION-LEVEL PROCEDURE LINK ---

                % --- SUBSECTIONS (Mantras) ---
                
                % Conditional rendering based on output_mode
                \par\filbreak

\phantomsection

\addcontentsline{toc}{subsection}{ऋक् १}

{\centering \footnotesize \textcolor{AccentPurple}{॥ सुकक्षत्रुतकाै (सुकक्ष आंगिरसः) इन्द्रः गायत्री ॥} \par}

\vspace{0.12em}

{\centering \textbf{\textcolor{AccentBlue}{उ\raisebox{0.25ex}{\accentmark{15}{\char"1CD2}}त्घेदभी\raisebox{0.25ex}{\accentmark{15}{\char"1CD2}}श्रूता\raisebox{0.25ex}{\accentmark{15}{\char"1CD2}}मा\raisebox{0.25ex}{\accentmark{12}{\char"0951}}घम्\kern 0.05em ।\hskip 0.12em plus 0.04em minus 0.02em वृषाभ\raisebox{0.25ex}{\accentmark{15}{\char"1CD2}}न्नर्या\raisebox{0.25ex}{\accentmark{12}{\char"0951}}पसम्\kern 0.05em ।\hskip 0.12em plus 0.04em minus 0.02em अस्तारमेळिसूर्या \nolinebreak\hskip 0.08em \mbox{॥१॥}}} \par}

\vspace{0.8em}
                % --- SUBSECTIONS (Mantras) ---
                
                % Conditional rendering based on output_mode
                \par\filbreak

\phantomsection

\addcontentsline{toc}{subsection}{ऋक् २}

{\centering \footnotesize \textcolor{AccentPurple}{॥ सुकक्षत्रुतकाै (सुकक्ष आंगिरसः) इन्द्रः गायत्री ॥} \par}

\vspace{0.12em}

{\centering \textbf{\textcolor{AccentBlue}{यादद्याकाच्चावृत्रहन्\kern 0.05em ।\hskip 0.12em plus 0.04em minus 0.02em ऊदगाआभी\raisebox{0.25ex}{\accentmark{15}{\char"1CD2}}सूर्या\raisebox{0.25ex}{\accentmark{15}{\char"1CD2}}\kern 0.05em ।\hskip 0.12em plus 0.04em minus 0.02em स\raisebox{0.25ex}{\accentmark{15}{\char"1CD2}}र्वन्ता\raisebox{0.25ex}{\accentmark{15}{\char"1CD2}}दिन्द्राते\raisebox{0.25ex}{\accentmark{12}{\char"0951}}वा\raisebox{0.25ex}{\accentmark{15}{\char"1CD2}}शे \nolinebreak\hskip 0.08em \mbox{॥२॥}}} \par}

\vspace{0.8em}
                % --- SUBSECTIONS (Mantras) ---
                
                % Conditional rendering based on output_mode
                \par\filbreak

\phantomsection

\addcontentsline{toc}{subsection}{ऋक् ३}

{\centering \footnotesize \textcolor{AccentPurple}{॥ भारद्वाजः (शय्युबार्हस्पत्यः) इन्द्रः गायत्री ॥} \par}

\vspace{0.12em}

{\centering \textbf{\textcolor{AccentBlue}{य\raisebox{0.25ex}{\accentmark{15}{\char"1CD2}}आ\raisebox{0.25ex}{\accentmark{15}{\char"1CD2}}ना\raisebox{0.25ex}{\accentmark{12}{\char"0951}}यत्पारावा\raisebox{0.25ex}{\accentmark{15}{\char"1CD2}}ताः\kern 0.05em ।\hskip 0.12em plus 0.04em minus 0.02em सू\raisebox{0.25ex}{\accentmark{15}{\char"1CD2}}नी\raisebox{0.25ex}{\accentmark{12}{\char"0951}}तीतूर्वा\raisebox{0.25ex}{\accentmark{15}{\char"1CD2}}शं\raisebox{0.25ex}{\accentmark{15}{\char"1CD2}}या\raisebox{0.25ex}{\accentmark{15}{\char"1CD2}}दुम्\raisebox{0.25ex}{\accentmark{12}{\char"0951}}\kern 0.05em ।\hskip 0.12em plus 0.04em minus 0.02em इ\raisebox{0.25ex}{\accentmark{15}{\char"1CD2}}न्द्रस्सा\raisebox{0.25ex}{\accentmark{15}{\char"1CD2}}नोयू\raisebox{0.25ex}{\accentmark{12}{\char"0951}}वासा\raisebox{0.25ex}{\accentmark{15}{\char"1CD2}}खा\raisebox{0.25ex}{\accentmark{12}{\char"0951}} \nolinebreak\hskip 0.08em \mbox{॥३॥}}} \par}

\vspace{0.8em}
                % --- SUBSECTIONS (Mantras) ---
                
                % Conditional rendering based on output_mode
                \par\filbreak

\phantomsection

\addcontentsline{toc}{subsection}{ऋक् ४}

{\centering \footnotesize \textcolor{AccentPurple}{॥ श्रुतकक्षः (सुकक्षो वा आंगिरसः) इन्द्रः गायत्री ॥} \par}

\vspace{0.12em}

{\centering \textbf{\textcolor{AccentBlue}{मा\raisebox{0.25ex}{\accentmark{15}{\char"1CD2}}ना\raisebox{0.25ex}{\accentmark{12}{\char"0951}}इन्द्राभ्याटट्यदीश\raisebox{0.25ex}{\accentmark{12}{\char"0951}}ः\kern 0.05em ।\hskip 0.12em plus 0.04em minus 0.02em सू\raisebox{0.15ex}{\accentmark{12}{\char"1CF9}}रोअत्तुष्वाया\raisebox{0.25ex}{\accentmark{12}{\char"0951}}मात्\kern 0.05em ।\hskip 0.12em plus 0.04em minus 0.02em त्वा\raisebox{0.25ex}{\accentmark{15}{\char"1CD2}}यूजा\raisebox{0.25ex}{\accentmark{15}{\char"1CD2}}वानेमाता\raisebox{0.25ex}{\accentmark{15}{\char"1CD2}}त् \nolinebreak\hskip 0.08em \mbox{॥४॥}}} \par}

\vspace{0.8em}
                % --- SUBSECTIONS (Mantras) ---
                
                % Conditional rendering based on output_mode
                \par\filbreak

\phantomsection

\addcontentsline{toc}{subsection}{ऋक् ५}

{\centering \footnotesize \textcolor{AccentPurple}{॥ मधुच्छन्दा वैश्वामित्रः इन्द्रः गायत्री ॥} \par}

\vspace{0.12em}

{\centering \textbf{\textcolor{AccentBlue}{ए\raisebox{0.25ex}{\accentmark{15}{\char"1CD2}}न्द्रा\raisebox{0.25ex}{\accentmark{12}{\char"0951}}सानासिं\raisebox{0.25ex}{\accentmark{15}{\char"1CD2}}रायिम्\raisebox{0.25ex}{\accentmark{15}{\char"1CD2}}\kern 0.05em ।\hskip 0.12em plus 0.04em minus 0.02em सा\raisebox{0.25ex}{\accentmark{15}{\char"1CD2}}जी\raisebox{0.25ex}{\accentmark{15}{\char"1CD2}}त्वा\raisebox{0.25ex}{\accentmark{12}{\char"0951}}नंसादासा\raisebox{0.25ex}{\accentmark{15}{\char"1CD2}}हम्\kern 0.05em ।\hskip 0.12em plus 0.04em minus 0.02em वर्षी\raisebox{0.25ex}{\accentmark{12}{\char"0951}}ष्ठामू\raisebox{0.25ex}{\accentmark{15}{\char"1CD2}}ताये\raisebox{0.25ex}{\accentmark{15}{\char"1CD2}}भरा \nolinebreak\hskip 0.08em \mbox{॥५॥}}} \par}

\vspace{0.8em}
                % --- SUBSECTIONS (Mantras) ---
                
                % Conditional rendering based on output_mode
                \par\filbreak

\phantomsection

\addcontentsline{toc}{subsection}{ऋक् ६}

{\centering \footnotesize \textcolor{AccentPurple}{॥ मधुच्छन्दा वैश्वामित्रः इन्द्रः गायत्री ॥} \par}

\vspace{0.12em}

{\centering \textbf{\textcolor{AccentBlue}{इ\raisebox{0.15ex}{\accentmark{12}{\char"1CF8}}न्द्रं\raisebox{0.25ex}{\accentmark{15}{\char"1CD2}}वाय\raisebox{0.25ex}{\accentmark{15}{\char"1CD2}}म्माहाधाने\raisebox{0.25ex}{\accentmark{15}{\char"1CD2}}\kern 0.05em ।\hskip 0.12em plus 0.04em minus 0.02em इन्द्रा\raisebox{0.25ex}{\accentmark{15}{\char"1CD2}}मर्भे\raisebox{0.25ex}{\accentmark{12}{\char"0951}}हवामहे\kern 0.05em ।\hskip 0.12em plus 0.04em minus 0.02em यूजं\raisebox{0.25ex}{\accentmark{12}{\char"0951}}वृ\raisebox{0.15ex}{\accentmark{12}{\char"1CF8}}त्रे\raisebox{0.15ex}{\accentmark{12}{\char"1CF9}}षुव\raisebox{0.25ex}{\accentmark{15}{\char"1CD2}}ज्रीण\raisebox{0.25ex}{\accentmark{12}{\char"0951}}म् \nolinebreak\hskip 0.08em \mbox{॥६॥}}} \par}

\vspace{0.8em}
                % --- SUBSECTIONS (Mantras) ---
                
                % Conditional rendering based on output_mode
                \par\filbreak

\phantomsection

\addcontentsline{toc}{subsection}{ऋक् ७}

{\centering \footnotesize \textcolor{AccentPurple}{॥ त्रिशोकः काण्वः इन्द्रः गायत्री ॥} \par}

\vspace{0.12em}

{\centering \textbf{\textcolor{AccentBlue}{आपी\raisebox{0.25ex}{\accentmark{12}{\char"0951}}बत्काद्रू\raisebox{0.15ex}{\accentmark{12}{\char"1CF9}}वा\raisebox{0.25ex}{\accentmark{12}{\char"0951}}स्सुतम्\raisebox{0.25ex}{\accentmark{15}{\char"1CD2}}\kern 0.05em ।\hskip 0.12em plus 0.04em minus 0.02em इन्द्रा\raisebox{0.15ex}{\accentmark{12}{\char"1CF8}}सा\raisebox{0.25ex}{\accentmark{12}{\char"0951}}हा\raisebox{0.25ex}{\accentmark{15}{\char"1CD2}}स्रा\raisebox{0.25ex}{\accentmark{12}{\char"0951}}वांभे\kern 0.05em ।\hskip 0.12em plus 0.04em minus 0.02em तत्राददी\raisebox{0.25ex}{\accentmark{12}{\char"0951}}ष्टापौंठस्यम् \nolinebreak\hskip 0.08em \mbox{॥७॥}}} \par}

\vspace{0.8em}
                % --- SUBSECTIONS (Mantras) ---
                
                % Conditional rendering based on output_mode
                \par\filbreak

\phantomsection

\addcontentsline{toc}{subsection}{ऋक् ८}

{\centering \footnotesize \textcolor{AccentPurple}{॥ वसिष्ठोः मैत्रावरुणिः इन्द्रः गायत्री ॥} \par}

\vspace{0.12em}

{\centering \textbf{\textcolor{AccentBlue}{वा\raisebox{0.25ex}{\accentmark{15}{\char"1CD2}}यामिन्द्रात्वाया\raisebox{0.25ex}{\accentmark{15}{\char"1CD2}}वा\raisebox{0.25ex}{\accentmark{12}{\char"0951}}ः\kern 0.05em ।\hskip 0.12em plus 0.04em minus 0.02em आभि\raisebox{0.25ex}{\accentmark{15}{\char"1CD2}}प्रानो\raisebox{0.25ex}{\accentmark{12}{\char"0951}}नुमोवृषन्\kern 0.05em ।\hskip 0.12em plus 0.04em minus 0.02em विद्धीत्वाटट्यस्या\raisebox{0.25ex}{\accentmark{15}{\char"1CD2}}नोवसो \nolinebreak\hskip 0.08em \mbox{॥८॥}}} \par}

\vspace{0.8em}
                % --- SUBSECTIONS (Mantras) ---
                
                % Conditional rendering based on output_mode
                \par\filbreak

\phantomsection

\addcontentsline{toc}{subsection}{ऋक् ९}

{\centering \footnotesize \textcolor{AccentPurple}{॥ त्रिशोकः काण्वः इन्द्रः (अग्नीन्द्रौ) गायत्री ॥} \par}

\vspace{0.12em}

{\centering \textbf{\textcolor{AccentBlue}{आ\raisebox{0.25ex}{\accentmark{15}{\char"1CD2}}घाये\raisebox{0.25ex}{\accentmark{15}{\char"1CD2}}अग्नी\raisebox{0.25ex}{\accentmark{15}{\char"1CD2}}मिन्धा\raisebox{0.25ex}{\accentmark{15}{\char"1CD2}}ते\kern 0.05em ।\hskip 0.12em plus 0.04em minus 0.02em स्तृणन्तीबही\raisebox{0.15ex}{\accentmark{12}{\char"1CF9}}रानूषाक्\kern 0.05em ।\hskip 0.12em plus 0.04em minus 0.02em ए\raisebox{0.25ex}{\accentmark{15}{\char"1CD2}}षा\raisebox{0.25ex}{\accentmark{15}{\char"1CD2}}मि\raisebox{0.25ex}{\accentmark{15}{\char"1CD2}}न्द्रोयूवासाखा\raisebox{0.25ex}{\accentmark{15}{\char"1CD2}} \nolinebreak\hskip 0.08em \mbox{॥९॥}}} \par}

\vspace{0.8em}
                % --- SUBSECTIONS (Mantras) ---
                
                % Conditional rendering based on output_mode
                \par\filbreak

\phantomsection

\addcontentsline{toc}{subsection}{ऋक् १०}

{\centering \footnotesize \textcolor{AccentPurple}{॥ त्रिशोकः काण्वः इन्द्रः गायत्री ॥} \par}

\vspace{0.12em}

{\centering \textbf{\textcolor{AccentBlue}{भि\raisebox{0.25ex}{\accentmark{15}{\char"1CD2}}न्धि\raisebox{0.25ex}{\accentmark{12}{\char"0951}}विश्वा\raisebox{0.25ex}{\accentmark{15}{\char"1CD2}}आपद्वी\raisebox{0.25ex}{\accentmark{15}{\char"1CD2}}षाः\kern 0.05em ।\hskip 0.12em plus 0.04em minus 0.02em पा\raisebox{0.25ex}{\accentmark{12}{\char"0951}}रीबाधो\raisebox{0.15ex}{\accentmark{12}{\char"1CF9}}जही\raisebox{0.25ex}{\accentmark{12}{\char"0951}}मृ\raisebox{0.25ex}{\accentmark{15}{\char"1CD2}}धा\raisebox{0.25ex}{\accentmark{12}{\char"0951}}ः\kern 0.05em ।\hskip 0.12em plus 0.04em minus 0.02em वा\raisebox{0.15ex}{\accentmark{12}{\char"1CF8}}सु\raisebox{0.25ex}{\accentmark{12}{\char"0951}}स्पार्हंतादाभा\raisebox{0.25ex}{\accentmark{12}{\char"0951}}रा \nolinebreak\hskip 0.08em \mbox{॥१०॥}}} \par}

\vspace{0.8em}

            % --- End of Kandah Closing Floral Line & Counts ---
            \par\nopagebreak\vspace{0.9em plus 0.3em minus 0.1em}
            \centerline{\textcolor{black!35}{\rule{1.8cm}{0.35pt}}\;\raisebox{-0.22ex}{\textcolor{vedicmaroon}{\large\SixFlowerPetalDotted}}\;\textcolor{black!35}{\rule{1.8cm}{0.35pt}}}
            \nopagebreak\vspace{0.45em}
            \centerline{\devafont\textbf{\textcolor{AccentBlue}{॥ ऋक् सङ्ख्या : १० ॥}}}
        
            \clearpage

            % --- SECTION ---
            
            \phantomsection  

            \section*{\centering तृतीय खण्डः}
            \addcontentsline{toc}{section}{तृतीय खण्डः} 
            \markright{तृतीय खण्डः}
            \nopagebreak % Keep section title with first mantra
            
            % --- SECTION-LEVEL PROCEDURE LINK ---

                % --- SUBSECTIONS (Mantras) ---
                
                % Conditional rendering based on output_mode
                \par\filbreak

\phantomsection

\addcontentsline{toc}{subsection}{ऋक् १}

{\centering \footnotesize \textcolor{AccentPurple}{॥ कण्वो घौरः इन्द्रः (मरुतः) गायत्री ॥} \par}

\vspace{0.12em}

{\centering \textbf{\textcolor{AccentBlue}{ईहे\raisebox{0.25ex}{\accentmark{15}{\char"1CD2}}वा\raisebox{0.25ex}{\accentmark{12}{\char"0951}}शृण्वाएषाम्\kern 0.05em ।\hskip 0.12em plus 0.04em minus 0.02em काशाह\raisebox{0.15ex}{\accentmark{12}{\char"1CF9}}स्तेषूयद्वादा\raisebox{0.25ex}{\accentmark{12}{\char"0951}}न्\kern 0.05em ।\hskip 0.12em plus 0.04em minus 0.02em निया\raisebox{0.15ex}{\accentmark{12}{\char"1CF8}}मैःचित्रा\raisebox{0.25ex}{\accentmark{15}{\char"1CD2}}मृञ्जते \nolinebreak\hskip 0.08em \mbox{॥१॥}}} \par}

\vspace{0.8em}
                % --- SUBSECTIONS (Mantras) ---
                
                % Conditional rendering based on output_mode
                \par\filbreak

\phantomsection

\addcontentsline{toc}{subsection}{ऋक् २}

{\centering \footnotesize \textcolor{AccentPurple}{॥ त्रिशोकः इन्द्रः गायत्री ॥} \par}

\vspace{0.12em}

{\centering \textbf{\textcolor{AccentBlue}{ईमा\raisebox{0.15ex}{\accentmark{12}{\char"1CF9}}ऊत्वावीचा\raisebox{0.25ex}{\accentmark{12}{\char"0951}}क्षते\kern 0.05em ।\hskip 0.12em plus 0.04em minus 0.02em सा\raisebox{0.25ex}{\accentmark{15}{\char"1CD2}}खा\raisebox{0.25ex}{\accentmark{12}{\char"0951}}याइन्द्रासोमीना\raisebox{0.25ex}{\accentmark{12}{\char"0951}}ः\kern 0.05em ।\hskip 0.12em plus 0.04em minus 0.02em पुष्टावन्तोया\raisebox{0.15ex}{\accentmark{12}{\char"1CF9}}था\raisebox{0.25ex}{\accentmark{12}{\char"0951}}पाशु\raisebox{0.25ex}{\accentmark{15}{\char"1CD2}}म् \nolinebreak\hskip 0.08em \mbox{॥२॥}}} \par}

\vspace{0.8em}
                % --- SUBSECTIONS (Mantras) ---
                
                % Conditional rendering based on output_mode
                \par\filbreak

\phantomsection

\addcontentsline{toc}{subsection}{ऋक् ३}

{\centering \footnotesize \textcolor{AccentPurple}{॥ वत्सः काण्वः इन्द्रः गायत्री ॥} \par}

\vspace{0.12em}

{\centering \textbf{\textcolor{AccentBlue}{सा\raisebox{0.25ex}{\accentmark{15}{\char"1CD2}}मा\raisebox{0.25ex}{\accentmark{12}{\char"0951}}स्यामन्या\raisebox{0.25ex}{\accentmark{15}{\char"1CD2}}वेवीशा\raisebox{0.25ex}{\accentmark{12}{\char"0951}}ः\kern 0.05em ।\hskip 0.12em plus 0.04em minus 0.02em विश्वे\raisebox{0.25ex}{\accentmark{12}{\char"0951}}नमन्ताकृष्टा\raisebox{0.25ex}{\accentmark{15}{\char"1CD2}}याः\kern 0.05em ।\hskip 0.12em plus 0.04em minus 0.02em समुद्रा\raisebox{0.15ex}{\accentmark{12}{\char"1CF9}}ये\raisebox{0.25ex}{\accentmark{12}{\char"0951}}वासिन्धावः \nolinebreak\hskip 0.08em \mbox{॥३॥}}} \par}

\vspace{0.8em}
                % --- SUBSECTIONS (Mantras) ---
                
                % Conditional rendering based on output_mode
                \par\filbreak

\phantomsection

\addcontentsline{toc}{subsection}{ऋक् ४}

{\centering \footnotesize \textcolor{AccentPurple}{॥ कुसीदी काण्वः इन्द्रः (विश्वेदेवाः) गायत्री ॥} \par}

\vspace{0.12em}

{\centering \textbf{\textcolor{AccentBlue}{देवा\raisebox{0.25ex}{\accentmark{15}{\char"1CD2}}नामी\raisebox{0.15ex}{\accentmark{12}{\char"1CF9}}दावो\raisebox{0.25ex}{\accentmark{12}{\char"0951}}माहा\raisebox{0.25ex}{\accentmark{15}{\char"1CD2}}त्\kern 0.05em ।\hskip 0.12em plus 0.04em minus 0.02em ता\raisebox{0.25ex}{\accentmark{15}{\char"1CD2}}दा\raisebox{0.25ex}{\accentmark{15}{\char"1CD2}}वृणीमहेवायम्\kern 0.05em ।\hskip 0.12em plus 0.04em minus 0.02em वृष्णा\raisebox{0.25ex}{\accentmark{12}{\char"0951}}मस्मा\raisebox{0.15ex}{\accentmark{12}{\char"1CF9}}भ्या\raisebox{0.25ex}{\accentmark{12}{\char"0951}}मूता\raisebox{0.25ex}{\accentmark{15}{\char"1CD2}}ये \nolinebreak\hskip 0.08em \mbox{॥४॥}}} \par}

\vspace{0.8em}
                % --- SUBSECTIONS (Mantras) ---
                
                % Conditional rendering based on output_mode
                \par\filbreak

\phantomsection

\addcontentsline{toc}{subsection}{ऋक् ५}

{\centering \footnotesize \textcolor{AccentPurple}{॥ मेधातिथिः काण्वः इन्द्रः (ब्रह्मणस्पतिः) गायत्री ॥} \par}

\vspace{0.12em}

{\centering \textbf{\textcolor{AccentBlue}{सोमानांस्वा\raisebox{0.25ex}{\accentmark{15}{\char"1CD2}}रा\raisebox{0.25ex}{\accentmark{12}{\char"0951}}णाम्\kern 0.05em ।\hskip 0.12em plus 0.04em minus 0.02em कृ\raisebox{0.25ex}{\accentmark{15}{\char"1CD2}}णू\raisebox{0.25ex}{\accentmark{15}{\char"1CD2}}ही\raisebox{0.25ex}{\accentmark{12}{\char"0951}}ब्रह्मणस्पते\kern 0.05em ।\hskip 0.12em plus 0.04em minus 0.02em कक्षीवन्तंया\raisebox{0.15ex}{\accentmark{12}{\char"1CF9}}औशीज\raisebox{0.25ex}{\accentmark{15}{\char"1CD2}}ः \nolinebreak\hskip 0.08em \mbox{॥५॥}}} \par}

\vspace{0.8em}
                % --- SUBSECTIONS (Mantras) ---
                
                % Conditional rendering based on output_mode
                \par\filbreak

\phantomsection

\addcontentsline{toc}{subsection}{ऋक् ६}

{\centering \footnotesize \textcolor{AccentPurple}{॥ श्रुतकक्षः (सुकक्षः) आंगिरसः इन्द्रः गायत्री ॥} \par}

\vspace{0.12em}

{\centering \textbf{\textcolor{AccentBlue}{बोध\raisebox{0.25ex}{\accentmark{15}{\char"1CD2}}न्माना\raisebox{0.25ex}{\accentmark{15}{\char"1CD2}}इद\raisebox{0.25ex}{\accentmark{15}{\char"1CD2}}स्तुनः\kern 0.05em ।\hskip 0.12em plus 0.04em minus 0.02em वृत्राहा\raisebox{0.15ex}{\accentmark{12}{\char"1CF8}}भू\raisebox{0.25ex}{\accentmark{15}{\char"1CD2}}र्या\raisebox{0.25ex}{\accentmark{12}{\char"0951}}सुतिः\kern 0.05em ।\hskip 0.12em plus 0.04em minus 0.02em श्रृणो\raisebox{0.15ex}{\accentmark{12}{\char"1CF9}}तू\raisebox{0.25ex}{\accentmark{12}{\char"0951}}शक्रा\raisebox{0.25ex}{\accentmark{15}{\char"1CD2}}आशी\raisebox{0.25ex}{\accentmark{15}{\char"1CD2}}ष\raisebox{0.25ex}{\accentmark{12}{\char"0951}}म् \nolinebreak\hskip 0.08em \mbox{॥६॥}}} \par}

\vspace{0.8em}
                % --- SUBSECTIONS (Mantras) ---
                
                % Conditional rendering based on output_mode
                \par\filbreak

\phantomsection

\addcontentsline{toc}{subsection}{ऋक् ७}

{\centering \footnotesize \textcolor{AccentPurple}{॥ श्यावाश्व आत्रयेः इन्द्रः (सविता) गायत्री ॥} \par}

\vspace{0.12em}

{\centering \textbf{\textcolor{AccentBlue}{अ\raisebox{0.25ex}{\accentmark{15}{\char"1CD2}}द्या\raisebox{0.25ex}{\accentmark{15}{\char"1CD2}}नो\raisebox{0.25ex}{\accentmark{12}{\char"0951}}देवसवितः\kern 0.05em ।\hskip 0.12em plus 0.04em minus 0.02em प्रा\raisebox{0.25ex}{\accentmark{15}{\char"1CD2}}जा\raisebox{0.25ex}{\accentmark{15}{\char"1CD2}}वा\raisebox{0.25ex}{\accentmark{12}{\char"0951}}त्सा\raisebox{0.25ex}{\accentmark{15}{\char"1CD2}}वीसौ\raisebox{0.25ex}{\accentmark{15}{\char"1CD2}}भा\raisebox{0.25ex}{\accentmark{12}{\char"0951}}गम्\kern 0.05em ।\hskip 0.12em plus 0.04em minus 0.02em पा\raisebox{0.15ex}{\accentmark{12}{\char"1CF8}}रा\raisebox{0.15ex}{\accentmark{12}{\char"1CF8}}दुष्वा\raisebox{0.25ex}{\accentmark{15}{\char"1CD2}}प्न्यिं\raisebox{0.25ex}{\accentmark{12}{\char"0951}}सुवा \nolinebreak\hskip 0.08em \mbox{॥७॥}}} \par}

\vspace{0.8em}
                % --- SUBSECTIONS (Mantras) ---
                
                % Conditional rendering based on output_mode
                \par\filbreak

\phantomsection

\addcontentsline{toc}{subsection}{ऋक् ८}

{\centering \footnotesize \textcolor{AccentPurple}{॥ प्रगाथः इन्द्रः गायत्री ॥} \par}

\vspace{0.12em}

{\centering \textbf{\textcolor{AccentBlue}{क्वाट्ट्यस्या\raisebox{0.15ex}{\accentmark{12}{\char"1CF9}}वृ\raisebox{0.25ex}{\accentmark{12}{\char"0951}}षाभोयू\raisebox{0.25ex}{\accentmark{12}{\char"0951}}वा\kern 0.05em ।\hskip 0.12em plus 0.04em minus 0.02em तुविग्री\raisebox{0.25ex}{\accentmark{15}{\char"1CD2}}वोआ\raisebox{0.25ex}{\accentmark{15}{\char"1CD2}}नानतः\kern 0.05em ।\hskip 0.12em plus 0.04em minus 0.02em ब्रह्मा\raisebox{0.25ex}{\accentmark{15}{\char"1CD2}}कस्तंसा\raisebox{0.25ex}{\accentmark{15}{\char"1CD2}}पर्यति \nolinebreak\hskip 0.08em \mbox{॥८॥}}} \par}

\vspace{0.8em}
                % --- SUBSECTIONS (Mantras) ---
                
                % Conditional rendering based on output_mode
                \par\filbreak

\phantomsection

\addcontentsline{toc}{subsection}{ऋक् ९}

{\centering \footnotesize \textcolor{AccentPurple}{॥ वत्सः काण्वः इन्द्रः गायत्री ॥} \par}

\vspace{0.12em}

{\centering \textbf{\textcolor{AccentBlue}{ऊपह्वारे\raisebox{0.15ex}{\accentmark{12}{\char"1CF9}}गी\raisebox{0.25ex}{\accentmark{12}{\char"0951}}रीणाम्\raisebox{0.25ex}{\accentmark{15}{\char"1CD2}}\kern 0.05em ।\hskip 0.12em plus 0.04em minus 0.02em संगमे\raisebox{0.15ex}{\accentmark{12}{\char"1CF9}}चा\raisebox{0.25ex}{\accentmark{12}{\char"0951}}नादी\raisebox{0.25ex}{\accentmark{15}{\char"1CD2}}नाम्\raisebox{0.25ex}{\accentmark{12}{\char"0951}}\kern 0.05em ।\hskip 0.12em plus 0.04em minus 0.02em धी\raisebox{0.25ex}{\accentmark{12}{\char"0951}}या\raisebox{0.25ex}{\accentmark{15}{\char"1CD2}}विप्रो\raisebox{0.25ex}{\accentmark{12}{\char"0951}}अजायत \nolinebreak\hskip 0.08em \mbox{॥९॥}}} \par}

\vspace{0.8em}
                % --- SUBSECTIONS (Mantras) ---
                
                % Conditional rendering based on output_mode
                \par\filbreak

\phantomsection

\addcontentsline{toc}{subsection}{ऋक् १०}

{\centering \footnotesize \textcolor{AccentPurple}{॥ इरिम्बिठः काण्वः इन्द्रः गायत्री ॥} \par}

\vspace{0.12em}

{\centering \textbf{\textcolor{AccentBlue}{प्रा\raisebox{0.25ex}{\accentmark{15}{\char"1CD2}}सम्म्रा\raisebox{0.25ex}{\accentmark{15}{\char"1CD2}}जं\raisebox{0.25ex}{\accentmark{15}{\char"1CD2}}चऋषाणीनाम्\raisebox{0.25ex}{\accentmark{15}{\char"1CD2}}\kern 0.05em ।\hskip 0.12em plus 0.04em minus 0.02em इ\raisebox{0.15ex}{\accentmark{12}{\char"1CF8}}न्द्रं\raisebox{0.25ex}{\accentmark{12}{\char"0951}}स्तोतान\raisebox{0.15ex}{\accentmark{12}{\char"1CF9}}व्यंगीर्भि\raisebox{0.25ex}{\accentmark{12}{\char"0951}}ः\kern 0.05em ।\hskip 0.12em plus 0.04em minus 0.02em ना\raisebox{0.15ex}{\accentmark{12}{\char"1CF8}}रन्नृळा\raisebox{0.25ex}{\accentmark{15}{\char"1CD2}}हम्मंही\raisebox{0.25ex}{\accentmark{12}{\char"0951}}ष्ठम् \nolinebreak\hskip 0.08em \mbox{॥१०॥}}} \par}

\vspace{0.8em}

            % --- End of Kandah Closing Floral Line & Counts ---
            \par\nopagebreak\vspace{0.9em plus 0.3em minus 0.1em}
            \centerline{\textcolor{black!35}{\rule{1.8cm}{0.35pt}}\;\raisebox{-0.22ex}{\textcolor{vedicmaroon}{\large\SixFlowerPetalDotted}}\;\textcolor{black!35}{\rule{1.8cm}{0.35pt}}}
            \nopagebreak\vspace{0.45em}
            \centerline{\devafont\textbf{\textcolor{AccentBlue}{॥ ऋक् सङ्ख्या : १० ॥}}}
        
            \clearpage

            % --- SECTION ---
            
            \phantomsection  

            \section*{\centering चतुर्थ खण्डः}
            \addcontentsline{toc}{section}{चतुर्थ खण्डः} 
            \markright{चतुर्थ खण्डः}
            \nopagebreak % Keep section title with first mantra
            
            % --- SECTION-LEVEL PROCEDURE LINK ---

                % --- SUBSECTIONS (Mantras) ---
                
                % Conditional rendering based on output_mode
                \par\filbreak

\phantomsection

\addcontentsline{toc}{subsection}{ऋक् १}

{\centering \footnotesize \textcolor{AccentPurple}{॥ श्रुतकक्षः (सुकक्षः) आंगिरसः ॥} \par}

\vspace{0.12em}

{\centering \textbf{\textcolor{AccentBlue}{आपा\raisebox{0.25ex}{\accentmark{12}{\char"0951}}दूशिप्रि\raisebox{0.25ex}{\accentmark{15}{\char"1CD2}}यन्धा\raisebox{0.25ex}{\accentmark{12}{\char"0951}}साः\kern 0.05em ।\hskip 0.12em plus 0.04em minus 0.02em सूदक्षा\raisebox{0.25ex}{\accentmark{12}{\char"0951}}स्याप्राहोषीणा\raisebox{0.25ex}{\accentmark{12}{\char"0951}}ः\kern 0.05em ।\hskip 0.12em plus 0.04em minus 0.02em इ\raisebox{0.25ex}{\accentmark{15}{\char"1CD2}}न्दोरिन्द्रोया\raisebox{0.25ex}{\accentmark{15}{\char"1CD2}}वा\raisebox{0.25ex}{\accentmark{12}{\char"0951}}शिरः \nolinebreak\hskip 0.08em \mbox{॥१॥}}} \par}

\vspace{0.8em}
                % --- SUBSECTIONS (Mantras) ---
                
                % Conditional rendering based on output_mode
                \par\filbreak

\phantomsection

\addcontentsline{toc}{subsection}{ऋक् २}

{\centering \footnotesize \textcolor{AccentPurple}{॥ मेधातिथिः (शय्युबार्हस्पत्यः) काण्व ॥} \par}

\vspace{0.12em}

{\centering \textbf{\textcolor{AccentBlue}{ईमा\raisebox{0.25ex}{\accentmark{15}{\char"1CD2}}उत्वापूरोवासो\kern 0.05em ।\hskip 0.12em plus 0.04em minus 0.02em आ\raisebox{0.25ex}{\accentmark{15}{\char"1CD2}}भिप्रा\raisebox{0.25ex}{\accentmark{15}{\char"1CD2}}नो\raisebox{0.25ex}{\accentmark{12}{\char"0951}}नावुत्गीरा\raisebox{0.25ex}{\accentmark{12}{\char"0951}}ः\kern 0.05em ।\hskip 0.12em plus 0.04em minus 0.02em गा\raisebox{0.15ex}{\accentmark{12}{\char"1CF8}}वो\raisebox{0.25ex}{\accentmark{12}{\char"0951}}वत्स\raisebox{0.25ex}{\accentmark{15}{\char"1CD2}}न्ना\raisebox{0.25ex}{\accentmark{15}{\char"1CD2}}धेना\raisebox{0.25ex}{\accentmark{15}{\char"1CD2}}वा\raisebox{0.25ex}{\accentmark{12}{\char"0951}}ः \nolinebreak\hskip 0.08em \mbox{॥२॥}}} \par}

\vspace{0.8em}
                % --- SUBSECTIONS (Mantras) ---
                
                % Conditional rendering based on output_mode
                \par\filbreak

\phantomsection

\addcontentsline{toc}{subsection}{ऋक् ३}

{\centering \textbf{\textcolor{AccentBlue}{अ\raisebox{0.25ex}{\accentmark{15}{\char"1CD2}}त्रा\raisebox{0.25ex}{\accentmark{15}{\char"1CD2}}हागो\raisebox{0.25ex}{\accentmark{15}{\char"1CD2}}रा\raisebox{0.25ex}{\accentmark{12}{\char"0951}}मन्वत\kern 0.05em ।\hskip 0.12em plus 0.04em minus 0.02em नामत्वष्टू\raisebox{0.25ex}{\accentmark{12}{\char"0951}}रा\raisebox{0.15ex}{\accentmark{12}{\char"1CF9}}पीठच्यम्\kern 0.05em ।\hskip 0.12em plus 0.04em minus 0.02em इत्था\raisebox{0.25ex}{\accentmark{15}{\char"1CD2}}चन्द्रा\raisebox{0.25ex}{\accentmark{15}{\char"1CD2}}मा\raisebox{0.25ex}{\accentmark{12}{\char"0951}}सोगृ\raisebox{0.25ex}{\accentmark{15}{\char"1CD2}}हे \nolinebreak\hskip 0.08em \mbox{॥३॥}}} \par}

\vspace{0.8em}
                % --- SUBSECTIONS (Mantras) ---
                
                % Conditional rendering based on output_mode
                \par\filbreak

\phantomsection

\addcontentsline{toc}{subsection}{ऋक् ४}

{\centering \textbf{\textcolor{AccentBlue}{य\raisebox{0.25ex}{\accentmark{15}{\char"1CD2}}दि\raisebox{0.25ex}{\accentmark{15}{\char"1CD2}}न्द्रोआ\raisebox{0.15ex}{\accentmark{12}{\char"1CF9}}ना\raisebox{0.25ex}{\accentmark{12}{\char"0951}}यद्रीताः\kern 0.05em ।\hskip 0.12em plus 0.04em minus 0.02em माही\raisebox{0.25ex}{\accentmark{15}{\char"1CD2}}रापो\raisebox{0.25ex}{\accentmark{15}{\char"1CD2}}वृ\raisebox{0.25ex}{\accentmark{15}{\char"1CD2}}ष\raisebox{0.25ex}{\accentmark{12}{\char"0951}}न्तमः\kern 0.05em ।\hskip 0.12em plus 0.04em minus 0.02em तत्रा\raisebox{0.25ex}{\accentmark{12}{\char"0951}}पूषा\raisebox{0.15ex}{\accentmark{12}{\char"1CF9}}भू\raisebox{0.25ex}{\accentmark{12}{\char"0951}}वत्सा\raisebox{0.25ex}{\accentmark{15}{\char"1CD2}}चा \nolinebreak\hskip 0.08em \mbox{॥४॥}}} \par}

\vspace{0.8em}
                % --- SUBSECTIONS (Mantras) ---
                
                % Conditional rendering based on output_mode
                \par\filbreak

\phantomsection

\addcontentsline{toc}{subsection}{ऋक् ५}

{\centering \textbf{\textcolor{AccentBlue}{गौ\raisebox{0.25ex}{\accentmark{15}{\char"1CD2}}र्धायतीमारुताम्\raisebox{0.25ex}{\accentmark{12}{\char"0951}}\kern 0.05em ।\hskip 0.12em plus 0.04em minus 0.02em श्रवस्यु\raisebox{0.25ex}{\accentmark{15}{\char"1CD2}}र्मातामाघो\raisebox{0.25ex}{\accentmark{15}{\char"1CD2}}नाम्\raisebox{0.25ex}{\accentmark{12}{\char"0951}}\kern 0.05em ।\hskip 0.12em plus 0.04em minus 0.02em युक्ता\raisebox{0.25ex}{\accentmark{15}{\char"1CD2}}व\raisebox{0.25ex}{\accentmark{15}{\char"1CD2}}न्ही\raisebox{0.15ex}{\accentmark{12}{\char"1CF8}}रा\raisebox{0.25ex}{\accentmark{15}{\char"1CD2}}था\raisebox{0.25ex}{\accentmark{12}{\char"0951}}नाम् \nolinebreak\hskip 0.08em \mbox{॥५॥}}} \par}

\vspace{0.8em}
                % --- SUBSECTIONS (Mantras) ---
                
                % Conditional rendering based on output_mode
                \par\filbreak

\phantomsection

\addcontentsline{toc}{subsection}{ऋक् ६}

{\centering \textbf{\textcolor{AccentBlue}{ऊ\raisebox{0.25ex}{\accentmark{15}{\char"1CD2}}पा\raisebox{0.25ex}{\accentmark{12}{\char"0951}}नोहा\raisebox{0.25ex}{\accentmark{15}{\char"1CD2}}रीभीस्सू\raisebox{0.25ex}{\accentmark{15}{\char"1CD2}}तं\kern 0.05em ।\hskip 0.12em plus 0.04em minus 0.02em याही\raisebox{0.25ex}{\accentmark{15}{\char"1CD2}}मा\raisebox{0.25ex}{\accentmark{15}{\char"1CD2}}दानांपते\raisebox{0.25ex}{\accentmark{12}{\char"0951}}\kern 0.05em ।\hskip 0.12em plus 0.04em minus 0.02em ऊ\raisebox{0.15ex}{\accentmark{12}{\char"1CF8}}पानोहा\raisebox{0.25ex}{\accentmark{15}{\char"1CD2}}री\raisebox{0.25ex}{\accentmark{12}{\char"0951}}भीस्सू\raisebox{0.25ex}{\accentmark{15}{\char"1CD2}}तम् \nolinebreak\hskip 0.08em \mbox{॥६॥}}} \par}

\vspace{0.8em}
                % --- SUBSECTIONS (Mantras) ---
                
                % Conditional rendering based on output_mode
                \par\filbreak

\phantomsection

\addcontentsline{toc}{subsection}{ऋक् ७}

{\centering \textbf{\textcolor{AccentBlue}{इष्टाहोत्राअसृक्षत\kern 0.05em ।\hskip 0.12em plus 0.04em minus 0.02em इ\raisebox{0.15ex}{\accentmark{12}{\char"1CF8}}न्द्रं\raisebox{0.25ex}{\accentmark{12}{\char"0951}}वृधन्तो\raisebox{0.25ex}{\accentmark{12}{\char"0951}}आध्वा\raisebox{0.25ex}{\accentmark{12}{\char"0951}}रे\raisebox{0.25ex}{\accentmark{15}{\char"1CD2}}\kern 0.05em ।\hskip 0.12em plus 0.04em minus 0.02em अच्छावाभृथमो\raisebox{0.25ex}{\accentmark{15}{\char"1CD2}}ज\raisebox{0.25ex}{\accentmark{12}{\char"0951}}सा \nolinebreak\hskip 0.08em \mbox{॥७॥}}} \par}

\vspace{0.8em}
                % --- SUBSECTIONS (Mantras) ---
                
                % Conditional rendering based on output_mode
                \par\filbreak

\phantomsection

\addcontentsline{toc}{subsection}{ऋक् ८}

{\centering \textbf{\textcolor{AccentBlue}{अ\raisebox{0.25ex}{\accentmark{15}{\char"1CD2}}ह\raisebox{0.25ex}{\accentmark{15}{\char"1CD2}}मि\raisebox{0.25ex}{\accentmark{15}{\char"1CD2}}द्धि\raisebox{0.25ex}{\accentmark{15}{\char"1CD2}}पीतुःपा\raisebox{0.25ex}{\accentmark{12}{\char"0951}}री\kern 0.05em ।\hskip 0.12em plus 0.04em minus 0.02em मेधा\raisebox{0.25ex}{\accentmark{15}{\char"1CD2}}मृता\raisebox{0.15ex}{\accentmark{12}{\char"1CF9}}स्या\raisebox{0.25ex}{\accentmark{12}{\char"0951}}जग्रा\raisebox{0.25ex}{\accentmark{15}{\char"1CD2}}हा\raisebox{0.25ex}{\accentmark{12}{\char"0951}}\kern 0.05em ।\hskip 0.12em plus 0.04em minus 0.02em अहं\raisebox{0.25ex}{\accentmark{15}{\char"1CD2}}सू\raisebox{0.25ex}{\accentmark{12}{\char"0951}}र्या\raisebox{0.25ex}{\accentmark{15}{\char"1CD2}}इवा\raisebox{0.25ex}{\accentmark{15}{\char"1CD2}}जनी \nolinebreak\hskip 0.08em \mbox{॥८॥}}} \par}

\vspace{0.8em}
                % --- SUBSECTIONS (Mantras) ---
                
                % Conditional rendering based on output_mode
                \par\filbreak

\phantomsection

\addcontentsline{toc}{subsection}{ऋक् ९}

{\centering \textbf{\textcolor{AccentBlue}{रेवातीर्ण\raisebox{0.25ex}{\accentmark{12}{\char"0951}}स्साधामादे\raisebox{0.25ex}{\accentmark{12}{\char"0951}}\kern 0.05em ।\hskip 0.12em plus 0.04em minus 0.02em इन्द्रे\raisebox{0.25ex}{\accentmark{12}{\char"0951}}सन्तूतूवी\raisebox{0.25ex}{\accentmark{15}{\char"1CD2}}वा\raisebox{0.25ex}{\accentmark{12}{\char"0951}}जाः\kern 0.05em ।\hskip 0.12em plus 0.04em minus 0.02em क्षुम\raisebox{0.25ex}{\accentmark{15}{\char"1CD2}}न्तोयाभिर्मा\raisebox{0.25ex}{\accentmark{12}{\char"0951}}दे\raisebox{0.25ex}{\accentmark{12}{\char"0951}}मा \nolinebreak\hskip 0.08em \mbox{॥९॥}}} \par}

\vspace{0.8em}
                % --- SUBSECTIONS (Mantras) ---
                
                % Conditional rendering based on output_mode
                \par\filbreak

\phantomsection

\addcontentsline{toc}{subsection}{ऋक् १०}

{\centering \textbf{\textcolor{AccentBlue}{सो\raisebox{0.15ex}{\accentmark{12}{\char"1CF9}}मा\raisebox{0.25ex}{\accentmark{12}{\char"0951}}ःपूषा\raisebox{0.25ex}{\accentmark{15}{\char"1CD2}}चा\raisebox{0.25ex}{\accentmark{12}{\char"0951}}चेततुः\kern 0.05em ।\hskip 0.12em plus 0.04em minus 0.02em वि\raisebox{0.25ex}{\accentmark{15}{\char"1CD2}}श्वा\raisebox{0.15ex}{\accentmark{12}{\char"1CF8}}सांसुक्षी\raisebox{0.25ex}{\accentmark{15}{\char"1CD2}}तीनाम्\raisebox{0.25ex}{\accentmark{15}{\char"1CD2}}\kern 0.05em ।\hskip 0.12em plus 0.04em minus 0.02em देवात्रा\raisebox{0.25ex}{\accentmark{15}{\char"1CD2}}रथ्योरहित\raisebox{0.25ex}{\accentmark{12}{\char"0951}}ः \nolinebreak\hskip 0.08em \mbox{॥१०॥}}} \par}

\vspace{0.8em}

            % --- End of Kandah Closing Floral Line & Counts ---
            \par\nopagebreak\vspace{0.9em plus 0.3em minus 0.1em}
            \centerline{\textcolor{black!35}{\rule{1.8cm}{0.35pt}}\;\raisebox{-0.22ex}{\textcolor{vedicmaroon}{\large\SixFlowerPetalDotted}}\;\textcolor{black!35}{\rule{1.8cm}{0.35pt}}}
            \nopagebreak\vspace{0.45em}
            \centerline{\devafont\textbf{\textcolor{AccentBlue}{॥ ऋक् सङ्ख्या : १० ॥}}}
        
            \clearpage

            % --- SECTION ---
            
            \phantomsection  

            \section*{\centering पञ्चम खण्डः}
            \addcontentsline{toc}{section}{पञ्चम खण्डः} 
            \markright{पञ्चम खण्डः}
            \nopagebreak % Keep section title with first mantra
            
            % --- SECTION-LEVEL PROCEDURE LINK ---

                % --- SUBSECTIONS (Mantras) ---
                
                % Conditional rendering based on output_mode
                \par\filbreak

\phantomsection

\addcontentsline{toc}{subsection}{ऋक् १}

{\centering \footnotesize \textcolor{AccentPurple}{॥ श्रुतकक्षः सुकक्षो वा आंगिरसः इन्द्रः गायत्री ॥} \par}

\vspace{0.12em}

{\centering \textbf{\textcolor{AccentBlue}{पा\raisebox{0.25ex}{\accentmark{15}{\char"1CD2}}न्तामा\raisebox{0.25ex}{\accentmark{15}{\char"1CD2}}वोअ\raisebox{0.25ex}{\accentmark{15}{\char"1CD2}}न्धा\raisebox{0.25ex}{\accentmark{12}{\char"0951}}सः\kern 0.05em ।\hskip 0.12em plus 0.04em minus 0.02em इ\raisebox{0.15ex}{\accentmark{12}{\char"1CF8}}न्द्रा\raisebox{0.25ex}{\accentmark{12}{\char"0951}}माभिप्रा\raisebox{0.25ex}{\accentmark{15}{\char"1CD2}}गा\raisebox{0.25ex}{\accentmark{12}{\char"0951}}यत\kern 0.05em ।\hskip 0.12em plus 0.04em minus 0.02em विश्वा\raisebox{0.25ex}{\accentmark{15}{\char"1CD2}}सा\raisebox{0.15ex}{\accentmark{12}{\char"1CF9}}हं\raisebox{0.25ex}{\accentmark{12}{\char"0951}}शाता\raisebox{0.25ex}{\accentmark{15}{\char"1CD2}}क्रातुम्\kern 0.05em ।\hskip 0.12em plus 0.04em minus 0.02em मंहीं\raisebox{0.25ex}{\accentmark{12}{\char"0951}}ष्ठंचऋषाणीना\raisebox{0.25ex}{\accentmark{15}{\char"1CD2}}म् \nolinebreak\hskip 0.08em \mbox{॥१॥}}} \par}

\vspace{0.8em}
                % --- SUBSECTIONS (Mantras) ---
                
                % Conditional rendering based on output_mode
                \par\filbreak

\phantomsection

\addcontentsline{toc}{subsection}{ऋक् २}

{\centering \footnotesize \textcolor{AccentPurple}{॥ वसिष्ठो मैत्रावरुणिः इन्द्रः गायत्री ॥} \par}

\vspace{0.12em}

{\centering \textbf{\textcolor{AccentBlue}{प्रा\raisebox{0.25ex}{\accentmark{15}{\char"1CD2}}वाइ\raisebox{0.15ex}{\accentmark{12}{\char"1CF9}}न्द्रा\raisebox{0.25ex}{\accentmark{12}{\char"0951}}यामा\raisebox{0.25ex}{\accentmark{15}{\char"1CD2}}दा\raisebox{0.25ex}{\accentmark{12}{\char"0951}}नम्\kern 0.05em ।\hskip 0.12em plus 0.04em minus 0.02em ह\raisebox{0.25ex}{\accentmark{15}{\char"1CD2}}रीया\raisebox{0.25ex}{\accentmark{12}{\char"0951}}श्वायगायता\kern 0.05em ।\hskip 0.12em plus 0.04em minus 0.02em साखा\raisebox{0.25ex}{\accentmark{12}{\char"0951}}यस्सोमापा\raisebox{0.25ex}{\accentmark{12}{\char"0951}}व्ने \nolinebreak\hskip 0.08em \mbox{॥२॥}}} \par}

\vspace{0.8em}
                % --- SUBSECTIONS (Mantras) ---
                
                % Conditional rendering based on output_mode
                \par\filbreak

\phantomsection

\addcontentsline{toc}{subsection}{ऋक् ३}

{\centering \footnotesize \textcolor{AccentPurple}{॥ मेधातिथिः काण्वः प्रियमेधश्चांगिरसः इन्द्रः गायत्री ॥} \par}

\vspace{0.12em}

{\centering \textbf{\textcolor{AccentBlue}{वाया\raisebox{0.25ex}{\accentmark{15}{\char"1CD2}}मू\raisebox{0.25ex}{\accentmark{12}{\char"0951}}त्वातादी\raisebox{0.25ex}{\accentmark{12}{\char"0951}}दर्थाः\kern 0.05em ।\hskip 0.12em plus 0.04em minus 0.02em इ\raisebox{0.15ex}{\accentmark{12}{\char"1CF8}}न्द्रत्वायन्तस्स\raisebox{0.25ex}{\accentmark{12}{\char"0951}}खायः\kern 0.05em ।\hskip 0.12em plus 0.04em minus 0.02em क\raisebox{0.15ex}{\accentmark{12}{\char"1CF8}}ण्वा\raisebox{0.25ex}{\accentmark{12}{\char"0951}}उक्थे\raisebox{0.25ex}{\accentmark{15}{\char"1CD2}}\raisebox{0.25ex}{\accentmark{12}{\char"0951}}भी\raisebox{0.25ex}{\accentmark{12}{\char"0951}}र्ज\raisebox{0.25ex}{\accentmark{15}{\char"1CD2}}रन्ते \nolinebreak\hskip 0.08em \mbox{॥३॥}}} \par}

\vspace{0.8em}
                % --- SUBSECTIONS (Mantras) ---
                
                % Conditional rendering based on output_mode
                \par\filbreak

\phantomsection

\addcontentsline{toc}{subsection}{ऋक् ४}

{\centering \footnotesize \textcolor{AccentPurple}{॥ श्रुतकक्षः सुकक्षो वा आंगिरसः इन्द्रः गायत्री ॥} \par}

\vspace{0.12em}

{\centering \textbf{\textcolor{AccentBlue}{इ\raisebox{0.15ex}{\accentmark{12}{\char"1CF8}}न्द्राया\raisebox{0.25ex}{\accentmark{12}{\char"0951}}माद्वा\raisebox{0.25ex}{\accentmark{12}{\char"0951}}नेसूतम्\kern 0.05em ।\hskip 0.12em plus 0.04em minus 0.02em पा\raisebox{0.25ex}{\accentmark{15}{\char"1CD2}}री\raisebox{0.25ex}{\accentmark{12}{\char"0951}}ष्टोभन्तूनो\raisebox{0.25ex}{\accentmark{15}{\char"1CD2}}गी\raisebox{0.25ex}{\accentmark{15}{\char"1CD2}}राः\kern 0.05em ।\hskip 0.12em plus 0.04em minus 0.02em अर्का\raisebox{0.25ex}{\accentmark{15}{\char"1CD2}}मा\raisebox{0.25ex}{\accentmark{15}{\char"1CD2}}चन्तूका\raisebox{0.25ex}{\accentmark{15}{\char"1CD2}}रा\raisebox{0.25ex}{\accentmark{15}{\char"1CD2}}वाः \nolinebreak\hskip 0.08em \mbox{॥४॥}}} \par}

\vspace{0.8em}
                % --- SUBSECTIONS (Mantras) ---
                
                % Conditional rendering based on output_mode
                \par\filbreak

\phantomsection

\addcontentsline{toc}{subsection}{ऋक् ५}

{\centering \footnotesize \textcolor{AccentPurple}{॥ इरिम्बिठः काण्वः इन्द्रः गायत्री ॥} \par}

\vspace{0.12em}

{\centering \textbf{\textcolor{AccentBlue}{आ\raisebox{0.25ex}{\accentmark{15}{\char"1CD2}}य\raisebox{0.25ex}{\accentmark{15}{\char"1CD2}}न्ताइन्द्रासो\raisebox{0.25ex}{\accentmark{15}{\char"1CD2}}माः\kern 0.05em ।\hskip 0.12em plus 0.04em minus 0.02em नीपू\raisebox{0.25ex}{\accentmark{12}{\char"0951}}तोआ\raisebox{0.15ex}{\accentmark{12}{\char"1CF9}}धी\raisebox{0.25ex}{\accentmark{12}{\char"0951}}बर्हीषि\kern 0.05em ।\hskip 0.12em plus 0.04em minus 0.02em ए\raisebox{0.15ex}{\accentmark{12}{\char"1CF8}}ही\raisebox{0.25ex}{\accentmark{12}{\char"0951}}मास्या\raisebox{0.25ex}{\accentmark{15}{\char"1CD2}}द्रावा\raisebox{0.25ex}{\accentmark{15}{\char"1CD2}}पी\raisebox{0.25ex}{\accentmark{15}{\char"1CD2}}बा \nolinebreak\hskip 0.08em \mbox{॥५॥}}} \par}

\vspace{0.8em}
                % --- SUBSECTIONS (Mantras) ---
                
                % Conditional rendering based on output_mode
                \par\filbreak

\phantomsection

\addcontentsline{toc}{subsection}{ऋक् ६}

{\centering \footnotesize \textcolor{AccentPurple}{॥ मधुच्छन्दा वैश्वामित्रः इन्द्रः गायत्री ॥} \par}

\vspace{0.12em}

{\centering \textbf{\textcolor{AccentBlue}{सुरूपकृ\raisebox{0.25ex}{\accentmark{15}{\char"1CD2}}त्नूमूता\raisebox{0.25ex}{\accentmark{15}{\char"1CD2}}ये\kern 0.05em ।\hskip 0.12em plus 0.04em minus 0.02em सूदूघा\raisebox{0.25ex}{\accentmark{15}{\char"1CD2}}मीवागोदूहे\kern 0.05em ।\hskip 0.12em plus 0.04em minus 0.02em जुहूमा\raisebox{0.25ex}{\accentmark{15}{\char"1CD2}}सी\raisebox{0.25ex}{\accentmark{12}{\char"0951}}द्यावी\raisebox{0.25ex}{\accentmark{15}{\char"1CD2}}द्यवि \nolinebreak\hskip 0.08em \mbox{॥६॥}}} \par}

\vspace{0.8em}
                % --- SUBSECTIONS (Mantras) ---
                
                % Conditional rendering based on output_mode
                \par\filbreak

\phantomsection

\addcontentsline{toc}{subsection}{ऋक् ७}

{\centering \footnotesize \textcolor{AccentPurple}{॥ त्रिशोकः काण्वः इन्द्रः गायत्री ॥} \par}

\vspace{0.12em}

{\centering \textbf{\textcolor{AccentBlue}{आभित्वावृषभासूते\raisebox{0.25ex}{\accentmark{15}{\char"1CD2}}\kern 0.05em ।\hskip 0.12em plus 0.04em minus 0.02em सूतं\raisebox{0.25ex}{\accentmark{15}{\char"1CD2}}सृजामीपीताये\raisebox{0.25ex}{\accentmark{12}{\char"0951}}\kern 0.05em ।\hskip 0.12em plus 0.04em minus 0.02em तृ\raisebox{0.25ex}{\accentmark{15}{\char"1CD2}}म्पा\raisebox{0.25ex}{\accentmark{15}{\char"1CD2}}व्याश्नूहीमा\raisebox{0.25ex}{\accentmark{15}{\char"1CD2}}दम्\raisebox{0.25ex}{\accentmark{15}{\char"1CD2}} \nolinebreak\hskip 0.08em \mbox{॥७॥}}} \par}

\vspace{0.8em}
                % --- SUBSECTIONS (Mantras) ---
                
                % Conditional rendering based on output_mode
                \par\filbreak

\phantomsection

\addcontentsline{toc}{subsection}{ऋक् ८}

{\centering \footnotesize \textcolor{AccentPurple}{॥ कुसीदी काण्वः इन्द्रः गायत्री ॥} \par}

\vspace{0.12em}

{\centering \textbf{\textcolor{AccentBlue}{या\raisebox{0.25ex}{\accentmark{15}{\char"1CD2}}ई\raisebox{0.25ex}{\accentmark{15}{\char"1CD2}}न्द्रचा\raisebox{0.25ex}{\accentmark{15}{\char"1CD2}}मा\raisebox{0.25ex}{\accentmark{12}{\char"0951}}स्सेष्वा\kern 0.05em ।\hskip 0.12em plus 0.04em minus 0.02em सो\raisebox{0.15ex}{\accentmark{12}{\char"1CF8}}मा\raisebox{0.25ex}{\accentmark{12}{\char"0951}}श्चामू\raisebox{0.25ex}{\accentmark{15}{\char"1CD2}}षू\raisebox{0.25ex}{\accentmark{12}{\char"0951}}तेसूतः\kern 0.05em ।\hskip 0.12em plus 0.04em minus 0.02em पीबे\raisebox{0.25ex}{\accentmark{15}{\char"1CD2}}द\raisebox{0.25ex}{\accentmark{12}{\char"0951}}स्यत्वा\raisebox{0.25ex}{\accentmark{15}{\char"1CD2}}मी\raisebox{0.25ex}{\accentmark{12}{\char"0951}}शिषे \nolinebreak\hskip 0.08em \mbox{॥८॥}}} \par}

\vspace{0.8em}
                % --- SUBSECTIONS (Mantras) ---
                
                % Conditional rendering based on output_mode
                \par\filbreak

\phantomsection

\addcontentsline{toc}{subsection}{ऋक् ९}

{\centering \footnotesize \textcolor{AccentPurple}{॥ शुनः शेप आजीर्गितः इन्द्रः गायत्री ॥} \par}

\vspace{0.12em}

{\centering \textbf{\textcolor{AccentBlue}{यो\raisebox{0.25ex}{\accentmark{15}{\char"1CD2}}गे\raisebox{0.25ex}{\accentmark{12}{\char"0951}}योगेतावास्ता\raisebox{0.25ex}{\accentmark{12}{\char"0951}}रम्\kern 0.05em ।\hskip 0.12em plus 0.04em minus 0.02em वाजेवा\raisebox{0.25ex}{\accentmark{15}{\char"1CD2}}जे\raisebox{0.25ex}{\accentmark{12}{\char"0951}}हवामाहे\kern 0.05em ।\hskip 0.12em plus 0.04em minus 0.02em साखा\raisebox{0.25ex}{\accentmark{12}{\char"0951}}याइन्द्रामूता\raisebox{0.25ex}{\accentmark{15}{\char"1CD2}}ये \nolinebreak\hskip 0.08em \mbox{॥९॥}}} \par}

\vspace{0.8em}
                % --- SUBSECTIONS (Mantras) ---
                
                % Conditional rendering based on output_mode
                \par\filbreak

\phantomsection

\addcontentsline{toc}{subsection}{ऋक् १०}

{\centering \footnotesize \textcolor{AccentPurple}{॥ मधुच्छन्दा वैश्वामित्रः इन्द्रः गायत्री ॥} \par}

\vspace{0.12em}

{\centering \textbf{\textcolor{AccentBlue}{आ\raisebox{0.15ex}{\accentmark{12}{\char"1CF8}}त्वे\raisebox{0.25ex}{\accentmark{12}{\char"0951}}तानीषी\raisebox{0.25ex}{\accentmark{12}{\char"0951}}दत\kern 0.05em ।\hskip 0.12em plus 0.04em minus 0.02em इन्द्रा\raisebox{0.25ex}{\accentmark{12}{\char"0951}}माभिप्रा\raisebox{0.25ex}{\accentmark{15}{\char"1CD2}}गा\raisebox{0.25ex}{\accentmark{12}{\char"0951}}यता\kern 0.05em ।\hskip 0.12em plus 0.04em minus 0.02em सा\raisebox{0.15ex}{\accentmark{12}{\char"1CF8}}खा\raisebox{0.25ex}{\accentmark{12}{\char"0951}}यस्तो\raisebox{0.25ex}{\accentmark{15}{\char"1CD2}}मा\raisebox{0.25ex}{\accentmark{12}{\char"0951}}वाहसः \nolinebreak\hskip 0.08em \mbox{॥१०॥}}} \par}

\vspace{0.8em}

            % --- End of Kandah Closing Floral Line & Counts ---
            \par\nopagebreak\vspace{0.9em plus 0.3em minus 0.1em}
            \centerline{\textcolor{black!35}{\rule{1.8cm}{0.35pt}}\;\raisebox{-0.22ex}{\textcolor{vedicmaroon}{\large\SixFlowerPetalDotted}}\;\textcolor{black!35}{\rule{1.8cm}{0.35pt}}}
            \nopagebreak\vspace{0.45em}
            \centerline{\devafont\textbf{\textcolor{AccentBlue}{॥ ऋक् सङ्ख्या : १० ॥}}}
        
            \clearpage

            % --- SECTION ---
            
            \phantomsection  

            \section*{\centering षष्ठ खण्डः}
            \addcontentsline{toc}{section}{षष्ठ खण्डः} 
            \markright{षष्ठ खण्डः}
            \nopagebreak % Keep section title with first mantra
            
            % --- SECTION-LEVEL PROCEDURE LINK ---

                % --- SUBSECTIONS (Mantras) ---
                
                % Conditional rendering based on output_mode
                \par\filbreak

\phantomsection

\addcontentsline{toc}{subsection}{ऋक् १}

{\centering \footnotesize \textcolor{AccentPurple}{॥ विश्वामित्रो गाथिनः इन्द्रः गायत्री ॥} \par}

\vspace{0.12em}

{\centering \textbf{\textcolor{AccentBlue}{ईदं\raisebox{0.25ex}{\accentmark{15}{\char"1CD2}}ह्यन्वोज\raisebox{0.25ex}{\accentmark{12}{\char"0951}}सा\kern 0.05em ।\hskip 0.12em plus 0.04em minus 0.02em सूतं\raisebox{0.25ex}{\accentmark{15}{\char"1CD2}}रा\raisebox{0.25ex}{\accentmark{12}{\char"0951}}धानांपते\kern 0.05em ।\hskip 0.12em plus 0.04em minus 0.02em पीबा\raisebox{0.25ex}{\accentmark{15}{\char"1CD2}}त्वाटट्यस्या\raisebox{0.25ex}{\accentmark{15}{\char"1CD2}}गी\raisebox{0.25ex}{\accentmark{12}{\char"0951}}र्वणः \nolinebreak\hskip 0.08em \mbox{॥१॥}}} \par}

\vspace{0.8em}
                % --- SUBSECTIONS (Mantras) ---
                
                % Conditional rendering based on output_mode
                \par\filbreak

\phantomsection

\addcontentsline{toc}{subsection}{ऋक् २}

{\centering \footnotesize \textcolor{AccentPurple}{॥ मधुच्छन्दा वैश्वामित्रः इन्द्रः गायत्री ॥} \par}

\vspace{0.12em}

{\centering \textbf{\textcolor{AccentBlue}{माहं\raisebox{0.15ex}{\accentmark{12}{\char"1CF9}}इ\raisebox{0.15ex}{\accentmark{12}{\char"1CF9}}न्द्रा\raisebox{0.25ex}{\accentmark{12}{\char"0951}}पूराश्चा\raisebox{0.25ex}{\accentmark{12}{\char"0951}}नः\kern 0.05em ।\hskip 0.12em plus 0.04em minus 0.02em महित्वामा\raisebox{0.25ex}{\accentmark{15}{\char"1CD2}}स्तू\raisebox{0.25ex}{\accentmark{15}{\char"1CD2}}वज्री\raisebox{0.25ex}{\accentmark{15}{\char"1CD2}}णे\raisebox{0.25ex}{\accentmark{12}{\char"0951}}\kern 0.05em ।\hskip 0.12em plus 0.04em minus 0.02em द्यौर्ना\raisebox{0.25ex}{\accentmark{15}{\char"1CD2}}प्रा\raisebox{0.15ex}{\accentmark{12}{\char"1CF9}}थी\raisebox{0.25ex}{\accentmark{12}{\char"0951}}नाशा\raisebox{0.25ex}{\accentmark{15}{\char"1CD2}}वाः \nolinebreak\hskip 0.08em \mbox{॥२॥}}} \par}

\vspace{0.8em}
                % --- SUBSECTIONS (Mantras) ---
                
                % Conditional rendering based on output_mode
                \par\filbreak

\phantomsection

\addcontentsline{toc}{subsection}{ऋक् ३}

{\centering \footnotesize \textcolor{AccentPurple}{॥ कुसीदी काण्वः इन्द्रः गायत्री ॥} \par}

\vspace{0.12em}

{\centering \textbf{\textcolor{AccentBlue}{आ\raisebox{0.25ex}{\accentmark{15}{\char"1CD2}}तूना\raisebox{0.25ex}{\accentmark{12}{\char"0951}}इन्द्राक्षूम\raisebox{0.25ex}{\accentmark{15}{\char"1CD2}}न्तम्\raisebox{0.25ex}{\accentmark{12}{\char"0951}}\kern 0.05em ।\hskip 0.12em plus 0.04em minus 0.02em चित्रं\raisebox{0.25ex}{\accentmark{15}{\char"1CD2}}रा\raisebox{0.25ex}{\accentmark{12}{\char"0951}}भंस\raisebox{0.25ex}{\accentmark{15}{\char"1CD2}}ङ्गृ\raisebox{0.25ex}{\accentmark{12}{\char"0951}}भाय\kern 0.05em ।\hskip 0.12em plus 0.04em minus 0.02em माहा\raisebox{0.25ex}{\accentmark{15}{\char"1CD2}}हास्ती\raisebox{0.25ex}{\accentmark{15}{\char"1CD2}}द\raisebox{0.25ex}{\accentmark{15}{\char"1CD2}}क्षी\raisebox{0.25ex}{\accentmark{15}{\char"1CD2}}णेन \nolinebreak\hskip 0.08em \mbox{॥३॥}}} \par}

\vspace{0.8em}
                % --- SUBSECTIONS (Mantras) ---
                
                % Conditional rendering based on output_mode
                \par\filbreak

\phantomsection

\addcontentsline{toc}{subsection}{ऋक् ४}

{\centering \footnotesize \textcolor{AccentPurple}{॥ प्रियमेध आंगिरसः इन्द्रः गायत्री ॥} \par}

\vspace{0.12em}

{\centering \textbf{\textcolor{AccentBlue}{आभिप्रा\raisebox{0.15ex}{\accentmark{12}{\char"1CF9}}गो\raisebox{0.15ex}{\accentmark{12}{\char"1CF9}}पा\raisebox{0.25ex}{\accentmark{12}{\char"0951}}तिङ्गीरा\raisebox{0.25ex}{\accentmark{12}{\char"0951}}ः\kern 0.05em ।\hskip 0.12em plus 0.04em minus 0.02em इन्द्रामश्चाया\raisebox{0.15ex}{\accentmark{12}{\char"1CF9}}था\raisebox{0.25ex}{\accentmark{12}{\char"0951}}वीदे\kern 0.05em ।\hskip 0.12em plus 0.04em minus 0.02em सूनूंसत्या\raisebox{0.25ex}{\accentmark{15}{\char"1CD2}}स्यासा\raisebox{0.25ex}{\accentmark{15}{\char"1CD2}}त्पा\raisebox{0.25ex}{\accentmark{12}{\char"0951}}तिम् \nolinebreak\hskip 0.08em \mbox{॥४॥}}} \par}

\vspace{0.8em}
                % --- SUBSECTIONS (Mantras) ---
                
                % Conditional rendering based on output_mode
                \par\filbreak

\phantomsection

\addcontentsline{toc}{subsection}{ऋक् ५}

{\centering \footnotesize \textcolor{AccentPurple}{॥ वामदेवो गौतमः इन्द्रः गायत्री ॥} \par}

\vspace{0.12em}

{\centering \textbf{\textcolor{AccentBlue}{का\raisebox{0.25ex}{\accentmark{15}{\char"1CD2}}या\raisebox{0.25ex}{\accentmark{12}{\char"0951}}नश्चित्रआ\raisebox{0.25ex}{\accentmark{15}{\char"1CD2}}भु\raisebox{0.25ex}{\accentmark{12}{\char"0951}}वात्\kern 0.05em ।\hskip 0.12em plus 0.04em minus 0.02em ऊतीस्सा\raisebox{0.15ex}{\accentmark{12}{\char"1CF9}}दावृ\raisebox{0.25ex}{\accentmark{12}{\char"0951}}धस्सा\raisebox{0.25ex}{\accentmark{15}{\char"1CD2}}खा\raisebox{0.25ex}{\accentmark{12}{\char"0951}}\kern 0.05em ।\hskip 0.12em plus 0.04em minus 0.02em का\raisebox{0.25ex}{\accentmark{15}{\char"1CD2}}या\raisebox{0.25ex}{\accentmark{15}{\char"1CD2}}शाची\raisebox{0.25ex}{\accentmark{12}{\char"0951}}ष्ठयावृतः \nolinebreak\hskip 0.08em \mbox{॥५॥}}} \par}

\vspace{0.8em}
                % --- SUBSECTIONS (Mantras) ---
                
                % Conditional rendering based on output_mode
                \par\filbreak

\phantomsection

\addcontentsline{toc}{subsection}{ऋक् ६}

{\centering \footnotesize \textcolor{AccentPurple}{॥ श्रुतकक्षः सुकक्षोः वा आंगिरसः इन्द्रः गायत्री ॥} \par}

\vspace{0.12em}

{\centering \textbf{\textcolor{AccentBlue}{त्या\raisebox{0.25ex}{\accentmark{15}{\char"1CD2}}मूवस्सत्रासा\raisebox{0.25ex}{\accentmark{15}{\char"1CD2}}हम्\raisebox{0.25ex}{\accentmark{12}{\char"0951}}\kern 0.05em ।\hskip 0.12em plus 0.04em minus 0.02em विश्वा\raisebox{0.25ex}{\accentmark{12}{\char"0951}}सूगीरिष्वा\raisebox{0.25ex}{\accentmark{15}{\char"1CD2}}यातम्\kern 0.05em ।\hskip 0.12em plus 0.04em minus 0.02em आच्या\raisebox{0.25ex}{\accentmark{12}{\char"0951}}वयास्यूता\raisebox{0.25ex}{\accentmark{15}{\char"1CD2}}ये \nolinebreak\hskip 0.08em \mbox{॥६॥}}} \par}

\vspace{0.8em}
                % --- SUBSECTIONS (Mantras) ---
                
                % Conditional rendering based on output_mode
                \par\filbreak

\phantomsection

\addcontentsline{toc}{subsection}{ऋक् ७}

{\centering \footnotesize \textcolor{AccentPurple}{॥ मेधातिथिः काण्वः इन्द्रः (सदसस्पतिः) गायत्री ॥} \par}

\vspace{0.12em}

{\centering \textbf{\textcolor{AccentBlue}{साद\raisebox{0.25ex}{\accentmark{12}{\char"0951}}सस्पातीमा\raisebox{0.25ex}{\accentmark{15}{\char"1CD2}}त्भू\raisebox{0.25ex}{\accentmark{12}{\char"0951}}तम्\kern 0.05em ।\hskip 0.12em plus 0.04em minus 0.02em प्री\raisebox{0.15ex}{\accentmark{12}{\char"1CF9}}य\raisebox{0.15ex}{\accentmark{12}{\char"1CF9}}मि\raisebox{0.15ex}{\accentmark{12}{\char"1CF9}}न्द्रा\raisebox{0.25ex}{\accentmark{12}{\char"0951}}स्याकाम्यम्\raisebox{0.25ex}{\accentmark{12}{\char"0951}}\kern 0.05em ।\hskip 0.12em plus 0.04em minus 0.02em सा\raisebox{0.25ex}{\accentmark{12}{\char"0951}}निंमेधा\raisebox{0.25ex}{\accentmark{15}{\char"1CD2}}मा\raisebox{0.25ex}{\accentmark{12}{\char"0951}}या\raisebox{0.25ex}{\accentmark{15}{\char"1CD2}}सिषम् \nolinebreak\hskip 0.08em \mbox{॥७॥}}} \par}

\vspace{0.8em}
                % --- SUBSECTIONS (Mantras) ---
                
                % Conditional rendering based on output_mode
                \par\filbreak

\phantomsection

\addcontentsline{toc}{subsection}{ऋक् ८}

{\centering \footnotesize \textcolor{AccentPurple}{॥ वामदेवो गौतमः इन्द्रः गायत्री ॥} \par}

\vspace{0.12em}

{\centering \textbf{\textcolor{AccentBlue}{एतेप\raisebox{0.15ex}{\accentmark{12}{\char"1CF9}}न्थाआथो\raisebox{0.25ex}{\accentmark{12}{\char"0951}}दीवः\kern 0.05em ।\hskip 0.12em plus 0.04em minus 0.02em एभीर्व्या\raisebox{0.25ex}{\accentmark{15}{\char"1CD2}}श्वामै\raisebox{0.25ex}{\accentmark{15}{\char"1CD2}}राया\kern 0.05em ।\hskip 0.12em plus 0.04em minus 0.02em ऊता\raisebox{0.25ex}{\accentmark{15}{\char"1CD2}}श्रो\raisebox{0.25ex}{\accentmark{12}{\char"0951}}षन्तूनोभुवाः \nolinebreak\hskip 0.08em \mbox{॥८॥}}} \par}

\vspace{0.8em}
                % --- SUBSECTIONS (Mantras) ---
                
                % Conditional rendering based on output_mode
                \par\filbreak

\phantomsection

\addcontentsline{toc}{subsection}{ऋक् ९}

{\centering \footnotesize \textcolor{AccentPurple}{॥ श्रुतकक्षः सुकक्षोः वा आंगिरसः (सुकक्ष आंगिरसः) इन्द्रः गायत्री ॥} \par}

\vspace{0.12em}

{\centering \textbf{\textcolor{AccentBlue}{भ\raisebox{0.25ex}{\accentmark{15}{\char"1CD2}}द्रं\raisebox{0.25ex}{\accentmark{15}{\char"1CD2}}भ\raisebox{0.25ex}{\accentmark{12}{\char"0951}}द्रन्नआभा\raisebox{0.25ex}{\accentmark{12}{\char"0951}}रा\kern 0.05em ।\hskip 0.12em plus 0.04em minus 0.02em ई\raisebox{0.25ex}{\accentmark{15}{\char"1CD2}}ळामूर्जं\raisebox{0.25ex}{\accentmark{12}{\char"0951}}शतक्रतो\kern 0.05em ।\hskip 0.12em plus 0.04em minus 0.02em यादिन्द्रा\raisebox{0.25ex}{\accentmark{12}{\char"0951}}मृढ\raisebox{0.25ex}{\accentmark{15}{\char"1CD2}}या\raisebox{0.25ex}{\accentmark{12}{\char"0951}}सिनः \nolinebreak\hskip 0.08em \mbox{॥९॥}}} \par}

\vspace{0.8em}
                % --- SUBSECTIONS (Mantras) ---
                
                % Conditional rendering based on output_mode
                \par\filbreak

\phantomsection

\addcontentsline{toc}{subsection}{ऋक् १०}

{\centering \footnotesize \textcolor{AccentPurple}{॥ बिन्दुः पूतदक्षो वा आंगिरसः इन्द्रः (मरुतः) गायत्री ॥} \par}

\vspace{0.12em}

{\centering \textbf{\textcolor{AccentBlue}{अस्तीसो\raisebox{0.25ex}{\accentmark{15}{\char"1CD2}}मोआयं\raisebox{0.25ex}{\accentmark{15}{\char"1CD2}}सूतः\kern 0.05em ।\hskip 0.12em plus 0.04em minus 0.02em पी\raisebox{0.25ex}{\accentmark{15}{\char"1CD2}}ब\raisebox{0.25ex}{\accentmark{15}{\char"1CD2}}न्त्यस्या\raisebox{0.25ex}{\accentmark{15}{\char"1CD2}}मारुता\raisebox{0.15ex}{\accentmark{12}{\char"1CF8}}ः\kern 0.05em ।\hskip 0.12em plus 0.04em minus 0.02em उतस्वारो\raisebox{0.15ex}{\accentmark{12}{\char"1CF9}}जो\raisebox{0.25ex}{\accentmark{12}{\char"0951}}अ\raisebox{0.25ex}{\accentmark{15}{\char"1CD2}}श्वीनाः \nolinebreak\hskip 0.08em \mbox{॥१०॥}}} \par}

\vspace{0.8em}

            % --- End of Kandah Closing Floral Line & Counts ---
            \par\nopagebreak\vspace{0.9em plus 0.3em minus 0.1em}
            \centerline{\textcolor{black!35}{\rule{1.8cm}{0.35pt}}\;\raisebox{-0.22ex}{\textcolor{vedicmaroon}{\large\SixFlowerPetalDotted}}\;\textcolor{black!35}{\rule{1.8cm}{0.35pt}}}
            \nopagebreak\vspace{0.45em}
            \centerline{\devafont\textbf{\textcolor{AccentBlue}{॥ ऋक् सङ्ख्या : १० ॥}}}
        
            \clearpage

            % --- SECTION ---
            
            \phantomsection  

            \section*{\centering सप्तम खण्डः}
            \addcontentsline{toc}{section}{सप्तम खण्डः} 
            \markright{सप्तम खण्डः}
            \nopagebreak % Keep section title with first mantra
            
            % --- SECTION-LEVEL PROCEDURE LINK ---

                % --- SUBSECTIONS (Mantras) ---
                
                % Conditional rendering based on output_mode
                \par\filbreak

\phantomsection

\addcontentsline{toc}{subsection}{ऋक् १}

{\centering \footnotesize \textcolor{AccentPurple}{॥ देवजामय इन्द्रमातरः इन्द्रः गायत्री ॥} \par}

\vspace{0.12em}

{\centering \textbf{\textcolor{AccentBlue}{ईं\raisebox{0.25ex}{\accentmark{15}{\char"1CD2}}ख्या\raisebox{0.25ex}{\accentmark{15}{\char"1CD2}}यन्तीरापस्यू\raisebox{0.25ex}{\accentmark{15}{\char"1CD2}}वा\raisebox{0.25ex}{\accentmark{12}{\char"0951}}ः\kern 0.05em ।\hskip 0.12em plus 0.04em minus 0.02em इ\raisebox{0.25ex}{\accentmark{15}{\char"1CD2}}न्द्रं\raisebox{0.25ex}{\accentmark{12}{\char"0951}}जातामू\raisebox{0.25ex}{\accentmark{15}{\char"1CD2}}पासते\kern 0.05em ।\hskip 0.12em plus 0.04em minus 0.02em वन्वाना\raisebox{0.15ex}{\accentmark{12}{\char"1CF9}}सा\raisebox{0.25ex}{\accentmark{12}{\char"0951}}ःसू\raisebox{0.25ex}{\accentmark{15}{\char"1CD2}}वी\raisebox{0.25ex}{\accentmark{12}{\char"0951}}र्यम् \nolinebreak\hskip 0.08em \mbox{॥१॥}}} \par}

\vspace{0.8em}
                % --- SUBSECTIONS (Mantras) ---
                
                % Conditional rendering based on output_mode
                \par\filbreak

\phantomsection

\addcontentsline{toc}{subsection}{ऋक् २}

{\centering \footnotesize \textcolor{AccentPurple}{॥ गोधा ऋषिका इन्द्रः गायत्री ॥} \par}

\vspace{0.12em}

{\centering \textbf{\textcolor{AccentBlue}{ना\raisebox{0.25ex}{\accentmark{15}{\char"1CD2}}की\raisebox{0.25ex}{\accentmark{12}{\char"0951}}देवाइनीमसि\kern 0.05em ।\hskip 0.12em plus 0.04em minus 0.02em न\raisebox{0.25ex}{\accentmark{15}{\char"1CD2}}क्या\raisebox{0.25ex}{\accentmark{15}{\char"1CD2}}यो\raisebox{0.25ex}{\accentmark{12}{\char"0951}}पयामसि\kern 0.05em ।\hskip 0.12em plus 0.04em minus 0.02em म\raisebox{0.25ex}{\accentmark{15}{\char"1CD2}}न्द्रश्रु\raisebox{0.25ex}{\accentmark{15}{\char"1CD2}}त्य\raisebox{0.25ex}{\accentmark{12}{\char"0951}}ङ्चरामसि \nolinebreak\hskip 0.08em \mbox{॥२॥}}} \par}

\vspace{0.8em}
                % --- SUBSECTIONS (Mantras) ---
                
                % Conditional rendering based on output_mode
                \par\filbreak

\phantomsection

\addcontentsline{toc}{subsection}{ऋक् ३}

{\centering \footnotesize \textcolor{AccentPurple}{॥ दध्यङ्गाथर्वणः इन्द्रः गायत्री ॥} \par}

\vspace{0.12em}

{\centering \textbf{\textcolor{AccentBlue}{दोषो\raisebox{0.15ex}{\accentmark{12}{\char"1CF9}}आ\raisebox{0.15ex}{\accentmark{12}{\char"1CF9}}गा\raisebox{0.25ex}{\accentmark{12}{\char"0951}}त्ब\raisebox{0.25ex}{\accentmark{15}{\char"1CD2}}हा\raisebox{0.25ex}{\accentmark{15}{\char"1CD2}}त्गा\raisebox{0.25ex}{\accentmark{12}{\char"0951}}या\kern 0.05em ।\hskip 0.12em plus 0.04em minus 0.02em द्धूमा\raisebox{0.25ex}{\accentmark{12}{\char"0951}}द्गमाथर्वणा\kern 0.05em ।\hskip 0.12em plus 0.04em minus 0.02em स्तू\raisebox{0.15ex}{\accentmark{12}{\char"1CF9}}हिदे\raisebox{0.25ex}{\accentmark{12}{\char"0951}}वंसा\raisebox{0.25ex}{\accentmark{15}{\char"1CD2}}वीता\raisebox{0.25ex}{\accentmark{12}{\char"0951}}रम् \nolinebreak\hskip 0.08em \mbox{॥३॥}}} \par}

\vspace{0.8em}
                % --- SUBSECTIONS (Mantras) ---
                
                % Conditional rendering based on output_mode
                \par\filbreak

\phantomsection

\addcontentsline{toc}{subsection}{ऋक् ४}

{\centering \footnotesize \textcolor{AccentPurple}{॥ प्रस्कण्वः काण्वः इन्द्रः (अश्विनौ) गायत्री ॥} \par}

\vspace{0.12em}

{\centering \textbf{\textcolor{AccentBlue}{एषो\raisebox{0.25ex}{\accentmark{15}{\char"1CD2}}ऊषा\raisebox{0.25ex}{\accentmark{15}{\char"1CD2}}आ\raisebox{0.25ex}{\accentmark{12}{\char"0951}}पूर्व्या\kern 0.05em ।\hskip 0.12em plus 0.04em minus 0.02em व्यूं३\raisebox{0.25ex}{\accentmark{12}{\char"0951}}च्छतीप्रीया\raisebox{0.25ex}{\accentmark{15}{\char"1CD2}}दीवः\kern 0.05em ।\hskip 0.12em plus 0.04em minus 0.02em स्तू\raisebox{0.25ex}{\accentmark{15}{\char"1CD2}}षेवा\raisebox{0.25ex}{\accentmark{12}{\char"0951}}मश्विनाबृहा\raisebox{0.25ex}{\accentmark{15}{\char"1CD2}}त् \nolinebreak\hskip 0.08em \mbox{॥४॥}}} \par}

\vspace{0.8em}
                % --- SUBSECTIONS (Mantras) ---
                
                % Conditional rendering based on output_mode
                \par\filbreak

\phantomsection

\addcontentsline{toc}{subsection}{ऋक् ५}

{\centering \footnotesize \textcolor{AccentPurple}{॥ गौतमो राहूगणः इन्द्रः गायत्री ॥} \par}

\vspace{0.12em}

{\centering \textbf{\textcolor{AccentBlue}{इन्द्रो\raisebox{0.25ex}{\accentmark{12}{\char"0951}}दधीचो\raisebox{0.25ex}{\accentmark{15}{\char"1CD2}}अस्था\raisebox{0.25ex}{\accentmark{15}{\char"1CD2}}भि\raisebox{0.25ex}{\accentmark{12}{\char"0951}}ः\kern 0.05em ।\hskip 0.12em plus 0.04em minus 0.02em वृत्रा\raisebox{0.25ex}{\accentmark{15}{\char"1CD2}}ण्य\raisebox{0.25ex}{\accentmark{15}{\char"1CD2}}प्रा\raisebox{0.25ex}{\accentmark{15}{\char"1CD2}}तिष्कृतः\kern 0.05em ।\hskip 0.12em plus 0.04em minus 0.02em जघा\raisebox{0.25ex}{\accentmark{15}{\char"1CD2}}ना\raisebox{0.25ex}{\accentmark{12}{\char"0951}}नावा\raisebox{0.25ex}{\accentmark{12}{\char"0951}}तीर्नावा\raisebox{0.25ex}{\accentmark{12}{\char"0951}} \nolinebreak\hskip 0.08em \mbox{॥५॥}}} \par}

\vspace{0.8em}
                % --- SUBSECTIONS (Mantras) ---
                
                % Conditional rendering based on output_mode
                \par\filbreak

\phantomsection

\addcontentsline{toc}{subsection}{ऋक् ६}

{\centering \footnotesize \textcolor{AccentPurple}{॥ मधुच्छन्दा वैश्वामित्रः इन्द्रः गायत्री ॥} \par}

\vspace{0.12em}

{\centering \textbf{\textcolor{AccentBlue}{इ\raisebox{0.25ex}{\accentmark{15}{\char"1CD2}}न्द्रे\raisebox{0.25ex}{\accentmark{15}{\char"1CD2}}हीम\raisebox{0.25ex}{\accentmark{15}{\char"1CD2}}त्स्य\raisebox{0.25ex}{\accentmark{15}{\char"1CD2}}न्धा\raisebox{0.25ex}{\accentmark{12}{\char"0951}}सः\kern 0.05em ।\hskip 0.12em plus 0.04em minus 0.02em विश्वेभिस्सोमापा\raisebox{0.25ex}{\accentmark{15}{\char"1CD2}}वा\raisebox{0.25ex}{\accentmark{12}{\char"0951}}भीः\kern 0.05em ।\hskip 0.12em plus 0.04em minus 0.02em मा\raisebox{0.15ex}{\accentmark{12}{\char"1CF9}}हं\raisebox{0.15ex}{\accentmark{12}{\char"1CF9}}आ\raisebox{0.25ex}{\accentmark{12}{\char"0951}}भी\raisebox{0.25ex}{\accentmark{15}{\char"1CD2}}ष्ठीरोज\raisebox{0.25ex}{\accentmark{12}{\char"0951}}सा \nolinebreak\hskip 0.08em \mbox{॥६॥}}} \par}

\vspace{0.8em}
                % --- SUBSECTIONS (Mantras) ---
                
                % Conditional rendering based on output_mode
                \par\filbreak

\phantomsection

\addcontentsline{toc}{subsection}{ऋक् ७}

{\centering \footnotesize \textcolor{AccentPurple}{॥ वामदेवो गौतमः इन्द्रः गायत्री ॥} \par}

\vspace{0.12em}

{\centering \textbf{\textcolor{AccentBlue}{आ\raisebox{0.25ex}{\accentmark{15}{\char"1CD2}}तू\raisebox{0.25ex}{\accentmark{15}{\char"1CD2}}नो\raisebox{0.25ex}{\accentmark{12}{\char"0951}}इन्द्रवृत्रहन्\kern 0.05em ।\hskip 0.12em plus 0.04em minus 0.02em अस्मा\raisebox{0.25ex}{\accentmark{15}{\char"1CD2}}कामा\raisebox{0.25ex}{\accentmark{15}{\char"1CD2}}र्धामाग\raisebox{0.25ex}{\accentmark{12}{\char"0951}}ही\kern 0.05em ।\hskip 0.12em plus 0.04em minus 0.02em माहा\raisebox{0.25ex}{\accentmark{15}{\char"1CD2}}न्मा\raisebox{0.15ex}{\accentmark{12}{\char"1CF9}}ही\raisebox{0.15ex}{\accentmark{12}{\char"1CF9}}भी\raisebox{0.25ex}{\accentmark{12}{\char"0951}}रुतीभी\raisebox{0.25ex}{\accentmark{12}{\char"0951}}ः \nolinebreak\hskip 0.08em \mbox{॥७॥}}} \par}

\vspace{0.8em}
                % --- SUBSECTIONS (Mantras) ---
                
                % Conditional rendering based on output_mode
                \par\filbreak

\phantomsection

\addcontentsline{toc}{subsection}{ऋक् ८}

{\centering \footnotesize \textcolor{AccentPurple}{॥ वत्सः काण्वः इन्द्रः गायत्री ॥} \par}

\vspace{0.12em}

{\centering \textbf{\textcolor{AccentBlue}{ओजस्ता\raisebox{0.25ex}{\accentmark{15}{\char"1CD2}}द\raisebox{0.25ex}{\accentmark{12}{\char"0951}}स्यतिद्विषे\kern 0.05em ।\hskip 0.12em plus 0.04em minus 0.02em ऊ\raisebox{0.25ex}{\accentmark{15}{\char"1CD2}}भेया\raisebox{0.25ex}{\accentmark{15}{\char"1CD2}}त्सामा\raisebox{0.25ex}{\accentmark{15}{\char"1CD2}}वा\raisebox{0.25ex}{\accentmark{12}{\char"0951}}र्क्तयात्\kern 0.05em ।\hskip 0.12em plus 0.04em minus 0.02em इन्द्राश्चर्मेवा\raisebox{0.25ex}{\accentmark{12}{\char"0951}}रो\raisebox{0.25ex}{\accentmark{15}{\char"1CD2}}द\raisebox{0.25ex}{\accentmark{12}{\char"0951}}सी \nolinebreak\hskip 0.08em \mbox{॥८॥}}} \par}

\vspace{0.8em}
                % --- SUBSECTIONS (Mantras) ---
                
                % Conditional rendering based on output_mode
                \par\filbreak

\phantomsection

\addcontentsline{toc}{subsection}{ऋक् ९}

{\centering \footnotesize \textcolor{AccentPurple}{॥ शुनः शेप आजीर्गितः इन्द्रः गायत्री ॥} \par}

\vspace{0.12em}

{\centering \textbf{\textcolor{AccentBlue}{आया\raisebox{0.25ex}{\accentmark{15}{\char"1CD2}}मू\raisebox{0.25ex}{\accentmark{12}{\char"0951}}तेसा\raisebox{0.25ex}{\accentmark{15}{\char"1CD2}}मा\raisebox{0.25ex}{\accentmark{12}{\char"0951}}तसि\kern 0.05em ।\hskip 0.12em plus 0.04em minus 0.02em का\raisebox{0.25ex}{\accentmark{15}{\char"1CD2}}पो\raisebox{0.25ex}{\accentmark{12}{\char"0951}}ताइवगर्भा\raisebox{0.25ex}{\accentmark{12}{\char"0951}}धिम्\kern 0.05em ।\hskip 0.12em plus 0.04em minus 0.02em वा\raisebox{0.25ex}{\accentmark{15}{\char"1CD2}}चस्ता\raisebox{0.25ex}{\accentmark{15}{\char"1CD2}}च्चिन्नओहसे\raisebox{0.25ex}{\accentmark{15}{\char"1CD2}} \nolinebreak\hskip 0.08em \mbox{॥९॥}}} \par}

\vspace{0.8em}
                % --- SUBSECTIONS (Mantras) ---
                
                % Conditional rendering based on output_mode
                \par\filbreak

\phantomsection

\addcontentsline{toc}{subsection}{ऋक् १०}

{\centering \footnotesize \textcolor{AccentPurple}{॥ उलो वातायनः इन्द्रः (वायुः) गायत्री ॥} \par}

\vspace{0.12em}

{\centering \textbf{\textcolor{AccentBlue}{वातआ\raisebox{0.25ex}{\accentmark{15}{\char"1CD2}}वा\raisebox{0.25ex}{\accentmark{12}{\char"0951}}तूभेषाजम्\raisebox{0.25ex}{\accentmark{15}{\char"1CD2}}\kern 0.05em ।\hskip 0.12em plus 0.04em minus 0.02em शंभूमा\raisebox{0.25ex}{\accentmark{15}{\char"1CD2}}योभूनोहृदे\kern 0.05em ।\hskip 0.12em plus 0.04em minus 0.02em प्राणआयूंषि\raisebox{0.25ex}{\accentmark{12}{\char"0951}}तारि\raisebox{0.25ex}{\accentmark{12}{\char"0951}}षात् \nolinebreak\hskip 0.08em \mbox{॥१०॥}}} \par}

\vspace{0.8em}

            % --- End of Kandah Closing Floral Line & Counts ---
            \par\nopagebreak\vspace{0.9em plus 0.3em minus 0.1em}
            \centerline{\textcolor{black!35}{\rule{1.8cm}{0.35pt}}\;\raisebox{-0.22ex}{\textcolor{vedicmaroon}{\large\SixFlowerPetalDotted}}\;\textcolor{black!35}{\rule{1.8cm}{0.35pt}}}
            \nopagebreak\vspace{0.45em}
            \centerline{\devafont\textbf{\textcolor{AccentBlue}{॥ ऋक् सङ्ख्या : १० ॥}}}
        
            \clearpage

            % --- SECTION ---
            
            \phantomsection  

            \section*{\centering अष्टम खण्डः}
            \addcontentsline{toc}{section}{अष्टम खण्डः} 
            \markright{अष्टम खण्डः}
            \nopagebreak % Keep section title with first mantra
            
            % --- SECTION-LEVEL PROCEDURE LINK ---

                % --- SUBSECTIONS (Mantras) ---
                
                % Conditional rendering based on output_mode
                \par\filbreak

\phantomsection

\addcontentsline{toc}{subsection}{ऋक् १}

{\centering \footnotesize \textcolor{AccentPurple}{॥ कण्वो घौरः इन्द्रः (वरुणमित्रा{-}यर्मणः) गायत्री ॥} \par}

\vspace{0.12em}

{\centering \textbf{\textcolor{AccentBlue}{यं\raisebox{0.15ex}{\accentmark{12}{\char"1CF8}}रा\raisebox{0.25ex}{\accentmark{12}{\char"0951}}क्ष\raisebox{0.15ex}{\accentmark{12}{\char"1CF9}}न्तिप्रा\raisebox{0.25ex}{\accentmark{15}{\char"1CD2}}चे\raisebox{0.25ex}{\accentmark{12}{\char"0951}}तसःवा\raisebox{0.25ex}{\accentmark{15}{\char"1CD2}}रू\raisebox{0.25ex}{\accentmark{12}{\char"0951}}णोमित्रो\raisebox{0.15ex}{\accentmark{12}{\char"1CF9}}अर्यामा\raisebox{0.25ex}{\accentmark{15}{\char"1CD2}}\kern 0.05em ।\hskip 0.12em plus 0.04em minus 0.02em ना\raisebox{0.25ex}{\accentmark{15}{\char"1CD2}}कीस्सा\raisebox{0.25ex}{\accentmark{15}{\char"1CD2}}दभ्यातेज\raisebox{0.25ex}{\accentmark{15}{\char"1CD2}}ना\raisebox{0.25ex}{\accentmark{12}{\char"0951}}ः \nolinebreak\hskip 0.08em \mbox{॥१॥}}} \par}

\vspace{0.8em}
                % --- SUBSECTIONS (Mantras) ---
                
                % Conditional rendering based on output_mode
                \par\filbreak

\phantomsection

\addcontentsline{toc}{subsection}{ऋक् २}

{\centering \footnotesize \textcolor{AccentPurple}{॥ वत्सः (वशोऽश्वयः) काण्वः इन्द्रः गायत्री ॥} \par}

\vspace{0.12em}

{\centering \textbf{\textcolor{AccentBlue}{ग\raisebox{0.25ex}{\accentmark{15}{\char"1CD2}}व्योषू\raisebox{0.25ex}{\accentmark{15}{\char"1CD2}}णो\raisebox{0.25ex}{\accentmark{12}{\char"0951}}या\raisebox{0.15ex}{\accentmark{12}{\char"1CF9}}था\raisebox{0.25ex}{\accentmark{12}{\char"0951}}पूरा\kern 0.05em ।\hskip 0.12em plus 0.04em minus 0.02em अश्वायो\raisebox{0.15ex}{\accentmark{12}{\char"1CF9}}तारा\raisebox{0.25ex}{\accentmark{12}{\char"0951}}थाया\raisebox{0.25ex}{\accentmark{15}{\char"1CD2}}\kern 0.05em ।\hskip 0.12em plus 0.04em minus 0.02em वरिवास्या\raisebox{0.25ex}{\accentmark{15}{\char"1CD2}}माहो\raisebox{0.25ex}{\accentmark{15}{\char"1CD2}}ना\raisebox{0.25ex}{\accentmark{12}{\char"0951}}म् \nolinebreak\hskip 0.08em \mbox{॥२॥}}} \par}

\vspace{0.8em}
                % --- SUBSECTIONS (Mantras) ---
                
                % Conditional rendering based on output_mode
                \par\filbreak

\phantomsection

\addcontentsline{toc}{subsection}{ऋक् ३}

{\centering \footnotesize \textcolor{AccentPurple}{॥ वत्सः काण्वः इन्द्रः गायत्री ॥} \par}

\vspace{0.12em}

{\centering \textbf{\textcolor{AccentBlue}{ईमा\raisebox{0.25ex}{\accentmark{15}{\char"1CD2}}स्ता\raisebox{0.25ex}{\accentmark{12}{\char"0951}}इन्द्रापृ\raisebox{0.25ex}{\accentmark{15}{\char"1CD2}}श्ना\raisebox{0.25ex}{\accentmark{12}{\char"0951}}यः\kern 0.05em ।\hskip 0.12em plus 0.04em minus 0.02em घृतन्दू\raisebox{0.25ex}{\accentmark{12}{\char"0951}}हाता\raisebox{0.25ex}{\accentmark{15}{\char"1CD2}}आशी\raisebox{0.25ex}{\accentmark{15}{\char"1CD2}}रम्\kern 0.05em ।\hskip 0.12em plus 0.04em minus 0.02em एनामृ\raisebox{0.15ex}{\accentmark{12}{\char"1CF9}}ता\raisebox{0.15ex}{\accentmark{12}{\char"1CF9}}स्या\raisebox{0.25ex}{\accentmark{12}{\char"0951}}पिष्यूषीः \nolinebreak\hskip 0.08em \mbox{॥३॥}}} \par}

\vspace{0.8em}
                % --- SUBSECTIONS (Mantras) ---
                
                % Conditional rendering based on output_mode
                \par\filbreak

\phantomsection

\addcontentsline{toc}{subsection}{ऋक् ४}

{\centering \footnotesize \textcolor{AccentPurple}{॥ श्रुतकक्षः सुकक्षो वा आंगिरसः इन्द्रः गायत्री ॥} \par}

\vspace{0.12em}

{\centering \textbf{\textcolor{AccentBlue}{आ\raisebox{0.25ex}{\accentmark{15}{\char"1CD2}}या\raisebox{0.25ex}{\accentmark{15}{\char"1CD2}}धीया\raisebox{0.25ex}{\accentmark{15}{\char"1CD2}}चर्चा\raisebox{0.25ex}{\accentmark{12}{\char"0951}}गव्याया\raisebox{0.25ex}{\accentmark{15}{\char"1CD2}}\kern 0.05em ।\hskip 0.12em plus 0.04em minus 0.02em पू\raisebox{0.25ex}{\accentmark{15}{\char"1CD2}}रू\raisebox{0.25ex}{\accentmark{12}{\char"0951}}नामन्पुरूष्ठूत\kern 0.05em ।\hskip 0.12em plus 0.04em minus 0.02em यत्सो\raisebox{0.25ex}{\accentmark{15}{\char"1CD2}}मे\raisebox{0.25ex}{\accentmark{12}{\char"0951}}सोमआ\raisebox{0.25ex}{\accentmark{15}{\char"1CD2}}भू\raisebox{0.25ex}{\accentmark{12}{\char"0951}}वः \nolinebreak\hskip 0.08em \mbox{॥४॥}}} \par}

\vspace{0.8em}
                % --- SUBSECTIONS (Mantras) ---
                
                % Conditional rendering based on output_mode
                \par\filbreak

\phantomsection

\addcontentsline{toc}{subsection}{ऋक् ५}

{\centering \footnotesize \textcolor{AccentPurple}{॥ मधुच्छन्दा वैश्वामित्रः इन्द्रः गायत्री ॥} \par}

\vspace{0.12em}

{\centering \textbf{\textcolor{AccentBlue}{पवाका\raisebox{0.25ex}{\accentmark{15}{\char"1CD2}}नस्सा\raisebox{0.25ex}{\accentmark{15}{\char"1CD2}}रा\raisebox{0.25ex}{\accentmark{15}{\char"1CD2}}स्वती\kern 0.05em ।\hskip 0.12em plus 0.04em minus 0.02em वा\raisebox{0.25ex}{\accentmark{15}{\char"1CD2}}जे\raisebox{0.25ex}{\accentmark{12}{\char"0951}}भिर्वाजि\raisebox{0.25ex}{\accentmark{12}{\char"0951}}नीवति\kern 0.05em ।\hskip 0.12em plus 0.04em minus 0.02em यज्ञां\raisebox{0.25ex}{\accentmark{15}{\char"1CD2}}वा\raisebox{0.25ex}{\accentmark{12}{\char"0951}}ष्टूधीया\raisebox{0.25ex}{\accentmark{15}{\char"1CD2}}वा\raisebox{0.25ex}{\accentmark{12}{\char"0951}}सुः \nolinebreak\hskip 0.08em \mbox{॥५॥}}} \par}

\vspace{0.8em}
                % --- SUBSECTIONS (Mantras) ---
                
                % Conditional rendering based on output_mode
                \par\filbreak

\phantomsection

\addcontentsline{toc}{subsection}{ऋक् ६}

{\centering \footnotesize \textcolor{AccentPurple}{॥ वामदेवो गौतमः इन्द्रः गायत्री ॥} \par}

\vspace{0.12em}

{\centering \textbf{\textcolor{AccentBlue}{का\raisebox{0.25ex}{\accentmark{15}{\char"1CD2}}ईम\raisebox{0.15ex}{\accentmark{12}{\char"1CF9}}न्ना\raisebox{0.15ex}{\accentmark{12}{\char"1CF9}}हू\raisebox{0.25ex}{\accentmark{12}{\char"0951}}षीष्वा\kern 0.05em ।\hskip 0.12em plus 0.04em minus 0.02em इ\raisebox{0.25ex}{\accentmark{15}{\char"1CD2}}न्द्रंसो\raisebox{0.25ex}{\accentmark{15}{\char"1CD2}}मा\raisebox{0.25ex}{\accentmark{15}{\char"1CD2}}स्यतर्पयात्\kern 0.05em ।\hskip 0.12em plus 0.04em minus 0.02em सा\raisebox{0.25ex}{\accentmark{15}{\char"1CD2}}नो\raisebox{0.25ex}{\accentmark{12}{\char"0951}}वासून्या\raisebox{0.25ex}{\accentmark{15}{\char"1CD2}}भा\raisebox{0.25ex}{\accentmark{12}{\char"0951}}रात् \nolinebreak\hskip 0.08em \mbox{॥६॥}}} \par}

\vspace{0.8em}
                % --- SUBSECTIONS (Mantras) ---
                
                % Conditional rendering based on output_mode
                \par\filbreak

\phantomsection

\addcontentsline{toc}{subsection}{ऋक् ७}

{\centering \footnotesize \textcolor{AccentPurple}{॥ इरिम्बिठः काण्वः इन्द्रः गायत्री ॥} \par}

\vspace{0.12em}

{\centering \textbf{\textcolor{AccentBlue}{आया\raisebox{0.25ex}{\accentmark{12}{\char"0951}}हीसूषूमा\raisebox{0.25ex}{\accentmark{15}{\char"1CD2}}ही\raisebox{0.25ex}{\accentmark{12}{\char"0951}}ते\raisebox{0.25ex}{\accentmark{15}{\char"1CD2}}\kern 0.05em ।\hskip 0.12em plus 0.04em minus 0.02em इ\raisebox{0.25ex}{\accentmark{15}{\char"1CD2}}न्द्रासो\raisebox{0.15ex}{\accentmark{12}{\char"1CF9}}मंपीबा\raisebox{0.25ex}{\accentmark{15}{\char"1CD2}}इमम्\kern 0.05em ।\hskip 0.12em plus 0.04em minus 0.02em ए\raisebox{0.25ex}{\accentmark{15}{\char"1CD2}}दंबर्हि\raisebox{0.15ex}{\accentmark{12}{\char"1CF9}}स्सादोमा\raisebox{0.25ex}{\accentmark{15}{\char"1CD2}}मा \nolinebreak\hskip 0.08em \mbox{॥७॥}}} \par}

\vspace{0.8em}
                % --- SUBSECTIONS (Mantras) ---
                
                % Conditional rendering based on output_mode
                \par\filbreak

\phantomsection

\addcontentsline{toc}{subsection}{ऋक् ८}

{\centering \footnotesize \textcolor{AccentPurple}{॥ सत्यव्रतबाऱुणिः इन्द्रः (आदित्यः) गायत्री ॥} \par}

\vspace{0.12em}

{\centering \textbf{\textcolor{AccentBlue}{माहि\raisebox{0.25ex}{\accentmark{12}{\char"0951}}त्रीणा\raisebox{0.25ex}{\accentmark{15}{\char"1CD2}}मा\raisebox{0.25ex}{\accentmark{15}{\char"1CD2}}वा\raisebox{0.25ex}{\accentmark{12}{\char"0951}}रस्तु\kern 0.05em ।\hskip 0.12em plus 0.04em minus 0.02em धुक्षंमित्र\raisebox{0.15ex}{\accentmark{12}{\char"1CF9}}स्या\raisebox{0.25ex}{\accentmark{12}{\char"0951}}र्यम्णः\kern 0.05em ।\hskip 0.12em plus 0.04em minus 0.02em दूराध\raisebox{0.25ex}{\accentmark{15}{\char"1CD2}}ऋषंवा\raisebox{0.25ex}{\accentmark{15}{\char"1CD2}}रूणस्या \nolinebreak\hskip 0.08em \mbox{॥८॥}}} \par}

\vspace{0.8em}
                % --- SUBSECTIONS (Mantras) ---
                
                % Conditional rendering based on output_mode
                \par\filbreak

\phantomsection

\addcontentsline{toc}{subsection}{ऋक् ९}

{\centering \footnotesize \textcolor{AccentPurple}{॥ वत्सः (वशोऽश्वयः) काण्वः इन्द्रः गायत्री ॥} \par}

\vspace{0.12em}

{\centering \textbf{\textcolor{AccentBlue}{त्वा\raisebox{0.25ex}{\accentmark{15}{\char"1CD2}}वा\raisebox{0.25ex}{\accentmark{12}{\char"0951}}तःपुरोवसो\kern 0.05em ।\hskip 0.12em plus 0.04em minus 0.02em वा\raisebox{0.25ex}{\accentmark{15}{\char"1CD2}}या\raisebox{0.25ex}{\accentmark{15}{\char"1CD2}}मिन्द्र\raisebox{0.25ex}{\accentmark{12}{\char"0951}}ःप्रणेता\kern 0.05em ।\hskip 0.12em plus 0.04em minus 0.02em अस्मासी\raisebox{0.25ex}{\accentmark{12}{\char"0951}}स्थातर्हरीणाम् \nolinebreak\hskip 0.08em \mbox{॥९॥}}} \par}

\vspace{0.8em}

            % --- End of Kandah Closing Floral Line & Counts ---
            \par\nopagebreak\vspace{0.9em plus 0.3em minus 0.1em}
            \centerline{\textcolor{black!35}{\rule{1.8cm}{0.35pt}}\;\raisebox{-0.22ex}{\textcolor{vedicmaroon}{\large\SixFlowerPetalDotted}}\;\textcolor{black!35}{\rule{1.8cm}{0.35pt}}}
            \nopagebreak\vspace{0.45em}
            \centerline{\devafont\textbf{\textcolor{AccentBlue}{॥ ऋक् सङ्ख्या : ९ ॥}}}
        
            \clearpage

            % --- SECTION ---
            
            \phantomsection  

            \section*{\centering नवम खण्डः}
            \addcontentsline{toc}{section}{नवम खण्डः} 
            \markright{नवम खण्डः}
            \nopagebreak % Keep section title with first mantra
            
            % --- SECTION-LEVEL PROCEDURE LINK ---

                % --- SUBSECTIONS (Mantras) ---
                
                % Conditional rendering based on output_mode
                \par\filbreak

\phantomsection

\addcontentsline{toc}{subsection}{ऋक् १}

{\centering \footnotesize \textcolor{AccentPurple}{॥ प्रगाथः काण्वः इन्द्रः गायत्री ॥} \par}

\vspace{0.12em}

{\centering \textbf{\textcolor{AccentBlue}{ऊत्वा\raisebox{0.25ex}{\accentmark{12}{\char"0951}}मन्ततूसो\raisebox{0.25ex}{\accentmark{15}{\char"1CD2}}माः\kern 0.05em ।\hskip 0.12em plus 0.04em minus 0.02em कृणुष्वाराथोअद्रिवः\kern 0.05em ।\hskip 0.12em plus 0.04em minus 0.02em आ\raisebox{0.25ex}{\accentmark{15}{\char"1CD2}}व\raisebox{0.25ex}{\accentmark{12}{\char"0951}}ब्रह्मद्दी\raisebox{0.25ex}{\accentmark{15}{\char"1CD2}}षो\raisebox{0.25ex}{\accentmark{12}{\char"0951}}जहि \nolinebreak\hskip 0.08em \mbox{॥१॥}}} \par}

\vspace{0.8em}
                % --- SUBSECTIONS (Mantras) ---
                
                % Conditional rendering based on output_mode
                \par\filbreak

\phantomsection

\addcontentsline{toc}{subsection}{ऋक् २}

{\centering \footnotesize \textcolor{AccentPurple}{॥ विश्वामित्रो गाथिनः इन्द्रः गायत्री ॥} \par}

\vspace{0.12em}

{\centering \textbf{\textcolor{AccentBlue}{गिर्वा\raisebox{0.25ex}{\accentmark{12}{\char"0951}}णःपाही\raisebox{0.15ex}{\accentmark{12}{\char"1CF9}}ना\raisebox{0.25ex}{\accentmark{12}{\char"0951}}स्सूतम्\raisebox{0.25ex}{\accentmark{15}{\char"1CD2}}\kern 0.05em ।\hskip 0.12em plus 0.04em minus 0.02em माधोर्धारा\raisebox{0.25ex}{\accentmark{12}{\char"0951}}भीरज्यसे\kern 0.05em ।\hskip 0.12em plus 0.04em minus 0.02em इन्द्रत्वा\raisebox{0.15ex}{\accentmark{12}{\char"1CF9}}दा\raisebox{0.25ex}{\accentmark{12}{\char"0951}}तामिद्या\raisebox{0.15ex}{\accentmark{12}{\char"1CF9}}शा\raisebox{0.25ex}{\accentmark{12}{\char"0951}}ः \nolinebreak\hskip 0.08em \mbox{॥२॥}}} \par}

\vspace{0.8em}
                % --- SUBSECTIONS (Mantras) ---
                
                % Conditional rendering based on output_mode
                \par\filbreak

\phantomsection

\addcontentsline{toc}{subsection}{ऋक् ३}

{\centering \footnotesize \textcolor{AccentPurple}{॥ वामदेवो गौतमः इन्द्रः गायत्री ॥} \par}

\vspace{0.12em}

{\centering \textbf{\textcolor{AccentBlue}{सादा\raisebox{0.15ex}{\accentmark{12}{\char"1CF9}}वा\raisebox{0.25ex}{\accentmark{12}{\char"0951}}इन्द्रश्चकृ\raisebox{0.25ex}{\accentmark{15}{\char"1CD2}}षात्\kern 0.05em ।\hskip 0.12em plus 0.04em minus 0.02em ऊ\raisebox{0.25ex}{\accentmark{15}{\char"1CD2}}पोनू\raisebox{0.15ex}{\accentmark{12}{\char"1CF9}}सा\raisebox{0.15ex}{\accentmark{12}{\char"1CF9}}सा\raisebox{0.25ex}{\accentmark{12}{\char"0951}}पर्यन्\kern 0.05em ।\hskip 0.12em plus 0.04em minus 0.02em ना\raisebox{0.25ex}{\accentmark{15}{\char"1CD2}}दे\raisebox{0.25ex}{\accentmark{15}{\char"1CD2}}वा\raisebox{0.25ex}{\accentmark{15}{\char"1CD2}}ऋतश्शूराइन्द्राः \nolinebreak\hskip 0.08em \mbox{॥३॥}}} \par}

\vspace{0.8em}
                % --- SUBSECTIONS (Mantras) ---
                
                % Conditional rendering based on output_mode
                \par\filbreak

\phantomsection

\addcontentsline{toc}{subsection}{ऋक् ४}

{\centering \footnotesize \textcolor{AccentPurple}{॥ श्रुतकक्षः आंगिरसः (सुकक्षोः वा) इन्द्रः गायत्री ॥} \par}

\vspace{0.12em}

{\centering \textbf{\textcolor{AccentBlue}{आ\raisebox{0.25ex}{\accentmark{15}{\char"1CD2}}त्वा\raisebox{0.25ex}{\accentmark{12}{\char"0951}}वीशन्तू\raisebox{0.25ex}{\accentmark{12}{\char"0951}}विन्दा\raisebox{0.25ex}{\accentmark{12}{\char"0951}}वः\kern 0.05em ।\hskip 0.12em plus 0.04em minus 0.02em समुद्रा\raisebox{0.15ex}{\accentmark{12}{\char"1CF9}}मी\raisebox{0.25ex}{\accentmark{12}{\char"0951}}वासिन्धा\raisebox{0.25ex}{\accentmark{12}{\char"0951}}वः\kern 0.05em ।\hskip 0.12em plus 0.04em minus 0.02em नत्वा\raisebox{0.25ex}{\accentmark{12}{\char"0951}}मिन्द्राति\raisebox{0.25ex}{\accentmark{12}{\char"0951}}रिच्यते \nolinebreak\hskip 0.08em \mbox{॥४॥}}} \par}

\vspace{0.8em}
                % --- SUBSECTIONS (Mantras) ---
                
                % Conditional rendering based on output_mode
                \par\filbreak

\phantomsection

\addcontentsline{toc}{subsection}{ऋक् ५}

{\centering \footnotesize \textcolor{AccentPurple}{॥ मधुच्छन्दा वैश्वामित्रः इन्द्रः गायत्री ॥} \par}

\vspace{0.12em}

{\centering \textbf{\textcolor{AccentBlue}{इ\raisebox{0.25ex}{\accentmark{15}{\char"1CD2}}न्द्रा\raisebox{0.15ex}{\accentmark{12}{\char"1CF8}}मित्गा\raisebox{0.25ex}{\accentmark{12}{\char"0951}}थीनो\raisebox{0.25ex}{\accentmark{12}{\char"0951}}बृहात्\kern 0.05em ।\hskip 0.12em plus 0.04em minus 0.02em इ\raisebox{0.15ex}{\accentmark{12}{\char"1CF8}}न्द्रा\raisebox{0.25ex}{\accentmark{12}{\char"0951}}मर्केभीरर्कीणा\raisebox{0.25ex}{\accentmark{12}{\char"0951}}ः\kern 0.05em ।\hskip 0.12em plus 0.04em minus 0.02em इ\raisebox{0.25ex}{\accentmark{15}{\char"1CD2}}न्द्रंवाणी\raisebox{0.25ex}{\accentmark{12}{\char"0951}}रनूषत \nolinebreak\hskip 0.08em \mbox{॥५॥}}} \par}

\vspace{0.8em}
                % --- SUBSECTIONS (Mantras) ---
                
                % Conditional rendering based on output_mode
                \par\filbreak

\phantomsection

\addcontentsline{toc}{subsection}{ऋक् ६}

{\centering \footnotesize \textcolor{AccentPurple}{॥ श्रुतकक्षः आंगिरसः (सुकक्ष आंगिरसः) इन्द्रः गायत्री ॥} \par}

\vspace{0.12em}

{\centering \textbf{\textcolor{AccentBlue}{इन्द्राईषे\raisebox{0.25ex}{\accentmark{15}{\char"1CD2}}द\raisebox{0.25ex}{\accentmark{12}{\char"0951}}दातुनः\kern 0.05em ।\hskip 0.12em plus 0.04em minus 0.02em ऋभुक्षा\raisebox{0.15ex}{\accentmark{12}{\char"1CF9}}णा\raisebox{0.25ex}{\accentmark{12}{\char"0951}}म्भुंरायिम्\raisebox{0.25ex}{\accentmark{15}{\char"1CD2}}\kern 0.05em ।\hskip 0.12em plus 0.04em minus 0.02em वाजीद\raisebox{0.25ex}{\accentmark{12}{\char"0951}}दातूवाजी\raisebox{0.25ex}{\accentmark{15}{\char"1CD2}}नम् \nolinebreak\hskip 0.08em \mbox{॥६॥}}} \par}

\vspace{0.8em}
                % --- SUBSECTIONS (Mantras) ---
                
                % Conditional rendering based on output_mode
                \par\filbreak

\phantomsection

\addcontentsline{toc}{subsection}{ऋक् ७}

{\centering \footnotesize \textcolor{AccentPurple}{॥ गृत्समदः शौनकः इन्द्रः गायत्री ॥} \par}

\vspace{0.12em}

{\centering \textbf{\textcolor{AccentBlue}{इन्द्रोअङ्गा\raisebox{0.25ex}{\accentmark{15}{\char"1CD2}}माहा\raisebox{0.25ex}{\accentmark{15}{\char"1CD2}}त्भाय\raisebox{0.25ex}{\accentmark{15}{\char"1CD2}}म्\kern 0.05em ।\hskip 0.12em plus 0.04em minus 0.02em आभीषाद\raisebox{0.25ex}{\accentmark{15}{\char"1CD2}}पा\raisebox{0.25ex}{\accentmark{12}{\char"0951}}चुच्यवा\raisebox{0.25ex}{\accentmark{12}{\char"0951}}त्\kern 0.05em ।\hskip 0.12em plus 0.04em minus 0.02em सहिस्थीरोवीच\raisebox{0.25ex}{\accentmark{12}{\char"0951}}र्षणिः \nolinebreak\hskip 0.08em \mbox{॥७॥}}} \par}

\vspace{0.8em}
                % --- SUBSECTIONS (Mantras) ---
                
                % Conditional rendering based on output_mode
                \par\filbreak

\phantomsection

\addcontentsline{toc}{subsection}{ऋक् ८}

{\centering \textbf{\textcolor{AccentBlue}{॥ भरद्वाजः(–शयुः)बार्हस्पत्यःइन्द्रःगायत्री ॥ ॥ भरद्वाजः(–)बार्हस्पत्यःइन्द्रः(इन्द्रापषूणौ)गायत्री ॥ ईमा\raisebox{0.25ex}{\accentmark{15}{\char"1CD2}}ऊ\raisebox{0.25ex}{\accentmark{12}{\char"0951}}त्वासूते\raisebox{0.25ex}{\accentmark{15}{\char"1CD2}}सू\raisebox{0.25ex}{\accentmark{12}{\char"0951}}ते\kern 0.05em ।\hskip 0.12em plus 0.04em minus 0.02em नक्ष\raisebox{0.25ex}{\accentmark{12}{\char"0951}}न्ते\raisebox{0.25ex}{\accentmark{12}{\char"0951}}गिर्वा\raisebox{0.25ex}{\accentmark{15}{\char"1CD2}}णोगीरा\raisebox{0.25ex}{\accentmark{15}{\char"1CD2}}ः\kern 0.05em ।\hskip 0.12em plus 0.04em minus 0.02em गा\raisebox{0.15ex}{\accentmark{12}{\char"1CF8}}वो\raisebox{0.25ex}{\accentmark{12}{\char"0951}}वत्सन्नाधे\raisebox{0.25ex}{\accentmark{15}{\char"1CD2}}नावा\raisebox{0.25ex}{\accentmark{12}{\char"0951}}ः \nolinebreak\hskip 0.08em \mbox{॥८॥}}} \par}

\vspace{0.8em}
                % --- SUBSECTIONS (Mantras) ---
                
                % Conditional rendering based on output_mode
                \par\filbreak

\phantomsection

\addcontentsline{toc}{subsection}{ऋक् ९}

{\centering \textbf{\textcolor{AccentBlue}{इन्द्रानूपूषा\raisebox{0.15ex}{\accentmark{12}{\char"1CF9}}णा\raisebox{0.25ex}{\accentmark{12}{\char"0951}}वायम्\kern 0.05em ।\hskip 0.12em plus 0.04em minus 0.02em सख्या\raisebox{0.15ex}{\accentmark{12}{\char"1CF9}}या\raisebox{0.25ex}{\accentmark{12}{\char"0951}}स्वस्ताये\kern 0.05em ।\hskip 0.12em plus 0.04em minus 0.02em हूवे\raisebox{0.25ex}{\accentmark{15}{\char"1CD2}}मावा\raisebox{0.25ex}{\accentmark{15}{\char"1CD2}}ज\raisebox{0.25ex}{\accentmark{12}{\char"0951}}सातये \nolinebreak\hskip 0.08em \mbox{॥९॥}}} \par}

\vspace{0.8em}
                % --- SUBSECTIONS (Mantras) ---
                
                % Conditional rendering based on output_mode
                \par\filbreak

\phantomsection

\addcontentsline{toc}{subsection}{ऋक् १०}

{\centering \footnotesize \textcolor{AccentPurple}{॥ वामदेवो गौतमः इन्द्रः गायत्री ॥} \par}

\vspace{0.12em}

{\centering \textbf{\textcolor{AccentBlue}{ना\raisebox{0.15ex}{\accentmark{12}{\char"1CF8}}कीन्द्रात्वा\raisebox{0.25ex}{\accentmark{15}{\char"1CD2}}दू\raisebox{0.25ex}{\accentmark{15}{\char"1CD2}}क्ता\raisebox{0.25ex}{\accentmark{12}{\char"0951}}रम्\kern 0.05em ।\hskip 0.12em plus 0.04em minus 0.02em न\raisebox{0.25ex}{\accentmark{15}{\char"1CD2}}ज्या\raisebox{0.25ex}{\accentmark{15}{\char"1CD2}}यो\raisebox{0.25ex}{\accentmark{12}{\char"0951}}अस्तिवृत्रहन्\kern 0.05em ।\hskip 0.12em plus 0.04em minus 0.02em न\raisebox{0.25ex}{\accentmark{15}{\char"1CD2}}क्येवंयाथात्व\raisebox{0.25ex}{\accentmark{15}{\char"1CD2}}म् \nolinebreak\hskip 0.08em \mbox{॥१०॥}}} \par}

\vspace{0.8em}

            % --- End of Kandah Closing Floral Line & Counts ---
            \par\nopagebreak\vspace{0.9em plus 0.3em minus 0.1em}
            \centerline{\textcolor{black!35}{\rule{1.8cm}{0.35pt}}\;\raisebox{-0.22ex}{\textcolor{vedicmaroon}{\large\SixFlowerPetalDotted}}\;\textcolor{black!35}{\rule{1.8cm}{0.35pt}}}
            \nopagebreak\vspace{0.45em}
            \centerline{\devafont\textbf{\textcolor{AccentBlue}{॥ ऋक् सङ्ख्या : १० ॥}}}
        
            \clearpage

            % --- SECTION ---
            
            \phantomsection  

            \section*{\centering दशम खण्डः}
            \addcontentsline{toc}{section}{दशम खण्डः} 
            \markright{दशम खण्डः}
            \nopagebreak % Keep section title with first mantra
            
            % --- SECTION-LEVEL PROCEDURE LINK ---

                % --- SUBSECTIONS (Mantras) ---
                
                % Conditional rendering based on output_mode
                \par\filbreak

\phantomsection

\addcontentsline{toc}{subsection}{ऋक् १}

{\centering \footnotesize \textcolor{AccentPurple}{॥ त्रिशोकः काण्वः इन्द्रः गायत्री ॥} \par}

\vspace{0.12em}

{\centering \textbf{\textcolor{AccentBlue}{ता\raisebox{0.15ex}{\accentmark{12}{\char"1CF9}}रा\raisebox{0.15ex}{\accentmark{12}{\char"1CF9}}णिवोज\raisebox{0.25ex}{\accentmark{15}{\char"1CD2}}ना\raisebox{0.25ex}{\accentmark{12}{\char"0951}}नाम्\kern 0.05em ।\hskip 0.12em plus 0.04em minus 0.02em त्रादं\raisebox{0.15ex}{\accentmark{12}{\char"1CF9}}वा\raisebox{0.15ex}{\accentmark{12}{\char"1CF9}}ज\raisebox{0.25ex}{\accentmark{12}{\char"0951}}स्यागो\raisebox{0.25ex}{\accentmark{15}{\char"1CD2}}मा\raisebox{0.25ex}{\accentmark{12}{\char"0951}}तः\kern 0.05em ।\hskip 0.12em plus 0.04em minus 0.02em समाना\raisebox{0.25ex}{\accentmark{15}{\char"1CD2}}मुप्रा\raisebox{0.25ex}{\accentmark{15}{\char"1CD2}}शंसिषः \nolinebreak\hskip 0.08em \mbox{॥१॥}}} \par}

\vspace{0.8em}
                % --- SUBSECTIONS (Mantras) ---
                
                % Conditional rendering based on output_mode
                \par\filbreak

\phantomsection

\addcontentsline{toc}{subsection}{ऋक् २}

{\centering \footnotesize \textcolor{AccentPurple}{॥ मधुच्छन्दा वैश्वामित्रः इन्द्रः गायत्री ॥} \par}

\vspace{0.12em}

{\centering \textbf{\textcolor{AccentBlue}{आसृ\raisebox{0.25ex}{\accentmark{12}{\char"0951}}ग्रमिन्द्राते\raisebox{0.25ex}{\accentmark{15}{\char"1CD2}}गीरा\raisebox{0.25ex}{\accentmark{12}{\char"0951}}ः\kern 0.05em ।\hskip 0.12em plus 0.04em minus 0.02em प्रा\raisebox{0.25ex}{\accentmark{15}{\char"1CD2}}तीत्वा\raisebox{0.25ex}{\accentmark{15}{\char"1CD2}}मू\raisebox{0.25ex}{\accentmark{15}{\char"1CD2}}दा\raisebox{0.25ex}{\accentmark{12}{\char"0951}}हासता\kern 0.05em ।\hskip 0.12em plus 0.04em minus 0.02em सा\raisebox{0.25ex}{\accentmark{15}{\char"1CD2}}जो\raisebox{0.25ex}{\accentmark{15}{\char"1CD2}}षा\raisebox{0.25ex}{\accentmark{12}{\char"0951}}वृषा\raisebox{0.15ex}{\accentmark{12}{\char"1CF8}}भं\raisebox{0.25ex}{\accentmark{15}{\char"1CD2}}पा\raisebox{0.25ex}{\accentmark{15}{\char"1CD2}}तिम् \nolinebreak\hskip 0.08em \mbox{॥२॥}}} \par}

\vspace{0.8em}
                % --- SUBSECTIONS (Mantras) ---
                
                % Conditional rendering based on output_mode
                \par\filbreak

\phantomsection

\addcontentsline{toc}{subsection}{ऋक् ३}

{\centering \footnotesize \textcolor{AccentPurple}{॥ वत्सः काण्वः (वशोऽश्वयः) इन्द्रः गायत्री ॥} \par}

\vspace{0.12em}

{\centering \textbf{\textcolor{AccentBlue}{सुनीथोघासामार्त्या\raisebox{0.25ex}{\accentmark{12}{\char"0951}}ः\kern 0.05em ।\hskip 0.12em plus 0.04em minus 0.02em य\raisebox{0.25ex}{\accentmark{15}{\char"1CD2}}म्मा\raisebox{0.15ex}{\accentmark{12}{\char"1CF8}}रूतोया\raisebox{0.15ex}{\accentmark{12}{\char"1CF9}}मर्यामा\kern 0.05em ।\hskip 0.12em plus 0.04em minus 0.02em मित्रःपा\raisebox{0.25ex}{\accentmark{15}{\char"1CD2}}न्त्यद्रू\raisebox{0.25ex}{\accentmark{15}{\char"1CD2}}होआ\raisebox{0.25ex}{\accentmark{15}{\char"1CD2}}तीद्वीषा\raisebox{0.25ex}{\accentmark{12}{\char"0951}}ः \nolinebreak\hskip 0.08em \mbox{॥३॥}}} \par}

\vspace{0.8em}
                % --- SUBSECTIONS (Mantras) ---
                
                % Conditional rendering based on output_mode
                \par\filbreak

\phantomsection

\addcontentsline{toc}{subsection}{ऋक् ४}

{\centering \footnotesize \textcolor{AccentPurple}{॥ त्रिशोकः काण्वः इन्द्रः गायत्री ॥} \par}

\vspace{0.12em}

{\centering \textbf{\textcolor{AccentBlue}{य\raisebox{0.25ex}{\accentmark{15}{\char"1CD2}}द्वीळा\raisebox{0.15ex}{\accentmark{12}{\char"1CF9}}वि\raisebox{0.25ex}{\accentmark{12}{\char"0951}}न्द्राया\raisebox{0.25ex}{\accentmark{15}{\char"1CD2}}त्स्थि\raisebox{0.25ex}{\accentmark{12}{\char"0951}}रे\kern 0.05em ।\hskip 0.12em plus 0.04em minus 0.02em य\raisebox{0.15ex}{\accentmark{12}{\char"1CF8}}त्प\raisebox{0.15ex}{\accentmark{12}{\char"1CF8}}र्शा\raisebox{0.25ex}{\accentmark{12}{\char"0951}}नेपा\raisebox{0.25ex}{\accentmark{15}{\char"1CD2}}रा\raisebox{0.25ex}{\accentmark{12}{\char"0951}}भृतम्\kern 0.05em ।\hskip 0.12em plus 0.04em minus 0.02em वा\raisebox{0.15ex}{\accentmark{12}{\char"1CF8}}सु\raisebox{0.25ex}{\accentmark{12}{\char"0951}}स्पार्हं\raisebox{0.25ex}{\accentmark{12}{\char"0951}}ता\raisebox{0.25ex}{\accentmark{15}{\char"1CD2}}दा\raisebox{0.25ex}{\accentmark{15}{\char"1CD2}}भा\raisebox{0.25ex}{\accentmark{12}{\char"0951}}रा \nolinebreak\hskip 0.08em \mbox{॥४॥}}} \par}

\vspace{0.8em}
                % --- SUBSECTIONS (Mantras) ---
                
                % Conditional rendering based on output_mode
                \par\filbreak

\phantomsection

\addcontentsline{toc}{subsection}{ऋक् ५}

{\centering \footnotesize \textcolor{AccentPurple}{॥ सुकक्ष आंगिरसः इन्द्रः गायत्री ॥} \par}

\vspace{0.12em}

{\centering \textbf{\textcolor{AccentBlue}{श्रूतं\raisebox{0.25ex}{\accentmark{15}{\char"1CD2}}वो\raisebox{0.25ex}{\accentmark{12}{\char"0951}}वृत्रा\raisebox{0.25ex}{\accentmark{12}{\char"0951}}ह\raisebox{0.25ex}{\accentmark{15}{\char"1CD2}}न्ता\raisebox{0.25ex}{\accentmark{12}{\char"0951}}मम्\kern 0.05em ।\hskip 0.12em plus 0.04em minus 0.02em प्र\raisebox{0.25ex}{\accentmark{15}{\char"1CD2}}श\raisebox{0.25ex}{\accentmark{15}{\char"1CD2}}र्धञ्चर्षो\raisebox{0.25ex}{\accentmark{12}{\char"0951}}णीनाम्\kern 0.05em ।\hskip 0.12em plus 0.04em minus 0.02em आ\raisebox{0.25ex}{\accentmark{15}{\char"1CD2}}शी\raisebox{0.25ex}{\accentmark{15}{\char"1CD2}}षेरा\raisebox{0.25ex}{\accentmark{15}{\char"1CD2}}धा\raisebox{0.25ex}{\accentmark{12}{\char"0951}}सेमाहे \nolinebreak\hskip 0.08em \mbox{॥५॥}}} \par}

\vspace{0.8em}
                % --- SUBSECTIONS (Mantras) ---
                
                % Conditional rendering based on output_mode
                \par\filbreak

\phantomsection

\addcontentsline{toc}{subsection}{ऋक् ६}

{\centering \footnotesize \textcolor{AccentPurple}{॥ वामदेवो गौतमः इन्द्रः गायत्री ॥} \par}

\vspace{0.12em}

{\centering \textbf{\textcolor{AccentBlue}{आर\raisebox{0.25ex}{\accentmark{12}{\char"0951}}न्तइन्द्राश्रा\raisebox{0.25ex}{\accentmark{15}{\char"1CD2}}वा\raisebox{0.25ex}{\accentmark{12}{\char"0951}}से\kern 0.05em ।\hskip 0.12em plus 0.04em minus 0.02em गमे\raisebox{0.25ex}{\accentmark{15}{\char"1CD2}}मा\raisebox{0.25ex}{\accentmark{12}{\char"0951}}शूरत्वा\raisebox{0.25ex}{\accentmark{15}{\char"1CD2}}वातः\kern 0.05em ।\hskip 0.12em plus 0.04em minus 0.02em आ\raisebox{0.25ex}{\accentmark{15}{\char"1CD2}}रं\raisebox{0.25ex}{\accentmark{12}{\char"0951}}शाक्रापा\raisebox{0.25ex}{\accentmark{15}{\char"1CD2}}रेमणि \nolinebreak\hskip 0.08em \mbox{॥६॥}}} \par}

\vspace{0.8em}
                % --- SUBSECTIONS (Mantras) ---
                
                % Conditional rendering based on output_mode
                \par\filbreak

\phantomsection

\addcontentsline{toc}{subsection}{ऋक् ७}

{\centering \footnotesize \textcolor{AccentPurple}{॥ विश्वामित्रो गाथिनः इन्द्रः गायत्री ॥} \par}

\vspace{0.12em}

{\centering \textbf{\textcolor{AccentBlue}{धानाव\raisebox{0.25ex}{\accentmark{12}{\char"0951}}न्तंकारंभीणम्\kern 0.05em ।\hskip 0.12em plus 0.04em minus 0.02em आ\raisebox{0.25ex}{\accentmark{15}{\char"1CD2}}पू\raisebox{0.25ex}{\accentmark{15}{\char"1CD2}}पाव\raisebox{0.25ex}{\accentmark{12}{\char"0951}}न्तामुक्थी\raisebox{0.25ex}{\accentmark{15}{\char"1CD2}}न\raisebox{0.25ex}{\accentmark{12}{\char"0951}}म्\kern 0.05em ।\hskip 0.12em plus 0.04em minus 0.02em इ\raisebox{0.15ex}{\accentmark{12}{\char"1CF8}}न्द्राप्राता\raisebox{0.25ex}{\accentmark{15}{\char"1CD2}}र्जु\raisebox{0.25ex}{\accentmark{12}{\char"0951}}षस्वनः \nolinebreak\hskip 0.08em \mbox{॥७॥}}} \par}

\vspace{0.8em}
                % --- SUBSECTIONS (Mantras) ---
                
                % Conditional rendering based on output_mode
                \par\filbreak

\phantomsection

\addcontentsline{toc}{subsection}{ऋक् ८}

{\centering \footnotesize \textcolor{AccentPurple}{॥ गोषूकश्वसूनूकाण्वायनाै इन्द्रः गायत्री ॥} \par}

\vspace{0.12em}

{\centering \textbf{\textcolor{AccentBlue}{आ\raisebox{0.25ex}{\accentmark{15}{\char"1CD2}}पांफे\raisebox{0.15ex}{\accentmark{12}{\char"1CF9}}ने\raisebox{0.25ex}{\accentmark{12}{\char"0951}}नानामू\raisebox{0.25ex}{\accentmark{12}{\char"0951}}चेः\kern 0.05em ।\hskip 0.12em plus 0.04em minus 0.02em शी\raisebox{0.25ex}{\accentmark{12}{\char"0951}}रा\raisebox{0.25ex}{\accentmark{12}{\char"0951}}इ\raisebox{0.25ex}{\accentmark{12}{\char"0951}}न्द्रोद\raisebox{0.25ex}{\accentmark{12}{\char"0951}}वर्क्तया\kern 0.05em ।\hskip 0.12em plus 0.04em minus 0.02em वि\raisebox{0.25ex}{\accentmark{15}{\char"1CD2}}श्वा\raisebox{0.25ex}{\accentmark{15}{\char"1CD2}}या\raisebox{0.15ex}{\accentmark{12}{\char"1CF9}}दजयस्पृधा\raisebox{0.25ex}{\accentmark{12}{\char"0951}}ः \nolinebreak\hskip 0.08em \mbox{॥८॥}}} \par}

\vspace{0.8em}
                % --- SUBSECTIONS (Mantras) ---
                
                % Conditional rendering based on output_mode
                \par\filbreak

\phantomsection

\addcontentsline{toc}{subsection}{ऋक् ९}

{\centering \textbf{\textcolor{AccentBlue}{॥ वामदेवोगौतमःइन्द्रःगायत्री ॥ ॥ श्रुतकक्षःसुकक्षोवाआंगिरसःइन्द्रःगायत्री ॥ ईमे\raisebox{0.25ex}{\accentmark{15}{\char"1CD2}}ताइन्द्रासो\raisebox{0.25ex}{\accentmark{15}{\char"1CD2}}माः\kern 0.05em ।\hskip 0.12em plus 0.04em minus 0.02em सूता\raisebox{0.25ex}{\accentmark{15}{\char"1CD2}}सोयेचासो\raisebox{0.25ex}{\accentmark{15}{\char"1CD2}}त्वा\raisebox{0.25ex}{\accentmark{12}{\char"0951}}ः\kern 0.05em ।\hskip 0.12em plus 0.04em minus 0.02em तेषांमत्स्वप्रभूवसो \nolinebreak\hskip 0.08em \mbox{॥९॥}}} \par}

\vspace{0.8em}
                % --- SUBSECTIONS (Mantras) ---
                
                % Conditional rendering based on output_mode
                \par\filbreak

\phantomsection

\addcontentsline{toc}{subsection}{ऋक् १०}

{\centering \textbf{\textcolor{AccentBlue}{तुभ्यंसूतासस्सोमा\raisebox{0.25ex}{\accentmark{12}{\char"0951}}ः\kern 0.05em ।\hskip 0.12em plus 0.04em minus 0.02em स्तीर्ण\raisebox{0.25ex}{\accentmark{12}{\char"0951}}म्बर्ही\raisebox{0.25ex}{\accentmark{12}{\char"0951}}र्वी\raisebox{0.25ex}{\accentmark{12}{\char"0951}}भवसो\kern 0.05em ।\hskip 0.12em plus 0.04em minus 0.02em स्तो\raisebox{0.25ex}{\accentmark{15}{\char"1CD2}}तृभ्या\raisebox{0.25ex}{\accentmark{12}{\char"0951}}इन्द्रमृडया \nolinebreak\hskip 0.08em \mbox{॥१०॥}}} \par}

\vspace{0.8em}

            % --- End of Kandah Closing Floral Line & Counts ---
            \par\nopagebreak\vspace{0.9em plus 0.3em minus 0.1em}
            \centerline{\textcolor{black!35}{\rule{1.8cm}{0.35pt}}\;\raisebox{-0.22ex}{\textcolor{vedicmaroon}{\large\SixFlowerPetalDotted}}\;\textcolor{black!35}{\rule{1.8cm}{0.35pt}}}
            \nopagebreak\vspace{0.45em}
            \centerline{\devafont\textbf{\textcolor{AccentBlue}{॥ ऋक् सङ्ख्या : १० ॥}}}
        
            \clearpage

            % --- SECTION ---
            
            \phantomsection  

            \section*{\centering एकादश खण्डः}
            \addcontentsline{toc}{section}{एकादश खण्डः} 
            \markright{एकादश खण्डः}
            \nopagebreak % Keep section title with first mantra
            
            % --- SECTION-LEVEL PROCEDURE LINK ---

                % --- SUBSECTIONS (Mantras) ---
                
                % Conditional rendering based on output_mode
                \par\filbreak

\phantomsection

\addcontentsline{toc}{subsection}{ऋक् १}

{\centering \footnotesize \textcolor{AccentPurple}{॥ शुनः शेप आजीर्गितः इन्द्रः गायत्री ॥} \par}

\vspace{0.12em}

{\centering \textbf{\textcolor{AccentBlue}{आ\raisebox{0.25ex}{\accentmark{15}{\char"1CD2}}वाइन्द्रङ्कृवींया\raisebox{0.25ex}{\accentmark{15}{\char"1CD2}}था\raisebox{0.25ex}{\accentmark{12}{\char"0951}}\kern 0.05em ।\hskip 0.12em plus 0.04em minus 0.02em वजय\raisebox{0.25ex}{\accentmark{15}{\char"1CD2}}न्ताश्शा\raisebox{0.25ex}{\accentmark{12}{\char"0951}}ताक्रातुम्\kern 0.05em ।\hskip 0.12em plus 0.04em minus 0.02em मंहीं\raisebox{0.25ex}{\accentmark{12}{\char"0951}}ष्ठंसिञ्चाइन्दुभिः \nolinebreak\hskip 0.08em \mbox{॥१॥}}} \par}

\vspace{0.8em}
                % --- SUBSECTIONS (Mantras) ---
                
                % Conditional rendering based on output_mode
                \par\filbreak

\phantomsection

\addcontentsline{toc}{subsection}{ऋक् २}

{\centering \footnotesize \textcolor{AccentPurple}{॥ श्रुतकक्ष आंगिरसः (सुकक्षो आंगिरसो वा) इन्द्रः गायत्री ॥} \par}

\vspace{0.12em}

{\centering \textbf{\textcolor{AccentBlue}{आताश्चिदिन्द्राना\raisebox{0.25ex}{\accentmark{15}{\char"1CD2}}ऊ\raisebox{0.25ex}{\accentmark{15}{\char"1CD2}}पा\raisebox{0.25ex}{\accentmark{12}{\char"0951}}\kern 0.05em ।\hskip 0.12em plus 0.04em minus 0.02em आ\raisebox{0.25ex}{\accentmark{15}{\char"1CD2}}या\raisebox{0.25ex}{\accentmark{12}{\char"0951}}हीसातावा\raisebox{0.25ex}{\accentmark{12}{\char"0951}}जया\kern 0.05em ।\hskip 0.12em plus 0.04em minus 0.02em इळा\raisebox{0.25ex}{\accentmark{15}{\char"1CD2}}साहा\raisebox{0.25ex}{\accentmark{15}{\char"1CD2}}स्रा\raisebox{0.25ex}{\accentmark{12}{\char"0951}}वाजया \nolinebreak\hskip 0.08em \mbox{॥२॥}}} \par}

\vspace{0.8em}
                % --- SUBSECTIONS (Mantras) ---
                
                % Conditional rendering based on output_mode
                \par\filbreak

\phantomsection

\addcontentsline{toc}{subsection}{ऋक् ३}

{\centering \footnotesize \textcolor{AccentPurple}{॥ त्रिशोकः काण्वः इन्द्रः गायत्री ॥} \par}

\vspace{0.12em}

{\centering \textbf{\textcolor{AccentBlue}{आबुन्दं\raisebox{0.15ex}{\accentmark{12}{\char"1CF9}}वृ\raisebox{0.15ex}{\accentmark{12}{\char"1CF9}}त्राहाददे\kern 0.05em ।\hskip 0.12em plus 0.04em minus 0.02em जात\raisebox{0.15ex}{\accentmark{12}{\char"1CF9}}ःप्र\raisebox{0.25ex}{\accentmark{12}{\char"0951}}च्छद्वीमाता\raisebox{0.25ex}{\accentmark{15}{\char"1CD2}}रम्\kern 0.05em ।\hskip 0.12em plus 0.04em minus 0.02em का\raisebox{0.25ex}{\accentmark{15}{\char"1CD2}}उग्रःकेहा\raisebox{0.25ex}{\accentmark{12}{\char"0951}}शृण्विरे\raisebox{0.25ex}{\accentmark{15}{\char"1CD2}} \nolinebreak\hskip 0.08em \mbox{॥३॥}}} \par}

\vspace{0.8em}
                % --- SUBSECTIONS (Mantras) ---
                
                % Conditional rendering based on output_mode
                \par\filbreak

\phantomsection

\addcontentsline{toc}{subsection}{ऋक् ४}

{\centering \footnotesize \textcolor{AccentPurple}{॥ मेधातिथिः काण्वः इन्द्रः गायत्री ॥} \par}

\vspace{0.12em}

{\centering \textbf{\textcolor{AccentBlue}{बृबदू\raisebox{0.25ex}{\accentmark{12}{\char"0951}}क्थंहवा\raisebox{0.25ex}{\accentmark{12}{\char"0951}}महे\kern 0.05em ।\hskip 0.12em plus 0.04em minus 0.02em सृप्रा\raisebox{0.25ex}{\accentmark{15}{\char"1CD2}}का\raisebox{0.25ex}{\accentmark{12}{\char"0951}}रस्नामू\raisebox{0.25ex}{\accentmark{15}{\char"1CD2}}ता\raisebox{0.25ex}{\accentmark{15}{\char"1CD2}}ये\kern 0.05em ।\hskip 0.12em plus 0.04em minus 0.02em सा\raisebox{0.15ex}{\accentmark{12}{\char"1CF8}}धा\raisebox{0.25ex}{\accentmark{12}{\char"0951}}ःकृण्व\raisebox{0.25ex}{\accentmark{15}{\char"1CD2}}न्ता\raisebox{0.25ex}{\accentmark{15}{\char"1CD2}}मा\raisebox{0.25ex}{\accentmark{15}{\char"1CD2}}वा\raisebox{0.25ex}{\accentmark{12}{\char"0951}}से \nolinebreak\hskip 0.08em \mbox{॥४॥}}} \par}

\vspace{0.8em}
                % --- SUBSECTIONS (Mantras) ---
                
                % Conditional rendering based on output_mode
                \par\filbreak

\phantomsection

\addcontentsline{toc}{subsection}{ऋक् ५}

{\centering \footnotesize \textcolor{AccentPurple}{॥ गौतमो राहूगणः इन्द्रः (विश्वे देवाः) गायत्री ॥} \par}

\vspace{0.12em}

{\centering \textbf{\textcolor{AccentBlue}{ऋजु\raisebox{0.25ex}{\accentmark{12}{\char"0951}}नीतीनोवारू\raisebox{0.25ex}{\accentmark{12}{\char"0951}}णः\kern 0.05em ।\hskip 0.12em plus 0.04em minus 0.02em मि\raisebox{0.25ex}{\accentmark{15}{\char"1CD2}}त्रो\raisebox{0.25ex}{\accentmark{15}{\char"1CD2}}नायति\raisebox{0.25ex}{\accentmark{15}{\char"1CD2}}विद्वांसः\kern 0.05em ।\hskip 0.12em plus 0.04em minus 0.02em अर्यामा\raisebox{0.25ex}{\accentmark{12}{\char"0951}}देवै\raisebox{0.25ex}{\accentmark{15}{\char"1CD2}}स्साजो\raisebox{0.25ex}{\accentmark{15}{\char"1CD2}}षाः \nolinebreak\hskip 0.08em \mbox{॥५॥}}} \par}

\vspace{0.8em}
                % --- SUBSECTIONS (Mantras) ---
                
                % Conditional rendering based on output_mode
                \par\filbreak

\phantomsection

\addcontentsline{toc}{subsection}{ऋक् ६}

{\centering \footnotesize \textcolor{AccentPurple}{॥ ब्रहमातिथिः काण्वः अश्विनौ मित्रावरुणौ गायत्री ॥} \par}

\vspace{0.12em}

{\centering \textbf{\textcolor{AccentBlue}{दूरा\raisebox{0.25ex}{\accentmark{15}{\char"1CD2}}दीहे\raisebox{0.25ex}{\accentmark{15}{\char"1CD2}}वायात्सातः\kern 0.05em ।\hskip 0.12em plus 0.04em minus 0.02em आरूणा\raisebox{0.25ex}{\accentmark{15}{\char"1CD2}}स्पूरा\raisebox{0.25ex}{\accentmark{15}{\char"1CD2}}शिश्वीतात्\kern 0.05em ।\hskip 0.12em plus 0.04em minus 0.02em विभा\raisebox{0.15ex}{\accentmark{12}{\char"1CF9}}नूंविश्वाथातनात् \nolinebreak\hskip 0.08em \mbox{॥६॥}}} \par}

\vspace{0.8em}
                % --- SUBSECTIONS (Mantras) ---
                
                % Conditional rendering based on output_mode
                \par\filbreak

\phantomsection

\addcontentsline{toc}{subsection}{ऋक् ७}

{\centering \footnotesize \textcolor{AccentPurple}{॥ विश्वामित्रो गाथिनो जमदग्निवा इन्द्रः गायत्री ॥} \par}

\vspace{0.12em}

{\centering \textbf{\textcolor{AccentBlue}{आनो\raisebox{0.25ex}{\accentmark{12}{\char"0951}}मित्रावरुणा\kern 0.05em ।\hskip 0.12em plus 0.04em minus 0.02em घृतै\raisebox{0.25ex}{\accentmark{15}{\char"1CD2}}ःगव्यूतिमुक्षतम्\kern 0.05em ।\hskip 0.12em plus 0.04em minus 0.02em मध्वारा\raisebox{0.25ex}{\accentmark{15}{\char"1CD2}}जां\raisebox{0.25ex}{\accentmark{12}{\char"0951}}सिसुक्रतू \nolinebreak\hskip 0.08em \mbox{॥७॥}}} \par}

\vspace{0.8em}
                % --- SUBSECTIONS (Mantras) ---
                
                % Conditional rendering based on output_mode
                \par\filbreak

\phantomsection

\addcontentsline{toc}{subsection}{ऋक् ८}

{\centering \footnotesize \textcolor{AccentPurple}{॥ प्रस्कण्वः काण्वः (कण्वो घौराः) मरुतः गायत्री ॥} \par}

\vspace{0.12em}

{\centering \textbf{\textcolor{AccentBlue}{ऊ\raisebox{0.25ex}{\accentmark{15}{\char"1CD2}}दूत्ये\raisebox{0.25ex}{\accentmark{15}{\char"1CD2}}सूना\raisebox{0.25ex}{\accentmark{15}{\char"1CD2}}वोगी\raisebox{0.25ex}{\accentmark{15}{\char"1CD2}}रा\raisebox{0.25ex}{\accentmark{12}{\char"0951}}ः\kern 0.05em ।\hskip 0.12em plus 0.04em minus 0.02em काष्ठार्यज्ञेष्वात्नता\kern 0.05em ।\hskip 0.12em plus 0.04em minus 0.02em वा\raisebox{0.25ex}{\accentmark{15}{\char"1CD2}}श्रा\raisebox{0.15ex}{\accentmark{12}{\char"1CF9}}आ\raisebox{0.25ex}{\accentmark{12}{\char"0951}}भीजुया\raisebox{0.25ex}{\accentmark{15}{\char"1CD2}}ता\raisebox{0.25ex}{\accentmark{12}{\char"0951}}वे \nolinebreak\hskip 0.08em \mbox{॥८॥}}} \par}

\vspace{0.8em}
                % --- SUBSECTIONS (Mantras) ---
                
                % Conditional rendering based on output_mode
                \par\filbreak

\phantomsection

\addcontentsline{toc}{subsection}{ऋक् ९}

{\centering \footnotesize \textcolor{AccentPurple}{॥ मेधातिथिः काण्वः विष्णुः गायत्री ॥} \par}

\vspace{0.12em}

{\centering \textbf{\textcolor{AccentBlue}{ईदंविष्णुर्विचा\raisebox{0.25ex}{\accentmark{12}{\char"0951}}क्रमे\kern 0.05em ।\hskip 0.12em plus 0.04em minus 0.02em त्रेधा\raisebox{0.25ex}{\accentmark{15}{\char"1CD2}}नी\raisebox{0.25ex}{\accentmark{15}{\char"1CD2}}दे\raisebox{0.25ex}{\accentmark{12}{\char"0951}}थेपादम्\kern 0.05em ।\hskip 0.12em plus 0.04em minus 0.02em सा\raisebox{0.25ex}{\accentmark{15}{\char"1CD2}}मू\raisebox{0.25ex}{\accentmark{12}{\char"0951}}ढमस्यपांसूले\raisebox{0.25ex}{\accentmark{15}{\char"1CD2}} \nolinebreak\hskip 0.08em \mbox{॥९॥}}} \par}

\vspace{0.8em}

            % --- End of Kandah Closing Floral Line & Counts ---
            \par\nopagebreak\vspace{0.9em plus 0.3em minus 0.1em}
            \centerline{\textcolor{black!35}{\rule{1.8cm}{0.35pt}}\;\raisebox{-0.22ex}{\textcolor{vedicmaroon}{\large\SixFlowerPetalDotted}}\;\textcolor{black!35}{\rule{1.8cm}{0.35pt}}}
            \nopagebreak\vspace{0.45em}
            \centerline{\devafont\textbf{\textcolor{AccentBlue}{॥ ऋक् सङ्ख्या : ९ ॥}}}
        
            \clearpage

            % --- SECTION ---
            
            \phantomsection  

            \section*{\centering द्वादश खण्डः}
            \addcontentsline{toc}{section}{द्वादश खण्डः} 
            \markright{द्वादश खण्डः}
            \nopagebreak % Keep section title with first mantra
            
            % --- SECTION-LEVEL PROCEDURE LINK ---

                % --- SUBSECTIONS (Mantras) ---
                
                % Conditional rendering based on output_mode
                \par\filbreak

\phantomsection

\addcontentsline{toc}{subsection}{ऋक् १}

{\centering \footnotesize \textcolor{AccentPurple}{॥ मेधातिथिः काण्वः इन्द्रः गायत्री ॥} \par}

\vspace{0.12em}

{\centering \textbf{\textcolor{AccentBlue}{आतीहीमन्यूषावीणम्\kern 0.05em ।\hskip 0.12em plus 0.04em minus 0.02em सूषूवां\raisebox{0.25ex}{\accentmark{15}{\char"1CD2}}सामू\raisebox{0.25ex}{\accentmark{15}{\char"1CD2}}पै\raisebox{0.25ex}{\accentmark{12}{\char"0951}}रय\kern 0.05em ।\hskip 0.12em plus 0.04em minus 0.02em अस्या\raisebox{0.25ex}{\accentmark{12}{\char"0951}}रा\raisebox{0.15ex}{\accentmark{12}{\char"1CF9}}तौ\raisebox{0.25ex}{\accentmark{12}{\char"0951}}सूतंपी\raisebox{0.25ex}{\accentmark{12}{\char"0951}}बा \nolinebreak\hskip 0.08em \mbox{॥१॥}}} \par}

\vspace{0.8em}
                % --- SUBSECTIONS (Mantras) ---
                
                % Conditional rendering based on output_mode
                \par\filbreak

\phantomsection

\addcontentsline{toc}{subsection}{ऋक् २}

{\centering \footnotesize \textcolor{AccentPurple}{॥ वामदेवो गौतमः इन्द्रः गायत्री ॥} \par}

\vspace{0.12em}

{\centering \textbf{\textcolor{AccentBlue}{का\raisebox{0.25ex}{\accentmark{15}{\char"1CD2}}दूप्रा\raisebox{0.25ex}{\accentmark{15}{\char"1CD2}}चे\raisebox{0.25ex}{\accentmark{12}{\char"0951}}तसेमहे\raisebox{0.25ex}{\accentmark{15}{\char"1CD2}}\kern 0.05em ।\hskip 0.12em plus 0.04em minus 0.02em वा\raisebox{0.15ex}{\accentmark{12}{\char"1CF8}}चो\raisebox{0.25ex}{\accentmark{12}{\char"0951}}देवाया\raisebox{0.25ex}{\accentmark{12}{\char"0951}}शस्यते\kern 0.05em ।\hskip 0.12em plus 0.04em minus 0.02em तदी\raisebox{0.15ex}{\accentmark{12}{\char"1CF9}}द्ध्या\raisebox{0.25ex}{\accentmark{12}{\char"0951}}स्यावार्धा\raisebox{0.25ex}{\accentmark{15}{\char"1CD2}}नम् \nolinebreak\hskip 0.08em \mbox{॥२॥}}} \par}

\vspace{0.8em}
                % --- SUBSECTIONS (Mantras) ---
                
                % Conditional rendering based on output_mode
                \par\filbreak

\phantomsection

\addcontentsline{toc}{subsection}{ऋक् ३}

{\centering \textbf{\textcolor{AccentBlue}{॥ मेधातिथिःकाण्वः,प्रियमेधश्चांगिरसःइन्द्रःगायत्री ॥ ॥ विश्वामित्रोगाथिनःइन्द्रःगायत्री ॥ उक्थञ्चाना\raisebox{0.25ex}{\accentmark{15}{\char"1CD2}}शस्या\raisebox{0.25ex}{\accentmark{15}{\char"1CD2}}मा\raisebox{0.25ex}{\accentmark{12}{\char"0951}}नम्\kern 0.05em ।\hskip 0.12em plus 0.04em minus 0.02em ना\raisebox{0.15ex}{\accentmark{12}{\char"1CF8}}गो\raisebox{0.25ex}{\accentmark{12}{\char"0951}}रा\raisebox{0.25ex}{\accentmark{15}{\char"1CD2}}यिरा\raisebox{0.25ex}{\accentmark{15}{\char"1CD2}}ची\raisebox{0.25ex}{\accentmark{12}{\char"0951}}केता\kern 0.05em ।\hskip 0.12em plus 0.04em minus 0.02em ना\raisebox{0.15ex}{\accentmark{12}{\char"1CF8}}गा\raisebox{0.25ex}{\accentmark{15}{\char"1CD2}}यत्रं\raisebox{0.25ex}{\accentmark{12}{\char"0951}}गीयामानम् \nolinebreak\hskip 0.08em \mbox{॥३॥}}} \par}

\vspace{0.8em}
                % --- SUBSECTIONS (Mantras) ---
                
                % Conditional rendering based on output_mode
                \par\filbreak

\phantomsection

\addcontentsline{toc}{subsection}{ऋक् ४}

{\centering \textbf{\textcolor{AccentBlue}{इ\raisebox{0.15ex}{\accentmark{12}{\char"1CF8}}न्द्रा\raisebox{0.25ex}{\accentmark{12}{\char"0951}}उक्थे\raisebox{0.25ex}{\accentmark{15}{\char"1CD2}}भिर्म\raisebox{0.25ex}{\accentmark{12}{\char"0951}}न्दीष्ठः\kern 0.05em ।\hskip 0.12em plus 0.04em minus 0.02em वा\raisebox{0.25ex}{\accentmark{15}{\char"1CD2}}जा\raisebox{0.25ex}{\accentmark{12}{\char"0951}}नांचावा\raisebox{0.25ex}{\accentmark{12}{\char"0951}}ज\raisebox{0.25ex}{\accentmark{15}{\char"1CD2}}प\raisebox{0.25ex}{\accentmark{12}{\char"0951}}तिः\kern 0.05em ।\hskip 0.12em plus 0.04em minus 0.02em हा\raisebox{0.25ex}{\accentmark{15}{\char"1CD2}}रीवा\raisebox{0.25ex}{\accentmark{12}{\char"0951}}न्सूतानांसाखा\raisebox{0.25ex}{\accentmark{12}{\char"0951}} \nolinebreak\hskip 0.08em \mbox{॥४॥}}} \par}

\vspace{0.8em}
                % --- SUBSECTIONS (Mantras) ---
                
                % Conditional rendering based on output_mode
                \par\filbreak

\phantomsection

\addcontentsline{toc}{subsection}{ऋक् ५}

{\centering \textbf{\textcolor{AccentBlue}{॥ मेधातिथिःकाण्वः,प्रियमेधश्चांगिरसःइन्द्रःगायत्री ॥ ॥ दुमित्रः(सुमित्रोवा)कौत्सःइन्द्रःगायत्री ॥ ब्र\raisebox{0.25ex}{\accentmark{15}{\char"1CD2}}ह्मा\raisebox{0.25ex}{\accentmark{12}{\char"0951}}णादीन्द्रारा\raisebox{0.25ex}{\accentmark{12}{\char"0951}}धा\raisebox{0.25ex}{\accentmark{12}{\char"0951}}साः\kern 0.05em ।\hskip 0.12em plus 0.04em minus 0.02em पी\raisebox{0.25ex}{\accentmark{15}{\char"1CD2}}बा\raisebox{0.15ex}{\accentmark{12}{\char"1CF9}}सो\raisebox{0.15ex}{\accentmark{12}{\char"1CF9}}मा\raisebox{0.25ex}{\accentmark{12}{\char"0951}}मृतूंरा\raisebox{0.25ex}{\accentmark{15}{\char"1CD2}}नू\kern 0.05em ।\hskip 0.12em plus 0.04em minus 0.02em तवे\raisebox{0.15ex}{\accentmark{12}{\char"1CF9}}दु\raisebox{0.25ex}{\accentmark{12}{\char"0951}}सख्या\raisebox{0.25ex}{\accentmark{15}{\char"1CD2}}मस्तृतम् \nolinebreak\hskip 0.08em \mbox{॥५॥}}} \par}

\vspace{0.8em}
                % --- SUBSECTIONS (Mantras) ---
                
                % Conditional rendering based on output_mode
                \par\filbreak

\phantomsection

\addcontentsline{toc}{subsection}{ऋक् ६}

{\centering \textbf{\textcolor{AccentBlue}{वा\raisebox{0.15ex}{\accentmark{12}{\char"1CF8}}यङ्घा\raisebox{0.25ex}{\accentmark{12}{\char"0951}}तेआपी\raisebox{0.25ex}{\accentmark{12}{\char"0951}}स्मसि\kern 0.05em ।\hskip 0.12em plus 0.04em minus 0.02em स्तोता\raisebox{0.25ex}{\accentmark{15}{\char"1CD2}}र\raisebox{0.25ex}{\accentmark{12}{\char"0951}}इन्द्र\raisebox{0.15ex}{\accentmark{12}{\char"1CF8}}गिर्वणः\kern 0.05em ।\hskip 0.12em plus 0.04em minus 0.02em त्व\raisebox{0.25ex}{\accentmark{15}{\char"1CD2}}न्नोजिन्वसोमपः \nolinebreak\hskip 0.08em \mbox{॥६॥}}} \par}

\vspace{0.8em}
                % --- SUBSECTIONS (Mantras) ---
                
                % Conditional rendering based on output_mode
                \par\filbreak

\phantomsection

\addcontentsline{toc}{subsection}{ऋक् ७}

{\centering \textbf{\textcolor{AccentBlue}{॥ मेधातिथिःकाण्वःइन्द्रःगायत्री ॥ ॥ मेधातिथिःकाण्वःइन्द्रःगायत्री ॥ आ\raisebox{0.25ex}{\accentmark{12}{\char"0951}}याह्यू\raisebox{0.25ex}{\accentmark{15}{\char"1CD2}}पा\raisebox{0.25ex}{\accentmark{12}{\char"0951}}नास्सू\raisebox{0.25ex}{\accentmark{15}{\char"1CD2}}तम्\kern 0.05em ।\hskip 0.12em plus 0.04em minus 0.02em वा\raisebox{0.15ex}{\accentmark{12}{\char"1CF8}}जेभिर्मा\raisebox{0.25ex}{\accentmark{12}{\char"0951}}हृ\raisebox{0.25ex}{\accentmark{12}{\char"0951}}णीयथा\kern 0.05em ।\hskip 0.12em plus 0.04em minus 0.02em माहं\raisebox{0.15ex}{\accentmark{12}{\char"1CF9}}ईवा\raisebox{0.25ex}{\accentmark{12}{\char"0951}}यूवा\raisebox{0.25ex}{\accentmark{15}{\char"1CD2}}जानि \nolinebreak\hskip 0.08em \mbox{॥७॥}}} \par}

\vspace{0.8em}
                % --- SUBSECTIONS (Mantras) ---
                
                % Conditional rendering based on output_mode
                \par\filbreak

\phantomsection

\addcontentsline{toc}{subsection}{ऋक् ८}

{\centering \textbf{\textcolor{AccentBlue}{का\raisebox{0.15ex}{\accentmark{12}{\char"1CF8}}दा\raisebox{0.25ex}{\accentmark{15}{\char"1CD2}}वासोस्तोत्रं\raisebox{0.15ex}{\accentmark{12}{\char"1CF9}}ह\raisebox{0.15ex}{\accentmark{12}{\char"1CF9}}र्याताआ\kern 0.05em ।\hskip 0.12em plus 0.04em minus 0.02em आव\raisebox{0.25ex}{\accentmark{12}{\char"0951}}श्मशा\raisebox{0.15ex}{\accentmark{12}{\char"1CF9}}रूधाद्वा\raisebox{0.25ex}{\accentmark{15}{\char"1CD2}}ः\kern 0.05em ।\hskip 0.12em plus 0.04em minus 0.02em दीर्घं\raisebox{0.25ex}{\accentmark{15}{\char"1CD2}}सूतं\raisebox{0.25ex}{\accentmark{15}{\char"1CD2}}वा\raisebox{0.25ex}{\accentmark{15}{\char"1CD2}}ता\raisebox{0.25ex}{\accentmark{15}{\char"1CD2}}प्या\raisebox{0.25ex}{\accentmark{12}{\char"0951}}या \nolinebreak\hskip 0.08em \mbox{॥८॥}}} \par}

\vspace{0.8em}
                % --- SUBSECTIONS (Mantras) ---
                
                % Conditional rendering based on output_mode
                \par\filbreak

\phantomsection

\addcontentsline{toc}{subsection}{ऋक् ९}

{\centering \footnotesize \textcolor{AccentPurple}{॥ विश्वामित्रो गाथिनोऽभीपाद उदलो वा इन्द्रः गायत्री ॥} \par}

\vspace{0.12em}

{\centering \textbf{\textcolor{AccentBlue}{ए\raisebox{0.15ex}{\accentmark{12}{\char"1CF8}}न्द्रा\raisebox{0.25ex}{\accentmark{12}{\char"0951}}पृक्षूका\raisebox{0.25ex}{\accentmark{12}{\char"0951}}सूचीत्\kern 0.05em ।\hskip 0.12em plus 0.04em minus 0.02em नृम्णं\raisebox{0.25ex}{\accentmark{15}{\char"1CD2}}तानू\raisebox{0.25ex}{\accentmark{15}{\char"1CD2}}षू\raisebox{0.25ex}{\accentmark{12}{\char"0951}}धेहिनः\kern 0.05em ।\hskip 0.12em plus 0.04em minus 0.02em स\raisebox{0.25ex}{\accentmark{15}{\char"1CD2}}त्रा\raisebox{0.25ex}{\accentmark{15}{\char"1CD2}}जिदुग्रापौंठ\raisebox{0.25ex}{\accentmark{12}{\char"0951}}स्यम् \nolinebreak\hskip 0.08em \mbox{॥९॥}}} \par}

\vspace{0.8em}
                % --- SUBSECTIONS (Mantras) ---
                
                % Conditional rendering based on output_mode
                \par\filbreak

\phantomsection

\addcontentsline{toc}{subsection}{ऋक् १०}

{\centering \footnotesize \textcolor{AccentPurple}{॥ श्रुतकक्षः (सुकेवा) आंगिरसः इन्द्रः गायत्री ॥} \par}

\vspace{0.12em}

{\centering \textbf{\textcolor{AccentBlue}{आ\raisebox{0.25ex}{\accentmark{15}{\char"1CD2}}यामेनंसचतां\raisebox{0.25ex}{\accentmark{15}{\char"1CD2}}सुत\raisebox{0.25ex}{\accentmark{15}{\char"1CD2}}ः\kern 0.05em ।\hskip 0.12em plus 0.04em minus 0.02em मन्दी\raisebox{0.25ex}{\accentmark{15}{\char"1CD2}}मिन्द्रा\raisebox{0.25ex}{\accentmark{12}{\char"0951}}यामन्दी\raisebox{0.25ex}{\accentmark{15}{\char"1CD2}}ने\kern 0.05em ।\hskip 0.12em plus 0.04em minus 0.02em च\raisebox{0.25ex}{\accentmark{15}{\char"1CD2}}क्रिंवि\raisebox{0.25ex}{\accentmark{12}{\char"0951}}श्वा\raisebox{0.25ex}{\accentmark{12}{\char"0951}}नी\raisebox{0.25ex}{\accentmark{12}{\char"0951}}चाक्रा\raisebox{0.25ex}{\accentmark{12}{\char"0951}}ये \nolinebreak\hskip 0.08em \mbox{॥१०॥}}} \par}

\vspace{0.8em}

            % --- End of Kandah Closing Floral Line & Counts ---
            \par\nopagebreak\vspace{0.9em plus 0.3em minus 0.1em}
            \centerline{\textcolor{black!35}{\rule{1.8cm}{0.35pt}}\;\raisebox{-0.22ex}{\textcolor{vedicmaroon}{\large\SixFlowerPetalDotted}}\;\textcolor{black!35}{\rule{1.8cm}{0.35pt}}}
            \nopagebreak\vspace{0.45em}
            \centerline{\devafont\textbf{\textcolor{AccentBlue}{॥ ऋक् सङ्ख्या : १० ॥}}}


    % --- SUPERSECTION ---
    \clearpage           
    \phantomsection              
    \chapter*{\centering बृहतिपाठः}
    \addcontentsline{toc}{chapter}{बृहतिपाठः}
    \markboth{बृहतिपाठः}{बृहतिपाठः}
    
    % GLUE FIX: Prevent page break between Chapter and First Section
    \nopagebreak

        

            % --- SECTION ---
            
            \phantomsection  

            \section*{\centering प्रथम खण्डः}
            \addcontentsline{toc}{section}{प्रथम खण्डः} 
            \markright{प्रथम खण्डः}
            \nopagebreak % Keep section title with first mantra
            
            % --- SECTION-LEVEL PROCEDURE LINK ---

                % --- SUBSECTIONS (Mantras) ---
                
                % Conditional rendering based on output_mode
                \par\filbreak

\phantomsection

\addcontentsline{toc}{subsection}{ऋक् १}

{\centering \footnotesize \textcolor{AccentPurple}{॥ वसिष्ठो मैत्रावरुणिः इन्द्रः बृहती ॥} \par}

\vspace{0.12em}

{\centering \textbf{\textcolor{AccentBlue}{आभित्वा\raisebox{0.25ex}{\accentmark{15}{\char"1CD2}}शू\raisebox{0.25ex}{\accentmark{15}{\char"1CD2}}र\raisebox{0.25ex}{\accentmark{12}{\char"0951}}नोनुमः\kern 0.05em ।\hskip 0.12em plus 0.04em minus 0.02em आदू\raisebox{0.15ex}{\accentmark{12}{\char"1CF8}}ग्धाइवाधेना\raisebox{0.25ex}{\accentmark{15}{\char"1CD2}}वाः\kern 0.05em ।\hskip 0.12em plus 0.04em minus 0.02em ई\raisebox{0.25ex}{\accentmark{15}{\char"1CD2}}शा\raisebox{0.25ex}{\accentmark{12}{\char"0951}}ना\raisebox{0.25ex}{\accentmark{15}{\char"1CD2}}मा\raisebox{0.25ex}{\accentmark{12}{\char"0951}}स्याजग\raisebox{0.25ex}{\accentmark{12}{\char"0951}}तःस्वदृशम्\kern 0.05em ।\hskip 0.12em plus 0.04em minus 0.02em ईशा\raisebox{0.25ex}{\accentmark{12}{\char"0951}}नमिन्द्रातस्थू\raisebox{0.25ex}{\accentmark{15}{\char"1CD2}}षा\raisebox{0.25ex}{\accentmark{12}{\char"0951}}ः \nolinebreak\hskip 0.08em \mbox{॥१॥}}} \par}

\vspace{0.8em}
                % --- SUBSECTIONS (Mantras) ---
                
                % Conditional rendering based on output_mode
                \par\filbreak

\phantomsection

\addcontentsline{toc}{subsection}{ऋक् २}

{\centering \footnotesize \textcolor{AccentPurple}{॥ भरद्वाजः (शयुः) बार्हस्पत्यः इन्द्रः बृहती ॥} \par}

\vspace{0.12em}

{\centering \textbf{\textcolor{AccentBlue}{त्वा\raisebox{0.25ex}{\accentmark{15}{\char"1CD2}}मि\raisebox{0.25ex}{\accentmark{15}{\char"1CD2}}द्धि\raisebox{0.25ex}{\accentmark{15}{\char"1CD2}}हा\raisebox{0.25ex}{\accentmark{15}{\char"1CD2}}वा\raisebox{0.25ex}{\accentmark{12}{\char"0951}}महे\raisebox{0.25ex}{\accentmark{12}{\char"0951}}\kern 0.05em ।\hskip 0.12em plus 0.04em minus 0.02em सा\raisebox{0.25ex}{\accentmark{15}{\char"1CD2}}तौवा\raisebox{0.25ex}{\accentmark{15}{\char"1CD2}}ज\raisebox{0.25ex}{\accentmark{12}{\char"0951}}स्याकारा\raisebox{0.25ex}{\accentmark{15}{\char"1CD2}}वा\raisebox{0.25ex}{\accentmark{12}{\char"0951}}ः\kern 0.05em ।\hskip 0.12em plus 0.04em minus 0.02em त्वांवृ\raisebox{0.25ex}{\accentmark{12}{\char"0951}}त्रे\raisebox{0.25ex}{\accentmark{15}{\char"1CD2}}ष्विन्द्रा\raisebox{0.25ex}{\accentmark{12}{\char"0951}}सात्पातिन्ना\raisebox{0.25ex}{\accentmark{15}{\char"1CD2}}रा\raisebox{0.25ex}{\accentmark{12}{\char"0951}}ः\kern 0.05em ।\hskip 0.12em plus 0.04em minus 0.02em त्वांकाष्ठास्व\raisebox{0.25ex}{\accentmark{12}{\char"0951}}र्वातः \nolinebreak\hskip 0.08em \mbox{॥२॥}}} \par}

\vspace{0.8em}
                % --- SUBSECTIONS (Mantras) ---
                
                % Conditional rendering based on output_mode
                \par\filbreak

\phantomsection

\addcontentsline{toc}{subsection}{ऋक् ३}

{\centering \footnotesize \textcolor{AccentPurple}{॥ प्रस्कण्वः इन्द्रः बृहती ॥} \par}

\vspace{0.12em}

{\centering \textbf{\textcolor{AccentBlue}{आभी\raisebox{0.15ex}{\accentmark{12}{\char"1CF9}}प्रा\raisebox{0.15ex}{\accentmark{12}{\char"1CF9}}वास्सूरा\raisebox{0.25ex}{\accentmark{15}{\char"1CD2}}धा\raisebox{0.25ex}{\accentmark{12}{\char"0951}}सम्\kern 0.05em ।\hskip 0.12em plus 0.04em minus 0.02em इन्द्रा\raisebox{0.25ex}{\accentmark{12}{\char"0951}}मर्चाया\raisebox{0.15ex}{\accentmark{12}{\char"1CF9}}था\raisebox{0.25ex}{\accentmark{12}{\char"0951}}वीदे\kern 0.05em ।\hskip 0.12em plus 0.04em minus 0.02em योज\raisebox{0.25ex}{\accentmark{12}{\char"0951}}रितृ\raisebox{0.15ex}{\accentmark{12}{\char"1CF9}}भ्यो\raisebox{0.25ex}{\accentmark{12}{\char"0951}}माघा\raisebox{0.25ex}{\accentmark{15}{\char"1CD2}}वा\raisebox{0.25ex}{\accentmark{12}{\char"0951}}पूरो\raisebox{0.15ex}{\accentmark{12}{\char"1CF8}}वासू\raisebox{0.15ex}{\accentmark{12}{\char"1CF9}}ः\kern 0.05em ।\hskip 0.12em plus 0.04em minus 0.02em सा\raisebox{0.25ex}{\accentmark{15}{\char"1CD2}}हा\raisebox{0.25ex}{\accentmark{12}{\char"0951}}स्रेणेवाशी\raisebox{0.25ex}{\accentmark{15}{\char"1CD2}}क्षा\raisebox{0.25ex}{\accentmark{12}{\char"0951}}ति \nolinebreak\hskip 0.08em \mbox{॥३॥}}} \par}

\vspace{0.8em}
                % --- SUBSECTIONS (Mantras) ---
                
                % Conditional rendering based on output_mode
                \par\filbreak

\phantomsection

\addcontentsline{toc}{subsection}{ऋक् ४}

{\centering \footnotesize \textcolor{AccentPurple}{॥ नोधा गौतमः इन्द्रः बृहती ॥} \par}

\vspace{0.12em}

{\centering \textbf{\textcolor{AccentBlue}{तंवो\raisebox{0.25ex}{\accentmark{12}{\char"0951}}दस्मा\raisebox{0.15ex}{\accentmark{12}{\char"1CF9}}मृ\raisebox{0.25ex}{\accentmark{12}{\char"0951}}तीषाहम्\kern 0.05em ।\hskip 0.12em plus 0.04em minus 0.02em वा\raisebox{0.25ex}{\accentmark{15}{\char"1CD2}}सो\raisebox{0.25ex}{\accentmark{12}{\char"0951}}र्मन्दानमन्धा\raisebox{0.25ex}{\accentmark{12}{\char"0951}}सा\kern 0.05em ।\hskip 0.12em plus 0.04em minus 0.02em आ\raisebox{0.25ex}{\accentmark{15}{\char"1CD2}}भी\raisebox{0.25ex}{\accentmark{15}{\char"1CD2}}वत्स\raisebox{0.25ex}{\accentmark{15}{\char"1CD2}}न्नस्वा\raisebox{0.25ex}{\accentmark{15}{\char"1CD2}}सारेषुधेना\raisebox{0.25ex}{\accentmark{15}{\char"1CD2}}वा\raisebox{0.25ex}{\accentmark{12}{\char"0951}}ः\kern 0.05em ।\hskip 0.12em plus 0.04em minus 0.02em इ\raisebox{0.15ex}{\accentmark{12}{\char"1CF8}}न्द्रं\raisebox{0.25ex}{\accentmark{12}{\char"0951}}गीर्भि\raisebox{0.25ex}{\accentmark{12}{\char"0951}}र्ना\raisebox{0.25ex}{\accentmark{12}{\char"0951}}वामहे \nolinebreak\hskip 0.08em \mbox{॥४॥}}} \par}

\vspace{0.8em}
                % --- SUBSECTIONS (Mantras) ---
                
                % Conditional rendering based on output_mode
                \par\filbreak

\phantomsection

\addcontentsline{toc}{subsection}{ऋक् ५}

{\centering \footnotesize \textcolor{AccentPurple}{॥ कलिः प्रागाथः इन्द्रः बृहती ॥} \par}

\vspace{0.12em}

{\centering \textbf{\textcolor{AccentBlue}{ता\raisebox{0.25ex}{\accentmark{15}{\char"1CD2}}रो\raisebox{0.25ex}{\accentmark{12}{\char"0951}}भिर्वोवीद\raisebox{0.25ex}{\accentmark{15}{\char"1CD2}}द्वा\raisebox{0.25ex}{\accentmark{12}{\char"0951}}सुम्\kern 0.05em ।\hskip 0.12em plus 0.04em minus 0.02em इन्द्रंसाबा\raisebox{0.15ex}{\accentmark{12}{\char"1CF9}}धा\raisebox{0.25ex}{\accentmark{12}{\char"0951}}ऊता\raisebox{0.25ex}{\accentmark{15}{\char"1CD2}}ये\kern 0.05em ।\hskip 0.12em plus 0.04em minus 0.02em बृह\raisebox{0.25ex}{\accentmark{15}{\char"1CD2}}त्गा\raisebox{0.25ex}{\accentmark{15}{\char"1CD2}}यन्तास्सूता\raisebox{0.25ex}{\accentmark{15}{\char"1CD2}}सो\raisebox{0.25ex}{\accentmark{15}{\char"1CD2}}मेआध्वारे\kern 0.05em ।\hskip 0.12em plus 0.04em minus 0.02em हूवे\raisebox{0.25ex}{\accentmark{15}{\char"1CD2}}भारन्नाकारीण\raisebox{0.25ex}{\accentmark{12}{\char"0951}}म् \nolinebreak\hskip 0.08em \mbox{॥५॥}}} \par}

\vspace{0.8em}
                % --- SUBSECTIONS (Mantras) ---
                
                % Conditional rendering based on output_mode
                \par\filbreak

\phantomsection

\addcontentsline{toc}{subsection}{ऋक् ६}

{\centering \footnotesize \textcolor{AccentPurple}{॥ वसिष्ठो मैत्रावरुणिः इन्द्रः बृहती ॥} \par}

\vspace{0.12em}

{\centering \textbf{\textcolor{AccentBlue}{ता\raisebox{0.25ex}{\accentmark{15}{\char"1CD2}}रा\raisebox{0.25ex}{\accentmark{15}{\char"1CD2}}णीरित्सी\raisebox{0.25ex}{\accentmark{15}{\char"1CD2}}षासती\kern 0.05em ।\hskip 0.12em plus 0.04em minus 0.02em वाजं\raisebox{0.25ex}{\accentmark{15}{\char"1CD2}}पू\raisebox{0.25ex}{\accentmark{15}{\char"1CD2}}रन्ध्यायूजा\kern 0.05em ।\hskip 0.12em plus 0.04em minus 0.02em आ\raisebox{0.25ex}{\accentmark{15}{\char"1CD2}}वाइन्द्रं\raisebox{0.25ex}{\accentmark{12}{\char"0951}}पूरुहूत\raisebox{0.25ex}{\accentmark{15}{\char"1CD2}}न्ना\raisebox{0.25ex}{\accentmark{12}{\char"0951}}मेगीरा\raisebox{0.25ex}{\accentmark{15}{\char"1CD2}}\kern 0.05em ।\hskip 0.12em plus 0.04em minus 0.02em नेमि\raisebox{0.25ex}{\accentmark{15}{\char"1CD2}}न्तष्टे\raisebox{0.25ex}{\accentmark{12}{\char"0951}}वासुद्रूव\raisebox{0.25ex}{\accentmark{12}{\char"0951}}म् \nolinebreak\hskip 0.08em \mbox{॥६॥}}} \par}

\vspace{0.8em}
                % --- SUBSECTIONS (Mantras) ---
                
                % Conditional rendering based on output_mode
                \par\filbreak

\phantomsection

\addcontentsline{toc}{subsection}{ऋक् ७}

{\centering \footnotesize \textcolor{AccentPurple}{॥ मेधातिथिः काण्वः इन्द्रः बृहती ॥} \par}

\vspace{0.12em}

{\centering \textbf{\textcolor{AccentBlue}{पीबा\raisebox{0.15ex}{\accentmark{12}{\char"1CF8}}सू\raisebox{0.15ex}{\accentmark{12}{\char"1CF9}}तस्या\raisebox{0.25ex}{\accentmark{12}{\char"0951}}रासी\raisebox{0.25ex}{\accentmark{15}{\char"1CD2}}नाः\kern 0.05em ।\hskip 0.12em plus 0.04em minus 0.02em मत्स्वा\raisebox{0.25ex}{\accentmark{12}{\char"0951}}नइन्द्रागो\raisebox{0.25ex}{\accentmark{15}{\char"1CD2}}मा\raisebox{0.25ex}{\accentmark{12}{\char"0951}}तः\kern 0.05em ।\hskip 0.12em plus 0.04em minus 0.02em आपिर्नोबोधीसा\raisebox{0.25ex}{\accentmark{15}{\char"1CD2}}धा\raisebox{0.25ex}{\accentmark{12}{\char"0951}}मा\raisebox{0.15ex}{\accentmark{12}{\char"1CF9}}द्ये\raisebox{0.25ex}{\accentmark{12}{\char"0951}}वृधे\raisebox{0.25ex}{\accentmark{15}{\char"1CD2}}\kern 0.05em ।\hskip 0.12em plus 0.04em minus 0.02em अस्मंआ\raisebox{0.25ex}{\accentmark{12}{\char"0951}}वन्तूते\raisebox{0.25ex}{\accentmark{15}{\char"1CD2}}धी\raisebox{0.25ex}{\accentmark{15}{\char"1CD2}}या\raisebox{0.25ex}{\accentmark{12}{\char"0951}}ः \nolinebreak\hskip 0.08em \mbox{॥७॥}}} \par}

\vspace{0.8em}
                % --- SUBSECTIONS (Mantras) ---
                
                % Conditional rendering based on output_mode
                \par\filbreak

\phantomsection

\addcontentsline{toc}{subsection}{ऋक् ८}

{\centering \footnotesize \textcolor{AccentPurple}{॥ भगः प्रागाथः इन्द्रः बृहती ॥} \par}

\vspace{0.12em}

{\centering \textbf{\textcolor{AccentBlue}{त्वं\raisebox{0.25ex}{\accentmark{15}{\char"1CD2}}ह्ये\raisebox{0.25ex}{\accentmark{15}{\char"1CD2}}हीचेरा\raisebox{0.25ex}{\accentmark{15}{\char"1CD2}}वे\kern 0.05em ।\hskip 0.12em plus 0.04em minus 0.02em वीदा\raisebox{0.15ex}{\accentmark{12}{\char"1CF9}}भा\raisebox{0.25ex}{\accentmark{12}{\char"0951}}गंवा\raisebox{0.25ex}{\accentmark{15}{\char"1CD2}}सू\raisebox{0.25ex}{\accentmark{12}{\char"0951}}क्तये\kern 0.05em ।\hskip 0.12em plus 0.04em minus 0.02em उद्वा\raisebox{0.25ex}{\accentmark{12}{\char"0951}}वृष\raisebox{0.25ex}{\accentmark{12}{\char"0951}}स्वमा\raisebox{0.25ex}{\accentmark{12}{\char"0951}}घावन्ग\raisebox{0.25ex}{\accentmark{12}{\char"0951}}वी\raisebox{0.25ex}{\accentmark{15}{\char"1CD2}}ष्टये\kern 0.05em ।\hskip 0.12em plus 0.04em minus 0.02em उदि\raisebox{0.25ex}{\accentmark{15}{\char"1CD2}}न्द्रो\raisebox{0.25ex}{\accentmark{12}{\char"0951}}अश्वा\raisebox{0.25ex}{\accentmark{12}{\char"0951}}मिष्ठ\raisebox{0.25ex}{\accentmark{12}{\char"0951}}ये \nolinebreak\hskip 0.08em \mbox{॥८॥}}} \par}

\vspace{0.8em}
                % --- SUBSECTIONS (Mantras) ---
                
                % Conditional rendering based on output_mode
                \par\filbreak

\phantomsection

\addcontentsline{toc}{subsection}{ऋक् ९}

{\centering \footnotesize \textcolor{AccentPurple}{॥ वसिष्ठो मैत्रावरुणिः मरुतः बृहती ॥} \par}

\vspace{0.12em}

{\centering \textbf{\textcolor{AccentBlue}{ना\raisebox{0.25ex}{\accentmark{15}{\char"1CD2}}ही\raisebox{0.25ex}{\accentmark{15}{\char"1CD2}}वा\raisebox{0.25ex}{\accentmark{12}{\char"0951}}श्वाराम\raisebox{0.25ex}{\accentmark{15}{\char"1CD2}}ञ्चाना\raisebox{0.25ex}{\accentmark{15}{\char"1CD2}}\kern 0.05em ।\hskip 0.12em plus 0.04em minus 0.02em वासी\raisebox{0.25ex}{\accentmark{12}{\char"0951}}ष्ठःपारीमं\raisebox{0.25ex}{\accentmark{15}{\char"1CD2}}सा\raisebox{0.25ex}{\accentmark{12}{\char"0951}}ते\kern 0.05em ।\hskip 0.12em plus 0.04em minus 0.02em अस्माका\raisebox{0.25ex}{\accentmark{12}{\char"0951}}माद्यामारूतस्सूते\raisebox{0.25ex}{\accentmark{15}{\char"1CD2}}सा\raisebox{0.25ex}{\accentmark{15}{\char"1CD2}}चा\raisebox{0.25ex}{\accentmark{12}{\char"0951}}\kern 0.05em ।\hskip 0.12em plus 0.04em minus 0.02em विश्वे\raisebox{0.25ex}{\accentmark{12}{\char"0951}}पिबन्तूका\raisebox{0.25ex}{\accentmark{15}{\char"1CD2}}मी\raisebox{0.25ex}{\accentmark{15}{\char"1CD2}}ना\raisebox{0.25ex}{\accentmark{12}{\char"0951}}ः \nolinebreak\hskip 0.08em \mbox{॥९॥}}} \par}

\vspace{0.8em}
                % --- SUBSECTIONS (Mantras) ---
                
                % Conditional rendering based on output_mode
                \par\filbreak

\phantomsection

\addcontentsline{toc}{subsection}{ऋक् १०}

{\centering \footnotesize \textcolor{AccentPurple}{॥ प्रगाथो घौरः काण्वः इन्द्रः बृहती ॥} \par}

\vspace{0.12em}

{\centering \textbf{\textcolor{AccentBlue}{माची\raisebox{0.25ex}{\accentmark{12}{\char"0951}}दन्यद्वी\raisebox{0.25ex}{\accentmark{15}{\char"1CD2}}शंसता\kern 0.05em ।\hskip 0.12em plus 0.04em minus 0.02em सा\raisebox{0.15ex}{\accentmark{12}{\char"1CF8}}खायोमा\raisebox{0.25ex}{\accentmark{15}{\char"1CD2}}रि\raisebox{0.25ex}{\accentmark{12}{\char"0951}}षण्यत\kern 0.05em ।\hskip 0.12em plus 0.04em minus 0.02em इ\raisebox{0.15ex}{\accentmark{12}{\char"1CF8}}न्द्रामित्स्तो\raisebox{0.15ex}{\accentmark{12}{\char"1CF8}}ता\raisebox{0.15ex}{\accentmark{12}{\char"1CF8}}वृषाणंसा\raisebox{0.15ex}{\accentmark{12}{\char"1CF9}}चा\raisebox{0.25ex}{\accentmark{12}{\char"0951}}सूते\raisebox{0.25ex}{\accentmark{15}{\char"1CD2}}\kern 0.05em ।\hskip 0.12em plus 0.04em minus 0.02em मुहू\raisebox{0.25ex}{\accentmark{12}{\char"0951}}रुत्थाचाशंसत \nolinebreak\hskip 0.08em \mbox{॥१०॥}}} \par}

\vspace{0.8em}

            % --- End of Kandah Closing Floral Line & Counts ---
            \par\nopagebreak\vspace{0.9em plus 0.3em minus 0.1em}
            \centerline{\textcolor{black!35}{\rule{1.8cm}{0.35pt}}\;\raisebox{-0.22ex}{\textcolor{vedicmaroon}{\large\SixFlowerPetalDotted}}\;\textcolor{black!35}{\rule{1.8cm}{0.35pt}}}
            \nopagebreak\vspace{0.45em}
            \centerline{\devafont\textbf{\textcolor{AccentBlue}{॥ ऋक् सङ्ख्या : १० ॥}}}
        
            \clearpage

            % --- SECTION ---
            
            \phantomsection  

            \section*{\centering द्वितीय खण्डः}
            \addcontentsline{toc}{section}{द्वितीय खण्डः} 
            \markright{द्वितीय खण्डः}
            \nopagebreak % Keep section title with first mantra
            
            % --- SECTION-LEVEL PROCEDURE LINK ---

                % --- SUBSECTIONS (Mantras) ---
                
                % Conditional rendering based on output_mode
                \par\filbreak

\phantomsection

\addcontentsline{toc}{subsection}{ऋक् १}

{\centering \footnotesize \textcolor{AccentPurple}{॥ पुरुहन्मा आंगिरसः इन्द्रः बृहती ॥} \par}

\vspace{0.12em}

{\centering \textbf{\textcolor{AccentBlue}{ना\raisebox{0.15ex}{\accentmark{12}{\char"1CF8}}की\raisebox{0.15ex}{\accentmark{12}{\char"1CF8}}ष्टंकर्माणानशात्\kern 0.05em ।\hskip 0.12em plus 0.04em minus 0.02em यश्चाका\raisebox{0.15ex}{\accentmark{12}{\char"1CF8}}रा\raisebox{0.25ex}{\accentmark{12}{\char"0951}}स्सा\raisebox{0.25ex}{\accentmark{15}{\char"1CD2}}दावृधम्\kern 0.05em ।\hskip 0.12em plus 0.04em minus 0.02em इ\raisebox{0.25ex}{\accentmark{15}{\char"1CD2}}न्द्रन्नाय\raisebox{0.15ex}{\accentmark{12}{\char"1CF9}}ज्ञैर्वि\raisebox{0.15ex}{\accentmark{12}{\char"1CF9}}श्वा\raisebox{0.15ex}{\accentmark{12}{\char"1CF9}}र्गू\raisebox{0.25ex}{\accentmark{12}{\char"0951}}र्तामृ\raisebox{0.25ex}{\accentmark{12}{\char"0951}}भ्वा\raisebox{0.25ex}{\accentmark{12}{\char"0951}}सम्\kern 0.05em ।\hskip 0.12em plus 0.04em minus 0.02em आधृष्ट\raisebox{0.25ex}{\accentmark{15}{\char"1CD2}}न्धृष्णु\raisebox{0.25ex}{\accentmark{15}{\char"1CD2}}मोज\raisebox{0.25ex}{\accentmark{12}{\char"0951}}सा \nolinebreak\hskip 0.08em \mbox{॥१॥}}} \par}

\vspace{0.8em}
                % --- SUBSECTIONS (Mantras) ---
                
                % Conditional rendering based on output_mode
                \par\filbreak

\phantomsection

\addcontentsline{toc}{subsection}{ऋक् २}

{\centering \footnotesize \textcolor{AccentPurple}{॥ मेधातिथि – मध्येआतिथी काण्वौ इन्द्रः बृहती ॥} \par}

\vspace{0.12em}

{\centering \textbf{\textcolor{AccentBlue}{या\raisebox{0.25ex}{\accentmark{12}{\char"0951}}ऋते\raisebox{0.25ex}{\accentmark{15}{\char"1CD2}}चीदभिश्री\raisebox{0.25ex}{\accentmark{15}{\char"1CD2}}षा\raisebox{0.25ex}{\accentmark{12}{\char"0951}}ः\kern 0.05em ।\hskip 0.12em plus 0.04em minus 0.02em पूरा\raisebox{0.25ex}{\accentmark{15}{\char"1CD2}}जत्रु\raisebox{0.15ex}{\accentmark{12}{\char"1CF9}}भ्या\raisebox{0.25ex}{\accentmark{12}{\char"0951}}आतृ\raisebox{0.25ex}{\accentmark{15}{\char"1CD2}}दा\raisebox{0.25ex}{\accentmark{12}{\char"0951}}ः\kern 0.05em ।\hskip 0.12em plus 0.04em minus 0.02em स\raisebox{0.25ex}{\accentmark{15}{\char"1CD2}}न्धा\raisebox{0.25ex}{\accentmark{12}{\char"0951}}तासंधींमा\raisebox{0.25ex}{\accentmark{12}{\char"0951}}घावा\raisebox{0.25ex}{\accentmark{12}{\char"0951}}पूरोवा\raisebox{0.25ex}{\accentmark{15}{\char"1CD2}}सूः\kern 0.05em ।\hskip 0.12em plus 0.04em minus 0.02em निष्क\raisebox{0.25ex}{\accentmark{12}{\char"0951}}र्ता\raisebox{0.25ex}{\accentmark{12}{\char"0951}}वी\raisebox{0.15ex}{\accentmark{12}{\char"1CF9}}ह्रूतंपू\raisebox{0.25ex}{\accentmark{15}{\char"1CD2}}नाः \nolinebreak\hskip 0.08em \mbox{॥२॥}}} \par}

\vspace{0.8em}
                % --- SUBSECTIONS (Mantras) ---
                
                % Conditional rendering based on output_mode
                \par\filbreak

\phantomsection

\addcontentsline{toc}{subsection}{ऋक् ३}

{\centering \footnotesize \textcolor{AccentPurple}{॥ मेधातिथि – मध्येआतिथी काण्वौ इन्द्रः बृहती ॥} \par}

\vspace{0.12em}

{\centering \textbf{\textcolor{AccentBlue}{आ\raisebox{0.15ex}{\accentmark{12}{\char"1CF8}}त्वासाहस्रा\raisebox{0.25ex}{\accentmark{12}{\char"0951}}माशात\raisebox{0.25ex}{\accentmark{15}{\char"1CD2}}म्\kern 0.05em ।\hskip 0.12em plus 0.04em minus 0.02em युक्ता\raisebox{0.25ex}{\accentmark{15}{\char"1CD2}}राथेहीरण्याये\raisebox{0.25ex}{\accentmark{12}{\char"0951}}\kern 0.05em ।\hskip 0.12em plus 0.04em minus 0.02em ब्रह्मायू\raisebox{0.25ex}{\accentmark{15}{\char"1CD2}}जोहा\raisebox{0.25ex}{\accentmark{15}{\char"1CD2}}रा\raisebox{0.25ex}{\accentmark{12}{\char"0951}}यिन्द्राकेशी\raisebox{0.25ex}{\accentmark{15}{\char"1CD2}}नाः\kern 0.05em ।\hskip 0.12em plus 0.04em minus 0.02em वाहन्तूसो\raisebox{0.25ex}{\accentmark{15}{\char"1CD2}}मापीतये \nolinebreak\hskip 0.08em \mbox{॥३॥}}} \par}

\vspace{0.8em}
                % --- SUBSECTIONS (Mantras) ---
                
                % Conditional rendering based on output_mode
                \par\filbreak

\phantomsection

\addcontentsline{toc}{subsection}{ऋक् ४}

{\centering \footnotesize \textcolor{AccentPurple}{॥ विश्वामित्रो गाथिनः इन्द्रः बृहती ॥} \par}

\vspace{0.12em}

{\centering \textbf{\textcolor{AccentBlue}{आमन्द्रै\raisebox{0.15ex}{\accentmark{12}{\char"1CF9}}रि\raisebox{0.25ex}{\accentmark{12}{\char"0951}}न्द्राहा\raisebox{0.25ex}{\accentmark{15}{\char"1CD2}}री\raisebox{0.25ex}{\accentmark{12}{\char"0951}}भिः\kern 0.05em ।\hskip 0.12em plus 0.04em minus 0.02em याही\raisebox{0.25ex}{\accentmark{15}{\char"1CD2}}मा\raisebox{0.25ex}{\accentmark{12}{\char"0951}}यूरा\raisebox{0.25ex}{\accentmark{12}{\char"0951}}रोम\raisebox{0.25ex}{\accentmark{15}{\char"1CD2}}भिः\kern 0.05em ।\hskip 0.12em plus 0.04em minus 0.02em मा\raisebox{0.25ex}{\accentmark{15}{\char"1CD2}}त्वा\raisebox{0.25ex}{\accentmark{12}{\char"0951}}केचिन्नी\raisebox{0.15ex}{\accentmark{12}{\char"1CF9}}ये\raisebox{0.25ex}{\accentmark{12}{\char"0951}}मुर्वि\raisebox{0.25ex}{\accentmark{15}{\char"1CD2}}न्ना\raisebox{0.25ex}{\accentmark{15}{\char"1CD2}}पाशी\raisebox{0.25ex}{\accentmark{15}{\char"1CD2}}ना\raisebox{0.25ex}{\accentmark{12}{\char"0951}}ः\kern 0.05em ।\hskip 0.12em plus 0.04em minus 0.02em आतीधन्वेवतंई\raisebox{0.25ex}{\accentmark{12}{\char"0951}}हि \nolinebreak\hskip 0.08em \mbox{॥४॥}}} \par}

\vspace{0.8em}
                % --- SUBSECTIONS (Mantras) ---
                
                % Conditional rendering based on output_mode
                \par\filbreak

\phantomsection

\addcontentsline{toc}{subsection}{ऋक् ५}

{\centering \footnotesize \textcolor{AccentPurple}{॥ गौतमो (गौतमो वा) राहूगणः इन्द्रः बृहती ॥} \par}

\vspace{0.12em}

{\centering \textbf{\textcolor{AccentBlue}{त्वा\raisebox{0.25ex}{\accentmark{15}{\char"1CD2}}मङ्गाप्रा\raisebox{0.25ex}{\accentmark{15}{\char"1CD2}}शं\raisebox{0.25ex}{\accentmark{12}{\char"0951}}सिषः\kern 0.05em ।\hskip 0.12em plus 0.04em minus 0.02em देवा\raisebox{0.25ex}{\accentmark{15}{\char"1CD2}}श्शा\raisebox{0.25ex}{\accentmark{12}{\char"0951}}वीष्ठमा\raisebox{0.25ex}{\accentmark{15}{\char"1CD2}}र्क्त्यम्\kern 0.05em ।\hskip 0.12em plus 0.04em minus 0.02em न\raisebox{0.15ex}{\accentmark{12}{\char"1CF8}}त्वा\raisebox{0.25ex}{\accentmark{12}{\char"0951}}दन्यो\raisebox{0.25ex}{\accentmark{15}{\char"1CD2}}मा\raisebox{0.25ex}{\accentmark{12}{\char"0951}}घावन्नस्तिमार्डीत\kern 0.05em ।\hskip 0.12em plus 0.04em minus 0.02em इन्द्रब्रा\raisebox{0.25ex}{\accentmark{15}{\char"1CD2}}वी\raisebox{0.25ex}{\accentmark{12}{\char"0951}}मीतेवा\raisebox{0.25ex}{\accentmark{15}{\char"1CD2}}चा\raisebox{0.25ex}{\accentmark{12}{\char"0951}}ः \nolinebreak\hskip 0.08em \mbox{॥५॥}}} \par}

\vspace{0.8em}
                % --- SUBSECTIONS (Mantras) ---
                
                % Conditional rendering based on output_mode
                \par\filbreak

\phantomsection

\addcontentsline{toc}{subsection}{ऋक् ६}

{\centering \footnotesize \textcolor{AccentPurple}{॥ नृमेध पुरुमेधावांगिरसौ इन्द्रः बृहती ॥} \par}

\vspace{0.12em}

{\centering \textbf{\textcolor{AccentBlue}{त्वा\raisebox{0.25ex}{\accentmark{15}{\char"1CD2}}मि\raisebox{0.25ex}{\accentmark{15}{\char"1CD2}}न्द्रायाशा\raisebox{0.25ex}{\accentmark{15}{\char"1CD2}}आ\raisebox{0.25ex}{\accentmark{12}{\char"0951}}सि\kern 0.05em ।\hskip 0.12em plus 0.04em minus 0.02em ऋजीषीशावासस्पा\raisebox{0.25ex}{\accentmark{15}{\char"1CD2}}तिः\kern 0.05em ।\hskip 0.12em plus 0.04em minus 0.02em त्वंवृत्रा\raisebox{0.25ex}{\accentmark{15}{\char"1CD2}}णी\raisebox{0.25ex}{\accentmark{12}{\char"0951}}हंस्यप्रातीन्ये\raisebox{0.25ex}{\accentmark{15}{\char"1CD2}}काई\raisebox{0.25ex}{\accentmark{15}{\char"1CD2}}त्पूरू\raisebox{0.25ex}{\accentmark{15}{\char"1CD2}}ः\kern 0.05em ।\hskip 0.12em plus 0.04em minus 0.02em आनूक्तस्चर्षाणिधृ\raisebox{0.25ex}{\accentmark{15}{\char"1CD2}}त्तिः \nolinebreak\hskip 0.08em \mbox{॥६॥}}} \par}

\vspace{0.8em}
                % --- SUBSECTIONS (Mantras) ---
                
                % Conditional rendering based on output_mode
                \par\filbreak

\phantomsection

\addcontentsline{toc}{subsection}{ऋक् ७}

{\centering \footnotesize \textcolor{AccentPurple}{॥ मेधातिमध्येआतिथी (मध्येआतिथिः) काण्वः इन्द्रः बृहती ॥} \par}

\vspace{0.12em}

{\centering \textbf{\textcolor{AccentBlue}{इ\raisebox{0.25ex}{\accentmark{15}{\char"1CD2}}न्द्रामीद्देवातातये\kern 0.05em ।\hskip 0.12em plus 0.04em minus 0.02em इन्द्रं\raisebox{0.25ex}{\accentmark{12}{\char"0951}}प्रायत्याध्वारे\raisebox{0.25ex}{\accentmark{15}{\char"1CD2}}\kern 0.05em ।\hskip 0.12em plus 0.04em minus 0.02em इन्द्रंसामीके\raisebox{0.25ex}{\accentmark{15}{\char"1CD2}}वानी\raisebox{0.25ex}{\accentmark{15}{\char"1CD2}}नो\raisebox{0.25ex}{\accentmark{12}{\char"0951}}हवामहे\kern 0.05em ।\hskip 0.12em plus 0.04em minus 0.02em इ\raisebox{0.25ex}{\accentmark{15}{\char"1CD2}}न्द्र\raisebox{0.25ex}{\accentmark{12}{\char"0951}}न्धानास्यासा\raisebox{0.25ex}{\accentmark{15}{\char"1CD2}}ता\raisebox{0.25ex}{\accentmark{15}{\char"1CD2}}ये\raisebox{0.25ex}{\accentmark{12}{\char"0951}} \nolinebreak\hskip 0.08em \mbox{॥७॥}}} \par}

\vspace{0.8em}
                % --- SUBSECTIONS (Mantras) ---
                
                % Conditional rendering based on output_mode
                \par\filbreak

\phantomsection

\addcontentsline{toc}{subsection}{ऋक् ८}

{\centering \footnotesize \textcolor{AccentPurple}{॥ मेधातिमध्येआतिथी (मध्येआतिथिः) काण्वः इन्द्रः बृहती ॥} \par}

\vspace{0.12em}

{\centering \textbf{\textcolor{AccentBlue}{ईमा\raisebox{0.25ex}{\accentmark{15}{\char"1CD2}}ऊ\raisebox{0.25ex}{\accentmark{12}{\char"0951}}त्वापूरोवसो\kern 0.05em ।\hskip 0.12em plus 0.04em minus 0.02em गी\raisebox{0.25ex}{\accentmark{15}{\char"1CD2}}रो\raisebox{0.25ex}{\accentmark{12}{\char"0951}}वर्धन्तूया\raisebox{0.25ex}{\accentmark{15}{\char"1CD2}}मा\raisebox{0.25ex}{\accentmark{15}{\char"1CD2}}मा\raisebox{0.25ex}{\accentmark{12}{\char"0951}}\kern 0.05em ।\hskip 0.12em plus 0.04em minus 0.02em पवाका\raisebox{0.15ex}{\accentmark{12}{\char"1CF9}}व\raisebox{0.25ex}{\accentmark{12}{\char"0951}}र्णश्शू\raisebox{0.25ex}{\accentmark{15}{\char"1CD2}}चा\raisebox{0.25ex}{\accentmark{15}{\char"1CD2}}योवीप\raisebox{0.25ex}{\accentmark{15}{\char"1CD2}}श्ची\raisebox{0.25ex}{\accentmark{15}{\char"1CD2}}ता\raisebox{0.25ex}{\accentmark{15}{\char"1CD2}}ः\kern 0.05em ।\hskip 0.12em plus 0.04em minus 0.02em आभी\raisebox{0.25ex}{\accentmark{15}{\char"1CD2}}स्तोमै\raisebox{0.25ex}{\accentmark{12}{\char"0951}}रनूषत \nolinebreak\hskip 0.08em \mbox{॥८॥}}} \par}

\vspace{0.8em}
                % --- SUBSECTIONS (Mantras) ---
                
                % Conditional rendering based on output_mode
                \par\filbreak

\phantomsection

\addcontentsline{toc}{subsection}{ऋक् ९}

{\centering \footnotesize \textcolor{AccentPurple}{॥ मेधातिमध्येआतिथी (मध्येआतिथिः) काण्वः इन्द्रः बृहती ॥} \par}

\vspace{0.12em}

{\centering \textbf{\textcolor{AccentBlue}{ऊ\raisebox{0.25ex}{\accentmark{15}{\char"1CD2}}दूत्येमाधू\raisebox{0.25ex}{\accentmark{12}{\char"0951}}मक्तमः\kern 0.05em ।\hskip 0.12em plus 0.04em minus 0.02em गीरस्तोमा\raisebox{0.25ex}{\accentmark{12}{\char"0951}}सैरते\kern 0.05em ।\hskip 0.12em plus 0.04em minus 0.02em स\raisebox{0.25ex}{\accentmark{15}{\char"1CD2}}त्राजितो\raisebox{0.25ex}{\accentmark{12}{\char"0951}}धानासाआक्षी\raisebox{0.25ex}{\accentmark{12}{\char"0951}}तोतयः\kern 0.05em ।\hskip 0.12em plus 0.04em minus 0.02em वजय\raisebox{0.25ex}{\accentmark{12}{\char"0951}}न्तोरा\raisebox{0.25ex}{\accentmark{12}{\char"0951}}था\raisebox{0.25ex}{\accentmark{12}{\char"0951}}इवा \nolinebreak\hskip 0.08em \mbox{॥९॥}}} \par}

\vspace{0.8em}
                % --- SUBSECTIONS (Mantras) ---
                
                % Conditional rendering based on output_mode
                \par\filbreak

\phantomsection

\addcontentsline{toc}{subsection}{ऋक् १०}

{\centering \footnotesize \textcolor{AccentPurple}{॥ देवातिथिः काण्वः इन्द्रः बृहती ॥} \par}

\vspace{0.12em}

{\centering \textbf{\textcolor{AccentBlue}{या\raisebox{0.15ex}{\accentmark{12}{\char"1CF8}}था\raisebox{0.25ex}{\accentmark{12}{\char"0951}}गौरो\raisebox{0.25ex}{\accentmark{15}{\char"1CD2}}आपा\raisebox{0.25ex}{\accentmark{15}{\char"1CD2}}कृ\raisebox{0.25ex}{\accentmark{12}{\char"0951}}त\raisebox{0.25ex}{\accentmark{15}{\char"1CD2}}म्\kern 0.05em ।\hskip 0.12em plus 0.04em minus 0.02em तृ\raisebox{0.25ex}{\accentmark{15}{\char"1CD2}}ष्यन्ने\raisebox{0.25ex}{\accentmark{15}{\char"1CD2}}त्यवे\raisebox{0.25ex}{\accentmark{15}{\char"1CD2}}रि\raisebox{0.25ex}{\accentmark{15}{\char"1CD2}}णम्\kern 0.05em ।\hskip 0.12em plus 0.04em minus 0.02em आपि\raisebox{0.25ex}{\accentmark{15}{\char"1CD2}}त्वे\raisebox{0.25ex}{\accentmark{12}{\char"0951}}नप्रापित्वेतू\raisebox{0.25ex}{\accentmark{15}{\char"1CD2}}या\raisebox{0.25ex}{\accentmark{15}{\char"1CD2}}मा\raisebox{0.25ex}{\accentmark{15}{\char"1CD2}}ग\raisebox{0.25ex}{\accentmark{12}{\char"0951}}हि\kern 0.05em ।\hskip 0.12em plus 0.04em minus 0.02em क\raisebox{0.15ex}{\accentmark{12}{\char"1CF8}}ण्वेषूसू\raisebox{0.25ex}{\accentmark{15}{\char"1CD2}}साचा\raisebox{0.25ex}{\accentmark{15}{\char"1CD2}}पी\raisebox{0.25ex}{\accentmark{15}{\char"1CD2}}बा \nolinebreak\hskip 0.08em \mbox{॥१०॥}}} \par}

\vspace{0.8em}

            % --- End of Kandah Closing Floral Line & Counts ---
            \par\nopagebreak\vspace{0.9em plus 0.3em minus 0.1em}
            \centerline{\textcolor{black!35}{\rule{1.8cm}{0.35pt}}\;\raisebox{-0.22ex}{\textcolor{vedicmaroon}{\large\SixFlowerPetalDotted}}\;\textcolor{black!35}{\rule{1.8cm}{0.35pt}}}
            \nopagebreak\vspace{0.45em}
            \centerline{\devafont\textbf{\textcolor{AccentBlue}{॥ ऋक् सङ्ख्या : १० ॥}}}
        
            \clearpage

            % --- SECTION ---
            
            \phantomsection  

            \section*{\centering तृतीय खण्डः}
            \addcontentsline{toc}{section}{तृतीय खण्डः} 
            \markright{तृतीय खण्डः}
            \nopagebreak % Keep section title with first mantra
            
            % --- SECTION-LEVEL PROCEDURE LINK ---

                % --- SUBSECTIONS (Mantras) ---
                
                % Conditional rendering based on output_mode
                \par\filbreak

\phantomsection

\addcontentsline{toc}{subsection}{ऋक् १}

{\centering \footnotesize \textcolor{AccentPurple}{॥ भर्गः प्रागाथः इन्द्रः बृहती ॥} \par}

\vspace{0.12em}

{\centering \textbf{\textcolor{AccentBlue}{श\raisebox{0.25ex}{\accentmark{15}{\char"1CD2}}ग्धूटट्यषूशा\raisebox{0.25ex}{\accentmark{12}{\char"0951}}चीपते\kern 0.05em ।\hskip 0.12em plus 0.04em minus 0.02em इन्द्रावि\raisebox{0.25ex}{\accentmark{12}{\char"0951}}श्वा\raisebox{0.25ex}{\accentmark{15}{\char"1CD2}}भीरू\raisebox{0.15ex}{\accentmark{12}{\char"1CF8}}तीभिः\kern 0.05em ।\hskip 0.12em plus 0.04em minus 0.02em भा\raisebox{0.25ex}{\accentmark{15}{\char"1CD2}}गन्न\raisebox{0.15ex}{\accentmark{12}{\char"1CF9}}हित्वा\raisebox{0.25ex}{\accentmark{12}{\char"0951}}याशा\raisebox{0.25ex}{\accentmark{15}{\char"1CD2}}सं\raisebox{0.25ex}{\accentmark{12}{\char"0951}}वासूवी\raisebox{0.25ex}{\accentmark{15}{\char"1CD2}}दम्\kern 0.05em ।\hskip 0.12em plus 0.04em minus 0.02em आनू\raisebox{0.25ex}{\accentmark{12}{\char"0951}}शूराचा\raisebox{0.25ex}{\accentmark{15}{\char"1CD2}}रा\raisebox{0.25ex}{\accentmark{12}{\char"0951}}मसि \nolinebreak\hskip 0.08em \mbox{॥१॥}}} \par}

\vspace{0.8em}
                % --- SUBSECTIONS (Mantras) ---
                
                % Conditional rendering based on output_mode
                \par\filbreak

\phantomsection

\addcontentsline{toc}{subsection}{ऋक् २}

{\centering \footnotesize \textcolor{AccentPurple}{॥ रेभः कश्यपः इन्द्रः बृहती ॥} \par}

\vspace{0.12em}

{\centering \textbf{\textcolor{AccentBlue}{या\raisebox{0.15ex}{\accentmark{12}{\char"1CF8}}इ\raisebox{0.25ex}{\accentmark{12}{\char"0951}}न्द्राभू\raisebox{0.25ex}{\accentmark{15}{\char"1CD2}}जआ\raisebox{0.25ex}{\accentmark{15}{\char"1CD2}}भा\raisebox{0.25ex}{\accentmark{12}{\char"0951}}रा\kern 0.05em ।\hskip 0.12em plus 0.04em minus 0.02em स्वठ\raisebox{0.25ex}{\accentmark{12}{\char"0951}}र्वंआसू\raisebox{0.25ex}{\accentmark{12}{\char"0951}}रेभ्यः\kern 0.05em ।\hskip 0.12em plus 0.04em minus 0.02em स्तोता\raisebox{0.25ex}{\accentmark{15}{\char"1CD2}}रामिन्मा\raisebox{0.25ex}{\accentmark{12}{\char"0951}}घवन्नस्यावा\raisebox{0.25ex}{\accentmark{12}{\char"0951}}र्धया\kern 0.05em ।\hskip 0.12em plus 0.04em minus 0.02em ये\raisebox{0.25ex}{\accentmark{15}{\char"1CD2}}चत्वे\raisebox{0.25ex}{\accentmark{12}{\char"0951}}वृक्ता\raisebox{0.25ex}{\accentmark{15}{\char"1CD2}}ब\raisebox{0.25ex}{\accentmark{12}{\char"0951}}र्हीषः \nolinebreak\hskip 0.08em \mbox{॥२॥}}} \par}

\vspace{0.8em}
                % --- SUBSECTIONS (Mantras) ---
                
                % Conditional rendering based on output_mode
                \par\filbreak

\phantomsection

\addcontentsline{toc}{subsection}{ऋक् ३}

{\centering \footnotesize \textcolor{AccentPurple}{॥ जमदग्निभागर्वः मित्रावरुणादित्याः बृहती ॥} \par}

\vspace{0.12em}

{\centering \textbf{\textcolor{AccentBlue}{प्रा\raisebox{0.25ex}{\accentmark{15}{\char"1CD2}}मित्रा\raisebox{0.25ex}{\accentmark{15}{\char"1CD2}}यप्रार्यम्णे\raisebox{0.25ex}{\accentmark{12}{\char"0951}}\kern 0.05em ।\hskip 0.12em plus 0.04em minus 0.02em सा\raisebox{0.25ex}{\accentmark{15}{\char"1CD2}}चात्थ्याठ\raisebox{0.25ex}{\accentmark{12}{\char"0951}}मृतावसो\kern 0.05em ।\hskip 0.12em plus 0.04em minus 0.02em वरु\raisebox{0.25ex}{\accentmark{15}{\char"1CD2}}त्थ्येटट्य\raisebox{0.25ex}{\accentmark{15}{\char"1CD2}}वा\raisebox{0.15ex}{\accentmark{12}{\char"1CF9}}रूणेछन्द्यंवा\raisebox{0.25ex}{\accentmark{15}{\char"1CD2}}चा\raisebox{0.25ex}{\accentmark{12}{\char"0951}}ः\kern 0.05em ।\hskip 0.12em plus 0.04em minus 0.02em स्तोत्रं\raisebox{0.25ex}{\accentmark{12}{\char"0951}}रा\raisebox{0.25ex}{\accentmark{15}{\char"1CD2}}ज\raisebox{0.25ex}{\accentmark{12}{\char"0951}}सुगायता \nolinebreak\hskip 0.08em \mbox{॥३॥}}} \par}

\vspace{0.8em}
                % --- SUBSECTIONS (Mantras) ---
                
                % Conditional rendering based on output_mode
                \par\filbreak

\phantomsection

\addcontentsline{toc}{subsection}{ऋक् ४}

{\centering \footnotesize \textcolor{AccentPurple}{॥ मेधातिथिः काण्वः (मध्येआतिथिः काण्वः) इन्द्रः बृहती ॥} \par}

\vspace{0.12em}

{\centering \textbf{\textcolor{AccentBlue}{आ\raisebox{0.25ex}{\accentmark{15}{\char"1CD2}}भी\raisebox{0.15ex}{\accentmark{12}{\char"1CF9}}त्वा\raisebox{0.25ex}{\accentmark{12}{\char"0951}}पूर्वापी\raisebox{0.25ex}{\accentmark{12}{\char"0951}}तये\kern 0.05em ।\hskip 0.12em plus 0.04em minus 0.02em इ\raisebox{0.25ex}{\accentmark{15}{\char"1CD2}}न्द्रास्तो\raisebox{0.25ex}{\accentmark{12}{\char"0951}}मेभीरायावाः\kern 0.05em ।\hskip 0.12em plus 0.04em minus 0.02em समी\raisebox{0.25ex}{\accentmark{15}{\char"1CD2}}चीना\raisebox{0.15ex}{\accentmark{12}{\char"1CF9}}साऋभा\raisebox{0.25ex}{\accentmark{15}{\char"1CD2}}वस्सा\raisebox{0.25ex}{\accentmark{15}{\char"1CD2}}मा\raisebox{0.25ex}{\accentmark{15}{\char"1CD2}}स्वरन्\kern 0.05em ।\hskip 0.12em plus 0.04em minus 0.02em रुद्रागृ\raisebox{0.25ex}{\accentmark{12}{\char"0951}}णन्तापूर्व्यम् \nolinebreak\hskip 0.08em \mbox{॥४॥}}} \par}

\vspace{0.8em}
                % --- SUBSECTIONS (Mantras) ---
                
                % Conditional rendering based on output_mode
                \par\filbreak

\phantomsection

\addcontentsline{toc}{subsection}{ऋक् ५}

{\centering \footnotesize \textcolor{AccentPurple}{॥ नृमेध पुरुमेधावांगिरसौ इन्द्रः बृहती ॥} \par}

\vspace{0.12em}

{\centering \textbf{\textcolor{AccentBlue}{प्रा\raisebox{0.25ex}{\accentmark{15}{\char"1CD2}}वाइ\raisebox{0.25ex}{\accentmark{15}{\char"1CD2}}न्द्रा\raisebox{0.25ex}{\accentmark{12}{\char"0951}}यबृहा\raisebox{0.25ex}{\accentmark{12}{\char"0951}}ते\kern 0.05em ।\hskip 0.12em plus 0.04em minus 0.02em मारूतोब्र\raisebox{0.25ex}{\accentmark{15}{\char"1CD2}}ह्मा\raisebox{0.25ex}{\accentmark{12}{\char"0951}}र्च्चता\kern 0.05em ।\hskip 0.12em plus 0.04em minus 0.02em वृत्रंहा\raisebox{0.25ex}{\accentmark{15}{\char"1CD2}}नातिवृत्रा\raisebox{0.25ex}{\accentmark{15}{\char"1CD2}}हाशाता\raisebox{0.25ex}{\accentmark{15}{\char"1CD2}}क्रा\raisebox{0.25ex}{\accentmark{12}{\char"0951}}तु\raisebox{0.25ex}{\accentmark{12}{\char"0951}}ः\kern 0.05em ।\hskip 0.12em plus 0.04em minus 0.02em व\raisebox{0.25ex}{\accentmark{15}{\char"1CD2}}ज्ज्रे\raisebox{0.25ex}{\accentmark{12}{\char"0951}}णाशाता\raisebox{0.25ex}{\accentmark{15}{\char"1CD2}}पा\raisebox{0.25ex}{\accentmark{12}{\char"0951}}र्वणा \nolinebreak\hskip 0.08em \mbox{॥५॥}}} \par}

\vspace{0.8em}
                % --- SUBSECTIONS (Mantras) ---
                
                % Conditional rendering based on output_mode
                \par\filbreak

\phantomsection

\addcontentsline{toc}{subsection}{ऋक् ६}

{\centering \footnotesize \textcolor{AccentPurple}{॥ नृमेध पुरुमेधावांगिरसौ इन्द्रः बृहती ॥} \par}

\vspace{0.12em}

{\centering \textbf{\textcolor{AccentBlue}{बृहा\raisebox{0.25ex}{\accentmark{15}{\char"1CD2}}दिन्द्रा\raisebox{0.25ex}{\accentmark{12}{\char"0951}}यगायत\kern 0.05em ।\hskip 0.12em plus 0.04em minus 0.02em मारू\raisebox{0.25ex}{\accentmark{12}{\char"0951}}तोवृ\raisebox{0.25ex}{\accentmark{15}{\char"1CD2}}त्राहन्ता\raisebox{0.25ex}{\accentmark{12}{\char"0951}}मं\kern 0.05em ।\hskip 0.12em plus 0.04em minus 0.02em ये\raisebox{0.25ex}{\accentmark{15}{\char"1CD2}}नज्यो\raisebox{0.25ex}{\accentmark{15}{\char"1CD2}}तीरा\raisebox{0.25ex}{\accentmark{15}{\char"1CD2}}ज\raisebox{0.25ex}{\accentmark{12}{\char"0951}}नयन्नृता\raisebox{0.25ex}{\accentmark{15}{\char"1CD2}}वृधाः\raisebox{0.25ex}{\accentmark{12}{\char"0951}}\kern 0.05em ।\hskip 0.12em plus 0.04em minus 0.02em देव\raisebox{0.25ex}{\accentmark{15}{\char"1CD2}}न्देवा\raisebox{0.25ex}{\accentmark{15}{\char"1CD2}}याजा\raisebox{0.25ex}{\accentmark{15}{\char"1CD2}}ग्रवीम् \nolinebreak\hskip 0.08em \mbox{॥६॥}}} \par}

\vspace{0.8em}
                % --- SUBSECTIONS (Mantras) ---
                
                % Conditional rendering based on output_mode
                \par\filbreak

\phantomsection

\addcontentsline{toc}{subsection}{ऋक् ७}

{\centering \footnotesize \textcolor{AccentPurple}{॥ वसिष्ठो मैत्रावरुणिः इन्द्रः बृहती ॥} \par}

\vspace{0.12em}

{\centering \textbf{\textcolor{AccentBlue}{इन्द्राक्रा\raisebox{0.15ex}{\accentmark{12}{\char"1CF9}}तुं\raisebox{0.25ex}{\accentmark{12}{\char"0951}}नआ\raisebox{0.25ex}{\accentmark{15}{\char"1CD2}}भा\raisebox{0.25ex}{\accentmark{12}{\char"0951}}रा\kern 0.05em ।\hskip 0.12em plus 0.04em minus 0.02em पिता\raisebox{0.25ex}{\accentmark{15}{\char"1CD2}}पुत्रे\raisebox{0.25ex}{\accentmark{15}{\char"1CD2}}भ्योया\raisebox{0.25ex}{\accentmark{15}{\char"1CD2}}था\kern 0.05em ।\hskip 0.12em plus 0.04em minus 0.02em शिक्षा\raisebox{0.25ex}{\accentmark{12}{\char"0951}}णोअस्मिन्पु\raisebox{0.25ex}{\accentmark{12}{\char"0951}}रुहूताया\raisebox{0.25ex}{\accentmark{15}{\char"1CD2}}मा\raisebox{0.25ex}{\accentmark{12}{\char"0951}}नी\kern 0.05em ।\hskip 0.12em plus 0.04em minus 0.02em जी\raisebox{0.25ex}{\accentmark{15}{\char"1CD2}}वाज्ज्योती\raisebox{0.25ex}{\accentmark{12}{\char"0951}}रशीमहि \nolinebreak\hskip 0.08em \mbox{॥७॥}}} \par}

\vspace{0.8em}
                % --- SUBSECTIONS (Mantras) ---
                
                % Conditional rendering based on output_mode
                \par\filbreak

\phantomsection

\addcontentsline{toc}{subsection}{ऋक् ८}

{\centering \footnotesize \textcolor{AccentPurple}{॥ रेभः कश्यपः इन्द्रः बृहती ॥} \par}

\vspace{0.12em}

{\centering \textbf{\textcolor{AccentBlue}{माना\raisebox{0.25ex}{\accentmark{12}{\char"0951}}इन्द्रापारा\raisebox{0.25ex}{\accentmark{12}{\char"0951}}वृणात्\kern 0.05em ।\hskip 0.12em plus 0.04em minus 0.02em भा\raisebox{0.25ex}{\accentmark{15}{\char"1CD2}}वा\raisebox{0.25ex}{\accentmark{12}{\char"0951}}नस्साधामाद्याः\raisebox{0.25ex}{\accentmark{12}{\char"0951}}\kern 0.05em ।\hskip 0.12em plus 0.04em minus 0.02em त्व\raisebox{0.15ex}{\accentmark{12}{\char"1CF9}}न्ना\raisebox{0.25ex}{\accentmark{12}{\char"0951}}ऊतीस्त्वामिन्ना\raisebox{0.25ex}{\accentmark{12}{\char"0951}}आप्यं\raisebox{0.25ex}{\accentmark{12}{\char"0951}}\kern 0.05em ।\hskip 0.12em plus 0.04em minus 0.02em मानाइन्द्रापा\raisebox{0.25ex}{\accentmark{15}{\char"1CD2}}रा\raisebox{0.25ex}{\accentmark{12}{\char"0951}}वृणात् \nolinebreak\hskip 0.08em \mbox{॥८॥}}} \par}

\vspace{0.8em}
                % --- SUBSECTIONS (Mantras) ---
                
                % Conditional rendering based on output_mode
                \par\filbreak

\phantomsection

\addcontentsline{toc}{subsection}{ऋक् ९}

{\centering \footnotesize \textcolor{AccentPurple}{॥ मेधातिथिः काण्वः (मध्येआतिथिः काण्वः) इन्द्रः बृहती ॥} \par}

\vspace{0.12em}

{\centering \textbf{\textcolor{AccentBlue}{वा\raisebox{0.25ex}{\accentmark{12}{\char"0951}}य\raisebox{0.25ex}{\accentmark{15}{\char"1CD2}}ङ्घात्वा\raisebox{0.25ex}{\accentmark{12}{\char"0951}}सूता\raisebox{0.25ex}{\accentmark{15}{\char"1CD2}}वन्तः\kern 0.05em ।\hskip 0.12em plus 0.04em minus 0.02em आ\raisebox{0.25ex}{\accentmark{15}{\char"1CD2}}पोना\raisebox{0.25ex}{\accentmark{15}{\char"1CD2}}वृक्ताब\raisebox{0.25ex}{\accentmark{12}{\char"0951}}र्हिषः\kern 0.05em ।\hskip 0.12em plus 0.04em minus 0.02em पा\raisebox{0.25ex}{\accentmark{15}{\char"1CD2}}वी\raisebox{0.25ex}{\accentmark{15}{\char"1CD2}}त्रास्याप्रस्रा\raisebox{0.25ex}{\accentmark{15}{\char"1CD2}}वा\raisebox{0.25ex}{\accentmark{12}{\char"0951}}णेषुवृत्रहन्\kern 0.05em ।\hskip 0.12em plus 0.04em minus 0.02em पा\raisebox{0.15ex}{\accentmark{12}{\char"1CF8}}रीस्तोता\raisebox{0.25ex}{\accentmark{15}{\char"1CD2}}रा\raisebox{0.25ex}{\accentmark{12}{\char"0951}}आसते \nolinebreak\hskip 0.08em \mbox{॥९॥}}} \par}

\vspace{0.8em}
                % --- SUBSECTIONS (Mantras) ---
                
                % Conditional rendering based on output_mode
                \par\filbreak

\phantomsection

\addcontentsline{toc}{subsection}{ऋक् १०}

{\centering \footnotesize \textcolor{AccentPurple}{॥ भरद्वाजः (शयुः) बार्हस्पत्यः इन्द्रः बृहती ॥} \par}

\vspace{0.12em}

{\centering \textbf{\textcolor{AccentBlue}{यादिन्द्रा\raisebox{0.25ex}{\accentmark{12}{\char"0951}}ना\raisebox{0.15ex}{\accentmark{12}{\char"1CF9}}हू\raisebox{0.25ex}{\accentmark{12}{\char"0951}}षीष्वा\kern 0.05em ।\hskip 0.12em plus 0.04em minus 0.02em ओजो\raisebox{0.25ex}{\accentmark{12}{\char"0951}}निम्णं\raisebox{0.15ex}{\accentmark{12}{\char"1CF9}}चा\raisebox{0.25ex}{\accentmark{12}{\char"0951}}कृष्टीषु\raisebox{0.25ex}{\accentmark{12}{\char"0951}}\kern 0.05em ।\hskip 0.12em plus 0.04em minus 0.02em य\raisebox{0.25ex}{\accentmark{15}{\char"1CD2}}द्वाप\raisebox{0.25ex}{\accentmark{15}{\char"1CD2}}ञ्चा\raisebox{0.25ex}{\accentmark{12}{\char"0951}}क्षीतीनां\raisebox{0.25ex}{\accentmark{15}{\char"1CD2}}द्युम्नमाभा\raisebox{0.25ex}{\accentmark{12}{\char"0951}}रा\kern 0.05em ।\hskip 0.12em plus 0.04em minus 0.02em सत्रा\raisebox{0.15ex}{\accentmark{12}{\char"1CF9}}वि\raisebox{0.15ex}{\accentmark{12}{\char"1CF9}}श्वा\raisebox{0.25ex}{\accentmark{12}{\char"0951}}नीपौं\raisebox{0.25ex}{\accentmark{15}{\char"1CD2}}१स्यम् \nolinebreak\hskip 0.08em \mbox{॥१०॥}}} \par}

\vspace{0.8em}

            % --- End of Kandah Closing Floral Line & Counts ---
            \par\nopagebreak\vspace{0.9em plus 0.3em minus 0.1em}
            \centerline{\textcolor{black!35}{\rule{1.8cm}{0.35pt}}\;\raisebox{-0.22ex}{\textcolor{vedicmaroon}{\large\SixFlowerPetalDotted}}\;\textcolor{black!35}{\rule{1.8cm}{0.35pt}}}
            \nopagebreak\vspace{0.45em}
            \centerline{\devafont\textbf{\textcolor{AccentBlue}{॥ ऋक् सङ्ख्या : १० ॥}}}
        
            \clearpage

            % --- SECTION ---
            
            \phantomsection  

            \section*{\centering चतुर्थ खण्डः}
            \addcontentsline{toc}{section}{चतुर्थ खण्डः} 
            \markright{चतुर्थ खण्डः}
            \nopagebreak % Keep section title with first mantra
            
            % --- SECTION-LEVEL PROCEDURE LINK ---

                % --- SUBSECTIONS (Mantras) ---
                
                % Conditional rendering based on output_mode
                \par\filbreak

\phantomsection

\addcontentsline{toc}{subsection}{ऋक् १}

{\centering \footnotesize \textcolor{AccentPurple}{॥ मेधातिथिः (मध्येआतिथिः) काण्वः इन्द्रः बृहती ॥} \par}

\vspace{0.12em}

{\centering \textbf{\textcolor{AccentBlue}{सत्य\raisebox{0.25ex}{\accentmark{15}{\char"1CD2}}मित्थावृषेद\raisebox{0.25ex}{\accentmark{12}{\char"0951}}सि\kern 0.05em ।\hskip 0.12em plus 0.04em minus 0.02em वृ\raisebox{0.25ex}{\accentmark{15}{\char"1CD2}}र्षा\raisebox{0.25ex}{\accentmark{12}{\char"0951}}जूतिर्नोवीता\kern 0.05em ।\hskip 0.12em plus 0.04em minus 0.02em वृ\raisebox{0.25ex}{\accentmark{15}{\char"1CD2}}षाह्युठे\raisebox{0.25ex}{\accentmark{12}{\char"0951}}ग्रश्रृण्वीषे\raisebox{0.15ex}{\accentmark{12}{\char"1CF8}}पा\raisebox{0.25ex}{\accentmark{12}{\char"0951}}रावा\raisebox{0.25ex}{\accentmark{12}{\char"0951}}ती\kern 0.05em ।\hskip 0.12em plus 0.04em minus 0.02em वृषोअर्वावा\raisebox{0.15ex}{\accentmark{12}{\char"1CF9}}ती\raisebox{0.25ex}{\accentmark{12}{\char"0951}}श्रुता\raisebox{0.15ex}{\accentmark{12}{\char"1CF8}}ः \nolinebreak\hskip 0.08em \mbox{॥१॥}}} \par}

\vspace{0.8em}
                % --- SUBSECTIONS (Mantras) ---
                
                % Conditional rendering based on output_mode
                \par\filbreak

\phantomsection

\addcontentsline{toc}{subsection}{ऋक् २}

{\centering \footnotesize \textcolor{AccentPurple}{॥ रेभः कश्यपः इन्द्रः बृहती ॥} \par}

\vspace{0.12em}

{\centering \textbf{\textcolor{AccentBlue}{या\raisebox{0.25ex}{\accentmark{15}{\char"1CD2}}च्चाक्रासी\raisebox{0.25ex}{\accentmark{12}{\char"0951}}पा\raisebox{0.25ex}{\accentmark{12}{\char"0951}}रावा\raisebox{0.25ex}{\accentmark{15}{\char"1CD2}}ति\kern 0.05em ।\hskip 0.12em plus 0.04em minus 0.02em याद\raisebox{0.25ex}{\accentmark{12}{\char"0951}}र्वावा\raisebox{0.25ex}{\accentmark{15}{\char"1CD2}}तीवृत्रहन्\kern 0.05em ।\hskip 0.12em plus 0.04em minus 0.02em आ\raisebox{0.25ex}{\accentmark{15}{\char"1CD2}}ता\raisebox{0.25ex}{\accentmark{12}{\char"0951}}स्त्वागीर्भिद्युगदिन्द्राकेशी\raisebox{0.25ex}{\accentmark{15}{\char"1CD2}}भीः\kern 0.05em ।\hskip 0.12em plus 0.04em minus 0.02em सूता\raisebox{0.25ex}{\accentmark{15}{\char"1CD2}}वंआ\raisebox{0.25ex}{\accentmark{15}{\char"1CD2}}वी\raisebox{0.25ex}{\accentmark{15}{\char"1CD2}}वा\raisebox{0.25ex}{\accentmark{12}{\char"0951}}सति \nolinebreak\hskip 0.08em \mbox{॥२॥}}} \par}

\vspace{0.8em}
                % --- SUBSECTIONS (Mantras) ---
                
                % Conditional rendering based on output_mode
                \par\filbreak

\phantomsection

\addcontentsline{toc}{subsection}{ऋक् ३}

{\centering \footnotesize \textcolor{AccentPurple}{॥ वत्सः (वशोऽश्वयः) इन्द्रः बृहती ॥} \par}

\vspace{0.12em}

{\centering \textbf{\textcolor{AccentBlue}{आभीवो\raisebox{0.25ex}{\accentmark{12}{\char"0951}}वीर\raisebox{0.15ex}{\accentmark{12}{\char"1CF9}}म\raisebox{0.25ex}{\accentmark{15}{\char"1CD2}}न्धा\raisebox{0.25ex}{\accentmark{12}{\char"0951}}सामादेषुगाया\kern 0.05em ।\hskip 0.12em plus 0.04em minus 0.02em गीरा\raisebox{0.25ex}{\accentmark{15}{\char"1CD2}}माहावी\raisebox{0.25ex}{\accentmark{15}{\char"1CD2}}चे\raisebox{0.25ex}{\accentmark{12}{\char"0951}}तसम्\kern 0.05em ।\hskip 0.12em plus 0.04em minus 0.02em इ\raisebox{0.25ex}{\accentmark{15}{\char"1CD2}}न्द्रन्नामश्रू\raisebox{0.15ex}{\accentmark{12}{\char"1CF9}}त्यंशाकीनं\raisebox{0.15ex}{\accentmark{12}{\char"1CF9}}वाचोयाथा \nolinebreak\hskip 0.08em \mbox{॥३॥}}} \par}

\vspace{0.8em}
                % --- SUBSECTIONS (Mantras) ---
                
                % Conditional rendering based on output_mode
                \par\filbreak

\phantomsection

\addcontentsline{toc}{subsection}{ऋक् ४}

{\centering \footnotesize \textcolor{AccentPurple}{॥ भरद्वाजः (शयुः) बार्हस्पत्यः इन्द्रः बृहती ॥} \par}

\vspace{0.12em}

{\centering \textbf{\textcolor{AccentBlue}{इ\raisebox{0.15ex}{\accentmark{12}{\char"1CF8}}न्द्रा\raisebox{0.25ex}{\accentmark{12}{\char"0951}}त्रीधा\raisebox{0.25ex}{\accentmark{15}{\char"1CD2}}तू\raisebox{0.25ex}{\accentmark{12}{\char"0951}}शाराणम्\raisebox{0.25ex}{\accentmark{15}{\char"1CD2}}\kern 0.05em ।\hskip 0.12em plus 0.04em minus 0.02em त्रीवा\raisebox{0.25ex}{\accentmark{15}{\char"1CD2}}रू\raisebox{0.25ex}{\accentmark{15}{\char"1CD2}}थंस्वस्ताये\raisebox{0.15ex}{\accentmark{12}{\char"1CF9}}\kern 0.05em ।\hskip 0.12em plus 0.04em minus 0.02em छर्दिया\raisebox{0.25ex}{\accentmark{12}{\char"0951}}च्छामाघा\raisebox{0.25ex}{\accentmark{12}{\char"0951}}वा\raisebox{0.25ex}{\accentmark{12}{\char"0951}}भ्या\raisebox{0.25ex}{\accentmark{12}{\char"0951}}श्चामंह्य\raisebox{0.25ex}{\accentmark{12}{\char"0951}}न्च\kern 0.05em ।\hskip 0.12em plus 0.04em minus 0.02em यावा\raisebox{0.15ex}{\accentmark{12}{\char"1CF9}}यादिद्यू\raisebox{0.25ex}{\accentmark{15}{\char"1CD2}}मे\raisebox{0.25ex}{\accentmark{12}{\char"0951}}भ्यः \nolinebreak\hskip 0.08em \mbox{॥४॥}}} \par}

\vspace{0.8em}
                % --- SUBSECTIONS (Mantras) ---
                
                % Conditional rendering based on output_mode
                \par\filbreak

\phantomsection

\addcontentsline{toc}{subsection}{ऋक् ५}

{\centering \footnotesize \textcolor{AccentPurple}{॥ नृमेध आंगिरसः इन्द्रः बृहती ॥} \par}

\vspace{0.12em}

{\centering \textbf{\textcolor{AccentBlue}{श्रा\raisebox{0.25ex}{\accentmark{15}{\char"1CD2}}यन्ता\raisebox{0.25ex}{\accentmark{12}{\char"0951}}ईवासूर्य\raisebox{0.25ex}{\accentmark{12}{\char"0951}}\kern 0.05em ।\hskip 0.12em plus 0.04em minus 0.02em विश्वे\raisebox{0.25ex}{\accentmark{15}{\char"1CD2}}दिन्द्रा\raisebox{0.25ex}{\accentmark{12}{\char"0951}}स्यभक्षत\kern 0.05em ।\hskip 0.12em plus 0.04em minus 0.02em वासू\raisebox{0.25ex}{\accentmark{15}{\char"1CD2}}नीजातो\raisebox{0.15ex}{\accentmark{12}{\char"1CF9}}ज\raisebox{0.25ex}{\accentmark{12}{\char"0951}}निमान्यो\raisebox{0.25ex}{\accentmark{15}{\char"1CD2}}ज\raisebox{0.25ex}{\accentmark{12}{\char"0951}}सा\kern 0.05em ।\hskip 0.12em plus 0.04em minus 0.02em प्रा\raisebox{0.15ex}{\accentmark{12}{\char"1CF8}}ती\raisebox{0.25ex}{\accentmark{12}{\char"0951}}भा\raisebox{0.25ex}{\accentmark{12}{\char"0951}}ग\raisebox{0.25ex}{\accentmark{15}{\char"1CD2}}न्ना\raisebox{0.25ex}{\accentmark{15}{\char"1CD2}}दी\raisebox{0.25ex}{\accentmark{12}{\char"0951}}धीमः \nolinebreak\hskip 0.08em \mbox{॥५॥}}} \par}

\vspace{0.8em}
                % --- SUBSECTIONS (Mantras) ---
                
                % Conditional rendering based on output_mode
                \par\filbreak

\phantomsection

\addcontentsline{toc}{subsection}{ऋक् ६}

{\centering \footnotesize \textcolor{AccentPurple}{॥ पुरुहन्मा आंगिरसः इन्द्रः बृहती ॥} \par}

\vspace{0.12em}

{\centering \textbf{\textcolor{AccentBlue}{ना\raisebox{0.25ex}{\accentmark{15}{\char"1CD2}}सीमादे\raisebox{0.25ex}{\accentmark{12}{\char"0951}}वआपाता\raisebox{0.25ex}{\accentmark{15}{\char"1CD2}}त्\kern 0.05em ।\hskip 0.12em plus 0.04em minus 0.02em ईळन्दी\raisebox{0.25ex}{\accentmark{12}{\char"0951}}र्घा\raisebox{0.25ex}{\accentmark{12}{\char"0951}}योमर्क्त्याः\kern 0.05em ।\hskip 0.12em plus 0.04em minus 0.02em एता\raisebox{0.25ex}{\accentmark{12}{\char"0951}}ग्वाचिद्या\raisebox{0.25ex}{\accentmark{15}{\char"1CD2}}एता\raisebox{0.25ex}{\accentmark{12}{\char"0951}}शोयूयो\raisebox{0.25ex}{\accentmark{15}{\char"1CD2}}जाते\kern 0.05em ।\hskip 0.12em plus 0.04em minus 0.02em इन्द्रोहा\raisebox{0.15ex}{\accentmark{12}{\char"1CF9}}रीयूयो\raisebox{0.25ex}{\accentmark{15}{\char"1CD2}}जा\raisebox{0.25ex}{\accentmark{12}{\char"0951}}ते \nolinebreak\hskip 0.08em \mbox{॥६॥}}} \par}

\vspace{0.8em}
                % --- SUBSECTIONS (Mantras) ---
                
                % Conditional rendering based on output_mode
                \par\filbreak

\phantomsection

\addcontentsline{toc}{subsection}{ऋक् ७}

{\centering \footnotesize \textcolor{AccentPurple}{॥ नृमेध{-} पुरुमेधाषांगिरसौ इन्द्रः बृहती ॥} \par}

\vspace{0.12em}

{\centering \textbf{\textcolor{AccentBlue}{आ\raisebox{0.25ex}{\accentmark{15}{\char"1CD2}}नोवि\raisebox{0.15ex}{\accentmark{12}{\char"1CF8}}श्वासुहा\raisebox{0.25ex}{\accentmark{15}{\char"1CD2}}व्यम्\raisebox{0.25ex}{\accentmark{12}{\char"0951}}\kern 0.05em ।\hskip 0.12em plus 0.04em minus 0.02em इ\raisebox{0.15ex}{\accentmark{12}{\char"1CF8}}न्द्रं\raisebox{0.25ex}{\accentmark{12}{\char"0951}}सामात्सू\raisebox{0.25ex}{\accentmark{12}{\char"0951}}भूषत\kern 0.05em ।\hskip 0.12em plus 0.04em minus 0.02em ऊ\raisebox{0.25ex}{\accentmark{15}{\char"1CD2}}पाब्र\raisebox{0.15ex}{\accentmark{12}{\char"1CF9}}ह्णा\raisebox{0.25ex}{\accentmark{12}{\char"0951}}णीसा\raisebox{0.25ex}{\accentmark{15}{\char"1CD2}}वा\raisebox{0.25ex}{\accentmark{12}{\char"0951}}नानिवृत्रहन्\kern 0.05em ।\hskip 0.12em plus 0.04em minus 0.02em परमाज्यऋचीषमा \nolinebreak\hskip 0.08em \mbox{॥७॥}}} \par}

\vspace{0.8em}
                % --- SUBSECTIONS (Mantras) ---
                
                % Conditional rendering based on output_mode
                \par\filbreak

\phantomsection

\addcontentsline{toc}{subsection}{ऋक् ८}

{\centering \footnotesize \textcolor{AccentPurple}{॥ वसिष्ठो मैत्रावरुणिः इन्द्रः बृहती ॥} \par}

\vspace{0.12em}

{\centering \textbf{\textcolor{AccentBlue}{त\raisebox{0.15ex}{\accentmark{12}{\char"1CF8}}वेदि\raisebox{0.25ex}{\accentmark{12}{\char"0951}}न्द्रावा\raisebox{0.25ex}{\accentmark{15}{\char"1CD2}}मंवा\raisebox{0.25ex}{\accentmark{15}{\char"1CD2}}सू\kern 0.05em ।\hskip 0.12em plus 0.04em minus 0.02em त्वंपू\raisebox{0.25ex}{\accentmark{12}{\char"0951}}ष्यसिमाध्यामम्\kern 0.05em ।\hskip 0.12em plus 0.04em minus 0.02em सत्रा\raisebox{0.25ex}{\accentmark{15}{\char"1CD2}}विश्वास्यापारामास्या\raisebox{0.25ex}{\accentmark{12}{\char"0951}}राजसि\kern 0.05em ।\hskip 0.12em plus 0.04em minus 0.02em ना\raisebox{0.15ex}{\accentmark{12}{\char"1CF8}}कीष्ट्वा\raisebox{0.15ex}{\accentmark{12}{\char"1CF8}}गोषू\raisebox{0.25ex}{\accentmark{12}{\char"0951}}वृण्वते \nolinebreak\hskip 0.08em \mbox{॥८॥}}} \par}

\vspace{0.8em}
                % --- SUBSECTIONS (Mantras) ---
                
                % Conditional rendering based on output_mode
                \par\filbreak

\phantomsection

\addcontentsline{toc}{subsection}{ऋक् ९}

{\centering \footnotesize \textcolor{AccentPurple}{॥ मेधातिथि{-}मध्येआतिथी काण्वौ इन्द्रः बृहती ॥} \par}

\vspace{0.12em}

{\centering \textbf{\textcolor{AccentBlue}{क्वे\raisebox{0.15ex}{\accentmark{12}{\char"1CF8}}या\raisebox{0.25ex}{\accentmark{12}{\char"0951}}थक्वे\raisebox{0.25ex}{\accentmark{15}{\char"1CD2}}दसि\kern 0.05em ।\hskip 0.12em plus 0.04em minus 0.02em पुरूत्रा\raisebox{0.25ex}{\accentmark{15}{\char"1CD2}}चिद्धितेमानाः\kern 0.05em ।\hskip 0.12em plus 0.04em minus 0.02em आलाऋषीयुध्मखजकृत्पुरन्दरा\kern 0.05em ।\hskip 0.12em plus 0.04em minus 0.02em प्रा\raisebox{0.15ex}{\accentmark{12}{\char"1CF8}}गा\raisebox{0.25ex}{\accentmark{12}{\char"0951}}यत्राआ\raisebox{0.25ex}{\accentmark{12}{\char"0951}}गासिषुः \nolinebreak\hskip 0.08em \mbox{॥९॥}}} \par}

\vspace{0.8em}
                % --- SUBSECTIONS (Mantras) ---
                
                % Conditional rendering based on output_mode
                \par\filbreak

\phantomsection

\addcontentsline{toc}{subsection}{ऋक् १०}

{\centering \footnotesize \textcolor{AccentPurple}{॥ कलिः प्रागाथः इन्द्रः बृहती ॥} \par}

\vspace{0.12em}

{\centering \textbf{\textcolor{AccentBlue}{वायामे\raisebox{0.25ex}{\accentmark{12}{\char"0951}}नामीदाह्यः\raisebox{0.25ex}{\accentmark{12}{\char"0951}}\kern 0.05em ।\hskip 0.12em plus 0.04em minus 0.02em आ\raisebox{0.25ex}{\accentmark{15}{\char"1CD2}}पी\raisebox{0.25ex}{\accentmark{12}{\char"0951}}पेमेहा\raisebox{0.25ex}{\accentmark{15}{\char"1CD2}}वज्री\raisebox{0.25ex}{\accentmark{12}{\char"0951}}णं\kern 0.05em ।\hskip 0.12em plus 0.04em minus 0.02em तस्मा\raisebox{0.25ex}{\accentmark{12}{\char"0951}}ऊआद्या\raisebox{0.25ex}{\accentmark{15}{\char"1CD2}}सा\raisebox{0.25ex}{\accentmark{15}{\char"1CD2}}वानेसूतं\raisebox{0.25ex}{\accentmark{15}{\char"1CD2}}भा\raisebox{0.25ex}{\accentmark{12}{\char"0951}}रा\kern 0.05em ।\hskip 0.12em plus 0.04em minus 0.02em आनू\raisebox{0.25ex}{\accentmark{15}{\char"1CD2}}नंभूषताश्रूते\raisebox{0.25ex}{\accentmark{15}{\char"1CD2}} \nolinebreak\hskip 0.08em \mbox{॥१०॥}}} \par}

\vspace{0.8em}

            % --- End of Kandah Closing Floral Line & Counts ---
            \par\nopagebreak\vspace{0.9em plus 0.3em minus 0.1em}
            \centerline{\textcolor{black!35}{\rule{1.8cm}{0.35pt}}\;\raisebox{-0.22ex}{\textcolor{vedicmaroon}{\large\SixFlowerPetalDotted}}\;\textcolor{black!35}{\rule{1.8cm}{0.35pt}}}
            \nopagebreak\vspace{0.45em}
            \centerline{\devafont\textbf{\textcolor{AccentBlue}{॥ ऋक् सङ्ख्या : १० ॥}}}
        
            \clearpage

            % --- SECTION ---
            
            \phantomsection  

            \section*{\centering पञ्चम खण्डः}
            \addcontentsline{toc}{section}{पञ्चम खण्डः} 
            \markright{पञ्चम खण्डः}
            \nopagebreak % Keep section title with first mantra
            
            % --- SECTION-LEVEL PROCEDURE LINK ---

                % --- SUBSECTIONS (Mantras) ---
                
                % Conditional rendering based on output_mode
                \par\filbreak

\phantomsection

\addcontentsline{toc}{subsection}{ऋक् १}

{\centering \footnotesize \textcolor{AccentPurple}{॥ पुरुहन्मा आंगिरसः इन्द्रः बृहती ॥} \par}

\vspace{0.12em}

{\centering \textbf{\textcolor{AccentBlue}{योरा\raisebox{0.25ex}{\accentmark{15}{\char"1CD2}}जा\raisebox{0.25ex}{\accentmark{12}{\char"0951}}चर्षा\raisebox{0.25ex}{\accentmark{12}{\char"0951}}णीनाम्\kern 0.05em ।\hskip 0.12em plus 0.04em minus 0.02em या\raisebox{0.25ex}{\accentmark{15}{\char"1CD2}}तारा\raisebox{0.15ex}{\accentmark{12}{\char"1CF9}}थेभीरा\raisebox{0.25ex}{\accentmark{15}{\char"1CD2}}ध्री\raisebox{0.25ex}{\accentmark{12}{\char"0951}}गुः\kern 0.05em ।\hskip 0.12em plus 0.04em minus 0.02em वि\raisebox{0.25ex}{\accentmark{15}{\char"1CD2}}श्वा\raisebox{0.25ex}{\accentmark{12}{\char"0951}}सान्तारु\raisebox{0.25ex}{\accentmark{15}{\char"1CD2}}तापृ\raisebox{0.25ex}{\accentmark{15}{\char"1CD2}}तानानां\kern 0.05em ।\hskip 0.12em plus 0.04em minus 0.02em जे\raisebox{0.25ex}{\accentmark{15}{\char"1CD2}}ष्ठंयो\raisebox{0.15ex}{\accentmark{12}{\char"1CF9}}वृत्राहा\raisebox{0.25ex}{\accentmark{15}{\char"1CD2}}गृणे \nolinebreak\hskip 0.08em \mbox{॥१॥}}} \par}

\vspace{0.8em}
                % --- SUBSECTIONS (Mantras) ---
                
                % Conditional rendering based on output_mode
                \par\filbreak

\phantomsection

\addcontentsline{toc}{subsection}{ऋक् २}

{\centering \footnotesize \textcolor{AccentPurple}{॥ भर्गः प्रागाथः इन्द्रः बृहती ॥} \par}

\vspace{0.12em}

{\centering \textbf{\textcolor{AccentBlue}{याता\raisebox{0.25ex}{\accentmark{12}{\char"0951}}इन्द्राभाया\raisebox{0.25ex}{\accentmark{12}{\char"0951}}महे\kern 0.05em ।\hskip 0.12em plus 0.04em minus 0.02em ता\raisebox{0.15ex}{\accentmark{12}{\char"1CF8}}तो\raisebox{0.25ex}{\accentmark{12}{\char"0951}}नोआ\raisebox{0.25ex}{\accentmark{15}{\char"1CD2}}भा\raisebox{0.25ex}{\accentmark{12}{\char"0951}}यंकृ\raisebox{0.25ex}{\accentmark{12}{\char"0951}}धि\kern 0.05em ।\hskip 0.12em plus 0.04em minus 0.02em मा\raisebox{0.25ex}{\accentmark{15}{\char"1CD2}}घा\raisebox{0.25ex}{\accentmark{12}{\char"0951}}वन्शग्धिता\raisebox{0.15ex}{\accentmark{12}{\char"1CF9}}वा\raisebox{0.25ex}{\accentmark{12}{\char"0951}}तन्नाऊता\raisebox{0.25ex}{\accentmark{15}{\char"1CD2}}ये\raisebox{0.25ex}{\accentmark{12}{\char"0951}}\kern 0.05em ।\hskip 0.12em plus 0.04em minus 0.02em वि\raisebox{0.25ex}{\accentmark{15}{\char"1CD2}}द्वि\raisebox{0.25ex}{\accentmark{15}{\char"1CD2}}षोविमृधो\raisebox{0.25ex}{\accentmark{12}{\char"0951}}जहि \nolinebreak\hskip 0.08em \mbox{॥२॥}}} \par}

\vspace{0.8em}
                % --- SUBSECTIONS (Mantras) ---
                
                % Conditional rendering based on output_mode
                \par\filbreak

\phantomsection

\addcontentsline{toc}{subsection}{ऋक् ३}

{\centering \footnotesize \textcolor{AccentPurple}{॥ इरिम्बिठिः काण्वः इन्द्रः (वास्तोष्पति वा) बृहती ॥} \par}

\vspace{0.12em}

{\centering \textbf{\textcolor{AccentBlue}{वास्तो\raisebox{0.25ex}{\accentmark{12}{\char"0951}}ष्पतेध्रूवा\raisebox{0.25ex}{\accentmark{15}{\char"1CD2}}स्थूणाम्\kern 0.05em ।\hskip 0.12em plus 0.04em minus 0.02em अंसत्रं\raisebox{0.25ex}{\accentmark{15}{\char"1CD2}}सौ\raisebox{0.25ex}{\accentmark{12}{\char"0951}}म्या\raisebox{0.25ex}{\accentmark{15}{\char"1CD2}}नाम्\kern 0.05em ।\hskip 0.12em plus 0.04em minus 0.02em द्रप्स\raisebox{0.25ex}{\accentmark{15}{\char"1CD2}}ःपूरां\raisebox{0.25ex}{\accentmark{15}{\char"1CD2}}भेक्ताशश्वा\raisebox{0.25ex}{\accentmark{12}{\char"0951}}तीनाम्\kern 0.05em ।\hskip 0.12em plus 0.04em minus 0.02em इ\raisebox{0.25ex}{\accentmark{15}{\char"1CD2}}न्द्रो\raisebox{0.25ex}{\accentmark{15}{\char"1CD2}}मू\raisebox{0.15ex}{\accentmark{12}{\char"1CF9}}नीनां\raisebox{0.25ex}{\accentmark{12}{\char"0951}}सा\raisebox{0.25ex}{\accentmark{15}{\char"1CD2}}खा\raisebox{0.25ex}{\accentmark{12}{\char"0951}} \nolinebreak\hskip 0.08em \mbox{॥३॥}}} \par}

\vspace{0.8em}
                % --- SUBSECTIONS (Mantras) ---
                
                % Conditional rendering based on output_mode
                \par\filbreak

\phantomsection

\addcontentsline{toc}{subsection}{ऋक् ४}

{\centering \footnotesize \textcolor{AccentPurple}{॥ जमदग्निभागर्वः इन्द्रः (सूर्यः) बृहती ॥} \par}

\vspace{0.12em}

{\centering \textbf{\textcolor{AccentBlue}{ब\raisebox{0.25ex}{\accentmark{15}{\char"1CD2}}ण्माहं\raisebox{0.25ex}{\accentmark{15}{\char"1CD2}}आ\raisebox{0.25ex}{\accentmark{12}{\char"0951}}सीसूर्या\kern 0.05em ।\hskip 0.12em plus 0.04em minus 0.02em बा\raisebox{0.25ex}{\accentmark{15}{\char"1CD2}}ळा\raisebox{0.25ex}{\accentmark{12}{\char"0951}}दित्यामाहंआसि\kern 0.05em ।\hskip 0.12em plus 0.04em minus 0.02em माह\raisebox{0.25ex}{\accentmark{15}{\char"1CD2}}स्ते\raisebox{0.25ex}{\accentmark{12}{\char"0951}}सातो\raisebox{0.15ex}{\accentmark{12}{\char"1CF9}}मा\raisebox{0.25ex}{\accentmark{12}{\char"0951}}हीमा\raisebox{0.25ex}{\accentmark{15}{\char"1CD2}}पा\raisebox{0.25ex}{\accentmark{12}{\char"0951}}निष्ठमा\kern 0.05em ।\hskip 0.12em plus 0.04em minus 0.02em मन्हा\raisebox{0.25ex}{\accentmark{15}{\char"1CD2}}दे\raisebox{0.25ex}{\accentmark{12}{\char"0951}}वामाहं\raisebox{0.25ex}{\accentmark{15}{\char"1CD2}}आसि \nolinebreak\hskip 0.08em \mbox{॥४॥}}} \par}

\vspace{0.8em}
                % --- SUBSECTIONS (Mantras) ---
                
                % Conditional rendering based on output_mode
                \par\filbreak

\phantomsection

\addcontentsline{toc}{subsection}{ऋक् ५}

{\centering \footnotesize \textcolor{AccentPurple}{॥ देवातिथिः काण्वः इन्द्रः बृहती ॥} \par}

\vspace{0.12em}

{\centering \textbf{\textcolor{AccentBlue}{या\raisebox{0.15ex}{\accentmark{12}{\char"1CF8}}दि\raisebox{0.25ex}{\accentmark{12}{\char"0951}}न्द्राप्रा\raisebox{0.25ex}{\accentmark{15}{\char"1CD2}}गपागूद\raisebox{0.25ex}{\accentmark{12}{\char"0951}}त्\kern 0.05em ।\hskip 0.12em plus 0.04em minus 0.02em न्या\raisebox{0.25ex}{\accentmark{12}{\char"0951}}ग्वाहूयासेनृभि\raisebox{0.25ex}{\accentmark{12}{\char"0951}}ः\kern 0.05em ।\hskip 0.12em plus 0.04em minus 0.02em सीमा\raisebox{0.25ex}{\accentmark{12}{\char"0951}}पूरूनृ\raisebox{0.25ex}{\accentmark{15}{\char"1CD2}}षू\raisebox{0.25ex}{\accentmark{12}{\char"0951}}तोअस्या\raisebox{0.25ex}{\accentmark{15}{\char"1CD2}}ना\raisebox{0.15ex}{\accentmark{12}{\char"1CF8}}वे\kern 0.05em ।\hskip 0.12em plus 0.04em minus 0.02em आ\raisebox{0.25ex}{\accentmark{15}{\char"1CD2}}सी\raisebox{0.25ex}{\accentmark{12}{\char"0951}}प्रशार्धातुर्वा\raisebox{0.25ex}{\accentmark{15}{\char"1CD2}}शे \nolinebreak\hskip 0.08em \mbox{॥५॥}}} \par}

\vspace{0.8em}
                % --- SUBSECTIONS (Mantras) ---
                
                % Conditional rendering based on output_mode
                \par\filbreak

\phantomsection

\addcontentsline{toc}{subsection}{ऋक् ६}

{\centering \footnotesize \textcolor{AccentPurple}{॥ वसिष्ठो मैत्रावरुणिः इन्द्रः बृहती ॥} \par}

\vspace{0.12em}

{\centering \textbf{\textcolor{AccentBlue}{क\raisebox{0.25ex}{\accentmark{15}{\char"1CD2}}स्तामि\raisebox{0.25ex}{\accentmark{12}{\char"0951}}न्द्रत्वावसो\kern 0.05em ।\hskip 0.12em plus 0.04em minus 0.02em आ\raisebox{0.25ex}{\accentmark{15}{\char"1CD2}}मा\raisebox{0.25ex}{\accentmark{15}{\char"1CD2}}र्क्त्यो\raisebox{0.25ex}{\accentmark{12}{\char"0951}}दधर्षते\kern 0.05em ।\hskip 0.12em plus 0.04em minus 0.02em श्र\raisebox{0.25ex}{\accentmark{15}{\char"1CD2}}द्धाही\raisebox{0.25ex}{\accentmark{15}{\char"1CD2}}तेमाघा\raisebox{0.25ex}{\accentmark{15}{\char"1CD2}}वन्पा\raisebox{0.15ex}{\accentmark{12}{\char"1CF9}}र्येदीवी\kern 0.05em ।\hskip 0.12em plus 0.04em minus 0.02em वाजी\raisebox{0.25ex}{\accentmark{15}{\char"1CD2}}वाजंसिषासति \nolinebreak\hskip 0.08em \mbox{॥६॥}}} \par}

\vspace{0.8em}
                % --- SUBSECTIONS (Mantras) ---
                
                % Conditional rendering based on output_mode
                \par\filbreak

\phantomsection

\addcontentsline{toc}{subsection}{ऋक् ७}

{\centering \footnotesize \textcolor{AccentPurple}{॥ देवातिथिः काण्वः इन्द्रः बृहती ॥} \par}

\vspace{0.12em}

{\centering \textbf{\textcolor{AccentBlue}{अश्वीराथीसुरूपायु\raisebox{0.25ex}{\accentmark{15}{\char"1CD2}}ः\kern 0.05em ।\hskip 0.12em plus 0.04em minus 0.02em गोम\raisebox{0.25ex}{\accentmark{15}{\char"1CD2}}य्यादिन्द्राते\raisebox{0.25ex}{\accentmark{15}{\char"1CD2}}साखा\kern 0.05em ।\hskip 0.12em plus 0.04em minus 0.02em श्वात्राभा\raisebox{0.25ex}{\accentmark{15}{\char"1CD2}}जावा\raisebox{0.25ex}{\accentmark{15}{\char"1CD2}}या\raisebox{0.25ex}{\accentmark{12}{\char"0951}}सासचाते\raisebox{0.25ex}{\accentmark{15}{\char"1CD2}}सादा\raisebox{0.25ex}{\accentmark{12}{\char"0951}}\kern 0.05em ।\hskip 0.12em plus 0.04em minus 0.02em चन्द्रै\raisebox{0.25ex}{\accentmark{15}{\char"1CD2}}र्यातीसाभा\raisebox{0.25ex}{\accentmark{15}{\char"1CD2}}ऊ\raisebox{0.25ex}{\accentmark{15}{\char"1CD2}}पा\raisebox{0.25ex}{\accentmark{12}{\char"0951}} \nolinebreak\hskip 0.08em \mbox{॥७॥}}} \par}

\vspace{0.8em}
                % --- SUBSECTIONS (Mantras) ---
                
                % Conditional rendering based on output_mode
                \par\filbreak

\phantomsection

\addcontentsline{toc}{subsection}{ऋक् ८}

{\centering \footnotesize \textcolor{AccentPurple}{॥ पुरुहन्मा आंगिरसः इन्द्रः बृहती ॥} \par}

\vspace{0.12em}

{\centering \textbf{\textcolor{AccentBlue}{य\raisebox{0.25ex}{\accentmark{15}{\char"1CD2}}द्यावा\raisebox{0.25ex}{\accentmark{15}{\char"1CD2}}इन्द्रतेशातम्\raisebox{0.25ex}{\accentmark{15}{\char"1CD2}}\kern 0.05em ।\hskip 0.12em plus 0.04em minus 0.02em शातं\raisebox{0.15ex}{\accentmark{12}{\char"1CF9}}भूमी\raisebox{0.25ex}{\accentmark{12}{\char"0951}}रूतास्युः\kern 0.05em ।\hskip 0.12em plus 0.04em minus 0.02em न\raisebox{0.25ex}{\accentmark{15}{\char"1CD2}}त्वावज्ज्रिन्साहा\raisebox{0.25ex}{\accentmark{15}{\char"1CD2}}स्रंसू\raisebox{0.25ex}{\accentmark{15}{\char"1CD2}}र्याआनू\raisebox{0.25ex}{\accentmark{12}{\char"0951}}\kern 0.05em ।\hskip 0.12em plus 0.04em minus 0.02em ना\raisebox{0.25ex}{\accentmark{15}{\char"1CD2}}जाता\raisebox{0.15ex}{\accentmark{12}{\char"1CF9}}माष्टारोद\raisebox{0.25ex}{\accentmark{12}{\char"0951}}सी \nolinebreak\hskip 0.08em \mbox{॥८॥}}} \par}

\vspace{0.8em}
                % --- SUBSECTIONS (Mantras) ---
                
                % Conditional rendering based on output_mode
                \par\filbreak

\phantomsection

\addcontentsline{toc}{subsection}{ऋक् ९}

{\centering \footnotesize \textcolor{AccentPurple}{॥ भरद्वाजो बार्हस्पत्यः इन्द्रः (इन्द्राग्नी) बृहती ॥} \par}

\vspace{0.12em}

{\centering \textbf{\textcolor{AccentBlue}{इन्द्रा\raisebox{0.25ex}{\accentmark{12}{\char"0951}}ग्नीआपा\raisebox{0.25ex}{\accentmark{15}{\char"1CD2}}दीय\raisebox{0.25ex}{\accentmark{15}{\char"1CD2}}म्\kern 0.05em ।\hskip 0.12em plus 0.04em minus 0.02em पूर्वागा\raisebox{0.15ex}{\accentmark{12}{\char"1CF8}}त्पद्वा\raisebox{0.25ex}{\accentmark{15}{\char"1CD2}}ती\raisebox{0.25ex}{\accentmark{12}{\char"0951}}भ्यः\kern 0.05em ।\hskip 0.12em plus 0.04em minus 0.02em हित्वा\raisebox{0.15ex}{\accentmark{12}{\char"1CF9}}शि\raisebox{0.15ex}{\accentmark{12}{\char"1CF9}}रो\raisebox{0.25ex}{\accentmark{12}{\char"0951}}जिह्वा\raisebox{0.25ex}{\accentmark{15}{\char"1CD2}}या\raisebox{0.15ex}{\accentmark{12}{\char"1CF9}}रा\raisebox{0.25ex}{\accentmark{12}{\char"0951}}रापच्चारात्\kern 0.05em ।\hskip 0.12em plus 0.04em minus 0.02em त्रिंशा\raisebox{0.25ex}{\accentmark{15}{\char"1CD2}}त्पादा\raisebox{0.25ex}{\accentmark{15}{\char"1CD2}}न्या\raisebox{0.25ex}{\accentmark{12}{\char"0951}}क्रमीत् \nolinebreak\hskip 0.08em \mbox{॥९॥}}} \par}

\vspace{0.8em}
                % --- SUBSECTIONS (Mantras) ---
                
                % Conditional rendering based on output_mode
                \par\filbreak

\phantomsection

\addcontentsline{toc}{subsection}{ऋक् १०}

{\centering \footnotesize \textcolor{AccentPurple}{॥ मेध्यः काण्वः इन्द्रः बृहती ॥} \par}

\vspace{0.12em}

{\centering \textbf{\textcolor{AccentBlue}{इ\raisebox{0.25ex}{\accentmark{15}{\char"1CD2}}न्द्राने\raisebox{0.15ex}{\accentmark{12}{\char"1CF9}}दी\raisebox{0.25ex}{\accentmark{15}{\char"1CD2}}याए\raisebox{0.25ex}{\accentmark{15}{\char"1CD2}}दिहि\kern 0.05em ।\hskip 0.12em plus 0.04em minus 0.02em मीता\raisebox{0.25ex}{\accentmark{15}{\char"1CD2}}मे\raisebox{0.25ex}{\accentmark{15}{\char"1CD2}}धाभीरू\raisebox{0.25ex}{\accentmark{15}{\char"1CD2}}ती\raisebox{0.25ex}{\accentmark{15}{\char"1CD2}}भी\raisebox{0.25ex}{\accentmark{12}{\char"0951}}ः\kern 0.05em ।\hskip 0.12em plus 0.04em minus 0.02em आ\raisebox{0.25ex}{\accentmark{15}{\char"1CD2}}श\raisebox{0.25ex}{\accentmark{12}{\char"0951}}न्तामाश\raisebox{0.25ex}{\accentmark{15}{\char"1CD2}}न्ता\raisebox{0.25ex}{\accentmark{12}{\char"0951}}माभी\raisebox{0.25ex}{\accentmark{15}{\char"1CD2}}राभीष्टीभिः\kern 0.05em ।\hskip 0.12em plus 0.04em minus 0.02em आस्वा\raisebox{0.25ex}{\accentmark{15}{\char"1CD2}}पेस्वापि\raisebox{0.25ex}{\accentmark{15}{\char"1CD2}}भि\raisebox{0.25ex}{\accentmark{12}{\char"0951}}ः \nolinebreak\hskip 0.08em \mbox{॥१०॥}}} \par}

\vspace{0.8em}

            % --- End of Kandah Closing Floral Line & Counts ---
            \par\nopagebreak\vspace{0.9em plus 0.3em minus 0.1em}
            \centerline{\textcolor{black!35}{\rule{1.8cm}{0.35pt}}\;\raisebox{-0.22ex}{\textcolor{vedicmaroon}{\large\SixFlowerPetalDotted}}\;\textcolor{black!35}{\rule{1.8cm}{0.35pt}}}
            \nopagebreak\vspace{0.45em}
            \centerline{\devafont\textbf{\textcolor{AccentBlue}{॥ ऋक् सङ्ख्या : १० ॥}}}
        
            \clearpage

            % --- SECTION ---
            
            \phantomsection  

            \section*{\centering षष्ठ खण्डः}
            \addcontentsline{toc}{section}{षष्ठ खण्डः} 
            \markright{षष्ठ खण्डः}
            \nopagebreak % Keep section title with first mantra
            
            % --- SECTION-LEVEL PROCEDURE LINK ---

                % --- SUBSECTIONS (Mantras) ---
                
                % Conditional rendering based on output_mode
                \par\filbreak

\phantomsection

\addcontentsline{toc}{subsection}{ऋक् १}

{\centering \footnotesize \textcolor{AccentPurple}{॥ नृमेध आंगिरसः इन्द्रः बृहती ॥} \par}

\vspace{0.12em}

{\centering \textbf{\textcolor{AccentBlue}{ई\raisebox{0.25ex}{\accentmark{15}{\char"1CD2}}ताऊती\raisebox{0.15ex}{\accentmark{12}{\char"1CF9}}वो\raisebox{0.25ex}{\accentmark{12}{\char"0951}}आज\raisebox{0.25ex}{\accentmark{15}{\char"1CD2}}रम्\raisebox{0.25ex}{\accentmark{12}{\char"0951}}\kern 0.05em ।\hskip 0.12em plus 0.04em minus 0.02em प्रे\raisebox{0.25ex}{\accentmark{15}{\char"1CD2}}हेता\raisebox{0.25ex}{\accentmark{15}{\char"1CD2}}रामा\raisebox{0.25ex}{\accentmark{15}{\char"1CD2}}प्रा\raisebox{0.25ex}{\accentmark{12}{\char"0951}}हितम्\kern 0.05em ।\hskip 0.12em plus 0.04em minus 0.02em आशू\raisebox{0.15ex}{\accentmark{12}{\char"1CF9}}ञ्जे\raisebox{0.15ex}{\accentmark{12}{\char"1CF9}}ता\raisebox{0.25ex}{\accentmark{12}{\char"0951}}रंहोता\raisebox{0.25ex}{\accentmark{12}{\char"0951}}रंराथी\raisebox{0.25ex}{\accentmark{15}{\char"1CD2}}तामम्\kern 0.05em ।\hskip 0.12em plus 0.04em minus 0.02em आ\raisebox{0.25ex}{\accentmark{15}{\char"1CD2}}तूर्तन्तूग्रीयावृ\raisebox{0.25ex}{\accentmark{15}{\char"1CD2}}धम् \nolinebreak\hskip 0.08em \mbox{॥१॥}}} \par}

\vspace{0.8em}
                % --- SUBSECTIONS (Mantras) ---
                
                % Conditional rendering based on output_mode
                \par\filbreak

\phantomsection

\addcontentsline{toc}{subsection}{ऋक् २}

{\centering \footnotesize \textcolor{AccentPurple}{॥ वसिष्ठो मैत्रावरुणिः इन्द्रः बृहती ॥} \par}

\vspace{0.12em}

{\centering \textbf{\textcolor{AccentBlue}{मोषु\raisebox{0.25ex}{\accentmark{15}{\char"1CD2}}त्वा\raisebox{0.25ex}{\accentmark{15}{\char"1CD2}}वाघा\raisebox{0.15ex}{\accentmark{12}{\char"1CF9}}ता\raisebox{0.25ex}{\accentmark{12}{\char"0951}}ःचाना\kern 0.05em ।\hskip 0.12em plus 0.04em minus 0.02em आ\raisebox{0.25ex}{\accentmark{15}{\char"1CD2}}रे\raisebox{0.25ex}{\accentmark{15}{\char"1CD2}}आस्मिन्नी\raisebox{0.25ex}{\accentmark{15}{\char"1CD2}}री\raisebox{0.25ex}{\accentmark{12}{\char"0951}}रमन्\kern 0.05em ।\hskip 0.12em plus 0.04em minus 0.02em आ\raisebox{0.25ex}{\accentmark{15}{\char"1CD2}}रा\raisebox{0.25ex}{\accentmark{15}{\char"1CD2}}त्ता\raisebox{0.25ex}{\accentmark{15}{\char"1CD2}}द्वासाधा\raisebox{0.25ex}{\accentmark{12}{\char"0951}}मा\raisebox{0.15ex}{\accentmark{12}{\char"1CF9}}द\raisebox{0.25ex}{\accentmark{12}{\char"0951}}न्नाआ\raisebox{0.25ex}{\accentmark{15}{\char"1CD2}}गहि\kern 0.05em ।\hskip 0.12em plus 0.04em minus 0.02em ईहा\raisebox{0.25ex}{\accentmark{15}{\char"1CD2}}वासन्नूपा\raisebox{0.25ex}{\accentmark{12}{\char"0951}}श्रुधि \nolinebreak\hskip 0.08em \mbox{॥२॥}}} \par}

\vspace{0.8em}
                % --- SUBSECTIONS (Mantras) ---
                
                % Conditional rendering based on output_mode
                \par\filbreak

\phantomsection

\addcontentsline{toc}{subsection}{ऋक् ३}

{\centering \footnotesize \textcolor{AccentPurple}{॥ वसिष्ठो मैत्रावरुणिः इन्द्रः बृहती ॥} \par}

\vspace{0.12em}

{\centering \textbf{\textcolor{AccentBlue}{सू\raisebox{0.25ex}{\accentmark{15}{\char"1CD2}}नोता\raisebox{0.25ex}{\accentmark{12}{\char"0951}}सो\raisebox{0.25ex}{\accentmark{15}{\char"1CD2}}मा\raisebox{0.25ex}{\accentmark{15}{\char"1CD2}}पाव्ने\kern 0.05em ।\hskip 0.12em plus 0.04em minus 0.02em सो\raisebox{0.25ex}{\accentmark{15}{\char"1CD2}}मामि\raisebox{0.25ex}{\accentmark{15}{\char"1CD2}}न्द्रा\raisebox{0.25ex}{\accentmark{12}{\char"0951}}यावज्रीणे\kern 0.05em ।\hskip 0.12em plus 0.04em minus 0.02em पा\raisebox{0.25ex}{\accentmark{15}{\char"1CD2}}चा\raisebox{0.25ex}{\accentmark{15}{\char"1CD2}}तापक्तिरावा\raisebox{0.25ex}{\accentmark{12}{\char"0951}}सेकृ\raisebox{0.25ex}{\accentmark{12}{\char"0951}}णुध्वमित्\kern 0.05em ।\hskip 0.12em plus 0.04em minus 0.02em पृ\raisebox{0.15ex}{\accentmark{12}{\char"1CF9}}णन्नि\raisebox{0.15ex}{\accentmark{12}{\char"1CF9}}त्पृणातेमाया\raisebox{0.25ex}{\accentmark{12}{\char"0951}}ः \nolinebreak\hskip 0.08em \mbox{॥३॥}}} \par}

\vspace{0.8em}
                % --- SUBSECTIONS (Mantras) ---
                
                % Conditional rendering based on output_mode
                \par\filbreak

\phantomsection

\addcontentsline{toc}{subsection}{ऋक् ४}

{\centering \footnotesize \textcolor{AccentPurple}{॥ भरद्वाजः (शयुः) बार्हस्पत्यः इन्द्रः बृहती ॥} \par}

\vspace{0.12em}

{\centering \textbf{\textcolor{AccentBlue}{य\raisebox{0.25ex}{\accentmark{15}{\char"1CD2}}स्सा\raisebox{0.25ex}{\accentmark{12}{\char"0951}}त्राहा\raisebox{0.25ex}{\accentmark{15}{\char"1CD2}}वीच\raisebox{0.25ex}{\accentmark{15}{\char"1CD2}}ऋषणिः\kern 0.05em ।\hskip 0.12em plus 0.04em minus 0.02em इ\raisebox{0.25ex}{\accentmark{15}{\char"1CD2}}न्द्रन्तंहू\raisebox{0.15ex}{\accentmark{12}{\char"1CF8}}माहेवायम्\raisebox{0.25ex}{\accentmark{15}{\char"1CD2}}\kern 0.05em ।\hskip 0.12em plus 0.04em minus 0.02em सा\raisebox{0.25ex}{\accentmark{15}{\char"1CD2}}हा\raisebox{0.25ex}{\accentmark{12}{\char"0951}}स्रमन्योतुविनिर्म\raisebox{0.25ex}{\accentmark{12}{\char"0951}}णसत्पते\kern 0.05em ।\hskip 0.12em plus 0.04em minus 0.02em भा\raisebox{0.15ex}{\accentmark{12}{\char"1CF8}}वा\raisebox{0.25ex}{\accentmark{12}{\char"0951}}सामात्सू\raisebox{0.25ex}{\accentmark{12}{\char"0951}}नोवृधे\raisebox{0.25ex}{\accentmark{15}{\char"1CD2}} \nolinebreak\hskip 0.08em \mbox{॥४॥}}} \par}

\vspace{0.8em}
                % --- SUBSECTIONS (Mantras) ---
                
                % Conditional rendering based on output_mode
                \par\filbreak

\phantomsection

\addcontentsline{toc}{subsection}{ऋक् ५}

{\centering \footnotesize \textcolor{AccentPurple}{॥ पुरुच्छेपो देवोदासिः इन्द्रः (शिवनौ) बृहती ॥} \par}

\vspace{0.12em}

{\centering \textbf{\textcolor{AccentBlue}{शा\raisebox{0.25ex}{\accentmark{15}{\char"1CD2}}चिभिर्नश्शची\raisebox{0.25ex}{\accentmark{15}{\char"1CD2}}वसू\kern 0.05em ।\hskip 0.12em plus 0.04em minus 0.02em दीवान\raisebox{0.25ex}{\accentmark{15}{\char"1CD2}}क्तंदिशस्यतम्\kern 0.05em ।\hskip 0.12em plus 0.04em minus 0.02em मा\raisebox{0.15ex}{\accentmark{12}{\char"1CF8}}वां\raisebox{0.25ex}{\accentmark{12}{\char"0951}}रातिरूपा\raisebox{0.25ex}{\accentmark{12}{\char"0951}}द\raisebox{0.25ex}{\accentmark{15}{\char"1CD2}}सात्कादा\raisebox{0.25ex}{\accentmark{15}{\char"1CD2}}चा\raisebox{0.25ex}{\accentmark{15}{\char"1CD2}}ना\kern 0.05em ।\hskip 0.12em plus 0.04em minus 0.02em अस्म\raisebox{0.25ex}{\accentmark{15}{\char"1CD2}}द्रा\raisebox{0.25ex}{\accentmark{15}{\char"1CD2}}तीःकादाचाना\raisebox{0.25ex}{\accentmark{15}{\char"1CD2}} \nolinebreak\hskip 0.08em \mbox{॥५॥}}} \par}

\vspace{0.8em}
                % --- SUBSECTIONS (Mantras) ---
                
                % Conditional rendering based on output_mode
                \par\filbreak

\phantomsection

\addcontentsline{toc}{subsection}{ऋक् ६}

{\centering \footnotesize \textcolor{AccentPurple}{॥ वामदेवो गौतमः इन्द्रः बृहती ॥} \par}

\vspace{0.12em}

{\centering \textbf{\textcolor{AccentBlue}{यादाकादा\raisebox{0.15ex}{\accentmark{12}{\char"1CF9}}चामीढूषे\raisebox{0.25ex}{\accentmark{12}{\char"0951}}\kern 0.05em ।\hskip 0.12em plus 0.04em minus 0.02em स्तोता\raisebox{0.25ex}{\accentmark{15}{\char"1CD2}}जरेतमार्क्त्याः\raisebox{0.25ex}{\accentmark{12}{\char"0951}}\kern 0.05em ।\hskip 0.12em plus 0.04em minus 0.02em आदीद्व\raisebox{0.25ex}{\accentmark{15}{\char"1CD2}}न्दे\raisebox{0.25ex}{\accentmark{15}{\char"1CD2}}तावारुणंवीपा\raisebox{0.25ex}{\accentmark{15}{\char"1CD2}}गीरा\raisebox{0.25ex}{\accentmark{12}{\char"0951}}\kern 0.05em ।\hskip 0.12em plus 0.04em minus 0.02em धक्ता\raisebox{0.25ex}{\accentmark{15}{\char"1CD2}}रंवि\raisebox{0.25ex}{\accentmark{15}{\char"1CD2}}प्रा\raisebox{0.25ex}{\accentmark{15}{\char"1CD2}}तानाम् \nolinebreak\hskip 0.08em \mbox{॥६॥}}} \par}

\vspace{0.8em}
                % --- SUBSECTIONS (Mantras) ---
                
                % Conditional rendering based on output_mode
                \par\filbreak

\phantomsection

\addcontentsline{toc}{subsection}{ऋक् ७}

{\centering \footnotesize \textcolor{AccentPurple}{॥ मध्येआतिथिः काण्वः इन्द्रः बृहती ॥} \par}

\vspace{0.12em}

{\centering \textbf{\textcolor{AccentBlue}{पाहि\raisebox{0.25ex}{\accentmark{15}{\char"1CD2}}गाया\raisebox{0.15ex}{\accentmark{12}{\char"1CF9}}न्धा\raisebox{0.25ex}{\accentmark{12}{\char"0951}}सामा\raisebox{0.25ex}{\accentmark{12}{\char"0951}}दे\kern 0.05em ।\hskip 0.12em plus 0.04em minus 0.02em इ\raisebox{0.25ex}{\accentmark{15}{\char"1CD2}}न्द्रा\raisebox{0.15ex}{\accentmark{12}{\char"1CF8}}यमेधियातिथे\kern 0.05em ।\hskip 0.12em plus 0.04em minus 0.02em य\raisebox{0.15ex}{\accentmark{12}{\char"1CF8}}स\raisebox{0.15ex}{\accentmark{12}{\char"1CF8}}म्मिश्लोहर्यो\raisebox{0.25ex}{\accentmark{12}{\char"0951}}र्यो\raisebox{0.25ex}{\accentmark{12}{\char"0951}}हीरण्या\raisebox{0.25ex}{\accentmark{15}{\char"1CD2}}या\kern 0.05em ।\hskip 0.12em plus 0.04em minus 0.02em इन्द्रोवज्री\raisebox{0.15ex}{\accentmark{12}{\char"1CF9}}हीरण्या\raisebox{0.25ex}{\accentmark{15}{\char"1CD2}}या \nolinebreak\hskip 0.08em \mbox{॥७॥}}} \par}

\vspace{0.8em}
                % --- SUBSECTIONS (Mantras) ---
                
                % Conditional rendering based on output_mode
                \par\filbreak

\phantomsection

\addcontentsline{toc}{subsection}{ऋक् ८}

{\centering \footnotesize \textcolor{AccentPurple}{॥ भगः प्रागाथः इन्द्रः बृहती ॥} \par}

\vspace{0.12em}

{\centering \textbf{\textcolor{AccentBlue}{ऊभा\raisebox{0.15ex}{\accentmark{12}{\char"1CF9}}यं\raisebox{0.25ex}{\accentmark{15}{\char"1CD2}}श्रृणा\raisebox{0.25ex}{\accentmark{15}{\char"1CD2}}वा\raisebox{0.25ex}{\accentmark{12}{\char"0951}}चनः\kern 0.05em ।\hskip 0.12em plus 0.04em minus 0.02em इ\raisebox{0.15ex}{\accentmark{12}{\char"1CF8}}न्द्रो\raisebox{0.25ex}{\accentmark{12}{\char"0951}}अर्वा\raisebox{0.25ex}{\accentmark{15}{\char"1CD2}}गीदं\raisebox{0.25ex}{\accentmark{15}{\char"1CD2}}वा\raisebox{0.25ex}{\accentmark{15}{\char"1CD2}}चाः\kern 0.05em ।\hskip 0.12em plus 0.04em minus 0.02em सत्रा\raisebox{0.15ex}{\accentmark{12}{\char"1CF9}}च्या\raisebox{0.15ex}{\accentmark{12}{\char"1CF9}}मा\raisebox{0.25ex}{\accentmark{12}{\char"0951}}घावा\raisebox{0.25ex}{\accentmark{15}{\char"1CD2}}सो\raisebox{0.25ex}{\accentmark{15}{\char"1CD2}}मा\raisebox{0.25ex}{\accentmark{12}{\char"0951}}पीतये\kern 0.05em ।\hskip 0.12em plus 0.04em minus 0.02em धीया\raisebox{0.15ex}{\accentmark{12}{\char"1CF9}}श्शावी\raisebox{0.25ex}{\accentmark{12}{\char"0951}}ष्ठआग\raisebox{0.25ex}{\accentmark{12}{\char"0951}}मात् \nolinebreak\hskip 0.08em \mbox{॥८॥}}} \par}

\vspace{0.8em}
                % --- SUBSECTIONS (Mantras) ---
                
                % Conditional rendering based on output_mode
                \par\filbreak

\phantomsection

\addcontentsline{toc}{subsection}{ऋक् ९}

{\centering \footnotesize \textcolor{AccentPurple}{॥ मेधातिथि –मध्येआतिथी काण्वौ इन्द्रः बृहती ॥} \par}

\vspace{0.12em}

{\centering \textbf{\textcolor{AccentBlue}{माहे\raisebox{0.25ex}{\accentmark{15}{\char"1CD2}}चा\raisebox{0.25ex}{\accentmark{15}{\char"1CD2}}ना\raisebox{0.25ex}{\accentmark{15}{\char"1CD2}}त्वा\raisebox{0.25ex}{\accentmark{12}{\char"0951}}आद्रीवः\kern 0.05em ।\hskip 0.12em plus 0.04em minus 0.02em पा\raisebox{0.15ex}{\accentmark{12}{\char"1CF9}}रा\raisebox{0.25ex}{\accentmark{12}{\char"0951}}शुल्काया\raisebox{0.25ex}{\accentmark{12}{\char"0951}}दीयसे\kern 0.05em ।\hskip 0.12em plus 0.04em minus 0.02em ना\raisebox{0.25ex}{\accentmark{15}{\char"1CD2}}साह\raisebox{0.15ex}{\accentmark{12}{\char"1CF9}}स्रा\raisebox{0.25ex}{\accentmark{12}{\char"0951}}यानायू\raisebox{0.25ex}{\accentmark{15}{\char"1CD2}}ता\raisebox{0.25ex}{\accentmark{12}{\char"0951}}वज्रिवः\kern 0.05em ।\hskip 0.12em plus 0.04em minus 0.02em ना\raisebox{0.25ex}{\accentmark{15}{\char"1CD2}}शा\raisebox{0.25ex}{\accentmark{15}{\char"1CD2}}ता\raisebox{0.25ex}{\accentmark{12}{\char"0951}}याशतामघ \nolinebreak\hskip 0.08em \mbox{॥९॥}}} \par}

\vspace{0.8em}
                % --- SUBSECTIONS (Mantras) ---
                
                % Conditional rendering based on output_mode
                \par\filbreak

\phantomsection

\addcontentsline{toc}{subsection}{ऋक् १०}

{\centering \footnotesize \textcolor{AccentPurple}{॥ मेधातिथि –मध्येआतिथी काण्वौ इन्द्रः बृहती ॥} \par}

\vspace{0.12em}

{\centering \textbf{\textcolor{AccentBlue}{व\raisebox{0.25ex}{\accentmark{15}{\char"1CD2}}स्यठ\raisebox{0.25ex}{\accentmark{12}{\char"0951}}इन्द्रासिमेपीतू\raisebox{0.25ex}{\accentmark{15}{\char"1CD2}}ः\kern 0.05em ।\hskip 0.12em plus 0.04em minus 0.02em ऊताभ्रा\raisebox{0.25ex}{\accentmark{15}{\char"1CD2}}तूरा\raisebox{0.25ex}{\accentmark{15}{\char"1CD2}}भु\raisebox{0.25ex}{\accentmark{12}{\char"0951}}ञ्जतः\kern 0.05em ।\hskip 0.12em plus 0.04em minus 0.02em मा\raisebox{0.25ex}{\accentmark{15}{\char"1CD2}}ता\raisebox{0.25ex}{\accentmark{15}{\char"1CD2}}चा\raisebox{0.25ex}{\accentmark{12}{\char"0951}}मेछदयथास्सामा\raisebox{0.25ex}{\accentmark{15}{\char"1CD2}}वासो\kern 0.05em ।\hskip 0.12em plus 0.04em minus 0.02em व\raisebox{0.25ex}{\accentmark{15}{\char"1CD2}}सुत्वाना\raisebox{0.25ex}{\accentmark{15}{\char"1CD2}}यारा\raisebox{0.25ex}{\accentmark{15}{\char"1CD2}}धा\raisebox{0.25ex}{\accentmark{12}{\char"0951}}से \nolinebreak\hskip 0.08em \mbox{॥१०॥}}} \par}

\vspace{0.8em}

            % --- End of Kandah Closing Floral Line & Counts ---
            \par\nopagebreak\vspace{0.9em plus 0.3em minus 0.1em}
            \centerline{\textcolor{black!35}{\rule{1.8cm}{0.35pt}}\;\raisebox{-0.22ex}{\textcolor{vedicmaroon}{\large\SixFlowerPetalDotted}}\;\textcolor{black!35}{\rule{1.8cm}{0.35pt}}}
            \nopagebreak\vspace{0.45em}
            \centerline{\devafont\textbf{\textcolor{AccentBlue}{॥ ऋक् सङ्ख्या : १० ॥}}}
        
            \clearpage

            % --- SECTION ---
            
            \phantomsection  

            \section*{\centering सप्तम खण्डः}
            \addcontentsline{toc}{section}{सप्तम खण्डः} 
            \markright{सप्तम खण्डः}
            \nopagebreak % Keep section title with first mantra
            
            % --- SECTION-LEVEL PROCEDURE LINK ---

                % --- SUBSECTIONS (Mantras) ---
                
                % Conditional rendering based on output_mode
                \par\filbreak

\phantomsection

\addcontentsline{toc}{subsection}{ऋक् १}

{\centering \footnotesize \textcolor{AccentPurple}{॥ वसिष्ठो मैत्रावरुणिः इन्द्रः बृहती ॥} \par}

\vspace{0.12em}

{\centering \textbf{\textcolor{AccentBlue}{इमइ\raisebox{0.25ex}{\accentmark{15}{\char"1CD2}}न्द्रायसुन्विरे\kern 0.05em ।\hskip 0.12em plus 0.04em minus 0.02em सो\raisebox{0.15ex}{\accentmark{12}{\char"1CF8}}मा\raisebox{0.25ex}{\accentmark{12}{\char"0951}}सोदध्या\raisebox{0.25ex}{\accentmark{15}{\char"1CD2}}शिरः\kern 0.05em ।\hskip 0.12em plus 0.04em minus 0.02em तं\raisebox{0.25ex}{\accentmark{15}{\char"1CD2}}आ\raisebox{0.25ex}{\accentmark{15}{\char"1CD2}}मा\raisebox{0.25ex}{\accentmark{15}{\char"1CD2}}दा\raisebox{0.25ex}{\accentmark{12}{\char"0951}}यवज्रहास्तापीता\raisebox{0.25ex}{\accentmark{15}{\char"1CD2}}ये\kern 0.05em ।\hskip 0.12em plus 0.04em minus 0.02em हा\raisebox{0.25ex}{\accentmark{15}{\char"1CD2}}रीभ्याय्याह्योकाआ \nolinebreak\hskip 0.08em \mbox{॥१॥}}} \par}

\vspace{0.8em}
                % --- SUBSECTIONS (Mantras) ---
                
                % Conditional rendering based on output_mode
                \par\filbreak

\phantomsection

\addcontentsline{toc}{subsection}{ऋक् २}

{\centering \footnotesize \textcolor{AccentPurple}{॥ वामदेवो गौतमः इन्द्रः बृहती ॥} \par}

\vspace{0.12em}

{\centering \textbf{\textcolor{AccentBlue}{ईमा\raisebox{0.15ex}{\accentmark{12}{\char"1CF9}}इ\raisebox{0.25ex}{\accentmark{12}{\char"0951}}न्द्रामा\raisebox{0.25ex}{\accentmark{15}{\char"1CD2}}दा\raisebox{0.25ex}{\accentmark{12}{\char"0951}}यते\kern 0.05em ।\hskip 0.12em plus 0.04em minus 0.02em सो\raisebox{0.25ex}{\accentmark{15}{\char"1CD2}}मा\raisebox{0.25ex}{\accentmark{12}{\char"0951}}श्चिकीत्राउक्थी\raisebox{0.25ex}{\accentmark{15}{\char"1CD2}}ना\raisebox{0.25ex}{\accentmark{12}{\char"0951}}ः\kern 0.05em ।\hskip 0.12em plus 0.04em minus 0.02em माधोपावाना\raisebox{0.15ex}{\accentmark{12}{\char"1CF9}}ऊ\raisebox{0.15ex}{\accentmark{12}{\char"1CF9}}पा\raisebox{0.25ex}{\accentmark{12}{\char"0951}}नोगी\raisebox{0.25ex}{\accentmark{15}{\char"1CD2}}रा\raisebox{0.25ex}{\accentmark{12}{\char"0951}}ःश्रुणु\kern 0.05em ।\hskip 0.12em plus 0.04em minus 0.02em रा\raisebox{0.15ex}{\accentmark{12}{\char"1CF8}}स्वा\raisebox{0.25ex}{\accentmark{12}{\char"0951}}स्तोत्रा\raisebox{0.25ex}{\accentmark{15}{\char"1CD2}}या\raisebox{0.25ex}{\accentmark{12}{\char"0951}}गिर्वणः \nolinebreak\hskip 0.08em \mbox{॥२॥}}} \par}

\vspace{0.8em}
                % --- SUBSECTIONS (Mantras) ---
                
                % Conditional rendering based on output_mode
                \par\filbreak

\phantomsection

\addcontentsline{toc}{subsection}{ऋक् ३}

{\centering \footnotesize \textcolor{AccentPurple}{॥ मेधातिथि{-}मध्येआतिथी काण्वौ, विश्वामित्र इत्यकेे इन्द्रः बृहती ॥} \par}

\vspace{0.12em}

{\centering \textbf{\textcolor{AccentBlue}{आत्वाटट्यद्या\raisebox{0.15ex}{\accentmark{12}{\char"1CF9}}सा\raisebox{0.25ex}{\accentmark{12}{\char"0951}}बर्दूघाम्\raisebox{0.25ex}{\accentmark{12}{\char"0951}}\kern 0.05em ।\hskip 0.12em plus 0.04em minus 0.02em हूवे\raisebox{0.15ex}{\accentmark{12}{\char"1CF9}}गा\raisebox{0.25ex}{\accentmark{12}{\char"0951}}यत्रावे\raisebox{0.25ex}{\accentmark{12}{\char"0951}}पसम्\kern 0.05em ।\hskip 0.12em plus 0.04em minus 0.02em इ\raisebox{0.15ex}{\accentmark{12}{\char"1CF8}}न्द्रं\raisebox{0.25ex}{\accentmark{12}{\char"0951}}धेनुं\raisebox{0.25ex}{\accentmark{15}{\char"1CD2}}सूदू\raisebox{0.25ex}{\accentmark{15}{\char"1CD2}}घा\raisebox{0.25ex}{\accentmark{12}{\char"0951}}म\raisebox{0.25ex}{\accentmark{12}{\char"0951}}न्यामीष\raisebox{0.25ex}{\accentmark{12}{\char"0951}}म्\kern 0.05em ।\hskip 0.12em plus 0.04em minus 0.02em ऊरू\raisebox{0.25ex}{\accentmark{15}{\char"1CD2}}धारामारंकृ\raisebox{0.25ex}{\accentmark{12}{\char"0951}}तम् \nolinebreak\hskip 0.08em \mbox{॥३॥}}} \par}

\vspace{0.8em}
                % --- SUBSECTIONS (Mantras) ---
                
                % Conditional rendering based on output_mode
                \par\filbreak

\phantomsection

\addcontentsline{toc}{subsection}{ऋक् ४}

{\centering \footnotesize \textcolor{AccentPurple}{॥ नोधा गौतमः इन्द्रः बृहती ॥} \par}

\vspace{0.12em}

{\centering \textbf{\textcolor{AccentBlue}{नं\raisebox{0.15ex}{\accentmark{12}{\char"1CF8}}त्वा\raisebox{0.25ex}{\accentmark{12}{\char"0951}}बृह\raisebox{0.25ex}{\accentmark{15}{\char"1CD2}}न्तोअ\raisebox{0.25ex}{\accentmark{15}{\char"1CD2}}द्रा\raisebox{0.25ex}{\accentmark{12}{\char"0951}}यः\kern 0.05em ।\hskip 0.12em plus 0.04em minus 0.02em वा\raisebox{0.25ex}{\accentmark{15}{\char"1CD2}}रन्त\raisebox{0.25ex}{\accentmark{12}{\char"0951}}इन्द्रावीळावा\raisebox{0.25ex}{\accentmark{12}{\char"0951}}ः\kern 0.05em ।\hskip 0.12em plus 0.04em minus 0.02em य\raisebox{0.25ex}{\accentmark{15}{\char"1CD2}}च्छि\raisebox{0.25ex}{\accentmark{15}{\char"1CD2}}क्षा\raisebox{0.25ex}{\accentmark{12}{\char"0951}}सिस्तूवाते\raisebox{0.25ex}{\accentmark{15}{\char"1CD2}}मावातेवा\raisebox{0.25ex}{\accentmark{15}{\char"1CD2}}सू\raisebox{0.25ex}{\accentmark{12}{\char"0951}}\kern 0.05em ।\hskip 0.12em plus 0.04em minus 0.02em नाकीष्ट\raisebox{0.25ex}{\accentmark{15}{\char"1CD2}}दा\raisebox{0.25ex}{\accentmark{15}{\char"1CD2}}मी\raisebox{0.25ex}{\accentmark{12}{\char"0951}}नातिते \nolinebreak\hskip 0.08em \mbox{॥४॥}}} \par}

\vspace{0.8em}
                % --- SUBSECTIONS (Mantras) ---
                
                % Conditional rendering based on output_mode
                \par\filbreak

\phantomsection

\addcontentsline{toc}{subsection}{ऋक् ५}

{\centering \footnotesize \textcolor{AccentPurple}{॥ मेधातिथिः (मध्येआतिथिः) काण्वः इन्द्रः बृहती ॥} \par}

\vspace{0.12em}

{\centering \textbf{\textcolor{AccentBlue}{का\raisebox{0.25ex}{\accentmark{15}{\char"1CD2}}यीं\raisebox{0.25ex}{\accentmark{12}{\char"0951}}वे\raisebox{0.25ex}{\accentmark{15}{\char"1CD2}}दस्सूते\raisebox{0.25ex}{\accentmark{15}{\char"1CD2}}सा\raisebox{0.25ex}{\accentmark{15}{\char"1CD2}}चा\raisebox{0.25ex}{\accentmark{15}{\char"1CD2}}\kern 0.05em ।\hskip 0.12em plus 0.04em minus 0.02em पी\raisebox{0.15ex}{\accentmark{12}{\char"1CF8}}बन्तःकद्वायो\raisebox{0.25ex}{\accentmark{12}{\char"0951}}दधुः\kern 0.05em ।\hskip 0.12em plus 0.04em minus 0.02em आयं\raisebox{0.25ex}{\accentmark{15}{\char"1CD2}}यःपू\raisebox{0.25ex}{\accentmark{12}{\char"0951}}रो\raisebox{0.25ex}{\accentmark{12}{\char"0951}}वीभीनत्यो\raisebox{0.25ex}{\accentmark{15}{\char"1CD2}}ज\raisebox{0.25ex}{\accentmark{12}{\char"0951}}सा\kern 0.05em ।\hskip 0.12em plus 0.04em minus 0.02em मन्दा\raisebox{0.25ex}{\accentmark{15}{\char"1CD2}}ना\raisebox{0.25ex}{\accentmark{15}{\char"1CD2}}ःशिप्री\raisebox{0.25ex}{\accentmark{15}{\char"1CD2}}य\raisebox{0.25ex}{\accentmark{15}{\char"1CD2}}न्धा\raisebox{0.25ex}{\accentmark{12}{\char"0951}}सा \nolinebreak\hskip 0.08em \mbox{॥५॥}}} \par}

\vspace{0.8em}
                % --- SUBSECTIONS (Mantras) ---
                
                % Conditional rendering based on output_mode
                \par\filbreak

\phantomsection

\addcontentsline{toc}{subsection}{ऋक् ६}

{\centering \footnotesize \textcolor{AccentPurple}{॥ वामदेवो गौतमः इन्द्रः बृहती ॥} \par}

\vspace{0.12em}

{\centering \textbf{\textcolor{AccentBlue}{यादिन्द्राशासो\raisebox{0.25ex}{\accentmark{12}{\char"0951}}आव्रातम्\raisebox{0.25ex}{\accentmark{15}{\char"1CD2}}\kern 0.05em ।\hskip 0.12em plus 0.04em minus 0.02em च्यावा\raisebox{0.25ex}{\accentmark{15}{\char"1CD2}}यासा\raisebox{0.15ex}{\accentmark{12}{\char"1CF9}}द\raisebox{0.25ex}{\accentmark{12}{\char"0951}}सस्पारी\raisebox{0.25ex}{\accentmark{12}{\char"0951}}\kern 0.05em ।\hskip 0.12em plus 0.04em minus 0.02em अस्मा\raisebox{0.15ex}{\accentmark{12}{\char"1CF9}}का\raisebox{0.25ex}{\accentmark{12}{\char"0951}}मंशुर्म्मा\raisebox{0.25ex}{\accentmark{12}{\char"0951}}धावन्पूरूस्पृहम्\raisebox{0.25ex}{\accentmark{15}{\char"1CD2}}\kern 0.05em ।\hskip 0.12em plus 0.04em minus 0.02em वसव्येटट्यआ\raisebox{0.25ex}{\accentmark{15}{\char"1CD2}}धी\raisebox{0.25ex}{\accentmark{12}{\char"0951}}बर्हय \nolinebreak\hskip 0.08em \mbox{॥६॥}}} \par}

\vspace{0.8em}
                % --- SUBSECTIONS (Mantras) ---
                
                % Conditional rendering based on output_mode
                \par\filbreak

\phantomsection

\addcontentsline{toc}{subsection}{ऋक् ७}

{\centering \footnotesize \textcolor{AccentPurple}{॥ वामदेवो गौतमः बहुः बृहती ॥} \par}

\vspace{0.12em}

{\centering \textbf{\textcolor{AccentBlue}{त्व\raisebox{0.15ex}{\accentmark{12}{\char"1CF8}}ष्टानोदै\raisebox{0.25ex}{\accentmark{15}{\char"1CD2}}व्यंवा\raisebox{0.25ex}{\accentmark{15}{\char"1CD2}}चा\raisebox{0.25ex}{\accentmark{12}{\char"0951}}ः\kern 0.05em ।\hskip 0.12em plus 0.04em minus 0.02em पर्ज\raisebox{0.25ex}{\accentmark{15}{\char"1CD2}}न्योब्र\raisebox{0.15ex}{\accentmark{12}{\char"1CF9}}ह्मा\raisebox{0.25ex}{\accentmark{12}{\char"0951}}ण\raisebox{0.25ex}{\accentmark{15}{\char"1CD2}}स्पा\raisebox{0.25ex}{\accentmark{15}{\char"1CD2}}तीः\kern 0.05em ।\hskip 0.12em plus 0.04em minus 0.02em पुत्रै\raisebox{0.15ex}{\accentmark{12}{\char"1CF9}}भ्रा\raisebox{0.15ex}{\accentmark{12}{\char"1CF9}}तृ\raisebox{0.25ex}{\accentmark{12}{\char"0951}}भिरा\raisebox{0.15ex}{\accentmark{12}{\char"1CF9}}दीतीर्नूपातु\raisebox{0.25ex}{\accentmark{12}{\char"0951}}नः\kern 0.05em ।\hskip 0.12em plus 0.04em minus 0.02em दुष्टा\raisebox{0.25ex}{\accentmark{12}{\char"0951}}रन्द्रा\raisebox{0.15ex}{\accentmark{12}{\char"1CF9}}मा\raisebox{0.25ex}{\accentmark{12}{\char"0951}}णंवा\raisebox{0.25ex}{\accentmark{15}{\char"1CD2}}चाः \nolinebreak\hskip 0.08em \mbox{॥७॥}}} \par}

\vspace{0.8em}
                % --- SUBSECTIONS (Mantras) ---
                
                % Conditional rendering based on output_mode
                \par\filbreak

\phantomsection

\addcontentsline{toc}{subsection}{ऋक् ८}

{\centering \footnotesize \textcolor{AccentPurple}{॥ श्रुष्टिगुः काण्वः इन्द्रः बृहती ॥} \par}

\vspace{0.12em}

{\centering \textbf{\textcolor{AccentBlue}{कादा\raisebox{0.25ex}{\accentmark{15}{\char"1CD2}}चाना\raisebox{0.25ex}{\accentmark{15}{\char"1CD2}}स्तारी\raisebox{0.25ex}{\accentmark{15}{\char"1CD2}}रा\raisebox{0.25ex}{\accentmark{12}{\char"0951}}सी\kern 0.05em ।\hskip 0.12em plus 0.04em minus 0.02em नेद्रं\raisebox{0.25ex}{\accentmark{12}{\char"0951}}सश्वसीदाशू\raisebox{0.25ex}{\accentmark{15}{\char"1CD2}}षे\kern 0.05em ।\hskip 0.12em plus 0.04em minus 0.02em ऊ\raisebox{0.15ex}{\accentmark{12}{\char"1CF8}}पो\raisebox{0.25ex}{\accentmark{12}{\char"0951}}पेन्नूमाघावन्भू\raisebox{0.25ex}{\accentmark{15}{\char"1CD2}}याइ\raisebox{0.25ex}{\accentmark{15}{\char"1CD2}}न्नूते\kern 0.05em ।\hskip 0.12em plus 0.04em minus 0.02em दा\raisebox{0.15ex}{\accentmark{12}{\char"1CF8}}नन्देवा\raisebox{0.25ex}{\accentmark{15}{\char"1CD2}}स्या\raisebox{0.25ex}{\accentmark{12}{\char"0951}}प्रच्यते \nolinebreak\hskip 0.08em \mbox{॥८॥}}} \par}

\vspace{0.8em}
                % --- SUBSECTIONS (Mantras) ---
                
                % Conditional rendering based on output_mode
                \par\filbreak

\phantomsection

\addcontentsline{toc}{subsection}{ऋक् ९}

{\centering \footnotesize \textcolor{AccentPurple}{॥ मध्येआतिथिः (मेधातिथिर्वा) काण्वः इन्द्रः बृहती ॥} \par}

\vspace{0.12em}

{\centering \textbf{\textcolor{AccentBlue}{युंक्ष्वा\raisebox{0.25ex}{\accentmark{15}{\char"1CD2}}हीवृ\raisebox{0.25ex}{\accentmark{12}{\char"0951}}त्रहन्तमः\kern 0.05em ।\hskip 0.12em plus 0.04em minus 0.02em हारी\raisebox{0.25ex}{\accentmark{12}{\char"0951}}इन्द्रापारावा\raisebox{0.25ex}{\accentmark{15}{\char"1CD2}}ता\raisebox{0.25ex}{\accentmark{12}{\char"0951}}ः\kern 0.05em ।\hskip 0.12em plus 0.04em minus 0.02em अर्वाची\raisebox{0.25ex}{\accentmark{15}{\char"1CD2}}नो\raisebox{0.25ex}{\accentmark{15}{\char"1CD2}}मा\raisebox{0.25ex}{\accentmark{12}{\char"0951}}घावन्सो\raisebox{0.25ex}{\accentmark{15}{\char"1CD2}}मा\raisebox{0.25ex}{\accentmark{12}{\char"0951}}पीतये\kern 0.05em ।\hskip 0.12em plus 0.04em minus 0.02em उग्रा\raisebox{0.15ex}{\accentmark{12}{\char"1CF9}}रि\raisebox{0.15ex}{\accentmark{12}{\char"1CF9}}ष्वाभीरा\raisebox{0.25ex}{\accentmark{15}{\char"1CD2}}ग\raisebox{0.25ex}{\accentmark{12}{\char"0951}}हि \nolinebreak\hskip 0.08em \mbox{॥९॥}}} \par}

\vspace{0.8em}
                % --- SUBSECTIONS (Mantras) ---
                
                % Conditional rendering based on output_mode
                \par\filbreak

\phantomsection

\addcontentsline{toc}{subsection}{ऋक् १०}

{\centering \footnotesize \textcolor{AccentPurple}{॥ नृमेध आंगिरसः इन्द्रः बृहती ॥} \par}

\vspace{0.12em}

{\centering \textbf{\textcolor{AccentBlue}{त्वा\raisebox{0.25ex}{\accentmark{15}{\char"1CD2}}मीदा\raisebox{0.25ex}{\accentmark{15}{\char"1CD2}}ह्योना\raisebox{0.25ex}{\accentmark{15}{\char"1CD2}}राः\kern 0.05em ।\hskip 0.12em plus 0.04em minus 0.02em आ\raisebox{0.25ex}{\accentmark{15}{\char"1CD2}}पीप्यन्वज्ज्रिन्भू\raisebox{0.25ex}{\accentmark{12}{\char"0951}}र्णा\raisebox{0.25ex}{\accentmark{12}{\char"0951}}यः\kern 0.05em ।\hskip 0.12em plus 0.04em minus 0.02em सा\raisebox{0.15ex}{\accentmark{12}{\char"1CF8}}इ\raisebox{0.25ex}{\accentmark{12}{\char"0951}}न्द्रस्तोमा\raisebox{0.25ex}{\accentmark{12}{\char"0951}}वाहसाई\raisebox{0.25ex}{\accentmark{15}{\char"1CD2}}हा\raisebox{0.25ex}{\accentmark{15}{\char"1CD2}}श्रू\raisebox{0.25ex}{\accentmark{12}{\char"0951}}धिः\kern 0.05em ।\hskip 0.12em plus 0.04em minus 0.02em ऊपा\raisebox{0.15ex}{\accentmark{12}{\char"1CF9}}स्वा\raisebox{0.15ex}{\accentmark{12}{\char"1CF9}}सा\raisebox{0.25ex}{\accentmark{12}{\char"0951}}रामा\raisebox{0.25ex}{\accentmark{15}{\char"1CD2}}ग\raisebox{0.25ex}{\accentmark{12}{\char"0951}}हि \nolinebreak\hskip 0.08em \mbox{॥१०॥}}} \par}

\vspace{0.8em}

            % --- End of Kandah Closing Floral Line & Counts ---
            \par\nopagebreak\vspace{0.9em plus 0.3em minus 0.1em}
            \centerline{\textcolor{black!35}{\rule{1.8cm}{0.35pt}}\;\raisebox{-0.22ex}{\textcolor{vedicmaroon}{\large\SixFlowerPetalDotted}}\;\textcolor{black!35}{\rule{1.8cm}{0.35pt}}}
            \nopagebreak\vspace{0.45em}
            \centerline{\devafont\textbf{\textcolor{AccentBlue}{॥ ऋक् सङ्ख्या : १० ॥}}}
        
            \clearpage

            % --- SECTION ---
            
            \phantomsection  

            \section*{\centering अष्टम खण्डः}
            \addcontentsline{toc}{section}{अष्टम खण्डः} 
            \markright{अष्टम खण्डः}
            \nopagebreak % Keep section title with first mantra
            
            % --- SECTION-LEVEL PROCEDURE LINK ---

                % --- SUBSECTIONS (Mantras) ---
                
                % Conditional rendering based on output_mode
                \par\filbreak

\phantomsection

\addcontentsline{toc}{subsection}{ऋक् १}

{\centering \footnotesize \textcolor{AccentPurple}{॥ वसिष्ठो मैत्रावरुणिः उषा बृहती ॥} \par}

\vspace{0.12em}

{\centering \textbf{\textcolor{AccentBlue}{प्रा\raisebox{0.25ex}{\accentmark{15}{\char"1CD2}}ती\raisebox{0.25ex}{\accentmark{12}{\char"0951}}वोदृश्याया\raisebox{0.25ex}{\accentmark{15}{\char"1CD2}}ती\kern 0.05em ।\hskip 0.12em plus 0.04em minus 0.02em उच्छन्ती\raisebox{0.25ex}{\accentmark{12}{\char"0951}}दूहीता\raisebox{0.25ex}{\accentmark{15}{\char"1CD2}}दीवः\kern 0.05em ।\hskip 0.12em plus 0.04em minus 0.02em आ\raisebox{0.15ex}{\accentmark{12}{\char"1CF8}}पो\raisebox{0.25ex}{\accentmark{12}{\char"0951}}माही\raisebox{0.25ex}{\accentmark{12}{\char"0951}}वृणूतेचा\raisebox{0.15ex}{\accentmark{12}{\char"1CF9}}क्षू\raisebox{0.25ex}{\accentmark{12}{\char"0951}}षाता\raisebox{0.25ex}{\accentmark{15}{\char"1CD2}}माः\kern 0.05em ।\hskip 0.12em plus 0.04em minus 0.02em ज्यो\raisebox{0.25ex}{\accentmark{15}{\char"1CD2}}तीःकृ\raisebox{0.25ex}{\accentmark{12}{\char"0951}}णोतीसूना\raisebox{0.25ex}{\accentmark{15}{\char"1CD2}}री \nolinebreak\hskip 0.08em \mbox{॥१॥}}} \par}

\vspace{0.8em}
                % --- SUBSECTIONS (Mantras) ---
                
                % Conditional rendering based on output_mode
                \par\filbreak

\phantomsection

\addcontentsline{toc}{subsection}{ऋक् २}

{\centering \footnotesize \textcolor{AccentPurple}{॥ वसिष्ठो मैत्रावरुणिः अश्विनौ बृहती ॥} \par}

\vspace{0.12em}

{\centering \textbf{\textcolor{AccentBlue}{इमा\raisebox{0.15ex}{\accentmark{12}{\char"1CF9}}ऊ\raisebox{0.25ex}{\accentmark{12}{\char"0951}}वा\raisebox{0.25ex}{\accentmark{15}{\char"1CD2}}न्दिविष्टयः\kern 0.05em ।\hskip 0.12em plus 0.04em minus 0.02em उस्रा\raisebox{0.25ex}{\accentmark{15}{\char"1CD2}}हा\raisebox{0.25ex}{\accentmark{15}{\char"1CD2}}वन्तेअश्विना\kern 0.05em ।\hskip 0.12em plus 0.04em minus 0.02em आ\raisebox{0.15ex}{\accentmark{12}{\char"1CF8}}यं\raisebox{0.25ex}{\accentmark{15}{\char"1CD2}}वा\raisebox{0.25ex}{\accentmark{12}{\char"0951}}मंह्वाया\raisebox{0.25ex}{\accentmark{15}{\char"1CD2}}वा\raisebox{0.25ex}{\accentmark{12}{\char"0951}}सेशचीवसू\kern 0.05em ।\hskip 0.12em plus 0.04em minus 0.02em वीशं\raisebox{0.25ex}{\accentmark{12}{\char"0951}}वीशंही\raisebox{0.25ex}{\accentmark{15}{\char"1CD2}}ग\raisebox{0.25ex}{\accentmark{15}{\char"1CD2}}च्छा\raisebox{0.25ex}{\accentmark{12}{\char"0951}}थः \nolinebreak\hskip 0.08em \mbox{॥२॥}}} \par}

\vspace{0.8em}
                % --- SUBSECTIONS (Mantras) ---
                
                % Conditional rendering based on output_mode
                \par\filbreak

\phantomsection

\addcontentsline{toc}{subsection}{ऋक् ३}

{\centering \footnotesize \textcolor{AccentPurple}{॥ अश्विनौ वैवस्वतौ अश्विनौ बृहती ॥} \par}

\vspace{0.12em}

{\centering \textbf{\textcolor{AccentBlue}{कु\raisebox{0.25ex}{\accentmark{15}{\char"1CD2}}ष्ठःकोवा\raisebox{0.25ex}{\accentmark{12}{\char"0951}}मश्विना\kern 0.05em ।\hskip 0.12em plus 0.04em minus 0.02em तपानो\raisebox{0.15ex}{\accentmark{12}{\char"1CF9}}दे\raisebox{0.25ex}{\accentmark{12}{\char"0951}}वामार्क्त्या\raisebox{0.25ex}{\accentmark{12}{\char"0951}}ः\kern 0.05em ।\hskip 0.12em plus 0.04em minus 0.02em घ्नाता\raisebox{0.25ex}{\accentmark{15}{\char"1CD2}}वा\raisebox{0.25ex}{\accentmark{15}{\char"1CD2}}मश्नायाक्षा\raisebox{0.25ex}{\accentmark{12}{\char"0951}}पा\raisebox{0.25ex}{\accentmark{12}{\char"0951}}माणोंशू\raisebox{0.25ex}{\accentmark{15}{\char"1CD2}}ना\kern 0.05em ।\hskip 0.12em plus 0.04em minus 0.02em इ\raisebox{0.25ex}{\accentmark{15}{\char"1CD2}}त्था\raisebox{0.25ex}{\accentmark{15}{\char"1CD2}}मूवा\raisebox{0.25ex}{\accentmark{15}{\char"1CD2}}द्वन्या\raisebox{0.25ex}{\accentmark{15}{\char"1CD2}}था\raisebox{0.25ex}{\accentmark{12}{\char"0951}} \nolinebreak\hskip 0.08em \mbox{॥३॥}}} \par}

\vspace{0.8em}
                % --- SUBSECTIONS (Mantras) ---
                
                % Conditional rendering based on output_mode
                \par\filbreak

\phantomsection

\addcontentsline{toc}{subsection}{ऋक् ४}

{\centering \footnotesize \textcolor{AccentPurple}{॥ प्रस्कण्वः काण्वः इन्द्रः बृहती ॥} \par}

\vspace{0.12em}

{\centering \textbf{\textcolor{AccentBlue}{आयं\raisebox{0.25ex}{\accentmark{15}{\char"1CD2}}वामा\raisebox{0.25ex}{\accentmark{15}{\char"1CD2}}धूमक्तमः\kern 0.05em ।\hskip 0.12em plus 0.04em minus 0.02em सू\raisebox{0.25ex}{\accentmark{15}{\char"1CD2}}तस्सो\raisebox{0.25ex}{\accentmark{15}{\char"1CD2}}मोदीवि\raisebox{0.25ex}{\accentmark{12}{\char"0951}}ष्टिषु\kern 0.05em ।\hskip 0.12em plus 0.04em minus 0.02em तामा\raisebox{0.25ex}{\accentmark{12}{\char"0951}}श्विनापीबतन्तीरो\raisebox{0.25ex}{\accentmark{15}{\char"1CD2}}अन्हं\kern 0.05em ।\hskip 0.12em plus 0.04em minus 0.02em धक्तं\raisebox{0.25ex}{\accentmark{12}{\char"0951}}र\raisebox{0.25ex}{\accentmark{15}{\char"1CD2}}त्ना\raisebox{0.25ex}{\accentmark{12}{\char"0951}}नीदाशू\raisebox{0.25ex}{\accentmark{12}{\char"0951}}षे \nolinebreak\hskip 0.08em \mbox{॥४॥}}} \par}

\vspace{0.8em}
                % --- SUBSECTIONS (Mantras) ---
                
                % Conditional rendering based on output_mode
                \par\filbreak

\phantomsection

\addcontentsline{toc}{subsection}{ऋक् ५}

{\centering \footnotesize \textcolor{AccentPurple}{॥ मेधातिथि{-}मध्येआतिथी काण्वौ इन्द्रः बृहती ॥} \par}

\vspace{0.12em}

{\centering \textbf{\textcolor{AccentBlue}{आत्वासो\raisebox{0.15ex}{\accentmark{12}{\char"1CF9}}मा\raisebox{0.25ex}{\accentmark{12}{\char"0951}}स्याग\raisebox{0.25ex}{\accentmark{15}{\char"1CD2}}ल्दा\raisebox{0.25ex}{\accentmark{12}{\char"0951}}या\kern 0.05em ।\hskip 0.12em plus 0.04em minus 0.02em सा\raisebox{0.25ex}{\accentmark{15}{\char"1CD2}}दाया\raisebox{0.15ex}{\accentmark{12}{\char"1CF9}}चन्नाहं\raisebox{0.25ex}{\accentmark{15}{\char"1CD2}}ज्या\kern 0.05em ।\hskip 0.12em plus 0.04em minus 0.02em भू\raisebox{0.25ex}{\accentmark{15}{\char"1CD2}}र्णिं\raisebox{0.25ex}{\accentmark{15}{\char"1CD2}}मृग\raisebox{0.25ex}{\accentmark{15}{\char"1CD2}}न्नासा\raisebox{0.25ex}{\accentmark{15}{\char"1CD2}}वा\raisebox{0.25ex}{\accentmark{12}{\char"0951}}नेषुचुक्रुधम्\kern 0.05em ।\hskip 0.12em plus 0.04em minus 0.02em कै\raisebox{0.15ex}{\accentmark{12}{\char"1CF8}}शा\raisebox{0.25ex}{\accentmark{12}{\char"0951}}नन्नयाचिषात् \nolinebreak\hskip 0.08em \mbox{॥५॥}}} \par}

\vspace{0.8em}
                % --- SUBSECTIONS (Mantras) ---
                
                % Conditional rendering based on output_mode
                \par\filbreak

\phantomsection

\addcontentsline{toc}{subsection}{ऋक् ६}

{\centering \footnotesize \textcolor{AccentPurple}{॥ देवातिथिः काण्वः इन्द्रः बृहती ॥} \par}

\vspace{0.12em}

{\centering \textbf{\textcolor{AccentBlue}{अध्वर्यो\raisebox{0.15ex}{\accentmark{12}{\char"1CF9}}आ\raisebox{0.25ex}{\accentmark{12}{\char"0951}}द्रावा\raisebox{0.25ex}{\accentmark{15}{\char"1CD2}}यात्वम्\raisebox{0.25ex}{\accentmark{15}{\char"1CD2}}\kern 0.05em ।\hskip 0.12em plus 0.04em minus 0.02em सो\raisebox{0.25ex}{\accentmark{15}{\char"1CD2}}मामिन्द्रा\raisebox{0.25ex}{\accentmark{12}{\char"0951}}ःपिपासति\kern 0.05em ।\hskip 0.12em plus 0.04em minus 0.02em ऊ\raisebox{0.15ex}{\accentmark{12}{\char"1CF8}}पो\raisebox{0.25ex}{\accentmark{12}{\char"0951}}नूनं\raisebox{0.25ex}{\accentmark{15}{\char"1CD2}}यूयूजे\raisebox{0.25ex}{\accentmark{12}{\char"0951}}वृ\raisebox{0.15ex}{\accentmark{12}{\char"1CF9}}षाणाहा\raisebox{0.25ex}{\accentmark{15}{\char"1CD2}}रि\kern 0.05em ।\hskip 0.12em plus 0.04em minus 0.02em आचा\raisebox{0.25ex}{\accentmark{15}{\char"1CD2}}जगा\raisebox{0.25ex}{\accentmark{15}{\char"1CD2}}मा\raisebox{0.25ex}{\accentmark{12}{\char"0951}}वृत्राहा\raisebox{0.25ex}{\accentmark{15}{\char"1CD2}} \nolinebreak\hskip 0.08em \mbox{॥६॥}}} \par}

\vspace{0.8em}
                % --- SUBSECTIONS (Mantras) ---
                
                % Conditional rendering based on output_mode
                \par\filbreak

\phantomsection

\addcontentsline{toc}{subsection}{ऋक् ७}

{\centering \footnotesize \textcolor{AccentPurple}{॥ वसिष्ठो मैत्रावरुणिः इन्द्रः बृहती ॥} \par}

\vspace{0.12em}

{\centering \textbf{\textcolor{AccentBlue}{आ\raisebox{0.25ex}{\accentmark{15}{\char"1CD2}}भीषादस्ता\raisebox{0.25ex}{\accentmark{15}{\char"1CD2}}दा\raisebox{0.25ex}{\accentmark{12}{\char"0951}}भा\raisebox{0.25ex}{\accentmark{12}{\char"0951}}रा\kern 0.05em ।\hskip 0.12em plus 0.04em minus 0.02em इ\raisebox{0.25ex}{\accentmark{15}{\char"1CD2}}न्द्रज्यायःकानी\raisebox{0.25ex}{\accentmark{12}{\char"0951}}यसः\kern 0.05em ।\hskip 0.12em plus 0.04em minus 0.02em पुरोवा\raisebox{0.25ex}{\accentmark{15}{\char"1CD2}}सूहिमाघावन्बभू\raisebox{0.25ex}{\accentmark{15}{\char"1CD2}}वी\raisebox{0.25ex}{\accentmark{12}{\char"0951}}था\kern 0.05em ।\hskip 0.12em plus 0.04em minus 0.02em भारेभरेचाहा\raisebox{0.25ex}{\accentmark{15}{\char"1CD2}}व्या\raisebox{0.25ex}{\accentmark{12}{\char"0951}}ः \nolinebreak\hskip 0.08em \mbox{॥७॥}}} \par}

\vspace{0.8em}
                % --- SUBSECTIONS (Mantras) ---
                
                % Conditional rendering based on output_mode
                \par\filbreak

\phantomsection

\addcontentsline{toc}{subsection}{ऋक् ८}

{\centering \footnotesize \textcolor{AccentPurple}{॥ वसिष्ठो मैत्रावरुणिः इन्द्रः बृहती ॥} \par}

\vspace{0.12em}

{\centering \textbf{\textcolor{AccentBlue}{यादिन्द्राया\raisebox{0.15ex}{\accentmark{12}{\char"1CF8}}वा\raisebox{0.25ex}{\accentmark{12}{\char"0951}}तस्त्वम्\kern 0.05em ।\hskip 0.12em plus 0.04em minus 0.02em एता\raisebox{0.15ex}{\accentmark{12}{\char"1CF9}}वा\raisebox{0.25ex}{\accentmark{12}{\char"0951}}दाह\raisebox{0.25ex}{\accentmark{15}{\char"1CD2}}मीशी\raisebox{0.25ex}{\accentmark{15}{\char"1CD2}}या\kern 0.05em ।\hskip 0.12em plus 0.04em minus 0.02em स्तोतारामी\raisebox{0.25ex}{\accentmark{15}{\char"1CD2}}द्दधिषेरदावसो\kern 0.05em ।\hskip 0.12em plus 0.04em minus 0.02em ना\raisebox{0.15ex}{\accentmark{12}{\char"1CF8}}पा\raisebox{0.25ex}{\accentmark{12}{\char"0951}}पात्वा\raisebox{0.25ex}{\accentmark{15}{\char"1CD2}}यारंसिषः \nolinebreak\hskip 0.08em \mbox{॥८॥}}} \par}

\vspace{0.8em}
                % --- SUBSECTIONS (Mantras) ---
                
                % Conditional rendering based on output_mode
                \par\filbreak

\phantomsection

\addcontentsline{toc}{subsection}{ऋक् ९}

{\centering \footnotesize \textcolor{AccentPurple}{॥ नृमेध आंगिरसः इन्द्रः बृहती ॥} \par}

\vspace{0.12em}

{\centering \textbf{\textcolor{AccentBlue}{त्वा\raisebox{0.15ex}{\accentmark{12}{\char"1CF8}}मि\raisebox{0.25ex}{\accentmark{12}{\char"0951}}न्द्राप्रातू\raisebox{0.25ex}{\accentmark{12}{\char"0951}}र्तिषु\kern 0.05em ।\hskip 0.12em plus 0.04em minus 0.02em आ\raisebox{0.25ex}{\accentmark{15}{\char"1CD2}}भीविश्वा\raisebox{0.25ex}{\accentmark{12}{\char"0951}}आसिस्पृर्धा\raisebox{0.25ex}{\accentmark{12}{\char"0951}}ः\kern 0.05em ।\hskip 0.12em plus 0.04em minus 0.02em आशस्ती\raisebox{0.15ex}{\accentmark{12}{\char"1CF9}}हाज\raisebox{0.25ex}{\accentmark{12}{\char"0951}}नीता\raisebox{0.15ex}{\accentmark{12}{\char"1CF9}}वृ\raisebox{0.25ex}{\accentmark{12}{\char"0951}}त्रातूरा\raisebox{0.25ex}{\accentmark{12}{\char"0951}}सी\kern 0.05em ।\hskip 0.12em plus 0.04em minus 0.02em त्व\raisebox{0.25ex}{\accentmark{15}{\char"1CD2}}न्तू\raisebox{0.25ex}{\accentmark{12}{\char"0951}}र्यतरूष्यातः \nolinebreak\hskip 0.08em \mbox{॥९॥}}} \par}

\vspace{0.8em}
                % --- SUBSECTIONS (Mantras) ---
                
                % Conditional rendering based on output_mode
                \par\filbreak

\phantomsection

\addcontentsline{toc}{subsection}{ऋक् १०}

{\centering \footnotesize \textcolor{AccentPurple}{॥ नोधाः गौतमः इन्द्रः बृहती ॥} \par}

\vspace{0.12em}

{\centering \textbf{\textcolor{AccentBlue}{प्र\raisebox{0.15ex}{\accentmark{12}{\char"1CF8}}योरीरिक्षाओ\raisebox{0.25ex}{\accentmark{15}{\char"1CD2}}ज\raisebox{0.25ex}{\accentmark{12}{\char"0951}}सा\kern 0.05em ।\hskip 0.12em plus 0.04em minus 0.02em दीव\raisebox{0.15ex}{\accentmark{12}{\char"1CF9}}ःस\raisebox{0.15ex}{\accentmark{12}{\char"1CF9}}दो\raisebox{0.25ex}{\accentmark{12}{\char"0951}}ठ\raisebox{0.25ex}{\accentmark{12}{\char"0951}}भ्यस्पा\raisebox{0.25ex}{\accentmark{15}{\char"1CD2}}री\kern 0.05em ।\hskip 0.12em plus 0.04em minus 0.02em न\raisebox{0.25ex}{\accentmark{15}{\char"1CD2}}त्वा\raisebox{0.25ex}{\accentmark{12}{\char"0951}}विव्या\raisebox{0.25ex}{\accentmark{15}{\char"1CD2}}चारा\raisebox{0.25ex}{\accentmark{15}{\char"1CD2}}जाइन्द्रापार्थी\raisebox{0.25ex}{\accentmark{12}{\char"0951}}वम्\kern 0.05em ।\hskip 0.12em plus 0.04em minus 0.02em आ\raisebox{0.25ex}{\accentmark{15}{\char"1CD2}}तीविश्वववक्षिथ \nolinebreak\hskip 0.08em \mbox{॥१०॥}}} \par}

\vspace{0.8em}

            % --- End of Kandah Closing Floral Line & Counts ---
            \par\nopagebreak\vspace{0.9em plus 0.3em minus 0.1em}
            \centerline{\textcolor{black!35}{\rule{1.8cm}{0.35pt}}\;\raisebox{-0.22ex}{\textcolor{vedicmaroon}{\large\SixFlowerPetalDotted}}\;\textcolor{black!35}{\rule{1.8cm}{0.35pt}}}
            \nopagebreak\vspace{0.45em}
            \centerline{\devafont\textbf{\textcolor{AccentBlue}{॥ ऋक् सङ्ख्या : १० ॥}}}


    % --- SUPERSECTION ---
    \clearpage           
    \phantomsection              
    \chapter*{\centering असाविपाठः}
    \addcontentsline{toc}{chapter}{असाविपाठः}
    \markboth{असाविपाठः}{असाविपाठः}
    
    % GLUE FIX: Prevent page break between Chapter and First Section
    \nopagebreak

        

            % --- SECTION ---
            
            \phantomsection  

            \section*{\centering प्रथम खण्डः}
            \addcontentsline{toc}{section}{प्रथम खण्डः} 
            \markright{प्रथम खण्डः}
            \nopagebreak % Keep section title with first mantra
            
            % --- SECTION-LEVEL PROCEDURE LINK ---

                % --- SUBSECTIONS (Mantras) ---
                
                % Conditional rendering based on output_mode
                \par\filbreak

\phantomsection

\addcontentsline{toc}{subsection}{ऋक् १}

{\centering \footnotesize \textcolor{AccentPurple}{॥ वसिष्ठो मैत्रावरुणिः इन्द्रः त्रिष्टुप् ॥} \par}

\vspace{0.12em}

{\centering \textbf{\textcolor{AccentBlue}{आ\raisebox{0.25ex}{\accentmark{15}{\char"1CD2}}सा\raisebox{0.25ex}{\accentmark{12}{\char"0951}}वीदेवङ्गो\raisebox{0.25ex}{\accentmark{12}{\char"0951}}ऋ\raisebox{0.15ex}{\accentmark{12}{\char"1CF8}}जीका\raisebox{0.15ex}{\accentmark{12}{\char"1CF8}}मन्धा\raisebox{0.25ex}{\accentmark{12}{\char"0951}}ः\kern 0.05em ।\hskip 0.12em plus 0.04em minus 0.02em न्य\raisebox{0.25ex}{\accentmark{15}{\char"1CD2}}स्मिन्नि\raisebox{0.25ex}{\accentmark{15}{\char"1CD2}}न्द्रोजनु\raisebox{0.25ex}{\accentmark{15}{\char"1CD2}}षेमूवोच\kern 0.05em ।\hskip 0.12em plus 0.04em minus 0.02em बो\raisebox{0.25ex}{\accentmark{15}{\char"1CD2}}धा\raisebox{0.25ex}{\accentmark{15}{\char"1CD2}}मसित्वाहर्यश्वार्यज्ञैः\kern 0.05em ।\hskip 0.12em plus 0.04em minus 0.02em बो\raisebox{0.25ex}{\accentmark{15}{\char"1CD2}}धा\raisebox{0.25ex}{\accentmark{12}{\char"0951}}नस्तोमाम\raisebox{0.15ex}{\accentmark{12}{\char"1CF9}}न्धा\raisebox{0.25ex}{\accentmark{12}{\char"0951}}सामा\raisebox{0.25ex}{\accentmark{15}{\char"1CD2}}दे\raisebox{0.25ex}{\accentmark{12}{\char"0951}}षु \nolinebreak\hskip 0.08em \mbox{॥१॥}}} \par}

\vspace{0.8em}
                % --- SUBSECTIONS (Mantras) ---
                
                % Conditional rendering based on output_mode
                \par\filbreak

\phantomsection

\addcontentsline{toc}{subsection}{ऋक् २}

{\centering \footnotesize \textcolor{AccentPurple}{॥ वसिष्ठो मैत्रावरुणिः इन्द्रः त्रिष्टुप् ॥} \par}

\vspace{0.12em}

{\centering \textbf{\textcolor{AccentBlue}{यो\raisebox{0.25ex}{\accentmark{15}{\char"1CD2}}नीष्टइन्द्रासा\raisebox{0.25ex}{\accentmark{15}{\char"1CD2}}द\raisebox{0.25ex}{\accentmark{12}{\char"0951}}नेयकारी\kern 0.05em ।\hskip 0.12em plus 0.04em minus 0.02em ता\raisebox{0.25ex}{\accentmark{15}{\char"1CD2}}मा\raisebox{0.25ex}{\accentmark{15}{\char"1CD2}}नृभी\raisebox{0.15ex}{\accentmark{12}{\char"1CF8}}ःपू\raisebox{0.15ex}{\accentmark{12}{\char"1CF8}}रूहूत\raisebox{0.25ex}{\accentmark{15}{\char"1CD2}}प्राया\raisebox{0.25ex}{\accentmark{12}{\char"0951}}हि\kern 0.05em ।\hskip 0.12em plus 0.04em minus 0.02em आ\raisebox{0.25ex}{\accentmark{15}{\char"1CD2}}सोयाथा\raisebox{0.25ex}{\accentmark{15}{\char"1CD2}}नोवीता\raisebox{0.25ex}{\accentmark{15}{\char"1CD2}}वृधा\raisebox{0.25ex}{\accentmark{12}{\char"0951}}श्चीत्\kern 0.05em ।\hskip 0.12em plus 0.04em minus 0.02em ददो\raisebox{0.25ex}{\accentmark{12}{\char"0951}}वा\raisebox{0.25ex}{\accentmark{15}{\char"1CD2}}सू\raisebox{0.25ex}{\accentmark{12}{\char"0951}}नीमामा\raisebox{0.15ex}{\accentmark{12}{\char"1CF9}}दश्चासो\raisebox{0.25ex}{\accentmark{15}{\char"1CD2}}मैः \nolinebreak\hskip 0.08em \mbox{॥२॥}}} \par}

\vspace{0.8em}
                % --- SUBSECTIONS (Mantras) ---
                
                % Conditional rendering based on output_mode
                \par\filbreak

\phantomsection

\addcontentsline{toc}{subsection}{ऋक् ३}

{\centering \footnotesize \textcolor{AccentPurple}{॥ गातरात्रयेः इन्द्रः त्रिष्टुप् ॥} \par}

\vspace{0.12em}

{\centering \textbf{\textcolor{AccentBlue}{आ\raisebox{0.15ex}{\accentmark{12}{\char"1CF8}}द\raisebox{0.25ex}{\accentmark{12}{\char"0951}}दरुत्सामा\raisebox{0.15ex}{\accentmark{12}{\char"1CF9}}सृ\raisebox{0.25ex}{\accentmark{12}{\char"0951}}जोविखानि\kern 0.05em ।\hskip 0.12em plus 0.04em minus 0.02em त्वा\raisebox{0.25ex}{\accentmark{15}{\char"1CD2}}मर्णावान्बल्बधानंआराम्णः\kern 0.05em ।\hskip 0.12em plus 0.04em minus 0.02em मा\raisebox{0.25ex}{\accentmark{15}{\char"1CD2}}हा\raisebox{0.25ex}{\accentmark{15}{\char"1CD2}}न्ता\raisebox{0.25ex}{\accentmark{12}{\char"0951}}मिन्द्राप\raisebox{0.15ex}{\accentmark{12}{\char"1CF9}}र्वातंवीयद्वाः\kern 0.05em ।\hskip 0.12em plus 0.04em minus 0.02em सृ\raisebox{0.25ex}{\accentmark{15}{\char"1CD2}}ज\raisebox{0.25ex}{\accentmark{15}{\char"1CD2}}द्धाराआ\raisebox{0.25ex}{\accentmark{15}{\char"1CD2}}वाय\raisebox{0.15ex}{\accentmark{12}{\char"1CF9}}द\raisebox{0.25ex}{\accentmark{12}{\char"0951}}दानावान्हन् \nolinebreak\hskip 0.08em \mbox{॥३॥}}} \par}

\vspace{0.8em}
                % --- SUBSECTIONS (Mantras) ---
                
                % Conditional rendering based on output_mode
                \par\filbreak

\phantomsection

\addcontentsline{toc}{subsection}{ऋक् ४}

{\centering \footnotesize \textcolor{AccentPurple}{॥ पृतवु्न्यः इन्द्रः त्रिष्टुप् ॥} \par}

\vspace{0.12em}

{\centering \textbf{\textcolor{AccentBlue}{सुष्वाणा\raisebox{0.25ex}{\accentmark{15}{\char"1CD2}}सा\raisebox{0.25ex}{\accentmark{12}{\char"0951}}इन्द्रा\raisebox{0.15ex}{\accentmark{12}{\char"1CF8}}स्तुमा\raisebox{0.25ex}{\accentmark{15}{\char"1CD2}}सीत्वा\kern 0.05em ।\hskip 0.12em plus 0.04em minus 0.02em सनिष्य\raisebox{0.25ex}{\accentmark{15}{\char"1CD2}}न्ता\raisebox{0.25ex}{\accentmark{12}{\char"0951}}श्चीत्तुविनृम्णावा\raisebox{0.25ex}{\accentmark{12}{\char"0951}}जम्\kern 0.05em ।\hskip 0.12em plus 0.04em minus 0.02em आनो\raisebox{0.25ex}{\accentmark{12}{\char"0951}}भरासूवीतं\raisebox{0.15ex}{\accentmark{12}{\char"1CF9}}य\raisebox{0.15ex}{\accentmark{12}{\char"1CF9}}स्याकोना\raisebox{0.25ex}{\accentmark{15}{\char"1CD2}}\kern 0.05em ।\hskip 0.12em plus 0.04em minus 0.02em ताना\raisebox{0.25ex}{\accentmark{15}{\char"1CD2}}त्मनासह्यामात्वो\raisebox{0.25ex}{\accentmark{15}{\char"1CD2}}ता\raisebox{0.25ex}{\accentmark{12}{\char"0951}}ः \nolinebreak\hskip 0.08em \mbox{॥४॥}}} \par}

\vspace{0.8em}
                % --- SUBSECTIONS (Mantras) ---
                
                % Conditional rendering based on output_mode
                \par\filbreak

\phantomsection

\addcontentsline{toc}{subsection}{ऋक् ५}

{\centering \footnotesize \textcolor{AccentPurple}{॥ सप्त्तगुरांगिरसः इन्द्रः (इन्द्रो वैकुण्ठः) त्रिष्टुप् ॥} \par}

\vspace{0.12em}

{\centering \textbf{\textcolor{AccentBlue}{ज\raisebox{0.25ex}{\accentmark{15}{\char"1CD2}}गृ\raisebox{0.25ex}{\accentmark{15}{\char"1CD2}}ह्यातेद\raisebox{0.25ex}{\accentmark{15}{\char"1CD2}}क्षीणमिन्द्राहा\raisebox{0.25ex}{\accentmark{15}{\char"1CD2}}स्तम्\raisebox{0.25ex}{\accentmark{12}{\char"0951}}\kern 0.05em ।\hskip 0.12em plus 0.04em minus 0.02em वसूया\raisebox{0.25ex}{\accentmark{15}{\char"1CD2}}वो\raisebox{0.25ex}{\accentmark{12}{\char"0951}}वसूपातेवा\raisebox{0.25ex}{\accentmark{15}{\char"1CD2}}सूनाम्\kern 0.05em ।\hskip 0.12em plus 0.04em minus 0.02em वि\raisebox{0.25ex}{\accentmark{15}{\char"1CD2}}द्या\raisebox{0.25ex}{\accentmark{15}{\char"1CD2}}हि\raisebox{0.25ex}{\accentmark{15}{\char"1CD2}}त्वागो\raisebox{0.25ex}{\accentmark{15}{\char"1CD2}}पा\raisebox{0.25ex}{\accentmark{15}{\char"1CD2}}तिंसूरागो\raisebox{0.25ex}{\accentmark{15}{\char"1CD2}}नाम्\kern 0.05em ।\hskip 0.12em plus 0.04em minus 0.02em अस्म\raisebox{0.15ex}{\accentmark{12}{\char"1CF9}}भ्य\raisebox{0.25ex}{\accentmark{12}{\char"0951}}ञ्चित्रं\raisebox{0.25ex}{\accentmark{15}{\char"1CD2}}वृ\raisebox{0.25ex}{\accentmark{15}{\char"1CD2}}षाणंरायि\raisebox{0.25ex}{\accentmark{15}{\char"1CD2}}न्दा\raisebox{0.25ex}{\accentmark{15}{\char"1CD2}}ः \nolinebreak\hskip 0.08em \mbox{॥५॥}}} \par}

\vspace{0.8em}
                % --- SUBSECTIONS (Mantras) ---
                
                % Conditional rendering based on output_mode
                \par\filbreak

\phantomsection

\addcontentsline{toc}{subsection}{ऋक् ६}

{\centering \footnotesize \textcolor{AccentPurple}{॥ वसिष्ठो मैत्रावरुणिः इन्द्रः त्रिष्टुप् ॥} \par}

\vspace{0.12em}

{\centering \textbf{\textcolor{AccentBlue}{इ\raisebox{0.25ex}{\accentmark{15}{\char"1CD2}}न्द्रन्ना\raisebox{0.15ex}{\accentmark{12}{\char"1CF9}}रो\raisebox{0.25ex}{\accentmark{12}{\char"0951}}नेमा\raisebox{0.25ex}{\accentmark{15}{\char"1CD2}}धी\raisebox{0.25ex}{\accentmark{12}{\char"0951}}ताहवन्ते\kern 0.05em ।\hskip 0.12em plus 0.04em minus 0.02em यत्पार्यायू\raisebox{0.25ex}{\accentmark{12}{\char"0951}}ना\raisebox{0.15ex}{\accentmark{12}{\char"1CF9}}ज\raisebox{0.25ex}{\accentmark{12}{\char"0951}}न्तेधीयस्ता\raisebox{0.25ex}{\accentmark{15}{\char"1CD2}}ः\kern 0.05em ।\hskip 0.12em plus 0.04em minus 0.02em शू\raisebox{0.25ex}{\accentmark{15}{\char"1CD2}}रोनृ\raisebox{0.15ex}{\accentmark{12}{\char"1CF9}}षा\raisebox{0.25ex}{\accentmark{12}{\char"0951}}ताश्रावा\raisebox{0.25ex}{\accentmark{12}{\char"0951}}साश्चका\raisebox{0.25ex}{\accentmark{15}{\char"1CD2}}मे\kern 0.05em ।\hskip 0.12em plus 0.04em minus 0.02em आ\raisebox{0.25ex}{\accentmark{15}{\char"1CD2}}गो\raisebox{0.25ex}{\accentmark{15}{\char"1CD2}}मा\raisebox{0.25ex}{\accentmark{12}{\char"0951}}तीव्रा\raisebox{0.15ex}{\accentmark{12}{\char"1CF8}}जे\raisebox{0.15ex}{\accentmark{12}{\char"1CF9}}भाजात्व\raisebox{0.25ex}{\accentmark{15}{\char"1CD2}}न्ना\raisebox{0.25ex}{\accentmark{12}{\char"0951}}ः \nolinebreak\hskip 0.08em \mbox{॥६॥}}} \par}

\vspace{0.8em}
                % --- SUBSECTIONS (Mantras) ---
                
                % Conditional rendering based on output_mode
                \par\filbreak

\phantomsection

\addcontentsline{toc}{subsection}{ऋक् ७}

{\centering \footnotesize \textcolor{AccentPurple}{॥ गौिरवीितः शाक्ः इन्द्रः त्रिष्टुप् ॥} \par}

\vspace{0.12em}

{\centering \textbf{\textcolor{AccentBlue}{वायाःसूपर्णाऊपासे\raisebox{0.25ex}{\accentmark{15}{\char"1CD2}}दूरिन्द्रम्\kern 0.05em ।\hskip 0.12em plus 0.04em minus 0.02em प्रीयामेधा\raisebox{0.25ex}{\accentmark{15}{\char"1CD2}}र्षायोना\raisebox{0.25ex}{\accentmark{15}{\char"1CD2}}धा\raisebox{0.25ex}{\accentmark{12}{\char"0951}}मा\raisebox{0.15ex}{\accentmark{12}{\char"1CF9}}नः\kern 0.05em ।\hskip 0.12em plus 0.04em minus 0.02em आपध्वान्तामूर्णहिपूर्धी\raisebox{0.25ex}{\accentmark{12}{\char"0951}}चाक्षुः\kern 0.05em ।\hskip 0.12em plus 0.04em minus 0.02em मुमूग्धाटट्यस्मान्नीधाये\raisebox{0.25ex}{\accentmark{12}{\char"0951}}वाबदध्नान् \nolinebreak\hskip 0.08em \mbox{॥७॥}}} \par}

\vspace{0.8em}
                % --- SUBSECTIONS (Mantras) ---
                
                % Conditional rendering based on output_mode
                \par\filbreak

\phantomsection

\addcontentsline{toc}{subsection}{ऋक् ८}

{\centering \footnotesize \textcolor{AccentPurple}{॥ वेनो भागर्वः वेनः त्रिष्टुप् ॥} \par}

\vspace{0.12em}

{\centering \textbf{\textcolor{AccentBlue}{नाकेसूपर्णा\raisebox{0.25ex}{\accentmark{12}{\char"0951}}मूपा\raisebox{0.25ex}{\accentmark{12}{\char"0951}}यत्पा\raisebox{0.25ex}{\accentmark{12}{\char"0951}}त\raisebox{0.25ex}{\accentmark{12}{\char"0951}}न्तम्\kern 0.05em ।\hskip 0.12em plus 0.04em minus 0.02em हृ\raisebox{0.25ex}{\accentmark{15}{\char"1CD2}}दा\raisebox{0.25ex}{\accentmark{15}{\char"1CD2}}वेन\raisebox{0.25ex}{\accentmark{12}{\char"0951}}न्तोअभ्या\raisebox{0.25ex}{\accentmark{15}{\char"1CD2}}चा\raisebox{0.25ex}{\accentmark{12}{\char"0951}}क्षतत्वा\kern 0.05em ।\hskip 0.12em plus 0.04em minus 0.02em हिर\raisebox{0.25ex}{\accentmark{12}{\char"0951}}ण्यपा\raisebox{0.15ex}{\accentmark{12}{\char"1CF8}}क्षंवारूणस्यादूतम्\raisebox{0.25ex}{\accentmark{15}{\char"1CD2}}\kern 0.05em ।\hskip 0.12em plus 0.04em minus 0.02em यामास्यायो\raisebox{0.25ex}{\accentmark{15}{\char"1CD2}}नू\raisebox{0.25ex}{\accentmark{12}{\char"0951}}शाकूनं\raisebox{0.15ex}{\accentmark{12}{\char"1CF9}}भूरण्यूम् \nolinebreak\hskip 0.08em \mbox{॥८॥}}} \par}

\vspace{0.8em}
                % --- SUBSECTIONS (Mantras) ---
                
                % Conditional rendering based on output_mode
                \par\filbreak

\phantomsection

\addcontentsline{toc}{subsection}{ऋक् ९}

{\centering \footnotesize \textcolor{AccentPurple}{॥ बृहस्पतिकुलो वा इन्द्रः त्रिष्टुप् ॥} \par}

\vspace{0.12em}

{\centering \textbf{\textcolor{AccentBlue}{ब्र\raisebox{0.25ex}{\accentmark{15}{\char"1CD2}}ह्मा\raisebox{0.25ex}{\accentmark{12}{\char"0951}}जज्ञानं\raisebox{0.15ex}{\accentmark{12}{\char"1CF9}}प्राथामंपुरास्ता\raisebox{0.25ex}{\accentmark{12}{\char"0951}}त्\kern 0.05em ।\hskip 0.12em plus 0.04em minus 0.02em वीसीमातस्सूरुचोवेनआ\raisebox{0.25ex}{\accentmark{12}{\char"0951}}वात्\kern 0.05em ।\hskip 0.12em plus 0.04em minus 0.02em सा\raisebox{0.25ex}{\accentmark{15}{\char"1CD2}}बुध्न्याऊ\raisebox{0.25ex}{\accentmark{12}{\char"0951}}पामाआस्याविष्ठाः\kern 0.05em ।\hskip 0.12em plus 0.04em minus 0.02em साताश्चा\raisebox{0.15ex}{\accentmark{12}{\char"1CF9}}योनी\raisebox{0.25ex}{\accentmark{15}{\char"1CD2}}मा\raisebox{0.25ex}{\accentmark{15}{\char"1CD2}}साताश्चावी\raisebox{0.25ex}{\accentmark{15}{\char"1CD2}}वाः \nolinebreak\hskip 0.08em \mbox{॥९॥}}} \par}

\vspace{0.8em}
                % --- SUBSECTIONS (Mantras) ---
                
                % Conditional rendering based on output_mode
                \par\filbreak

\phantomsection

\addcontentsline{toc}{subsection}{ऋक् १०}

{\centering \footnotesize \textcolor{AccentPurple}{॥ सुहोत्रो भारद्वाजः इन्द्रः त्रिष्टुप् ॥} \par}

\vspace{0.12em}

{\centering \textbf{\textcolor{AccentBlue}{आपू\raisebox{0.25ex}{\accentmark{12}{\char"0951}}र्व्यापूरूता\raisebox{0.15ex}{\accentmark{12}{\char"1CF9}}मा\raisebox{0.25ex}{\accentmark{12}{\char"0951}}न्यस्मै\kern 0.05em ।\hskip 0.12em plus 0.04em minus 0.02em माहे\raisebox{0.25ex}{\accentmark{15}{\char"1CD2}}वीराया\raisebox{0.25ex}{\accentmark{12}{\char"0951}}तावा\raisebox{0.15ex}{\accentmark{12}{\char"1CF9}}सेतूराया\raisebox{0.25ex}{\accentmark{12}{\char"0951}}\kern 0.05em ।\hskip 0.12em plus 0.04em minus 0.02em वीरप्सीनेवज्रीणेशन्ता\raisebox{0.25ex}{\accentmark{12}{\char"0951}}मानी\kern 0.05em ।\hskip 0.12em plus 0.04em minus 0.02em वा\raisebox{0.25ex}{\accentmark{15}{\char"1CD2}}चां\raisebox{0.25ex}{\accentmark{12}{\char"0951}}स्यास्मयस्थावी\raisebox{0.25ex}{\accentmark{12}{\char"0951}}रायतक्षुः \nolinebreak\hskip 0.08em \mbox{॥१०॥}}} \par}

\vspace{0.8em}

            % --- End of Kandah Closing Floral Line & Counts ---
            \par\nopagebreak\vspace{0.9em plus 0.3em minus 0.1em}
            \centerline{\textcolor{black!35}{\rule{1.8cm}{0.35pt}}\;\raisebox{-0.22ex}{\textcolor{vedicmaroon}{\large\SixFlowerPetalDotted}}\;\textcolor{black!35}{\rule{1.8cm}{0.35pt}}}
            \nopagebreak\vspace{0.45em}
            \centerline{\devafont\textbf{\textcolor{AccentBlue}{॥ ऋक् सङ्ख्या : १० ॥}}}
        
            \clearpage

            % --- SECTION ---
            
            \phantomsection  

            \section*{\centering द्वितीय खण्डः}
            \addcontentsline{toc}{section}{द्वितीय खण्डः} 
            \markright{द्वितीय खण्डः}
            \nopagebreak % Keep section title with first mantra
            
            % --- SECTION-LEVEL PROCEDURE LINK ---

                % --- SUBSECTIONS (Mantras) ---
                
                % Conditional rendering based on output_mode
                \par\filbreak

\phantomsection

\addcontentsline{toc}{subsection}{ऋक् १}

{\centering \footnotesize \textcolor{AccentPurple}{॥ द्युतानो मारुतः (तिरश्चीरांगिरसः) इन्द्रः त्रिष्टुप् ॥} \par}

\vspace{0.12em}

{\centering \textbf{\textcolor{AccentBlue}{आ\raisebox{0.25ex}{\accentmark{15}{\char"1CD2}}व\raisebox{0.25ex}{\accentmark{12}{\char"0951}}द्रप्सो\raisebox{0.15ex}{\accentmark{12}{\char"1CF9}}अं\raisebox{0.25ex}{\accentmark{12}{\char"0951}}शूमा\raisebox{0.25ex}{\accentmark{15}{\char"1CD2}}तीमतिष्ठात्\kern 0.05em ।\hskip 0.12em plus 0.04em minus 0.02em इया\raisebox{0.25ex}{\accentmark{15}{\char"1CD2}}न\raisebox{0.25ex}{\accentmark{15}{\char"1CD2}}ःकृष्णो\raisebox{0.25ex}{\accentmark{15}{\char"1CD2}}दशा\raisebox{0.15ex}{\accentmark{12}{\char"1CF9}}भीस्साहा\raisebox{0.25ex}{\accentmark{15}{\char"1CD2}}स्रै\raisebox{0.25ex}{\accentmark{12}{\char"0951}}ः\kern 0.05em ।\hskip 0.12em plus 0.04em minus 0.02em आवक्ता\raisebox{0.25ex}{\accentmark{15}{\char"1CD2}}मिन्द्र\raisebox{0.25ex}{\accentmark{12}{\char"0951}}श्शच्या\raisebox{0.25ex}{\accentmark{15}{\char"1CD2}}धामन्तम्\kern 0.05em ।\hskip 0.12em plus 0.04em minus 0.02em आ\raisebox{0.25ex}{\accentmark{15}{\char"1CD2}}पस्नीही\raisebox{0.25ex}{\accentmark{12}{\char"0951}}तिंनृमा\raisebox{0.25ex}{\accentmark{15}{\char"1CD2}}णा\raisebox{0.25ex}{\accentmark{12}{\char"0951}}आधद्रा\raisebox{0.25ex}{\accentmark{15}{\char"1CD2}}ः \nolinebreak\hskip 0.08em \mbox{॥१॥}}} \par}

\vspace{0.8em}
                % --- SUBSECTIONS (Mantras) ---
                
                % Conditional rendering based on output_mode
                \par\filbreak

\phantomsection

\addcontentsline{toc}{subsection}{ऋक् २}

{\centering \footnotesize \textcolor{AccentPurple}{॥ द्युतानो मारुतः (तिरश्चीरांगिरसः) इन्द्रः त्रिष्टुप् ॥} \par}

\vspace{0.12em}

{\centering \textbf{\textcolor{AccentBlue}{वृत्रास्या\raisebox{0.25ex}{\accentmark{15}{\char"1CD2}}त्वश्वासाथादीषामाणः\kern 0.05em ।\hskip 0.12em plus 0.04em minus 0.02em विश्वेदेवा\raisebox{0.25ex}{\accentmark{15}{\char"1CD2}}आ\raisebox{0.25ex}{\accentmark{12}{\char"0951}}जाहुर्ये\raisebox{0.25ex}{\accentmark{15}{\char"1CD2}}साखायः\kern 0.05em ।\hskip 0.12em plus 0.04em minus 0.02em मारू\raisebox{0.25ex}{\accentmark{15}{\char"1CD2}}द्धी\raisebox{0.25ex}{\accentmark{12}{\char"0951}}रिन्द्रासख्य\raisebox{0.25ex}{\accentmark{15}{\char"1CD2}}न्ते\raisebox{0.25ex}{\accentmark{12}{\char"0951}}अस्तु\kern 0.05em ।\hskip 0.12em plus 0.04em minus 0.02em आ\raisebox{0.25ex}{\accentmark{15}{\char"1CD2}}थेमा\raisebox{0.25ex}{\accentmark{15}{\char"1CD2}}विश्वाःपृ\raisebox{0.25ex}{\accentmark{15}{\char"1CD2}}नाजयासि \nolinebreak\hskip 0.08em \mbox{॥२॥}}} \par}

\vspace{0.8em}
                % --- SUBSECTIONS (Mantras) ---
                
                % Conditional rendering based on output_mode
                \par\filbreak

\phantomsection

\addcontentsline{toc}{subsection}{ऋक् ३}

{\centering \footnotesize \textcolor{AccentPurple}{॥ बृहदुक्तो वामदेव्यः इन्द्रः त्रिष्टुप् ॥} \par}

\vspace{0.12em}

{\centering \textbf{\textcolor{AccentBlue}{विधु\raisebox{0.15ex}{\accentmark{12}{\char"1CF9}}न्दद्रा\raisebox{0.25ex}{\accentmark{12}{\char"0951}}णंसा\raisebox{0.25ex}{\accentmark{15}{\char"1CD2}}मा\raisebox{0.25ex}{\accentmark{15}{\char"1CD2}}नेबहूनाम्\kern 0.05em ।\hskip 0.12em plus 0.04em minus 0.02em यू\raisebox{0.25ex}{\accentmark{15}{\char"1CD2}}वा\raisebox{0.25ex}{\accentmark{15}{\char"1CD2}}नंस\raisebox{0.25ex}{\accentmark{15}{\char"1CD2}}न्तं\raisebox{0.25ex}{\accentmark{12}{\char"0951}}पालीतो\raisebox{0.25ex}{\accentmark{15}{\char"1CD2}}ज\raisebox{0.25ex}{\accentmark{12}{\char"0951}}गारा\kern 0.05em ।\hskip 0.12em plus 0.04em minus 0.02em दे\raisebox{0.25ex}{\accentmark{15}{\char"1CD2}}वा\raisebox{0.25ex}{\accentmark{15}{\char"1CD2}}स्या\raisebox{0.25ex}{\accentmark{12}{\char"0951}}पाश्या\raisebox{0.25ex}{\accentmark{12}{\char"0951}}का\raisebox{0.25ex}{\accentmark{15}{\char"1CD2}}व्यं\raisebox{0.25ex}{\accentmark{12}{\char"0951}}मा\raisebox{0.25ex}{\accentmark{12}{\char"0951}}हित्वा\raisebox{0.25ex}{\accentmark{12}{\char"0951}}\kern 0.05em ।\hskip 0.12em plus 0.04em minus 0.02em अ\raisebox{0.25ex}{\accentmark{15}{\char"1CD2}}द्या\raisebox{0.25ex}{\accentmark{15}{\char"1CD2}}मा\raisebox{0.25ex}{\accentmark{15}{\char"1CD2}}मारास्स\raisebox{0.25ex}{\accentmark{15}{\char"1CD2}}ह्य\raisebox{0.25ex}{\accentmark{15}{\char"1CD2}}स्सामा\raisebox{0.25ex}{\accentmark{15}{\char"1CD2}}ना \nolinebreak\hskip 0.08em \mbox{॥३॥}}} \par}

\vspace{0.8em}
                % --- SUBSECTIONS (Mantras) ---
                
                % Conditional rendering based on output_mode
                \par\filbreak

\phantomsection

\addcontentsline{toc}{subsection}{ऋक् ४}

{\centering \footnotesize \textcolor{AccentPurple}{॥ द्युतानो मारुतः (तिरश्चीरांगिरसः) इन्द्रः त्रिष्टुप् ॥} \par}

\vspace{0.12em}

{\centering \textbf{\textcolor{AccentBlue}{त्वंहत्यस्सप्त\raisebox{0.25ex}{\accentmark{15}{\char"1CD2}}भ्योजा\raisebox{0.25ex}{\accentmark{15}{\char"1CD2}}या\raisebox{0.25ex}{\accentmark{12}{\char"0951}}मानः\kern 0.05em ।\hskip 0.12em plus 0.04em minus 0.02em अ\raisebox{0.25ex}{\accentmark{15}{\char"1CD2}}श\raisebox{0.25ex}{\accentmark{15}{\char"1CD2}}त्रुभ्योअभावश्श\raisebox{0.25ex}{\accentmark{15}{\char"1CD2}}त्रू\raisebox{0.25ex}{\accentmark{12}{\char"0951}}रिन्द्रा\kern 0.05em ।\hskip 0.12em plus 0.04em minus 0.02em गू\raisebox{0.25ex}{\accentmark{15}{\char"1CD2}}ढे\raisebox{0.25ex}{\accentmark{15}{\char"1CD2}}द्या\raisebox{0.25ex}{\accentmark{15}{\char"1CD2}}वा\raisebox{0.25ex}{\accentmark{12}{\char"0951}}पृथिवीअ\raisebox{0.25ex}{\accentmark{15}{\char"1CD2}}न्वा\raisebox{0.25ex}{\accentmark{15}{\char"1CD2}}विन्दः\kern 0.05em ।\hskip 0.12em plus 0.04em minus 0.02em वि\raisebox{0.25ex}{\accentmark{15}{\char"1CD2}}भूमा\raisebox{0.25ex}{\accentmark{15}{\char"1CD2}}द्भ्योभू\raisebox{0.25ex}{\accentmark{15}{\char"1CD2}}वा\raisebox{0.25ex}{\accentmark{12}{\char"0951}}नेभ्योरा\raisebox{0.25ex}{\accentmark{15}{\char"1CD2}}ण\raisebox{0.25ex}{\accentmark{15}{\char"1CD2}}न्धा\raisebox{0.25ex}{\accentmark{12}{\char"0951}}ः \nolinebreak\hskip 0.08em \mbox{॥४॥}}} \par}

\vspace{0.8em}
                % --- SUBSECTIONS (Mantras) ---
                
                % Conditional rendering based on output_mode
                \par\filbreak

\phantomsection

\addcontentsline{toc}{subsection}{ऋक् ५}

{\centering \footnotesize \textcolor{AccentPurple}{॥ वामदेवोः गोतमः इन्द्रः त्रिष्टुप् ॥} \par}

\vspace{0.12em}

{\centering \textbf{\textcolor{AccentBlue}{मेळि\raisebox{0.15ex}{\accentmark{12}{\char"1CF9}}न्न\raisebox{0.15ex}{\accentmark{12}{\char"1CF9}}त्वा\raisebox{0.25ex}{\accentmark{12}{\char"0951}}वज्रीणं\raisebox{0.25ex}{\accentmark{12}{\char"0951}}भृष्टीम\raisebox{0.25ex}{\accentmark{15}{\char"1CD2}}न्तम्\kern 0.05em ।\hskip 0.12em plus 0.04em minus 0.02em पु\raisebox{0.25ex}{\accentmark{15}{\char"1CD2}}रूधा\raisebox{0.25ex}{\accentmark{15}{\char"1CD2}}स्मानं\raisebox{0.25ex}{\accentmark{12}{\char"0951}}वृषाभंस्थीरा\raisebox{0.25ex}{\accentmark{15}{\char"1CD2}}प्सु\raisebox{0.25ex}{\accentmark{12}{\char"0951}}नुम्\kern 0.05em ।\hskip 0.12em plus 0.04em minus 0.02em का\raisebox{0.25ex}{\accentmark{15}{\char"1CD2}}रो\raisebox{0.25ex}{\accentmark{12}{\char"0951}}ष्यर्यस्तारुषीर्दूवस्यूः\kern 0.05em ।\hskip 0.12em plus 0.04em minus 0.02em डन्दइन्द्र\raisebox{0.25ex}{\accentmark{12}{\char"0951}}द्युक्षंवृ\raisebox{0.25ex}{\accentmark{12}{\char"0951}}त्राहा\raisebox{0.25ex}{\accentmark{15}{\char"1CD2}}णंगृणीषे \nolinebreak\hskip 0.08em \mbox{॥५॥}}} \par}

\vspace{0.8em}
                % --- SUBSECTIONS (Mantras) ---
                
                % Conditional rendering based on output_mode
                \par\filbreak

\phantomsection

\addcontentsline{toc}{subsection}{ऋक् ६}

{\centering \footnotesize \textcolor{AccentPurple}{॥ वसिष्ठो मैत्रावरुणिः इन्द्रः ऋषिः विराट् ॥} \par}

\vspace{0.12em}

{\centering \textbf{\textcolor{AccentBlue}{प्रावो\raisebox{0.25ex}{\accentmark{12}{\char"0951}}माहे\raisebox{0.15ex}{\accentmark{12}{\char"1CF9}}मा\raisebox{0.25ex}{\accentmark{12}{\char"0951}}हेवृ\raisebox{0.25ex}{\accentmark{15}{\char"1CD2}}धे\raisebox{0.25ex}{\accentmark{15}{\char"1CD2}}भरध्वम्\kern 0.05em ।\hskip 0.12em plus 0.04em minus 0.02em प्रा\raisebox{0.25ex}{\accentmark{15}{\char"1CD2}}चे\raisebox{0.25ex}{\accentmark{12}{\char"0951}}तासेप्रासू\raisebox{0.15ex}{\accentmark{12}{\char"1CF8}}मातिंकृ\raisebox{0.25ex}{\accentmark{12}{\char"0951}}णुध्वम्\kern 0.05em ।\hskip 0.12em plus 0.04em minus 0.02em वीशःपूर्वीप्रा\raisebox{0.25ex}{\accentmark{15}{\char"1CD2}}चा\raisebox{0.25ex}{\accentmark{12}{\char"0951}}राचर्षाणिप्रा\raisebox{0.25ex}{\accentmark{12}{\char"0951}}ः \nolinebreak\hskip 0.08em \mbox{॥६॥}}} \par}

\vspace{0.8em}
                % --- SUBSECTIONS (Mantras) ---
                
                % Conditional rendering based on output_mode
                \par\filbreak

\phantomsection

\addcontentsline{toc}{subsection}{ऋक् ७}

{\centering \footnotesize \textcolor{AccentPurple}{॥ विश्वामित्रो गाथिनः इन्द्रः त्रिष्टुप् ॥} \par}

\vspace{0.12em}

{\centering \textbf{\textcolor{AccentBlue}{शू\raisebox{0.25ex}{\accentmark{15}{\char"1CD2}}नं\raisebox{0.25ex}{\accentmark{15}{\char"1CD2}}हूवेमामाघा\raisebox{0.15ex}{\accentmark{12}{\char"1CF9}}वानामिन्द्रम्\kern 0.05em ।\hskip 0.12em plus 0.04em minus 0.02em अ\raisebox{0.25ex}{\accentmark{15}{\char"1CD2}}स्मिन्भा\raisebox{0.25ex}{\accentmark{15}{\char"1CD2}}रेनृ\raisebox{0.15ex}{\accentmark{12}{\char"1CF9}}तामंवा\raisebox{0.25ex}{\accentmark{15}{\char"1CD2}}ज\raisebox{0.25ex}{\accentmark{12}{\char"0951}}सातौ\kern 0.05em ।\hskip 0.12em plus 0.04em minus 0.02em श्रृण्वन्ता\raisebox{0.25ex}{\accentmark{15}{\char"1CD2}}मुग्रा\raisebox{0.25ex}{\accentmark{15}{\char"1CD2}}मूता\raisebox{0.15ex}{\accentmark{12}{\char"1CF9}}येसामात्सु\raisebox{0.25ex}{\accentmark{15}{\char"1CD2}}\kern 0.05em ।\hskip 0.12em plus 0.04em minus 0.02em घ्नन्तं\raisebox{0.25ex}{\accentmark{12}{\char"0951}}वृत्राणि\raisebox{0.25ex}{\accentmark{12}{\char"0951}}संची\raisebox{0.25ex}{\accentmark{15}{\char"1CD2}}तन्धा\raisebox{0.25ex}{\accentmark{15}{\char"1CD2}}ना\raisebox{0.25ex}{\accentmark{12}{\char"0951}}नी \nolinebreak\hskip 0.08em \mbox{॥७॥}}} \par}

\vspace{0.8em}
                % --- SUBSECTIONS (Mantras) ---
                
                % Conditional rendering based on output_mode
                \par\filbreak

\phantomsection

\addcontentsline{toc}{subsection}{ऋक् ८}

{\centering \footnotesize \textcolor{AccentPurple}{॥ वसिष्ठो मैत्रावरुणिः इन्द्रः त्रिष्टुप् ॥} \par}

\vspace{0.12em}

{\centering \textbf{\textcolor{AccentBlue}{ऊदुब्र\raisebox{0.25ex}{\accentmark{12}{\char"0951}}ह्माणैरतश्रवस्या\raisebox{0.25ex}{\accentmark{15}{\char"1CD2}}\kern 0.05em ।\hskip 0.12em plus 0.04em minus 0.02em इन्द्रंसामर्येमाहयावसिष्ठा\kern 0.05em ।\hskip 0.12em plus 0.04em minus 0.02em आ\raisebox{0.15ex}{\accentmark{12}{\char"1CF8}}यो\raisebox{0.15ex}{\accentmark{12}{\char"1CF8}}र्विश्वा\raisebox{0.25ex}{\accentmark{12}{\char"0951}}नीश्रा\raisebox{0.25ex}{\accentmark{15}{\char"1CD2}}वासताताना\raisebox{0.25ex}{\accentmark{12}{\char"0951}}\kern 0.05em ।\hskip 0.12em plus 0.04em minus 0.02em ऊपश्रो\raisebox{0.25ex}{\accentmark{15}{\char"1CD2}}ता\raisebox{0.15ex}{\accentmark{12}{\char"1CF9}}मा\raisebox{0.25ex}{\accentmark{15}{\char"1CD2}}इवा\raisebox{0.25ex}{\accentmark{15}{\char"1CD2}}तोवा\raisebox{0.25ex}{\accentmark{15}{\char"1CD2}}चां\raisebox{0.25ex}{\accentmark{12}{\char"0951}}सि \nolinebreak\hskip 0.08em \mbox{॥८॥}}} \par}

\vspace{0.8em}
                % --- SUBSECTIONS (Mantras) ---
                
                % Conditional rendering based on output_mode
                \par\filbreak

\phantomsection

\addcontentsline{toc}{subsection}{ऋक् ९}

{\centering \footnotesize \textcolor{AccentPurple}{॥ गौिरवीितः शाक्ः इन्द्रः त्रिष्टुप् ॥} \par}

\vspace{0.12em}

{\centering \textbf{\textcolor{AccentBlue}{च\raisebox{0.25ex}{\accentmark{15}{\char"1CD2}}क्रं\raisebox{0.15ex}{\accentmark{12}{\char"1CF9}}याद\raisebox{0.25ex}{\accentmark{12}{\char"0951}}स्यप्सू\raisebox{0.25ex}{\accentmark{15}{\char"1CD2}}आनिषाक्तम्\kern 0.05em ।\hskip 0.12em plus 0.04em minus 0.02em ऊतो\raisebox{0.15ex}{\accentmark{12}{\char"1CF9}}ता\raisebox{0.15ex}{\accentmark{12}{\char"1CF9}}द\raisebox{0.25ex}{\accentmark{12}{\char"0951}}स्मै\raisebox{0.25ex}{\accentmark{15}{\char"1CD2}}म\raisebox{0.25ex}{\accentmark{15}{\char"1CD2}}ध्वी\raisebox{0.25ex}{\accentmark{15}{\char"1CD2}}चा\raisebox{0.25ex}{\accentmark{12}{\char"0951}}च्छद्यात्\kern 0.05em ।\hskip 0.12em plus 0.04em minus 0.02em पृ\raisebox{0.25ex}{\accentmark{15}{\char"1CD2}}थि\raisebox{0.25ex}{\accentmark{15}{\char"1CD2}}व्या\raisebox{0.25ex}{\accentmark{12}{\char"0951}}मा\raisebox{0.25ex}{\accentmark{15}{\char"1CD2}}तीषितंया\raisebox{0.25ex}{\accentmark{15}{\char"1CD2}}दूधाः\kern 0.05em ।\hskip 0.12em plus 0.04em minus 0.02em पायोगो\raisebox{0.15ex}{\accentmark{12}{\char"1CF9}}ष्वा\raisebox{0.15ex}{\accentmark{12}{\char"1CF9}}द\raisebox{0.25ex}{\accentmark{12}{\char"0951}}धाओ\raisebox{0.25ex}{\accentmark{15}{\char"1CD2}}षाधीषु \nolinebreak\hskip 0.08em \mbox{॥९॥}}} \par}

\vspace{0.8em}

            % --- End of Kandah Closing Floral Line & Counts ---
            \par\nopagebreak\vspace{0.9em plus 0.3em minus 0.1em}
            \centerline{\textcolor{black!35}{\rule{1.8cm}{0.35pt}}\;\raisebox{-0.22ex}{\textcolor{vedicmaroon}{\large\SixFlowerPetalDotted}}\;\textcolor{black!35}{\rule{1.8cm}{0.35pt}}}
            \nopagebreak\vspace{0.45em}
            \centerline{\devafont\textbf{\textcolor{AccentBlue}{॥ ऋक् सङ्ख्या : ९ ॥}}}
        
            \clearpage

            % --- SECTION ---
            
            \phantomsection  

            \section*{\centering तृतीय खण्डः}
            \addcontentsline{toc}{section}{तृतीय खण्डः} 
            \markright{तृतीय खण्डः}
            \nopagebreak % Keep section title with first mantra
            
            % --- SECTION-LEVEL PROCEDURE LINK ---

                % --- SUBSECTIONS (Mantras) ---
                
                % Conditional rendering based on output_mode
                \par\filbreak

\phantomsection

\addcontentsline{toc}{subsection}{ऋक् १}

{\centering \footnotesize \textcolor{AccentPurple}{॥ अरिष्टनेमिस्ताक्ष्र्यः इन्द्रः (ताक्ष्र्यः) त्रिष्टुप् ॥} \par}

\vspace{0.12em}

{\centering \textbf{\textcolor{AccentBlue}{त्या\raisebox{0.25ex}{\accentmark{15}{\char"1CD2}}मूषुवाजी\raisebox{0.25ex}{\accentmark{15}{\char"1CD2}}नन्देवा\raisebox{0.25ex}{\accentmark{15}{\char"1CD2}}जू\raisebox{0.25ex}{\accentmark{12}{\char"0951}}तम्\kern 0.05em ।\hskip 0.12em plus 0.04em minus 0.02em स\raisebox{0.25ex}{\accentmark{15}{\char"1CD2}}होवा\raisebox{0.25ex}{\accentmark{15}{\char"1CD2}}नं\raisebox{0.25ex}{\accentmark{12}{\char"0951}}ता\raisebox{0.25ex}{\accentmark{15}{\char"1CD2}}रूता\raisebox{0.25ex}{\accentmark{15}{\char"1CD2}}रंराथा\raisebox{0.25ex}{\accentmark{12}{\char"0951}}नाम्\kern 0.05em ।\hskip 0.12em plus 0.04em minus 0.02em आ\raisebox{0.25ex}{\accentmark{15}{\char"1CD2}}री\raisebox{0.25ex}{\accentmark{12}{\char"0951}}ष्टनेमिंपृतना\raisebox{0.15ex}{\accentmark{12}{\char"1CF9}}ज\raisebox{0.25ex}{\accentmark{12}{\char"0951}}माशू\raisebox{0.25ex}{\accentmark{15}{\char"1CD2}}म्\kern 0.05em ।\hskip 0.12em plus 0.04em minus 0.02em स्व\raisebox{0.25ex}{\accentmark{15}{\char"1CD2}}स्ताये\raisebox{0.25ex}{\accentmark{15}{\char"1CD2}}ता\raisebox{0.15ex}{\accentmark{12}{\char"1CF9}}क्ष्यामीहाहू\raisebox{0.25ex}{\accentmark{12}{\char"0951}}वेम \nolinebreak\hskip 0.08em \mbox{॥१॥}}} \par}

\vspace{0.8em}
                % --- SUBSECTIONS (Mantras) ---
                
                % Conditional rendering based on output_mode
                \par\filbreak

\phantomsection

\addcontentsline{toc}{subsection}{ऋक् २}

{\centering \footnotesize \textcolor{AccentPurple}{॥ भरद्वाजः (गर्गो भारद्वाजः) इन्द्रः त्रिष्टुप् ॥} \par}

\vspace{0.12em}

{\centering \textbf{\textcolor{AccentBlue}{त्रा\raisebox{0.25ex}{\accentmark{15}{\char"1CD2}}ता\raisebox{0.25ex}{\accentmark{15}{\char"1CD2}}रमि\raisebox{0.25ex}{\accentmark{15}{\char"1CD2}}न्द्रामा\raisebox{0.25ex}{\accentmark{15}{\char"1CD2}}वीता\raisebox{0.25ex}{\accentmark{15}{\char"1CD2}}रामि\raisebox{0.25ex}{\accentmark{15}{\char"1CD2}}न्द्रम्\raisebox{0.25ex}{\accentmark{12}{\char"0951}}\kern 0.05em ।\hskip 0.12em plus 0.04em minus 0.02em हा\raisebox{0.25ex}{\accentmark{15}{\char"1CD2}}वे\raisebox{0.25ex}{\accentmark{12}{\char"0951}}हवेसूहावंशू\raisebox{0.25ex}{\accentmark{15}{\char"1CD2}}रामि\raisebox{0.25ex}{\accentmark{12}{\char"0951}}न्द्रम्\kern 0.05em ।\hskip 0.12em plus 0.04em minus 0.02em हूवे\raisebox{0.25ex}{\accentmark{15}{\char"1CD2}}नूशक्रं\raisebox{0.25ex}{\accentmark{15}{\char"1CD2}}पू\raisebox{0.25ex}{\accentmark{12}{\char"0951}}रुहूतामिन्द्रम्\kern 0.05em ।\hskip 0.12em plus 0.04em minus 0.02em ई\raisebox{0.25ex}{\accentmark{15}{\char"1CD2}}दंहा\raisebox{0.25ex}{\accentmark{15}{\char"1CD2}}वि\raisebox{0.25ex}{\accentmark{15}{\char"1CD2}}र्मा\raisebox{0.25ex}{\accentmark{12}{\char"0951}}घा\raisebox{0.15ex}{\accentmark{12}{\char"1CF9}}वावे\raisebox{0.25ex}{\accentmark{12}{\char"0951}}त्विन्द्रा\raisebox{0.25ex}{\accentmark{12}{\char"0951}}ः \nolinebreak\hskip 0.08em \mbox{॥२॥}}} \par}

\vspace{0.8em}
                % --- SUBSECTIONS (Mantras) ---
                
                % Conditional rendering based on output_mode
                \par\filbreak

\phantomsection

\addcontentsline{toc}{subsection}{ऋक् ३}

{\centering \footnotesize \textcolor{AccentPurple}{॥ विमद ऐन्द्र वसुकृद्वा वासुकः (प्रजापत्यो वा) इन्द्रः त्रिष्टुप् ॥} \par}

\vspace{0.12em}

{\centering \textbf{\textcolor{AccentBlue}{या\raisebox{0.25ex}{\accentmark{15}{\char"1CD2}}जा\raisebox{0.25ex}{\accentmark{15}{\char"1CD2}}माहाइ\raisebox{0.25ex}{\accentmark{15}{\char"1CD2}}न्द्रंव\raisebox{0.25ex}{\accentmark{15}{\char"1CD2}}ज्रा\raisebox{0.25ex}{\accentmark{12}{\char"0951}}दक्षिणम्\kern 0.05em ।\hskip 0.12em plus 0.04em minus 0.02em हा\raisebox{0.25ex}{\accentmark{15}{\char"1CD2}}री\raisebox{0.25ex}{\accentmark{12}{\char"0951}}णांरथ्यांटट्यवि\raisebox{0.25ex}{\accentmark{15}{\char"1CD2}}व्रा\raisebox{0.25ex}{\accentmark{12}{\char"0951}}तानाम्\kern 0.05em ।\hskip 0.12em plus 0.04em minus 0.02em प्र\raisebox{0.15ex}{\accentmark{12}{\char"1CF8}}श्म\raisebox{0.15ex}{\accentmark{12}{\char"1CF8}}श्रू\raisebox{0.25ex}{\accentmark{12}{\char"0951}}भिर्दोधुवदू\raisebox{0.25ex}{\accentmark{15}{\char"1CD2}}र्ध्वाधा\raisebox{0.25ex}{\accentmark{12}{\char"0951}}भुवात्\kern 0.05em ।\hskip 0.12em plus 0.04em minus 0.02em वी\raisebox{0.15ex}{\accentmark{12}{\char"1CF8}}से\raisebox{0.15ex}{\accentmark{12}{\char"1CF8}}ना\raisebox{0.25ex}{\accentmark{12}{\char"0951}}भिर्भा\raisebox{0.25ex}{\accentmark{15}{\char"1CD2}}या\raisebox{0.25ex}{\accentmark{12}{\char"0951}}मानोवीरा\raisebox{0.25ex}{\accentmark{15}{\char"1CD2}}धा\raisebox{0.25ex}{\accentmark{12}{\char"0951}}सा \nolinebreak\hskip 0.08em \mbox{॥३॥}}} \par}

\vspace{0.8em}
                % --- SUBSECTIONS (Mantras) ---
                
                % Conditional rendering based on output_mode
                \par\filbreak

\phantomsection

\addcontentsline{toc}{subsection}{ऋक् ४}

{\centering \footnotesize \textcolor{AccentPurple}{॥ वामदेवो गौतमः इन्द्रः त्रिष्टुप् ॥} \par}

\vspace{0.12em}

{\centering \textbf{\textcolor{AccentBlue}{सत्रा\raisebox{0.25ex}{\accentmark{15}{\char"1CD2}}हा\raisebox{0.25ex}{\accentmark{12}{\char"0951}}णंदाधृ\raisebox{0.25ex}{\accentmark{12}{\char"0951}}षिन्तूम्रामिन्द्रम्\raisebox{0.25ex}{\accentmark{12}{\char"0951}}\kern 0.05em ।\hskip 0.12em plus 0.04em minus 0.02em मा\raisebox{0.25ex}{\accentmark{15}{\char"1CD2}}हामा\raisebox{0.25ex}{\accentmark{12}{\char"0951}}पारं\raisebox{0.15ex}{\accentmark{12}{\char"1CF9}}वृ\raisebox{0.25ex}{\accentmark{12}{\char"0951}}षाभं\raisebox{0.25ex}{\accentmark{15}{\char"1CD2}}सूवा\raisebox{0.25ex}{\accentmark{15}{\char"1CD2}}ज्र\raisebox{0.25ex}{\accentmark{12}{\char"0951}}म्\kern 0.05em ।\hskip 0.12em plus 0.04em minus 0.02em हन्तायो\raisebox{0.25ex}{\accentmark{15}{\char"1CD2}}वृत्रं\raisebox{0.15ex}{\accentmark{12}{\char"1CF9}}सानीतोता\raisebox{0.15ex}{\accentmark{12}{\char"1CF9}}वा\raisebox{0.25ex}{\accentmark{15}{\char"1CD2}}जम्\raisebox{0.25ex}{\accentmark{12}{\char"0951}}\kern 0.05em ।\hskip 0.12em plus 0.04em minus 0.02em दातामाघा\raisebox{0.15ex}{\accentmark{12}{\char"1CF9}}नीमाघा\raisebox{0.15ex}{\accentmark{12}{\char"1CF9}}वा\raisebox{0.25ex}{\accentmark{12}{\char"0951}}सूरा\raisebox{0.25ex}{\accentmark{15}{\char"1CD2}}धा\raisebox{0.25ex}{\accentmark{12}{\char"0951}}ः \nolinebreak\hskip 0.08em \mbox{॥४॥}}} \par}

\vspace{0.8em}
                % --- SUBSECTIONS (Mantras) ---
                
                % Conditional rendering based on output_mode
                \par\filbreak

\phantomsection

\addcontentsline{toc}{subsection}{ऋक् ५}

{\centering \footnotesize \textcolor{AccentPurple}{॥ वामदेवो गौतमः इन्द्रः त्रिष्टुप् ॥} \par}

\vspace{0.12em}

{\centering \textbf{\textcolor{AccentBlue}{यो\raisebox{0.25ex}{\accentmark{15}{\char"1CD2}}नो\raisebox{0.25ex}{\accentmark{12}{\char"0951}}वनूष्य\raisebox{0.15ex}{\accentmark{12}{\char"1CF9}}न्ना\raisebox{0.25ex}{\accentmark{15}{\char"1CD2}}भीदा\raisebox{0.25ex}{\accentmark{15}{\char"1CD2}}तीमा\raisebox{0.25ex}{\accentmark{15}{\char"1CD2}}र्क्त्या\raisebox{0.25ex}{\accentmark{12}{\char"0951}}ः\kern 0.05em ।\hskip 0.12em plus 0.04em minus 0.02em ऊगणावा\raisebox{0.15ex}{\accentmark{12}{\char"1CF9}}म\raisebox{0.25ex}{\accentmark{15}{\char"1CD2}}न्यामानास्तुरो\raisebox{0.25ex}{\accentmark{15}{\char"1CD2}}वा\kern 0.05em ।\hskip 0.12em plus 0.04em minus 0.02em क्षी\raisebox{0.25ex}{\accentmark{15}{\char"1CD2}}धीयूधा\raisebox{0.25ex}{\accentmark{15}{\char"1CD2}}शावा\raisebox{0.25ex}{\accentmark{12}{\char"0951}}सावा\raisebox{0.15ex}{\accentmark{12}{\char"1CF9}}ता\raisebox{0.25ex}{\accentmark{12}{\char"0951}}मिन्द्रा\raisebox{0.25ex}{\accentmark{15}{\char"1CD2}}ः\kern 0.05em ।\hskip 0.12em plus 0.04em minus 0.02em आ\raisebox{0.25ex}{\accentmark{15}{\char"1CD2}}भि\raisebox{0.25ex}{\accentmark{15}{\char"1CD2}}ष्या\raisebox{0.25ex}{\accentmark{12}{\char"0951}}मवृषामाणस्त्वोताः \nolinebreak\hskip 0.08em \mbox{॥५॥}}} \par}

\vspace{0.8em}
                % --- SUBSECTIONS (Mantras) ---
                
                % Conditional rendering based on output_mode
                \par\filbreak

\phantomsection

\addcontentsline{toc}{subsection}{ऋक् ६}

{\centering \footnotesize \textcolor{AccentPurple}{॥ वामदेवो गौतमः इन्द्रः त्रिष्टुप् ॥} \par}

\vspace{0.12em}

{\centering \textbf{\textcolor{AccentBlue}{यं\raisebox{0.25ex}{\accentmark{15}{\char"1CD2}}वृत्रे\raisebox{0.15ex}{\accentmark{12}{\char"1CF9}}षु\raisebox{0.25ex}{\accentmark{12}{\char"0951}}क्षीता\raisebox{0.25ex}{\accentmark{15}{\char"1CD2}}यास्प\raisebox{0.25ex}{\accentmark{15}{\char"1CD2}}र्धा\raisebox{0.25ex}{\accentmark{15}{\char"1CD2}}मानः\kern 0.05em ।\hskip 0.12em plus 0.04em minus 0.02em यं\raisebox{0.25ex}{\accentmark{15}{\char"1CD2}}युक्तेषु\raisebox{0.25ex}{\accentmark{12}{\char"0951}}तूरा\raisebox{0.15ex}{\accentmark{12}{\char"1CF9}}य॑\raisebox{0.25ex}{\accentmark{12}{\char"0951}}न्तोहा\raisebox{0.25ex}{\accentmark{15}{\char"1CD2}}व\raisebox{0.25ex}{\accentmark{12}{\char"0951}}न्ते\kern 0.05em ।\hskip 0.12em plus 0.04em minus 0.02em यं\raisebox{0.25ex}{\accentmark{15}{\char"1CD2}}शू\raisebox{0.25ex}{\accentmark{15}{\char"1CD2}}रा\raisebox{0.25ex}{\accentmark{12}{\char"0951}}सातौयामापा\raisebox{0.25ex}{\accentmark{12}{\char"0951}}मूपा\raisebox{0.25ex}{\accentmark{12}{\char"0951}}ज्मन्\kern 0.05em ।\hskip 0.12em plus 0.04em minus 0.02em यंवि\raisebox{0.25ex}{\accentmark{15}{\char"1CD2}}प्रा\raisebox{0.25ex}{\accentmark{12}{\char"0951}}सो\raisebox{0.25ex}{\accentmark{12}{\char"0951}}वाज\raisebox{0.15ex}{\accentmark{12}{\char"1CF9}}यन्तेसइन्द्रा\raisebox{0.25ex}{\accentmark{12}{\char"0951}}ः \nolinebreak\hskip 0.08em \mbox{॥६॥}}} \par}

\vspace{0.8em}
                % --- SUBSECTIONS (Mantras) ---
                
                % Conditional rendering based on output_mode
                \par\filbreak

\phantomsection

\addcontentsline{toc}{subsection}{ऋक् ७}

{\centering \footnotesize \textcolor{AccentPurple}{॥ विश्वामित्रो गाथिनः इन्द्रः (पर्वतान्द्रौ) त्रिष्टुप् ॥} \par}

\vspace{0.12em}

{\centering \textbf{\textcolor{AccentBlue}{इन्द्रा\raisebox{0.25ex}{\accentmark{12}{\char"0951}}पर्वताबृहा\raisebox{0.25ex}{\accentmark{15}{\char"1CD2}}तारा\raisebox{0.25ex}{\accentmark{15}{\char"1CD2}}थेना\kern 0.05em ।\hskip 0.12em plus 0.04em minus 0.02em वा\raisebox{0.25ex}{\accentmark{15}{\char"1CD2}}मी\raisebox{0.25ex}{\accentmark{15}{\char"1CD2}}रीळआ\raisebox{0.25ex}{\accentmark{15}{\char"1CD2}}वाहतंसूवी\raisebox{0.25ex}{\accentmark{15}{\char"1CD2}}राः\kern 0.05em ।\hskip 0.12em plus 0.04em minus 0.02em वी\raisebox{0.25ex}{\accentmark{15}{\char"1CD2}}तं\raisebox{0.25ex}{\accentmark{15}{\char"1CD2}}हाव्या\raisebox{0.15ex}{\accentmark{12}{\char"1CF9}}न्याध्वारेषुदेवा\raisebox{0.25ex}{\accentmark{15}{\char"1CD2}}ः\kern 0.05em ।\hskip 0.12em plus 0.04em minus 0.02em वर्धेथांगीर्भि\raisebox{0.15ex}{\accentmark{12}{\char"1CF9}}रिळा\raisebox{0.25ex}{\accentmark{12}{\char"0951}}यामाद\raisebox{0.25ex}{\accentmark{12}{\char"0951}}न्ताः \nolinebreak\hskip 0.08em \mbox{॥७॥}}} \par}

\vspace{0.8em}
                % --- SUBSECTIONS (Mantras) ---
                
                % Conditional rendering based on output_mode
                \par\filbreak

\phantomsection

\addcontentsline{toc}{subsection}{ऋक् ८}

{\centering \footnotesize \textcolor{AccentPurple}{॥ रेणुर्वैश्वामित्रः इन्द्रः त्रिष्टुप् ॥} \par}

\vspace{0.12em}

{\centering \textbf{\textcolor{AccentBlue}{इ\raisebox{0.15ex}{\accentmark{12}{\char"1CF8}}न्द्रा\raisebox{0.25ex}{\accentmark{12}{\char"0951}}यागीरोआनी\raisebox{0.25ex}{\accentmark{12}{\char"0951}}शीतसर्ग\raisebox{0.25ex}{\accentmark{12}{\char"0951}}ः\kern 0.05em ।\hskip 0.12em plus 0.04em minus 0.02em आ\raisebox{0.25ex}{\accentmark{15}{\char"1CD2}}पप्रै\raisebox{0.15ex}{\accentmark{12}{\char"1CF9}}रा\raisebox{0.25ex}{\accentmark{12}{\char"0951}}यत्सा\raisebox{0.25ex}{\accentmark{15}{\char"1CD2}}ग\raisebox{0.25ex}{\accentmark{12}{\char"0951}}रस्याबुध्ना\raisebox{0.25ex}{\accentmark{15}{\char"1CD2}}त्\kern 0.05em ।\hskip 0.12em plus 0.04em minus 0.02em यो\raisebox{0.25ex}{\accentmark{15}{\char"1CD2}}आ\raisebox{0.25ex}{\accentmark{15}{\char"1CD2}}क्षेणेवाचा\raisebox{0.25ex}{\accentmark{12}{\char"0951}}क्रीयौशा\raisebox{0.25ex}{\accentmark{15}{\char"1CD2}}ची\raisebox{0.25ex}{\accentmark{12}{\char"0951}}भिः\kern 0.05em ।\hskip 0.12em plus 0.04em minus 0.02em वि\raisebox{0.15ex}{\accentmark{12}{\char"1CF8}}ष्वात्तस्तंभा\raisebox{0.25ex}{\accentmark{12}{\char"0951}}पृथि\raisebox{0.25ex}{\accentmark{15}{\char"1CD2}}वी\raisebox{0.25ex}{\accentmark{15}{\char"1CD2}}मूतद्याम् \nolinebreak\hskip 0.08em \mbox{॥८॥}}} \par}

\vspace{0.8em}
                % --- SUBSECTIONS (Mantras) ---
                
                % Conditional rendering based on output_mode
                \par\filbreak

\phantomsection

\addcontentsline{toc}{subsection}{ऋक् ९}

{\centering \footnotesize \textcolor{AccentPurple}{॥ वामदेवो गौतमः (यमी वैवस्वती) इन्द्रः (यमो वैवस्वतः) त्रिष्टुप् ॥} \par}

\vspace{0.12em}

{\centering \textbf{\textcolor{AccentBlue}{आ\raisebox{0.25ex}{\accentmark{15}{\char"1CD2}}त्वासा\raisebox{0.25ex}{\accentmark{15}{\char"1CD2}}खा\raisebox{0.25ex}{\accentmark{12}{\char"0951}}यासख्यावा\raisebox{0.25ex}{\accentmark{12}{\char"0951}}वृत्यु\raisebox{0.25ex}{\accentmark{12}{\char"0951}}ः\kern 0.05em ।\hskip 0.12em plus 0.04em minus 0.02em ती\raisebox{0.15ex}{\accentmark{12}{\char"1CF8}}रा\raisebox{0.25ex}{\accentmark{15}{\char"1CD2}}ःपूरू\raisebox{0.25ex}{\accentmark{15}{\char"1CD2}}ची\raisebox{0.25ex}{\accentmark{12}{\char"0951}}दर्णावा\raisebox{0.25ex}{\accentmark{15}{\char"1CD2}}ञ्ज\raisebox{0.25ex}{\accentmark{12}{\char"0951}}गम्या\kern 0.05em ।\hskip 0.12em plus 0.04em minus 0.02em पि\raisebox{0.25ex}{\accentmark{15}{\char"1CD2}}तु\raisebox{0.15ex}{\accentmark{12}{\char"1CF9}}र्नापातामा\raisebox{0.25ex}{\accentmark{15}{\char"1CD2}}दधीतावेधाः\kern 0.05em ।\hskip 0.12em plus 0.04em minus 0.02em अस्मि\raisebox{0.25ex}{\accentmark{15}{\char"1CD2}}न्क्षायेप्रातारा\raisebox{0.25ex}{\accentmark{15}{\char"1CD2}}न्दीद्यानः \nolinebreak\hskip 0.08em \mbox{॥९॥}}} \par}

\vspace{0.8em}
                % --- SUBSECTIONS (Mantras) ---
                
                % Conditional rendering based on output_mode
                \par\filbreak

\phantomsection

\addcontentsline{toc}{subsection}{ऋक् १०}

{\centering \footnotesize \textcolor{AccentPurple}{॥ गौतमो राहूगणः इन्द्रः त्रिष्टुप् ॥} \par}

\vspace{0.12em}

{\centering \textbf{\textcolor{AccentBlue}{को\raisebox{0.25ex}{\accentmark{15}{\char"1CD2}}अद्यायुक्तेधुरिगा\raisebox{0.25ex}{\accentmark{15}{\char"1CD2}}ऋता\raisebox{0.25ex}{\accentmark{15}{\char"1CD2}}स्य\kern 0.05em ।\hskip 0.12em plus 0.04em minus 0.02em शी\raisebox{0.25ex}{\accentmark{15}{\char"1CD2}}मी\raisebox{0.25ex}{\accentmark{12}{\char"0951}}वतोभामीनो\raisebox{0.25ex}{\accentmark{12}{\char"0951}}दूर्ह्रणायुन्\kern 0.05em ।\hskip 0.12em plus 0.04em minus 0.02em आ\raisebox{0.25ex}{\accentmark{15}{\char"1CD2}}स\raisebox{0.25ex}{\accentmark{15}{\char"1CD2}}न्नेषामप्सूवा\raisebox{0.25ex}{\accentmark{15}{\char"1CD2}}होमायोभून्\kern 0.05em ।\hskip 0.12em plus 0.04em minus 0.02em या\raisebox{0.25ex}{\accentmark{15}{\char"1CD2}}ए\raisebox{0.25ex}{\accentmark{12}{\char"0951}}षांभृत्या\raisebox{0.25ex}{\accentmark{15}{\char"1CD2}}मृणा\raisebox{0.25ex}{\accentmark{15}{\char"1CD2}}धत्सा\raisebox{0.25ex}{\accentmark{15}{\char"1CD2}}जी\raisebox{0.25ex}{\accentmark{12}{\char"0951}}वान् \nolinebreak\hskip 0.08em \mbox{॥१०॥}}} \par}

\vspace{0.8em}

            % --- End of Kandah Closing Floral Line & Counts ---
            \par\nopagebreak\vspace{0.9em plus 0.3em minus 0.1em}
            \centerline{\textcolor{black!35}{\rule{1.8cm}{0.35pt}}\;\raisebox{-0.22ex}{\textcolor{vedicmaroon}{\large\SixFlowerPetalDotted}}\;\textcolor{black!35}{\rule{1.8cm}{0.35pt}}}
            \nopagebreak\vspace{0.45em}
            \centerline{\devafont\textbf{\textcolor{AccentBlue}{॥ ऋक् सङ्ख्या : १० ॥}}}
        
            \clearpage

            % --- SECTION ---
            
            \phantomsection  

            \section*{\centering चतुर्थ खण्डः}
            \addcontentsline{toc}{section}{चतुर्थ खण्डः} 
            \markright{चतुर्थ खण्डः}
            \nopagebreak % Keep section title with first mantra
            
            % --- SECTION-LEVEL PROCEDURE LINK ---

                % --- SUBSECTIONS (Mantras) ---
                
                % Conditional rendering based on output_mode
                \par\filbreak

\phantomsection

\addcontentsline{toc}{subsection}{ऋक् १}

{\centering \footnotesize \textcolor{AccentPurple}{॥ मधुच्छन्दा वैश्वामित्रः इन्द्रः अनुष्टुप् ॥} \par}

\vspace{0.12em}

{\centering \textbf{\textcolor{AccentBlue}{गा\raisebox{0.25ex}{\accentmark{15}{\char"1CD2}}यन्तित्वागायत्रीणा\raisebox{0.25ex}{\accentmark{12}{\char"0951}}ः\kern 0.05em ।\hskip 0.12em plus 0.04em minus 0.02em अ\raisebox{0.15ex}{\accentmark{12}{\char"1CF8}}र्चन्त्यर्का\raisebox{0.25ex}{\accentmark{15}{\char"1CD2}}मर्कीणाः\kern 0.05em ।\hskip 0.12em plus 0.04em minus 0.02em ब्र\raisebox{0.25ex}{\accentmark{15}{\char"1CD2}}ह्मा\raisebox{0.25ex}{\accentmark{15}{\char"1CD2}}णस्त्वाशतक्रतो\kern 0.05em ।\hskip 0.12em plus 0.04em minus 0.02em उ\raisebox{0.25ex}{\accentmark{15}{\char"1CD2}}द्वं\raisebox{0.25ex}{\accentmark{12}{\char"0951}}शामीवयेमिरे \nolinebreak\hskip 0.08em \mbox{॥१॥}}} \par}

\vspace{0.8em}
                % --- SUBSECTIONS (Mantras) ---
                
                % Conditional rendering based on output_mode
                \par\filbreak

\phantomsection

\addcontentsline{toc}{subsection}{ऋक् २}

{\centering \footnotesize \textcolor{AccentPurple}{॥ क्शेता माधुछन्दसः इन्द्रः अनुष्टुप् ॥} \par}

\vspace{0.12em}

{\centering \textbf{\textcolor{AccentBlue}{इ\raisebox{0.25ex}{\accentmark{15}{\char"1CD2}}न्द्रंविश्वा\raisebox{0.25ex}{\accentmark{12}{\char"0951}}अविवृधन्\kern 0.05em ।\hskip 0.12em plus 0.04em minus 0.02em समूद्रा\raisebox{0.25ex}{\accentmark{15}{\char"1CD2}}व्या\raisebox{0.25ex}{\accentmark{12}{\char"0951}}चासंगीरा\raisebox{0.25ex}{\accentmark{12}{\char"0951}}ः\kern 0.05em ।\hskip 0.12em plus 0.04em minus 0.02em रा\raisebox{0.25ex}{\accentmark{15}{\char"1CD2}}थीता\raisebox{0.25ex}{\accentmark{12}{\char"0951}}मंराथी\raisebox{0.25ex}{\accentmark{15}{\char"1CD2}}नाम्\raisebox{0.25ex}{\accentmark{12}{\char"0951}}\kern 0.05em ।\hskip 0.12em plus 0.04em minus 0.02em वाजानां\raisebox{0.25ex}{\accentmark{15}{\char"1CD2}}सात्पा\raisebox{0.25ex}{\accentmark{12}{\char"0951}}तिंपा\raisebox{0.25ex}{\accentmark{15}{\char"1CD2}}तिम् \nolinebreak\hskip 0.08em \mbox{॥२॥}}} \par}

\vspace{0.8em}
                % --- SUBSECTIONS (Mantras) ---
                
                % Conditional rendering based on output_mode
                \par\filbreak

\phantomsection

\addcontentsline{toc}{subsection}{ऋक् ३}

{\centering \footnotesize \textcolor{AccentPurple}{॥ गौतमो राहूगणः इन्द्रः अनुष्टुप् ॥} \par}

\vspace{0.12em}

{\centering \textbf{\textcolor{AccentBlue}{ईमा\raisebox{0.25ex}{\accentmark{15}{\char"1CD2}}मिन्द्रासूतंपी\raisebox{0.25ex}{\accentmark{15}{\char"1CD2}}बा\kern 0.05em ।\hskip 0.12em plus 0.04em minus 0.02em ज्ये\raisebox{0.25ex}{\accentmark{12}{\char"0951}}ष्ठमा\raisebox{0.15ex}{\accentmark{12}{\char"1CF9}}मार्त्य\raisebox{0.25ex}{\accentmark{12}{\char"0951}}म्मादम्\raisebox{0.25ex}{\accentmark{15}{\char"1CD2}}\kern 0.05em ।\hskip 0.12em plus 0.04em minus 0.02em शुक्र\raisebox{0.15ex}{\accentmark{12}{\char"1CF9}}स्या\raisebox{0.25ex}{\accentmark{12}{\char"0951}}त्वाभ्याठ\raisebox{0.25ex}{\accentmark{12}{\char"0951}}क्ष\raisebox{0.15ex}{\accentmark{12}{\char"1CF9}}रन्\kern 0.05em ।\hskip 0.12em plus 0.04em minus 0.02em धा\raisebox{0.15ex}{\accentmark{12}{\char"1CF9}}रा\raisebox{0.25ex}{\accentmark{12}{\char"0951}}ऋतस्यासाद\raisebox{0.25ex}{\accentmark{12}{\char"0951}}ने \nolinebreak\hskip 0.08em \mbox{॥३॥}}} \par}

\vspace{0.8em}
                % --- SUBSECTIONS (Mantras) ---
                
                % Conditional rendering based on output_mode
                \par\filbreak

\phantomsection

\addcontentsline{toc}{subsection}{ऋक् ४}

{\centering \footnotesize \textcolor{AccentPurple}{॥ अत्रिभर्गः इन्द्रः अनुष्टुप् ॥} \par}

\vspace{0.12em}

{\centering \textbf{\textcolor{AccentBlue}{या\raisebox{0.25ex}{\accentmark{15}{\char"1CD2}}दि\raisebox{0.25ex}{\accentmark{12}{\char"0951}}न्द्रचित्रमाह्याना\raisebox{0.25ex}{\accentmark{15}{\char"1CD2}}\kern 0.05em ।\hskip 0.12em plus 0.04em minus 0.02em अ\raisebox{0.25ex}{\accentmark{15}{\char"1CD2}}स्ति\raisebox{0.15ex}{\accentmark{12}{\char"1CF8}}त्वादा\raisebox{0.25ex}{\accentmark{12}{\char"0951}}तमद्रिवः\kern 0.05em ।\hskip 0.12em plus 0.04em minus 0.02em रा\raisebox{0.25ex}{\accentmark{15}{\char"1CD2}}धास्त\raisebox{0.25ex}{\accentmark{15}{\char"1CD2}}न्नोवि\raisebox{0.25ex}{\accentmark{15}{\char"1CD2}}दद्वसो\kern 0.05em ।\hskip 0.12em plus 0.04em minus 0.02em उ\raisebox{0.25ex}{\accentmark{15}{\char"1CD2}}भ\raisebox{0.25ex}{\accentmark{15}{\char"1CD2}}याहस्त्याभा\raisebox{0.25ex}{\accentmark{12}{\char"0951}}रा \nolinebreak\hskip 0.08em \mbox{॥४॥}}} \par}

\vspace{0.8em}
                % --- SUBSECTIONS (Mantras) ---
                
                % Conditional rendering based on output_mode
                \par\filbreak

\phantomsection

\addcontentsline{toc}{subsection}{ऋक् ५}

{\centering \footnotesize \textcolor{AccentPurple}{॥ तिरश्चीरांगिरसः इन्द्रः अनुष्टुप् ॥} \par}

\vspace{0.12em}

{\centering \textbf{\textcolor{AccentBlue}{श्रू\raisebox{0.25ex}{\accentmark{15}{\char"1CD2}}धीहा\raisebox{0.25ex}{\accentmark{15}{\char"1CD2}}व\raisebox{0.25ex}{\accentmark{12}{\char"0951}}न्तीरश्या\raisebox{0.25ex}{\accentmark{12}{\char"0951}}ः\kern 0.05em ।\hskip 0.12em plus 0.04em minus 0.02em इन्द्राय\raisebox{0.25ex}{\accentmark{15}{\char"1CD2}}स्त्वासापर्या\raisebox{0.25ex}{\accentmark{15}{\char"1CD2}}ति\kern 0.05em ।\hskip 0.12em plus 0.04em minus 0.02em सू\raisebox{0.25ex}{\accentmark{15}{\char"1CD2}}र्वी\raisebox{0.15ex}{\accentmark{12}{\char"1CF9}}र्या\raisebox{0.25ex}{\accentmark{12}{\char"0951}}स्यागोमा\raisebox{0.25ex}{\accentmark{12}{\char"0951}}तः\kern 0.05em ।\hskip 0.12em plus 0.04em minus 0.02em रा\raisebox{0.25ex}{\accentmark{15}{\char"1CD2}}या\raisebox{0.25ex}{\accentmark{15}{\char"1CD2}}स्पूर्धी\raisebox{0.25ex}{\accentmark{12}{\char"0951}}माहं\raisebox{0.25ex}{\accentmark{15}{\char"1CD2}}आ\raisebox{0.25ex}{\accentmark{12}{\char"0951}}सि \nolinebreak\hskip 0.08em \mbox{॥५॥}}} \par}

\vspace{0.8em}
                % --- SUBSECTIONS (Mantras) ---
                
                % Conditional rendering based on output_mode
                \par\filbreak

\phantomsection

\addcontentsline{toc}{subsection}{ऋक् ६}

{\centering \footnotesize \textcolor{AccentPurple}{॥ गौतमो राहूगणः इन्द्रः अनुष्टुप् ॥} \par}

\vspace{0.12em}

{\centering \textbf{\textcolor{AccentBlue}{आ\raisebox{0.15ex}{\accentmark{12}{\char"1CF8}}सा\raisebox{0.25ex}{\accentmark{12}{\char"0951}}वीसो\raisebox{0.25ex}{\accentmark{15}{\char"1CD2}}मा\raisebox{0.25ex}{\accentmark{15}{\char"1CD2}}इन्द्रते\kern 0.05em ।\hskip 0.12em plus 0.04em minus 0.02em शा\raisebox{0.25ex}{\accentmark{15}{\char"1CD2}}वी\raisebox{0.25ex}{\accentmark{12}{\char"0951}}ष्ठधृष्णा\raisebox{0.25ex}{\accentmark{15}{\char"1CD2}}वाग\raisebox{0.25ex}{\accentmark{12}{\char"0951}}हि\kern 0.05em ।\hskip 0.12em plus 0.04em minus 0.02em आत्वा॑\raisebox{0.25ex}{\accentmark{12}{\char"0951}}प्रणत्विन्द्रियम्\kern 0.05em ।\hskip 0.12em plus 0.04em minus 0.02em रा\raisebox{0.25ex}{\accentmark{15}{\char"1CD2}}जस्सू\raisebox{0.25ex}{\accentmark{15}{\char"1CD2}}र्यो\raisebox{0.25ex}{\accentmark{12}{\char"0951}}नारश्मी\raisebox{0.25ex}{\accentmark{15}{\char"1CD2}}भीः \nolinebreak\hskip 0.08em \mbox{॥६॥}}} \par}

\vspace{0.8em}
                % --- SUBSECTIONS (Mantras) ---
                
                % Conditional rendering based on output_mode
                \par\filbreak

\phantomsection

\addcontentsline{toc}{subsection}{ऋक् ७}

{\centering \footnotesize \textcolor{AccentPurple}{॥ नोपातिथिः काण्वः इन्द्रः अनुष्टुप् ॥} \par}

\vspace{0.12em}

{\centering \textbf{\textcolor{AccentBlue}{ए\raisebox{0.15ex}{\accentmark{12}{\char"1CF8}}न्द्रायाहिहारी\raisebox{0.25ex}{\accentmark{12}{\char"0951}}भिः\kern 0.05em ।\hskip 0.12em plus 0.04em minus 0.02em ऊ\raisebox{0.25ex}{\accentmark{15}{\char"1CD2}}पा\raisebox{0.25ex}{\accentmark{15}{\char"1CD2}}कण्वा\raisebox{0.25ex}{\accentmark{12}{\char"0951}}स्यसूष्ट्रतिम्\kern 0.05em ।\hskip 0.12em plus 0.04em minus 0.02em देवा\raisebox{0.25ex}{\accentmark{15}{\char"1CD2}}आमू\raisebox{0.25ex}{\accentmark{15}{\char"1CD2}}ष्याशासातः\kern 0.05em ।\hskip 0.12em plus 0.04em minus 0.02em दीवं\raisebox{0.25ex}{\accentmark{15}{\char"1CD2}}याया\raisebox{0.25ex}{\accentmark{15}{\char"1CD2}}दिवावसो \nolinebreak\hskip 0.08em \mbox{॥७॥}}} \par}

\vspace{0.8em}
                % --- SUBSECTIONS (Mantras) ---
                
                % Conditional rendering based on output_mode
                \par\filbreak

\phantomsection

\addcontentsline{toc}{subsection}{ऋक् ८}

{\centering \footnotesize \textcolor{AccentPurple}{॥ तिरश्चीरांगिरसः इन्द्रः अनुष्टुप् ॥} \par}

\vspace{0.12em}

{\centering \textbf{\textcolor{AccentBlue}{आत्वागी\raisebox{0.15ex}{\accentmark{12}{\char"1CF9}}रो\raisebox{0.25ex}{\accentmark{12}{\char"0951}}राथीरिव\kern 0.05em ।\hskip 0.12em plus 0.04em minus 0.02em अ\raisebox{0.15ex}{\accentmark{12}{\char"1CF8}}स्तु\raisebox{0.25ex}{\accentmark{12}{\char"0951}}स्सूते\raisebox{0.25ex}{\accentmark{15}{\char"1CD2}}षू\raisebox{0.25ex}{\accentmark{12}{\char"0951}}गिर्वणः\kern 0.05em ।\hskip 0.12em plus 0.04em minus 0.02em आभीत्वासामानूषत\kern 0.05em ।\hskip 0.12em plus 0.04em minus 0.02em गा\raisebox{0.15ex}{\accentmark{12}{\char"1CF8}}वो\raisebox{0.25ex}{\accentmark{12}{\char"0951}}वत्सन्ना\raisebox{0.25ex}{\accentmark{15}{\char"1CD2}}धेना\raisebox{0.25ex}{\accentmark{15}{\char"1CD2}}वाः \nolinebreak\hskip 0.08em \mbox{॥८॥}}} \par}

\vspace{0.8em}
                % --- SUBSECTIONS (Mantras) ---
                
                % Conditional rendering based on output_mode
                \par\filbreak

\phantomsection

\addcontentsline{toc}{subsection}{ऋक् ९}

{\centering \footnotesize \textcolor{AccentPurple}{॥ विश्वामित्रो गाथिनः इन्द्रः अनुष्टुप् ॥} \par}

\vspace{0.12em}

{\centering \textbf{\textcolor{AccentBlue}{ए\raisebox{0.25ex}{\accentmark{15}{\char"1CD2}}तोन्वी\raisebox{0.25ex}{\accentmark{12}{\char"0951}}न्द्रंस्ता\raisebox{0.25ex}{\accentmark{15}{\char"1CD2}}वामा\kern 0.05em ।\hskip 0.12em plus 0.04em minus 0.02em शुद्धं\raisebox{0.25ex}{\accentmark{15}{\char"1CD2}}शुद्धे\raisebox{0.25ex}{\accentmark{12}{\char"0951}}नासाम्ना\kern 0.05em ।\hskip 0.12em plus 0.04em minus 0.02em शुद्धै\raisebox{0.25ex}{\accentmark{12}{\char"0951}}रुक्थै\raisebox{0.15ex}{\accentmark{12}{\char"1CF9}}र्वावृध्वां\raisebox{0.25ex}{\accentmark{15}{\char"1CD2}}सम्\kern 0.05em ।\hskip 0.12em plus 0.04em minus 0.02em शुद्धैराशी\raisebox{0.25ex}{\accentmark{12}{\char"0951}}र्वान्ममक्तु \nolinebreak\hskip 0.08em \mbox{॥९॥}}} \par}

\vspace{0.8em}
                % --- SUBSECTIONS (Mantras) ---
                
                % Conditional rendering based on output_mode
                \par\filbreak

\phantomsection

\addcontentsline{toc}{subsection}{ऋक् १०}

{\centering \footnotesize \textcolor{AccentPurple}{॥ तिरश्चोरांगिरसः शय्युबार्हस्पत्यो वा इन्द्रः अनुष्टुप् ॥} \par}

\vspace{0.12em}

{\centering \textbf{\textcolor{AccentBlue}{योरा\raisebox{0.25ex}{\accentmark{15}{\char"1CD2}}यिं\raisebox{0.15ex}{\accentmark{12}{\char"1CF9}}वोरायिं\raisebox{0.25ex}{\accentmark{12}{\char"0951}}तामः\kern 0.05em ।\hskip 0.12em plus 0.04em minus 0.02em योद्यु\raisebox{0.25ex}{\accentmark{15}{\char"1CD2}}म्नैद्यु\raisebox{0.25ex}{\accentmark{15}{\char"1CD2}}म्ना\raisebox{0.25ex}{\accentmark{15}{\char"1CD2}}वाक्तमः\kern 0.05em ।\hskip 0.12em plus 0.04em minus 0.02em सोम\raisebox{0.25ex}{\accentmark{12}{\char"0951}}स्सूतस्सा\raisebox{0.25ex}{\accentmark{15}{\char"1CD2}}इ\raisebox{0.25ex}{\accentmark{12}{\char"0951}}न्द्रते\kern 0.05em ।\hskip 0.12em plus 0.04em minus 0.02em अस्ति\raisebox{0.25ex}{\accentmark{15}{\char"1CD2}}स्वाधा\raisebox{0.25ex}{\accentmark{12}{\char"0951}}पातेमा\raisebox{0.25ex}{\accentmark{15}{\char"1CD2}}दाः \nolinebreak\hskip 0.08em \mbox{॥१०॥}}} \par}

\vspace{0.8em}

            % --- End of Kandah Closing Floral Line & Counts ---
            \par\nopagebreak\vspace{0.9em plus 0.3em minus 0.1em}
            \centerline{\textcolor{black!35}{\rule{1.8cm}{0.35pt}}\;\raisebox{-0.22ex}{\textcolor{vedicmaroon}{\large\SixFlowerPetalDotted}}\;\textcolor{black!35}{\rule{1.8cm}{0.35pt}}}
            \nopagebreak\vspace{0.45em}
            \centerline{\devafont\textbf{\textcolor{AccentBlue}{॥ ऋक् सङ्ख्या : १० ॥}}}
        
            \clearpage

            % --- SECTION ---
            
            \phantomsection  

            \section*{\centering पञ्चम खण्डः}
            \addcontentsline{toc}{section}{पञ्चम खण्डः} 
            \markright{पञ्चम खण्डः}
            \nopagebreak % Keep section title with first mantra
            
            % --- SECTION-LEVEL PROCEDURE LINK ---

                % --- SUBSECTIONS (Mantras) ---
                
                % Conditional rendering based on output_mode
                \par\filbreak

\phantomsection

\addcontentsline{toc}{subsection}{ऋक् १}

{\centering \footnotesize \textcolor{AccentPurple}{॥ भरद्वाजो बार्हस्पत्यः इन्द्रः अनुष्टुप् ॥} \par}

\vspace{0.12em}

{\centering \textbf{\textcolor{AccentBlue}{प्र\raisebox{0.15ex}{\accentmark{12}{\char"1CF8}}त्य\raisebox{0.25ex}{\accentmark{12}{\char"0951}}स्मैपी\raisebox{0.25ex}{\accentmark{15}{\char"1CD2}}पी\raisebox{0.25ex}{\accentmark{12}{\char"0951}}षते\kern 0.05em ।\hskip 0.12em plus 0.04em minus 0.02em वि\raisebox{0.25ex}{\accentmark{15}{\char"1CD2}}श्वा\raisebox{0.25ex}{\accentmark{12}{\char"0951}}नीविदू\raisebox{0.25ex}{\accentmark{15}{\char"1CD2}}षे\raisebox{0.25ex}{\accentmark{12}{\char"0951}}भरा\kern 0.05em ।\hskip 0.12em plus 0.04em minus 0.02em अ\raisebox{0.15ex}{\accentmark{12}{\char"1CF8}}रं\raisebox{0.25ex}{\accentmark{12}{\char"0951}}गमा\raisebox{0.25ex}{\accentmark{15}{\char"1CD2}}याज\raisebox{0.25ex}{\accentmark{15}{\char"1CD2}}ग्मा\raisebox{0.25ex}{\accentmark{12}{\char"0951}}ये\kern 0.05em ।\hskip 0.12em plus 0.04em minus 0.02em आ\raisebox{0.25ex}{\accentmark{15}{\char"1CD2}}पा\raisebox{0.25ex}{\accentmark{12}{\char"0951}}श्चाद्दाघ्वाने\raisebox{0.25ex}{\accentmark{15}{\char"1CD2}}नाराः \nolinebreak\hskip 0.08em \mbox{॥१॥}}} \par}

\vspace{0.8em}
                % --- SUBSECTIONS (Mantras) ---
                
                % Conditional rendering based on output_mode
                \par\filbreak

\phantomsection

\addcontentsline{toc}{subsection}{ऋक् २}

{\centering \footnotesize \textcolor{AccentPurple}{॥ वामदेवो गौतमः, शाकपतूो वा इन्द्रः अनुष्टुप् ॥} \par}

\vspace{0.12em}

{\centering \textbf{\textcolor{AccentBlue}{आ\raisebox{0.25ex}{\accentmark{15}{\char"1CD2}}नोवयोवयाश्शायम्\raisebox{0.25ex}{\accentmark{15}{\char"1CD2}}\kern 0.05em ।\hskip 0.12em plus 0.04em minus 0.02em माहा\raisebox{0.25ex}{\accentmark{12}{\char"0951}}न्तंगंभारेष्ठाम्\kern 0.05em ।\hskip 0.12em plus 0.04em minus 0.02em माहा\raisebox{0.25ex}{\accentmark{15}{\char"1CD2}}न्तं\raisebox{0.25ex}{\accentmark{12}{\char"0951}}पूर्वी\raisebox{0.25ex}{\accentmark{12}{\char"0951}}नेष्टाम्\raisebox{0.25ex}{\accentmark{15}{\char"1CD2}}\kern 0.05em ।\hskip 0.12em plus 0.04em minus 0.02em उग्रं\raisebox{0.25ex}{\accentmark{15}{\char"1CD2}}वाचोआ\raisebox{0.25ex}{\accentmark{15}{\char"1CD2}}पा\raisebox{0.25ex}{\accentmark{12}{\char"0951}}वधीः \nolinebreak\hskip 0.08em \mbox{॥२॥}}} \par}

\vspace{0.8em}
                % --- SUBSECTIONS (Mantras) ---
                
                % Conditional rendering based on output_mode
                \par\filbreak

\phantomsection

\addcontentsline{toc}{subsection}{ऋक् ३}

{\centering \footnotesize \textcolor{AccentPurple}{॥ प्रियमेध आंगिरसः इन्द्रः अनुष्टुप् ॥} \par}

\vspace{0.12em}

{\centering \textbf{\textcolor{AccentBlue}{आ\raisebox{0.25ex}{\accentmark{15}{\char"1CD2}}त्वाराथं\raisebox{0.25ex}{\accentmark{12}{\char"0951}}याथोता\raisebox{0.25ex}{\accentmark{15}{\char"1CD2}}ये\kern 0.05em ।\hskip 0.12em plus 0.04em minus 0.02em सुम्नाया\raisebox{0.25ex}{\accentmark{12}{\char"0951}}वर्क्तया\raisebox{0.25ex}{\accentmark{12}{\char"0951}}मसि\kern 0.05em ।\hskip 0.12em plus 0.04em minus 0.02em तु\raisebox{0.25ex}{\accentmark{15}{\char"1CD2}}विकूर्मि\raisebox{0.25ex}{\accentmark{12}{\char"0951}}मृतीषाहम्\kern 0.05em ।\hskip 0.12em plus 0.04em minus 0.02em इ\raisebox{0.25ex}{\accentmark{15}{\char"1CD2}}न्द्रं\raisebox{0.25ex}{\accentmark{12}{\char"0951}}शवीष्ठासा\raisebox{0.25ex}{\accentmark{15}{\char"1CD2}}त्पा\raisebox{0.25ex}{\accentmark{12}{\char"0951}}तिम् \nolinebreak\hskip 0.08em \mbox{॥३॥}}} \par}

\vspace{0.8em}
                % --- SUBSECTIONS (Mantras) ---
                
                % Conditional rendering based on output_mode
                \par\filbreak

\phantomsection

\addcontentsline{toc}{subsection}{ऋक् ४}

{\centering \footnotesize \textcolor{AccentPurple}{॥ प्रगाथः काण्वः इन्द्रः अनुष्टुप् ॥} \par}

\vspace{0.12em}

{\centering \textbf{\textcolor{AccentBlue}{सा\raisebox{0.25ex}{\accentmark{15}{\char"1CD2}}पूर्व्यो\raisebox{0.25ex}{\accentmark{15}{\char"1CD2}}माहो\raisebox{0.25ex}{\accentmark{15}{\char"1CD2}}नाम्\kern 0.05em ।\hskip 0.12em plus 0.04em minus 0.02em वेनःक्रा\raisebox{0.25ex}{\accentmark{15}{\char"1CD2}}तूभिरानजे\kern 0.05em ।\hskip 0.12em plus 0.04em minus 0.02em यस्या\raisebox{0.25ex}{\accentmark{15}{\char"1CD2}}द्वारामा\raisebox{0.15ex}{\accentmark{12}{\char"1CF9}}नू\raisebox{0.25ex}{\accentmark{12}{\char"0951}}ःपीता\raisebox{0.25ex}{\accentmark{15}{\char"1CD2}}\kern 0.05em ।\hskip 0.12em plus 0.04em minus 0.02em देवे\raisebox{0.25ex}{\accentmark{15}{\char"1CD2}}षूधी\raisebox{0.25ex}{\accentmark{15}{\char"1CD2}}याआनाजे \nolinebreak\hskip 0.08em \mbox{॥४॥}}} \par}

\vspace{0.8em}
                % --- SUBSECTIONS (Mantras) ---
                
                % Conditional rendering based on output_mode
                \par\filbreak

\phantomsection

\addcontentsline{toc}{subsection}{ऋक् ५}

{\centering \footnotesize \textcolor{AccentPurple}{॥ श्यावाश्व आत्रयेः मरुतः अनुष्टुप् ॥} \par}

\vspace{0.12em}

{\centering \textbf{\textcolor{AccentBlue}{या\raisebox{0.25ex}{\accentmark{15}{\char"1CD2}}दीवा\raisebox{0.15ex}{\accentmark{12}{\char"1CF9}}ह\raisebox{0.25ex}{\accentmark{12}{\char"0951}}न्त्याशा\raisebox{0.25ex}{\accentmark{15}{\char"1CD2}}वाः\kern 0.05em ।\hskip 0.12em plus 0.04em minus 0.02em भ्राज\raisebox{0.25ex}{\accentmark{12}{\char"0951}}मानाराथेष्वा\raisebox{0.25ex}{\accentmark{15}{\char"1CD2}}\kern 0.05em ।\hskip 0.12em plus 0.04em minus 0.02em पीब\raisebox{0.25ex}{\accentmark{12}{\char"0951}}न्तोमादीरम्माधू\raisebox{0.25ex}{\accentmark{12}{\char"0951}}\kern 0.05em ।\hskip 0.12em plus 0.04em minus 0.02em तत्रश्रा\raisebox{0.25ex}{\accentmark{15}{\char"1CD2}}वां\raisebox{0.25ex}{\accentmark{12}{\char"0951}}सिकृण्वते \nolinebreak\hskip 0.08em \mbox{॥५॥}}} \par}

\vspace{0.8em}
                % --- SUBSECTIONS (Mantras) ---
                
                % Conditional rendering based on output_mode
                \par\filbreak

\phantomsection

\addcontentsline{toc}{subsection}{ऋक् ६}

{\centering \footnotesize \textcolor{AccentPurple}{॥ शय्युबार्हस्पत्यः इन्द्रः अनुष्टुप् ॥} \par}

\vspace{0.12em}

{\centering \textbf{\textcolor{AccentBlue}{त्या\raisebox{0.15ex}{\accentmark{12}{\char"1CF8}}मू\raisebox{0.25ex}{\accentmark{12}{\char"0951}}वोअप्रा\raisebox{0.25ex}{\accentmark{12}{\char"0951}}हणम्\kern 0.05em ।\hskip 0.12em plus 0.04em minus 0.02em गृणीषे\raisebox{0.15ex}{\accentmark{12}{\char"1CF9}}शा\raisebox{0.15ex}{\accentmark{12}{\char"1CF9}}वासस्पा\raisebox{0.25ex}{\accentmark{15}{\char"1CD2}}तिम्\kern 0.05em ।\hskip 0.12em plus 0.04em minus 0.02em इन्द्रं\raisebox{0.25ex}{\accentmark{12}{\char"0951}}विश्वासाहन्नारम्\kern 0.05em ।\hskip 0.12em plus 0.04em minus 0.02em शाची\raisebox{0.25ex}{\accentmark{12}{\char"0951}}ष्ठंविश्वा\raisebox{0.25ex}{\accentmark{12}{\char"0951}}वे\raisebox{0.25ex}{\accentmark{12}{\char"0951}}दसम् \nolinebreak\hskip 0.08em \mbox{॥६॥}}} \par}

\vspace{0.8em}
                % --- SUBSECTIONS (Mantras) ---
                
                % Conditional rendering based on output_mode
                \par\filbreak

\phantomsection

\addcontentsline{toc}{subsection}{ऋक् ७}

{\centering \footnotesize \textcolor{AccentPurple}{॥ वामदेवो गौतमः दधिक्रा वा अनुष्टुप् ॥} \par}

\vspace{0.12em}

{\centering \textbf{\textcolor{AccentBlue}{दधिक्रा\raisebox{0.25ex}{\accentmark{15}{\char"1CD2}}विण्णो\raisebox{0.25ex}{\accentmark{12}{\char"0951}}अकार्षम्\kern 0.05em ।\hskip 0.12em plus 0.04em minus 0.02em जिष्णोरश्वा\raisebox{0.25ex}{\accentmark{12}{\char"0951}}स्यावाजीना\raisebox{0.25ex}{\accentmark{12}{\char"0951}}ः\kern 0.05em ।\hskip 0.12em plus 0.04em minus 0.02em सूराभी\raisebox{0.25ex}{\accentmark{15}{\char"1CD2}}नोमू\raisebox{0.25ex}{\accentmark{15}{\char"1CD2}}खा\raisebox{0.25ex}{\accentmark{12}{\char"0951}}करात्\kern 0.05em ।\hskip 0.12em plus 0.04em minus 0.02em प्राणआयूंषितारिषात् \nolinebreak\hskip 0.08em \mbox{॥७॥}}} \par}

\vspace{0.8em}
                % --- SUBSECTIONS (Mantras) ---
                
                % Conditional rendering based on output_mode
                \par\filbreak

\phantomsection

\addcontentsline{toc}{subsection}{ऋक् ८}

{\centering \footnotesize \textcolor{AccentPurple}{॥ जेता माधुच्छन्दसः इन्द्रः अनुष्टुप् ॥} \par}

\vspace{0.12em}

{\centering \textbf{\textcolor{AccentBlue}{पूरांभिन्दू\raisebox{0.25ex}{\accentmark{15}{\char"1CD2}}रियूवा\raisebox{0.25ex}{\accentmark{15}{\char"1CD2}}ःकविः\kern 0.05em ।\hskip 0.12em plus 0.04em minus 0.02em आमी\raisebox{0.25ex}{\accentmark{12}{\char"0951}}तौजाअजायत\kern 0.05em ।\hskip 0.12em plus 0.04em minus 0.02em इ\raisebox{0.25ex}{\accentmark{15}{\char"1CD2}}न्द्रो\raisebox{0.15ex}{\accentmark{12}{\char"1CF9}}विश्वास्याक\raisebox{0.25ex}{\accentmark{15}{\char"1CD2}}र्मा\raisebox{0.25ex}{\accentmark{12}{\char"0951}}णः\kern 0.05em ।\hskip 0.12em plus 0.04em minus 0.02em धर्ता\raisebox{0.25ex}{\accentmark{12}{\char"0951}}वज्रीपू\raisebox{0.25ex}{\accentmark{12}{\char"0951}}रूष्ट्रतः \nolinebreak\hskip 0.08em \mbox{॥८॥}}} \par}

\vspace{0.8em}

            % --- End of Kandah Closing Floral Line & Counts ---
            \par\nopagebreak\vspace{0.9em plus 0.3em minus 0.1em}
            \centerline{\textcolor{black!35}{\rule{1.8cm}{0.35pt}}\;\raisebox{-0.22ex}{\textcolor{vedicmaroon}{\large\SixFlowerPetalDotted}}\;\textcolor{black!35}{\rule{1.8cm}{0.35pt}}}
            \nopagebreak\vspace{0.45em}
            \centerline{\devafont\textbf{\textcolor{AccentBlue}{॥ ऋक् सङ्ख्या : ८ ॥}}}
        
            \clearpage

            % --- SECTION ---
            
            \phantomsection  

            \section*{\centering षष्ठ खण्डः}
            \addcontentsline{toc}{section}{षष्ठ खण्डः} 
            \markright{षष्ठ खण्डः}
            \nopagebreak % Keep section title with first mantra
            
            % --- SECTION-LEVEL PROCEDURE LINK ---

                % --- SUBSECTIONS (Mantras) ---
                
                % Conditional rendering based on output_mode
                \par\filbreak

\phantomsection

\addcontentsline{toc}{subsection}{ऋक् १}

{\centering \footnotesize \textcolor{AccentPurple}{॥ प्रियमेध आंगिरसः इन्द्रः अनुष्टुप् ॥} \par}

\vspace{0.12em}

{\centering \textbf{\textcolor{AccentBlue}{प्रप्रा\raisebox{0.25ex}{\accentmark{12}{\char"0951}}वास्त्रष्टू\raisebox{0.25ex}{\accentmark{15}{\char"1CD2}}भामीषम्\kern 0.05em ।\hskip 0.12em plus 0.04em minus 0.02em वन्द\raisebox{0.15ex}{\accentmark{12}{\char"1CF9}}द्वी\raisebox{0.25ex}{\accentmark{12}{\char"0951}}राये\raisebox{0.25ex}{\accentmark{15}{\char"1CD2}}न्दा\raisebox{0.25ex}{\accentmark{12}{\char"0951}}वे\kern 0.05em ।\hskip 0.12em plus 0.04em minus 0.02em धी\raisebox{0.25ex}{\accentmark{15}{\char"1CD2}}या\raisebox{0.15ex}{\accentmark{12}{\char"1CF9}}वो\raisebox{0.25ex}{\accentmark{12}{\char"0951}}मेधा\raisebox{0.25ex}{\accentmark{15}{\char"1CD2}}सातये\kern 0.05em ।\hskip 0.12em plus 0.04em minus 0.02em पूरन्ध्या\raisebox{0.25ex}{\accentmark{12}{\char"0951}}वीवासती \nolinebreak\hskip 0.08em \mbox{॥१॥}}} \par}

\vspace{0.8em}
                % --- SUBSECTIONS (Mantras) ---
                
                % Conditional rendering based on output_mode
                \par\filbreak

\phantomsection

\addcontentsline{toc}{subsection}{ऋक् २}

{\centering \footnotesize \textcolor{AccentPurple}{॥ वामदेवो गौतमः इन्द्रः अनुष्टुप् ॥} \par}

\vspace{0.12em}

{\centering \textbf{\textcolor{AccentBlue}{कश्यापा\raisebox{0.25ex}{\accentmark{12}{\char"0951}}स्यास्वर्वीदाः\kern 0.05em ।\hskip 0.12em plus 0.04em minus 0.02em या\raisebox{0.25ex}{\accentmark{15}{\char"1CD2}}आहुस्सायूजावी\raisebox{0.25ex}{\accentmark{12}{\char"0951}}ति\kern 0.05em ।\hskip 0.12em plus 0.04em minus 0.02em या\raisebox{0.25ex}{\accentmark{15}{\char"1CD2}}योर्वीश्वामा\raisebox{0.15ex}{\accentmark{12}{\char"1CF9}}पी\raisebox{0.25ex}{\accentmark{12}{\char"0951}}व्राते\kern 0.05em ।\hskip 0.12em plus 0.04em minus 0.02em यज्ञ\raisebox{0.15ex}{\accentmark{12}{\char"1CF9}}न्धी\raisebox{0.15ex}{\accentmark{12}{\char"1CF9}}रो\raisebox{0.25ex}{\accentmark{12}{\char"0951}}नीचा\raisebox{0.25ex}{\accentmark{15}{\char"1CD2}}य्या \nolinebreak\hskip 0.08em \mbox{॥२॥}}} \par}

\vspace{0.8em}
                % --- SUBSECTIONS (Mantras) ---
                
                % Conditional rendering based on output_mode
                \par\filbreak

\phantomsection

\addcontentsline{toc}{subsection}{ऋक् ३}

{\centering \footnotesize \textcolor{AccentPurple}{॥ प्रियमेध आंगिरसः इन्द्रः अनुष्टुप् ॥} \par}

\vspace{0.12em}

{\centering \textbf{\textcolor{AccentBlue}{विश्वा\raisebox{0.25ex}{\accentmark{15}{\char"1CD2}}नरा\raisebox{0.25ex}{\accentmark{12}{\char"0951}}स्यवास्पातिम्\kern 0.05em ।\hskip 0.12em plus 0.04em minus 0.02em आनानतास्याशा\raisebox{0.25ex}{\accentmark{15}{\char"1CD2}}वासः\kern 0.05em ।\hskip 0.12em plus 0.04em minus 0.02em एवै\raisebox{0.25ex}{\accentmark{12}{\char"0951}}श्वचर्षा\raisebox{0.25ex}{\accentmark{12}{\char"0951}}णीनाम्\raisebox{0.25ex}{\accentmark{15}{\char"1CD2}}\kern 0.05em ।\hskip 0.12em plus 0.04em minus 0.02em ऊतीहू\raisebox{0.25ex}{\accentmark{12}{\char"0951}}वेराथानाम् \nolinebreak\hskip 0.08em \mbox{॥३॥}}} \par}

\vspace{0.8em}
                % --- SUBSECTIONS (Mantras) ---
                
                % Conditional rendering based on output_mode
                \par\filbreak

\phantomsection

\addcontentsline{toc}{subsection}{ऋक् ४}

{\centering \footnotesize \textcolor{AccentPurple}{॥ मधुच्छन्दा विश्वामित्रः इन्द्रः अनुष्टुप् ॥} \par}

\vspace{0.12em}

{\centering \textbf{\textcolor{AccentBlue}{सा\raisebox{0.25ex}{\accentmark{15}{\char"1CD2}}खाय\raisebox{0.15ex}{\accentmark{12}{\char"1CF9}}स्ते\raisebox{0.25ex}{\accentmark{12}{\char"0951}}दीवोना\raisebox{0.25ex}{\accentmark{12}{\char"0951}}राः\kern 0.05em ।\hskip 0.12em plus 0.04em minus 0.02em धीया\raisebox{0.15ex}{\accentmark{12}{\char"1CF9}}मर्तास्याश\raisebox{0.25ex}{\accentmark{15}{\char"1CD2}}र्मा\raisebox{0.25ex}{\accentmark{12}{\char"0951}}तः\kern 0.05em ।\hskip 0.12em plus 0.04em minus 0.02em ऊ\raisebox{0.25ex}{\accentmark{15}{\char"1CD2}}तीसा\raisebox{0.15ex}{\accentmark{12}{\char"1CF9}}बृ\raisebox{0.15ex}{\accentmark{12}{\char"1CF9}}हातो\raisebox{0.25ex}{\accentmark{15}{\char"1CD2}}दीवः\kern 0.05em ।\hskip 0.12em plus 0.04em minus 0.02em द्वीषो\raisebox{0.25ex}{\accentmark{15}{\char"1CD2}}अंहोना\raisebox{0.25ex}{\accentmark{15}{\char"1CD2}}ता\raisebox{0.25ex}{\accentmark{12}{\char"0951}}रती \nolinebreak\hskip 0.08em \mbox{॥४॥}}} \par}

\vspace{0.8em}
                % --- SUBSECTIONS (Mantras) ---
                
                % Conditional rendering based on output_mode
                \par\filbreak

\phantomsection

\addcontentsline{toc}{subsection}{ऋक् ५}

{\centering \footnotesize \textcolor{AccentPurple}{॥ प्रियमेध आंगिरसः इन्द्रः अनुष्टुप् ॥} \par}

\vspace{0.12em}

{\centering \textbf{\textcolor{AccentBlue}{अ\raisebox{0.25ex}{\accentmark{15}{\char"1CD2}}र्चातप्रा\raisebox{0.25ex}{\accentmark{15}{\char"1CD2}}र्चातानरः\kern 0.05em ।\hskip 0.12em plus 0.04em minus 0.02em प्री\raisebox{0.25ex}{\accentmark{15}{\char"1CD2}}यामे\raisebox{0.25ex}{\accentmark{12}{\char"0951}}धासोअ\raisebox{0.25ex}{\accentmark{15}{\char"1CD2}}र्चातः\kern 0.05em ।\hskip 0.12em plus 0.04em minus 0.02em अर्चन्तूपूत्राका\raisebox{0.25ex}{\accentmark{15}{\char"1CD2}}ऊता\kern 0.05em ।\hskip 0.12em plus 0.04em minus 0.02em पू\raisebox{0.25ex}{\accentmark{15}{\char"1CD2}}रामिधृष्णुंठ\raisebox{0.25ex}{\accentmark{12}{\char"0951}}वार्चत \nolinebreak\hskip 0.08em \mbox{॥५॥}}} \par}

\vspace{0.8em}
                % --- SUBSECTIONS (Mantras) ---
                
                % Conditional rendering based on output_mode
                \par\filbreak

\phantomsection

\addcontentsline{toc}{subsection}{ऋक् ६}

{\centering \footnotesize \textcolor{AccentPurple}{॥ भरद्वाजो बार्हस्पथः इन्द्रः (अग्निः) अनुष्टुप् ॥} \par}

\vspace{0.12em}

{\centering \textbf{\textcolor{AccentBlue}{उक्थ\raisebox{0.15ex}{\accentmark{12}{\char"1CF9}}मिन्द्रा\raisebox{0.25ex}{\accentmark{12}{\char"0951}}याशं\raisebox{0.25ex}{\accentmark{15}{\char"1CD2}}स्यम्\raisebox{0.25ex}{\accentmark{12}{\char"0951}}\kern 0.05em ।\hskip 0.12em plus 0.04em minus 0.02em व\raisebox{0.25ex}{\accentmark{15}{\char"1CD2}}र्धानंपूरूनिष्षिधे\kern 0.05em ।\hskip 0.12em plus 0.04em minus 0.02em शक्रो\raisebox{0.15ex}{\accentmark{12}{\char"1CF9}}या\raisebox{0.15ex}{\accentmark{12}{\char"1CF9}}था\raisebox{0.25ex}{\accentmark{12}{\char"0951}}सूते\raisebox{0.25ex}{\accentmark{12}{\char"0951}}षू\raisebox{0.25ex}{\accentmark{12}{\char"0951}}णः\kern 0.05em ।\hskip 0.12em plus 0.04em minus 0.02em रारा\raisebox{0.15ex}{\accentmark{12}{\char"1CF9}}णात्सख्ये\raisebox{0.25ex}{\accentmark{15}{\char"1CD2}}षु\raisebox{0.25ex}{\accentmark{12}{\char"0951}}च \nolinebreak\hskip 0.08em \mbox{॥६॥}}} \par}

\vspace{0.8em}
                % --- SUBSECTIONS (Mantras) ---
                
                % Conditional rendering based on output_mode
                \par\filbreak

\phantomsection

\addcontentsline{toc}{subsection}{ऋक् ७}

{\centering \footnotesize \textcolor{AccentPurple}{॥ अत्रिभूमः इन्द्रः अनुष्टुप् ॥} \par}

\vspace{0.12em}

{\centering \textbf{\textcolor{AccentBlue}{विभोष्टइन्द्ररा\raisebox{0.25ex}{\accentmark{15}{\char"1CD2}}धा\raisebox{0.25ex}{\accentmark{12}{\char"0951}}सः\kern 0.05em ।\hskip 0.12em plus 0.04em minus 0.02em विभ्वीरातीश्शातक्रतो\kern 0.05em ।\hskip 0.12em plus 0.04em minus 0.02em अ\raisebox{0.25ex}{\accentmark{15}{\char"1CD2}}था\raisebox{0.25ex}{\accentmark{12}{\char"0951}}नोविश्वचर्ष\raisebox{0.25ex}{\accentmark{12}{\char"0951}}णे\kern 0.05em ।\hskip 0.12em plus 0.04em minus 0.02em द्युम्नं\raisebox{0.25ex}{\accentmark{15}{\char"1CD2}}सूदत्रमंहया \nolinebreak\hskip 0.08em \mbox{॥७॥}}} \par}

\vspace{0.8em}
                % --- SUBSECTIONS (Mantras) ---
                
                % Conditional rendering based on output_mode
                \par\filbreak

\phantomsection

\addcontentsline{toc}{subsection}{ऋक् ८}

{\centering \footnotesize \textcolor{AccentPurple}{॥ प्रस्कण्वः काण्वः उषाः अनुष्टुप् ॥} \par}

\vspace{0.12em}

{\centering \textbf{\textcolor{AccentBlue}{वा\raisebox{0.25ex}{\accentmark{15}{\char"1CD2}}या\raisebox{0.25ex}{\accentmark{12}{\char"0951}}श्चीक्तेपा\raisebox{0.25ex}{\accentmark{15}{\char"1CD2}}तत्रीणा\raisebox{0.25ex}{\accentmark{12}{\char"0951}}ः\kern 0.05em ।\hskip 0.12em plus 0.04em minus 0.02em द्वीपाश्चा\raisebox{0.25ex}{\accentmark{15}{\char"1CD2}}तू\raisebox{0.25ex}{\accentmark{12}{\char"0951}}ष्पदर्जूनी\kern 0.05em ।\hskip 0.12em plus 0.04em minus 0.02em ऊषप्रा\raisebox{0.15ex}{\accentmark{12}{\char"1CF9}}र\raisebox{0.25ex}{\accentmark{12}{\char"0951}}न्मृतूंरा\raisebox{0.25ex}{\accentmark{15}{\char"1CD2}}नू\raisebox{0.25ex}{\accentmark{12}{\char"0951}}\kern 0.05em ।\hskip 0.12em plus 0.04em minus 0.02em दीवो\raisebox{0.15ex}{\accentmark{12}{\char"1CF9}}अ\raisebox{0.25ex}{\accentmark{12}{\char"0951}}न्तभ्यस्पारी\raisebox{0.25ex}{\accentmark{12}{\char"0951}} \nolinebreak\hskip 0.08em \mbox{॥८॥}}} \par}

\vspace{0.8em}
                % --- SUBSECTIONS (Mantras) ---
                
                % Conditional rendering based on output_mode
                \par\filbreak

\phantomsection

\addcontentsline{toc}{subsection}{ऋक् ९}

{\centering \footnotesize \textcolor{AccentPurple}{॥ त्रित आत्यः (कुत्स आंगिरसो वा) विश्वेदेवाः अनुष्टुप् ॥} \par}

\vspace{0.12em}

{\centering \textbf{\textcolor{AccentBlue}{आ\raisebox{0.25ex}{\accentmark{15}{\char"1CD2}}मीये\raisebox{0.25ex}{\accentmark{15}{\char"1CD2}}देवा\raisebox{0.25ex}{\accentmark{15}{\char"1CD2}}स्था\raisebox{0.25ex}{\accentmark{15}{\char"1CD2}}ना\raisebox{0.25ex}{\accentmark{12}{\char"0951}}\kern 0.05em ।\hskip 0.12em plus 0.04em minus 0.02em मध्येये\raisebox{0.15ex}{\accentmark{12}{\char"1CF9}}रो\raisebox{0.15ex}{\accentmark{12}{\char"1CF9}}चा\raisebox{0.25ex}{\accentmark{12}{\char"0951}}नेदीवः\kern 0.05em ।\hskip 0.12em plus 0.04em minus 0.02em कद्व\raisebox{0.15ex}{\accentmark{12}{\char"1CF8}}ऋ\raisebox{0.25ex}{\accentmark{12}{\char"0951}}तंकादमृ\raisebox{0.25ex}{\accentmark{15}{\char"1CD2}}तम्\kern 0.05em ।\hskip 0.12em plus 0.04em minus 0.02em का\raisebox{0.25ex}{\accentmark{15}{\char"1CD2}}प्रत्ना\raisebox{0.25ex}{\accentmark{15}{\char"1CD2}}आहू\raisebox{0.25ex}{\accentmark{12}{\char"0951}}तिः \nolinebreak\hskip 0.08em \mbox{॥९॥}}} \par}

\vspace{0.8em}
                % --- SUBSECTIONS (Mantras) ---
                
                % Conditional rendering based on output_mode
                \par\filbreak

\phantomsection

\addcontentsline{toc}{subsection}{ऋक् १०}

{\centering \footnotesize \textcolor{AccentPurple}{॥ वामदेवो गौतमः इन्द्रः अनुष्टुप् ॥} \par}

\vspace{0.12em}

{\centering \textbf{\textcolor{AccentBlue}{ऋ\raisebox{0.25ex}{\accentmark{15}{\char"1CD2}}चं\raisebox{0.25ex}{\accentmark{15}{\char"1CD2}}सा\raisebox{0.25ex}{\accentmark{15}{\char"1CD2}}मायजामहे\kern 0.05em ।\hskip 0.12em plus 0.04em minus 0.02em याम्यांकर्माणीकृ\raisebox{0.25ex}{\accentmark{12}{\char"0951}}ण्वा\raisebox{0.25ex}{\accentmark{15}{\char"1CD2}}ते\kern 0.05em ।\hskip 0.12em plus 0.04em minus 0.02em वी\raisebox{0.25ex}{\accentmark{15}{\char"1CD2}}तेसाद\raisebox{0.25ex}{\accentmark{12}{\char"0951}}सिराजतः\kern 0.05em ।\hskip 0.12em plus 0.04em minus 0.02em य\raisebox{0.25ex}{\accentmark{15}{\char"1CD2}}ज्ञं\raisebox{0.25ex}{\accentmark{15}{\char"1CD2}}देवे\raisebox{0.25ex}{\accentmark{15}{\char"1CD2}}षु\raisebox{0.25ex}{\accentmark{12}{\char"0951}}वक्ष\raisebox{0.25ex}{\accentmark{12}{\char"0951}}तः \nolinebreak\hskip 0.08em \mbox{॥१०॥}}} \par}

\vspace{0.8em}

            % --- End of Kandah Closing Floral Line & Counts ---
            \par\nopagebreak\vspace{0.9em plus 0.3em minus 0.1em}
            \centerline{\textcolor{black!35}{\rule{1.8cm}{0.35pt}}\;\raisebox{-0.22ex}{\textcolor{vedicmaroon}{\large\SixFlowerPetalDotted}}\;\textcolor{black!35}{\rule{1.8cm}{0.35pt}}}
            \nopagebreak\vspace{0.45em}
            \centerline{\devafont\textbf{\textcolor{AccentBlue}{॥ ऋक् सङ्ख्या : १० ॥}}}


    % --- SUPERSECTION ---
    \clearpage           
    \phantomsection              
    \chapter*{\centering ऐन्द्रपाठः}
    \addcontentsline{toc}{chapter}{ऐन्द्रपाठः}
    \markboth{ऐन्द्रपाठः}{ऐन्द्रपाठः}
    
    % GLUE FIX: Prevent page break between Chapter and First Section
    \nopagebreak

        

            % --- SECTION ---
            
            \phantomsection  

            \section*{\centering प्रथम खण्डः}
            \addcontentsline{toc}{section}{प्रथम खण्डः} 
            \markright{प्रथम खण्डः}
            \nopagebreak % Keep section title with first mantra
            
            % --- SECTION-LEVEL PROCEDURE LINK ---

                % --- SUBSECTIONS (Mantras) ---
                
                % Conditional rendering based on output_mode
                \par\filbreak

\phantomsection

\addcontentsline{toc}{subsection}{ऋक् १}

{\centering \footnotesize \textcolor{AccentPurple}{॥ रेभः कश्यपः इन्द्रः अति जगती ॥} \par}

\vspace{0.12em}

{\centering \textbf{\textcolor{AccentBlue}{विश्वा\raisebox{0.25ex}{\accentmark{15}{\char"1CD2}}ःपृ\raisebox{0.25ex}{\accentmark{15}{\char"1CD2}}ता\raisebox{0.25ex}{\accentmark{12}{\char"0951}}ना\raisebox{0.25ex}{\accentmark{15}{\char"1CD2}}आभीभू\raisebox{0.15ex}{\accentmark{12}{\char"1CF9}}ता\raisebox{0.25ex}{\accentmark{12}{\char"0951}}रन्नाराः\kern 0.05em ।\hskip 0.12em plus 0.04em minus 0.02em साजू\raisebox{0.25ex}{\accentmark{15}{\char"1CD2}}स्ताता\raisebox{0.25ex}{\accentmark{12}{\char"0951}}क्षुरिन्द्रं\raisebox{0.25ex}{\accentmark{12}{\char"0951}}जजनु\raisebox{0.15ex}{\accentmark{12}{\char"1CF9}}श्चा\raisebox{0.25ex}{\accentmark{12}{\char"0951}}राजसे\kern 0.05em ।\hskip 0.12em plus 0.04em minus 0.02em क्र\raisebox{0.25ex}{\accentmark{15}{\char"1CD2}}त्वेवा\raisebox{0.15ex}{\accentmark{12}{\char"1CF9}}रेस्थेमन्यामू\raisebox{0.15ex}{\accentmark{12}{\char"1CF9}}री\raisebox{0.25ex}{\accentmark{12}{\char"0951}}मू\raisebox{0.25ex}{\accentmark{15}{\char"1CD2}}ता\kern 0.05em ।\hskip 0.12em plus 0.04em minus 0.02em उ\raisebox{0.25ex}{\accentmark{15}{\char"1CD2}}ग्र\raisebox{0.25ex}{\accentmark{15}{\char"1CD2}}मो\raisebox{0.25ex}{\accentmark{15}{\char"1CD2}}जी\raisebox{0.25ex}{\accentmark{12}{\char"0951}}ष्ठन्तारा\raisebox{0.25ex}{\accentmark{15}{\char"1CD2}}स\raisebox{0.25ex}{\accentmark{12}{\char"0951}}न्तारस्वीनम् \nolinebreak\hskip 0.08em \mbox{॥१॥}}} \par}

\vspace{0.8em}
                % --- SUBSECTIONS (Mantras) ---
                
                % Conditional rendering based on output_mode
                \par\filbreak

\phantomsection

\addcontentsline{toc}{subsection}{ऋक् २}

{\centering \footnotesize \textcolor{AccentPurple}{॥ सुवेदाः शैलूषिः इन्द्रः जगती ॥} \par}

\vspace{0.12em}

{\centering \textbf{\textcolor{AccentBlue}{श्र\raisebox{0.25ex}{\accentmark{15}{\char"1CD2}}क्ते\raisebox{0.25ex}{\accentmark{12}{\char"0951}}दधामिप्राथामा\raisebox{0.15ex}{\accentmark{12}{\char"1CF9}}या\raisebox{0.25ex}{\accentmark{12}{\char"0951}}मन्या\raisebox{0.25ex}{\accentmark{15}{\char"1CD2}}वे\kern 0.05em ।\hskip 0.12em plus 0.04em minus 0.02em आ\raisebox{0.25ex}{\accentmark{15}{\char"1CD2}}ह\raisebox{0.25ex}{\accentmark{15}{\char"1CD2}}न्यद्द\raisebox{0.25ex}{\accentmark{15}{\char"1CD2}}स्युन्न\raisebox{0.15ex}{\accentmark{12}{\char"1CF9}}र्यं\raisebox{0.25ex}{\accentmark{12}{\char"0951}}वीवे\raisebox{0.25ex}{\accentmark{15}{\char"1CD2}}रापः\kern 0.05em ।\hskip 0.12em plus 0.04em minus 0.02em ऊ\raisebox{0.25ex}{\accentmark{15}{\char"1CD2}}भेयत्वा\raisebox{0.25ex}{\accentmark{15}{\char"1CD2}}रोद\raisebox{0.25ex}{\accentmark{12}{\char"0951}}सीधावा\raisebox{0.25ex}{\accentmark{12}{\char"0951}}तामा\raisebox{0.25ex}{\accentmark{15}{\char"1CD2}}नू\raisebox{0.25ex}{\accentmark{12}{\char"0951}}\kern 0.05em ।\hskip 0.12em plus 0.04em minus 0.02em भ्या\raisebox{0.15ex}{\accentmark{12}{\char"1CF8}}सा\raisebox{0.25ex}{\accentmark{12}{\char"0951}}क्तेशू\raisebox{0.25ex}{\accentmark{15}{\char"1CD2}}ष्मात्पृथि\raisebox{0.25ex}{\accentmark{15}{\char"1CD2}}वी\raisebox{0.25ex}{\accentmark{12}{\char"0951}}ची\raisebox{0.25ex}{\accentmark{12}{\char"0951}}दद्रीवः \nolinebreak\hskip 0.08em \mbox{॥२॥}}} \par}

\vspace{0.8em}
                % --- SUBSECTIONS (Mantras) ---
                
                % Conditional rendering based on output_mode
                \par\filbreak

\phantomsection

\addcontentsline{toc}{subsection}{ऋक् ३}

{\centering \footnotesize \textcolor{AccentPurple}{॥ वामदेवो गौतमः इन्द्रः जगती ॥} \par}

\vspace{0.12em}

{\centering \textbf{\textcolor{AccentBlue}{सामेताविश्वा\raisebox{0.15ex}{\accentmark{12}{\char"1CF9}}ओज\raisebox{0.25ex}{\accentmark{12}{\char"0951}}सापा\raisebox{0.15ex}{\accentmark{12}{\char"1CF9}}तिन्दीव\raisebox{0.25ex}{\accentmark{15}{\char"1CD2}}ः\kern 0.05em ।\hskip 0.12em plus 0.04em minus 0.02em य\raisebox{0.25ex}{\accentmark{15}{\char"1CD2}}ए\raisebox{0.25ex}{\accentmark{15}{\char"1CD2}}का\raisebox{0.15ex}{\accentmark{12}{\char"1CF9}}इ\raisebox{0.15ex}{\accentmark{12}{\char"1CF8}}त्भू\raisebox{0.25ex}{\accentmark{12}{\char"0951}}रा\raisebox{0.25ex}{\accentmark{15}{\char"1CD2}}तीथिर्जना\raisebox{0.25ex}{\accentmark{12}{\char"0951}}नाम्\kern 0.05em ।\hskip 0.12em plus 0.04em minus 0.02em सा\raisebox{0.25ex}{\accentmark{15}{\char"1CD2}}पूर्व्यो\raisebox{0.25ex}{\accentmark{15}{\char"1CD2}}नूता\raisebox{0.25ex}{\accentmark{12}{\char"0951}}नामाजी\raisebox{0.25ex}{\accentmark{15}{\char"1CD2}}गीष\raisebox{0.25ex}{\accentmark{12}{\char"0951}}न्तंवार्क्तानीरानूवावृताएकाईत् \nolinebreak\hskip 0.08em \mbox{॥३॥}}} \par}

\vspace{0.8em}
                % --- SUBSECTIONS (Mantras) ---
                
                % Conditional rendering based on output_mode
                \par\filbreak

\phantomsection

\addcontentsline{toc}{subsection}{ऋक् ४}

{\centering \footnotesize \textcolor{AccentPurple}{॥ सव्य आंगिरसः इन्द्रः जगती ॥} \par}

\vspace{0.12em}

{\centering \textbf{\textcolor{AccentBlue}{ईमेताइन्द्रातेवा\raisebox{0.25ex}{\accentmark{12}{\char"0951}}यंपूरूष्टुत\kern 0.05em ।\hskip 0.12em plus 0.04em minus 0.02em येत्वाराभ्याचारामसिप्रभूवसो\kern 0.05em ।\hskip 0.12em plus 0.04em minus 0.02em नाहित्वादन्योगीर्वाणोगीरस्साघात्\kern 0.05em ।\hskip 0.12em plus 0.04em minus 0.02em क्षोणीरीवप्रातीताद्धरीयानो\raisebox{0.25ex}{\accentmark{15}{\char"1CD2}}वा\raisebox{0.25ex}{\accentmark{15}{\char"1CD2}}चा\raisebox{0.25ex}{\accentmark{12}{\char"0951}}ः \nolinebreak\hskip 0.08em \mbox{॥४॥}}} \par}

\vspace{0.8em}
                % --- SUBSECTIONS (Mantras) ---
                
                % Conditional rendering based on output_mode
                \par\filbreak

\phantomsection

\addcontentsline{toc}{subsection}{ऋक् ५}

{\centering \footnotesize \textcolor{AccentPurple}{॥ विश्वामित्रो गाथिनः इन्द्रः जगती ॥} \par}

\vspace{0.12em}

{\centering \textbf{\textcolor{AccentBlue}{चऋष\raisebox{0.25ex}{\accentmark{12}{\char"0951}}णीधृ\raisebox{0.15ex}{\accentmark{12}{\char"1CF9}}तंम्माघावानाठ\raisebox{0.25ex}{\accentmark{12}{\char"0951}}मुक्त्थ्यम्\kern 0.05em ।\hskip 0.12em plus 0.04em minus 0.02em इ\raisebox{0.25ex}{\accentmark{15}{\char"1CD2}}न्द्रंगीरोबृहातीरभ्यानूषत\kern 0.05em ।\hskip 0.12em plus 0.04em minus 0.02em वा\raisebox{0.25ex}{\accentmark{15}{\char"1CD2}}वृधा\raisebox{0.25ex}{\accentmark{15}{\char"1CD2}}नंपू\raisebox{0.25ex}{\accentmark{12}{\char"0951}}रूहूतं\raisebox{0.15ex}{\accentmark{12}{\char"1CF9}}सूवृक्ती\raisebox{0.25ex}{\accentmark{12}{\char"0951}}भीः\kern 0.05em ।\hskip 0.12em plus 0.04em minus 0.02em आ\raisebox{0.15ex}{\accentmark{12}{\char"1CF9}}मा\raisebox{0.25ex}{\accentmark{12}{\char"0951}}र्क्त्यंजरा\raisebox{0.25ex}{\accentmark{12}{\char"0951}}माणन्दीवे\raisebox{0.25ex}{\accentmark{15}{\char"1CD2}}दिवे \nolinebreak\hskip 0.08em \mbox{॥५॥}}} \par}

\vspace{0.8em}
                % --- SUBSECTIONS (Mantras) ---
                
                % Conditional rendering based on output_mode
                \par\filbreak

\phantomsection

\addcontentsline{toc}{subsection}{ऋक् ६}

{\centering \footnotesize \textcolor{AccentPurple}{॥ कृष्ण आंगिरसः इन्द्रः जगती ॥} \par}

\vspace{0.12em}

{\centering \textbf{\textcolor{AccentBlue}{अच्छावाइन्द्रम्माता\raisebox{0.15ex}{\accentmark{12}{\char"1CF9}}यास्वर्यू\raisebox{0.25ex}{\accentmark{15}{\char"1CD2}}वा\raisebox{0.25ex}{\accentmark{12}{\char"0951}}ः\kern 0.05em ।\hskip 0.12em plus 0.04em minus 0.02em स\raisebox{0.25ex}{\accentmark{15}{\char"1CD2}}ध्री\raisebox{0.25ex}{\accentmark{15}{\char"1CD2}}चीर्विश्वा\raisebox{0.25ex}{\accentmark{12}{\char"0951}}ऊशाती\raisebox{0.25ex}{\accentmark{15}{\char"1CD2}}रा\raisebox{0.25ex}{\accentmark{12}{\char"0951}}नूषत्\kern 0.05em ।\hskip 0.12em plus 0.04em minus 0.02em पा\raisebox{0.25ex}{\accentmark{15}{\char"1CD2}}री\raisebox{0.25ex}{\accentmark{12}{\char"0951}}ष्वजन्ताज\raisebox{0.15ex}{\accentmark{12}{\char"1CF9}}ना\raisebox{0.15ex}{\accentmark{12}{\char"1CF9}}यो\raisebox{0.25ex}{\accentmark{12}{\char"0951}}याथापातिम्\kern 0.05em ।\hskip 0.12em plus 0.04em minus 0.02em म\raisebox{0.25ex}{\accentmark{15}{\char"1CD2}}र्यान्ना\raisebox{0.25ex}{\accentmark{15}{\char"1CD2}}शुन्ध्यु\raisebox{0.25ex}{\accentmark{15}{\char"1CD2}}मा॑\raisebox{0.25ex}{\accentmark{12}{\char"0951}}घा\raisebox{0.15ex}{\accentmark{12}{\char"1CF8}}वा\raisebox{0.25ex}{\accentmark{12}{\char"0951}}नामूता\raisebox{0.25ex}{\accentmark{15}{\char"1CD2}}ये \nolinebreak\hskip 0.08em \mbox{॥६॥}}} \par}

\vspace{0.8em}
                % --- SUBSECTIONS (Mantras) ---
                
                % Conditional rendering based on output_mode
                \par\filbreak

\phantomsection

\addcontentsline{toc}{subsection}{ऋक् ७}

{\centering \textbf{\textcolor{AccentBlue}{॥ सव्यआंगिरसःइन्द्रःजगती ॥ ॥ सव्यआंगिरसःइन्द्रःजगती ॥ आ\raisebox{0.25ex}{\accentmark{15}{\char"1CD2}}भि\raisebox{0.25ex}{\accentmark{15}{\char"1CD2}}त्यं\raisebox{0.25ex}{\accentmark{15}{\char"1CD2}}मेषंपू\raisebox{0.25ex}{\accentmark{12}{\char"0951}}रुहूता\raisebox{0.25ex}{\accentmark{15}{\char"1CD2}}मृग्मी\raisebox{0.25ex}{\accentmark{15}{\char"1CD2}}यम्\kern 0.05em ।\hskip 0.12em plus 0.04em minus 0.02em इन्द्रं\raisebox{0.25ex}{\accentmark{12}{\char"0951}}गीर्भि\raisebox{0.25ex}{\accentmark{12}{\char"0951}}र्मादतावस्वोअर्णा\raisebox{0.25ex}{\accentmark{12}{\char"0951}}वम्\kern 0.05em ।\hskip 0.12em plus 0.04em minus 0.02em य\raisebox{0.25ex}{\accentmark{15}{\char"1CD2}}स्या\raisebox{0.25ex}{\accentmark{15}{\char"1CD2}}द्यावोना\raisebox{0.25ex}{\accentmark{15}{\char"1CD2}}वीचा\raisebox{0.15ex}{\accentmark{12}{\char"1CF9}}र\raisebox{0.25ex}{\accentmark{12}{\char"0951}}न्तीमानू\raisebox{0.25ex}{\accentmark{12}{\char"0951}}षम्\kern 0.05em ।\hskip 0.12em plus 0.04em minus 0.02em भु\raisebox{0.25ex}{\accentmark{15}{\char"1CD2}}जे\raisebox{0.25ex}{\accentmark{15}{\char"1CD2}}मं\raisebox{0.25ex}{\accentmark{15}{\char"1CD2}}हीष्ठामाभी\raisebox{0.25ex}{\accentmark{15}{\char"1CD2}}विप्रा\raisebox{0.25ex}{\accentmark{12}{\char"0951}}मर्चत \nolinebreak\hskip 0.08em \mbox{॥७॥}}} \par}

\vspace{0.8em}
                % --- SUBSECTIONS (Mantras) ---
                
                % Conditional rendering based on output_mode
                \par\filbreak

\phantomsection

\addcontentsline{toc}{subsection}{ऋक् ८}

{\centering \textbf{\textcolor{AccentBlue}{त्यं\raisebox{0.25ex}{\accentmark{15}{\char"1CD2}}सूमेष\raisebox{0.25ex}{\accentmark{15}{\char"1CD2}}म्मा\raisebox{0.25ex}{\accentmark{12}{\char"0951}}हयास्व\raisebox{0.25ex}{\accentmark{15}{\char"1CD2}}र्वी\raisebox{0.25ex}{\accentmark{15}{\char"1CD2}}दम्\raisebox{0.25ex}{\accentmark{12}{\char"0951}}\kern 0.05em ।\hskip 0.12em plus 0.04em minus 0.02em शा\raisebox{0.25ex}{\accentmark{15}{\char"1CD2}}तं\raisebox{0.15ex}{\accentmark{12}{\char"1CF9}}यस्या\raisebox{0.25ex}{\accentmark{12}{\char"0951}}सूभूवाठ\raisebox{0.25ex}{\accentmark{12}{\char"0951}}साक\raisebox{0.25ex}{\accentmark{15}{\char"1CD2}}मीरा\raisebox{0.25ex}{\accentmark{12}{\char"0951}}ते\kern 0.05em ।\hskip 0.12em plus 0.04em minus 0.02em अ\raisebox{0.25ex}{\accentmark{15}{\char"1CD2}}त्यन्ना\raisebox{0.25ex}{\accentmark{15}{\char"1CD2}}वा\raisebox{0.25ex}{\accentmark{15}{\char"1CD2}}जं\raisebox{0.25ex}{\accentmark{12}{\char"0951}}हवानस्यादंरा\raisebox{0.25ex}{\accentmark{15}{\char"1CD2}}थम्\raisebox{0.25ex}{\accentmark{12}{\char"0951}}\kern 0.05em ।\hskip 0.12em plus 0.04em minus 0.02em ए\raisebox{0.25ex}{\accentmark{15}{\char"1CD2}}न्द्रं\raisebox{0.15ex}{\accentmark{12}{\char"1CF8}}ववृत्यामा\raisebox{0.25ex}{\accentmark{15}{\char"1CD2}}वा\raisebox{0.25ex}{\accentmark{12}{\char"0951}}सेसू\raisebox{0.25ex}{\accentmark{15}{\char"1CD2}}वृक्तीभी\raisebox{0.25ex}{\accentmark{12}{\char"0951}}ः \nolinebreak\hskip 0.08em \mbox{॥८॥}}} \par}

\vspace{0.8em}
                % --- SUBSECTIONS (Mantras) ---
                
                % Conditional rendering based on output_mode
                \par\filbreak

\phantomsection

\addcontentsline{toc}{subsection}{ऋक् ९}

{\centering \footnotesize \textcolor{AccentPurple}{॥ भरद्वाजो बार्हस्पत्यः द्यावापृथिवी जगती ॥} \par}

\vspace{0.12em}

{\centering \textbf{\textcolor{AccentBlue}{घृ\raisebox{0.25ex}{\accentmark{15}{\char"1CD2}}ता\raisebox{0.15ex}{\accentmark{12}{\char"1CF9}}वा\raisebox{0.25ex}{\accentmark{12}{\char"0951}}तीभू\raisebox{0.25ex}{\accentmark{15}{\char"1CD2}}वानानामाभि\raisebox{0.25ex}{\accentmark{15}{\char"1CD2}}श्रीयाः\kern 0.05em ।\hskip 0.12em plus 0.04em minus 0.02em उ\raisebox{0.25ex}{\accentmark{15}{\char"1CD2}}र्वी\raisebox{0.25ex}{\accentmark{12}{\char"0951}}पृथ्वी\raisebox{0.15ex}{\accentmark{12}{\char"1CF9}}माधूदु\raisebox{0.15ex}{\accentmark{12}{\char"1CF9}}घे\raisebox{0.25ex}{\accentmark{12}{\char"0951}}सुपे\raisebox{0.25ex}{\accentmark{15}{\char"1CD2}}शा\raisebox{0.25ex}{\accentmark{12}{\char"0951}}सा\kern 0.05em ।\hskip 0.12em plus 0.04em minus 0.02em द्या\raisebox{0.25ex}{\accentmark{15}{\char"1CD2}}वा\raisebox{0.25ex}{\accentmark{12}{\char"0951}}पृथिवी\raisebox{0.25ex}{\accentmark{15}{\char"1CD2}}वारू\raisebox{0.25ex}{\accentmark{12}{\char"0951}}णस्याध\raisebox{0.25ex}{\accentmark{15}{\char"1CD2}}र्मा\raisebox{0.25ex}{\accentmark{12}{\char"0951}}णा\kern 0.05em ।\hskip 0.12em plus 0.04em minus 0.02em विष्का\raisebox{0.25ex}{\accentmark{12}{\char"0951}}भितेआजरेभू\raisebox{0.25ex}{\accentmark{15}{\char"1CD2}}रिरेतसा \nolinebreak\hskip 0.08em \mbox{॥९॥}}} \par}

\vspace{0.8em}
                % --- SUBSECTIONS (Mantras) ---
                
                % Conditional rendering based on output_mode
                \par\filbreak

\phantomsection

\addcontentsline{toc}{subsection}{ऋक् १०}

{\centering \footnotesize \textcolor{AccentPurple}{॥ मेधातिथिः काण्वः (मान्धाता यौवनाश्वः) इन्द्रः महापंक्तिः ॥} \par}

\vspace{0.12em}

{\centering \textbf{\textcolor{AccentBlue}{ऊर्भे\raisebox{0.15ex}{\accentmark{12}{\char"1CF9}}या\raisebox{0.15ex}{\accentmark{12}{\char"1CF9}}दिन्द्रारो\raisebox{0.25ex}{\accentmark{15}{\char"1CD2}}द\raisebox{0.25ex}{\accentmark{12}{\char"0951}}सि\kern 0.05em ।\hskip 0.12em plus 0.04em minus 0.02em आपःप्रा\raisebox{0.25ex}{\accentmark{15}{\char"1CD2}}थोषा\raisebox{0.25ex}{\accentmark{15}{\char"1CD2}}ई\raisebox{0.25ex}{\accentmark{12}{\char"0951}}वा\kern 0.05em ।\hskip 0.12em plus 0.04em minus 0.02em माहा\raisebox{0.25ex}{\accentmark{15}{\char"1CD2}}न्त\raisebox{0.25ex}{\accentmark{12}{\char"0951}}न्त्वामाही\raisebox{0.25ex}{\accentmark{15}{\char"1CD2}}नाम्\raisebox{0.25ex}{\accentmark{12}{\char"0951}}\kern 0.05em ।\hskip 0.12em plus 0.04em minus 0.02em स\raisebox{0.25ex}{\accentmark{15}{\char"1CD2}}म्रा\raisebox{0.25ex}{\accentmark{12}{\char"0951}}जन्वर्षाणीनाम्\raisebox{0.25ex}{\accentmark{15}{\char"1CD2}}\kern 0.05em ।\hskip 0.12em plus 0.04em minus 0.02em दे\raisebox{0.25ex}{\accentmark{15}{\char"1CD2}}वी\raisebox{0.25ex}{\accentmark{15}{\char"1CD2}}जनीत्र\raisebox{0.25ex}{\accentmark{12}{\char"0951}}यजीजनात्\kern 0.05em ।\hskip 0.12em plus 0.04em minus 0.02em भ\raisebox{0.25ex}{\accentmark{15}{\char"1CD2}}द्रा\raisebox{0.25ex}{\accentmark{15}{\char"1CD2}}ज\raisebox{0.25ex}{\accentmark{15}{\char"1CD2}}नीत्रयजीजना\raisebox{0.25ex}{\accentmark{15}{\char"1CD2}}त् \nolinebreak\hskip 0.08em \mbox{॥१०॥}}} \par}

\vspace{0.8em}
                % --- SUBSECTIONS (Mantras) ---
                
                % Conditional rendering based on output_mode
                \par\filbreak

\phantomsection

\addcontentsline{toc}{subsection}{ऋक् ११}

{\centering \footnotesize \textcolor{AccentPurple}{॥ कुत्स आंगिरसः इन्द्रः जगती ॥} \par}

\vspace{0.12em}

{\centering \textbf{\textcolor{AccentBlue}{प्रा\raisebox{0.25ex}{\accentmark{15}{\char"1CD2}}मन्दीने\raisebox{0.25ex}{\accentmark{12}{\char"0951}}पीतूमाद॑\raisebox{0.25ex}{\accentmark{12}{\char"0951}}र्चातावा\raisebox{0.25ex}{\accentmark{15}{\char"1CD2}}चा\raisebox{0.25ex}{\accentmark{12}{\char"0951}}ः\kern 0.05em ।\hskip 0.12em plus 0.04em minus 0.02em याःकृष्णागर्भा\raisebox{0.25ex}{\accentmark{12}{\char"0951}}णीरा\raisebox{0.15ex}{\accentmark{12}{\char"1CF9}}ह\raisebox{0.25ex}{\accentmark{12}{\char"0951}}न्नृजीश्वा\raisebox{0.25ex}{\accentmark{12}{\char"0951}}ना\kern 0.05em ।\hskip 0.12em plus 0.04em minus 0.02em अव\raisebox{0.25ex}{\accentmark{15}{\char"1CD2}}स्यावो\raisebox{0.25ex}{\accentmark{15}{\char"1CD2}}वृ\raisebox{0.25ex}{\accentmark{15}{\char"1CD2}}षा\raisebox{0.25ex}{\accentmark{12}{\char"0951}}णंव\raisebox{0.25ex}{\accentmark{15}{\char"1CD2}}ज्रा\raisebox{0.25ex}{\accentmark{12}{\char"0951}}दक्षिणम्\kern 0.05em ।\hskip 0.12em plus 0.04em minus 0.02em मा\raisebox{0.25ex}{\accentmark{15}{\char"1CD2}}रूत्व\raisebox{0.25ex}{\accentmark{12}{\char"0951}}न्तंसख्या\raisebox{0.25ex}{\accentmark{15}{\char"1CD2}}या\raisebox{0.25ex}{\accentmark{12}{\char"0951}}हुवेमही \nolinebreak\hskip 0.08em \mbox{॥११॥}}} \par}

\vspace{0.8em}

            % --- End of Kandah Closing Floral Line & Counts ---
            \par\nopagebreak\vspace{0.9em plus 0.3em minus 0.1em}
            \centerline{\textcolor{black!35}{\rule{1.8cm}{0.35pt}}\;\raisebox{-0.22ex}{\textcolor{vedicmaroon}{\large\SixFlowerPetalDotted}}\;\textcolor{black!35}{\rule{1.8cm}{0.35pt}}}
            \nopagebreak\vspace{0.45em}
            \centerline{\devafont\textbf{\textcolor{AccentBlue}{॥ ऋक् सङ्ख्या : ११ ॥}}}
        
            \clearpage

            % --- SECTION ---
            
            \phantomsection  

            \section*{\centering द्वितीय खण्डः}
            \addcontentsline{toc}{section}{द्वितीय खण्डः} 
            \markright{द्वितीय खण्डः}
            \nopagebreak % Keep section title with first mantra
            
            % --- SECTION-LEVEL PROCEDURE LINK ---

                % --- SUBSECTIONS (Mantras) ---
                
                % Conditional rendering based on output_mode
                \par\filbreak

\phantomsection

\addcontentsline{toc}{subsection}{ऋक् १}

{\centering \footnotesize \textcolor{AccentPurple}{॥ गोतमो (सम्मदो वा) राहूगणः इन्द्रः पङ्क्त्तिः ॥} \par}

\vspace{0.12em}

{\centering \textbf{\textcolor{AccentBlue}{स्वादो\raisebox{0.25ex}{\accentmark{15}{\char"1CD2}}रित्था\raisebox{0.15ex}{\accentmark{12}{\char"1CF9}}वी\raisebox{0.25ex}{\accentmark{12}{\char"0951}}षूवा\raisebox{0.25ex}{\accentmark{15}{\char"1CD2}}ताः\kern 0.05em ।\hskip 0.12em plus 0.04em minus 0.02em मा\raisebox{0.25ex}{\accentmark{15}{\char"1CD2}}धो\raisebox{0.25ex}{\accentmark{12}{\char"0951}}ःपिबन्तिगौर्याःठ\raisebox{0.25ex}{\accentmark{12}{\char"0951}}\kern 0.05em ।\hskip 0.12em plus 0.04em minus 0.02em या\raisebox{0.25ex}{\accentmark{15}{\char"1CD2}}इ\raisebox{0.25ex}{\accentmark{15}{\char"1CD2}}न्द्रे\raisebox{0.25ex}{\accentmark{12}{\char"0951}}णासा\raisebox{0.25ex}{\accentmark{15}{\char"1CD2}}याववा\raisebox{0.25ex}{\accentmark{12}{\char"0951}}रिः\kern 0.05em ।\hskip 0.12em plus 0.04em minus 0.02em वृष्णामाद\raisebox{0.25ex}{\accentmark{12}{\char"0951}}न्तीशो\raisebox{0.25ex}{\accentmark{15}{\char"1CD2}}भा\raisebox{0.25ex}{\accentmark{15}{\char"1CD2}}था\kern 0.05em ।\hskip 0.12em plus 0.04em minus 0.02em व\raisebox{0.25ex}{\accentmark{15}{\char"1CD2}}स्वीरा\raisebox{0.15ex}{\accentmark{12}{\char"1CF9}}नू\raisebox{0.25ex}{\accentmark{12}{\char"0951}}स्वारा\raisebox{0.25ex}{\accentmark{15}{\char"1CD2}}ज्य\raisebox{0.25ex}{\accentmark{12}{\char"0951}}म् \nolinebreak\hskip 0.08em \mbox{॥१॥}}} \par}

\vspace{0.8em}
                % --- SUBSECTIONS (Mantras) ---
                
                % Conditional rendering based on output_mode
                \par\filbreak

\phantomsection

\addcontentsline{toc}{subsection}{ऋक् २}

{\centering \footnotesize \textcolor{AccentPurple}{॥ गोतमो (सम्मदो वा) राहूगणः इन्द्रः पङ्क्त्तिः ॥} \par}

\vspace{0.12em}

{\centering \textbf{\textcolor{AccentBlue}{इत्था\raisebox{0.25ex}{\accentmark{15}{\char"1CD2}}हीसो\raisebox{0.25ex}{\accentmark{15}{\char"1CD2}}माइ\raisebox{0.25ex}{\accentmark{15}{\char"1CD2}}न्मा\raisebox{0.25ex}{\accentmark{15}{\char"1CD2}}दा\raisebox{0.25ex}{\accentmark{12}{\char"0951}}ः\kern 0.05em ।\hskip 0.12em plus 0.04em minus 0.02em ब्र\raisebox{0.15ex}{\accentmark{12}{\char"1CF8}}ह्मा\raisebox{0.25ex}{\accentmark{12}{\char"0951}}चाका\raisebox{0.25ex}{\accentmark{15}{\char"1CD2}}रावा\raisebox{0.25ex}{\accentmark{15}{\char"1CD2}}र्धा\raisebox{0.25ex}{\accentmark{12}{\char"0951}}नम्\kern 0.05em ।\hskip 0.12em plus 0.04em minus 0.02em शावी\raisebox{0.25ex}{\accentmark{12}{\char"0951}}ष्ठवाज्रिन्नो\raisebox{0.25ex}{\accentmark{12}{\char"0951}}ज\raisebox{0.25ex}{\accentmark{12}{\char"0951}}सा\kern 0.05em ।\hskip 0.12em plus 0.04em minus 0.02em पृ\raisebox{0.15ex}{\accentmark{12}{\char"1CF8}}थिव्या\raisebox{0.15ex}{\accentmark{12}{\char"1CF9}}नि\raisebox{0.15ex}{\accentmark{12}{\char"1CF9}}श्शा\raisebox{0.25ex}{\accentmark{12}{\char"0951}}शाआ\raisebox{0.25ex}{\accentmark{15}{\char"1CD2}}हिं\kern 0.05em ।\hskip 0.12em plus 0.04em minus 0.02em अर्चन्ना\raisebox{0.15ex}{\accentmark{12}{\char"1CF9}}नूस्वारा\raisebox{0.25ex}{\accentmark{15}{\char"1CD2}}ज्य\raisebox{0.25ex}{\accentmark{12}{\char"0951}}म् \nolinebreak\hskip 0.08em \mbox{॥२॥}}} \par}

\vspace{0.8em}
                % --- SUBSECTIONS (Mantras) ---
                
                % Conditional rendering based on output_mode
                \par\filbreak

\phantomsection

\addcontentsline{toc}{subsection}{ऋक् ३}

{\centering \footnotesize \textcolor{AccentPurple}{॥ गोतमो (सम्मदो वा) राहूगणः इन्द्रः पङ्क्त्तिः ॥} \par}

\vspace{0.12em}

{\centering \textbf{\textcolor{AccentBlue}{इ\raisebox{0.25ex}{\accentmark{15}{\char"1CD2}}न्द्रोमा\raisebox{0.25ex}{\accentmark{15}{\char"1CD2}}दायवावृधे\kern 0.05em ।\hskip 0.12em plus 0.04em minus 0.02em शावा\raisebox{0.25ex}{\accentmark{12}{\char"0951}}सेवृत्रा\raisebox{0.15ex}{\accentmark{12}{\char"1CF8}}हा\raisebox{0.25ex}{\accentmark{15}{\char"1CD2}}नृ\raisebox{0.25ex}{\accentmark{15}{\char"1CD2}}र्भी\raisebox{0.25ex}{\accentmark{12}{\char"0951}}ः\kern 0.05em ।\hskip 0.12em plus 0.04em minus 0.02em त\raisebox{0.25ex}{\accentmark{15}{\char"1CD2}}मिन्माहा\raisebox{0.25ex}{\accentmark{15}{\char"1CD2}}त्स्वाजीषु\kern 0.05em ।\hskip 0.12em plus 0.04em minus 0.02em ऊतीमर्भेहवामहे\kern 0.05em ।\hskip 0.12em plus 0.04em minus 0.02em सा\raisebox{0.15ex}{\accentmark{12}{\char"1CF8}}वा\raisebox{0.15ex}{\accentmark{12}{\char"1CF8}}जे\raisebox{0.25ex}{\accentmark{12}{\char"0951}}षुप्रा\raisebox{0.25ex}{\accentmark{15}{\char"1CD2}}नो\raisebox{0.25ex}{\accentmark{12}{\char"0951}}विषात् \nolinebreak\hskip 0.08em \mbox{॥३॥}}} \par}

\vspace{0.8em}
                % --- SUBSECTIONS (Mantras) ---
                
                % Conditional rendering based on output_mode
                \par\filbreak

\phantomsection

\addcontentsline{toc}{subsection}{ऋक् ४}

{\centering \footnotesize \textcolor{AccentPurple}{॥ गोतमो (सम्मदो वा) राहूगणः इन्द्रः पङ्क्त्तिः ॥} \par}

\vspace{0.12em}

{\centering \textbf{\textcolor{AccentBlue}{इन्द्रातू\raisebox{0.25ex}{\accentmark{15}{\char"1CD2}}भ्या\raisebox{0.25ex}{\accentmark{15}{\char"1CD2}}मीद\raisebox{0.25ex}{\accentmark{12}{\char"0951}}द्रीवः\kern 0.05em ।\hskip 0.12em plus 0.04em minus 0.02em आ\raisebox{0.25ex}{\accentmark{15}{\char"1CD2}}नु\raisebox{0.25ex}{\accentmark{12}{\char"0951}}क्तंवज्रिन्वीर्यम्\raisebox{0.25ex}{\accentmark{12}{\char"0951}}\kern 0.05em ।\hskip 0.12em plus 0.04em minus 0.02em यद्धत्यं\raisebox{0.25ex}{\accentmark{15}{\char"1CD2}}मायी\raisebox{0.15ex}{\accentmark{12}{\char"1CF9}}नं\raisebox{0.25ex}{\accentmark{12}{\char"0951}}मृगम्\kern 0.05em ।\hskip 0.12em plus 0.04em minus 0.02em ता\raisebox{0.25ex}{\accentmark{15}{\char"1CD2}}वत्य\raisebox{0.25ex}{\accentmark{15}{\char"1CD2}}म्माययावा\raisebox{0.25ex}{\accentmark{15}{\char"1CD2}}धीः\kern 0.05em ।\hskip 0.12em plus 0.04em minus 0.02em अर्चन्ना\raisebox{0.15ex}{\accentmark{12}{\char"1CF9}}नू\raisebox{0.25ex}{\accentmark{12}{\char"0951}}स्वारा\raisebox{0.25ex}{\accentmark{15}{\char"1CD2}}ज्य\raisebox{0.25ex}{\accentmark{12}{\char"0951}}म् \nolinebreak\hskip 0.08em \mbox{॥४॥}}} \par}

\vspace{0.8em}
                % --- SUBSECTIONS (Mantras) ---
                
                % Conditional rendering based on output_mode
                \par\filbreak

\phantomsection

\addcontentsline{toc}{subsection}{ऋक् ५}

{\centering \footnotesize \textcolor{AccentPurple}{॥ गोतमो (सम्मदो वा) राहूगणः इन्द्रः पङ्क्त्तिः ॥} \par}

\vspace{0.12em}

{\centering \textbf{\textcolor{AccentBlue}{प्रे\raisebox{0.25ex}{\accentmark{15}{\char"1CD2}}ह्याभीठ\raisebox{0.25ex}{\accentmark{12}{\char"0951}}हीधृष्णूहि\raisebox{0.25ex}{\accentmark{15}{\char"1CD2}}\kern 0.05em ।\hskip 0.12em plus 0.04em minus 0.02em नातेवाज्रोनी\raisebox{0.25ex}{\accentmark{15}{\char"1CD2}}यंसते\kern 0.05em ।\hskip 0.12em plus 0.04em minus 0.02em इन्द्रानृम्णं\raisebox{0.15ex}{\accentmark{12}{\char"1CF9}}ही\raisebox{0.25ex}{\accentmark{12}{\char"0951}}तेशावा\raisebox{0.25ex}{\accentmark{12}{\char"0951}}ः\kern 0.05em ।\hskip 0.12em plus 0.04em minus 0.02em हानो\raisebox{0.25ex}{\accentmark{12}{\char"0951}}वृत्र\raisebox{0.15ex}{\accentmark{12}{\char"1CF9}}ञ्ज\raisebox{0.15ex}{\accentmark{12}{\char"1CF9}}या\raisebox{0.25ex}{\accentmark{12}{\char"0951}}आ\raisebox{0.25ex}{\accentmark{12}{\char"0951}}पः\kern 0.05em ।\hskip 0.12em plus 0.04em minus 0.02em अर्चन्नानू\raisebox{0.25ex}{\accentmark{12}{\char"0951}}स्वारा\raisebox{0.25ex}{\accentmark{15}{\char"1CD2}}ज्य\raisebox{0.25ex}{\accentmark{12}{\char"0951}}म् \nolinebreak\hskip 0.08em \mbox{॥५॥}}} \par}

\vspace{0.8em}
                % --- SUBSECTIONS (Mantras) ---
                
                % Conditional rendering based on output_mode
                \par\filbreak

\phantomsection

\addcontentsline{toc}{subsection}{ऋक् ६}

{\centering \footnotesize \textcolor{AccentPurple}{॥ गोतमो (सम्मदो वा) राहूगणः इन्द्रः पङ्क्त्तिः ॥} \par}

\vspace{0.12em}

{\centering \textbf{\textcolor{AccentBlue}{या\raisebox{0.25ex}{\accentmark{15}{\char"1CD2}}दूदीराताआज\raisebox{0.25ex}{\accentmark{15}{\char"1CD2}}या\raisebox{0.25ex}{\accentmark{12}{\char"0951}}ः\kern 0.05em ।\hskip 0.12em plus 0.04em minus 0.02em धृ\raisebox{0.25ex}{\accentmark{15}{\char"1CD2}}ष्णा\raisebox{0.25ex}{\accentmark{15}{\char"1CD2}}वे\raisebox{0.25ex}{\accentmark{12}{\char"0951}}धीयातेधानम्\raisebox{0.25ex}{\accentmark{12}{\char"0951}}\kern 0.05em ।\hskip 0.12em plus 0.04em minus 0.02em युं\raisebox{0.25ex}{\accentmark{15}{\char"1CD2}}क्ष्वा\raisebox{0.15ex}{\accentmark{12}{\char"1CF9}}मा\raisebox{0.25ex}{\accentmark{12}{\char"0951}}दच्यू\raisebox{0.25ex}{\accentmark{15}{\char"1CD2}}ताहा\raisebox{0.25ex}{\accentmark{15}{\char"1CD2}}री\kern 0.05em ।\hskip 0.12em plus 0.04em minus 0.02em कंहा\raisebox{0.25ex}{\accentmark{15}{\char"1CD2}}नःकं\raisebox{0.25ex}{\accentmark{15}{\char"1CD2}}वा\raisebox{0.25ex}{\accentmark{15}{\char"1CD2}}सौदधः\kern 0.05em ।\hskip 0.12em plus 0.04em minus 0.02em अस्मं\raisebox{0.15ex}{\accentmark{12}{\char"1CF9}}इ\raisebox{0.25ex}{\accentmark{12}{\char"0951}}न्द्रावा\raisebox{0.25ex}{\accentmark{15}{\char"1CD2}}सौ\raisebox{0.25ex}{\accentmark{12}{\char"0951}}दधः \nolinebreak\hskip 0.08em \mbox{॥६॥}}} \par}

\vspace{0.8em}
                % --- SUBSECTIONS (Mantras) ---
                
                % Conditional rendering based on output_mode
                \par\filbreak

\phantomsection

\addcontentsline{toc}{subsection}{ऋक् ७}

{\centering \textbf{\textcolor{AccentBlue}{॥ गोतमो(सम्मदोवा)राहूगणःइन्द्रःपङ्क्त्तिः ॥ ॥ गोतमो(सम्मदोवा)राहूगणःइन्द्रःपङ्क्त्तिः ॥ अ\raisebox{0.25ex}{\accentmark{15}{\char"1CD2}}क्ष\raisebox{0.25ex}{\accentmark{15}{\char"1CD2}}न्ना\raisebox{0.25ex}{\accentmark{15}{\char"1CD2}}मी\raisebox{0.25ex}{\accentmark{12}{\char"0951}}मादन्ताही\raisebox{0.25ex}{\accentmark{15}{\char"1CD2}}\kern 0.05em ।\hskip 0.12em plus 0.04em minus 0.02em आ\raisebox{0.15ex}{\accentmark{12}{\char"1CF8}}वा\raisebox{0.25ex}{\accentmark{12}{\char"0951}}प्रियाआधूषत\kern 0.05em ।\hskip 0.12em plus 0.04em minus 0.02em अ\raisebox{0.25ex}{\accentmark{15}{\char"1CD2}}स्तो\raisebox{0.25ex}{\accentmark{12}{\char"0951}}षातस्वाभा\raisebox{0.25ex}{\accentmark{12}{\char"0951}}नवः\kern 0.05em ।\hskip 0.12em plus 0.04em minus 0.02em त्रिप्रान्नावीष्ठयामति\kern 0.05em ।\hskip 0.12em plus 0.04em minus 0.02em यो\raisebox{0.25ex}{\accentmark{15}{\char"1CD2}}जान्विन्द्रातेहारी\raisebox{0.25ex}{\accentmark{12}{\char"0951}} \nolinebreak\hskip 0.08em \mbox{॥७॥}}} \par}

\vspace{0.8em}
                % --- SUBSECTIONS (Mantras) ---
                
                % Conditional rendering based on output_mode
                \par\filbreak

\phantomsection

\addcontentsline{toc}{subsection}{ऋक् ८}

{\centering \textbf{\textcolor{AccentBlue}{ऊ\raisebox{0.25ex}{\accentmark{15}{\char"1CD2}}पोषू\raisebox{0.15ex}{\accentmark{12}{\char"1CF9}}श्रूणूहीगीरा\raisebox{0.25ex}{\accentmark{12}{\char"0951}}ः\kern 0.05em ।\hskip 0.12em plus 0.04em minus 0.02em मा\raisebox{0.15ex}{\accentmark{12}{\char"1CF8}}घा\raisebox{0.25ex}{\accentmark{12}{\char"0951}}वन्माताथाइव\kern 0.05em ।\hskip 0.12em plus 0.04em minus 0.02em कादा\raisebox{0.15ex}{\accentmark{12}{\char"1CF9}}ना\raisebox{0.25ex}{\accentmark{12}{\char"0951}}स्सू\raisebox{0.25ex}{\accentmark{12}{\char"0951}}नृता\raisebox{0.25ex}{\accentmark{12}{\char"0951}}वतः\kern 0.05em ।\hskip 0.12em plus 0.04em minus 0.02em काराई\raisebox{0.25ex}{\accentmark{15}{\char"1CD2}}दा\raisebox{0.25ex}{\accentmark{12}{\char"0951}}र्थार्यासाइत्\kern 0.05em ।\hskip 0.12em plus 0.04em minus 0.02em यो\raisebox{0.25ex}{\accentmark{15}{\char"1CD2}}जान्वि\raisebox{0.25ex}{\accentmark{12}{\char"0951}}न्द्राते\raisebox{0.25ex}{\accentmark{15}{\char"1CD2}}हा\raisebox{0.25ex}{\accentmark{15}{\char"1CD2}}री\raisebox{0.25ex}{\accentmark{12}{\char"0951}} \nolinebreak\hskip 0.08em \mbox{॥८॥}}} \par}

\vspace{0.8em}
                % --- SUBSECTIONS (Mantras) ---
                
                % Conditional rendering based on output_mode
                \par\filbreak

\phantomsection

\addcontentsline{toc}{subsection}{ऋक् ९}

{\centering \footnotesize \textcolor{AccentPurple}{॥ त्रित आप्त्यः (कुत्स आङ्गिरसो वा) विश्वेदेवाः पङ्क्त्तिः ॥} \par}

\vspace{0.12em}

{\centering \textbf{\textcolor{AccentBlue}{चन्द्रा\raisebox{0.15ex}{\accentmark{12}{\char"1CF9}}मा\raisebox{0.25ex}{\accentmark{12}{\char"0951}}प्सूआ\raisebox{0.25ex}{\accentmark{12}{\char"0951}}टट्यन्तारा\kern 0.05em ।\hskip 0.12em plus 0.04em minus 0.02em सूपर्णोधा\raisebox{0.25ex}{\accentmark{12}{\char"0951}}वतेदीवि\raisebox{0.25ex}{\accentmark{15}{\char"1CD2}}\kern 0.05em ।\hskip 0.12em plus 0.04em minus 0.02em नावो\raisebox{0.25ex}{\accentmark{12}{\char"0951}}हीरण्यनेमयः\kern 0.05em ।\hskip 0.12em plus 0.04em minus 0.02em पा\raisebox{0.25ex}{\accentmark{15}{\char"1CD2}}दं\raisebox{0.25ex}{\accentmark{15}{\char"1CD2}}वि\raisebox{0.25ex}{\accentmark{15}{\char"1CD2}}न्द\raisebox{0.25ex}{\accentmark{12}{\char"0951}}न्तीवीद्यूताः\kern 0.05em ।\hskip 0.12em plus 0.04em minus 0.02em विक्त\raisebox{0.15ex}{\accentmark{12}{\char"1CF9}}म्मेअस्य\raisebox{0.25ex}{\accentmark{15}{\char"1CD2}}रो\raisebox{0.25ex}{\accentmark{15}{\char"1CD2}}दसि \nolinebreak\hskip 0.08em \mbox{॥९॥}}} \par}

\vspace{0.8em}
                % --- SUBSECTIONS (Mantras) ---
                
                % Conditional rendering based on output_mode
                \par\filbreak

\phantomsection

\addcontentsline{toc}{subsection}{ऋक् १०}

{\centering \footnotesize \textcolor{AccentPurple}{॥ अवस्युरात्रयेः अश्विनौ पङ्क्त्तिः ॥} \par}

\vspace{0.12em}

{\centering \textbf{\textcolor{AccentBlue}{प्रा\raisebox{0.15ex}{\accentmark{12}{\char"1CF8}}तिप्री\raisebox{0.25ex}{\accentmark{12}{\char"0951}}यातामंरा\raisebox{0.25ex}{\accentmark{15}{\char"1CD2}}थम्\raisebox{0.25ex}{\accentmark{12}{\char"0951}}\kern 0.05em ।\hskip 0.12em plus 0.04em minus 0.02em वृषाणंवासूवाहा\raisebox{0.25ex}{\accentmark{12}{\char"0951}}नाम्\kern 0.05em ।\hskip 0.12em plus 0.04em minus 0.02em स्तोता\raisebox{0.25ex}{\accentmark{15}{\char"1CD2}}वा\raisebox{0.25ex}{\accentmark{12}{\char"0951}}मश्विनावृ\raisebox{0.25ex}{\accentmark{15}{\char"1CD2}}षी\raisebox{0.25ex}{\accentmark{12}{\char"0951}}ः\kern 0.05em ।\hskip 0.12em plus 0.04em minus 0.02em स्तोमेभिर्भूषातिप्राती\kern 0.05em ।\hskip 0.12em plus 0.04em minus 0.02em मध्वीमा\raisebox{0.25ex}{\accentmark{15}{\char"1CD2}}मा\raisebox{0.25ex}{\accentmark{12}{\char"0951}}श्रुतंहा\raisebox{0.25ex}{\accentmark{15}{\char"1CD2}}वम् \nolinebreak\hskip 0.08em \mbox{॥१०॥}}} \par}

\vspace{0.8em}

            % --- End of Kandah Closing Floral Line & Counts ---
            \par\nopagebreak\vspace{0.9em plus 0.3em minus 0.1em}
            \centerline{\textcolor{black!35}{\rule{1.8cm}{0.35pt}}\;\raisebox{-0.22ex}{\textcolor{vedicmaroon}{\large\SixFlowerPetalDotted}}\;\textcolor{black!35}{\rule{1.8cm}{0.35pt}}}
            \nopagebreak\vspace{0.45em}
            \centerline{\devafont\textbf{\textcolor{AccentBlue}{॥ ऋक् सङ्ख्या : १० ॥}}}
        
            \clearpage

            % --- SECTION ---
            
            \phantomsection  

            \section*{\centering तृतीय खण्डः}
            \addcontentsline{toc}{section}{तृतीय खण्डः} 
            \markright{तृतीय खण्डः}
            \nopagebreak % Keep section title with first mantra
            
            % --- SECTION-LEVEL PROCEDURE LINK ---

                % --- SUBSECTIONS (Mantras) ---
                
                % Conditional rendering based on output_mode
                \par\filbreak

\phantomsection

\addcontentsline{toc}{subsection}{ऋक् १}

{\centering \footnotesize \textcolor{AccentPurple}{॥ वसुश्रुत अत्रेयः अग्निः पङ्क्त्तिः ॥} \par}

\vspace{0.12em}

{\centering \textbf{\textcolor{AccentBlue}{आ\raisebox{0.25ex}{\accentmark{15}{\char"1CD2}}तेअग्नेइधीमही\raisebox{0.25ex}{\accentmark{15}{\char"1CD2}}\kern 0.05em ।\hskip 0.12em plus 0.04em minus 0.02em घूमन्तं\raisebox{0.25ex}{\accentmark{12}{\char"0951}}देवाआज\raisebox{0.25ex}{\accentmark{15}{\char"1CD2}}रम्\raisebox{0.25ex}{\accentmark{12}{\char"0951}}\kern 0.05em ।\hskip 0.12em plus 0.04em minus 0.02em यद्धस्या\raisebox{0.25ex}{\accentmark{15}{\char"1CD2}}तेपानी\raisebox{0.25ex}{\accentmark{12}{\char"0951}}यसी\kern 0.05em ।\hskip 0.12em plus 0.04em minus 0.02em सामी\raisebox{0.25ex}{\accentmark{15}{\char"1CD2}}द्दीदया\raisebox{0.25ex}{\accentmark{12}{\char"0951}}तिद्याविईळंस्तोतृ\raisebox{0.25ex}{\accentmark{15}{\char"1CD2}}भ्यआ\raisebox{0.25ex}{\accentmark{15}{\char"1CD2}}भारा \nolinebreak\hskip 0.08em \mbox{॥१॥}}} \par}

\vspace{0.8em}
                % --- SUBSECTIONS (Mantras) ---
                
                % Conditional rendering based on output_mode
                \par\filbreak

\phantomsection

\addcontentsline{toc}{subsection}{ऋक् २}

{\centering \footnotesize \textcolor{AccentPurple}{॥ विमद ऐन्द्रः (प्राजापत्यो वा, वसुकृद्वा वासुकः) अग्निः पङ्क्त्तिः ॥} \par}

\vspace{0.12em}

{\centering \textbf{\textcolor{AccentBlue}{अग्निन्नस्वा\raisebox{0.25ex}{\accentmark{15}{\char"1CD2}}वृ\raisebox{0.25ex}{\accentmark{12}{\char"0951}}क्तीभिः\kern 0.05em ।\hskip 0.12em plus 0.04em minus 0.02em हो\raisebox{0.25ex}{\accentmark{15}{\char"1CD2}}ता\raisebox{0.25ex}{\accentmark{12}{\char"0951}}रन्त्वावृणीमहे\kern 0.05em ।\hskip 0.12em plus 0.04em minus 0.02em शीरं\raisebox{0.15ex}{\accentmark{12}{\char"1CF9}}पा\raisebox{0.25ex}{\accentmark{12}{\char"0951}}वाका\raisebox{0.25ex}{\accentmark{12}{\char"0951}}शोची\raisebox{0.25ex}{\accentmark{15}{\char"1CD2}}षंवीवोमादे\kern 0.05em ।\hskip 0.12em plus 0.04em minus 0.02em यज्ञे\raisebox{0.15ex}{\accentmark{12}{\char"1CF9}}षु\raisebox{0.25ex}{\accentmark{12}{\char"0951}}स्तीर्णाब\raisebox{0.25ex}{\accentmark{12}{\char"0951}}ऋहीषंवीवाक्षसे \nolinebreak\hskip 0.08em \mbox{॥२॥}}} \par}

\vspace{0.8em}
                % --- SUBSECTIONS (Mantras) ---
                
                % Conditional rendering based on output_mode
                \par\filbreak

\phantomsection

\addcontentsline{toc}{subsection}{ऋक् ३}

{\centering \footnotesize \textcolor{AccentPurple}{॥ सत्वश्रवा आत्रेयः उषाः पङ्क्त्तिः ॥} \par}

\vspace{0.12em}

{\centering \textbf{\textcolor{AccentBlue}{माहे\raisebox{0.15ex}{\accentmark{12}{\char"1CF9}}नो\raisebox{0.25ex}{\accentmark{12}{\char"0951}}अद्याबो\raisebox{0.25ex}{\accentmark{12}{\char"0951}}धय\kern 0.05em ।\hskip 0.12em plus 0.04em minus 0.02em ऊ\raisebox{0.15ex}{\accentmark{12}{\char"1CF8}}षो\raisebox{0.25ex}{\accentmark{12}{\char"0951}}राये\raisebox{0.25ex}{\accentmark{15}{\char"1CD2}}दिवी\raisebox{0.25ex}{\accentmark{15}{\char"1CD2}}द्माती\kern 0.05em ।\hskip 0.12em plus 0.04em minus 0.02em या\raisebox{0.25ex}{\accentmark{15}{\char"1CD2}}थाचि\raisebox{0.25ex}{\accentmark{15}{\char"1CD2}}न्नोआ\raisebox{0.25ex}{\accentmark{15}{\char"1CD2}}बो\raisebox{0.25ex}{\accentmark{12}{\char"0951}}धयः\kern 0.05em ।\hskip 0.12em plus 0.04em minus 0.02em सत्या\raisebox{0.25ex}{\accentmark{15}{\char"1CD2}}श्रावसीवाय्ये\raisebox{0.25ex}{\accentmark{15}{\char"1CD2}}\kern 0.05em ।\hskip 0.12em plus 0.04em minus 0.02em सू\raisebox{0.15ex}{\accentmark{12}{\char"1CF8}}जा\raisebox{0.25ex}{\accentmark{12}{\char"0951}}तेअश्वा\raisebox{0.25ex}{\accentmark{12}{\char"0951}}सूनृते \nolinebreak\hskip 0.08em \mbox{॥३॥}}} \par}

\vspace{0.8em}
                % --- SUBSECTIONS (Mantras) ---
                
                % Conditional rendering based on output_mode
                \par\filbreak

\phantomsection

\addcontentsline{toc}{subsection}{ऋक् ४}

{\centering \footnotesize \textcolor{AccentPurple}{॥ विमद ऐन्द्रः (प्राजापत्यो वा, वसुकृद्वा वासुकः) सोमः पङ्क्त्तिः ॥} \par}

\vspace{0.12em}

{\centering \textbf{\textcolor{AccentBlue}{भद्र\raisebox{0.25ex}{\accentmark{15}{\char"1CD2}}न्नो\raisebox{0.25ex}{\accentmark{15}{\char"1CD2}}आपी\raisebox{0.25ex}{\accentmark{12}{\char"0951}}वातय\kern 0.05em ।\hskip 0.12em plus 0.04em minus 0.02em मानोद\raisebox{0.15ex}{\accentmark{12}{\char"1CF9}}क्षा\raisebox{0.25ex}{\accentmark{12}{\char"0951}}मूत\raisebox{0.25ex}{\accentmark{15}{\char"1CD2}}क्रातुम्\raisebox{0.25ex}{\accentmark{12}{\char"0951}}\kern 0.05em ।\hskip 0.12em plus 0.04em minus 0.02em अथा\raisebox{0.25ex}{\accentmark{15}{\char"1CD2}}तेसख्ये\raisebox{0.15ex}{\accentmark{12}{\char"1CF9}}अ\raisebox{0.25ex}{\accentmark{12}{\char"0951}}न्धासावी\raisebox{0.25ex}{\accentmark{15}{\char"1CD2}}वोमादे\kern 0.05em ।\hskip 0.12em plus 0.04em minus 0.02em रा\raisebox{0.25ex}{\accentmark{15}{\char"1CD2}}णागा\raisebox{0.25ex}{\accentmark{15}{\char"1CD2}}वोनया\raisebox{0.15ex}{\accentmark{12}{\char"1CF9}}वा\raisebox{0.25ex}{\accentmark{12}{\char"0951}}सेवीवा\raisebox{0.25ex}{\accentmark{12}{\char"0951}}क्षसे \nolinebreak\hskip 0.08em \mbox{॥४॥}}} \par}

\vspace{0.8em}
                % --- SUBSECTIONS (Mantras) ---
                
                % Conditional rendering based on output_mode
                \par\filbreak

\phantomsection

\addcontentsline{toc}{subsection}{ऋक् ५}

{\centering \footnotesize \textcolor{AccentPurple}{॥ गोतमो राहूगणः इन्द्रः पङ्क्त्तिः ॥} \par}

\vspace{0.12em}

{\centering \textbf{\textcolor{AccentBlue}{कृ\raisebox{0.15ex}{\accentmark{12}{\char"1CF8}}त्वामाहं\raisebox{0.25ex}{\accentmark{15}{\char"1CD2}}आनूष्वाध\raisebox{0.25ex}{\accentmark{15}{\char"1CD2}}ः\kern 0.05em ।\hskip 0.12em plus 0.04em minus 0.02em भीम\raisebox{0.25ex}{\accentmark{15}{\char"1CD2}}आ\raisebox{0.25ex}{\accentmark{15}{\char"1CD2}}वा\raisebox{0.25ex}{\accentmark{12}{\char"0951}}वृतेशा\raisebox{0.25ex}{\accentmark{15}{\char"1CD2}}वा\raisebox{0.25ex}{\accentmark{12}{\char"0951}}ः\kern 0.05em ।\hskip 0.12em plus 0.04em minus 0.02em श्रीय्या\raisebox{0.15ex}{\accentmark{12}{\char"1CF9}}रिष्वा\raisebox{0.15ex}{\accentmark{12}{\char"1CF9}}उपाका\raisebox{0.25ex}{\accentmark{15}{\char"1CD2}}योः\kern 0.05em ।\hskip 0.12em plus 0.04em minus 0.02em नी\raisebox{0.25ex}{\accentmark{15}{\char"1CD2}}शीप्रीहा\raisebox{0.25ex}{\accentmark{15}{\char"1CD2}}रीवान्दधे\kern 0.05em ।\hskip 0.12em plus 0.04em minus 0.02em ह\raisebox{0.15ex}{\accentmark{12}{\char"1CF8}}स्ता\raisebox{0.25ex}{\accentmark{15}{\char"1CD2}}योर्वज्रा\raisebox{0.25ex}{\accentmark{12}{\char"0951}}मायासम् \nolinebreak\hskip 0.08em \mbox{॥५॥}}} \par}

\vspace{0.8em}
                % --- SUBSECTIONS (Mantras) ---
                
                % Conditional rendering based on output_mode
                \par\filbreak

\phantomsection

\addcontentsline{toc}{subsection}{ऋक् ६}

{\centering \footnotesize \textcolor{AccentPurple}{॥ गोतमो राहूगणः इन्द्रः पङ्क्त्तिः ॥} \par}

\vspace{0.12em}

{\centering \textbf{\textcolor{AccentBlue}{साघातं\raisebox{0.15ex}{\accentmark{12}{\char"1CF9}}वृ\raisebox{0.25ex}{\accentmark{12}{\char"0951}}षाणंरा\raisebox{0.25ex}{\accentmark{15}{\char"1CD2}}थम्\raisebox{0.25ex}{\accentmark{12}{\char"0951}}\kern 0.05em ।\hskip 0.12em plus 0.04em minus 0.02em आधीतिष्ठाती\raisebox{0.25ex}{\accentmark{12}{\char"0951}}गोवी\raisebox{0.25ex}{\accentmark{12}{\char"0951}}दम्\kern 0.05em ।\hskip 0.12em plus 0.04em minus 0.02em यःपात्रं\raisebox{0.25ex}{\accentmark{12}{\char"0951}}हार्यो\raisebox{0.25ex}{\accentmark{15}{\char"1CD2}}जान\raisebox{0.25ex}{\accentmark{12}{\char"0951}}म्\kern 0.05em ।\hskip 0.12em plus 0.04em minus 0.02em पूर्णा\raisebox{0.15ex}{\accentmark{12}{\char"1CF9}}मि\raisebox{0.25ex}{\accentmark{12}{\char"0951}}न्द्राचि\raisebox{0.25ex}{\accentmark{15}{\char"1CD2}}के\raisebox{0.25ex}{\accentmark{12}{\char"0951}}ततिः\kern 0.05em ।\hskip 0.12em plus 0.04em minus 0.02em यो\raisebox{0.25ex}{\accentmark{15}{\char"1CD2}}जान्विन्द्रातेहा\raisebox{0.25ex}{\accentmark{15}{\char"1CD2}}री \nolinebreak\hskip 0.08em \mbox{॥६॥}}} \par}

\vspace{0.8em}
                % --- SUBSECTIONS (Mantras) ---
                
                % Conditional rendering based on output_mode
                \par\filbreak

\phantomsection

\addcontentsline{toc}{subsection}{ऋक् ७}

{\centering \footnotesize \textcolor{AccentPurple}{॥ वसुश्रुत अत्रेयः अग्निः पङ्क्त्तिः ॥} \par}

\vspace{0.12em}

{\centering \textbf{\textcolor{AccentBlue}{अ\raisebox{0.25ex}{\accentmark{15}{\char"1CD2}}ग्निन्त\raisebox{0.15ex}{\accentmark{12}{\char"1CF9}}म्मन्ये\raisebox{0.25ex}{\accentmark{15}{\char"1CD2}}योवासू\raisebox{0.25ex}{\accentmark{12}{\char"0951}}ः\kern 0.05em ।\hskip 0.12em plus 0.04em minus 0.02em अस्तंयय्य\raisebox{0.15ex}{\accentmark{12}{\char"1CF9}}न्ती\raisebox{0.25ex}{\accentmark{12}{\char"0951}}धेना\raisebox{0.25ex}{\accentmark{15}{\char"1CD2}}वाः\kern 0.05em ।\hskip 0.12em plus 0.04em minus 0.02em अ\raisebox{0.25ex}{\accentmark{15}{\char"1CD2}}स्ता\raisebox{0.25ex}{\accentmark{12}{\char"0951}}म\raisebox{0.25ex}{\accentmark{15}{\char"1CD2}}र्वतोआशा\raisebox{0.25ex}{\accentmark{15}{\char"1CD2}}वा\raisebox{0.25ex}{\accentmark{12}{\char"0951}}ः\kern 0.05em ।\hskip 0.12em plus 0.04em minus 0.02em अस्तन्नित्या\raisebox{0.25ex}{\accentmark{12}{\char"0951}}सोवाजी\raisebox{0.25ex}{\accentmark{15}{\char"1CD2}}नाः\kern 0.05em ।\hskip 0.12em plus 0.04em minus 0.02em ई\raisebox{0.15ex}{\accentmark{12}{\char"1CF8}}षं\raisebox{0.25ex}{\accentmark{12}{\char"0951}}स्तोतृभ्यआ\raisebox{0.25ex}{\accentmark{15}{\char"1CD2}}भा\raisebox{0.25ex}{\accentmark{12}{\char"0951}}रा \nolinebreak\hskip 0.08em \mbox{॥७॥}}} \par}

\vspace{0.8em}
                % --- SUBSECTIONS (Mantras) ---
                
                % Conditional rendering based on output_mode
                \par\filbreak

\phantomsection

\addcontentsline{toc}{subsection}{ऋक् ८}

{\centering \footnotesize \textcolor{AccentPurple}{॥ अङ्होमुग्वामदेव्यः (कुल्मलबर्हिषः शैलूषिर्वा) विश्वेदेवाः बृहती ॥} \par}

\vspace{0.12em}

{\centering \textbf{\textcolor{AccentBlue}{नतमं\raisebox{0.25ex}{\accentmark{15}{\char"1CD2}}हो\raisebox{0.25ex}{\accentmark{12}{\char"0951}}ना\raisebox{0.15ex}{\accentmark{12}{\char"1CF9}}दू\raisebox{0.25ex}{\accentmark{12}{\char"0951}}रीतम्\kern 0.05em ।\hskip 0.12em plus 0.04em minus 0.02em देवा\raisebox{0.25ex}{\accentmark{12}{\char"0951}}सोआष्टामा\raisebox{0.25ex}{\accentmark{15}{\char"1CD2}}र्क्त्यम्\kern 0.05em ।\hskip 0.12em plus 0.04em minus 0.02em सा\raisebox{0.25ex}{\accentmark{15}{\char"1CD2}}जो\raisebox{0.25ex}{\accentmark{15}{\char"1CD2}}षा\raisebox{0.25ex}{\accentmark{12}{\char"0951}}सोयामर्यामा\kern 0.05em ।\hskip 0.12em plus 0.04em minus 0.02em मित्रोनाय\raisebox{0.25ex}{\accentmark{12}{\char"0951}}न्तीवा\raisebox{0.15ex}{\accentmark{12}{\char"1CF9}}रूणोआ\raisebox{0.25ex}{\accentmark{15}{\char"1CD2}}तिद्वी\raisebox{0.25ex}{\accentmark{15}{\char"1CD2}}षाः \nolinebreak\hskip 0.08em \mbox{॥८॥}}} \par}

\vspace{0.8em}

            % --- End of Kandah Closing Floral Line & Counts ---
            \par\nopagebreak\vspace{0.9em plus 0.3em minus 0.1em}
            \centerline{\textcolor{black!35}{\rule{1.8cm}{0.35pt}}\;\raisebox{-0.22ex}{\textcolor{vedicmaroon}{\large\SixFlowerPetalDotted}}\;\textcolor{black!35}{\rule{1.8cm}{0.35pt}}}
            \nopagebreak\vspace{0.45em}
            \centerline{\devafont\textbf{\textcolor{AccentBlue}{॥ ऋक् सङ्ख्या : ८ ॥}}}
        
            \clearpage

            % --- SECTION ---
            
            \phantomsection  

            \section*{\centering चतुर्थ खण्डः}
            \addcontentsline{toc}{section}{चतुर्थ खण्डः} 
            \markright{चतुर्थ खण्डः}
            \nopagebreak % Keep section title with first mantra
            
            % --- SECTION-LEVEL PROCEDURE LINK ---

                % --- SUBSECTIONS (Mantras) ---
                
                % Conditional rendering based on output_mode
                \par\filbreak

\phantomsection

\addcontentsline{toc}{subsection}{ऋक् १}

{\centering \footnotesize \textcolor{AccentPurple}{॥ ऋण – त्रसदस्यू पवमानः सोमः द्विपदा विराट् ॥} \par}

\vspace{0.12em}

{\centering \textbf{\textcolor{AccentBlue}{पा\raisebox{0.25ex}{\accentmark{15}{\char"1CD2}}रिप्राध\raisebox{0.25ex}{\accentmark{12}{\char"0951}}न्वाइन्द्रा\raisebox{0.25ex}{\accentmark{12}{\char"0951}}यसोमा\raisebox{0.25ex}{\accentmark{15}{\char"1CD2}}\kern 0.05em ।\hskip 0.12em plus 0.04em minus 0.02em स्वादु\raisebox{0.25ex}{\accentmark{15}{\char"1CD2}}र्मित्रा\raisebox{0.25ex}{\accentmark{15}{\char"1CD2}}य\kern 0.05em ।\hskip 0.12em plus 0.04em minus 0.02em पू\raisebox{0.25ex}{\accentmark{15}{\char"1CD2}}ष्णेभा\raisebox{0.25ex}{\accentmark{15}{\char"1CD2}}गा\raisebox{0.25ex}{\accentmark{12}{\char"0951}}य \nolinebreak\hskip 0.08em \mbox{॥१॥}}} \par}

\vspace{0.8em}
                % --- SUBSECTIONS (Mantras) ---
                
                % Conditional rendering based on output_mode
                \par\filbreak

\phantomsection

\addcontentsline{toc}{subsection}{ऋक् २}

{\centering \footnotesize \textcolor{AccentPurple}{॥ ऋण – त्रसदस्यू (त्र्यरुणस्त्रैवृष्णः, त्रसदस्युः पौरुकुत्सः) पवमानः सोमः त्रिपदा अनुष्टुप्पिपीलिकामध्या ॥} \par}

\vspace{0.12em}

{\centering \textbf{\textcolor{AccentBlue}{पर्यू\raisebox{0.15ex}{\accentmark{12}{\char"1CF9}}षु\raisebox{0.15ex}{\accentmark{12}{\char"1CF9}}प्रा\raisebox{0.15ex}{\accentmark{12}{\char"1CF9}}ध\raisebox{0.25ex}{\accentmark{12}{\char"0951}}न्वा\raisebox{0.25ex}{\accentmark{15}{\char"1CD2}}वा\raisebox{0.25ex}{\accentmark{15}{\char"1CD2}}ज\raisebox{0.25ex}{\accentmark{12}{\char"0951}}सातये\kern 0.05em ।\hskip 0.12em plus 0.04em minus 0.02em पारी\raisebox{0.25ex}{\accentmark{12}{\char"0951}}वृ\raisebox{0.25ex}{\accentmark{12}{\char"0951}}त्रा\raisebox{0.15ex}{\accentmark{12}{\char"1CF9}}णी\raisebox{0.25ex}{\accentmark{12}{\char"0951}}सक्षा\raisebox{0.25ex}{\accentmark{15}{\char"1CD2}}णी\raisebox{0.25ex}{\accentmark{12}{\char"0951}}ः\kern 0.05em ।\hskip 0.12em plus 0.04em minus 0.02em द्वीष\raisebox{0.15ex}{\accentmark{12}{\char"1CF9}}स्ता\raisebox{0.25ex}{\accentmark{12}{\char"0951}}राधी\raisebox{0.25ex}{\accentmark{15}{\char"1CD2}}या\raisebox{0.15ex}{\accentmark{12}{\char"1CF8}}ऋणाया\raisebox{0.25ex}{\accentmark{15}{\char"1CD2}}ना\raisebox{0.25ex}{\accentmark{12}{\char"0951}}ईरसे \nolinebreak\hskip 0.08em \mbox{॥२॥}}} \par}

\vspace{0.8em}
                % --- SUBSECTIONS (Mantras) ---
                
                % Conditional rendering based on output_mode
                \par\filbreak

\phantomsection

\addcontentsline{toc}{subsection}{ऋक् ३}

{\centering \textbf{\textcolor{AccentBlue}{॥ ऋण–त्रसदस्यूपवमानःसोमःद्विपदाविराट् ॥ ॥ ऋण–त्रसदस्यूपवमानःसोमःद्विपदाविराट् ॥ पा\raisebox{0.25ex}{\accentmark{15}{\char"1CD2}}वा\raisebox{0.25ex}{\accentmark{12}{\char"0951}}स्वसोमा\raisebox{0.25ex}{\accentmark{15}{\char"1CD2}}\kern 0.05em ।\hskip 0.12em plus 0.04em minus 0.02em माहा\raisebox{0.15ex}{\accentmark{12}{\char"1CF9}}न्सा\raisebox{0.25ex}{\accentmark{12}{\char"0951}}मुद्रः\kern 0.05em ।\hskip 0.12em plus 0.04em minus 0.02em पीता\raisebox{0.25ex}{\accentmark{15}{\char"1CD2}}दे\raisebox{0.25ex}{\accentmark{12}{\char"0951}}वानाम्\raisebox{0.25ex}{\accentmark{12}{\char"0951}}\kern 0.05em ।\hskip 0.12em plus 0.04em minus 0.02em विश्वाभि\raisebox{0.25ex}{\accentmark{15}{\char"1CD2}}धामा \nolinebreak\hskip 0.08em \mbox{॥३॥}}} \par}

\vspace{0.8em}
                % --- SUBSECTIONS (Mantras) ---
                
                % Conditional rendering based on output_mode
                \par\filbreak

\phantomsection

\addcontentsline{toc}{subsection}{ऋक् ४}

{\centering \textbf{\textcolor{AccentBlue}{पा\raisebox{0.25ex}{\accentmark{15}{\char"1CD2}}वा\raisebox{0.25ex}{\accentmark{12}{\char"0951}}स्वसोमा\raisebox{0.25ex}{\accentmark{15}{\char"1CD2}}\kern 0.05em ।\hskip 0.12em plus 0.04em minus 0.02em माहे\raisebox{0.25ex}{\accentmark{15}{\char"1CD2}}दक्षा\raisebox{0.25ex}{\accentmark{12}{\char"0951}}य\kern 0.05em ।\hskip 0.12em plus 0.04em minus 0.02em अ\raisebox{0.25ex}{\accentmark{15}{\char"1CD2}}श्वो\raisebox{0.25ex}{\accentmark{15}{\char"1CD2}}ना\raisebox{0.25ex}{\accentmark{15}{\char"1CD2}}निक्तः\kern 0.05em ।\hskip 0.12em plus 0.04em minus 0.02em वाजी\raisebox{0.25ex}{\accentmark{15}{\char"1CD2}}धा\raisebox{0.25ex}{\accentmark{15}{\char"1CD2}}नाय \nolinebreak\hskip 0.08em \mbox{॥४॥}}} \par}

\vspace{0.8em}
                % --- SUBSECTIONS (Mantras) ---
                
                % Conditional rendering based on output_mode
                \par\filbreak

\phantomsection

\addcontentsline{toc}{subsection}{ऋक् ५}

{\centering \footnotesize \textcolor{AccentPurple}{॥ ऋण – त्रसदस्यू पवमानः सोमः द्विपदा विराट् ॥} \par}

\vspace{0.12em}

{\centering \textbf{\textcolor{AccentBlue}{इ\raisebox{0.25ex}{\accentmark{15}{\char"1CD2}}न्दू\raisebox{0.25ex}{\accentmark{12}{\char"0951}}ःपविष्टा\kern 0.05em ।\hskip 0.12em plus 0.04em minus 0.02em चारूर्मा\raisebox{0.25ex}{\accentmark{12}{\char"0951}}दा\raisebox{0.25ex}{\accentmark{12}{\char"0951}}य\kern 0.05em ।\hskip 0.12em plus 0.04em minus 0.02em आपा\raisebox{0.25ex}{\accentmark{15}{\char"1CD2}}मूपा\raisebox{0.25ex}{\accentmark{15}{\char"1CD2}}स्थे\raisebox{0.25ex}{\accentmark{12}{\char"0951}}\kern 0.05em ।\hskip 0.12em plus 0.04em minus 0.02em का\raisebox{0.25ex}{\accentmark{15}{\char"1CD2}}विर्भा\raisebox{0.25ex}{\accentmark{12}{\char"0951}}गा\raisebox{0.25ex}{\accentmark{12}{\char"0951}}य \nolinebreak\hskip 0.08em \mbox{॥५॥}}} \par}

\vspace{0.8em}
                % --- SUBSECTIONS (Mantras) ---
                
                % Conditional rendering based on output_mode
                \par\filbreak

\phantomsection

\addcontentsline{toc}{subsection}{ऋक् ६}

{\centering \footnotesize \textcolor{AccentPurple}{॥ ऋण – त्रसदस्यू (त्र्यरुणस्त्रैवृष्णः, त्रसदस्युः पौरुकुत्सः) पवमानः सोमः त्रिपदा अनुष्टुप्पिपीलिकामध्या ॥} \par}

\vspace{0.12em}

{\centering \textbf{\textcolor{AccentBlue}{आ\raisebox{0.25ex}{\accentmark{15}{\char"1CD2}}नूहित्वा\raisebox{0.15ex}{\accentmark{12}{\char"1CF8}}सूतं\raisebox{0.15ex}{\accentmark{12}{\char"1CF9}}सो\raisebox{0.25ex}{\accentmark{12}{\char"0951}}मा\raisebox{0.25ex}{\accentmark{15}{\char"1CD2}}मा\raisebox{0.25ex}{\accentmark{15}{\char"1CD2}}दा\raisebox{0.25ex}{\accentmark{12}{\char"0951}}मसि\kern 0.05em ।\hskip 0.12em plus 0.04em minus 0.02em मा\raisebox{0.25ex}{\accentmark{15}{\char"1CD2}}हं\raisebox{0.25ex}{\accentmark{15}{\char"1CD2}}सा\raisebox{0.25ex}{\accentmark{12}{\char"0951}}मरियारा\raisebox{0.25ex}{\accentmark{15}{\char"1CD2}}ज्ये\raisebox{0.25ex}{\accentmark{12}{\char"0951}}\kern 0.05em ।\hskip 0.12em plus 0.04em minus 0.02em वा\raisebox{0.15ex}{\accentmark{12}{\char"1CF8}}जं\raisebox{0.25ex}{\accentmark{12}{\char"0951}}आभि\raisebox{0.25ex}{\accentmark{15}{\char"1CD2}}पा\raisebox{0.25ex}{\accentmark{12}{\char"0951}}वमानप्रागा\raisebox{0.25ex}{\accentmark{12}{\char"0951}}हसे \nolinebreak\hskip 0.08em \mbox{॥६॥}}} \par}

\vspace{0.8em}
                % --- SUBSECTIONS (Mantras) ---
                
                % Conditional rendering based on output_mode
                \par\filbreak

\phantomsection

\addcontentsline{toc}{subsection}{ऋक् ७}

{\centering \footnotesize \textcolor{AccentPurple}{॥ वसिष्ठो मैत्रावरुणिः मरुतः द्विपदा विराट् ॥} \par}

\vspace{0.12em}

{\centering \textbf{\textcolor{AccentBlue}{का\raisebox{0.25ex}{\accentmark{15}{\char"1CD2}}ईंव्या१क्ताः\kern 0.05em ।\hskip 0.12em plus 0.04em minus 0.02em ना\raisebox{0.25ex}{\accentmark{15}{\char"1CD2}}रस्सा\raisebox{0.25ex}{\accentmark{15}{\char"1CD2}}नी\raisebox{0.25ex}{\accentmark{12}{\char"0951}}ळाः\kern 0.05em ।\hskip 0.12em plus 0.04em minus 0.02em रु\raisebox{0.25ex}{\accentmark{15}{\char"1CD2}}द्रा\raisebox{0.25ex}{\accentmark{15}{\char"1CD2}}स्साम\raisebox{0.25ex}{\accentmark{15}{\char"1CD2}}र्याः\kern 0.05em ।\hskip 0.12em plus 0.04em minus 0.02em या\raisebox{0.25ex}{\accentmark{15}{\char"1CD2}}थास्व\raisebox{0.25ex}{\accentmark{15}{\char"1CD2}}श्वाः\raisebox{0.25ex}{\accentmark{12}{\char"0951}} \nolinebreak\hskip 0.08em \mbox{॥७॥}}} \par}

\vspace{0.8em}
                % --- SUBSECTIONS (Mantras) ---
                
                % Conditional rendering based on output_mode
                \par\filbreak

\phantomsection

\addcontentsline{toc}{subsection}{ऋक् ८}

{\centering \footnotesize \textcolor{AccentPurple}{॥ वामदेवो गौतमः अग्निः पदपङ्क्तिः ॥} \par}

\vspace{0.12em}

{\centering \textbf{\textcolor{AccentBlue}{अ\raisebox{0.25ex}{\accentmark{15}{\char"1CD2}}ग्नेता\raisebox{0.25ex}{\accentmark{15}{\char"1CD2}}मद्या\raisebox{0.25ex}{\accentmark{15}{\char"1CD2}}\kern 0.05em ।\hskip 0.12em plus 0.04em minus 0.02em अश्वन्न\raisebox{0.25ex}{\accentmark{15}{\char"1CD2}}स्तोमैः\raisebox{0.25ex}{\accentmark{12}{\char"0951}}\kern 0.05em ।\hskip 0.12em plus 0.04em minus 0.02em क्रा\raisebox{0.25ex}{\accentmark{15}{\char"1CD2}}तून्ना\raisebox{0.25ex}{\accentmark{15}{\char"1CD2}}भद्रं\raisebox{0.25ex}{\accentmark{15}{\char"1CD2}}\kern 0.05em ।\hskip 0.12em plus 0.04em minus 0.02em हृदिस्पृ\raisebox{0.25ex}{\accentmark{15}{\char"1CD2}}शं\raisebox{0.25ex}{\accentmark{12}{\char"0951}}\kern 0.05em ।\hskip 0.12em plus 0.04em minus 0.02em ऋद्ध्या\raisebox{0.15ex}{\accentmark{12}{\char"1CF9}}मा\raisebox{0.25ex}{\accentmark{12}{\char"0951}}ताओ\raisebox{0.25ex}{\accentmark{15}{\char"1CD2}}हैः\raisebox{0.25ex}{\accentmark{12}{\char"0951}} \nolinebreak\hskip 0.08em \mbox{॥८॥}}} \par}

\vspace{0.8em}
                % --- SUBSECTIONS (Mantras) ---
                
                % Conditional rendering based on output_mode
                \par\filbreak

\phantomsection

\addcontentsline{toc}{subsection}{ऋक् ९}

{\centering \footnotesize \textcolor{AccentPurple}{॥ ऋण – त्रसदस्यू वाजिनः पुरउष्णिक् ॥} \par}

\vspace{0.12em}

{\centering \textbf{\textcolor{AccentBlue}{आवि\raisebox{0.15ex}{\accentmark{12}{\char"1CF8}}र्म\raisebox{0.25ex}{\accentmark{12}{\char"0951}}र्याआवा\raisebox{0.25ex}{\accentmark{12}{\char"0951}}जं\raisebox{0.25ex}{\accentmark{12}{\char"0951}}वाजी\raisebox{0.25ex}{\accentmark{15}{\char"1CD2}}नो\raisebox{0.25ex}{\accentmark{12}{\char"0951}}ग्मन्\kern 0.05em ।\hskip 0.12em plus 0.04em minus 0.02em देवा\raisebox{0.25ex}{\accentmark{15}{\char"1CD2}}स्या\raisebox{0.25ex}{\accentmark{12}{\char"0951}}सावीतु\raisebox{0.25ex}{\accentmark{15}{\char"1CD2}}स्सा\raisebox{0.25ex}{\accentmark{12}{\char"0951}}वे\kern 0.05em ।\hskip 0.12em plus 0.04em minus 0.02em स्वर्गा\raisebox{0.25ex}{\accentmark{12}{\char"0951}}अ\raisebox{0.15ex}{\accentmark{12}{\char"1CF9}}र्वन्तोजा\raisebox{0.25ex}{\accentmark{15}{\char"1CD2}}यर्या\raisebox{0.25ex}{\accentmark{12}{\char"0951}}ता \nolinebreak\hskip 0.08em \mbox{॥९॥}}} \par}

\vspace{0.8em}
                % --- SUBSECTIONS (Mantras) ---
                
                % Conditional rendering based on output_mode
                \par\filbreak

\phantomsection

\addcontentsline{toc}{subsection}{ऋक् १०}

{\centering \footnotesize \textcolor{AccentPurple}{॥ ऋण – त्रसदस्यू पवमानः सोमः द्विपदा विराट् ॥} \par}

\vspace{0.12em}

{\centering \textbf{\textcolor{AccentBlue}{पावा\raisebox{0.25ex}{\accentmark{12}{\char"0951}}स्वसोम\kern 0.05em ।\hskip 0.12em plus 0.04em minus 0.02em द्युम्नी\raisebox{0.15ex}{\accentmark{12}{\char"1CF9}}सूधारः\raisebox{0.25ex}{\accentmark{12}{\char"0951}}\kern 0.05em ।\hskip 0.12em plus 0.04em minus 0.02em माही\raisebox{0.25ex}{\accentmark{15}{\char"1CD2}}ना\raisebox{0.25ex}{\accentmark{12}{\char"0951}}मानूपूर्व्याः\raisebox{0.25ex}{\accentmark{15}{\char"1CD2}} \nolinebreak\hskip 0.08em \mbox{॥१०॥}}} \par}

\vspace{0.8em}

            % --- End of Kandah Closing Floral Line & Counts ---
            \par\nopagebreak\vspace{0.9em plus 0.3em minus 0.1em}
            \centerline{\textcolor{black!35}{\rule{1.8cm}{0.35pt}}\;\raisebox{-0.22ex}{\textcolor{vedicmaroon}{\large\SixFlowerPetalDotted}}\;\textcolor{black!35}{\rule{1.8cm}{0.35pt}}}
            \nopagebreak\vspace{0.45em}
            \centerline{\devafont\textbf{\textcolor{AccentBlue}{॥ ऋक् सङ्ख्या : १० ॥}}}
        
            \clearpage

            % --- SECTION ---
            
            \phantomsection  

            \section*{\centering पञ्चम खण्डः}
            \addcontentsline{toc}{section}{पञ्चम खण्डः} 
            \markright{पञ्चम खण्डः}
            \nopagebreak % Keep section title with first mantra
            
            % --- SECTION-LEVEL PROCEDURE LINK ---

                % --- SUBSECTIONS (Mantras) ---
                
                % Conditional rendering based on output_mode
                \par\filbreak

\phantomsection

\addcontentsline{toc}{subsection}{ऋक् १}

{\centering \footnotesize \textcolor{AccentPurple}{॥ नारदः काण्वः इन्द्रः उष्णिक् ॥} \par}

\vspace{0.12em}

{\centering \textbf{\textcolor{AccentBlue}{इन्द्रा\raisebox{0.25ex}{\accentmark{12}{\char"0951}}सूते\raisebox{0.25ex}{\accentmark{15}{\char"1CD2}}षुसो\raisebox{0.25ex}{\accentmark{15}{\char"1CD2}}मे\raisebox{0.25ex}{\accentmark{12}{\char"0951}}षु\kern 0.05em ।\hskip 0.12em plus 0.04em minus 0.02em क्रातुं\raisebox{0.25ex}{\accentmark{12}{\char"0951}}पुनीषाउक्त्थम्\kern 0.05em ।\hskip 0.12em plus 0.04em minus 0.02em वीदे\raisebox{0.25ex}{\accentmark{15}{\char"1CD2}}वृधा\raisebox{0.25ex}{\accentmark{15}{\char"1CD2}}स्याद\raisebox{0.25ex}{\accentmark{15}{\char"1CD2}}क्षा\raisebox{0.25ex}{\accentmark{12}{\char"0951}}स्यामाहं\raisebox{0.25ex}{\accentmark{15}{\char"1CD2}}हिष\raisebox{0.25ex}{\accentmark{15}{\char"1CD2}}ः \nolinebreak\hskip 0.08em \mbox{॥१॥}}} \par}

\vspace{0.8em}
                % --- SUBSECTIONS (Mantras) ---
                
                % Conditional rendering based on output_mode
                \par\filbreak

\phantomsection

\addcontentsline{toc}{subsection}{ऋक् २}

{\centering \footnotesize \textcolor{AccentPurple}{॥ गाेषूक्त्यश्वसूक्तिनाै काण्वायनो इन्द्रः उष्णिक् ॥} \par}

\vspace{0.12em}

{\centering \textbf{\textcolor{AccentBlue}{तमुवांठ\raisebox{0.25ex}{\accentmark{12}{\char"0951}}भिप्रा\raisebox{0.25ex}{\accentmark{15}{\char"1CD2}}गा\raisebox{0.25ex}{\accentmark{12}{\char"0951}}यत\kern 0.05em ।\hskip 0.12em plus 0.04em minus 0.02em पू\raisebox{0.25ex}{\accentmark{12}{\char"0951}}रुहूतं\raisebox{0.25ex}{\accentmark{15}{\char"1CD2}}पू\raisebox{0.25ex}{\accentmark{12}{\char"0951}}रूष्टुतम्\raisebox{0.25ex}{\accentmark{15}{\char"1CD2}}\kern 0.05em ।\hskip 0.12em plus 0.04em minus 0.02em इन्द्रं\raisebox{0.25ex}{\accentmark{12}{\char"0951}}गी\raisebox{0.15ex}{\accentmark{12}{\char"1CF9}}र्भीस्ता\raisebox{0.25ex}{\accentmark{12}{\char"0951}}वीषामा\raisebox{0.25ex}{\accentmark{15}{\char"1CD2}}वी\raisebox{0.25ex}{\accentmark{12}{\char"0951}}वासता \nolinebreak\hskip 0.08em \mbox{॥२॥}}} \par}

\vspace{0.8em}
                % --- SUBSECTIONS (Mantras) ---
                
                % Conditional rendering based on output_mode
                \par\filbreak

\phantomsection

\addcontentsline{toc}{subsection}{ऋक् ३}

{\centering \footnotesize \textcolor{AccentPurple}{॥ गाेषूक्त्यश्वसूक्तिनाै काण्वायनो इन्द्रः उष्णिक् ॥} \par}

\vspace{0.12em}

{\centering \textbf{\textcolor{AccentBlue}{तन्तेमा\raisebox{0.25ex}{\accentmark{15}{\char"1CD2}}दं\raisebox{0.25ex}{\accentmark{12}{\char"0951}}गृणीमसि\kern 0.05em ।\hskip 0.12em plus 0.04em minus 0.02em वृ\raisebox{0.25ex}{\accentmark{15}{\char"1CD2}}षा\raisebox{0.25ex}{\accentmark{12}{\char"0951}}णंप्रक्षु\raisebox{0.15ex}{\accentmark{12}{\char"1CF9}}सासा\raisebox{0.25ex}{\accentmark{12}{\char"0951}}हिम्\kern 0.05em ।\hskip 0.12em plus 0.04em minus 0.02em ऊलोकाकूत्नू\raisebox{0.25ex}{\accentmark{15}{\char"1CD2}}माद्रीवो\raisebox{0.25ex}{\accentmark{15}{\char"1CD2}}हारिश्रीय\raisebox{0.25ex}{\accentmark{12}{\char"0951}}म् \nolinebreak\hskip 0.08em \mbox{॥३॥}}} \par}

\vspace{0.8em}
                % --- SUBSECTIONS (Mantras) ---
                
                % Conditional rendering based on output_mode
                \par\filbreak

\phantomsection

\addcontentsline{toc}{subsection}{ऋक् ४}

{\centering \footnotesize \textcolor{AccentPurple}{॥ पर्वतः काण्वः इन्द्रः उष्णिक् ॥} \par}

\vspace{0.12em}

{\centering \textbf{\textcolor{AccentBlue}{य\raisebox{0.25ex}{\accentmark{15}{\char"1CD2}}त्सो\raisebox{0.25ex}{\accentmark{15}{\char"1CD2}}मा\raisebox{0.25ex}{\accentmark{12}{\char"0951}}मिन्द्रावि\raisebox{0.25ex}{\accentmark{15}{\char"1CD2}}ष्णा\raisebox{0.25ex}{\accentmark{12}{\char"0951}}वे\kern 0.05em ।\hskip 0.12em plus 0.04em minus 0.02em यद्वा\raisebox{0.25ex}{\accentmark{12}{\char"0951}}घात्रीता\raisebox{0.25ex}{\accentmark{15}{\char"1CD2}}आप्त्ये\raisebox{0.25ex}{\accentmark{15}{\char"1CD2}}\kern 0.05em ।\hskip 0.12em plus 0.04em minus 0.02em यद्वा\raisebox{0.15ex}{\accentmark{12}{\char"1CF9}}मा\raisebox{0.25ex}{\accentmark{12}{\char"0951}}रूत्सूम\raisebox{0.15ex}{\accentmark{12}{\char"1CF9}}न्दसेसमिन्दुभिः \nolinebreak\hskip 0.08em \mbox{॥४॥}}} \par}

\vspace{0.8em}
                % --- SUBSECTIONS (Mantras) ---
                
                % Conditional rendering based on output_mode
                \par\filbreak

\phantomsection

\addcontentsline{toc}{subsection}{ऋक् ५}

{\centering \textbf{\textcolor{AccentBlue}{॥ विश्वमनावैयश्वःइन्द्रःउष्णिक् ॥ ॥ विश्वमनावैयश्वःइन्द्रःउष्णिक् ॥ ए\raisebox{0.25ex}{\accentmark{15}{\char"1CD2}}दुमा\raisebox{0.15ex}{\accentmark{12}{\char"1CF9}}धोर्मा\raisebox{0.25ex}{\accentmark{12}{\char"0951}}दिन्ता\raisebox{0.25ex}{\accentmark{12}{\char"0951}}रम्\kern 0.05em ।\hskip 0.12em plus 0.04em minus 0.02em सि\raisebox{0.25ex}{\accentmark{15}{\char"1CD2}}ञ्चा\raisebox{0.15ex}{\accentmark{12}{\char"1CF9}}ध्व\raisebox{0.25ex}{\accentmark{12}{\char"0951}}र्योअ\raisebox{0.25ex}{\accentmark{15}{\char"1CD2}}न्धा\raisebox{0.25ex}{\accentmark{12}{\char"0951}}सा\kern 0.05em ।\hskip 0.12em plus 0.04em minus 0.02em ए\raisebox{0.25ex}{\accentmark{15}{\char"1CD2}}वा\raisebox{0.25ex}{\accentmark{15}{\char"1CD2}}हीवीर\raisebox{0.25ex}{\accentmark{15}{\char"1CD2}}स्ता\raisebox{0.25ex}{\accentmark{15}{\char"1CD2}}वा\raisebox{0.25ex}{\accentmark{12}{\char"0951}}तेसादा\raisebox{0.25ex}{\accentmark{15}{\char"1CD2}}वृ\raisebox{0.25ex}{\accentmark{12}{\char"0951}}धः \nolinebreak\hskip 0.08em \mbox{॥५॥}}} \par}

\vspace{0.8em}
                % --- SUBSECTIONS (Mantras) ---
                
                % Conditional rendering based on output_mode
                \par\filbreak

\phantomsection

\addcontentsline{toc}{subsection}{ऋक् ६}

{\centering \textbf{\textcolor{AccentBlue}{ए\raisebox{0.25ex}{\accentmark{15}{\char"1CD2}}न्दूमिन्द्रा\raisebox{0.25ex}{\accentmark{12}{\char"0951}}यसिञ्चत\kern 0.05em ।\hskip 0.12em plus 0.04em minus 0.02em पी\raisebox{0.25ex}{\accentmark{15}{\char"1CD2}}बातासोम्यम्माधू\raisebox{0.25ex}{\accentmark{12}{\char"0951}}\kern 0.05em ।\hskip 0.12em plus 0.04em minus 0.02em प्रारा\raisebox{0.25ex}{\accentmark{15}{\char"1CD2}}धां\raisebox{0.25ex}{\accentmark{12}{\char"0951}}सिचोदयतामहीत्वाना \nolinebreak\hskip 0.08em \mbox{॥६॥}}} \par}

\vspace{0.8em}
                % --- SUBSECTIONS (Mantras) ---
                
                % Conditional rendering based on output_mode
                \par\filbreak

\phantomsection

\addcontentsline{toc}{subsection}{ऋक् ७}

{\centering \footnotesize \textcolor{AccentPurple}{॥ विश्वमना वैयश्वः इन्द्रः उष्णिक् ॥} \par}

\vspace{0.12em}

{\centering \textbf{\textcolor{AccentBlue}{ए\raisebox{0.25ex}{\accentmark{15}{\char"1CD2}}तोन्विन्द्रंस्ता\raisebox{0.25ex}{\accentmark{15}{\char"1CD2}}वा\raisebox{0.25ex}{\accentmark{12}{\char"0951}}मा\kern 0.05em ।\hskip 0.12em plus 0.04em minus 0.02em सा\raisebox{0.15ex}{\accentmark{12}{\char"1CF8}}खा\raisebox{0.25ex}{\accentmark{12}{\char"0951}}यस्तोम्यन्ना\raisebox{0.25ex}{\accentmark{15}{\char"1CD2}}रम्\raisebox{0.25ex}{\accentmark{12}{\char"0951}}\kern 0.05em ।\hskip 0.12em plus 0.04em minus 0.02em कृष्टिर्यो\raisebox{0.15ex}{\accentmark{12}{\char"1CF9}}विश्वा\raisebox{0.25ex}{\accentmark{12}{\char"0951}}अभ्यास्ते\raisebox{0.25ex}{\accentmark{15}{\char"1CD2}}काई\raisebox{0.25ex}{\accentmark{15}{\char"1CD2}}त् \nolinebreak\hskip 0.08em \mbox{॥७॥}}} \par}

\vspace{0.8em}
                % --- SUBSECTIONS (Mantras) ---
                
                % Conditional rendering based on output_mode
                \par\filbreak

\phantomsection

\addcontentsline{toc}{subsection}{ऋक् ८}

{\centering \footnotesize \textcolor{AccentPurple}{॥ नृमेध आंगिरसः इन्द्रः उष्णिक् ॥} \par}

\vspace{0.12em}

{\centering \textbf{\textcolor{AccentBlue}{इ\raisebox{0.15ex}{\accentmark{12}{\char"1CF8}}न्द्रा\raisebox{0.25ex}{\accentmark{12}{\char"0951}}यासामा\raisebox{0.25ex}{\accentmark{12}{\char"0951}}गाय\raisebox{0.25ex}{\accentmark{15}{\char"1CD2}}ता\kern 0.05em ।\hskip 0.12em plus 0.04em minus 0.02em वि\raisebox{0.25ex}{\accentmark{15}{\char"1CD2}}प्रा\raisebox{0.25ex}{\accentmark{12}{\char"0951}}याबृहातेबृहा\raisebox{0.25ex}{\accentmark{15}{\char"1CD2}}त्\kern 0.05em ।\hskip 0.12em plus 0.04em minus 0.02em ब्र\raisebox{0.25ex}{\accentmark{15}{\char"1CD2}}ह्माकृ\raisebox{0.25ex}{\accentmark{15}{\char"1CD2}}ते\raisebox{0.25ex}{\accentmark{12}{\char"0951}}वीपश्ची\raisebox{0.25ex}{\accentmark{15}{\char"1CD2}}तेपानस्या\raisebox{0.25ex}{\accentmark{15}{\char"1CD2}}वे \nolinebreak\hskip 0.08em \mbox{॥८॥}}} \par}

\vspace{0.8em}
                % --- SUBSECTIONS (Mantras) ---
                
                % Conditional rendering based on output_mode
                \par\filbreak

\phantomsection

\addcontentsline{toc}{subsection}{ऋक् ९}

{\centering \footnotesize \textcolor{AccentPurple}{॥ गाेतमाे राहूगणः इन्द्रः उष्णिक् ॥} \par}

\vspace{0.12em}

{\centering \textbf{\textcolor{AccentBlue}{यए\raisebox{0.25ex}{\accentmark{15}{\char"1CD2}}काई\raisebox{0.25ex}{\accentmark{15}{\char"1CD2}}द्वीद\raisebox{0.25ex}{\accentmark{15}{\char"1CD2}}या\raisebox{0.25ex}{\accentmark{12}{\char"0951}}ते\kern 0.05em ।\hskip 0.12em plus 0.04em minus 0.02em वा\raisebox{0.25ex}{\accentmark{15}{\char"1CD2}}सूमर्क्ता\raisebox{0.25ex}{\accentmark{12}{\char"0951}}यादाशू\raisebox{0.25ex}{\accentmark{15}{\char"1CD2}}षे\raisebox{0.25ex}{\accentmark{12}{\char"0951}}\kern 0.05em ।\hskip 0.12em plus 0.04em minus 0.02em ई\raisebox{0.15ex}{\accentmark{12}{\char"1CF8}}शा\raisebox{0.25ex}{\accentmark{12}{\char"0951}}नोअ\raisebox{0.25ex}{\accentmark{15}{\char"1CD2}}प्रातिष्कुता\kern 0.05em ।\hskip 0.12em plus 0.04em minus 0.02em इन्द्रो\raisebox{0.15ex}{\accentmark{12}{\char"1CF9}}अङ्गा \nolinebreak\hskip 0.08em \mbox{॥९॥}}} \par}

\vspace{0.8em}
                % --- SUBSECTIONS (Mantras) ---
                
                % Conditional rendering based on output_mode
                \par\filbreak

\phantomsection

\addcontentsline{toc}{subsection}{ऋक् १०}

{\centering \footnotesize \textcolor{AccentPurple}{॥ विश्वमना वैयश्वः इन्द्रः उष्णिक् ॥} \par}

\vspace{0.12em}

{\centering \textbf{\textcolor{AccentBlue}{सा\raisebox{0.15ex}{\accentmark{12}{\char"1CF8}}खा\raisebox{0.25ex}{\accentmark{12}{\char"0951}}य\raisebox{0.25ex}{\accentmark{15}{\char"1CD2}}आ\raisebox{0.25ex}{\accentmark{15}{\char"1CD2}}शी\raisebox{0.25ex}{\accentmark{12}{\char"0951}}षामहे\kern 0.05em ।\hskip 0.12em plus 0.04em minus 0.02em ब्रह्मे\raisebox{0.25ex}{\accentmark{15}{\char"1CD2}}न्द्रा\raisebox{0.25ex}{\accentmark{12}{\char"0951}}यावज्रीणे\kern 0.05em ।\hskip 0.12em plus 0.04em minus 0.02em स्तूषा\raisebox{0.25ex}{\accentmark{15}{\char"1CD2}}ऊ\raisebox{0.25ex}{\accentmark{12}{\char"0951}}षूवोनृतामा\raisebox{0.25ex}{\accentmark{12}{\char"0951}}याधृष्णा\raisebox{0.25ex}{\accentmark{12}{\char"0951}}वे \nolinebreak\hskip 0.08em \mbox{॥१०॥}}} \par}

\vspace{0.8em}

            % --- End of Kandah Closing Floral Line & Counts ---
            \par\nopagebreak\vspace{0.9em plus 0.3em minus 0.1em}
            \centerline{\textcolor{black!35}{\rule{1.8cm}{0.35pt}}\;\raisebox{-0.22ex}{\textcolor{vedicmaroon}{\large\SixFlowerPetalDotted}}\;\textcolor{black!35}{\rule{1.8cm}{0.35pt}}}
            \nopagebreak\vspace{0.45em}
            \centerline{\devafont\textbf{\textcolor{AccentBlue}{॥ ऋक् सङ्ख्या : १० ॥}}}
        
            \clearpage

            % --- SECTION ---
            
            \phantomsection  

            \section*{\centering षष्ठ खण्डः}
            \addcontentsline{toc}{section}{षष्ठ खण्डः} 
            \markright{षष्ठ खण्डः}
            \nopagebreak % Keep section title with first mantra
            
            % --- SECTION-LEVEL PROCEDURE LINK ---

                % --- SUBSECTIONS (Mantras) ---
                
                % Conditional rendering based on output_mode
                \par\filbreak

\phantomsection

\addcontentsline{toc}{subsection}{ऋक् १}

{\centering \footnotesize \textcolor{AccentPurple}{॥ नृमेध आंगिरसः इन्द्रः उष्णिक् ॥} \par}

\vspace{0.12em}

{\centering \textbf{\textcolor{AccentBlue}{ए\raisebox{0.25ex}{\accentmark{15}{\char"1CD2}}न्द्रा\raisebox{0.25ex}{\accentmark{12}{\char"0951}}नोगधिप्रिया\kern 0.05em ।\hskip 0.12em plus 0.04em minus 0.02em सत्रा\raisebox{0.25ex}{\accentmark{12}{\char"0951}}जितगोह्या\kern 0.05em ।\hskip 0.12em plus 0.04em minus 0.02em गीरि\raisebox{0.25ex}{\accentmark{15}{\char"1CD2}}र्न्ना\raisebox{0.25ex}{\accentmark{12}{\char"0951}}विश्वात\raisebox{0.25ex}{\accentmark{12}{\char"0951}}ःपृथूःपा\raisebox{0.25ex}{\accentmark{12}{\char"0951}}तीर्दीवः \nolinebreak\hskip 0.08em \mbox{॥१॥}}} \par}

\vspace{0.8em}
                % --- SUBSECTIONS (Mantras) ---
                
                % Conditional rendering based on output_mode
                \par\filbreak

\phantomsection

\addcontentsline{toc}{subsection}{ऋक् २}

{\centering \footnotesize \textcolor{AccentPurple}{॥ भरद्वाजो बार्हस्पत्यः इन्द्रः उष्णिक् ॥} \par}

\vspace{0.12em}

{\centering \textbf{\textcolor{AccentBlue}{यस्यत्य\raisebox{0.15ex}{\accentmark{12}{\char"1CF9}}च्छां\raisebox{0.15ex}{\accentmark{12}{\char"1CF9}}बरम्मा\raisebox{0.25ex}{\accentmark{15}{\char"1CD2}}दे\raisebox{0.25ex}{\accentmark{12}{\char"0951}}\kern 0.05em ।\hskip 0.12em plus 0.04em minus 0.02em दीवो\raisebox{0.25ex}{\accentmark{15}{\char"1CD2}}दा\raisebox{0.25ex}{\accentmark{15}{\char"1CD2}}सा\raisebox{0.25ex}{\accentmark{12}{\char"0951}}यारन्धा\raisebox{0.25ex}{\accentmark{15}{\char"1CD2}}यन्\kern 0.05em ।\hskip 0.12em plus 0.04em minus 0.02em आयं\raisebox{0.25ex}{\accentmark{15}{\char"1CD2}}ससो\raisebox{0.25ex}{\accentmark{15}{\char"1CD2}}मा\raisebox{0.25ex}{\accentmark{12}{\char"0951}}इन्द्रतेसूतःपीबा \nolinebreak\hskip 0.08em \mbox{॥२॥}}} \par}

\vspace{0.8em}
                % --- SUBSECTIONS (Mantras) ---
                
                % Conditional rendering based on output_mode
                \par\filbreak

\phantomsection

\addcontentsline{toc}{subsection}{ऋक् ३}

{\centering \footnotesize \textcolor{AccentPurple}{॥ प्रगाथो घौरः काण्वः इन्द्रः उष्णिक् ॥} \par}

\vspace{0.12em}

{\centering \textbf{\textcolor{AccentBlue}{गृणेता\raisebox{0.25ex}{\accentmark{15}{\char"1CD2}}दिन्द्रा\raisebox{0.25ex}{\accentmark{12}{\char"0951}}तेशा\raisebox{0.25ex}{\accentmark{15}{\char"1CD2}}वा\raisebox{0.25ex}{\accentmark{12}{\char"0951}}ः\kern 0.05em ।\hskip 0.12em plus 0.04em minus 0.02em ऊपा\raisebox{0.25ex}{\accentmark{15}{\char"1CD2}}मान्देवा\raisebox{0.25ex}{\accentmark{15}{\char"1CD2}}ता\raisebox{0.25ex}{\accentmark{12}{\char"0951}}तये\kern 0.05em ।\hskip 0.12em plus 0.04em minus 0.02em यद्धं\raisebox{0.15ex}{\accentmark{12}{\char"1CF9}}सी\raisebox{0.25ex}{\accentmark{12}{\char"0951}}वृत्र\raisebox{0.25ex}{\accentmark{15}{\char"1CD2}}मोज\raisebox{0.25ex}{\accentmark{12}{\char"0951}}सा\raisebox{0.25ex}{\accentmark{15}{\char"1CD2}}शा\raisebox{0.25ex}{\accentmark{15}{\char"1CD2}}चीपते \nolinebreak\hskip 0.08em \mbox{॥३॥}}} \par}

\vspace{0.8em}
                % --- SUBSECTIONS (Mantras) ---
                
                % Conditional rendering based on output_mode
                \par\filbreak

\phantomsection

\addcontentsline{toc}{subsection}{ऋक् ४}

{\centering \footnotesize \textcolor{AccentPurple}{॥ पर्वतः काण्वः इन्द्रः उष्णिक् ॥} \par}

\vspace{0.12em}

{\centering \textbf{\textcolor{AccentBlue}{याइ\raisebox{0.25ex}{\accentmark{15}{\char"1CD2}}न्द्रासोमापाता\raisebox{0.25ex}{\accentmark{12}{\char"0951}}मः\kern 0.05em ।\hskip 0.12em plus 0.04em minus 0.02em मा\raisebox{0.25ex}{\accentmark{15}{\char"1CD2}}द\raisebox{0.25ex}{\accentmark{12}{\char"0951}}श्शवीष्ठाचे\raisebox{0.25ex}{\accentmark{15}{\char"1CD2}}ताती\kern 0.05em ।\hskip 0.12em plus 0.04em minus 0.02em ये\raisebox{0.25ex}{\accentmark{15}{\char"1CD2}}नाहंसिन्याटट्यत्रीणंता\raisebox{0.25ex}{\accentmark{15}{\char"1CD2}}मीमाहे \nolinebreak\hskip 0.08em \mbox{॥४॥}}} \par}

\vspace{0.8em}
                % --- SUBSECTIONS (Mantras) ---
                
                % Conditional rendering based on output_mode
                \par\filbreak

\phantomsection

\addcontentsline{toc}{subsection}{ऋक् ५}

{\centering \footnotesize \textcolor{AccentPurple}{॥ इरिम्बिठः काण्वः आदित्याः उष्णिक् ॥} \par}

\vspace{0.12em}

{\centering \textbf{\textcolor{AccentBlue}{तूचेतूना\raisebox{0.15ex}{\accentmark{12}{\char"1CF9}}ःयात\raisebox{0.25ex}{\accentmark{15}{\char"1CD2}}त्सू\raisebox{0.25ex}{\accentmark{12}{\char"0951}}नाः\kern 0.05em ।\hskip 0.12em plus 0.04em minus 0.02em द्राघी\raisebox{0.25ex}{\accentmark{12}{\char"0951}}याआ\raisebox{0.15ex}{\accentmark{12}{\char"1CF9}}यू\raisebox{0.25ex}{\accentmark{12}{\char"0951}}र्जीवा\raisebox{0.25ex}{\accentmark{15}{\char"1CD2}}से\raisebox{0.25ex}{\accentmark{12}{\char"0951}}\kern 0.05em ।\hskip 0.12em plus 0.04em minus 0.02em आ\raisebox{0.25ex}{\accentmark{12}{\char"0951}}दि\raisebox{0.25ex}{\accentmark{12}{\char"0951}}त्यासस्सम्महसाःकृणो\raisebox{0.25ex}{\accentmark{15}{\char"1CD2}}ता\raisebox{0.25ex}{\accentmark{12}{\char"0951}}नः \nolinebreak\hskip 0.08em \mbox{॥५॥}}} \par}

\vspace{0.8em}
                % --- SUBSECTIONS (Mantras) ---
                
                % Conditional rendering based on output_mode
                \par\filbreak

\phantomsection

\addcontentsline{toc}{subsection}{ऋक् ६}

{\centering \footnotesize \textcolor{AccentPurple}{॥ विश्वमना वैयश्वः इन्द्रः उष्णिक् ॥} \par}

\vspace{0.12em}

{\centering \textbf{\textcolor{AccentBlue}{वे\raisebox{0.25ex}{\accentmark{15}{\char"1CD2}}त्थाहि\raisebox{0.25ex}{\accentmark{15}{\char"1CD2}}नि\raisebox{0.25ex}{\accentmark{15}{\char"1CD2}}रिऋ\raisebox{0.25ex}{\accentmark{12}{\char"0951}}तीनाम्\kern 0.05em ।\hskip 0.12em plus 0.04em minus 0.02em वज्रा\raisebox{0.25ex}{\accentmark{12}{\char"0951}}हस्तापारिव्रा\raisebox{0.25ex}{\accentmark{15}{\char"1CD2}}जम्\raisebox{0.25ex}{\accentmark{12}{\char"0951}}\kern 0.05em ।\hskip 0.12em plus 0.04em minus 0.02em आहा\raisebox{0.25ex}{\accentmark{12}{\char"0951}}रहश्शुन्धू\raisebox{0.15ex}{\accentmark{12}{\char"1CF9}}ःपा\raisebox{0.25ex}{\accentmark{12}{\char"0951}}रीपादा\raisebox{0.25ex}{\accentmark{12}{\char"0951}}मिवा \nolinebreak\hskip 0.08em \mbox{॥६॥}}} \par}

\vspace{0.8em}
                % --- SUBSECTIONS (Mantras) ---
                
                % Conditional rendering based on output_mode
                \par\filbreak

\phantomsection

\addcontentsline{toc}{subsection}{ऋक् ७}

{\centering \footnotesize \textcolor{AccentPurple}{॥ इरिम्बिठः काण्वः आदित्याः उष्णिक् ॥} \par}

\vspace{0.12em}

{\centering \textbf{\textcolor{AccentBlue}{आ\raisebox{0.15ex}{\accentmark{12}{\char"1CF8}}पा\raisebox{0.15ex}{\accentmark{12}{\char"1CF8}}मी\raisebox{0.25ex}{\accentmark{12}{\char"0951}}वा\raisebox{0.25ex}{\accentmark{15}{\char"1CD2}}मा\raisebox{0.25ex}{\accentmark{15}{\char"1CD2}}पस्री\raisebox{0.25ex}{\accentmark{15}{\char"1CD2}}धम्\kern 0.05em ।\hskip 0.12em plus 0.04em minus 0.02em आ\raisebox{0.25ex}{\accentmark{15}{\char"1CD2}}पा\raisebox{0.25ex}{\accentmark{12}{\char"0951}}सेधतदुर्मा\raisebox{0.25ex}{\accentmark{12}{\char"0951}}तिम्\kern 0.05em ।\hskip 0.12em plus 0.04em minus 0.02em आदी\raisebox{0.25ex}{\accentmark{12}{\char"0951}}त्यासोयूयोतानानो\raisebox{0.25ex}{\accentmark{15}{\char"1CD2}}अं\raisebox{0.25ex}{\accentmark{15}{\char"1CD2}}हा\raisebox{0.25ex}{\accentmark{12}{\char"0951}}सः \nolinebreak\hskip 0.08em \mbox{॥७॥}}} \par}

\vspace{0.8em}
                % --- SUBSECTIONS (Mantras) ---
                
                % Conditional rendering based on output_mode
                \par\filbreak

\phantomsection

\addcontentsline{toc}{subsection}{ऋक् ८}

{\centering \footnotesize \textcolor{AccentPurple}{॥ वसिष्ठो मैत्रावरुणिः इन्द्रः विराटुष्णिक् ॥} \par}

\vspace{0.12em}

{\centering \textbf{\textcolor{AccentBlue}{पीबासो\raisebox{0.25ex}{\accentmark{15}{\char"1CD2}}मा\raisebox{0.25ex}{\accentmark{12}{\char"0951}}मिन्द्रा\raisebox{0.25ex}{\accentmark{15}{\char"1CD2}}मन्द\raisebox{0.25ex}{\accentmark{12}{\char"0951}}तुत्वा\kern 0.05em ।\hskip 0.12em plus 0.04em minus 0.02em य\raisebox{0.15ex}{\accentmark{12}{\char"1CF8}}न्ते\raisebox{0.25ex}{\accentmark{12}{\char"0951}}सूषा\raisebox{0.25ex}{\accentmark{15}{\char"1CD2}}वाहरिय्य\raisebox{0.25ex}{\accentmark{15}{\char"1CD2}}श्वाद्रि\raisebox{0.25ex}{\accentmark{15}{\char"1CD2}}ः\kern 0.05em ।\hskip 0.12em plus 0.04em minus 0.02em सोतु\raisebox{0.25ex}{\accentmark{15}{\char"1CD2}}र्बा\raisebox{0.25ex}{\accentmark{12}{\char"0951}}हुभ्यांसू\raisebox{0.15ex}{\accentmark{12}{\char"1CF9}}यातोना\raisebox{0.25ex}{\accentmark{15}{\char"1CD2}}र्वा \nolinebreak\hskip 0.08em \mbox{॥८॥}}} \par}

\vspace{0.8em}

            % --- End of Kandah Closing Floral Line & Counts ---
            \par\nopagebreak\vspace{0.9em plus 0.3em minus 0.1em}
            \centerline{\textcolor{black!35}{\rule{1.8cm}{0.35pt}}\;\raisebox{-0.22ex}{\textcolor{vedicmaroon}{\large\SixFlowerPetalDotted}}\;\textcolor{black!35}{\rule{1.8cm}{0.35pt}}}
            \nopagebreak\vspace{0.45em}
            \centerline{\devafont\textbf{\textcolor{AccentBlue}{॥ ऋक् सङ्ख्या : ८ ॥}}}
        
            \clearpage

            % --- SECTION ---
            
            \phantomsection  

            \section*{\centering सप्तम खण्डः}
            \addcontentsline{toc}{section}{सप्तम खण्डः} 
            \markright{सप्तम खण्डः}
            \nopagebreak % Keep section title with first mantra
            
            % --- SECTION-LEVEL PROCEDURE LINK ---

                % --- SUBSECTIONS (Mantras) ---
                
                % Conditional rendering based on output_mode
                \par\filbreak

\phantomsection

\addcontentsline{toc}{subsection}{ऋक् १}

{\centering \footnotesize \textcolor{AccentPurple}{॥ सोभिरः काण्वः इन्द्रः ककुप् ॥} \par}

\vspace{0.12em}

{\centering \textbf{\textcolor{AccentBlue}{अभ्रा\raisebox{0.25ex}{\accentmark{15}{\char"1CD2}}तृ\raisebox{0.25ex}{\accentmark{15}{\char"1CD2}}व्योआना\raisebox{0.25ex}{\accentmark{15}{\char"1CD2}}त्वम्\raisebox{0.25ex}{\accentmark{15}{\char"1CD2}}\kern 0.05em ।\hskip 0.12em plus 0.04em minus 0.02em आनापिरिन्द्राजनु\raisebox{0.15ex}{\accentmark{12}{\char"1CF9}}षा\raisebox{0.25ex}{\accentmark{12}{\char"0951}}साना\raisebox{0.25ex}{\accentmark{15}{\char"1CD2}}द\raisebox{0.25ex}{\accentmark{12}{\char"0951}}सि\kern 0.05em ।\hskip 0.12em plus 0.04em minus 0.02em यू\raisebox{0.15ex}{\accentmark{12}{\char"1CF8}}थे\raisebox{0.15ex}{\accentmark{12}{\char"1CF9}}दापित्वामी\raisebox{0.25ex}{\accentmark{12}{\char"0951}}च्छासे \nolinebreak\hskip 0.08em \mbox{॥१॥}}} \par}

\vspace{0.8em}
                % --- SUBSECTIONS (Mantras) ---
                
                % Conditional rendering based on output_mode
                \par\filbreak

\phantomsection

\addcontentsline{toc}{subsection}{ऋक् २}

{\centering \footnotesize \textcolor{AccentPurple}{॥ सोभिरः काण्वः इन्द्रः ककुप् ॥} \par}

\vspace{0.12em}

{\centering \textbf{\textcolor{AccentBlue}{यो\raisebox{0.15ex}{\accentmark{12}{\char"1CF8}}ना\raisebox{0.25ex}{\accentmark{12}{\char"0951}}ईदामी\raisebox{0.25ex}{\accentmark{12}{\char"0951}}दंपूरा\kern 0.05em ।\hskip 0.12em plus 0.04em minus 0.02em प्रव\raisebox{0.25ex}{\accentmark{15}{\char"1CD2}}स्याआ\raisebox{0.25ex}{\accentmark{12}{\char"0951}}नीना\raisebox{0.25ex}{\accentmark{15}{\char"1CD2}}याता\raisebox{0.25ex}{\accentmark{12}{\char"0951}}मू\raisebox{0.25ex}{\accentmark{12}{\char"0951}}वस्तूषे\kern 0.05em ।\hskip 0.12em plus 0.04em minus 0.02em साखा\raisebox{0.25ex}{\accentmark{15}{\char"1CD2}}याइ\raisebox{0.15ex}{\accentmark{12}{\char"1CF9}}न्दामूताये \nolinebreak\hskip 0.08em \mbox{॥२॥}}} \par}

\vspace{0.8em}
                % --- SUBSECTIONS (Mantras) ---
                
                % Conditional rendering based on output_mode
                \par\filbreak

\phantomsection

\addcontentsline{toc}{subsection}{ऋक् ३}

{\centering \footnotesize \textcolor{AccentPurple}{॥ सोभिरः काण्वः मरुतः ककुप् ॥} \par}

\vspace{0.12em}

{\centering \textbf{\textcolor{AccentBlue}{आ\raisebox{0.15ex}{\accentmark{12}{\char"1CF8}}ग\raisebox{0.25ex}{\accentmark{12}{\char"0951}}न्तामारिषण्यत\kern 0.05em ।\hskip 0.12em plus 0.04em minus 0.02em प्रस्था\raisebox{0.25ex}{\accentmark{12}{\char"0951}}वानोमापा\raisebox{0.25ex}{\accentmark{12}{\char"0951}}स्थातासमन्यवः\kern 0.05em ।\hskip 0.12em plus 0.04em minus 0.02em दृढा\raisebox{0.25ex}{\accentmark{15}{\char"1CD2}}श्चिद्यामायिष्णावाः \nolinebreak\hskip 0.08em \mbox{॥३॥}}} \par}

\vspace{0.8em}
                % --- SUBSECTIONS (Mantras) ---
                
                % Conditional rendering based on output_mode
                \par\filbreak

\phantomsection

\addcontentsline{toc}{subsection}{ऋक् ४}

{\centering \textbf{\textcolor{AccentBlue}{॥ सोभिरःकाण्वःइन्द्रःककुप् ॥ ॥ सोभिरःकाण्वःइन्द्रःककुप् ॥ आ\raisebox{0.15ex}{\accentmark{12}{\char"1CF8}}या\raisebox{0.25ex}{\accentmark{12}{\char"0951}}ह्यायामिन्दा\raisebox{0.25ex}{\accentmark{12}{\char"0951}}वे\kern 0.05em ।\hskip 0.12em plus 0.04em minus 0.02em आ\raisebox{0.25ex}{\accentmark{12}{\char"0951}}श्वा\raisebox{0.25ex}{\accentmark{12}{\char"0951}}पा\raisebox{0.25ex}{\accentmark{15}{\char"1CD2}}तेगो\raisebox{0.15ex}{\accentmark{12}{\char"1CF9}}पा\raisebox{0.25ex}{\accentmark{12}{\char"0951}}ताउर्वारापते\kern 0.05em ।\hskip 0.12em plus 0.04em minus 0.02em सो\raisebox{0.25ex}{\accentmark{15}{\char"1CD2}}मं\raisebox{0.25ex}{\accentmark{12}{\char"0951}}सोमपतेपीबा\raisebox{0.25ex}{\accentmark{12}{\char"0951}} \nolinebreak\hskip 0.08em \mbox{॥४॥}}} \par}

\vspace{0.8em}
                % --- SUBSECTIONS (Mantras) ---
                
                % Conditional rendering based on output_mode
                \par\filbreak

\phantomsection

\addcontentsline{toc}{subsection}{ऋक् ५}

{\centering \textbf{\textcolor{AccentBlue}{त्वाया\raisebox{0.25ex}{\accentmark{12}{\char"0951}}हस्वीद्युजावायम्\raisebox{0.25ex}{\accentmark{15}{\char"1CD2}}\kern 0.05em ।\hskip 0.12em plus 0.04em minus 0.02em प्रातिश्व\raisebox{0.25ex}{\accentmark{15}{\char"1CD2}}स\raisebox{0.25ex}{\accentmark{12}{\char"0951}}न्तंवृषभब्रूवीमही\raisebox{0.25ex}{\accentmark{12}{\char"0951}}\kern 0.05em ।\hskip 0.12em plus 0.04em minus 0.02em संस्थे\raisebox{0.15ex}{\accentmark{12}{\char"1CF9}}ज\raisebox{0.15ex}{\accentmark{12}{\char"1CF9}}नास्यागोमा\raisebox{0.25ex}{\accentmark{12}{\char"0951}}तः \nolinebreak\hskip 0.08em \mbox{॥५॥}}} \par}

\vspace{0.8em}
                % --- SUBSECTIONS (Mantras) ---
                
                % Conditional rendering based on output_mode
                \par\filbreak

\phantomsection

\addcontentsline{toc}{subsection}{ऋक् ६}

{\centering \footnotesize \textcolor{AccentPurple}{॥ सोभिरः काण्वः मरुतः ककुप् ॥} \par}

\vspace{0.12em}

{\centering \textbf{\textcolor{AccentBlue}{गावा\raisebox{0.25ex}{\accentmark{15}{\char"1CD2}}श्चित्घासमन्यवः\kern 0.05em ।\hskip 0.12em plus 0.04em minus 0.02em सा\raisebox{0.25ex}{\accentmark{15}{\char"1CD2}}जात्येनामारूतस्साब\raisebox{0.25ex}{\accentmark{12}{\char"0951}}न्धवः\kern 0.05em ।\hskip 0.12em plus 0.04em minus 0.02em रिहाते\raisebox{0.25ex}{\accentmark{12}{\char"0951}}काकू\raisebox{0.15ex}{\accentmark{12}{\char"1CF9}}भो\raisebox{0.25ex}{\accentmark{12}{\char"0951}}मीथः \nolinebreak\hskip 0.08em \mbox{॥६॥}}} \par}

\vspace{0.8em}
                % --- SUBSECTIONS (Mantras) ---
                
                % Conditional rendering based on output_mode
                \par\filbreak

\phantomsection

\addcontentsline{toc}{subsection}{ऋक् ७}

{\centering \footnotesize \textcolor{AccentPurple}{॥ नृमेध आंगिरसः इन्द्रः ककुप् ॥} \par}

\vspace{0.12em}

{\centering \textbf{\textcolor{AccentBlue}{त्व\raisebox{0.15ex}{\accentmark{12}{\char"1CF8}}न्ना\raisebox{0.25ex}{\accentmark{12}{\char"0951}}इन्द्रा\raisebox{0.25ex}{\accentmark{15}{\char"1CD2}}भा\raisebox{0.25ex}{\accentmark{12}{\char"0951}}रा\kern 0.05em ।\hskip 0.12em plus 0.04em minus 0.02em ओजो\raisebox{0.25ex}{\accentmark{12}{\char"0951}}निर्म्णं\raisebox{0.25ex}{\accentmark{15}{\char"1CD2}}शाताक्रतोविचर्षणे\kern 0.05em ।\hskip 0.12em plus 0.04em minus 0.02em आ\raisebox{0.25ex}{\accentmark{15}{\char"1CD2}}वीरं\raisebox{0.25ex}{\accentmark{15}{\char"1CD2}}पृ\raisebox{0.25ex}{\accentmark{12}{\char"0951}}तनासा\raisebox{0.25ex}{\accentmark{15}{\char"1CD2}}हम् \nolinebreak\hskip 0.08em \mbox{॥७॥}}} \par}

\vspace{0.8em}
                % --- SUBSECTIONS (Mantras) ---
                
                % Conditional rendering based on output_mode
                \par\filbreak

\phantomsection

\addcontentsline{toc}{subsection}{ऋक् ८}

{\centering \footnotesize \textcolor{AccentPurple}{॥ नृमेध आंगिरसः इन्द्रः ककुप् ॥} \par}

\vspace{0.12em}

{\centering \textbf{\textcolor{AccentBlue}{आ\raisebox{0.25ex}{\accentmark{15}{\char"1CD2}}थाहीन्द्रगिर्वणः\kern 0.05em ।\hskip 0.12em plus 0.04em minus 0.02em ऊपा\raisebox{0.25ex}{\accentmark{12}{\char"0951}}त्वाका\raisebox{0.15ex}{\accentmark{12}{\char"1CF9}}मा\raisebox{0.25ex}{\accentmark{12}{\char"0951}}ईमा\raisebox{0.25ex}{\accentmark{15}{\char"1CD2}}हे\raisebox{0.25ex}{\accentmark{12}{\char"0951}}सासृग्माहेऊदेवग्म\raisebox{0.15ex}{\accentmark{12}{\char"1CF9}}न्ता\raisebox{0.25ex}{\accentmark{12}{\char"0951}}ऊद\raisebox{0.25ex}{\accentmark{15}{\char"1CD2}}भीः \nolinebreak\hskip 0.08em \mbox{॥८॥}}} \par}

\vspace{0.8em}
                % --- SUBSECTIONS (Mantras) ---
                
                % Conditional rendering based on output_mode
                \par\filbreak

\phantomsection

\addcontentsline{toc}{subsection}{ऋक् ९}

{\centering \footnotesize \textcolor{AccentPurple}{॥ सोभिरः काण्वः इन्द्रः ककुप् ॥} \par}

\vspace{0.12em}

{\centering \textbf{\textcolor{AccentBlue}{सी\raisebox{0.25ex}{\accentmark{15}{\char"1CD2}}दन्त\raisebox{0.25ex}{\accentmark{12}{\char"0951}}स्ते\raisebox{0.15ex}{\accentmark{12}{\char"1CF9}}वायोया\raisebox{0.25ex}{\accentmark{15}{\char"1CD2}}था\kern 0.05em ।\hskip 0.12em plus 0.04em minus 0.02em गो\raisebox{0.25ex}{\accentmark{15}{\char"1CD2}}श्री\raisebox{0.25ex}{\accentmark{15}{\char"1CD2}}स्तेमा\raisebox{0.25ex}{\accentmark{15}{\char"1CD2}}धोर्मा\raisebox{0.25ex}{\accentmark{15}{\char"1CD2}}दीरोवी\raisebox{0.25ex}{\accentmark{12}{\char"0951}}वक्षणः\kern 0.05em ।\hskip 0.12em plus 0.04em minus 0.02em आ\raisebox{0.25ex}{\accentmark{15}{\char"1CD2}}भित्वा\raisebox{0.25ex}{\accentmark{15}{\char"1CD2}}मिन्द्रनोनुमः \nolinebreak\hskip 0.08em \mbox{॥९॥}}} \par}

\vspace{0.8em}
                % --- SUBSECTIONS (Mantras) ---
                
                % Conditional rendering based on output_mode
                \par\filbreak

\phantomsection

\addcontentsline{toc}{subsection}{ऋक् १०}

{\centering \footnotesize \textcolor{AccentPurple}{॥ सोभिरः काण्वः इन्द्रः ककुप् ॥} \par}

\vspace{0.12em}

{\centering \textbf{\textcolor{AccentBlue}{वा\raisebox{0.25ex}{\accentmark{15}{\char"1CD2}}या\raisebox{0.25ex}{\accentmark{15}{\char"1CD2}}मुत्वामा\raisebox{0.25ex}{\accentmark{12}{\char"0951}}पूर्व्या\raisebox{0.25ex}{\accentmark{12}{\char"0951}}\kern 0.05em ।\hskip 0.12em plus 0.04em minus 0.02em स्थूरन्न\raisebox{0.25ex}{\accentmark{15}{\char"1CD2}}का\raisebox{0.25ex}{\accentmark{15}{\char"1CD2}}च्चिद्भा\raisebox{0.25ex}{\accentmark{15}{\char"1CD2}}र\raisebox{0.25ex}{\accentmark{12}{\char"0951}}न्तोवस्या\raisebox{0.25ex}{\accentmark{12}{\char"0951}}वा\raisebox{0.15ex}{\accentmark{12}{\char"1CF8}}ः\kern 0.05em ।\hskip 0.12em plus 0.04em minus 0.02em वज्रि\raisebox{0.25ex}{\accentmark{15}{\char"1CD2}}ञ्चित्रंहा\raisebox{0.25ex}{\accentmark{12}{\char"0951}}वामहे \nolinebreak\hskip 0.08em \mbox{॥१०॥}}} \par}

\vspace{0.8em}

            % --- End of Kandah Closing Floral Line & Counts ---
            \par\nopagebreak\vspace{0.9em plus 0.3em minus 0.1em}
            \centerline{\textcolor{black!35}{\rule{1.8cm}{0.35pt}}\;\raisebox{-0.22ex}{\textcolor{vedicmaroon}{\large\SixFlowerPetalDotted}}\;\textcolor{black!35}{\rule{1.8cm}{0.35pt}}}
            \nopagebreak\vspace{0.45em}
            \centerline{\devafont\textbf{\textcolor{AccentBlue}{॥ ऋक् सङ्ख्या : १० ॥}}}
        
            \clearpage

            % --- SECTION ---
            
            \phantomsection  

            \section*{\centering अष्टम खण्डः}
            \addcontentsline{toc}{section}{अष्टम खण्डः} 
            \markright{अष्टम खण्डः}
            \nopagebreak % Keep section title with first mantra
            
            % --- SECTION-LEVEL PROCEDURE LINK ---

                % --- SUBSECTIONS (Mantras) ---
                
                % Conditional rendering based on output_mode
                \par\filbreak

\phantomsection

\addcontentsline{toc}{subsection}{ऋक् १}

{\centering \footnotesize \textcolor{AccentPurple}{॥ त्रसदस्युः इन्द्रः द्विपदा विराट् ॥} \par}

\vspace{0.12em}

{\centering \textbf{\textcolor{AccentBlue}{विश्वा\raisebox{0.25ex}{\accentmark{12}{\char"0951}}तोदाव\raisebox{0.25ex}{\accentmark{15}{\char"1CD2}}न्विश्वा\raisebox{0.15ex}{\accentmark{12}{\char"1CF9}}तो\raisebox{0.25ex}{\accentmark{12}{\char"0951}}नआभा\raisebox{0.25ex}{\accentmark{12}{\char"0951}}रा\kern 0.05em ।\hskip 0.12em plus 0.04em minus 0.02em य\raisebox{0.25ex}{\accentmark{15}{\char"1CD2}}न्त्वाशा\raisebox{0.15ex}{\accentmark{12}{\char"1CF9}}विष्ठामी\raisebox{0.25ex}{\accentmark{15}{\char"1CD2}}मा\raisebox{0.25ex}{\accentmark{12}{\char"0951}}हे\kern 0.05em ।\hskip 0.12em plus 0.04em minus 0.02em \nolinebreak\hskip 0.08em \mbox{॥१॥}}} \par}

\vspace{0.8em}
                % --- SUBSECTIONS (Mantras) ---
                
                % Conditional rendering based on output_mode
                \par\filbreak

\phantomsection

\addcontentsline{toc}{subsection}{ऋक् २}

{\centering \footnotesize \textcolor{AccentPurple}{॥ त्रसदस्युः इन्द्रः द्विपदा विराट् ॥} \par}

\vspace{0.12em}

{\centering \textbf{\textcolor{AccentBlue}{एष\raisebox{0.25ex}{\accentmark{15}{\char"1CD2}}ब्रह्मा\raisebox{0.25ex}{\accentmark{15}{\char"1CD2}}याऋत्वी\raisebox{0.25ex}{\accentmark{15}{\char"1CD2}}याः\kern 0.05em ।\hskip 0.12em plus 0.04em minus 0.02em इन्द्रो\raisebox{0.25ex}{\accentmark{15}{\char"1CD2}}ना\raisebox{0.25ex}{\accentmark{15}{\char"1CD2}}माश्रू\raisebox{0.15ex}{\accentmark{12}{\char"1CF9}}तो\raisebox{0.25ex}{\accentmark{12}{\char"0951}}गृणे\raisebox{0.25ex}{\accentmark{15}{\char"1CD2}} \nolinebreak\hskip 0.08em \mbox{॥२॥}}} \par}

\vspace{0.8em}
                % --- SUBSECTIONS (Mantras) ---
                
                % Conditional rendering based on output_mode
                \par\filbreak

\phantomsection

\addcontentsline{toc}{subsection}{ऋक् ३}

{\centering \textbf{\textcolor{AccentBlue}{॥ त्रसदस्युःइन्द्रःद्विपदाविराट् ॥ ॥ त्रसदस्युःइन्द्रःद्विपदाविराट् ॥ ब्रह्माणाइ\raisebox{0.25ex}{\accentmark{12}{\char"0951}}न्द्रं\raisebox{0.25ex}{\accentmark{12}{\char"0951}}मा\raisebox{0.25ex}{\accentmark{12}{\char"0951}}हा\raisebox{0.25ex}{\accentmark{15}{\char"1CD2}}य\raisebox{0.25ex}{\accentmark{12}{\char"0951}}न्तोअर्कैः\kern 0.05em ।\hskip 0.12em plus 0.04em minus 0.02em आ\raisebox{0.25ex}{\accentmark{15}{\char"1CD2}}वा\raisebox{0.25ex}{\accentmark{15}{\char"1CD2}}र्धायन्ना\raisebox{0.15ex}{\accentmark{12}{\char"1CF9}}हा\raisebox{0.25ex}{\accentmark{12}{\char"0951}}येहं\raisebox{0.25ex}{\accentmark{15}{\char"1CD2}}तावा\raisebox{0.25ex}{\accentmark{15}{\char"1CD2}}ऊ\raisebox{0.25ex}{\accentmark{12}{\char"0951}} \nolinebreak\hskip 0.08em \mbox{॥३॥}}} \par}

\vspace{0.8em}
                % --- SUBSECTIONS (Mantras) ---
                
                % Conditional rendering based on output_mode
                \par\filbreak

\phantomsection

\addcontentsline{toc}{subsection}{ऋक् ४}

{\centering \textbf{\textcolor{AccentBlue}{आ\raisebox{0.25ex}{\accentmark{15}{\char"1CD2}}नावस्ते\raisebox{0.25ex}{\accentmark{15}{\char"1CD2}}राथामश्वा\raisebox{0.25ex}{\accentmark{15}{\char"1CD2}}यतक्षुः\kern 0.05em ।\hskip 0.12em plus 0.04em minus 0.02em त्वष्टव\raisebox{0.25ex}{\accentmark{15}{\char"1CD2}}ज्रंपूरुहूताद्युम\raisebox{0.25ex}{\accentmark{15}{\char"1CD2}}न्तम् \nolinebreak\hskip 0.08em \mbox{॥४॥}}} \par}

\vspace{0.8em}
                % --- SUBSECTIONS (Mantras) ---
                
                % Conditional rendering based on output_mode
                \par\filbreak

\phantomsection

\addcontentsline{toc}{subsection}{ऋक् ५}

{\centering \textbf{\textcolor{AccentBlue}{॥ त्रसदस्युःइन्द्रःद्विपदाविराट् ॥ ॥ त्रसदस्युःविश्वेदेवाःद्विपदाविराट् ॥ शं\raisebox{0.25ex}{\accentmark{15}{\char"1CD2}}पादम्माघं\raisebox{0.15ex}{\accentmark{12}{\char"1CF9}}रायी\raisebox{0.25ex}{\accentmark{12}{\char"0951}}षीणे\kern 0.05em ।\hskip 0.12em plus 0.04em minus 0.02em नाका\raisebox{0.25ex}{\accentmark{15}{\char"1CD2}}मामाव्रातो\raisebox{0.25ex}{\accentmark{15}{\char"1CD2}}ही\raisebox{0.25ex}{\accentmark{12}{\char"0951}}नोतिना\raisebox{0.25ex}{\accentmark{15}{\char"1CD2}}स्पृशाद्रायिम् \nolinebreak\hskip 0.08em \mbox{॥५॥}}} \par}

\vspace{0.8em}
                % --- SUBSECTIONS (Mantras) ---
                
                % Conditional rendering based on output_mode
                \par\filbreak

\phantomsection

\addcontentsline{toc}{subsection}{ऋक् ६}

{\centering \textbf{\textcolor{AccentBlue}{सादागावश्शूचायोविश्वा\raisebox{0.25ex}{\accentmark{15}{\char"1CD2}}धा\raisebox{0.25ex}{\accentmark{12}{\char"0951}}यसः\kern 0.05em ।\hskip 0.12em plus 0.04em minus 0.02em सादा\raisebox{0.25ex}{\accentmark{12}{\char"0951}}देवा\raisebox{0.15ex}{\accentmark{12}{\char"1CF9}}आरे\raisebox{0.25ex}{\accentmark{15}{\char"1CD2}}पासा\raisebox{0.25ex}{\accentmark{12}{\char"0951}}ः \nolinebreak\hskip 0.08em \mbox{॥६॥}}} \par}

\vspace{0.8em}
                % --- SUBSECTIONS (Mantras) ---
                
                % Conditional rendering based on output_mode
                \par\filbreak

\phantomsection

\addcontentsline{toc}{subsection}{ऋक् ७}

{\centering \footnotesize \textcolor{AccentPurple}{॥ सम्वर्त आङ्गिरसः उषाः द्विपदा विराट् ॥} \par}

\vspace{0.12em}

{\centering \textbf{\textcolor{AccentBlue}{आ\raisebox{0.25ex}{\accentmark{15}{\char"1CD2}}याहीवा\raisebox{0.25ex}{\accentmark{15}{\char"1CD2}}नासासाहा\raisebox{0.25ex}{\accentmark{15}{\char"1CD2}}\kern 0.05em ।\hskip 0.12em plus 0.04em minus 0.02em गा\raisebox{0.25ex}{\accentmark{15}{\char"1CD2}}वा\raisebox{0.25ex}{\accentmark{12}{\char"0951}}स्सचन्तवा\raisebox{0.25ex}{\accentmark{12}{\char"0951}}र्तानींयादू\raisebox{0.25ex}{\accentmark{15}{\char"1CD2}}धा\raisebox{0.25ex}{\accentmark{12}{\char"0951}}भिः \nolinebreak\hskip 0.08em \mbox{॥७॥}}} \par}

\vspace{0.8em}
                % --- SUBSECTIONS (Mantras) ---
                
                % Conditional rendering based on output_mode
                \par\filbreak

\phantomsection

\addcontentsline{toc}{subsection}{ऋक् ८}

{\centering \footnotesize \textcolor{AccentPurple}{॥ त्रसदस्युः इन्द्रः द्विपदा विराट् ॥} \par}

\vspace{0.12em}

{\centering \textbf{\textcolor{AccentBlue}{उपप्रक्षेमा\raisebox{0.25ex}{\accentmark{15}{\char"1CD2}}धू\raisebox{0.25ex}{\accentmark{12}{\char"0951}}मतिक्षी\raisebox{0.25ex}{\accentmark{12}{\char"0951}}यन्ताः\kern 0.05em ।\hskip 0.12em plus 0.04em minus 0.02em पु\raisebox{0.25ex}{\accentmark{15}{\char"1CD2}}ष्ये\raisebox{0.25ex}{\accentmark{12}{\char"0951}}मारायिन्धीमा\raisebox{0.25ex}{\accentmark{15}{\char"1CD2}}हे\raisebox{0.25ex}{\accentmark{12}{\char"0951}}तइन्द्र \nolinebreak\hskip 0.08em \mbox{॥८॥}}} \par}

\vspace{0.8em}
                % --- SUBSECTIONS (Mantras) ---
                
                % Conditional rendering based on output_mode
                \par\filbreak

\phantomsection

\addcontentsline{toc}{subsection}{ऋक् ९}

{\centering \footnotesize \textcolor{AccentPurple}{॥ त्रसदस्युः इन्द्रः द्विपदा विराट् ॥} \par}

\vspace{0.12em}

{\centering \textbf{\textcolor{AccentBlue}{अ\raisebox{0.15ex}{\accentmark{12}{\char"1CF8}}र्चन्त्यर्कम्मारू\raisebox{0.15ex}{\accentmark{12}{\char"1CF9}}ता\raisebox{0.25ex}{\accentmark{12}{\char"0951}}स्त्वर्कः\kern 0.05em ।\hskip 0.12em plus 0.04em minus 0.02em आ\raisebox{0.25ex}{\accentmark{15}{\char"1CD2}}स्तोभतीश्रूतो\raisebox{0.25ex}{\accentmark{15}{\char"1CD2}}यूवासइ\raisebox{0.25ex}{\accentmark{15}{\char"1CD2}}न्द्राः \nolinebreak\hskip 0.08em \mbox{॥९॥}}} \par}

\vspace{0.8em}
                % --- SUBSECTIONS (Mantras) ---
                
                % Conditional rendering based on output_mode
                \par\filbreak

\phantomsection

\addcontentsline{toc}{subsection}{ऋक् १०}

{\centering \footnotesize \textcolor{AccentPurple}{॥ त्रसदस्युः इन्द्रः द्विपदा विराट् ॥} \par}

\vspace{0.12em}

{\centering \textbf{\textcolor{AccentBlue}{प्रा\raisebox{0.25ex}{\accentmark{15}{\char"1CD2}}वाइ\raisebox{0.25ex}{\accentmark{15}{\char"1CD2}}न्द्रायवृत्राह\raisebox{0.25ex}{\accentmark{15}{\char"1CD2}}न्ता\raisebox{0.25ex}{\accentmark{12}{\char"0951}}माया\kern 0.05em ।\hskip 0.12em plus 0.04em minus 0.02em वि\raisebox{0.25ex}{\accentmark{15}{\char"1CD2}}प्रा\raisebox{0.25ex}{\accentmark{12}{\char"0951}}या\raisebox{0.25ex}{\accentmark{15}{\char"1CD2}}गा\raisebox{0.25ex}{\accentmark{12}{\char"0951}}थंगायाताय\raisebox{0.25ex}{\accentmark{15}{\char"1CD2}}ञ्जुजो\raisebox{0.25ex}{\accentmark{15}{\char"1CD2}}षाते \nolinebreak\hskip 0.08em \mbox{॥१०॥}}} \par}

\vspace{0.8em}

            % --- End of Kandah Closing Floral Line & Counts ---
            \par\nopagebreak\vspace{0.9em plus 0.3em minus 0.1em}
            \centerline{\textcolor{black!35}{\rule{1.8cm}{0.35pt}}\;\raisebox{-0.22ex}{\textcolor{vedicmaroon}{\large\SixFlowerPetalDotted}}\;\textcolor{black!35}{\rule{1.8cm}{0.35pt}}}
            \nopagebreak\vspace{0.45em}
            \centerline{\devafont\textbf{\textcolor{AccentBlue}{॥ ऋक् सङ्ख्या : १० ॥}}}
        
            \clearpage

            % --- SECTION ---
            
            \phantomsection  

            \section*{\centering नवम खण्डः}
            \addcontentsline{toc}{section}{नवम खण्डः} 
            \markright{नवम खण्डः}
            \nopagebreak % Keep section title with first mantra
            
            % --- SECTION-LEVEL PROCEDURE LINK ---

                % --- SUBSECTIONS (Mantras) ---
                
                % Conditional rendering based on output_mode
                \par\filbreak

\phantomsection

\addcontentsline{toc}{subsection}{ऋक् १}

{\centering \footnotesize \textcolor{AccentPurple}{॥ पृषध्रः काण्वः अग्निः द्विपदा विराट् ॥} \par}

\vspace{0.12em}

{\centering \textbf{\textcolor{AccentBlue}{आ\raisebox{0.15ex}{\accentmark{12}{\char"1CF8}}चे\raisebox{0.25ex}{\accentmark{12}{\char"0951}}त्याग्नी\raisebox{0.25ex}{\accentmark{15}{\char"1CD2}}श्ची\raisebox{0.25ex}{\accentmark{15}{\char"1CD2}}की\raisebox{0.25ex}{\accentmark{12}{\char"0951}}तिः\kern 0.05em ।\hskip 0.12em plus 0.04em minus 0.02em ह\raisebox{0.25ex}{\accentmark{15}{\char"1CD2}}व्यावाण्णासू\raisebox{0.15ex}{\accentmark{12}{\char"1CF8}}मा\raisebox{0.25ex}{\accentmark{15}{\char"1CD2}}द्रा\raisebox{0.25ex}{\accentmark{12}{\char"0951}}थः \nolinebreak\hskip 0.08em \mbox{॥१॥}}} \par}

\vspace{0.8em}
                % --- SUBSECTIONS (Mantras) ---
                
                % Conditional rendering based on output_mode
                \par\filbreak

\phantomsection

\addcontentsline{toc}{subsection}{ऋक् २}

{\centering \footnotesize \textcolor{AccentPurple}{॥ बन्धुः सुबन्धुः श्रुतबन्धुर्विप्रबन्धुश्च क्रमेण गोपायना लोपायना वा अग्निः द्विपदा विराट् ॥} \par}

\vspace{0.12em}

{\centering \textbf{\textcolor{AccentBlue}{अ\raisebox{0.25ex}{\accentmark{15}{\char"1CD2}}ग्नेत्व\raisebox{0.25ex}{\accentmark{15}{\char"1CD2}}न्नो\raisebox{0.25ex}{\accentmark{12}{\char"0951}}अन्ता\raisebox{0.25ex}{\accentmark{12}{\char"0951}}मः\kern 0.05em ।\hskip 0.12em plus 0.04em minus 0.02em ऊता\raisebox{0.25ex}{\accentmark{15}{\char"1CD2}}त्राता\raisebox{0.25ex}{\accentmark{15}{\char"1CD2}}शीवो\raisebox{0.25ex}{\accentmark{15}{\char"1CD2}}भू\raisebox{0.25ex}{\accentmark{12}{\char"0951}}वोवारूत्थ्याःठ\raisebox{0.25ex}{\accentmark{12}{\char"0951}} \nolinebreak\hskip 0.08em \mbox{॥२॥}}} \par}

\vspace{0.8em}
                % --- SUBSECTIONS (Mantras) ---
                
                % Conditional rendering based on output_mode
                \par\filbreak

\phantomsection

\addcontentsline{toc}{subsection}{ऋक् ३}

{\centering \footnotesize \textcolor{AccentPurple}{॥ बन्धुः सुबन्धुः श्रुतबन्धुर्विप्रबन्धुश्च क्रमेण गोपायना लोपायना वा इन्द्रः द्विपदा विराट् ॥} \par}

\vspace{0.12em}

{\centering \textbf{\textcolor{AccentBlue}{भा\raisebox{0.25ex}{\accentmark{15}{\char"1CD2}}गोनाचित्रः\kern 0.05em ।\hskip 0.12em plus 0.04em minus 0.02em अग्निर्मा\raisebox{0.25ex}{\accentmark{15}{\char"1CD2}}हो\raisebox{0.25ex}{\accentmark{15}{\char"1CD2}}नाम्\kern 0.05em ।\hskip 0.12em plus 0.04em minus 0.02em द\raisebox{0.15ex}{\accentmark{12}{\char"1CF8}}धा\raisebox{0.25ex}{\accentmark{12}{\char"0951}}तीरा\raisebox{0.25ex}{\accentmark{15}{\char"1CD2}}त्नम् \nolinebreak\hskip 0.08em \mbox{॥३॥}}} \par}

\vspace{0.8em}
                % --- SUBSECTIONS (Mantras) ---
                
                % Conditional rendering based on output_mode
                \par\filbreak

\phantomsection

\addcontentsline{toc}{subsection}{ऋक् ४}

{\centering \footnotesize \textcolor{AccentPurple}{॥ बन्धुः सुबन्धुः श्रुतबन्धुर्विप्रबन्धुश्च क्रमेण गोपायना लोपायना वा इन्द्रः द्विपदा विराट् ॥} \par}

\vspace{0.12em}

{\centering \textbf{\textcolor{AccentBlue}{वि\raisebox{0.15ex}{\accentmark{12}{\char"1CF8}}श्व\raisebox{0.25ex}{\accentmark{12}{\char"0951}}स्यप्रस्तो\raisebox{0.25ex}{\accentmark{12}{\char"0951}}भा\kern 0.05em ।\hskip 0.12em plus 0.04em minus 0.02em पूरो\raisebox{0.25ex}{\accentmark{15}{\char"1CD2}}वासन्या\raisebox{0.15ex}{\accentmark{12}{\char"1CF9}}दीवेहा\raisebox{0.25ex}{\accentmark{15}{\char"1CD2}}नूनम् \nolinebreak\hskip 0.08em \mbox{॥४॥}}} \par}

\vspace{0.8em}
                % --- SUBSECTIONS (Mantras) ---
                
                % Conditional rendering based on output_mode
                \par\filbreak

\phantomsection

\addcontentsline{toc}{subsection}{ऋक् ५}

{\centering \footnotesize \textcolor{AccentPurple}{॥ सम्वर्त आङ्गिरसः उषाः द्विपदा विराट् ॥} \par}

\vspace{0.12em}

{\centering \textbf{\textcolor{AccentBlue}{ऊषा\raisebox{0.25ex}{\accentmark{15}{\char"1CD2}}आपा\raisebox{0.25ex}{\accentmark{15}{\char"1CD2}}स्वासुष्टा\raisebox{0.25ex}{\accentmark{15}{\char"1CD2}}माः\kern 0.05em ।\hskip 0.12em plus 0.04em minus 0.02em सं\raisebox{0.25ex}{\accentmark{15}{\char"1CD2}}वा\raisebox{0.25ex}{\accentmark{12}{\char"0951}}र्तयतिवार्क्तानिं\raisebox{0.15ex}{\accentmark{12}{\char"1CF9}}सूजाता\raisebox{0.25ex}{\accentmark{15}{\char"1CD2}}ता\raisebox{0.25ex}{\accentmark{12}{\char"0951}} \nolinebreak\hskip 0.08em \mbox{॥५॥}}} \par}

\vspace{0.8em}
                % --- SUBSECTIONS (Mantras) ---
                
                % Conditional rendering based on output_mode
                \par\filbreak

\phantomsection

\addcontentsline{toc}{subsection}{ऋक् ६}

{\centering \footnotesize \textcolor{AccentPurple}{॥ भुवन आप्त्यः साधनो वा भौवनः विश्वेदेवाः द्विपदा विराट् ॥} \par}

\vspace{0.12em}

{\centering \textbf{\textcolor{AccentBlue}{इमा\raisebox{0.25ex}{\accentmark{12}{\char"0951}}नूकं\raisebox{0.25ex}{\accentmark{12}{\char"0951}}भू\raisebox{0.25ex}{\accentmark{15}{\char"1CD2}}वा\raisebox{0.25ex}{\accentmark{12}{\char"0951}}नासिषधेम\kern 0.05em ।\hskip 0.12em plus 0.04em minus 0.02em इन्द्रा\raisebox{0.25ex}{\accentmark{12}{\char"0951}}श्वावि\raisebox{0.25ex}{\accentmark{15}{\char"1CD2}}श्वे\raisebox{0.25ex}{\accentmark{12}{\char"0951}}चादेवा\raisebox{0.25ex}{\accentmark{15}{\char"1CD2}}ः \nolinebreak\hskip 0.08em \mbox{॥६॥}}} \par}

\vspace{0.8em}
                % --- SUBSECTIONS (Mantras) ---
                
                % Conditional rendering based on output_mode
                \par\filbreak

\phantomsection

\addcontentsline{toc}{subsection}{ऋक् ७}

{\centering \footnotesize \textcolor{AccentPurple}{॥ आत्रेयः विश्वेदेवाः द्विपदा विराट् ॥} \par}

\vspace{0.12em}

{\centering \textbf{\textcolor{AccentBlue}{ऊर्जा\raisebox{0.25ex}{\accentmark{15}{\char"1CD2}}मित्रोवारू\raisebox{0.25ex}{\accentmark{12}{\char"0951}}णोपिन्वा\raisebox{0.25ex}{\accentmark{15}{\char"1CD2}}ताईळाठ\raisebox{0.25ex}{\accentmark{12}{\char"0951}}\kern 0.05em ।\hskip 0.12em plus 0.04em minus 0.02em पी\raisebox{0.15ex}{\accentmark{12}{\char"1CF8}}वा\raisebox{0.25ex}{\accentmark{12}{\char"0951}}रीमीषं\raisebox{0.25ex}{\accentmark{12}{\char"0951}}कृणू\raisebox{0.25ex}{\accentmark{12}{\char"0951}}हीना\raisebox{0.25ex}{\accentmark{12}{\char"0951}}इन्द्रा \nolinebreak\hskip 0.08em \mbox{॥७॥}}} \par}

\vspace{0.8em}
                % --- SUBSECTIONS (Mantras) ---
                
                % Conditional rendering based on output_mode
                \par\filbreak

\phantomsection

\addcontentsline{toc}{subsection}{ऋक् ८}

{\centering \footnotesize \textcolor{AccentPurple}{॥ कवत्र ऐलूवः विश्वेदेवाः द्विपदा विराट् ॥} \par}

\vspace{0.12em}

{\centering \textbf{\textcolor{AccentBlue}{विस्रूतायो\raisebox{0.15ex}{\accentmark{12}{\char"1CF9}}या\raisebox{0.15ex}{\accentmark{12}{\char"1CF9}}था\raisebox{0.25ex}{\accentmark{12}{\char"0951}}पा\raisebox{0.25ex}{\accentmark{15}{\char"1CD2}}था\raisebox{0.25ex}{\accentmark{12}{\char"0951}}\kern 0.05em ।\hskip 0.12em plus 0.04em minus 0.02em इन्द्रत्वद्य\raisebox{0.25ex}{\accentmark{12}{\char"0951}}न्तूराताया\raisebox{0.25ex}{\accentmark{12}{\char"0951}}ः \nolinebreak\hskip 0.08em \mbox{॥८॥}}} \par}

\vspace{0.8em}
                % --- SUBSECTIONS (Mantras) ---
                
                % Conditional rendering based on output_mode
                \par\filbreak

\phantomsection

\addcontentsline{toc}{subsection}{ऋक् ९}

{\centering \footnotesize \textcolor{AccentPurple}{॥ भरद्वाजो बार्हस्पत्यः इन्द्रः द्विपदा विराट् ॥} \par}

\vspace{0.12em}

{\centering \textbf{\textcolor{AccentBlue}{आया\raisebox{0.15ex}{\accentmark{12}{\char"1CF9}}वा\raisebox{0.15ex}{\accentmark{12}{\char"1CF9}}जं\raisebox{0.25ex}{\accentmark{12}{\char"0951}}देवा\raisebox{0.25ex}{\accentmark{15}{\char"1CD2}}ही\raisebox{0.25ex}{\accentmark{12}{\char"0951}}तंसनेम\raisebox{0.25ex}{\accentmark{15}{\char"1CD2}}\kern 0.05em ।\hskip 0.12em plus 0.04em minus 0.02em मा\raisebox{0.25ex}{\accentmark{12}{\char"0951}}देमाश्शा\raisebox{0.25ex}{\accentmark{15}{\char"1CD2}}ता\raisebox{0.25ex}{\accentmark{15}{\char"1CD2}}हीमासूवी\raisebox{0.25ex}{\accentmark{15}{\char"1CD2}}रा\raisebox{0.25ex}{\accentmark{12}{\char"0951}}ः \nolinebreak\hskip 0.08em \mbox{॥९॥}}} \par}

\vspace{0.8em}
                % --- SUBSECTIONS (Mantras) ---
                
                % Conditional rendering based on output_mode
                \par\filbreak

\phantomsection

\addcontentsline{toc}{subsection}{ऋक् १०}

{\centering \footnotesize \textcolor{AccentPurple}{॥ वसिष्ठो मैत्रावरुणिः इन्द्रः एकपदा ॥} \par}

\vspace{0.12em}

{\centering \textbf{\textcolor{AccentBlue}{इन्द्रोवि\raisebox{0.25ex}{\accentmark{15}{\char"1CD2}}श्वास्यराजति \nolinebreak\hskip 0.08em \mbox{॥१०॥}}} \par}

\vspace{0.8em}

            % --- End of Kandah Closing Floral Line & Counts ---
            \par\nopagebreak\vspace{0.9em plus 0.3em minus 0.1em}
            \centerline{\textcolor{black!35}{\rule{1.8cm}{0.35pt}}\;\raisebox{-0.22ex}{\textcolor{vedicmaroon}{\large\SixFlowerPetalDotted}}\;\textcolor{black!35}{\rule{1.8cm}{0.35pt}}}
            \nopagebreak\vspace{0.45em}
            \centerline{\devafont\textbf{\textcolor{AccentBlue}{॥ ऋक् सङ्ख्या : १० ॥}}}
        
            \clearpage

            % --- SECTION ---
            
            \phantomsection  

            \section*{\centering दशम खण्डः}
            \addcontentsline{toc}{section}{दशम खण्डः} 
            \markright{दशम खण्डः}
            \nopagebreak % Keep section title with first mantra
            
            % --- SECTION-LEVEL PROCEDURE LINK ---

                % --- SUBSECTIONS (Mantras) ---
                
                % Conditional rendering based on output_mode
                \par\filbreak

\phantomsection

\addcontentsline{toc}{subsection}{ऋक् १}

{\centering \footnotesize \textcolor{AccentPurple}{॥ गृत्समद शौनकः इन्द्रः अष्टिः ॥} \par}

\vspace{0.12em}

{\centering \textbf{\textcolor{AccentBlue}{त्रीका\raisebox{0.25ex}{\accentmark{15}{\char"1CD2}}दूकेषुमा\raisebox{0.25ex}{\accentmark{15}{\char"1CD2}}हीषो\raisebox{0.25ex}{\accentmark{15}{\char"1CD2}}या\raisebox{0.25ex}{\accentmark{15}{\char"1CD2}}वा\raisebox{0.25ex}{\accentmark{12}{\char"0951}}शीर\raisebox{0.25ex}{\accentmark{12}{\char"0951}}न्तूवी\raisebox{0.25ex}{\accentmark{15}{\char"1CD2}}शूष्मः\kern 0.05em ।\hskip 0.12em plus 0.04em minus 0.02em त्रम्पत्सो\raisebox{0.25ex}{\accentmark{15}{\char"1CD2}}मामपीबद्वीष्णू\raisebox{0.25ex}{\accentmark{12}{\char"0951}}ना\raisebox{0.25ex}{\accentmark{15}{\char"1CD2}}सूतंय्या\raisebox{0.25ex}{\accentmark{12}{\char"0951}}थावाशम्\kern 0.05em ।\hskip 0.12em plus 0.04em minus 0.02em साईं\raisebox{0.25ex}{\accentmark{12}{\char"0951}}ममादमाहीकर्मा\raisebox{0.25ex}{\accentmark{12}{\char"0951}}का\raisebox{0.25ex}{\accentmark{15}{\char"1CD2}}र्तावेमाहा\raisebox{0.25ex}{\accentmark{15}{\char"1CD2}}मूरुम्\raisebox{0.25ex}{\accentmark{15}{\char"1CD2}}\kern 0.05em ।\hskip 0.12em plus 0.04em minus 0.02em सै\raisebox{0.25ex}{\accentmark{15}{\char"1CD2}}नं\raisebox{0.25ex}{\accentmark{12}{\char"0951}}सश्चाद्देवो\raisebox{0.25ex}{\accentmark{15}{\char"1CD2}}देवंसत्य\raisebox{0.15ex}{\accentmark{12}{\char"1CF9}}इ\raisebox{0.15ex}{\accentmark{12}{\char"1CF9}}न्दू\raisebox{0.25ex}{\accentmark{12}{\char"0951}}ःसत्या\raisebox{0.25ex}{\accentmark{15}{\char"1CD2}}मिन्द्राः \nolinebreak\hskip 0.08em \mbox{॥१॥}}} \par}

\vspace{0.8em}
                % --- SUBSECTIONS (Mantras) ---
                
                % Conditional rendering based on output_mode
                \par\filbreak

\phantomsection

\addcontentsline{toc}{subsection}{ऋक् २}

{\centering \footnotesize \textcolor{AccentPurple}{॥ गौराङ्गिरसः सूर्यः अतिजगती (अष्टिर्वा) ॥} \par}

\vspace{0.12em}

{\centering \textbf{\textcolor{AccentBlue}{आ\raisebox{0.25ex}{\accentmark{15}{\char"1CD2}}यं\raisebox{0.25ex}{\accentmark{15}{\char"1CD2}}सा\raisebox{0.25ex}{\accentmark{12}{\char"0951}}हस्रामान\raisebox{0.25ex}{\accentmark{12}{\char"0951}}वः\kern 0.05em ।\hskip 0.12em plus 0.04em minus 0.02em दृशा\raisebox{0.15ex}{\accentmark{12}{\char"1CF9}}ःकावीना\raisebox{0.25ex}{\accentmark{15}{\char"1CD2}}म्मातिज्योतिर्वीध\raisebox{0.25ex}{\accentmark{12}{\char"0951}}र्म\kern 0.05em ।\hskip 0.12em plus 0.04em minus 0.02em बुध्नस्सामीचीरूषा\raisebox{0.25ex}{\accentmark{15}{\char"1CD2}}सस्सा\raisebox{0.25ex}{\accentmark{15}{\char"1CD2}}मी\raisebox{0.25ex}{\accentmark{12}{\char"0951}}रयात्\kern 0.05em ।\hskip 0.12em plus 0.04em minus 0.02em आरेपा\raisebox{0.25ex}{\accentmark{15}{\char"1CD2}}सस्साचे\raisebox{0.25ex}{\accentmark{12}{\char"0951}}तसः\kern 0.05em ।\hskip 0.12em plus 0.04em minus 0.02em स्व\raisebox{0.25ex}{\accentmark{15}{\char"1CD2}}सा\raisebox{0.25ex}{\accentmark{12}{\char"0951}}रेम\raisebox{0.15ex}{\accentmark{12}{\char"1CF9}}न्यूमन्ता\raisebox{0.25ex}{\accentmark{12}{\char"0951}}श्चि\raisebox{0.25ex}{\accentmark{15}{\char"1CD2}}तागो\raisebox{0.25ex}{\accentmark{12}{\char"0951}}ः \nolinebreak\hskip 0.08em \mbox{॥२॥}}} \par}

\vspace{0.8em}
                % --- SUBSECTIONS (Mantras) ---
                
                % Conditional rendering based on output_mode
                \par\filbreak

\phantomsection

\addcontentsline{toc}{subsection}{ऋक् ३}

{\centering \footnotesize \textcolor{AccentPurple}{॥ परुच्छेपो दैवोदासिः इन्द्रः अत्यष्टिः ॥} \par}

\vspace{0.12em}

{\centering \textbf{\textcolor{AccentBlue}{ए\raisebox{0.15ex}{\accentmark{12}{\char"1CF8}}न्द्रा\raisebox{0.25ex}{\accentmark{12}{\char"0951}}याह्यू\raisebox{0.25ex}{\accentmark{15}{\char"1CD2}}पा\raisebox{0.25ex}{\accentmark{12}{\char"0951}}नःपारावा\raisebox{0.25ex}{\accentmark{15}{\char"1CD2}}ता\raisebox{0.25ex}{\accentmark{12}{\char"0951}}ः\kern 0.05em ।\hskip 0.12em plus 0.04em minus 0.02em ना\raisebox{0.15ex}{\accentmark{12}{\char"1CF8}}यमा\raisebox{0.15ex}{\accentmark{12}{\char"1CF8}}च्छा\raisebox{0.25ex}{\accentmark{12}{\char"0951}}वीद\raisebox{0.25ex}{\accentmark{15}{\char"1CD2}}धा\raisebox{0.25ex}{\accentmark{12}{\char"0951}}नीवासात्पा\raisebox{0.25ex}{\accentmark{12}{\char"0951}}तिः\kern 0.05em ।\hskip 0.12em plus 0.04em minus 0.02em अस्तारा\raisebox{0.15ex}{\accentmark{12}{\char"1CF9}}जे\raisebox{0.25ex}{\accentmark{12}{\char"0951}}वासात्पा\raisebox{0.25ex}{\accentmark{12}{\char"0951}}तिः\kern 0.05em ।\hskip 0.12em plus 0.04em minus 0.02em हावा\raisebox{0.25ex}{\accentmark{12}{\char"0951}}महेत्वाप्राया\raisebox{0.25ex}{\accentmark{12}{\char"0951}}स्वन्तस्सू\raisebox{0.25ex}{\accentmark{15}{\char"1CD2}}तेष्वा\raisebox{0.25ex}{\accentmark{12}{\char"0951}}\kern 0.05em ।\hskip 0.12em plus 0.04em minus 0.02em पु\raisebox{0.25ex}{\accentmark{15}{\char"1CD2}}त्रासो\raisebox{0.25ex}{\accentmark{12}{\char"0951}}ना\raisebox{0.25ex}{\accentmark{15}{\char"1CD2}}पीता\raisebox{0.25ex}{\accentmark{15}{\char"1CD2}}रंवा\raisebox{0.25ex}{\accentmark{15}{\char"1CD2}}ज\raisebox{0.25ex}{\accentmark{12}{\char"0951}}सातये\kern 0.05em ।\hskip 0.12em plus 0.04em minus 0.02em मंहीष्ठंवाजसातये \nolinebreak\hskip 0.08em \mbox{॥३॥}}} \par}

\vspace{0.8em}
                % --- SUBSECTIONS (Mantras) ---
                
                % Conditional rendering based on output_mode
                \par\filbreak

\phantomsection

\addcontentsline{toc}{subsection}{ऋक् ४}

{\centering \footnotesize \textcolor{AccentPurple}{॥ रेभः काश्यपः इन्द्रः अतिजगती (अष्टिर्वा) ॥} \par}

\vspace{0.12em}

{\centering \textbf{\textcolor{AccentBlue}{त\raisebox{0.25ex}{\accentmark{15}{\char"1CD2}}मि\raisebox{0.25ex}{\accentmark{15}{\char"1CD2}}न्द्रं\raisebox{0.25ex}{\accentmark{12}{\char"0951}}ज्योहवीमिमाघा\raisebox{0.25ex}{\accentmark{15}{\char"1CD2}}वा\raisebox{0.25ex}{\accentmark{12}{\char"0951}}नामुग्रम्\kern 0.05em ।\hskip 0.12em plus 0.04em minus 0.02em सत्रा\raisebox{0.15ex}{\accentmark{12}{\char"1CF9}}दधा\raisebox{0.25ex}{\accentmark{12}{\char"0951}}नामा\raisebox{0.25ex}{\accentmark{15}{\char"1CD2}}प्रातिष्कृतम्\kern 0.05em ।\hskip 0.12em plus 0.04em minus 0.02em श्रा\raisebox{0.15ex}{\accentmark{12}{\char"1CF9}}वां\raisebox{0.25ex}{\accentmark{12}{\char"0951}}सिभू\raisebox{0.25ex}{\accentmark{15}{\char"1CD2}}रिम्\kern 0.05em ।\hskip 0.12em plus 0.04em minus 0.02em मंही\raisebox{0.25ex}{\accentmark{12}{\char"0951}}ष्ठोगीर्भि\raisebox{0.25ex}{\accentmark{15}{\char"1CD2}}राचायाज्ञीयो\raisebox{0.25ex}{\accentmark{12}{\char"0951}}वा\raisebox{0.25ex}{\accentmark{15}{\char"1CD2}}वा\raisebox{0.25ex}{\accentmark{15}{\char"1CD2}}र्क्ता\raisebox{0.25ex}{\accentmark{12}{\char"0951}}\kern 0.05em ।\hskip 0.12em plus 0.04em minus 0.02em रा\raisebox{0.25ex}{\accentmark{15}{\char"1CD2}}येनो\raisebox{0.15ex}{\accentmark{12}{\char"1CF9}}विश्वासूपाथाकृ\raisebox{0.25ex}{\accentmark{12}{\char"0951}}णो\raisebox{0.25ex}{\accentmark{15}{\char"1CD2}}तूवज्री\raisebox{0.25ex}{\accentmark{15}{\char"1CD2}} \nolinebreak\hskip 0.08em \mbox{॥४॥}}} \par}

\vspace{0.8em}
                % --- SUBSECTIONS (Mantras) ---
                
                % Conditional rendering based on output_mode
                \par\filbreak

\phantomsection

\addcontentsline{toc}{subsection}{ऋक् ५}

{\centering \footnotesize \textcolor{AccentPurple}{॥ नकुलः सविता अत्यष्टिः ॥} \par}

\vspace{0.12em}

{\centering \textbf{\textcolor{AccentBlue}{आ\raisebox{0.25ex}{\accentmark{15}{\char"1CD2}}भित्य\raisebox{0.25ex}{\accentmark{15}{\char"1CD2}}न्देवंसा\raisebox{0.25ex}{\accentmark{12}{\char"0951}}वीता\raisebox{0.15ex}{\accentmark{12}{\char"1CF9}}रा\raisebox{0.15ex}{\accentmark{12}{\char"1CF9}}मोण्यो\raisebox{0.25ex}{\accentmark{15}{\char"1CD2}}ःका\raisebox{0.25ex}{\accentmark{15}{\char"1CD2}}विक्रा\raisebox{0.25ex}{\accentmark{12}{\char"0951}}तुम्\kern 0.05em ।\hskip 0.12em plus 0.04em minus 0.02em अर्चा\raisebox{0.25ex}{\accentmark{12}{\char"0951}}मीसत्या\raisebox{0.25ex}{\accentmark{15}{\char"1CD2}}सावंरात्नाधा\raisebox{0.25ex}{\accentmark{15}{\char"1CD2}}मा\raisebox{0.25ex}{\accentmark{15}{\char"1CD2}}भीप्रीयम्मातिम्\kern 0.05em ।\hskip 0.12em plus 0.04em minus 0.02em ऊर्ध्वा\raisebox{0.25ex}{\accentmark{15}{\char"1CD2}}य\raisebox{0.25ex}{\accentmark{15}{\char"1CD2}}स्याआमा\raisebox{0.25ex}{\accentmark{12}{\char"0951}}ती\raisebox{0.25ex}{\accentmark{15}{\char"1CD2}}र्भाआतिद्युत\raisebox{0.25ex}{\accentmark{15}{\char"1CD2}}त्सावि\raisebox{0.25ex}{\accentmark{12}{\char"0951}}मनी\kern 0.05em ।\hskip 0.12em plus 0.04em minus 0.02em हिर\raisebox{0.25ex}{\accentmark{12}{\char"0951}}ण्यपाणीरमीमीता\raisebox{0.25ex}{\accentmark{15}{\char"1CD2}}सू\raisebox{0.15ex}{\accentmark{12}{\char"1CF8}}क्रा\raisebox{0.15ex}{\accentmark{12}{\char"1CF9}}तू\raisebox{0.25ex}{\accentmark{12}{\char"0951}}कृपा\raisebox{0.25ex}{\accentmark{15}{\char"1CD2}}स्वा\raisebox{0.25ex}{\accentmark{12}{\char"0951}}ः \nolinebreak\hskip 0.08em \mbox{॥५॥}}} \par}

\vspace{0.8em}
                % --- SUBSECTIONS (Mantras) ---
                
                % Conditional rendering based on output_mode
                \par\filbreak

\phantomsection

\addcontentsline{toc}{subsection}{ऋक् ६}

{\centering \footnotesize \textcolor{AccentPurple}{॥ गृत्समद शौनकः इन्द्रः अष्टिः (अतिशक्करी वा) ॥} \par}

\vspace{0.12em}

{\centering \textbf{\textcolor{AccentBlue}{ता\raisebox{0.25ex}{\accentmark{15}{\char"1CD2}}वत्य\raisebox{0.15ex}{\accentmark{12}{\char"1CF9}}न्न\raisebox{0.15ex}{\accentmark{12}{\char"1CF9}}री\raisebox{0.25ex}{\accentmark{12}{\char"0951}}यनृतो\raisebox{0.25ex}{\accentmark{15}{\char"1CD2}}पा\raisebox{0.25ex}{\accentmark{12}{\char"0951}}इन्द्रा\kern 0.05em ।\hskip 0.12em plus 0.04em minus 0.02em प्राथामंपूर्व्य\raisebox{0.25ex}{\accentmark{15}{\char"1CD2}}न्ती\raisebox{0.25ex}{\accentmark{12}{\char"0951}}विप्रावा\raisebox{0.15ex}{\accentmark{12}{\char"1CF9}}च्यंकृत\raisebox{0.25ex}{\accentmark{15}{\char"1CD2}}म्\kern 0.05em ।\hskip 0.12em plus 0.04em minus 0.02em यो\raisebox{0.25ex}{\accentmark{15}{\char"1CD2}}देव\raisebox{0.25ex}{\accentmark{15}{\char"1CD2}}स्याशा\raisebox{0.25ex}{\accentmark{15}{\char"1CD2}}वासप्रारी\raisebox{0.25ex}{\accentmark{12}{\char"0951}}णाया\raisebox{0.15ex}{\accentmark{12}{\char"1CF9}}सुं\raisebox{0.25ex}{\accentmark{12}{\char"0951}}रीणन्नापा\raisebox{0.25ex}{\accentmark{15}{\char"1CD2}}ः\kern 0.05em ।\hskip 0.12em plus 0.04em minus 0.02em भु\raisebox{0.25ex}{\accentmark{15}{\char"1CD2}}वोविश्वा\raisebox{0.25ex}{\accentmark{12}{\char"0951}}माभ्या\raisebox{0.15ex}{\accentmark{12}{\char"1CF9}}दे\raisebox{0.25ex}{\accentmark{12}{\char"0951}}वामो\raisebox{0.25ex}{\accentmark{15}{\char"1CD2}}ज\raisebox{0.25ex}{\accentmark{12}{\char"0951}}सा\kern 0.05em ।\hskip 0.12em plus 0.04em minus 0.02em वीर्द\raisebox{0.15ex}{\accentmark{12}{\char"1CF9}}दू\raisebox{0.15ex}{\accentmark{12}{\char"1CF9}}र्जंशा\raisebox{0.25ex}{\accentmark{15}{\char"1CD2}}ता\raisebox{0.25ex}{\accentmark{12}{\char"0951}}क्रातूर्वीदे\raisebox{0.25ex}{\accentmark{15}{\char"1CD2}}दीषम् \nolinebreak\hskip 0.08em \mbox{॥६॥}}} \par}

\vspace{0.8em}
                % --- SUBSECTIONS (Mantras) ---
                
                % Conditional rendering based on output_mode
                \par\filbreak

\phantomsection

\addcontentsline{toc}{subsection}{ऋक् ७}

{\centering \footnotesize \textcolor{AccentPurple}{॥ परुच्छेपो दैवोदासिः विश्वेदेवाः अत्यष्टिः ॥} \par}

\vspace{0.12em}

{\centering \textbf{\textcolor{AccentBlue}{अस्तुश्रौ\raisebox{0.25ex}{\accentmark{15}{\char"1CD2}}षा\raisebox{0.25ex}{\accentmark{12}{\char"0951}}ट्पूरो\raisebox{0.25ex}{\accentmark{15}{\char"1CD2}}आ\raisebox{0.25ex}{\accentmark{12}{\char"0951}}ग्नीन्धीया\raisebox{0.25ex}{\accentmark{15}{\char"1CD2}}दधे\kern 0.05em ।\hskip 0.12em plus 0.04em minus 0.02em आ\raisebox{0.15ex}{\accentmark{12}{\char"1CF8}}नूत्य\raisebox{0.15ex}{\accentmark{12}{\char"1CF8}}च्छर्धोदिव्यं\raisebox{0.25ex}{\accentmark{15}{\char"1CD2}}वृणीमहे\kern 0.05em ।\hskip 0.12em plus 0.04em minus 0.02em इन्द्रं\raisebox{0.25ex}{\accentmark{12}{\char"0951}}वायू\raisebox{0.25ex}{\accentmark{15}{\char"1CD2}}वृ\raisebox{0.25ex}{\accentmark{12}{\char"0951}}णीमहे\kern 0.05em ।\hskip 0.12em plus 0.04em minus 0.02em य\raisebox{0.15ex}{\accentmark{12}{\char"1CF8}}द्धा\raisebox{0.25ex}{\accentmark{12}{\char"0951}}क्राणा\raisebox{0.25ex}{\accentmark{15}{\char"1CD2}}वी\raisebox{0.25ex}{\accentmark{15}{\char"1CD2}}वा\raisebox{0.25ex}{\accentmark{15}{\char"1CD2}}स्वा\raisebox{0.25ex}{\accentmark{12}{\char"0951}}ते\kern 0.05em ।\hskip 0.12em plus 0.04em minus 0.02em ना\raisebox{0.15ex}{\accentmark{12}{\char"1CF8}}भा\raisebox{0.25ex}{\accentmark{12}{\char"0951}}सन्दायाना\raisebox{0.25ex}{\accentmark{15}{\char"1CD2}}व्या\raisebox{0.15ex}{\accentmark{12}{\char"1CF8}}सेआ\raisebox{0.25ex}{\accentmark{15}{\char"1CD2}}धःप्रा\raisebox{0.25ex}{\accentmark{15}{\char"1CD2}}नूना\raisebox{0.25ex}{\accentmark{15}{\char"1CD2}}मूपा\raisebox{0.25ex}{\accentmark{12}{\char"0951}}यन्तीधीता\raisebox{0.25ex}{\accentmark{15}{\char"1CD2}}या\raisebox{0.25ex}{\accentmark{12}{\char"0951}}ः\kern 0.05em ।\hskip 0.12em plus 0.04em minus 0.02em दे\raisebox{0.25ex}{\accentmark{15}{\char"1CD2}}वं\raisebox{0.25ex}{\accentmark{15}{\char"1CD2}}अ\raisebox{0.25ex}{\accentmark{15}{\char"1CD2}}च्छाना\raisebox{0.25ex}{\accentmark{15}{\char"1CD2}}धीता\raisebox{0.25ex}{\accentmark{15}{\char"1CD2}}याः\kern 0.05em ।\hskip 0.12em plus 0.04em minus 0.02em \nolinebreak\hskip 0.08em \mbox{॥७॥}}} \par}

\vspace{0.8em}
                % --- SUBSECTIONS (Mantras) ---
                
                % Conditional rendering based on output_mode
                \par\filbreak

\phantomsection

\addcontentsline{toc}{subsection}{ऋक् ८}

{\centering \footnotesize \textcolor{AccentPurple}{॥ एवयामरुदात्रेयः मरुतः अतिजगती (अष्टिर्वा) ॥} \par}

\vspace{0.12em}

{\centering \textbf{\textcolor{AccentBlue}{प्र\raisebox{0.15ex}{\accentmark{12}{\char"1CF8}}वोमाहे\raisebox{0.25ex}{\accentmark{15}{\char"1CD2}}माता\raisebox{0.25ex}{\accentmark{15}{\char"1CD2}}यो\raisebox{0.25ex}{\accentmark{12}{\char"0951}}यन्तूवी\raisebox{0.25ex}{\accentmark{15}{\char"1CD2}}ष्णा\raisebox{0.25ex}{\accentmark{12}{\char"0951}}वे\kern 0.05em ।\hskip 0.12em plus 0.04em minus 0.02em मारू\raisebox{0.25ex}{\accentmark{15}{\char"1CD2}}त्वा\raisebox{0.25ex}{\accentmark{12}{\char"0951}}तेगीरीजा\raisebox{0.15ex}{\accentmark{12}{\char"1CF9}}ये\raisebox{0.25ex}{\accentmark{12}{\char"0951}}वायामा\raisebox{0.25ex}{\accentmark{12}{\char"0951}}रूत्\kern 0.05em ।\hskip 0.12em plus 0.04em minus 0.02em प्र\raisebox{0.15ex}{\accentmark{12}{\char"1CF8}}शर्धायाप्रयज्यावेसूखाद\raisebox{0.25ex}{\accentmark{15}{\char"1CD2}}ये\raisebox{0.25ex}{\accentmark{12}{\char"0951}}\kern 0.05em ।\hskip 0.12em plus 0.04em minus 0.02em ता\raisebox{0.25ex}{\accentmark{15}{\char"1CD2}}वा\raisebox{0.15ex}{\accentmark{12}{\char"1CF8}}सेभन्द\raisebox{0.25ex}{\accentmark{15}{\char"1CD2}}दिष्टये\kern 0.05em ।\hskip 0.12em plus 0.04em minus 0.02em धू\raisebox{0.25ex}{\accentmark{15}{\char"1CD2}}नीव्रतायाशा\raisebox{0.25ex}{\accentmark{15}{\char"1CD2}}वासे \nolinebreak\hskip 0.08em \mbox{॥८॥}}} \par}

\vspace{0.8em}
                % --- SUBSECTIONS (Mantras) ---
                
                % Conditional rendering based on output_mode
                \par\filbreak

\phantomsection

\addcontentsline{toc}{subsection}{ऋक् ९}

{\centering \footnotesize \textcolor{AccentPurple}{॥ अनानतः पारुच्छेपिः पवमानः सोमः अत्यष्टिः ॥} \par}

\vspace{0.12em}

{\centering \textbf{\textcolor{AccentBlue}{आ\raisebox{0.25ex}{\accentmark{15}{\char"1CD2}}या\raisebox{0.25ex}{\accentmark{15}{\char"1CD2}}रूचा\raisebox{0.25ex}{\accentmark{15}{\char"1CD2}}हा\raisebox{0.25ex}{\accentmark{15}{\char"1CD2}}रिण्यापूना\raisebox{0.25ex}{\accentmark{15}{\char"1CD2}}न\raisebox{0.25ex}{\accentmark{15}{\char"1CD2}}ः\kern 0.05em ।\hskip 0.12em plus 0.04em minus 0.02em वि\raisebox{0.25ex}{\accentmark{15}{\char"1CD2}}श्वा\raisebox{0.25ex}{\accentmark{15}{\char"1CD2}}द्वेषांसितरतीसायूग्वा\raisebox{0.25ex}{\accentmark{12}{\char"0951}}भिः\kern 0.05em ।\hskip 0.12em plus 0.04em minus 0.02em सूरोना\raisebox{0.25ex}{\accentmark{15}{\char"1CD2}}सायूग्वा\raisebox{0.25ex}{\accentmark{12}{\char"0951}}भिः\kern 0.05em ।\hskip 0.12em plus 0.04em minus 0.02em धा\raisebox{0.15ex}{\accentmark{12}{\char"1CF8}}रा\raisebox{0.25ex}{\accentmark{12}{\char"0951}}पृष्ठास्या\raisebox{0.25ex}{\accentmark{12}{\char"0951}}रोचते\kern 0.05em ।\hskip 0.12em plus 0.04em minus 0.02em पु\raisebox{0.25ex}{\accentmark{15}{\char"1CD2}}नानो\raisebox{0.15ex}{\accentmark{12}{\char"1CF9}}आ\raisebox{0.25ex}{\accentmark{12}{\char"0951}}रुषोहा\raisebox{0.25ex}{\accentmark{15}{\char"1CD2}}रिःवि\raisebox{0.25ex}{\accentmark{15}{\char"1CD2}}श्वा\raisebox{0.25ex}{\accentmark{12}{\char"0951}}यद्रूपा\raisebox{0.15ex}{\accentmark{12}{\char"1CF9}}पा\raisebox{0.25ex}{\accentmark{12}{\char"0951}}र्या\raisebox{0.25ex}{\accentmark{15}{\char"1CD2}}स्यऋत्वा\raisebox{0.25ex}{\accentmark{12}{\char"0951}}भिः\kern 0.05em ।\hskip 0.12em plus 0.04em minus 0.02em सप्ता\raisebox{0.15ex}{\accentmark{12}{\char"1CF9}}स्ये\raisebox{0.25ex}{\accentmark{12}{\char"0951}}भिऋ\raisebox{0.25ex}{\accentmark{15}{\char"1CD2}}त्वा\raisebox{0.25ex}{\accentmark{12}{\char"0951}}भिः \nolinebreak\hskip 0.08em \mbox{॥९॥}}} \par}

\vspace{0.8em}
                % --- SUBSECTIONS (Mantras) ---
                
                % Conditional rendering based on output_mode
                \par\filbreak

\phantomsection

\addcontentsline{toc}{subsection}{ऋक् १०}

{\centering \footnotesize \textcolor{AccentPurple}{॥ परुच्छेपो दैवोदासिः अग्निः अत्यष्टिः ॥} \par}

\vspace{0.12em}

{\centering \textbf{\textcolor{AccentBlue}{अग्निंहो\raisebox{0.25ex}{\accentmark{15}{\char"1CD2}}ता\raisebox{0.25ex}{\accentmark{12}{\char"0951}}रम्मन्येदा\raisebox{0.25ex}{\accentmark{15}{\char"1CD2}}स्व\raisebox{0.25ex}{\accentmark{12}{\char"0951}}न्तम्\kern 0.05em ।\hskip 0.12em plus 0.04em minus 0.02em वा\raisebox{0.15ex}{\accentmark{12}{\char"1CF8}}सो\raisebox{0.25ex}{\accentmark{12}{\char"0951}}स्सूनु\raisebox{0.25ex}{\accentmark{12}{\char"0951}}सा\raisebox{0.25ex}{\accentmark{15}{\char"1CD2}}हा\raisebox{0.25ex}{\accentmark{12}{\char"0951}}साजाता\raisebox{0.25ex}{\accentmark{15}{\char"1CD2}}वेदसम्\kern 0.05em ।\hskip 0.12em plus 0.04em minus 0.02em विप्रान्ना\raisebox{0.25ex}{\accentmark{15}{\char"1CD2}}जाता\raisebox{0.25ex}{\accentmark{15}{\char"1CD2}}वे\raisebox{0.25ex}{\accentmark{12}{\char"0951}}दसम्\kern 0.05em ।\hskip 0.12em plus 0.04em minus 0.02em या\raisebox{0.25ex}{\accentmark{12}{\char"0951}}ऊर्ध्वायायास्वाध्वार\raisebox{0.25ex}{\accentmark{15}{\char"1CD2}}ः\kern 0.05em ।\hskip 0.12em plus 0.04em minus 0.02em देवोदेवा\raisebox{0.15ex}{\accentmark{12}{\char"1CF9}}च्य\raisebox{0.25ex}{\accentmark{12}{\char"0951}}कृ\raisebox{0.25ex}{\accentmark{12}{\char"0951}}पा\raisebox{0.25ex}{\accentmark{12}{\char"0951}}\kern 0.05em ।\hskip 0.12em plus 0.04em minus 0.02em घृ\raisebox{0.25ex}{\accentmark{15}{\char"1CD2}}त\raisebox{0.25ex}{\accentmark{15}{\char"1CD2}}स्या\raisebox{0.25ex}{\accentmark{12}{\char"0951}}बिभ्र\raisebox{0.25ex}{\accentmark{12}{\char"0951}}ष्टीमा\raisebox{0.15ex}{\accentmark{12}{\char"1CF9}}नू\raisebox{0.25ex}{\accentmark{12}{\char"0951}}शुक्रा\raisebox{0.25ex}{\accentmark{15}{\char"1CD2}}शो\raisebox{0.25ex}{\accentmark{12}{\char"0951}}चिषः\kern 0.05em ।\hskip 0.12em plus 0.04em minus 0.02em आ\raisebox{0.25ex}{\accentmark{15}{\char"1CD2}}जुं\raisebox{0.25ex}{\accentmark{15}{\char"1CD2}}भ्ह्वा\raisebox{0.25ex}{\accentmark{12}{\char"0951}}नास्यासर्पीषाः \nolinebreak\hskip 0.08em \mbox{॥१०॥}}} \par}

\vspace{0.8em}

            % --- End of Kandah Closing Floral Line & Counts ---
            \par\nopagebreak\vspace{0.9em plus 0.3em minus 0.1em}
            \centerline{\textcolor{black!35}{\rule{1.8cm}{0.35pt}}\;\raisebox{-0.22ex}{\textcolor{vedicmaroon}{\large\SixFlowerPetalDotted}}\;\textcolor{black!35}{\rule{1.8cm}{0.35pt}}}
            \nopagebreak\vspace{0.45em}
            \centerline{\devafont\textbf{\textcolor{AccentBlue}{॥ ऋक् सङ्ख्या : १० ॥}}}


    % --- SUPERSECTION ---
    \clearpage           
    \phantomsection              
    \chapter*{\centering पवमानपाठः}
    \addcontentsline{toc}{chapter}{पवमानपाठः}
    \markboth{पवमानपाठः}{पवमानपाठः}
    
    % GLUE FIX: Prevent page break between Chapter and First Section
    \nopagebreak

        

            % --- SECTION ---
            
            \phantomsection  

            \section*{\centering प्रथम खण्डः}
            \addcontentsline{toc}{section}{प्रथम खण्डः} 
            \markright{प्रथम खण्डः}
            \nopagebreak % Keep section title with first mantra
            
            % --- SECTION-LEVEL PROCEDURE LINK ---

                % --- SUBSECTIONS (Mantras) ---
                
                % Conditional rendering based on output_mode
                \par\filbreak

\phantomsection

\addcontentsline{toc}{subsection}{ऋक् १}

{\centering \footnotesize \textcolor{AccentPurple}{॥ अमहीयुराङ्गिरसः पवमानः सोमः गायत्री ॥} \par}

\vspace{0.12em}

{\centering \textbf{\textcolor{AccentBlue}{उच्चा\raisebox{0.25ex}{\accentmark{12}{\char"0951}}ते\raisebox{0.25ex}{\accentmark{12}{\char"0951}}जातमन्धा\raisebox{0.15ex}{\accentmark{12}{\char"1CF8}}सा\raisebox{0.25ex}{\accentmark{15}{\char"1CD2}}\kern 0.05em ।\hskip 0.12em plus 0.04em minus 0.02em दिविसा\raisebox{0.25ex}{\accentmark{15}{\char"1CD2}}त्भूम्याद\raisebox{0.25ex}{\accentmark{12}{\char"0951}}दे\kern 0.05em ।\hskip 0.12em plus 0.04em minus 0.02em उग्रं\raisebox{0.25ex}{\accentmark{15}{\char"1CD2}}शर्म्मा\raisebox{0.25ex}{\accentmark{15}{\char"1CD2}}माहिश्रा\raisebox{0.25ex}{\accentmark{15}{\char"1CD2}}वाः \nolinebreak\hskip 0.08em \mbox{॥१॥}}} \par}

\vspace{0.8em}
                % --- SUBSECTIONS (Mantras) ---
                
                % Conditional rendering based on output_mode
                \par\filbreak

\phantomsection

\addcontentsline{toc}{subsection}{ऋक् २}

{\centering \footnotesize \textcolor{AccentPurple}{॥ मधुच्छन्दा वैश्वामित्रः पवमानः सोमः गायत्री ॥} \par}

\vspace{0.12em}

{\centering \textbf{\textcolor{AccentBlue}{स्वा\raisebox{0.25ex}{\accentmark{12}{\char"0951}}दी\raisebox{0.25ex}{\accentmark{12}{\char"0951}}ष्ठायामा\raisebox{0.25ex}{\accentmark{15}{\char"1CD2}}दी\raisebox{0.25ex}{\accentmark{12}{\char"0951}}ष्ठाया\kern 0.05em ।\hskip 0.12em plus 0.04em minus 0.02em पा\raisebox{0.25ex}{\accentmark{15}{\char"1CD2}}वा\raisebox{0.25ex}{\accentmark{12}{\char"0951}}स्वसोमाधारा\raisebox{0.25ex}{\accentmark{12}{\char"0951}}या\kern 0.05em ।\hskip 0.12em plus 0.04em minus 0.02em इ\raisebox{0.15ex}{\accentmark{12}{\char"1CF8}}न्द्रा\raisebox{0.25ex}{\accentmark{12}{\char"0951}}यापा\raisebox{0.25ex}{\accentmark{15}{\char"1CD2}}तावेसूत\raisebox{0.25ex}{\accentmark{15}{\char"1CD2}}ः \nolinebreak\hskip 0.08em \mbox{॥२॥}}} \par}

\vspace{0.8em}
                % --- SUBSECTIONS (Mantras) ---
                
                % Conditional rendering based on output_mode
                \par\filbreak

\phantomsection

\addcontentsline{toc}{subsection}{ऋक् ३}

{\centering \footnotesize \textcolor{AccentPurple}{॥ भृगुर्वारुणिर्जमदग्निर्भार्गवो वा पवमानः सोमः गायत्री ॥} \par}

\vspace{0.12em}

{\centering \textbf{\textcolor{AccentBlue}{वृ\raisebox{0.25ex}{\accentmark{15}{\char"1CD2}}षा\raisebox{0.25ex}{\accentmark{12}{\char"0951}}पवास्वधा\raisebox{0.25ex}{\accentmark{15}{\char"1CD2}}रा\raisebox{0.25ex}{\accentmark{12}{\char"0951}}या\kern 0.05em ।\hskip 0.12em plus 0.04em minus 0.02em मा\raisebox{0.25ex}{\accentmark{15}{\char"1CD2}}रूत्वा\raisebox{0.25ex}{\accentmark{12}{\char"0951}}तेचमात्सार\raisebox{0.25ex}{\accentmark{15}{\char"1CD2}}ः\kern 0.05em ।\hskip 0.12em plus 0.04em minus 0.02em विश्वा\raisebox{0.25ex}{\accentmark{12}{\char"0951}}दधा\raisebox{0.25ex}{\accentmark{12}{\char"0951}}नओज\raisebox{0.25ex}{\accentmark{15}{\char"1CD2}}सा \nolinebreak\hskip 0.08em \mbox{॥३॥}}} \par}

\vspace{0.8em}
                % --- SUBSECTIONS (Mantras) ---
                
                % Conditional rendering based on output_mode
                \par\filbreak

\phantomsection

\addcontentsline{toc}{subsection}{ऋक् ४}

{\centering \footnotesize \textcolor{AccentPurple}{॥ अमहीयुराङ्गिरसः पवमानः सोमः गायत्री ॥} \par}

\vspace{0.12em}

{\centering \textbf{\textcolor{AccentBlue}{यस्तेमा\raisebox{0.25ex}{\accentmark{15}{\char"1CD2}}दोवा\raisebox{0.25ex}{\accentmark{15}{\char"1CD2}}रे\raisebox{0.25ex}{\accentmark{12}{\char"0951}}ण्यः\kern 0.05em ।\hskip 0.12em plus 0.04em minus 0.02em तेना\raisebox{0.25ex}{\accentmark{12}{\char"0951}}पावस्वा\raisebox{0.25ex}{\accentmark{15}{\char"1CD2}}न्ध\raisebox{0.25ex}{\accentmark{12}{\char"0951}}सा\kern 0.05em ।\hskip 0.12em plus 0.04em minus 0.02em देवावीरा\raisebox{0.25ex}{\accentmark{12}{\char"0951}}घशंसाहा\raisebox{0.25ex}{\accentmark{15}{\char"1CD2}} \nolinebreak\hskip 0.08em \mbox{॥४॥}}} \par}

\vspace{0.8em}
                % --- SUBSECTIONS (Mantras) ---
                
                % Conditional rendering based on output_mode
                \par\filbreak

\phantomsection

\addcontentsline{toc}{subsection}{ऋक् ५}

{\centering \footnotesize \textcolor{AccentPurple}{॥ त्रित आप्त्यः पवमानः सोमः गायत्री ॥} \par}

\vspace{0.12em}

{\centering \textbf{\textcolor{AccentBlue}{ति\raisebox{0.15ex}{\accentmark{12}{\char"1CF8}}स्रो\raisebox{0.25ex}{\accentmark{15}{\char"1CD2}}वाचाऊदीरते\kern 0.05em ।\hskip 0.12em plus 0.04em minus 0.02em गा\raisebox{0.25ex}{\accentmark{15}{\char"1CD2}}वो\raisebox{0.25ex}{\accentmark{12}{\char"0951}}मिमन्तीधेना\raisebox{0.25ex}{\accentmark{15}{\char"1CD2}}वा\raisebox{0.25ex}{\accentmark{12}{\char"0951}}ः\kern 0.05em ।\hskip 0.12em plus 0.04em minus 0.02em हा\raisebox{0.25ex}{\accentmark{15}{\char"1CD2}}री\raisebox{0.25ex}{\accentmark{12}{\char"0951}}रेतीका\raisebox{0.25ex}{\accentmark{15}{\char"1CD2}}नीक्रदत् \nolinebreak\hskip 0.08em \mbox{॥५॥}}} \par}

\vspace{0.8em}
                % --- SUBSECTIONS (Mantras) ---
                
                % Conditional rendering based on output_mode
                \par\filbreak

\phantomsection

\addcontentsline{toc}{subsection}{ऋक् ६}

{\centering \footnotesize \textcolor{AccentPurple}{॥ कश्यपोः मारीचः पवमानः सोमः गायत्री ॥} \par}

\vspace{0.12em}

{\centering \textbf{\textcolor{AccentBlue}{इन्द्रा\raisebox{0.25ex}{\accentmark{12}{\char"0951}}येन्दोमा\raisebox{0.15ex}{\accentmark{12}{\char"1CF9}}रू\raisebox{0.25ex}{\accentmark{15}{\char"1CD2}}त्वा\raisebox{0.25ex}{\accentmark{12}{\char"0951}}ते\kern 0.05em ।\hskip 0.12em plus 0.04em minus 0.02em पा\raisebox{0.15ex}{\accentmark{12}{\char"1CF8}}वास्वामा\raisebox{0.25ex}{\accentmark{15}{\char"1CD2}}धूमक्तमः\kern 0.05em ।\hskip 0.12em plus 0.04em minus 0.02em अर्का\raisebox{0.25ex}{\accentmark{15}{\char"1CD2}}स्यायो\raisebox{0.25ex}{\accentmark{15}{\char"1CD2}}नीमा\raisebox{0.25ex}{\accentmark{15}{\char"1CD2}}साद\raisebox{0.25ex}{\accentmark{12}{\char"0951}}ः \nolinebreak\hskip 0.08em \mbox{॥६॥}}} \par}

\vspace{0.8em}
                % --- SUBSECTIONS (Mantras) ---
                
                % Conditional rendering based on output_mode
                \par\filbreak

\phantomsection

\addcontentsline{toc}{subsection}{ऋक् ७}

{\centering \footnotesize \textcolor{AccentPurple}{॥ जमदग्निर्भार्गवः पवमानः सोमः गायत्री ॥} \par}

\vspace{0.12em}

{\centering \textbf{\textcolor{AccentBlue}{अ\raisebox{0.15ex}{\accentmark{12}{\char"1CF8}}स\raisebox{0.25ex}{\accentmark{12}{\char"0951}}व्यंशुर्मा\raisebox{0.25ex}{\accentmark{15}{\char"1CD2}}दाय\kern 0.05em ।\hskip 0.12em plus 0.04em minus 0.02em अ\raisebox{0.15ex}{\accentmark{12}{\char"1CF8}}प्सु\raisebox{0.25ex}{\accentmark{15}{\char"1CD2}}द\raisebox{0.25ex}{\accentmark{15}{\char"1CD2}}क्षोगीरीष्ठाः\kern 0.05em ।\hskip 0.12em plus 0.04em minus 0.02em श्ये\raisebox{0.25ex}{\accentmark{15}{\char"1CD2}}नो\raisebox{0.25ex}{\accentmark{15}{\char"1CD2}}नयोनीमा\raisebox{0.25ex}{\accentmark{15}{\char"1CD2}}सा\raisebox{0.25ex}{\accentmark{12}{\char"0951}}दत् \nolinebreak\hskip 0.08em \mbox{॥७॥}}} \par}

\vspace{0.8em}
                % --- SUBSECTIONS (Mantras) ---
                
                % Conditional rendering based on output_mode
                \par\filbreak

\phantomsection

\addcontentsline{toc}{subsection}{ऋक् ८}

{\centering \footnotesize \textcolor{AccentPurple}{॥ दृढच्युत आगस्त्यः पवमानः सोमः गायत्री ॥} \par}

\vspace{0.12em}

{\centering \textbf{\textcolor{AccentBlue}{पा\raisebox{0.15ex}{\accentmark{12}{\char"1CF8}}वा\raisebox{0.25ex}{\accentmark{12}{\char"0951}}स्वदक्षासाधा\raisebox{0.25ex}{\accentmark{12}{\char"0951}}नः\kern 0.05em ।\hskip 0.12em plus 0.04em minus 0.02em देवे\raisebox{0.15ex}{\accentmark{12}{\char"1CF9}}भ्य\raisebox{0.25ex}{\accentmark{12}{\char"0951}}ःपीता\raisebox{0.25ex}{\accentmark{15}{\char"1CD2}}ये\raisebox{0.25ex}{\accentmark{12}{\char"0951}}हरे\kern 0.05em ।\hskip 0.12em plus 0.04em minus 0.02em मारू\raisebox{0.15ex}{\accentmark{12}{\char"1CF9}}त्भ्योवाया\raisebox{0.25ex}{\accentmark{15}{\char"1CD2}}वेमा\raisebox{0.25ex}{\accentmark{15}{\char"1CD2}}दा\raisebox{0.25ex}{\accentmark{12}{\char"0951}}ः \nolinebreak\hskip 0.08em \mbox{॥८॥}}} \par}

\vspace{0.8em}
                % --- SUBSECTIONS (Mantras) ---
                
                % Conditional rendering based on output_mode
                \par\filbreak

\phantomsection

\addcontentsline{toc}{subsection}{ऋक् ९}

{\centering \footnotesize \textcolor{AccentPurple}{॥ असितः काश्यपो देवलो वा पवमानः सोमः गायत्री ॥} \par}

\vspace{0.12em}

{\centering \textbf{\textcolor{AccentBlue}{पारीस्वानो\raisebox{0.15ex}{\accentmark{12}{\char"1CF9}}गीरीष्ठा\raisebox{0.25ex}{\accentmark{15}{\char"1CD2}}ः\kern 0.05em ।\hskip 0.12em plus 0.04em minus 0.02em पावी\raisebox{0.25ex}{\accentmark{15}{\char"1CD2}}त्रेसोमो\raisebox{0.25ex}{\accentmark{12}{\char"0951}}अक्षरात्\kern 0.05em ।\hskip 0.12em plus 0.04em minus 0.02em मा\raisebox{0.25ex}{\accentmark{15}{\char"1CD2}}दे\raisebox{0.25ex}{\accentmark{12}{\char"0951}}षुसार्वाधा\raisebox{0.25ex}{\accentmark{15}{\char"1CD2}}आसि \nolinebreak\hskip 0.08em \mbox{॥९॥}}} \par}

\vspace{0.8em}
                % --- SUBSECTIONS (Mantras) ---
                
                % Conditional rendering based on output_mode
                \par\filbreak

\phantomsection

\addcontentsline{toc}{subsection}{ऋक् १०}

{\centering \footnotesize \textcolor{AccentPurple}{॥ असितः काश्यपो देवलो वा पवमानः सोमः गायत्री ॥} \par}

\vspace{0.12em}

{\centering \textbf{\textcolor{AccentBlue}{पा\raisebox{0.15ex}{\accentmark{12}{\char"1CF8}}रि\raisebox{0.25ex}{\accentmark{12}{\char"0951}}प्रीया\raisebox{0.25ex}{\accentmark{15}{\char"1CD2}}दीवा\raisebox{0.25ex}{\accentmark{15}{\char"1CD2}}ःकावी\raisebox{0.25ex}{\accentmark{15}{\char"1CD2}}ः\kern 0.05em ।\hskip 0.12em plus 0.04em minus 0.02em वा\raisebox{0.25ex}{\accentmark{15}{\char"1CD2}}यां\raisebox{0.25ex}{\accentmark{12}{\char"0951}}सिनप्त्यो\raisebox{0.25ex}{\accentmark{15}{\char"1CD2}}र्हितः\kern 0.05em ।\hskip 0.12em plus 0.04em minus 0.02em स्वानै\raisebox{0.25ex}{\accentmark{15}{\char"1CD2}}र्यातीःकावी\raisebox{0.25ex}{\accentmark{15}{\char"1CD2}}क्रातुः \nolinebreak\hskip 0.08em \mbox{॥१०॥}}} \par}

\vspace{0.8em}

            % --- End of Kandah Closing Floral Line & Counts ---
            \par\nopagebreak\vspace{0.9em plus 0.3em minus 0.1em}
            \centerline{\textcolor{black!35}{\rule{1.8cm}{0.35pt}}\;\raisebox{-0.22ex}{\textcolor{vedicmaroon}{\large\SixFlowerPetalDotted}}\;\textcolor{black!35}{\rule{1.8cm}{0.35pt}}}
            \nopagebreak\vspace{0.45em}
            \centerline{\devafont\textbf{\textcolor{AccentBlue}{॥ ऋक् सङ्ख्या : १० ॥}}}
        
            \clearpage

            % --- SECTION ---
            
            \phantomsection  

            \section*{\centering द्वितीय खण्डः}
            \addcontentsline{toc}{section}{द्वितीय खण्डः} 
            \markright{द्वितीय खण्डः}
            \nopagebreak % Keep section title with first mantra
            
            % --- SECTION-LEVEL PROCEDURE LINK ---

                % --- SUBSECTIONS (Mantras) ---
                
                % Conditional rendering based on output_mode
                \par\filbreak

\phantomsection

\addcontentsline{toc}{subsection}{ऋक् १}

{\centering \textbf{\textcolor{AccentBlue}{॥ (कविर्मेधावी)श्यावाश्वःआत्रेयःपवमानःसोमःगायत्री ॥ ॥ त्रितआप्त्यःपवमानःसोमःगायत्री ॥ प्र\raisebox{0.15ex}{\accentmark{12}{\char"1CF8}}सो\raisebox{0.15ex}{\accentmark{12}{\char"1CF8}}मासोमादच्यू\raisebox{0.25ex}{\accentmark{15}{\char"1CD2}}ताः\kern 0.05em ।\hskip 0.12em plus 0.04em minus 0.02em श्रा\raisebox{0.25ex}{\accentmark{15}{\char"1CD2}}वा\raisebox{0.25ex}{\accentmark{12}{\char"0951}}सेनोमा\raisebox{0.25ex}{\accentmark{12}{\char"0951}}घो\raisebox{0.25ex}{\accentmark{15}{\char"1CD2}}नाम्\raisebox{0.25ex}{\accentmark{12}{\char"0951}}\kern 0.05em ।\hskip 0.12em plus 0.04em minus 0.02em सू\raisebox{0.25ex}{\accentmark{15}{\char"1CD2}}तावीद\raisebox{0.25ex}{\accentmark{15}{\char"1CD2}}थेअक्रमुः \nolinebreak\hskip 0.08em \mbox{॥१॥}}} \par}

\vspace{0.8em}
                % --- SUBSECTIONS (Mantras) ---
                
                % Conditional rendering based on output_mode
                \par\filbreak

\phantomsection

\addcontentsline{toc}{subsection}{ऋक् २}

{\centering \textbf{\textcolor{AccentBlue}{प्र\raisebox{0.25ex}{\accentmark{15}{\char"1CD2}}सो\raisebox{0.25ex}{\accentmark{15}{\char"1CD2}}मा\raisebox{0.25ex}{\accentmark{12}{\char"0951}}सोवीपश्ची\raisebox{0.25ex}{\accentmark{15}{\char"1CD2}}ता\raisebox{0.25ex}{\accentmark{12}{\char"0951}}ः\kern 0.05em ।\hskip 0.12em plus 0.04em minus 0.02em आ\raisebox{0.25ex}{\accentmark{15}{\char"1CD2}}पो\raisebox{0.25ex}{\accentmark{12}{\char"0951}}नायन्ताऊर्म\raisebox{0.25ex}{\accentmark{12}{\char"0951}}या\raisebox{0.25ex}{\accentmark{12}{\char"0951}}ः\kern 0.05em ।\hskip 0.12em plus 0.04em minus 0.02em वाना\raisebox{0.25ex}{\accentmark{15}{\char"1CD2}}नी\raisebox{0.25ex}{\accentmark{12}{\char"0951}}माहीषा\raisebox{0.25ex}{\accentmark{12}{\char"0951}}ईवा \nolinebreak\hskip 0.08em \mbox{॥२॥}}} \par}

\vspace{0.8em}
                % --- SUBSECTIONS (Mantras) ---
                
                % Conditional rendering based on output_mode
                \par\filbreak

\phantomsection

\addcontentsline{toc}{subsection}{ऋक् ३}

{\centering \textbf{\textcolor{AccentBlue}{॥ अमहीयुराङ्गिरसःपवमानःसोमःगायत्री ॥ ॥ भृगुर्वारुणिर्जमदग्निर्भार्गवोवापवमानःसोमःगायत्री ॥ पा\raisebox{0.25ex}{\accentmark{15}{\char"1CD2}}वा\raisebox{0.25ex}{\accentmark{12}{\char"0951}}स्वेन्दोवृ\raisebox{0.15ex}{\accentmark{12}{\char"1CF9}}षा\raisebox{0.25ex}{\accentmark{12}{\char"0951}}सूतः\kern 0.05em ।\hskip 0.12em plus 0.04em minus 0.02em कृ\raisebox{0.25ex}{\accentmark{15}{\char"1CD2}}धी\raisebox{0.15ex}{\accentmark{12}{\char"1CF9}}नोयाशासोजने\kern 0.05em ।\hskip 0.12em plus 0.04em minus 0.02em विश्वा\raisebox{0.25ex}{\accentmark{15}{\char"1CD2}}आप\raisebox{0.25ex}{\accentmark{15}{\char"1CD2}}द्वी\raisebox{0.25ex}{\accentmark{15}{\char"1CD2}}षो\raisebox{0.25ex}{\accentmark{12}{\char"0951}}जहि \nolinebreak\hskip 0.08em \mbox{॥३॥}}} \par}

\vspace{0.8em}
                % --- SUBSECTIONS (Mantras) ---
                
                % Conditional rendering based on output_mode
                \par\filbreak

\phantomsection

\addcontentsline{toc}{subsection}{ऋक् ४}

{\centering \textbf{\textcolor{AccentBlue}{वृषाह्या\raisebox{0.15ex}{\accentmark{12}{\char"1CF9}}सी\raisebox{0.25ex}{\accentmark{12}{\char"0951}}भानु\raisebox{0.25ex}{\accentmark{15}{\char"1CD2}}ना\kern 0.05em ।\hskip 0.12em plus 0.04em minus 0.02em द्यु\raisebox{0.25ex}{\accentmark{15}{\char"1CD2}}मन्तं\raisebox{0.25ex}{\accentmark{12}{\char"0951}}त्वाहवामहे\kern 0.05em ।\hskip 0.12em plus 0.04em minus 0.02em पा\raisebox{0.25ex}{\accentmark{15}{\char"1CD2}}वा\raisebox{0.25ex}{\accentmark{12}{\char"0951}}मानास्वर्दृ\raisebox{0.25ex}{\accentmark{15}{\char"1CD2}}श\raisebox{0.25ex}{\accentmark{12}{\char"0951}}म् \nolinebreak\hskip 0.08em \mbox{॥४॥}}} \par}

\vspace{0.8em}
                % --- SUBSECTIONS (Mantras) ---
                
                % Conditional rendering based on output_mode
                \par\filbreak

\phantomsection

\addcontentsline{toc}{subsection}{ऋक् ५}

{\centering \footnotesize \textcolor{AccentPurple}{॥ पवमानः सोमः गायत्री ॥} \par}

\vspace{0.12em}

{\centering \textbf{\textcolor{AccentBlue}{इ\raisebox{0.25ex}{\accentmark{15}{\char"1CD2}}न्दू\raisebox{0.25ex}{\accentmark{12}{\char"0951}}ःपविष्टाचे\raisebox{0.25ex}{\accentmark{15}{\char"1CD2}}तानः\kern 0.05em ।\hskip 0.12em plus 0.04em minus 0.02em प्रीय\raisebox{0.15ex}{\accentmark{12}{\char"1CF9}}ःका\raisebox{0.25ex}{\accentmark{12}{\char"0951}}वीना\raisebox{0.25ex}{\accentmark{15}{\char"1CD2}}म्मातिः\kern 0.05em ।\hskip 0.12em plus 0.04em minus 0.02em सृजद\raisebox{0.15ex}{\accentmark{12}{\char"1CF9}}श्वं\raisebox{0.25ex}{\accentmark{12}{\char"0951}}राथिरिव \nolinebreak\hskip 0.08em \mbox{॥५॥}}} \par}

\vspace{0.8em}
                % --- SUBSECTIONS (Mantras) ---
                
                % Conditional rendering based on output_mode
                \par\filbreak

\phantomsection

\addcontentsline{toc}{subsection}{ऋक् ६}

{\centering \footnotesize \textcolor{AccentPurple}{॥ कश्यपो मारीचः पवमानः सोमः गायत्री ॥} \par}

\vspace{0.12em}

{\centering \textbf{\textcolor{AccentBlue}{आ\raisebox{0.25ex}{\accentmark{15}{\char"1CD2}}सृक्षातप्रावाजी\raisebox{0.25ex}{\accentmark{15}{\char"1CD2}}नाः\kern 0.05em ।\hskip 0.12em plus 0.04em minus 0.02em गव्या\raisebox{0.25ex}{\accentmark{15}{\char"1CD2}}सोमा\raisebox{0.25ex}{\accentmark{12}{\char"0951}}सोआश्वाया\raisebox{0.25ex}{\accentmark{15}{\char"1CD2}}\kern 0.05em ।\hskip 0.12em plus 0.04em minus 0.02em शुक्रा\raisebox{0.25ex}{\accentmark{15}{\char"1CD2}}सो\raisebox{0.25ex}{\accentmark{12}{\char"0951}}वीराया\raisebox{0.25ex}{\accentmark{15}{\char"1CD2}}शा\raisebox{0.25ex}{\accentmark{15}{\char"1CD2}}वा\raisebox{0.25ex}{\accentmark{15}{\char"1CD2}}ः \nolinebreak\hskip 0.08em \mbox{॥६॥}}} \par}

\vspace{0.8em}
                % --- SUBSECTIONS (Mantras) ---
                
                % Conditional rendering based on output_mode
                \par\filbreak

\phantomsection

\addcontentsline{toc}{subsection}{ऋक् ७}

{\centering \footnotesize \textcolor{AccentPurple}{॥ निध्रुविः काश्यपः पवमानः सोमः गायत्री ॥} \par}

\vspace{0.12em}

{\centering \textbf{\textcolor{AccentBlue}{पा\raisebox{0.25ex}{\accentmark{15}{\char"1CD2}}वा\raisebox{0.25ex}{\accentmark{12}{\char"0951}}स्वादेवा\raisebox{0.15ex}{\accentmark{12}{\char"1CF9}}आयूषा\raisebox{0.25ex}{\accentmark{12}{\char"0951}}त्\kern 0.05em ।\hskip 0.12em plus 0.04em minus 0.02em इन्द्रं\raisebox{0.25ex}{\accentmark{12}{\char"0951}}गच्छतूतेमादा\raisebox{0.25ex}{\accentmark{12}{\char"0951}}ः\kern 0.05em ।\hskip 0.12em plus 0.04em minus 0.02em वायुमा\raisebox{0.15ex}{\accentmark{12}{\char"1CF9}}रो\raisebox{0.25ex}{\accentmark{12}{\char"0951}}हाध\raisebox{0.25ex}{\accentmark{15}{\char"1CD2}}र्मा\raisebox{0.25ex}{\accentmark{12}{\char"0951}}णा \nolinebreak\hskip 0.08em \mbox{॥७॥}}} \par}

\vspace{0.8em}
                % --- SUBSECTIONS (Mantras) ---
                
                % Conditional rendering based on output_mode
                \par\filbreak

\phantomsection

\addcontentsline{toc}{subsection}{ऋक् ८}

{\centering \footnotesize \textcolor{AccentPurple}{॥ अमहीयुराङ्गिरसः पवमानः सोमः गायत्री ॥} \par}

\vspace{0.12em}

{\centering \textbf{\textcolor{AccentBlue}{पा\raisebox{0.25ex}{\accentmark{15}{\char"1CD2}}वा\raisebox{0.25ex}{\accentmark{12}{\char"0951}}मानोअजीजनात्\kern 0.05em ।\hskip 0.12em plus 0.04em minus 0.02em दीव\raisebox{0.25ex}{\accentmark{15}{\char"1CD2}}श्चि\raisebox{0.25ex}{\accentmark{15}{\char"1CD2}}त्रन्ना\raisebox{0.15ex}{\accentmark{12}{\char"1CF9}}तन्या\raisebox{0.25ex}{\accentmark{12}{\char"0951}}तुम्\raisebox{0.25ex}{\accentmark{15}{\char"1CD2}}\kern 0.05em ।\hskip 0.12em plus 0.04em minus 0.02em ज्योतिर्वैश्वानारं\raisebox{0.25ex}{\accentmark{15}{\char"1CD2}}बृहा\raisebox{0.25ex}{\accentmark{15}{\char"1CD2}}त् \nolinebreak\hskip 0.08em \mbox{॥८॥}}} \par}

\vspace{0.8em}
                % --- SUBSECTIONS (Mantras) ---
                
                % Conditional rendering based on output_mode
                \par\filbreak

\phantomsection

\addcontentsline{toc}{subsection}{ऋक् ९}

{\centering \footnotesize \textcolor{AccentPurple}{॥ असितः काश्यपो देवलो वा पवमानः सोमः गायत्री ॥} \par}

\vspace{0.12em}

{\centering \textbf{\textcolor{AccentBlue}{पा\raisebox{0.15ex}{\accentmark{12}{\char"1CF8}}री\raisebox{0.25ex}{\accentmark{12}{\char"0951}}स्वाना\raisebox{0.25ex}{\accentmark{15}{\char"1CD2}}साइन्दा\raisebox{0.25ex}{\accentmark{12}{\char"0951}}वः\kern 0.05em ।\hskip 0.12em plus 0.04em minus 0.02em मा\raisebox{0.25ex}{\accentmark{15}{\char"1CD2}}दा\raisebox{0.25ex}{\accentmark{12}{\char"0951}}याबर्ह\raisebox{0.15ex}{\accentmark{12}{\char"1CF9}}णा\raisebox{0.25ex}{\accentmark{12}{\char"0951}}गीरा\kern 0.05em ।\hskip 0.12em plus 0.04em minus 0.02em मा\raisebox{0.25ex}{\accentmark{15}{\char"1CD2}}धो\raisebox{0.25ex}{\accentmark{12}{\char"0951}}अर्षन्तीधारा\raisebox{0.25ex}{\accentmark{12}{\char"0951}}या\raisebox{0.25ex}{\accentmark{15}{\char"1CD2}} \nolinebreak\hskip 0.08em \mbox{॥९॥}}} \par}

\vspace{0.8em}
                % --- SUBSECTIONS (Mantras) ---
                
                % Conditional rendering based on output_mode
                \par\filbreak

\phantomsection

\addcontentsline{toc}{subsection}{ऋक् १०}

{\centering \footnotesize \textcolor{AccentPurple}{॥ असितः काश्यपो देवलो वा पवमानः सोमः गायत्री ॥} \par}

\vspace{0.12em}

{\centering \textbf{\textcolor{AccentBlue}{पा\raisebox{0.25ex}{\accentmark{15}{\char"1CD2}}रीप्रासी\raisebox{0.25ex}{\accentmark{12}{\char"0951}}ष्यदद्काविः\kern 0.05em ।\hskip 0.12em plus 0.04em minus 0.02em सिन्धो\raisebox{0.25ex}{\accentmark{12}{\char"0951}}रूर्म\raisebox{0.15ex}{\accentmark{12}{\char"1CF9}}वा\raisebox{0.25ex}{\accentmark{15}{\char"1CD2}}धि\raisebox{0.25ex}{\accentmark{12}{\char"0951}}श्रीता\raisebox{0.25ex}{\accentmark{15}{\char"1CD2}}ः\kern 0.05em ।\hskip 0.12em plus 0.04em minus 0.02em कारूं\raisebox{0.25ex}{\accentmark{15}{\char"1CD2}}बि\raisebox{0.25ex}{\accentmark{15}{\char"1CD2}}भ्राट्यूरूस्पृ\raisebox{0.25ex}{\accentmark{15}{\char"1CD2}}ह\raisebox{0.25ex}{\accentmark{12}{\char"0951}}म् \nolinebreak\hskip 0.08em \mbox{॥१०॥}}} \par}

\vspace{0.8em}

            % --- End of Kandah Closing Floral Line & Counts ---
            \par\nopagebreak\vspace{0.9em plus 0.3em minus 0.1em}
            \centerline{\textcolor{black!35}{\rule{1.8cm}{0.35pt}}\;\raisebox{-0.22ex}{\textcolor{vedicmaroon}{\large\SixFlowerPetalDotted}}\;\textcolor{black!35}{\rule{1.8cm}{0.35pt}}}
            \nopagebreak\vspace{0.45em}
            \centerline{\devafont\textbf{\textcolor{AccentBlue}{॥ ऋक् सङ्ख्या : १० ॥}}}
        
            \clearpage

            % --- SECTION ---
            
            \phantomsection  

            \section*{\centering तृतीय खण्डः}
            \addcontentsline{toc}{section}{तृतीय खण्डः} 
            \markright{तृतीय खण्डः}
            \nopagebreak % Keep section title with first mantra
            
            % --- SECTION-LEVEL PROCEDURE LINK ---

                % --- SUBSECTIONS (Mantras) ---
                
                % Conditional rendering based on output_mode
                \par\filbreak

\phantomsection

\addcontentsline{toc}{subsection}{ऋक् १}

{\centering \footnotesize \textcolor{AccentPurple}{॥ अमहीयुराङ्गिरसः पवमानः सोमः गायत्री ॥} \par}

\vspace{0.12em}

{\centering \textbf{\textcolor{AccentBlue}{ऊ\raisebox{0.25ex}{\accentmark{15}{\char"1CD2}}पोषू\raisebox{0.25ex}{\accentmark{15}{\char"1CD2}}जा\raisebox{0.25ex}{\accentmark{15}{\char"1CD2}}तामप्तू\raisebox{0.25ex}{\accentmark{15}{\char"1CD2}}रम्\kern 0.05em ।\hskip 0.12em plus 0.04em minus 0.02em गो\raisebox{0.15ex}{\accentmark{12}{\char"1CF8}}भि\raisebox{0.25ex}{\accentmark{12}{\char"0951}}र्भङ्ग\raisebox{0.25ex}{\accentmark{15}{\char"1CD2}}म्पारी\raisebox{0.25ex}{\accentmark{12}{\char"0951}}ष्कृतम्\kern 0.05em ।\hskip 0.12em plus 0.04em minus 0.02em इ\raisebox{0.15ex}{\accentmark{12}{\char"1CF8}}न्दु\raisebox{0.25ex}{\accentmark{12}{\char"0951}}न्देवा\raisebox{0.25ex}{\accentmark{15}{\char"1CD2}}आ\raisebox{0.25ex}{\accentmark{12}{\char"0951}}यासिषुः \nolinebreak\hskip 0.08em \mbox{॥१॥}}} \par}

\vspace{0.8em}
                % --- SUBSECTIONS (Mantras) ---
                
                % Conditional rendering based on output_mode
                \par\filbreak

\phantomsection

\addcontentsline{toc}{subsection}{ऋक् २}

{\centering \footnotesize \textcolor{AccentPurple}{॥ बृहन्मतिराङ्गिरसः पवमानः सोमः गायत्री ॥} \par}

\vspace{0.12em}

{\centering \textbf{\textcolor{AccentBlue}{पुनानो\raisebox{0.25ex}{\accentmark{15}{\char"1CD2}}आ\raisebox{0.25ex}{\accentmark{12}{\char"0951}}क्रामीदाभि\raisebox{0.25ex}{\accentmark{15}{\char"1CD2}}ः\kern 0.05em ।\hskip 0.12em plus 0.04em minus 0.02em वि\raisebox{0.25ex}{\accentmark{15}{\char"1CD2}}श्वामृधोवी\raisebox{0.25ex}{\accentmark{15}{\char"1CD2}}चऋ\raisebox{0.25ex}{\accentmark{15}{\char"1CD2}}षणिः\kern 0.05em ।\hskip 0.12em plus 0.04em minus 0.02em शुम्भ\raisebox{0.25ex}{\accentmark{15}{\char"1CD2}}न्तीवी\raisebox{0.15ex}{\accentmark{12}{\char"1CF9}}प्र\raisebox{0.25ex}{\accentmark{12}{\char"0951}}न्धीतीभी\raisebox{0.25ex}{\accentmark{12}{\char"0951}}ः \nolinebreak\hskip 0.08em \mbox{॥२॥}}} \par}

\vspace{0.8em}
                % --- SUBSECTIONS (Mantras) ---
                
                % Conditional rendering based on output_mode
                \par\filbreak

\phantomsection

\addcontentsline{toc}{subsection}{ऋक् ३}

{\centering \footnotesize \textcolor{AccentPurple}{॥ जमदग्निर्भार्गवः पवमानः सोमः गायत्री ॥} \par}

\vspace{0.12em}

{\centering \textbf{\textcolor{AccentBlue}{आवीश\raisebox{0.25ex}{\accentmark{15}{\char"1CD2}}ङ्काला\raisebox{0.15ex}{\accentmark{12}{\char"1CF9}}शं\raisebox{0.25ex}{\accentmark{12}{\char"0951}}सू\raisebox{0.25ex}{\accentmark{12}{\char"0951}}त\raisebox{0.25ex}{\accentmark{15}{\char"1CD2}}ः\kern 0.05em ।\hskip 0.12em plus 0.04em minus 0.02em विश्वाअ\raisebox{0.15ex}{\accentmark{12}{\char"1CF9}}र्ष\raisebox{0.25ex}{\accentmark{12}{\char"0951}}न्ताभि\raisebox{0.25ex}{\accentmark{15}{\char"1CD2}}श्री\raisebox{0.25ex}{\accentmark{15}{\char"1CD2}}या\raisebox{0.25ex}{\accentmark{12}{\char"0951}}ः\kern 0.05em ।\hskip 0.12em plus 0.04em minus 0.02em इन्दूरि\raisebox{0.25ex}{\accentmark{15}{\char"1CD2}}न्द्रा\raisebox{0.25ex}{\accentmark{12}{\char"0951}}यधीयते \nolinebreak\hskip 0.08em \mbox{॥३॥}}} \par}

\vspace{0.8em}
                % --- SUBSECTIONS (Mantras) ---
                
                % Conditional rendering based on output_mode
                \par\filbreak

\phantomsection

\addcontentsline{toc}{subsection}{ऋक् ४}

{\centering \footnotesize \textcolor{AccentPurple}{॥ प्रभूवसुराङ्गिरसः पवमानः सोमः गायत्री ॥} \par}

\vspace{0.12em}

{\centering \textbf{\textcolor{AccentBlue}{आ\raisebox{0.15ex}{\accentmark{12}{\char"1CF8}}स\raisebox{0.25ex}{\accentmark{12}{\char"0951}}र्जिरथ्यो\raisebox{0.25ex}{\accentmark{15}{\char"1CD2}}या\raisebox{0.25ex}{\accentmark{15}{\char"1CD2}}था\raisebox{0.25ex}{\accentmark{12}{\char"0951}}\kern 0.05em ।\hskip 0.12em plus 0.04em minus 0.02em पावी\raisebox{0.15ex}{\accentmark{12}{\char"1CF9}}त्रेचमूवोठ\raisebox{0.25ex}{\accentmark{12}{\char"0951}}सूतः\kern 0.05em ।\hskip 0.12em plus 0.04em minus 0.02em का\raisebox{0.15ex}{\accentmark{12}{\char"1CF9}}र्ष्ण\raisebox{0.25ex}{\accentmark{12}{\char"0951}}न्वा\raisebox{0.25ex}{\accentmark{12}{\char"0951}}जीन्या\raisebox{0.25ex}{\accentmark{12}{\char"0951}}क्रमीत् \nolinebreak\hskip 0.08em \mbox{॥४॥}}} \par}

\vspace{0.8em}
                % --- SUBSECTIONS (Mantras) ---
                
                % Conditional rendering based on output_mode
                \par\filbreak

\phantomsection

\addcontentsline{toc}{subsection}{ऋक् ५}

{\centering \footnotesize \textcolor{AccentPurple}{॥ मेध्यातिथिः काण्वः पवमानः सोमः गायत्री ॥} \par}

\vspace{0.12em}

{\centering \textbf{\textcolor{AccentBlue}{प्र()य\raisebox{0.15ex}{\accentmark{12}{\char"1CF8}}द्गा\raisebox{0.25ex}{\accentmark{12}{\char"0951}}वोनभू\raisebox{0.25ex}{\accentmark{15}{\char"1CD2}}र्णायः\kern 0.05em ।\hskip 0.12em plus 0.04em minus 0.02em त्वेषा\raisebox{0.25ex}{\accentmark{15}{\char"1CD2}}आया\raisebox{0.25ex}{\accentmark{15}{\char"1CD2}}सोआ\raisebox{0.25ex}{\accentmark{15}{\char"1CD2}}क्रा\raisebox{0.25ex}{\accentmark{15}{\char"1CD2}}मुः\kern 0.05em ।\hskip 0.12em plus 0.04em minus 0.02em घ्न\raisebox{0.15ex}{\accentmark{12}{\char"1CF8}}न्त\raisebox{0.25ex}{\accentmark{12}{\char"0951}}ःकृष्णा\raisebox{0.25ex}{\accentmark{15}{\char"1CD2}}मापत्वा\raisebox{0.25ex}{\accentmark{15}{\char"1CD2}}च\raisebox{0.25ex}{\accentmark{12}{\char"0951}}म् \nolinebreak\hskip 0.08em \mbox{॥५॥}}} \par}

\vspace{0.8em}
                % --- SUBSECTIONS (Mantras) ---
                
                % Conditional rendering based on output_mode
                \par\filbreak

\phantomsection

\addcontentsline{toc}{subsection}{ऋक् ६}

{\centering \footnotesize \textcolor{AccentPurple}{॥ निध्रुविः काश्यपः पवमानः सोमः गायत्री ॥} \par}

\vspace{0.12em}

{\centering \textbf{\textcolor{AccentBlue}{अपध्नन्पा\raisebox{0.25ex}{\accentmark{12}{\char"0951}}वासेमृ\raisebox{0.25ex}{\accentmark{15}{\char"1CD2}}धाः\kern 0.05em ।\hskip 0.12em plus 0.04em minus 0.02em ऋतूवित्सो\raisebox{0.25ex}{\accentmark{12}{\char"0951}}मा\raisebox{0.25ex}{\accentmark{15}{\char"1CD2}}मात्सार\raisebox{0.25ex}{\accentmark{15}{\char"1CD2}}ः\kern 0.05em ।\hskip 0.12em plus 0.04em minus 0.02em नुतस्वा\raisebox{0.25ex}{\accentmark{15}{\char"1CD2}}दे\raisebox{0.25ex}{\accentmark{12}{\char"0951}}वायुञ्ज\raisebox{0.25ex}{\accentmark{15}{\char"1CD2}}न\raisebox{0.25ex}{\accentmark{12}{\char"0951}}म् \nolinebreak\hskip 0.08em \mbox{॥६॥}}} \par}

\vspace{0.8em}
                % --- SUBSECTIONS (Mantras) ---
                
                % Conditional rendering based on output_mode
                \par\filbreak

\phantomsection

\addcontentsline{toc}{subsection}{ऋक् ७}

{\centering \footnotesize \textcolor{AccentPurple}{॥ निध्रुविः काश्यपः पवमानः सोमः गायत्री ॥} \par}

\vspace{0.12em}

{\centering \textbf{\textcolor{AccentBlue}{आया\raisebox{0.25ex}{\accentmark{15}{\char"1CD2}}पा\raisebox{0.25ex}{\accentmark{12}{\char"0951}}वास्वाधा\raisebox{0.25ex}{\accentmark{15}{\char"1CD2}}रा\raisebox{0.25ex}{\accentmark{12}{\char"0951}}या\kern 0.05em ।\hskip 0.12em plus 0.04em minus 0.02em या\raisebox{0.25ex}{\accentmark{15}{\char"1CD2}}या\raisebox{0.25ex}{\accentmark{15}{\char"1CD2}}सू\raisebox{0.25ex}{\accentmark{15}{\char"1CD2}}र्यामरो\raisebox{0.25ex}{\accentmark{12}{\char"0951}}चयः\kern 0.05em ।\hskip 0.12em plus 0.04em minus 0.02em हिन्वानो\raisebox{0.25ex}{\accentmark{15}{\char"1CD2}}मानू\raisebox{0.25ex}{\accentmark{12}{\char"0951}}षीरापा\raisebox{0.25ex}{\accentmark{15}{\char"1CD2}}ः \nolinebreak\hskip 0.08em \mbox{॥७॥}}} \par}

\vspace{0.8em}
                % --- SUBSECTIONS (Mantras) ---
                
                % Conditional rendering based on output_mode
                \par\filbreak

\phantomsection

\addcontentsline{toc}{subsection}{ऋक् ८}

{\centering \footnotesize \textcolor{AccentPurple}{॥ अमहीयुराङ्गिरसः पवमानः सोमः गायत्री ॥} \par}

\vspace{0.12em}

{\centering \textbf{\textcolor{AccentBlue}{सापा\raisebox{0.25ex}{\accentmark{12}{\char"0951}}वास्वायआ\raisebox{0.25ex}{\accentmark{12}{\char"0951}}वी\raisebox{0.25ex}{\accentmark{12}{\char"0951}}था\kern 0.05em ।\hskip 0.12em plus 0.04em minus 0.02em इ\raisebox{0.15ex}{\accentmark{12}{\char"1CF8}}न्द्रं\raisebox{0.15ex}{\accentmark{12}{\char"1CF8}}वृत्राया\raisebox{0.25ex}{\accentmark{15}{\char"1CD2}}ह\raisebox{0.25ex}{\accentmark{15}{\char"1CD2}}न्ता\raisebox{0.25ex}{\accentmark{12}{\char"0951}}वे\kern 0.05em ।\hskip 0.12em plus 0.04em minus 0.02em वब्रीवां\raisebox{0.15ex}{\accentmark{12}{\char"1CF9}}स\raisebox{0.25ex}{\accentmark{12}{\char"0951}}म्माहीराप\raisebox{0.25ex}{\accentmark{15}{\char"1CD2}}ः \nolinebreak\hskip 0.08em \mbox{॥८॥}}} \par}

\vspace{0.8em}
                % --- SUBSECTIONS (Mantras) ---
                
                % Conditional rendering based on output_mode
                \par\filbreak

\phantomsection

\addcontentsline{toc}{subsection}{ऋक् ९}

{\centering \footnotesize \textcolor{AccentPurple}{॥ अमहीयुराङ्गिरसः पवमानः सोमः गायत्री ॥} \par}

\vspace{0.12em}

{\centering \textbf{\textcolor{AccentBlue}{आया\raisebox{0.25ex}{\accentmark{15}{\char"1CD2}}वीतीपा\raisebox{0.25ex}{\accentmark{15}{\char"1CD2}}रिस्रावा\kern 0.05em ।\hskip 0.12em plus 0.04em minus 0.02em यस्ताइन्दो\raisebox{0.25ex}{\accentmark{12}{\char"0951}}मातेष्वा\kern 0.05em ।\hskip 0.12em plus 0.04em minus 0.02em आवा\raisebox{0.25ex}{\accentmark{15}{\char"1CD2}}ह\raisebox{0.25ex}{\accentmark{12}{\char"0951}}न्नावातीर्ना\raisebox{0.25ex}{\accentmark{15}{\char"1CD2}}वा \nolinebreak\hskip 0.08em \mbox{॥९॥}}} \par}

\vspace{0.8em}
                % --- SUBSECTIONS (Mantras) ---
                
                % Conditional rendering based on output_mode
                \par\filbreak

\phantomsection

\addcontentsline{toc}{subsection}{ऋक् १०}

{\centering \footnotesize \textcolor{AccentPurple}{॥ उचथ्य आङ्गिरसः पवमानः सोमः गायत्री ॥} \par}

\vspace{0.12em}

{\centering \textbf{\textcolor{AccentBlue}{पारी\raisebox{0.25ex}{\accentmark{12}{\char"0951}}धुक्षं\raisebox{0.15ex}{\accentmark{12}{\char"1CF9}}सा\raisebox{0.15ex}{\accentmark{12}{\char"1CF8}}ना\raisebox{0.25ex}{\accentmark{12}{\char"0951}}द्रायि\raisebox{0.25ex}{\accentmark{15}{\char"1CD2}}म्\kern 0.05em ।\hskip 0.12em plus 0.04em minus 0.02em भारद्वा\raisebox{0.15ex}{\accentmark{12}{\char"1CF9}}ज\raisebox{0.25ex}{\accentmark{12}{\char"0951}}न्नोअ\raisebox{0.25ex}{\accentmark{15}{\char"1CD2}}न्धा\raisebox{0.25ex}{\accentmark{12}{\char"0951}}सा\kern 0.05em ।\hskip 0.12em plus 0.04em minus 0.02em स्वानोअ\raisebox{0.25ex}{\accentmark{12}{\char"0951}}र्षपावी\raisebox{0.25ex}{\accentmark{15}{\char"1CD2}}त्राआ\raisebox{0.25ex}{\accentmark{15}{\char"1CD2}} \nolinebreak\hskip 0.08em \mbox{॥१०॥}}} \par}

\vspace{0.8em}

            % --- End of Kandah Closing Floral Line & Counts ---
            \par\nopagebreak\vspace{0.9em plus 0.3em minus 0.1em}
            \centerline{\textcolor{black!35}{\rule{1.8cm}{0.35pt}}\;\raisebox{-0.22ex}{\textcolor{vedicmaroon}{\large\SixFlowerPetalDotted}}\;\textcolor{black!35}{\rule{1.8cm}{0.35pt}}}
            \nopagebreak\vspace{0.45em}
            \centerline{\devafont\textbf{\textcolor{AccentBlue}{॥ ऋक् सङ्ख्या : १० ॥}}}
        
            \clearpage

            % --- SECTION ---
            
            \phantomsection  

            \section*{\centering चतुर्थ खण्डः}
            \addcontentsline{toc}{section}{चतुर्थ खण्डः} 
            \markright{चतुर्थ खण्डः}
            \nopagebreak % Keep section title with first mantra
            
            % --- SECTION-LEVEL PROCEDURE LINK ---

                % --- SUBSECTIONS (Mantras) ---
                
                % Conditional rendering based on output_mode
                \par\filbreak

\phantomsection

\addcontentsline{toc}{subsection}{ऋक् १}

{\centering \footnotesize \textcolor{AccentPurple}{॥ मेधातिथिः काण्वः पवमानः सोमः गायत्री ॥} \par}

\vspace{0.12em}

{\centering \textbf{\textcolor{AccentBlue}{आ\raisebox{0.25ex}{\accentmark{15}{\char"1CD2}}चिक्रादत्वृषाहा\raisebox{0.25ex}{\accentmark{15}{\char"1CD2}}री\raisebox{0.25ex}{\accentmark{12}{\char"0951}}ः\kern 0.05em ।\hskip 0.12em plus 0.04em minus 0.02em माहा\raisebox{0.25ex}{\accentmark{15}{\char"1CD2}}न्मित्रो\raisebox{0.15ex}{\accentmark{12}{\char"1CF9}}ना\raisebox{0.15ex}{\accentmark{12}{\char"1CF9}}द\raisebox{0.25ex}{\accentmark{12}{\char"0951}}र्शातः\kern 0.05em ।\hskip 0.12em plus 0.04em minus 0.02em संसू\raisebox{0.25ex}{\accentmark{15}{\char"1CD2}}र्येणदिद्युते \nolinebreak\hskip 0.08em \mbox{॥१॥}}} \par}

\vspace{0.8em}
                % --- SUBSECTIONS (Mantras) ---
                
                % Conditional rendering based on output_mode
                \par\filbreak

\phantomsection

\addcontentsline{toc}{subsection}{ऋक् २}

{\centering \footnotesize \textcolor{AccentPurple}{॥ भृगुर्वारुणिर्जमदग्निर्भार्गवो वा पवमानः सोमः गायत्री ॥} \par}

\vspace{0.12em}

{\centering \textbf{\textcolor{AccentBlue}{आ\raisebox{0.25ex}{\accentmark{15}{\char"1CD2}}तेदक्षं\raisebox{0.25ex}{\accentmark{12}{\char"0951}}मायोभू\raisebox{0.25ex}{\accentmark{15}{\char"1CD2}}व\raisebox{0.25ex}{\accentmark{12}{\char"0951}}म्\kern 0.05em ।\hskip 0.12em plus 0.04em minus 0.02em व\raisebox{0.25ex}{\accentmark{15}{\char"1CD2}}न्ही\raisebox{0.25ex}{\accentmark{15}{\char"1CD2}}मद्या\raisebox{0.25ex}{\accentmark{15}{\char"1CD2}}वृ\raisebox{0.25ex}{\accentmark{12}{\char"0951}}णीमहे\raisebox{0.25ex}{\accentmark{15}{\char"1CD2}}\kern 0.05em ।\hskip 0.12em plus 0.04em minus 0.02em पा\raisebox{0.25ex}{\accentmark{15}{\char"1CD2}}न्तामा\raisebox{0.15ex}{\accentmark{12}{\char"1CF9}}पू\raisebox{0.25ex}{\accentmark{12}{\char"0951}}रूस्पृह\raisebox{0.25ex}{\accentmark{12}{\char"0951}}म् \nolinebreak\hskip 0.08em \mbox{॥२॥}}} \par}

\vspace{0.8em}
                % --- SUBSECTIONS (Mantras) ---
                
                % Conditional rendering based on output_mode
                \par\filbreak

\phantomsection

\addcontentsline{toc}{subsection}{ऋक् ३}

{\centering \footnotesize \textcolor{AccentPurple}{॥ उचथ्य आङ्गिरसः पवमानः सोमः गायत्री ॥} \par}

\vspace{0.12em}

{\centering \textbf{\textcolor{AccentBlue}{अ\raisebox{0.15ex}{\accentmark{12}{\char"1CF8}}ध्व\raisebox{0.25ex}{\accentmark{12}{\char"0951}}र्योआ\raisebox{0.25ex}{\accentmark{15}{\char"1CD2}}द्री\raisebox{0.25ex}{\accentmark{12}{\char"0951}}भीस्सूतम्\kern 0.05em ।\hskip 0.12em plus 0.04em minus 0.02em सो\raisebox{0.15ex}{\accentmark{12}{\char"1CF8}}मं\raisebox{0.25ex}{\accentmark{12}{\char"0951}}पावी\raisebox{0.25ex}{\accentmark{15}{\char"1CD2}}त्रा\raisebox{0.25ex}{\accentmark{15}{\char"1CD2}}आ\raisebox{0.25ex}{\accentmark{15}{\char"1CD2}}नाया\kern 0.05em ।\hskip 0.12em plus 0.04em minus 0.02em पुनाही\raisebox{0.15ex}{\accentmark{12}{\char"1CF9}}न्द्रा\raisebox{0.25ex}{\accentmark{12}{\char"0951}}यापा\raisebox{0.25ex}{\accentmark{15}{\char"1CD2}}ता\raisebox{0.25ex}{\accentmark{12}{\char"0951}}वे \nolinebreak\hskip 0.08em \mbox{॥३॥}}} \par}

\vspace{0.8em}
                % --- SUBSECTIONS (Mantras) ---
                
                % Conditional rendering based on output_mode
                \par\filbreak

\phantomsection

\addcontentsline{toc}{subsection}{ऋक् ४}

{\centering \footnotesize \textcolor{AccentPurple}{॥ अवत्सारः काश्यपः पवमानः सोमः गायत्री ॥} \par}

\vspace{0.12em}

{\centering \textbf{\textcolor{AccentBlue}{ता\raisebox{0.25ex}{\accentmark{15}{\char"1CD2}}रत्सा\raisebox{0.25ex}{\accentmark{15}{\char"1CD2}}मन्दीधा\raisebox{0.25ex}{\accentmark{12}{\char"0951}}वती\kern 0.05em ।\hskip 0.12em plus 0.04em minus 0.02em धा\raisebox{0.15ex}{\accentmark{12}{\char"1CF8}}रासूतस्या\raisebox{0.25ex}{\accentmark{15}{\char"1CD2}}न्धा\raisebox{0.25ex}{\accentmark{12}{\char"0951}}सः\kern 0.05em ।\hskip 0.12em plus 0.04em minus 0.02em ता\raisebox{0.25ex}{\accentmark{15}{\char"1CD2}}रत्सा\raisebox{0.25ex}{\accentmark{15}{\char"1CD2}}मन्दीधावती \nolinebreak\hskip 0.08em \mbox{॥४॥}}} \par}

\vspace{0.8em}
                % --- SUBSECTIONS (Mantras) ---
                
                % Conditional rendering based on output_mode
                \par\filbreak

\phantomsection

\addcontentsline{toc}{subsection}{ऋक् ५}

{\centering \footnotesize \textcolor{AccentPurple}{॥ निध्रुविः काश्यपः पवमानः सोमः गायत्री ॥} \par}

\vspace{0.12em}

{\centering \textbf{\textcolor{AccentBlue}{आ\raisebox{0.25ex}{\accentmark{15}{\char"1CD2}}पा\raisebox{0.25ex}{\accentmark{12}{\char"0951}}वस्वा\raisebox{0.25ex}{\accentmark{15}{\char"1CD2}}साहास्री\raisebox{0.25ex}{\accentmark{15}{\char"1CD2}}णम्\kern 0.05em ।\hskip 0.12em plus 0.04em minus 0.02em रा\raisebox{0.25ex}{\accentmark{15}{\char"1CD2}}यिंसो\raisebox{0.25ex}{\accentmark{12}{\char"0951}}मासूवी\raisebox{0.25ex}{\accentmark{15}{\char"1CD2}}र्यम्\kern 0.05em ।\hskip 0.12em plus 0.04em minus 0.02em अस्मेश्रा\raisebox{0.25ex}{\accentmark{15}{\char"1CD2}}वां\raisebox{0.15ex}{\accentmark{12}{\char"1CF8}}सिधारया \nolinebreak\hskip 0.08em \mbox{॥५॥}}} \par}

\vspace{0.8em}
                % --- SUBSECTIONS (Mantras) ---
                
                % Conditional rendering based on output_mode
                \par\filbreak

\phantomsection

\addcontentsline{toc}{subsection}{ऋक् ६}

{\centering \footnotesize \textcolor{AccentPurple}{॥ असितः काश्यपो देवलो वा पवमानः सोमः गायत्री ॥} \par}

\vspace{0.12em}

{\centering \textbf{\textcolor{AccentBlue}{आनूप्रत्ना\raisebox{0.15ex}{\accentmark{12}{\char"1CF9}}सा\raisebox{0.25ex}{\accentmark{12}{\char"0951}}आया\raisebox{0.25ex}{\accentmark{15}{\char"1CD2}}वा\raisebox{0.25ex}{\accentmark{12}{\char"0951}}ः\kern 0.05em ।\hskip 0.12em plus 0.04em minus 0.02em पादन्ना\raisebox{0.25ex}{\accentmark{15}{\char"1CD2}}वी\raisebox{0.25ex}{\accentmark{12}{\char"0951}}योअक्रमुः\kern 0.05em ।\hskip 0.12em plus 0.04em minus 0.02em रुचे\raisebox{0.25ex}{\accentmark{15}{\char"1CD2}}ज\raisebox{0.25ex}{\accentmark{12}{\char"0951}}नान्तासूर्य\raisebox{0.25ex}{\accentmark{12}{\char"0951}}म् \nolinebreak\hskip 0.08em \mbox{॥६॥}}} \par}

\vspace{0.8em}
                % --- SUBSECTIONS (Mantras) ---
                
                % Conditional rendering based on output_mode
                \par\filbreak

\phantomsection

\addcontentsline{toc}{subsection}{ऋक् ७}

{\centering \footnotesize \textcolor{AccentPurple}{॥ भृगुर्वारुणिर्जमदग्निर्भार्गवो वा पवमानः सोमः गायत्री ॥} \par}

\vspace{0.12em}

{\centering \textbf{\textcolor{AccentBlue}{अ\raisebox{0.25ex}{\accentmark{12}{\char"0951}}र्षासोमा\raisebox{0.25ex}{\accentmark{15}{\char"1CD2}}धूमात्ता\raisebox{0.25ex}{\accentmark{12}{\char"0951}}मः\kern 0.05em ।\hskip 0.12em plus 0.04em minus 0.02em आभीद्रो\raisebox{0.15ex}{\accentmark{12}{\char"1CF9}}णा\raisebox{0.25ex}{\accentmark{12}{\char"0951}}नीरो\raisebox{0.25ex}{\accentmark{15}{\char"1CD2}}रूवात्\kern 0.05em ।\hskip 0.12em plus 0.04em minus 0.02em सी\raisebox{0.25ex}{\accentmark{15}{\char"1CD2}}दन्यो\raisebox{0.25ex}{\accentmark{15}{\char"1CD2}}नौवा\raisebox{0.25ex}{\accentmark{15}{\char"1CD2}}नेष्वा \nolinebreak\hskip 0.08em \mbox{॥७॥}}} \par}

\vspace{0.8em}
                % --- SUBSECTIONS (Mantras) ---
                
                % Conditional rendering based on output_mode
                \par\filbreak

\phantomsection

\addcontentsline{toc}{subsection}{ऋक् ८}

{\centering \footnotesize \textcolor{AccentPurple}{॥ कश्यपो मारीचः पवमानः सोमः गायत्री ॥} \par}

\vspace{0.12em}

{\centering \textbf{\textcolor{AccentBlue}{वृ\raisebox{0.25ex}{\accentmark{15}{\char"1CD2}}षा\raisebox{0.25ex}{\accentmark{12}{\char"0951}}सोमाधूमंआ\raisebox{0.25ex}{\accentmark{12}{\char"0951}}सी\kern 0.05em ।\hskip 0.12em plus 0.04em minus 0.02em वृषा\raisebox{0.25ex}{\accentmark{15}{\char"1CD2}}देवावृ\raisebox{0.25ex}{\accentmark{15}{\char"1CD2}}षाव्रा\raisebox{0.25ex}{\accentmark{12}{\char"0951}}तः\kern 0.05em ।\hskip 0.12em plus 0.04em minus 0.02em वृषाध\raisebox{0.25ex}{\accentmark{15}{\char"1CD2}}र्मा\raisebox{0.25ex}{\accentmark{12}{\char"0951}}णिदध्री\raisebox{0.25ex}{\accentmark{15}{\char"1CD2}}षे\raisebox{0.25ex}{\accentmark{15}{\char"1CD2}} \nolinebreak\hskip 0.08em \mbox{॥८॥}}} \par}

\vspace{0.8em}
                % --- SUBSECTIONS (Mantras) ---
                
                % Conditional rendering based on output_mode
                \par\filbreak

\phantomsection

\addcontentsline{toc}{subsection}{ऋक् ९}

{\centering \footnotesize \textcolor{AccentPurple}{॥ कश्यपो मारीचः पवमानः सोमः गायत्री ॥} \par}

\vspace{0.12em}

{\centering \textbf{\textcolor{AccentBlue}{ईषे\raisebox{0.25ex}{\accentmark{15}{\char"1CD2}}पा\raisebox{0.25ex}{\accentmark{12}{\char"0951}}वा\raisebox{0.15ex}{\accentmark{12}{\char"1CF8}}स्वाधा\raisebox{0.25ex}{\accentmark{15}{\char"1CD2}}रा\raisebox{0.25ex}{\accentmark{12}{\char"0951}}या\kern 0.05em ।\hskip 0.12em plus 0.04em minus 0.02em मृज्या\raisebox{0.25ex}{\accentmark{15}{\char"1CD2}}मा\raisebox{0.25ex}{\accentmark{12}{\char"0951}}नोमानीषी\raisebox{0.25ex}{\accentmark{15}{\char"1CD2}}भि\raisebox{0.25ex}{\accentmark{12}{\char"0951}}ः\kern 0.05em ।\hskip 0.12em plus 0.04em minus 0.02em इ\raisebox{0.15ex}{\accentmark{12}{\char"1CF8}}न्दो\raisebox{0.25ex}{\accentmark{12}{\char"0951}}रूचा\raisebox{0.25ex}{\accentmark{15}{\char"1CD2}}भीगा\raisebox{0.25ex}{\accentmark{15}{\char"1CD2}}ईहि \nolinebreak\hskip 0.08em \mbox{॥९॥}}} \par}

\vspace{0.8em}
                % --- SUBSECTIONS (Mantras) ---
                
                % Conditional rendering based on output_mode
                \par\filbreak

\phantomsection

\addcontentsline{toc}{subsection}{ऋक् १०}

{\centering \footnotesize \textcolor{AccentPurple}{॥ असितः काश्यपो देवलो वा पवमानः सोमः गायत्री ॥} \par}

\vspace{0.12em}

{\centering \textbf{\textcolor{AccentBlue}{मन्द्राया\raisebox{0.25ex}{\accentmark{12}{\char"0951}}सोमाधा\raisebox{0.25ex}{\accentmark{15}{\char"1CD2}}रा\raisebox{0.25ex}{\accentmark{12}{\char"0951}}य\kern 0.05em ।\hskip 0.12em plus 0.04em minus 0.02em वृषा\raisebox{0.25ex}{\accentmark{12}{\char"0951}}पवस्वदे\raisebox{0.25ex}{\accentmark{15}{\char"1CD2}}वा\raisebox{0.25ex}{\accentmark{15}{\char"1CD2}}युः\kern 0.05em ।\hskip 0.12em plus 0.04em minus 0.02em अव्यावा\raisebox{0.25ex}{\accentmark{15}{\char"1CD2}}रे\raisebox{0.25ex}{\accentmark{12}{\char"0951}}भिरास्मायुः \nolinebreak\hskip 0.08em \mbox{॥१०॥}}} \par}

\vspace{0.8em}
                % --- SUBSECTIONS (Mantras) ---
                
                % Conditional rendering based on output_mode
                \par\filbreak

\phantomsection

\addcontentsline{toc}{subsection}{ऋक् ११}

{\centering \footnotesize \textcolor{AccentPurple}{॥ कविर्भार्गवः पवमानः सोमः गायत्री ॥} \par}

\vspace{0.12em}

{\centering \textbf{\textcolor{AccentBlue}{आया\raisebox{0.25ex}{\accentmark{15}{\char"1CD2}}सोमसूकृत्या\raisebox{0.25ex}{\accentmark{15}{\char"1CD2}}या\kern 0.05em ।\hskip 0.12em plus 0.04em minus 0.02em माहा\raisebox{0.25ex}{\accentmark{15}{\char"1CD2}}न्सन्नाभ्या\raisebox{0.25ex}{\accentmark{12}{\char"0951}}वर्धथाः\kern 0.05em ।\hskip 0.12em plus 0.04em minus 0.02em म\raisebox{0.15ex}{\accentmark{12}{\char"1CF8}}न्दा\raisebox{0.25ex}{\accentmark{12}{\char"0951}}नाइ\raisebox{0.25ex}{\accentmark{12}{\char"0951}}द्वृषायसे \nolinebreak\hskip 0.08em \mbox{॥११॥}}} \par}

\vspace{0.8em}
                % --- SUBSECTIONS (Mantras) ---
                
                % Conditional rendering based on output_mode
                \par\filbreak

\phantomsection

\addcontentsline{toc}{subsection}{ऋक् १२}

{\centering \footnotesize \textcolor{AccentPurple}{॥ जमदग्निर्भार्गवः पवमानः सोमः गायत्री ॥} \par}

\vspace{0.12em}

{\centering \textbf{\textcolor{AccentBlue}{आया\raisebox{0.25ex}{\accentmark{15}{\char"1CD2}}वीचर्षणीर्हित\raisebox{0.25ex}{\accentmark{15}{\char"1CD2}}ः\kern 0.05em ।\hskip 0.12em plus 0.04em minus 0.02em पावा\raisebox{0.25ex}{\accentmark{12}{\char"0951}}मानस्साचे\raisebox{0.25ex}{\accentmark{12}{\char"0951}}तति\kern 0.05em ।\hskip 0.12em plus 0.04em minus 0.02em हिन्वाना\raisebox{0.15ex}{\accentmark{12}{\char"1CF9}}आ\raisebox{0.15ex}{\accentmark{12}{\char"1CF8}}प्यंबृहात् \nolinebreak\hskip 0.08em \mbox{॥१२॥}}} \par}

\vspace{0.8em}
                % --- SUBSECTIONS (Mantras) ---
                
                % Conditional rendering based on output_mode
                \par\filbreak

\phantomsection

\addcontentsline{toc}{subsection}{ऋक् १३}

{\centering \footnotesize \textcolor{AccentPurple}{॥ अयास्य आङ्गिरसः पवमानः सोमः गायत्री ॥} \par}

\vspace{0.12em}

{\centering \textbf{\textcolor{AccentBlue}{प्राणा\raisebox{0.25ex}{\accentmark{12}{\char"0951}}इन्दोमाहेतू\raisebox{0.25ex}{\accentmark{15}{\char"1CD2}}ना\raisebox{0.25ex}{\accentmark{12}{\char"0951}}ः\kern 0.05em ।\hskip 0.12em plus 0.04em minus 0.02em ऊर्म्मि\raisebox{0.25ex}{\accentmark{12}{\char"0951}}न्नाबिभ्रादर्ष\raisebox{0.25ex}{\accentmark{12}{\char"0951}}सि\kern 0.05em ।\hskip 0.12em plus 0.04em minus 0.02em आभिदेवं\raisebox{0.25ex}{\accentmark{15}{\char"1CD2}}आयास्या\raisebox{0.25ex}{\accentmark{12}{\char"0951}}ः \nolinebreak\hskip 0.08em \mbox{॥१३॥}}} \par}

\vspace{0.8em}
                % --- SUBSECTIONS (Mantras) ---
                
                % Conditional rendering based on output_mode
                \par\filbreak

\phantomsection

\addcontentsline{toc}{subsection}{ऋक् १४}

{\centering \footnotesize \textcolor{AccentPurple}{॥ आमहीयुराङ्गिरसः पवमानः सोमः गायत्री ॥} \par}

\vspace{0.12em}

{\centering \textbf{\textcolor{AccentBlue}{अपध्नन्पा\raisebox{0.25ex}{\accentmark{12}{\char"0951}}वाते\raisebox{0.25ex}{\accentmark{15}{\char"1CD2}}मृधा\raisebox{0.25ex}{\accentmark{12}{\char"0951}}ः\kern 0.05em ।\hskip 0.12em plus 0.04em minus 0.02em आ\raisebox{0.25ex}{\accentmark{15}{\char"1CD2}}प\raisebox{0.15ex}{\accentmark{12}{\char"1CF8}}सो\raisebox{0.25ex}{\accentmark{15}{\char"1CD2}}मोआ\raisebox{0.25ex}{\accentmark{15}{\char"1CD2}}राव्ण\raisebox{0.25ex}{\accentmark{15}{\char"1CD2}}ः\kern 0.05em ।\hskip 0.12em plus 0.04em minus 0.02em गच्छन्निन्द्रा\raisebox{0.25ex}{\accentmark{12}{\char"0951}}स्यनिष्कृतम् \nolinebreak\hskip 0.08em \mbox{॥१४॥}}} \par}

\vspace{0.8em}

            % --- End of Kandah Closing Floral Line & Counts ---
            \par\nopagebreak\vspace{0.9em plus 0.3em minus 0.1em}
            \centerline{\textcolor{black!35}{\rule{1.8cm}{0.35pt}}\;\raisebox{-0.22ex}{\textcolor{vedicmaroon}{\large\SixFlowerPetalDotted}}\;\textcolor{black!35}{\rule{1.8cm}{0.35pt}}}
            \nopagebreak\vspace{0.45em}
            \centerline{\devafont\textbf{\textcolor{AccentBlue}{॥ ऋक् सङ्ख्या : १४ ॥}}}
        
            \clearpage

            % --- SECTION ---
            
            \phantomsection  

            \section*{\centering पञ्चम खण्डः}
            \addcontentsline{toc}{section}{पञ्चम खण्डः} 
            \markright{पञ्चम खण्डः}
            \nopagebreak % Keep section title with first mantra
            
            % --- SECTION-LEVEL PROCEDURE LINK ---

                % --- SUBSECTIONS (Mantras) ---
                
                % Conditional rendering based on output_mode
                \par\filbreak

\phantomsection

\addcontentsline{toc}{subsection}{ऋक् १}

{\centering \footnotesize \textcolor{AccentPurple}{॥ सप्तर्षयः (भरद्वाजो बार्हस्पत्यः) पवमानः सोमः बृहती ॥} \par}

\vspace{0.12em}

{\centering \textbf{\textcolor{AccentBlue}{पुना\raisebox{0.15ex}{\accentmark{12}{\char"1CF9}}नस्सो\raisebox{0.25ex}{\accentmark{12}{\char"0951}}माधा\raisebox{0.25ex}{\accentmark{15}{\char"1CD2}}रा\raisebox{0.25ex}{\accentmark{12}{\char"0951}}या\kern 0.05em ।\hskip 0.12em plus 0.04em minus 0.02em आपो\raisebox{0.25ex}{\accentmark{12}{\char"0951}}वा\raisebox{0.25ex}{\accentmark{15}{\char"1CD2}}सा\raisebox{0.25ex}{\accentmark{12}{\char"0951}}नोअर्ष\raisebox{0.25ex}{\accentmark{12}{\char"0951}}सि\kern 0.05em ।\hskip 0.12em plus 0.04em minus 0.02em आर\raisebox{0.25ex}{\accentmark{12}{\char"0951}}त्नाधा\raisebox{0.15ex}{\accentmark{12}{\char"1CF9}}यो\raisebox{0.15ex}{\accentmark{12}{\char"1CF9}}नीमृतास्या\raisebox{0.25ex}{\accentmark{12}{\char"0951}}सीदसि\kern 0.05em ।\hskip 0.12em plus 0.04em minus 0.02em उत्सोदेवो\raisebox{0.15ex}{\accentmark{12}{\char"1CF9}}ही\raisebox{0.25ex}{\accentmark{12}{\char"0951}}रण्या\raisebox{0.25ex}{\accentmark{12}{\char"0951}}या \nolinebreak\hskip 0.08em \mbox{॥१॥}}} \par}

\vspace{0.8em}
                % --- SUBSECTIONS (Mantras) ---
                
                % Conditional rendering based on output_mode
                \par\filbreak

\phantomsection

\addcontentsline{toc}{subsection}{ऋक् २}

{\centering \footnotesize \textcolor{AccentPurple}{॥ सप्तर्षयः (कश्यपो मारीचः) पवमानः सोमः बृहती ॥} \par}

\vspace{0.12em}

{\centering \textbf{\textcolor{AccentBlue}{पा\raisebox{0.25ex}{\accentmark{15}{\char"1CD2}}रीतोषी\raisebox{0.25ex}{\accentmark{12}{\char"0951}}ञ्चतासू\raisebox{0.25ex}{\accentmark{15}{\char"1CD2}}तम्\kern 0.05em ।\hskip 0.12em plus 0.04em minus 0.02em सोमोयाऊ\raisebox{0.25ex}{\accentmark{12}{\char"0951}}त्तामंहा\raisebox{0.25ex}{\accentmark{15}{\char"1CD2}}विः\kern 0.05em ।\hskip 0.12em plus 0.04em minus 0.02em दधन्वायो\raisebox{0.15ex}{\accentmark{12}{\char"1CF9}}न\raisebox{0.15ex}{\accentmark{12}{\char"1CF9}}र्योअप्सू\raisebox{0.25ex}{\accentmark{15}{\char"1CD2}}आट्यन्तारा\raisebox{0.25ex}{\accentmark{15}{\char"1CD2}}\kern 0.05em ।\hskip 0.12em plus 0.04em minus 0.02em सूषा\raisebox{0.25ex}{\accentmark{15}{\char"1CD2}}वासो\raisebox{0.25ex}{\accentmark{15}{\char"1CD2}}मामा\raisebox{0.25ex}{\accentmark{15}{\char"1CD2}}द्री\raisebox{0.25ex}{\accentmark{12}{\char"0951}}भिः \nolinebreak\hskip 0.08em \mbox{॥२॥}}} \par}

\vspace{0.8em}
                % --- SUBSECTIONS (Mantras) ---
                
                % Conditional rendering based on output_mode
                \par\filbreak

\phantomsection

\addcontentsline{toc}{subsection}{ऋक् ३}

{\centering \footnotesize \textcolor{AccentPurple}{॥ सप्तर्षयः (गोतमो राहूगणः) पवमानः सोमः बृहती ॥} \par}

\vspace{0.12em}

{\centering \textbf{\textcolor{AccentBlue}{आ\raisebox{0.25ex}{\accentmark{15}{\char"1CD2}}सो\raisebox{0.25ex}{\accentmark{12}{\char"0951}}मस्वानो\raisebox{0.25ex}{\accentmark{15}{\char"1CD2}}अद्री\raisebox{0.25ex}{\accentmark{12}{\char"0951}}भिः\kern 0.05em ।\hskip 0.12em plus 0.04em minus 0.02em तीरोवा\raisebox{0.25ex}{\accentmark{12}{\char"0951}}राण्यव्या\raisebox{0.25ex}{\accentmark{12}{\char"0951}}या\kern 0.05em ।\hskip 0.12em plus 0.04em minus 0.02em जनोना\raisebox{0.25ex}{\accentmark{15}{\char"1CD2}}पूरीचम्मुवो१र्वीशद्धा\raisebox{0.25ex}{\accentmark{15}{\char"1CD2}}रिः\kern 0.05em ।\hskip 0.12em plus 0.04em minus 0.02em सादावा\raisebox{0.25ex}{\accentmark{15}{\char"1CD2}}नेषुद॑\raisebox{0.25ex}{\accentmark{12}{\char"0951}}ध्रीषे \nolinebreak\hskip 0.08em \mbox{॥३॥}}} \par}

\vspace{0.8em}
                % --- SUBSECTIONS (Mantras) ---
                
                % Conditional rendering based on output_mode
                \par\filbreak

\phantomsection

\addcontentsline{toc}{subsection}{ऋक् ४}

{\centering \footnotesize \textcolor{AccentPurple}{॥ सप्तर्षयः (अत्रिर्भीमः) पवमानः सोमः बृहती ॥} \par}

\vspace{0.12em}

{\centering \textbf{\textcolor{AccentBlue}{प्रा\raisebox{0.15ex}{\accentmark{12}{\char"1CF8}}सो\raisebox{0.25ex}{\accentmark{12}{\char"0951}}मादेवा\raisebox{0.25ex}{\accentmark{15}{\char"1CD2}}वी\raisebox{0.25ex}{\accentmark{12}{\char"0951}}तये\kern 0.05em ।\hskip 0.12em plus 0.04em minus 0.02em सिन्धुर्न्ना\raisebox{0.15ex}{\accentmark{12}{\char"1CF9}}पी\raisebox{0.25ex}{\accentmark{12}{\char"0951}}प्येअ\raisebox{0.25ex}{\accentmark{15}{\char"1CD2}}र्णासा\kern 0.05em ।\hskip 0.12em plus 0.04em minus 0.02em अं\raisebox{0.15ex}{\accentmark{12}{\char"1CF8}}शोःपा\raisebox{0.25ex}{\accentmark{15}{\char"1CD2}}या\raisebox{0.25ex}{\accentmark{12}{\char"0951}}सामादीरो\raisebox{0.25ex}{\accentmark{15}{\char"1CD2}}ना\raisebox{0.25ex}{\accentmark{15}{\char"1CD2}}जा\raisebox{0.25ex}{\accentmark{15}{\char"1CD2}}गृविः\kern 0.05em ।\hskip 0.12em plus 0.04em minus 0.02em अच्छाको\raisebox{0.25ex}{\accentmark{15}{\char"1CD2}}श\raisebox{0.25ex}{\accentmark{12}{\char"0951}}म्माधुश्चू\raisebox{0.25ex}{\accentmark{15}{\char"1CD2}}त\raisebox{0.25ex}{\accentmark{12}{\char"0951}}म् \nolinebreak\hskip 0.08em \mbox{॥४॥}}} \par}

\vspace{0.8em}
                % --- SUBSECTIONS (Mantras) ---
                
                % Conditional rendering based on output_mode
                \par\filbreak

\phantomsection

\addcontentsline{toc}{subsection}{ऋक् ५}

{\centering \footnotesize \textcolor{AccentPurple}{॥ सप्तर्षयः (विश्वामित्रो गाथिनः) पवमानः सोमः बृहती ॥} \par}

\vspace{0.12em}

{\centering \textbf{\textcolor{AccentBlue}{सोमा\raisebox{0.25ex}{\accentmark{15}{\char"1CD2}}उष्वाण\raisebox{0.25ex}{\accentmark{15}{\char"1CD2}}स्सोतृ\raisebox{0.25ex}{\accentmark{15}{\char"1CD2}}भीः\raisebox{0.25ex}{\accentmark{15}{\char"1CD2}}\kern 0.05em ।\hskip 0.12em plus 0.04em minus 0.02em आधीष्णू\raisebox{0.25ex}{\accentmark{15}{\char"1CD2}}भीरावी\raisebox{0.25ex}{\accentmark{12}{\char"0951}}नां\kern 0.05em ।\hskip 0.12em plus 0.04em minus 0.02em अश्वा\raisebox{0.25ex}{\accentmark{15}{\char"1CD2}}ये\raisebox{0.25ex}{\accentmark{12}{\char"0951}}वाहारी\raisebox{0.25ex}{\accentmark{15}{\char"1CD2}}ता\raisebox{0.25ex}{\accentmark{12}{\char"0951}}यातीधा\raisebox{0.25ex}{\accentmark{15}{\char"1CD2}}रा\raisebox{0.25ex}{\accentmark{12}{\char"0951}}या\kern 0.05em ।\hskip 0.12em plus 0.04em minus 0.02em मन्द्रा\raisebox{0.25ex}{\accentmark{15}{\char"1CD2}}या\raisebox{0.25ex}{\accentmark{12}{\char"0951}}यातीधा\raisebox{0.25ex}{\accentmark{15}{\char"1CD2}}रा\raisebox{0.25ex}{\accentmark{12}{\char"0951}}या \nolinebreak\hskip 0.08em \mbox{॥५॥}}} \par}

\vspace{0.8em}
                % --- SUBSECTIONS (Mantras) ---
                
                % Conditional rendering based on output_mode
                \par\filbreak

\phantomsection

\addcontentsline{toc}{subsection}{ऋक् ६}

{\centering \footnotesize \textcolor{AccentPurple}{॥ सप्तर्षयः (जमदग्निर्भार्गवः) पवमानः सोमः बृहती ॥} \par}

\vspace{0.12em}

{\centering \textbf{\textcolor{AccentBlue}{तावाहंसो\raisebox{0.15ex}{\accentmark{12}{\char"1CF8}}मरारणा\kern 0.05em ।\hskip 0.12em plus 0.04em minus 0.02em सख्या\raisebox{0.25ex}{\accentmark{15}{\char"1CD2}}इ\raisebox{0.25ex}{\accentmark{12}{\char"0951}}न्दोदीवेदीवे\kern 0.05em ।\hskip 0.12em plus 0.04em minus 0.02em पूरू\raisebox{0.25ex}{\accentmark{15}{\char"1CD2}}णी\raisebox{0.25ex}{\accentmark{12}{\char"0951}}बभ्रोनीचा\raisebox{0.25ex}{\accentmark{15}{\char"1CD2}}र\raisebox{0.25ex}{\accentmark{12}{\char"0951}}न्तीमा\raisebox{0.25ex}{\accentmark{15}{\char"1CD2}}मा\raisebox{0.25ex}{\accentmark{15}{\char"1CD2}}वा\kern 0.05em ।\hskip 0.12em plus 0.04em minus 0.02em पारीधीं\raisebox{0.25ex}{\accentmark{15}{\char"1CD2}}रा\raisebox{0.25ex}{\accentmark{15}{\char"1CD2}}तीतं\raisebox{0.25ex}{\accentmark{15}{\char"1CD2}}ई\raisebox{0.25ex}{\accentmark{12}{\char"0951}}हि \nolinebreak\hskip 0.08em \mbox{॥६॥}}} \par}

\vspace{0.8em}
                % --- SUBSECTIONS (Mantras) ---
                
                % Conditional rendering based on output_mode
                \par\filbreak

\phantomsection

\addcontentsline{toc}{subsection}{ऋक् ७}

{\centering \footnotesize \textcolor{AccentPurple}{॥ सप्तर्षयः (वसिष्ठो मैत्रावरुणिः) पवमानः सोमः बृहती ॥} \par}

\vspace{0.12em}

{\centering \textbf{\textcolor{AccentBlue}{मृ\raisebox{0.25ex}{\accentmark{15}{\char"1CD2}}ज्या\raisebox{0.25ex}{\accentmark{15}{\char"1CD2}}मा\raisebox{0.25ex}{\accentmark{12}{\char"0951}}नस्सूहस्त्या\raisebox{0.25ex}{\accentmark{15}{\char"1CD2}}\kern 0.05em ।\hskip 0.12em plus 0.04em minus 0.02em समुद्रेवा\raisebox{0.25ex}{\accentmark{15}{\char"1CD2}}चा\raisebox{0.25ex}{\accentmark{12}{\char"0951}}मिन्वसि\kern 0.05em ।\hskip 0.12em plus 0.04em minus 0.02em रायिंपी\raisebox{0.25ex}{\accentmark{15}{\char"1CD2}}श\raisebox{0.25ex}{\accentmark{12}{\char"0951}}ङ्गंबहूलं\raisebox{0.15ex}{\accentmark{12}{\char"1CF9}}पू\raisebox{0.25ex}{\accentmark{12}{\char"0951}}रूस्पृ\raisebox{0.25ex}{\accentmark{12}{\char"0951}}हम्\kern 0.05em ।\hskip 0.12em plus 0.04em minus 0.02em पावामानाभ्यर्षसि \nolinebreak\hskip 0.08em \mbox{॥७॥}}} \par}

\vspace{0.8em}
                % --- SUBSECTIONS (Mantras) ---
                
                % Conditional rendering based on output_mode
                \par\filbreak

\phantomsection

\addcontentsline{toc}{subsection}{ऋक् ८}

{\centering \footnotesize \textcolor{AccentPurple}{॥ सप्तर्षयः पवमानः सोमः बृहती ॥} \par}

\vspace{0.12em}

{\centering \textbf{\textcolor{AccentBlue}{अभीसो\raisebox{0.25ex}{\accentmark{15}{\char"1CD2}}मा\raisebox{0.25ex}{\accentmark{12}{\char"0951}}साआया\raisebox{0.25ex}{\accentmark{15}{\char"1CD2}}वाः\raisebox{0.25ex}{\accentmark{12}{\char"0951}}\kern 0.05em ।\hskip 0.12em plus 0.04em minus 0.02em पा\raisebox{0.15ex}{\accentmark{12}{\char"1CF8}}वन्तेमा\raisebox{0.25ex}{\accentmark{15}{\char"1CD2}}द्यम्मादं\raisebox{0.25ex}{\accentmark{12}{\char"0951}}\kern 0.05em ।\hskip 0.12em plus 0.04em minus 0.02em समुद्रा\raisebox{0.15ex}{\accentmark{12}{\char"1CF9}}स्याधीविष्टापेमानीषीणाः\raisebox{0.25ex}{\accentmark{12}{\char"0951}}\kern 0.05em ।\hskip 0.12em plus 0.04em minus 0.02em मत्सा\raisebox{0.25ex}{\accentmark{15}{\char"1CD2}}रासोमादच्यू\raisebox{0.25ex}{\accentmark{15}{\char"1CD2}}ताः \nolinebreak\hskip 0.08em \mbox{॥८॥}}} \par}

\vspace{0.8em}
                % --- SUBSECTIONS (Mantras) ---
                
                % Conditional rendering based on output_mode
                \par\filbreak

\phantomsection

\addcontentsline{toc}{subsection}{ऋक् ९}

{\centering \footnotesize \textcolor{AccentPurple}{॥ सप्तर्षयः पवमानः सोमः बृहती ॥} \par}

\vspace{0.12em}

{\centering \textbf{\textcolor{AccentBlue}{पुनान\raisebox{0.15ex}{\accentmark{12}{\char"1CF9}}स्सोमाजाग्र\raisebox{0.25ex}{\accentmark{12}{\char"0951}}विः\kern 0.05em ।\hskip 0.12em plus 0.04em minus 0.02em अव्या\raisebox{0.25ex}{\accentmark{15}{\char"1CD2}}वारैः\raisebox{0.25ex}{\accentmark{15}{\char"1CD2}}पा\raisebox{0.15ex}{\accentmark{12}{\char"1CF9}}री\raisebox{0.25ex}{\accentmark{12}{\char"0951}}प्रीयः\kern 0.05em ।\hskip 0.12em plus 0.04em minus 0.02em त्वांवि\raisebox{0.25ex}{\accentmark{15}{\char"1CD2}}प्रो\raisebox{0.25ex}{\accentmark{12}{\char"0951}}अभवोअंगिरस्तमः\kern 0.05em ।\hskip 0.12em plus 0.04em minus 0.02em मध्वा\raisebox{0.25ex}{\accentmark{12}{\char"0951}}य\raisebox{0.25ex}{\accentmark{15}{\char"1CD2}}ज्ञं\raisebox{0.25ex}{\accentmark{15}{\char"1CD2}}मीमीक्षणः \nolinebreak\hskip 0.08em \mbox{॥९॥}}} \par}

\vspace{0.8em}
                % --- SUBSECTIONS (Mantras) ---
                
                % Conditional rendering based on output_mode
                \par\filbreak

\phantomsection

\addcontentsline{toc}{subsection}{ऋक् १०}

{\centering \footnotesize \textcolor{AccentPurple}{॥ सप्तर्षयः पवमानः सोमः बृहती ॥} \par}

\vspace{0.12em}

{\centering \textbf{\textcolor{AccentBlue}{पा\raisebox{0.25ex}{\accentmark{15}{\char"1CD2}}वामानाअसृक्षता\raisebox{0.25ex}{\accentmark{12}{\char"0951}}\kern 0.05em ।\hskip 0.12em plus 0.04em minus 0.02em पावी\raisebox{0.25ex}{\accentmark{15}{\char"1CD2}}त्रामाती\raisebox{0.25ex}{\accentmark{15}{\char"1CD2}}धा\raisebox{0.25ex}{\accentmark{15}{\char"1CD2}}रा\raisebox{0.25ex}{\accentmark{12}{\char"0951}}या\kern 0.05em ।\hskip 0.12em plus 0.04em minus 0.02em मारू\raisebox{0.25ex}{\accentmark{15}{\char"1CD2}}त्व\raisebox{0.25ex}{\accentmark{12}{\char"0951}}न्तोमत्सारा\raisebox{0.15ex}{\accentmark{12}{\char"1CF9}}इन्द्रीया\raisebox{0.25ex}{\accentmark{15}{\char"1CD2}}हायाः\raisebox{0.25ex}{\accentmark{15}{\char"1CD2}}\kern 0.05em ।\hskip 0.12em plus 0.04em minus 0.02em मेधा\raisebox{0.25ex}{\accentmark{15}{\char"1CD2}}माभि\raisebox{0.25ex}{\accentmark{15}{\char"1CD2}}प्रा\raisebox{0.25ex}{\accentmark{15}{\char"1CD2}}यां\raisebox{0.25ex}{\accentmark{12}{\char"0951}}सिच\raisebox{0.25ex}{\accentmark{15}{\char"1CD2}} \nolinebreak\hskip 0.08em \mbox{॥१०॥}}} \par}

\vspace{0.8em}
                % --- SUBSECTIONS (Mantras) ---
                
                % Conditional rendering based on output_mode
                \par\filbreak

\phantomsection

\addcontentsline{toc}{subsection}{ऋक् ११}

{\centering \textbf{\textcolor{AccentBlue}{॥ सप्तर्षयःपवमानःसोमःबृहती ॥ ॥ सप्तर्षयःपवमानःसोमःबृहती ॥ पा\raisebox{0.25ex}{\accentmark{15}{\char"1CD2}}वा\raisebox{0.25ex}{\accentmark{12}{\char"0951}}स्ववाजसा\raisebox{0.25ex}{\accentmark{15}{\char"1CD2}}ता\raisebox{0.25ex}{\accentmark{12}{\char"0951}}मः\kern 0.05em ।\hskip 0.12em plus 0.04em minus 0.02em आ\raisebox{0.25ex}{\accentmark{15}{\char"1CD2}}भी\raisebox{0.15ex}{\accentmark{12}{\char"1CF8}}वि\raisebox{0.15ex}{\accentmark{12}{\char"1CF9}}श्वा\raisebox{0.25ex}{\accentmark{12}{\char"0951}}नीवार्याः\kern 0.05em ।\hskip 0.12em plus 0.04em minus 0.02em त्वंसा\raisebox{0.25ex}{\accentmark{12}{\char"0951}}मुद्रःप्रा\raisebox{0.25ex}{\accentmark{12}{\char"0951}}थामेवि\raisebox{0.25ex}{\accentmark{15}{\char"1CD2}}धर्मन्\kern 0.05em ।\hskip 0.12em plus 0.04em minus 0.02em देवे\raisebox{0.25ex}{\accentmark{15}{\char"1CD2}}भ्यासोममात्सा\raisebox{0.25ex}{\accentmark{15}{\char"1CD2}}रः \nolinebreak\hskip 0.08em \mbox{॥११॥}}} \par}

\vspace{0.8em}
                % --- SUBSECTIONS (Mantras) ---
                
                % Conditional rendering based on output_mode
                \par\filbreak

\phantomsection

\addcontentsline{toc}{subsection}{ऋक् १२}

{\centering \textbf{\textcolor{AccentBlue}{इन्द्रा\raisebox{0.25ex}{\accentmark{12}{\char"0951}}यापावातेमा\raisebox{0.25ex}{\accentmark{15}{\char"1CD2}}दा\raisebox{0.25ex}{\accentmark{12}{\char"0951}}ः\kern 0.05em ।\hskip 0.12em plus 0.04em minus 0.02em सो\raisebox{0.15ex}{\accentmark{12}{\char"1CF8}}मो\raisebox{0.25ex}{\accentmark{12}{\char"0951}}मारू\raisebox{0.25ex}{\accentmark{15}{\char"1CD2}}त्वा\raisebox{0.25ex}{\accentmark{12}{\char"0951}}तेसूत\raisebox{0.25ex}{\accentmark{15}{\char"1CD2}}ः\kern 0.05em ।\hskip 0.12em plus 0.04em minus 0.02em साहा\raisebox{0.25ex}{\accentmark{15}{\char"1CD2}}स्रा\raisebox{0.25ex}{\accentmark{12}{\char"0951}}धारोअत्या\raisebox{0.25ex}{\accentmark{15}{\char"1CD2}}व्याम\raisebox{0.25ex}{\accentmark{15}{\char"1CD2}}र्ष\raisebox{0.25ex}{\accentmark{12}{\char"0951}}ति\kern 0.05em ।\hskip 0.12em plus 0.04em minus 0.02em तामि\raisebox{0.25ex}{\accentmark{12}{\char"0951}}मृजन्त्यायावाः \nolinebreak\hskip 0.08em \mbox{॥१२॥}}} \par}

\vspace{0.8em}

            % --- End of Kandah Closing Floral Line & Counts ---
            \par\nopagebreak\vspace{0.9em plus 0.3em minus 0.1em}
            \centerline{\textcolor{black!35}{\rule{1.8cm}{0.35pt}}\;\raisebox{-0.22ex}{\textcolor{vedicmaroon}{\large\SixFlowerPetalDotted}}\;\textcolor{black!35}{\rule{1.8cm}{0.35pt}}}
            \nopagebreak\vspace{0.45em}
            \centerline{\devafont\textbf{\textcolor{AccentBlue}{॥ ऋक् सङ्ख्या : १२ ॥}}}
        
            \clearpage

            % --- SECTION ---
            
            \phantomsection  

            \section*{\centering षष्ठ खण्डः}
            \addcontentsline{toc}{section}{षष्ठ खण्डः} 
            \markright{षष्ठ खण्डः}
            \nopagebreak % Keep section title with first mantra
            
            % --- SECTION-LEVEL PROCEDURE LINK ---

                % --- SUBSECTIONS (Mantras) ---
                
                % Conditional rendering based on output_mode
                \par\filbreak

\phantomsection

\addcontentsline{toc}{subsection}{ऋक् १}

{\centering \footnotesize \textcolor{AccentPurple}{॥ उशना काव्यः पवमानः सोमः त्रिष्टुप् ॥} \par}

\vspace{0.12em}

{\centering \textbf{\textcolor{AccentBlue}{प्र\raisebox{0.15ex}{\accentmark{12}{\char"1CF8}}तूद्रा\raisebox{0.25ex}{\accentmark{15}{\char"1CD2}}वापा\raisebox{0.25ex}{\accentmark{15}{\char"1CD2}}रीकोश\raisebox{0.25ex}{\accentmark{15}{\char"1CD2}}न्नीषी\raisebox{0.25ex}{\accentmark{12}{\char"0951}}दा\kern 0.05em ।\hskip 0.12em plus 0.04em minus 0.02em नृभी\raisebox{0.25ex}{\accentmark{12}{\char"0951}}ःपूनानोआभीवाज\raisebox{0.25ex}{\accentmark{12}{\char"0951}}मर्ष\raisebox{0.25ex}{\accentmark{12}{\char"0951}}\kern 0.05em ।\hskip 0.12em plus 0.04em minus 0.02em अश्वन्न\raisebox{0.15ex}{\accentmark{12}{\char"1CF9}}त्वा\raisebox{0.25ex}{\accentmark{12}{\char"0951}}वाजीन\raisebox{0.25ex}{\accentmark{12}{\char"0951}}र्म्म\raisebox{0.25ex}{\accentmark{12}{\char"0951}}र्जयन्ते\kern 0.05em ।\hskip 0.12em plus 0.04em minus 0.02em अच्छा\raisebox{0.25ex}{\accentmark{12}{\char"0951}}बर्ही\raisebox{0.15ex}{\accentmark{12}{\char"1CF9}}रा\raisebox{0.25ex}{\accentmark{12}{\char"0951}}शाना\raisebox{0.25ex}{\accentmark{15}{\char"1CD2}}भिर्न्न\raisebox{0.25ex}{\accentmark{12}{\char"0951}}यन्ते \nolinebreak\hskip 0.08em \mbox{॥१॥}}} \par}

\vspace{0.8em}
                % --- SUBSECTIONS (Mantras) ---
                
                % Conditional rendering based on output_mode
                \par\filbreak

\phantomsection

\addcontentsline{toc}{subsection}{ऋक् २}

{\centering \footnotesize \textcolor{AccentPurple}{॥ वृषगणो वासिष्ठः पवमानः सोमः त्रिष्टुप् ॥} \par}

\vspace{0.12em}

{\centering \textbf{\textcolor{AccentBlue}{प्राकाव्यामूशानेवब्रूवाणः\kern 0.05em ।\hskip 0.12em plus 0.04em minus 0.02em देवोदेवानाञ्जनीमाविवक्ति\kern 0.05em ।\hskip 0.12em plus 0.04em minus 0.02em माहीव्रातश्शूचीबन्धुःपावाकाः\kern 0.05em ।\hskip 0.12em plus 0.04em minus 0.02em पादावाराहोअभ्येतीरेभन् \nolinebreak\hskip 0.08em \mbox{॥२॥}}} \par}

\vspace{0.8em}
                % --- SUBSECTIONS (Mantras) ---
                
                % Conditional rendering based on output_mode
                \par\filbreak

\phantomsection

\addcontentsline{toc}{subsection}{ऋक् ३}

{\centering \footnotesize \textcolor{AccentPurple}{॥ पराशरः शाक्त्यः पवमानः सोमः त्रिष्टुप् ॥} \par}

\vspace{0.12em}

{\centering \textbf{\textcolor{AccentBlue}{तिस्रोवाचा\raisebox{0.25ex}{\accentmark{12}{\char"0951}}ईरयातीप्रा\raisebox{0.25ex}{\accentmark{15}{\char"1CD2}}वह्रीः\kern 0.05em ।\hskip 0.12em plus 0.04em minus 0.02em ऋतस्या\raisebox{0.25ex}{\accentmark{12}{\char"0951}}धीतिंब्रह्मा\raisebox{0.25ex}{\accentmark{12}{\char"0951}}णोमानीषाम्\kern 0.05em ।\hskip 0.12em plus 0.04em minus 0.02em गावो\raisebox{0.25ex}{\accentmark{12}{\char"0951}}यन्तीगो\raisebox{0.25ex}{\accentmark{15}{\char"1CD2}}पातिंपृच्छामानाः\kern 0.05em ।\hskip 0.12em plus 0.04em minus 0.02em सो\raisebox{0.25ex}{\accentmark{15}{\char"1CD2}}मंयन्तिमा\raisebox{0.25ex}{\accentmark{15}{\char"1CD2}}ता\raisebox{0.25ex}{\accentmark{15}{\char"1CD2}}यो\raisebox{0.25ex}{\accentmark{12}{\char"0951}}वावाशाना\raisebox{0.25ex}{\accentmark{15}{\char"1CD2}}ः \nolinebreak\hskip 0.08em \mbox{॥३॥}}} \par}

\vspace{0.8em}
                % --- SUBSECTIONS (Mantras) ---
                
                % Conditional rendering based on output_mode
                \par\filbreak

\phantomsection

\addcontentsline{toc}{subsection}{ऋक् ४}

{\centering \footnotesize \textcolor{AccentPurple}{॥ वसिष्ठो मैत्रावरुणिः पवमानः सोमः त्रिष्टुप् ॥} \par}

\vspace{0.12em}

{\centering \textbf{\textcolor{AccentBlue}{अस्य\raisebox{0.25ex}{\accentmark{15}{\char"1CD2}}प्रेषा\raisebox{0.25ex}{\accentmark{15}{\char"1CD2}}हेमा\raisebox{0.15ex}{\accentmark{12}{\char"1CF9}}ना\raisebox{0.25ex}{\accentmark{12}{\char"0951}}पूया\raisebox{0.25ex}{\accentmark{15}{\char"1CD2}}मा\raisebox{0.25ex}{\accentmark{12}{\char"0951}}नः\kern 0.05em ।\hskip 0.12em plus 0.04em minus 0.02em देवो\raisebox{0.25ex}{\accentmark{15}{\char"1CD2}}देवे\raisebox{0.25ex}{\accentmark{15}{\char"1CD2}}भिस्सा\raisebox{0.25ex}{\accentmark{15}{\char"1CD2}}मा\raisebox{0.25ex}{\accentmark{12}{\char"0951}}पृक्तारा\raisebox{0.25ex}{\accentmark{15}{\char"1CD2}}स\raisebox{0.25ex}{\accentmark{12}{\char"0951}}म्\kern 0.05em ।\hskip 0.12em plus 0.04em minus 0.02em सू\raisebox{0.25ex}{\accentmark{15}{\char"1CD2}}तःपा\raisebox{0.15ex}{\accentmark{12}{\char"1CF9}}वित्रंपर्येतीरे\raisebox{0.25ex}{\accentmark{15}{\char"1CD2}}भ\raisebox{0.25ex}{\accentmark{12}{\char"0951}}न्\kern 0.05em ।\hskip 0.12em plus 0.04em minus 0.02em मीतेवासा\raisebox{0.25ex}{\accentmark{15}{\char"1CD2}}द्मा\raisebox{0.25ex}{\accentmark{12}{\char"0951}}पाशूम\raisebox{0.25ex}{\accentmark{12}{\char"0951}}न्तीहो\raisebox{0.25ex}{\accentmark{15}{\char"1CD2}}ता\raisebox{0.25ex}{\accentmark{12}{\char"0951}} \nolinebreak\hskip 0.08em \mbox{॥४॥}}} \par}

\vspace{0.8em}
                % --- SUBSECTIONS (Mantras) ---
                
                % Conditional rendering based on output_mode
                \par\filbreak

\phantomsection

\addcontentsline{toc}{subsection}{ऋक् ५}

{\centering \footnotesize \textcolor{AccentPurple}{॥ पराशरः शाक्त्यः पवमानः सोमः त्रिष्टुप् ॥} \par}

\vspace{0.12em}

{\centering \textbf{\textcolor{AccentBlue}{अक्रा\raisebox{0.25ex}{\accentmark{15}{\char"1CD2}}न्सामुद्रः\raisebox{0.15ex}{\accentmark{12}{\char"1CF9}}प्रा\raisebox{0.25ex}{\accentmark{12}{\char"0951}}थामे\raisebox{0.25ex}{\accentmark{15}{\char"1CD2}}वि\raisebox{0.25ex}{\accentmark{15}{\char"1CD2}}ध\raisebox{0.25ex}{\accentmark{15}{\char"1CD2}}र्म्म\raisebox{0.25ex}{\accentmark{12}{\char"0951}}न्\kern 0.05em ।\hskip 0.12em plus 0.04em minus 0.02em जना\raisebox{0.15ex}{\accentmark{12}{\char"1CF9}}य\raisebox{0.25ex}{\accentmark{12}{\char"0951}}न्प्राजा\raisebox{0.25ex}{\accentmark{15}{\char"1CD2}}भु\raisebox{0.25ex}{\accentmark{15}{\char"1CD2}}वानास्यागोपाः\raisebox{0.25ex}{\accentmark{15}{\char"1CD2}}\kern 0.05em ।\hskip 0.12em plus 0.04em minus 0.02em वृषा\raisebox{0.25ex}{\accentmark{12}{\char"0951}}पावी\raisebox{0.25ex}{\accentmark{15}{\char"1CD2}}त्रेआ\raisebox{0.25ex}{\accentmark{15}{\char"1CD2}}धीसा\raisebox{0.25ex}{\accentmark{15}{\char"1CD2}}नोअव्ये\raisebox{0.25ex}{\accentmark{12}{\char"0951}}\kern 0.05em ।\hskip 0.12em plus 0.04em minus 0.02em बृ\raisebox{0.25ex}{\accentmark{15}{\char"1CD2}}ह\raisebox{0.25ex}{\accentmark{15}{\char"1CD2}}त्सो\raisebox{0.25ex}{\accentmark{15}{\char"1CD2}}मो\raisebox{0.25ex}{\accentmark{12}{\char"0951}}वावृधेस्वानो\raisebox{0.25ex}{\accentmark{15}{\char"1CD2}}अ\raisebox{0.25ex}{\accentmark{15}{\char"1CD2}}द्रीः \nolinebreak\hskip 0.08em \mbox{॥५॥}}} \par}

\vspace{0.8em}
                % --- SUBSECTIONS (Mantras) ---
                
                % Conditional rendering based on output_mode
                \par\filbreak

\phantomsection

\addcontentsline{toc}{subsection}{ऋक् ६}

{\centering \footnotesize \textcolor{AccentPurple}{॥ प्रस्कण्वः काण्वः पवमानः सोमः त्रिष्टुप् ॥} \par}

\vspace{0.12em}

{\centering \textbf{\textcolor{AccentBlue}{का\raisebox{0.25ex}{\accentmark{15}{\char"1CD2}}नी\raisebox{0.25ex}{\accentmark{12}{\char"0951}}क्रन्तीहा\raisebox{0.25ex}{\accentmark{15}{\char"1CD2}}री\raisebox{0.15ex}{\accentmark{12}{\char"1CF8}}रासृज्या\raisebox{0.25ex}{\accentmark{15}{\char"1CD2}}मा\raisebox{0.25ex}{\accentmark{12}{\char"0951}}नः\kern 0.05em ।\hskip 0.12em plus 0.04em minus 0.02em सी\raisebox{0.25ex}{\accentmark{15}{\char"1CD2}}दन्वा\raisebox{0.25ex}{\accentmark{15}{\char"1CD2}}ना\raisebox{0.25ex}{\accentmark{12}{\char"0951}}स्याजठा\raisebox{0.25ex}{\accentmark{12}{\char"0951}}रेपूनानः\kern 0.05em ।\hskip 0.12em plus 0.04em minus 0.02em नृभी\raisebox{0.15ex}{\accentmark{12}{\char"1CF9}}रीयाताःकृणुतेनिर्न्नी\raisebox{0.25ex}{\accentmark{12}{\char"0951}}जङ्गाम्\kern 0.05em ।\hskip 0.12em plus 0.04em minus 0.02em आ\raisebox{0.15ex}{\accentmark{12}{\char"1CF8}}तो\raisebox{0.25ex}{\accentmark{12}{\char"0951}}मातीञ्जनायातास्वाधा\raisebox{0.25ex}{\accentmark{15}{\char"1CD2}}भिः \nolinebreak\hskip 0.08em \mbox{॥६॥}}} \par}

\vspace{0.8em}
                % --- SUBSECTIONS (Mantras) ---
                
                % Conditional rendering based on output_mode
                \par\filbreak

\phantomsection

\addcontentsline{toc}{subsection}{ऋक् ७}

{\centering \footnotesize \textcolor{AccentPurple}{॥ उशना काव्यः पवमानः सोमः त्रिष्टुप् ॥} \par}

\vspace{0.12em}

{\centering \textbf{\textcolor{AccentBlue}{ए\raisebox{0.25ex}{\accentmark{15}{\char"1CD2}}षस्या\raisebox{0.25ex}{\accentmark{15}{\char"1CD2}}तेमा\raisebox{0.25ex}{\accentmark{15}{\char"1CD2}}धू\raisebox{0.25ex}{\accentmark{12}{\char"0951}}मंइन्द्रासोमाः\raisebox{0.25ex}{\accentmark{12}{\char"0951}}\kern 0.05em ।\hskip 0.12em plus 0.04em minus 0.02em वृ\raisebox{0.25ex}{\accentmark{15}{\char"1CD2}}षावृष्णःपा\raisebox{0.15ex}{\accentmark{12}{\char"1CF9}}री\raisebox{0.25ex}{\accentmark{12}{\char"0951}}पावीत्रे\raisebox{0.25ex}{\accentmark{12}{\char"0951}}क्षाः\kern 0.05em ।\hskip 0.12em plus 0.04em minus 0.02em साहास्रादश्शा\raisebox{0.25ex}{\accentmark{12}{\char"0951}}तादा\raisebox{0.15ex}{\accentmark{12}{\char"1CF9}}भु\raisebox{0.25ex}{\accentmark{15}{\char"1CD2}}रीदावा\raisebox{0.25ex}{\accentmark{12}{\char"0951}}\kern 0.05em ।\hskip 0.12em plus 0.04em minus 0.02em शश्वक्तामं\raisebox{0.25ex}{\accentmark{15}{\char"1CD2}}बर्ही\raisebox{0.25ex}{\accentmark{15}{\char"1CD2}}रावाज्यास्थात् \nolinebreak\hskip 0.08em \mbox{॥७॥}}} \par}

\vspace{0.8em}
                % --- SUBSECTIONS (Mantras) ---
                
                % Conditional rendering based on output_mode
                \par\filbreak

\phantomsection

\addcontentsline{toc}{subsection}{ऋक् ८}

{\centering \footnotesize \textcolor{AccentPurple}{॥ प्रतर्दनो दैवोदासिः पवमानः सोमः त्रिष्टुप् ॥} \par}

\vspace{0.12em}

{\centering \textbf{\textcolor{AccentBlue}{पा\raisebox{0.25ex}{\accentmark{15}{\char"1CD2}}वा\raisebox{0.25ex}{\accentmark{12}{\char"0951}}स्वसो\raisebox{0.25ex}{\accentmark{15}{\char"1CD2}}मामाधूमंऋता\raisebox{0.25ex}{\accentmark{15}{\char"1CD2}}वा\kern 0.05em ।\hskip 0.12em plus 0.04em minus 0.02em आपो\raisebox{0.15ex}{\accentmark{12}{\char"1CF9}}वा\raisebox{0.15ex}{\accentmark{12}{\char"1CF9}}सा\raisebox{0.25ex}{\accentmark{12}{\char"0951}}नोआ\raisebox{0.25ex}{\accentmark{15}{\char"1CD2}}धीसा\raisebox{0.25ex}{\accentmark{15}{\char"1CD2}}नोअ\raisebox{0.25ex}{\accentmark{15}{\char"1CD2}}व्ये\raisebox{0.25ex}{\accentmark{12}{\char"0951}}\kern 0.05em ।\hskip 0.12em plus 0.04em minus 0.02em आ\raisebox{0.25ex}{\accentmark{15}{\char"1CD2}}वद्रो\raisebox{0.25ex}{\accentmark{15}{\char"1CD2}}णानीघृता\raisebox{0.25ex}{\accentmark{15}{\char"1CD2}}व\raisebox{0.15ex}{\accentmark{12}{\char"1CF8}}न्तीरोहा\kern 0.05em ।\hskip 0.12em plus 0.04em minus 0.02em मादिन्ता\raisebox{0.15ex}{\accentmark{12}{\char"1CF8}}मोमत्सारा\raisebox{0.15ex}{\accentmark{12}{\char"1CF9}}इन्द्रा\raisebox{0.25ex}{\accentmark{12}{\char"0951}}पानाः \nolinebreak\hskip 0.08em \mbox{॥८॥}}} \par}

\vspace{0.8em}
                % --- SUBSECTIONS (Mantras) ---
                
                % Conditional rendering based on output_mode
                \par\filbreak

\phantomsection

\addcontentsline{toc}{subsection}{ऋक् ९}

{\centering \footnotesize \textcolor{AccentPurple}{॥ प्रतर्दनो दैवोदासिः पवमानः सोमः त्रिष्टुप् ॥} \par}

\vspace{0.12em}

{\centering \textbf{\textcolor{AccentBlue}{सो\raisebox{0.25ex}{\accentmark{15}{\char"1CD2}}माःपवतेज\raisebox{0.15ex}{\accentmark{12}{\char"1CF9}}नीता\raisebox{0.25ex}{\accentmark{15}{\char"1CD2}}मा\raisebox{0.15ex}{\accentmark{12}{\char"1CF9}}तीनां\kern 0.05em ।\hskip 0.12em plus 0.04em minus 0.02em जनीता\raisebox{0.25ex}{\accentmark{15}{\char"1CD2}}दीवोज\raisebox{0.25ex}{\accentmark{12}{\char"0951}}नीता\raisebox{0.15ex}{\accentmark{12}{\char"1CF9}}पृ\raisebox{0.25ex}{\accentmark{12}{\char"0951}}थिव्याः\raisebox{0.25ex}{\accentmark{12}{\char"0951}}\kern 0.05em ।\hskip 0.12em plus 0.04em minus 0.02em जनीता\raisebox{0.25ex}{\accentmark{12}{\char"0951}}ग्ने\raisebox{0.25ex}{\accentmark{12}{\char"0951}}र्जनीतासूर्या\raisebox{0.25ex}{\accentmark{12}{\char"0951}}स्या\kern 0.05em ।\hskip 0.12em plus 0.04em minus 0.02em जनीते\raisebox{0.25ex}{\accentmark{15}{\char"1CD2}}न्द्रा\raisebox{0.25ex}{\accentmark{12}{\char"0951}}स्याजनीतोतवी\raisebox{0.25ex}{\accentmark{15}{\char"1CD2}}ष्णोः\raisebox{0.25ex}{\accentmark{12}{\char"0951}} \nolinebreak\hskip 0.08em \mbox{॥९॥}}} \par}

\vspace{0.8em}
                % --- SUBSECTIONS (Mantras) ---
                
                % Conditional rendering based on output_mode
                \par\filbreak

\phantomsection

\addcontentsline{toc}{subsection}{ऋक् १०}

{\centering \footnotesize \textcolor{AccentPurple}{॥ वसिष्ठो मैत्रावरुणिः पवमानः सोमः त्रिष्टुप् ॥} \par}

\vspace{0.12em}

{\centering \textbf{\textcolor{AccentBlue}{आ\raisebox{0.25ex}{\accentmark{15}{\char"1CD2}}भित्रीपृष्ठंवृ\raisebox{0.25ex}{\accentmark{15}{\char"1CD2}}षाणं\raisebox{0.25ex}{\accentmark{12}{\char"0951}}वायोधाम्\kern 0.05em ।\hskip 0.12em plus 0.04em minus 0.02em अ\raisebox{0.15ex}{\accentmark{12}{\char"1CF8}}ङ्गोषीणा\raisebox{0.15ex}{\accentmark{12}{\char"1CF9}}मवावाशन्तावाणीः\kern 0.05em ।\hskip 0.12em plus 0.04em minus 0.02em वानावसा\raisebox{0.25ex}{\accentmark{12}{\char"0951}}नोवा\raisebox{0.15ex}{\accentmark{12}{\char"1CF9}}रू\raisebox{0.25ex}{\accentmark{12}{\char"0951}}णोनसी\raisebox{0.25ex}{\accentmark{15}{\char"1CD2}}न्धूः\kern 0.05em ।\hskip 0.12em plus 0.04em minus 0.02em वीरात्ना\raisebox{0.25ex}{\accentmark{15}{\char"1CD2}}धा\raisebox{0.25ex}{\accentmark{15}{\char"1CD2}}द\raisebox{0.25ex}{\accentmark{15}{\char"1CD2}}यातेवा\raisebox{0.25ex}{\accentmark{15}{\char"1CD2}}र्याणि \nolinebreak\hskip 0.08em \mbox{॥१०॥}}} \par}

\vspace{0.8em}

            % --- End of Kandah Closing Floral Line & Counts ---
            \par\nopagebreak\vspace{0.9em plus 0.3em minus 0.1em}
            \centerline{\textcolor{black!35}{\rule{1.8cm}{0.35pt}}\;\raisebox{-0.22ex}{\textcolor{vedicmaroon}{\large\SixFlowerPetalDotted}}\;\textcolor{black!35}{\rule{1.8cm}{0.35pt}}}
            \nopagebreak\vspace{0.45em}
            \centerline{\devafont\textbf{\textcolor{AccentBlue}{॥ ऋक् सङ्ख्या : १० ॥}}}
        
            \clearpage

            % --- SECTION ---
            
            \phantomsection  

            \section*{\centering सप्तम खण्डः}
            \addcontentsline{toc}{section}{सप्तम खण्डः} 
            \markright{सप्तम खण्डः}
            \nopagebreak % Keep section title with first mantra
            
            % --- SECTION-LEVEL PROCEDURE LINK ---

                % --- SUBSECTIONS (Mantras) ---
                
                % Conditional rendering based on output_mode
                \par\filbreak

\phantomsection

\addcontentsline{toc}{subsection}{ऋक् १}

{\centering \footnotesize \textcolor{AccentPurple}{॥ प्रतर्दनो दैवोदासिः पवमानः सोमः त्रिष्टुप् ॥} \par}

\vspace{0.12em}

{\centering \textbf{\textcolor{AccentBlue}{प्रा\raisebox{0.15ex}{\accentmark{12}{\char"1CF8}}से\raisebox{0.25ex}{\accentmark{12}{\char"0951}}नानी\raisebox{0.25ex}{\accentmark{15}{\char"1CD2}}श्शू\raisebox{0.25ex}{\accentmark{15}{\char"1CD2}}रोअ\raisebox{0.25ex}{\accentmark{15}{\char"1CD2}}ग्रे\raisebox{0.25ex}{\accentmark{15}{\char"1CD2}}राथा\raisebox{0.25ex}{\accentmark{12}{\char"0951}}नाम्\kern 0.05em ।\hskip 0.12em plus 0.04em minus 0.02em गव्य\raisebox{0.15ex}{\accentmark{12}{\char"1CF9}}न्ने\raisebox{0.25ex}{\accentmark{12}{\char"0951}}तीह\raisebox{0.25ex}{\accentmark{15}{\char"1CD2}}र्षातेअस्या\raisebox{0.25ex}{\accentmark{15}{\char"1CD2}}से\raisebox{0.25ex}{\accentmark{15}{\char"1CD2}}ना\raisebox{0.25ex}{\accentmark{12}{\char"0951}}भद्रान्कृण्वान्निन्द्राहावान्सा\raisebox{0.25ex}{\accentmark{15}{\char"1CD2}}खी\raisebox{0.25ex}{\accentmark{12}{\char"0951}}भ्यः\kern 0.05em ।\hskip 0.12em plus 0.04em minus 0.02em आसोमोवा\raisebox{0.25ex}{\accentmark{15}{\char"1CD2}}स्त्रा\raisebox{0.25ex}{\accentmark{12}{\char"0951}}राभासा\raisebox{0.25ex}{\accentmark{15}{\char"1CD2}}नी\raisebox{0.25ex}{\accentmark{12}{\char"0951}}दक्ते \nolinebreak\hskip 0.08em \mbox{॥१॥}}} \par}

\vspace{0.8em}
                % --- SUBSECTIONS (Mantras) ---
                
                % Conditional rendering based on output_mode
                \par\filbreak

\phantomsection

\addcontentsline{toc}{subsection}{ऋक् २}

{\centering \textbf{\textcolor{AccentBlue}{॥ पराशरःशाक्त्यःपवमानःसोमःत्रिष्टुप् ॥ ॥ इन्द्रप्रमतिर्वासिष्ठःपवमानःसोमःत्रिष्टुप् ॥ ॥ वसिष्ठोमैत्रावरुणिःपवमानःसोमःत्रिष्टुप् ॥ प्रातेधा\raisebox{0.25ex}{\accentmark{15}{\char"1CD2}}रामा\raisebox{0.25ex}{\accentmark{15}{\char"1CD2}}धू\raisebox{0.25ex}{\accentmark{12}{\char"0951}}मतिरसृगृन्\kern 0.05em ।\hskip 0.12em plus 0.04em minus 0.02em वा\raisebox{0.25ex}{\accentmark{15}{\char"1CD2}}रंय\raisebox{0.25ex}{\accentmark{15}{\char"1CD2}}त्पूतोअत्ये\raisebox{0.25ex}{\accentmark{15}{\char"1CD2}}ष\raisebox{0.25ex}{\accentmark{15}{\char"1CD2}}व्यम्\kern 0.05em ।\hskip 0.12em plus 0.04em minus 0.02em पा\raisebox{0.25ex}{\accentmark{15}{\char"1CD2}}वा\raisebox{0.25ex}{\accentmark{12}{\char"0951}}मानापा\raisebox{0.15ex}{\accentmark{12}{\char"1CF9}}वा\raisebox{0.25ex}{\accentmark{12}{\char"0951}}सेधा\raisebox{0.25ex}{\accentmark{15}{\char"1CD2}}मागोना\raisebox{0.25ex}{\accentmark{12}{\char"0951}}म्\kern 0.05em ।\hskip 0.12em plus 0.04em minus 0.02em ज\raisebox{0.25ex}{\accentmark{15}{\char"1CD2}}ना\raisebox{0.25ex}{\accentmark{15}{\char"1CD2}}यन्सू\raisebox{0.25ex}{\accentmark{15}{\char"1CD2}}र्या\raisebox{0.25ex}{\accentmark{12}{\char"0951}}मपिन्वोअर्कै\raisebox{0.25ex}{\accentmark{15}{\char"1CD2}}ः \nolinebreak\hskip 0.08em \mbox{॥२॥}}} \par}

\vspace{0.8em}
                % --- SUBSECTIONS (Mantras) ---
                
                % Conditional rendering based on output_mode
                \par\filbreak

\phantomsection

\addcontentsline{toc}{subsection}{ऋक् ३}

{\centering \textbf{\textcolor{AccentBlue}{प्रा\raisebox{0.25ex}{\accentmark{15}{\char"1CD2}}गा\raisebox{0.25ex}{\accentmark{12}{\char"0951}}यताभ्या\raisebox{0.25ex}{\accentmark{12}{\char"0951}}र्चामादेवा\raisebox{0.25ex}{\accentmark{15}{\char"1CD2}}न्\kern 0.05em ।\hskip 0.12em plus 0.04em minus 0.02em सो\raisebox{0.25ex}{\accentmark{15}{\char"1CD2}}मं\raisebox{0.25ex}{\accentmark{12}{\char"0951}}हिनोतामाहातेधाना॑\raisebox{0.25ex}{\accentmark{12}{\char"0951}}या\kern 0.05em ।\hskip 0.12em plus 0.04em minus 0.02em स्वादु\raisebox{0.25ex}{\accentmark{15}{\char"1CD2}}ःपा\raisebox{0.25ex}{\accentmark{12}{\char"0951}}वातामा\raisebox{0.25ex}{\accentmark{15}{\char"1CD2}}तीवा\raisebox{0.25ex}{\accentmark{15}{\char"1CD2}}रामा\raisebox{0.25ex}{\accentmark{15}{\char"1CD2}}व्यम्\kern 0.05em ।\hskip 0.12em plus 0.04em minus 0.02em आ\raisebox{0.25ex}{\accentmark{15}{\char"1CD2}}सी\raisebox{0.25ex}{\accentmark{12}{\char"0951}}दतुकालाश\raisebox{0.25ex}{\accentmark{12}{\char"0951}}न्देव\raisebox{0.25ex}{\accentmark{15}{\char"1CD2}}इ\raisebox{0.25ex}{\accentmark{15}{\char"1CD2}}न्दू\raisebox{0.25ex}{\accentmark{12}{\char"0951}}ः \nolinebreak\hskip 0.08em \mbox{॥३॥}}} \par}

\vspace{0.8em}
                % --- SUBSECTIONS (Mantras) ---
                
                % Conditional rendering based on output_mode
                \par\filbreak

\phantomsection

\addcontentsline{toc}{subsection}{ऋक् ४}

{\centering \textbf{\textcolor{AccentBlue}{प्राहि\raisebox{0.25ex}{\accentmark{15}{\char"1CD2}}न्वानो\raisebox{0.15ex}{\accentmark{12}{\char"1CF9}}ज\raisebox{0.25ex}{\accentmark{12}{\char"0951}}नीतारो\raisebox{0.25ex}{\accentmark{15}{\char"1CD2}}दस्योः\kern 0.05em ।\hskip 0.12em plus 0.04em minus 0.02em रा\raisebox{0.25ex}{\accentmark{15}{\char"1CD2}}थोनवा\raisebox{0.25ex}{\accentmark{15}{\char"1CD2}}जं\raisebox{0.25ex}{\accentmark{12}{\char"0951}}सानीष\raisebox{0.25ex}{\accentmark{15}{\char"1CD2}}न्ना\raisebox{0.25ex}{\accentmark{12}{\char"0951}}यासीत्\kern 0.05em ।\hskip 0.12em plus 0.04em minus 0.02em इ\raisebox{0.25ex}{\accentmark{15}{\char"1CD2}}न्द्रं\raisebox{0.25ex}{\accentmark{12}{\char"0951}}गच्छन्ना\raisebox{0.25ex}{\accentmark{15}{\char"1CD2}}यु\raisebox{0.25ex}{\accentmark{12}{\char"0951}}धासंशी\raisebox{0.25ex}{\accentmark{15}{\char"1CD2}}शा\raisebox{0.25ex}{\accentmark{12}{\char"0951}}नाः\kern 0.05em ।\hskip 0.12em plus 0.04em minus 0.02em वि\raisebox{0.25ex}{\accentmark{15}{\char"1CD2}}श्वा\raisebox{0.25ex}{\accentmark{12}{\char"0951}}वा\raisebox{0.25ex}{\accentmark{12}{\char"0951}}सूहा\raisebox{0.25ex}{\accentmark{15}{\char"1CD2}}स्ता\raisebox{0.25ex}{\accentmark{12}{\char"0951}}योरादधा\raisebox{0.25ex}{\accentmark{12}{\char"0951}}नः \nolinebreak\hskip 0.08em \mbox{॥४॥}}} \par}

\vspace{0.8em}
                % --- SUBSECTIONS (Mantras) ---
                
                % Conditional rendering based on output_mode
                \par\filbreak

\phantomsection

\addcontentsline{toc}{subsection}{ऋक् ५}

{\centering \footnotesize \textcolor{AccentPurple}{॥ कर्णश्रुद्वासिष्ठः पवमानः सोमः त्रिष्टुप् ॥} \par}

\vspace{0.12em}

{\centering \textbf{\textcolor{AccentBlue}{त\raisebox{0.25ex}{\accentmark{15}{\char"1CD2}}क्षद्या\raisebox{0.25ex}{\accentmark{15}{\char"1CD2}}दिमा\raisebox{0.25ex}{\accentmark{12}{\char"0951}}नासावेना\raisebox{0.25ex}{\accentmark{12}{\char"0951}}तोवा\raisebox{0.25ex}{\accentmark{15}{\char"1CD2}}क्\kern 0.05em ।\hskip 0.12em plus 0.04em minus 0.02em ज्येष्ठा\raisebox{0.25ex}{\accentmark{12}{\char"0951}}स्यधर्मन्द्यूक्षोरानी\raisebox{0.25ex}{\accentmark{12}{\char"0951}}के\kern 0.05em ।\hskip 0.12em plus 0.04em minus 0.02em आ\raisebox{0.15ex}{\accentmark{12}{\char"1CF8}}दी\raisebox{0.25ex}{\accentmark{12}{\char"0951}}मायन्वा\raisebox{0.25ex}{\accentmark{15}{\char"1CD2}}रामा\raisebox{0.25ex}{\accentmark{15}{\char"1CD2}}वावाशान\raisebox{0.25ex}{\accentmark{15}{\char"1CD2}}ः\kern 0.05em ।\hskip 0.12em plus 0.04em minus 0.02em जुष्टंपा\raisebox{0.15ex}{\accentmark{12}{\char"1CF9}}तिं\raisebox{0.25ex}{\accentmark{12}{\char"0951}}कालाशेगा\raisebox{0.25ex}{\accentmark{15}{\char"1CD2}}वइ\raisebox{0.25ex}{\accentmark{15}{\char"1CD2}}न्दुम् \nolinebreak\hskip 0.08em \mbox{॥५॥}}} \par}

\vspace{0.8em}
                % --- SUBSECTIONS (Mantras) ---
                
                % Conditional rendering based on output_mode
                \par\filbreak

\phantomsection

\addcontentsline{toc}{subsection}{ऋक् ६}

{\centering \textbf{\textcolor{AccentBlue}{॥ नोधागौतमःपवमानःसोमःत्रिष्टुप् ॥ ॥ मन्युर्वासिष्ठःपवमानःसोमःत्रिष्टुप् ॥ सा\raisebox{0.25ex}{\accentmark{15}{\char"1CD2}}कामु\raisebox{0.25ex}{\accentmark{15}{\char"1CD2}}क्षोमर्जयन्तस्वासा\raisebox{0.25ex}{\accentmark{12}{\char"0951}}रः\kern 0.05em ।\hskip 0.12em plus 0.04em minus 0.02em दशधीरा\raisebox{0.25ex}{\accentmark{12}{\char"0951}}स्याधीता\raisebox{0.25ex}{\accentmark{15}{\char"1CD2}}योधा\raisebox{0.25ex}{\accentmark{15}{\char"1CD2}}नूत्रीः\kern 0.05em ।\hskip 0.12em plus 0.04em minus 0.02em हा\raisebox{0.25ex}{\accentmark{15}{\char"1CD2}}रीःप\raisebox{0.25ex}{\accentmark{15}{\char"1CD2}}र्याद्रावार्जस्सू\raisebox{0.25ex}{\accentmark{15}{\char"1CD2}}र्या\raisebox{0.25ex}{\accentmark{12}{\char"0951}}स्या\kern 0.05em ।\hskip 0.12em plus 0.04em minus 0.02em द्रो\raisebox{0.25ex}{\accentmark{15}{\char"1CD2}}ण\raisebox{0.25ex}{\accentmark{15}{\char"1CD2}}न्ननक्षेअत्योना\raisebox{0.25ex}{\accentmark{15}{\char"1CD2}}वाजी \nolinebreak\hskip 0.08em \mbox{॥६॥}}} \par}

\vspace{0.8em}
                % --- SUBSECTIONS (Mantras) ---
                
                % Conditional rendering based on output_mode
                \par\filbreak

\phantomsection

\addcontentsline{toc}{subsection}{ऋक् ७}

{\centering \textbf{\textcolor{AccentBlue}{इ\raisebox{0.15ex}{\accentmark{12}{\char"1CF8}}न्दू\raisebox{0.25ex}{\accentmark{12}{\char"0951}}र्वाजी\raisebox{0.25ex}{\accentmark{15}{\char"1CD2}}पा\raisebox{0.25ex}{\accentmark{12}{\char"0951}}वातेगो\raisebox{0.25ex}{\accentmark{15}{\char"1CD2}}न्योघाः\kern 0.05em ।\hskip 0.12em plus 0.04em minus 0.02em इ\raisebox{0.25ex}{\accentmark{15}{\char"1CD2}}न्द्रेसो\raisebox{0.25ex}{\accentmark{15}{\char"1CD2}}मस्सा\raisebox{0.25ex}{\accentmark{15}{\char"1CD2}}हाइ\raisebox{0.25ex}{\accentmark{15}{\char"1CD2}}न्वन्मा\raisebox{0.25ex}{\accentmark{15}{\char"1CD2}}दा\raisebox{0.25ex}{\accentmark{12}{\char"0951}}य\kern 0.05em ।\hskip 0.12em plus 0.04em minus 0.02em ह\raisebox{0.25ex}{\accentmark{15}{\char"1CD2}}न्तीरक्षो\raisebox{0.25ex}{\accentmark{12}{\char"0951}}बा\raisebox{0.15ex}{\accentmark{12}{\char"1CF9}}धा\raisebox{0.25ex}{\accentmark{12}{\char"0951}}तेपरियारा\raisebox{0.25ex}{\accentmark{12}{\char"0951}}तिं\kern 0.05em ।\hskip 0.12em plus 0.04em minus 0.02em वा\raisebox{0.25ex}{\accentmark{15}{\char"1CD2}}री\raisebox{0.25ex}{\accentmark{12}{\char"0951}}वःकृण्व\raisebox{0.25ex}{\accentmark{15}{\char"1CD2}}न्वृ\raisebox{0.15ex}{\accentmark{12}{\char"1CF9}}ज\raisebox{0.15ex}{\accentmark{12}{\char"1CF9}}ना\raisebox{0.25ex}{\accentmark{12}{\char"0951}}स्या\raisebox{0.25ex}{\accentmark{15}{\char"1CD2}}रा\raisebox{0.25ex}{\accentmark{15}{\char"1CD2}}जा\raisebox{0.25ex}{\accentmark{12}{\char"0951}} \nolinebreak\hskip 0.08em \mbox{॥७॥}}} \par}

\vspace{0.8em}
                % --- SUBSECTIONS (Mantras) ---
                
                % Conditional rendering based on output_mode
                \par\filbreak

\phantomsection

\addcontentsline{toc}{subsection}{ऋक् ८}

{\centering \footnotesize \textcolor{AccentPurple}{॥ कण्वो घौरः पवमानः सोमः त्रिष्टुप् ॥} \par}

\vspace{0.12em}

{\centering \textbf{\textcolor{AccentBlue}{आ\raisebox{0.25ex}{\accentmark{15}{\char"1CD2}}धी\raisebox{0.25ex}{\accentmark{15}{\char"1CD2}}यादस्मिन्वा\raisebox{0.15ex}{\accentmark{12}{\char"1CF9}}जीनी\raisebox{0.25ex}{\accentmark{12}{\char"0951}}वाशु\raisebox{0.25ex}{\accentmark{15}{\char"1CD2}}भाः\kern 0.05em ।\hskip 0.12em plus 0.04em minus 0.02em स्प\raisebox{0.15ex}{\accentmark{12}{\char"1CF8}}र्ध\raisebox{0.25ex}{\accentmark{12}{\char"0951}}न्तेधी\raisebox{0.25ex}{\accentmark{15}{\char"1CD2}}यस्सू\raisebox{0.25ex}{\accentmark{15}{\char"1CD2}}रेनवी\raisebox{0.25ex}{\accentmark{15}{\char"1CD2}}शाः\kern 0.05em ।\hskip 0.12em plus 0.04em minus 0.02em आ\raisebox{0.25ex}{\accentmark{15}{\char"1CD2}}पो\raisebox{0.25ex}{\accentmark{12}{\char"0951}}वृणा\raisebox{0.25ex}{\accentmark{15}{\char"1CD2}}नःपा\raisebox{0.25ex}{\accentmark{12}{\char"0951}}वातेका\raisebox{0.25ex}{\accentmark{15}{\char"1CD2}}वी\raisebox{0.25ex}{\accentmark{12}{\char"0951}}यान्\kern 0.05em ।\hskip 0.12em plus 0.04em minus 0.02em व्राजन्ना\raisebox{0.15ex}{\accentmark{12}{\char"1CF9}}प\raisebox{0.25ex}{\accentmark{12}{\char"0951}}शुवा\raisebox{0.25ex}{\accentmark{15}{\char"1CD2}}र्थानायाम\raisebox{0.25ex}{\accentmark{15}{\char"1CD2}}न्मा\raisebox{0.25ex}{\accentmark{12}{\char"0951}} \nolinebreak\hskip 0.08em \mbox{॥८॥}}} \par}

\vspace{0.8em}
                % --- SUBSECTIONS (Mantras) ---
                
                % Conditional rendering based on output_mode
                \par\filbreak

\phantomsection

\addcontentsline{toc}{subsection}{ऋक् ९}

{\centering \footnotesize \textcolor{AccentPurple}{॥ पराशरः शाक्त्यः पवमानः सोमः त्रिष्टुप् ॥} \par}

\vspace{0.12em}

{\centering \textbf{\textcolor{AccentBlue}{महत्तत्सो\raisebox{0.25ex}{\accentmark{15}{\char"1CD2}}मोमाहीषा\raisebox{0.25ex}{\accentmark{15}{\char"1CD2}}श्चा\raisebox{0.25ex}{\accentmark{12}{\char"0951}}कारा\kern 0.05em ।\hskip 0.12em plus 0.04em minus 0.02em आ\raisebox{0.25ex}{\accentmark{15}{\char"1CD2}}पां\raisebox{0.25ex}{\accentmark{15}{\char"1CD2}}यत्ग\raisebox{0.25ex}{\accentmark{15}{\char"1CD2}}र्भो\raisebox{0.25ex}{\accentmark{15}{\char"1CD2}}वृ\raisebox{0.25ex}{\accentmark{12}{\char"0951}}णीतादेवा\raisebox{0.25ex}{\accentmark{15}{\char"1CD2}}न्\kern 0.05em ।\hskip 0.12em plus 0.04em minus 0.02em आद\raisebox{0.25ex}{\accentmark{12}{\char"0951}}धादिन्द्रेपा\raisebox{0.15ex}{\accentmark{12}{\char"1CF8}}वा\raisebox{0.25ex}{\accentmark{12}{\char"0951}}मानाओ\raisebox{0.25ex}{\accentmark{15}{\char"1CD2}}जा\raisebox{0.25ex}{\accentmark{12}{\char"0951}}ः\kern 0.05em ।\hskip 0.12em plus 0.04em minus 0.02em आज\raisebox{0.25ex}{\accentmark{12}{\char"0951}}नायत्सू\raisebox{0.25ex}{\accentmark{15}{\char"1CD2}}र्येज्यो\raisebox{0.25ex}{\accentmark{15}{\char"1CD2}}तीरि\raisebox{0.25ex}{\accentmark{15}{\char"1CD2}}न्दू\raisebox{0.25ex}{\accentmark{12}{\char"0951}}ः \nolinebreak\hskip 0.08em \mbox{॥९॥}}} \par}

\vspace{0.8em}
                % --- SUBSECTIONS (Mantras) ---
                
                % Conditional rendering based on output_mode
                \par\filbreak

\phantomsection

\addcontentsline{toc}{subsection}{ऋक् १०}

{\centering \footnotesize \textcolor{AccentPurple}{॥ प्रस्कण्वः काण्वः पवमानः सोमः त्रिष्टुप् ॥} \par}

\vspace{0.12em}

{\centering \textbf{\textcolor{AccentBlue}{अपा\raisebox{0.25ex}{\accentmark{15}{\char"1CD2}}मीवेदूर्मायास्त\raisebox{0.25ex}{\accentmark{15}{\char"1CD2}}क्तूरा\raisebox{0.25ex}{\accentmark{12}{\char"0951}}णः\kern 0.05em ।\hskip 0.12em plus 0.04em minus 0.02em प्रा\raisebox{0.15ex}{\accentmark{12}{\char"1CF8}}मा\raisebox{0.25ex}{\accentmark{12}{\char"0951}}नीषा\raisebox{0.25ex}{\accentmark{12}{\char"0951}}ई\raisebox{0.25ex}{\accentmark{12}{\char"0951}}रातेसोमामा\raisebox{0.25ex}{\accentmark{15}{\char"1CD2}}च्छा\raisebox{0.25ex}{\accentmark{12}{\char"0951}}\kern 0.05em ।\hskip 0.12em plus 0.04em minus 0.02em नमस्य\raisebox{0.25ex}{\accentmark{15}{\char"1CD2}}न्तीरू\raisebox{0.15ex}{\accentmark{12}{\char"1CF9}}पाचायन्तीसञ्चा\raisebox{0.25ex}{\accentmark{12}{\char"0951}}\kern 0.05em ।\hskip 0.12em plus 0.04em minus 0.02em आचा\raisebox{0.25ex}{\accentmark{12}{\char"0951}}विशन्त्यू१शाती\raisebox{0.25ex}{\accentmark{15}{\char"1CD2}}रूश\raisebox{0.25ex}{\accentmark{15}{\char"1CD2}}न्त\raisebox{0.25ex}{\accentmark{12}{\char"0951}}म् \nolinebreak\hskip 0.08em \mbox{॥१०॥}}} \par}

\vspace{0.8em}
                % --- SUBSECTIONS (Mantras) ---
                
                % Conditional rendering based on output_mode
                \par\filbreak

\phantomsection

\addcontentsline{toc}{subsection}{ऋक् ११}

{\centering \footnotesize \textcolor{AccentPurple}{॥ कुत्स आङ्गिरसः पवमानः सोमः त्रिष्टुप् ॥} \par}

\vspace{0.12em}

{\centering \textbf{\textcolor{AccentBlue}{आया\raisebox{0.25ex}{\accentmark{15}{\char"1CD2}}पा\raisebox{0.25ex}{\accentmark{12}{\char"0951}}वापावस्वैनावा\raisebox{0.25ex}{\accentmark{15}{\char"1CD2}}सू\raisebox{0.15ex}{\accentmark{12}{\char"1CF8}}नी\kern 0.05em ।\hskip 0.12em plus 0.04em minus 0.02em मां\raisebox{0.25ex}{\accentmark{15}{\char"1CD2}}श्चत्वाइन्दोसा\raisebox{0.15ex}{\accentmark{12}{\char"1CF9}}रा\raisebox{0.25ex}{\accentmark{12}{\char"0951}}साप्राध\raisebox{0.25ex}{\accentmark{12}{\char"0951}}न्वा\kern 0.05em ।\hskip 0.12em plus 0.04em minus 0.02em ब्रध्न\raisebox{0.25ex}{\accentmark{12}{\char"0951}}श्चिद्यस्यावा\raisebox{0.25ex}{\accentmark{15}{\char"1CD2}}तोनाजूतिं\raisebox{0.25ex}{\accentmark{15}{\char"1CD2}}\kern 0.05em ।\hskip 0.12em plus 0.04em minus 0.02em पुरू\raisebox{0.25ex}{\accentmark{15}{\char"1CD2}}मेधा\raisebox{0.25ex}{\accentmark{12}{\char"0951}}श्चित्ताका\raisebox{0.25ex}{\accentmark{12}{\char"0951}}वेना\raisebox{0.25ex}{\accentmark{15}{\char"1CD2}}रन्धात् \nolinebreak\hskip 0.08em \mbox{॥११॥}}} \par}

\vspace{0.8em}
                % --- SUBSECTIONS (Mantras) ---
                
                % Conditional rendering based on output_mode
                \par\filbreak

\phantomsection

\addcontentsline{toc}{subsection}{ऋक् १२}

{\centering \footnotesize \textcolor{AccentPurple}{॥ कश्यपो मारीचः पवमानः सोमः त्रिष्टुप् ॥} \par}

\vspace{0.12em}

{\centering \textbf{\textcolor{AccentBlue}{आ\raisebox{0.25ex}{\accentmark{15}{\char"1CD2}}स\raisebox{0.25ex}{\accentmark{12}{\char"0951}}र्जीवक्वारय्थेया\raisebox{0.25ex}{\accentmark{15}{\char"1CD2}}थाजौ\raisebox{0.25ex}{\accentmark{15}{\char"1CD2}}\kern 0.05em ।\hskip 0.12em plus 0.04em minus 0.02em धीया\raisebox{0.15ex}{\accentmark{12}{\char"1CF9}}मानो\raisebox{0.25ex}{\accentmark{12}{\char"0951}}क्ताप्राथामा\raisebox{0.15ex}{\accentmark{12}{\char"1CF9}}मा\raisebox{0.25ex}{\accentmark{12}{\char"0951}}नीषा\kern 0.05em ।\hskip 0.12em plus 0.04em minus 0.02em दशस्वासा\raisebox{0.25ex}{\accentmark{12}{\char"0951}}रोआधीसा\raisebox{0.25ex}{\accentmark{12}{\char"0951}}नोअव्ये\kern 0.05em ।\hskip 0.12em plus 0.04em minus 0.02em मृज\raisebox{0.25ex}{\accentmark{15}{\char"1CD2}}न्ती\raisebox{0.25ex}{\accentmark{15}{\char"1CD2}}वहिंसा\raisebox{0.15ex}{\accentmark{12}{\char"1CF9}}द\raisebox{0.25ex}{\accentmark{12}{\char"0951}}नेष्व\raisebox{0.25ex}{\accentmark{15}{\char"1CD2}}च्छा\raisebox{0.25ex}{\accentmark{12}{\char"0951}} \nolinebreak\hskip 0.08em \mbox{॥१२॥}}} \par}

\vspace{0.8em}

            % --- End of Kandah Closing Floral Line & Counts ---
            \par\nopagebreak\vspace{0.9em plus 0.3em minus 0.1em}
            \centerline{\textcolor{black!35}{\rule{1.8cm}{0.35pt}}\;\raisebox{-0.22ex}{\textcolor{vedicmaroon}{\large\SixFlowerPetalDotted}}\;\textcolor{black!35}{\rule{1.8cm}{0.35pt}}}
            \nopagebreak\vspace{0.45em}
            \centerline{\devafont\textbf{\textcolor{AccentBlue}{॥ ऋक् सङ्ख्या : १२ ॥}}}
        
            \clearpage

            % --- SECTION ---
            
            \phantomsection  

            \section*{\centering अष्टम खण्डः}
            \addcontentsline{toc}{section}{अष्टम खण्डः} 
            \markright{अष्टम खण्डः}
            \nopagebreak % Keep section title with first mantra
            
            % --- SECTION-LEVEL PROCEDURE LINK ---

                % --- SUBSECTIONS (Mantras) ---
                
                % Conditional rendering based on output_mode
                \par\filbreak

\phantomsection

\addcontentsline{toc}{subsection}{ऋक् १}

{\centering \footnotesize \textcolor{AccentPurple}{॥ अन्भीगुः श्यावाश्विः पवमानः सोमः अनुष्टुप् ॥} \par}

\vspace{0.12em}

{\centering \textbf{\textcolor{AccentBlue}{पूरो\raisebox{0.25ex}{\accentmark{15}{\char"1CD2}}जी\raisebox{0.25ex}{\accentmark{12}{\char"0951}}तीवोअ\raisebox{0.25ex}{\accentmark{15}{\char"1CD2}}न्धासः\kern 0.05em ।\hskip 0.12em plus 0.04em minus 0.02em सूता\raisebox{0.25ex}{\accentmark{15}{\char"1CD2}}या\raisebox{0.25ex}{\accentmark{12}{\char"0951}}मादयित्ना\raisebox{0.25ex}{\accentmark{15}{\char"1CD2}}वे\kern 0.05em ।\hskip 0.12em plus 0.04em minus 0.02em आपश्वा\raisebox{0.25ex}{\accentmark{15}{\char"1CD2}}नंश्नथिष्टना\kern 0.05em ।\hskip 0.12em plus 0.04em minus 0.02em साखा\raisebox{0.25ex}{\accentmark{12}{\char"0951}}योदीर्घाजिठ\raisebox{0.25ex}{\accentmark{12}{\char"0951}}ह्व्यम् \nolinebreak\hskip 0.08em \mbox{॥१॥}}} \par}

\vspace{0.8em}
                % --- SUBSECTIONS (Mantras) ---
                
                % Conditional rendering based on output_mode
                \par\filbreak

\phantomsection

\addcontentsline{toc}{subsection}{ऋक् २}

{\centering \footnotesize \textcolor{AccentPurple}{॥ नहुषो मानवः पवमानः सोमः अनुष्टुप् ॥} \par}

\vspace{0.12em}

{\centering \textbf{\textcolor{AccentBlue}{आयंपूषा\raisebox{0.25ex}{\accentmark{15}{\char"1CD2}}रायि\raisebox{0.25ex}{\accentmark{15}{\char"1CD2}}र्भागा\raisebox{0.25ex}{\accentmark{15}{\char"1CD2}}ः\kern 0.05em ।\hskip 0.12em plus 0.04em minus 0.02em सोमा\raisebox{0.25ex}{\accentmark{12}{\char"0951}}ःपूनानोअ\raisebox{0.25ex}{\accentmark{12}{\char"0951}}र्ष\raisebox{0.25ex}{\accentmark{12}{\char"0951}}ति\kern 0.05em ।\hskip 0.12em plus 0.04em minus 0.02em पाति\raisebox{0.25ex}{\accentmark{12}{\char"0951}}र्वि\raisebox{0.15ex}{\accentmark{12}{\char"1CF9}}श्वा\raisebox{0.25ex}{\accentmark{12}{\char"0951}}स्याभू\raisebox{0.25ex}{\accentmark{15}{\char"1CD2}}म\raisebox{0.25ex}{\accentmark{12}{\char"0951}}नः\kern 0.05em ।\hskip 0.12em plus 0.04em minus 0.02em व्यख्यद्रोद\raisebox{0.25ex}{\accentmark{12}{\char"0951}}सीऊभे \nolinebreak\hskip 0.08em \mbox{॥२॥}}} \par}

\vspace{0.8em}
                % --- SUBSECTIONS (Mantras) ---
                
                % Conditional rendering based on output_mode
                \par\filbreak

\phantomsection

\addcontentsline{toc}{subsection}{ऋक् ३}

{\centering \footnotesize \textcolor{AccentPurple}{॥ रैभसूनू काश्यपौ पवमानः सोमः बृहती ॥} \par}

\vspace{0.12em}

{\centering \textbf{\textcolor{AccentBlue}{आह\raisebox{0.25ex}{\accentmark{12}{\char"0951}}रीयातायाधृ\raisebox{0.25ex}{\accentmark{15}{\char"1CD2}}ष्णावे\kern 0.05em ।\hskip 0.12em plus 0.04em minus 0.02em धा\raisebox{0.25ex}{\accentmark{15}{\char"1CD2}}नु\raisebox{0.25ex}{\accentmark{12}{\char"0951}}ष्टन्वन्तीपौंठ\raisebox{0.25ex}{\accentmark{12}{\char"0951}}स्यम्\kern 0.05em ।\hskip 0.12em plus 0.04em minus 0.02em शु\raisebox{0.25ex}{\accentmark{15}{\char"1CD2}}क्रावीय्यन्त्या\raisebox{0.25ex}{\accentmark{15}{\char"1CD2}}सू\raisebox{0.25ex}{\accentmark{15}{\char"1CD2}}रायानिर्णीजे\kern 0.05em ।\hskip 0.12em plus 0.04em minus 0.02em वी\raisebox{0.25ex}{\accentmark{12}{\char"0951}}पा\raisebox{0.25ex}{\accentmark{15}{\char"1CD2}}मग्रे\raisebox{0.25ex}{\accentmark{15}{\char"1CD2}}माहीयू\raisebox{0.25ex}{\accentmark{15}{\char"1CD2}}वा\raisebox{0.25ex}{\accentmark{12}{\char"0951}}ः \nolinebreak\hskip 0.08em \mbox{॥३॥}}} \par}

\vspace{0.8em}
                % --- SUBSECTIONS (Mantras) ---
                
                % Conditional rendering based on output_mode
                \par\filbreak

\phantomsection

\addcontentsline{toc}{subsection}{ऋक् ४}

{\centering \footnotesize \textcolor{AccentPurple}{॥ अम्बरीषो वार्षागिरः ऋजिष्वा भारद्वाजश्च पवमानः सोमः अनुष्टुप् ॥} \par}

\vspace{0.12em}

{\centering \textbf{\textcolor{AccentBlue}{पा\raisebox{0.25ex}{\accentmark{15}{\char"1CD2}}रित्यं\raisebox{0.15ex}{\accentmark{12}{\char"1CF9}}हा\raisebox{0.25ex}{\accentmark{12}{\char"0951}}रीयातंहारीम्\kern 0.05em ।\hskip 0.12em plus 0.04em minus 0.02em बभ्रं\raisebox{0.25ex}{\accentmark{15}{\char"1CD2}}पु\raisebox{0.25ex}{\accentmark{12}{\char"0951}}नन्तीवा\raisebox{0.25ex}{\accentmark{15}{\char"1CD2}}रे\raisebox{0.25ex}{\accentmark{12}{\char"0951}}ण\kern 0.05em ।\hskip 0.12em plus 0.04em minus 0.02em यो\raisebox{0.25ex}{\accentmark{15}{\char"1CD2}}देवान्वि\raisebox{0.25ex}{\accentmark{15}{\char"1CD2}}श्वंइत्पारी\raisebox{0.25ex}{\accentmark{12}{\char"0951}}\kern 0.05em ।\hskip 0.12em plus 0.04em minus 0.02em मा\raisebox{0.25ex}{\accentmark{15}{\char"1CD2}}दे\raisebox{0.25ex}{\accentmark{12}{\char"0951}}नासह\raisebox{0.25ex}{\accentmark{15}{\char"1CD2}}गच्छा\raisebox{0.25ex}{\accentmark{12}{\char"0951}}ति \nolinebreak\hskip 0.08em \mbox{॥४॥}}} \par}

\vspace{0.8em}
                % --- SUBSECTIONS (Mantras) ---
                
                % Conditional rendering based on output_mode
                \par\filbreak

\phantomsection

\addcontentsline{toc}{subsection}{ऋक् ५}

{\centering \footnotesize \textcolor{AccentPurple}{॥ ययातिर्नाहुषः पवमानः सोमः अनुष्टुप् ॥} \par}

\vspace{0.12em}

{\centering \textbf{\textcolor{AccentBlue}{सूतासोमा\raisebox{0.25ex}{\accentmark{15}{\char"1CD2}}धूमक्तामः\kern 0.05em ।\hskip 0.12em plus 0.04em minus 0.02em सोमाइ\raisebox{0.25ex}{\accentmark{15}{\char"1CD2}}न्द्रा\raisebox{0.25ex}{\accentmark{12}{\char"0951}}याम\raisebox{0.25ex}{\accentmark{15}{\char"1CD2}}न्दीनाः\kern 0.05em ।\hskip 0.12em plus 0.04em minus 0.02em पावी\raisebox{0.25ex}{\accentmark{15}{\char"1CD2}}त्रवन्तोअक्षरन्\kern 0.05em ।\hskip 0.12em plus 0.04em minus 0.02em देवा\raisebox{0.25ex}{\accentmark{15}{\char"1CD2}}न्ग\raisebox{0.25ex}{\accentmark{15}{\char"1CD2}}च्छन्तूवोमा\raisebox{0.25ex}{\accentmark{15}{\char"1CD2}}दा\raisebox{0.25ex}{\accentmark{12}{\char"0951}}ः \nolinebreak\hskip 0.08em \mbox{॥५॥}}} \par}

\vspace{0.8em}
                % --- SUBSECTIONS (Mantras) ---
                
                % Conditional rendering based on output_mode
                \par\filbreak

\phantomsection

\addcontentsline{toc}{subsection}{ऋक् ६}

{\centering \footnotesize \textcolor{AccentPurple}{॥ मनुः साम्वरणः पवमानः सोमः अनुष्टुप् ॥} \par}

\vspace{0.12em}

{\centering \textbf{\textcolor{AccentBlue}{सोमाःपवन्ताइन्दा\raisebox{0.25ex}{\accentmark{12}{\char"0951}}वः\kern 0.05em ।\hskip 0.12em plus 0.04em minus 0.02em अ\raisebox{0.25ex}{\accentmark{15}{\char"1CD2}}स्मा\raisebox{0.25ex}{\accentmark{15}{\char"1CD2}}भ्य\raisebox{0.25ex}{\accentmark{12}{\char"0951}}ङ्गातुवित्ता\raisebox{0.25ex}{\accentmark{12}{\char"0951}}मः\kern 0.05em ।\hskip 0.12em plus 0.04em minus 0.02em मित्रा\raisebox{0.25ex}{\accentmark{15}{\char"1CD2}}स्वाना\raisebox{0.15ex}{\accentmark{12}{\char"1CF9}}आ\raisebox{0.25ex}{\accentmark{12}{\char"0951}}रेपा\raisebox{0.25ex}{\accentmark{15}{\char"1CD2}}सा\raisebox{0.25ex}{\accentmark{12}{\char"0951}}ः\kern 0.05em ।\hskip 0.12em plus 0.04em minus 0.02em स्वाध्या\raisebox{0.25ex}{\accentmark{15}{\char"1CD2}}ठ\raisebox{0.25ex}{\accentmark{12}{\char"0951}}स्व\raisebox{0.25ex}{\accentmark{15}{\char"1CD2}}र्वीद\raisebox{0.25ex}{\accentmark{12}{\char"0951}}ः \nolinebreak\hskip 0.08em \mbox{॥६॥}}} \par}

\vspace{0.8em}
                % --- SUBSECTIONS (Mantras) ---
                
                % Conditional rendering based on output_mode
                \par\filbreak

\phantomsection

\addcontentsline{toc}{subsection}{ऋक् ७}

{\centering \footnotesize \textcolor{AccentPurple}{॥ अम्बरीषो वार्षागिरः ऋजिष्वा भारद्वाजश्च पवमानः सोमः अनुष्टुप् ॥} \par}

\vspace{0.12em}

{\centering \textbf{\textcolor{AccentBlue}{आ\raisebox{0.25ex}{\accentmark{15}{\char"1CD2}}भी\raisebox{0.25ex}{\accentmark{15}{\char"1CD2}}नोवाजसा\raisebox{0.25ex}{\accentmark{15}{\char"1CD2}}तामम्\kern 0.05em ।\hskip 0.12em plus 0.04em minus 0.02em रायीमऋ\raisebox{0.25ex}{\accentmark{12}{\char"0951}}षाशातस्पृह\raisebox{0.25ex}{\accentmark{12}{\char"0951}}म्\kern 0.05em ।\hskip 0.12em plus 0.04em minus 0.02em इ\raisebox{0.15ex}{\accentmark{12}{\char"1CF8}}न्दो\raisebox{0.25ex}{\accentmark{12}{\char"0951}}साहास्रा\raisebox{0.25ex}{\accentmark{12}{\char"0951}}भर्णा\raisebox{0.25ex}{\accentmark{12}{\char"0951}}सम्\kern 0.05em ।\hskip 0.12em plus 0.04em minus 0.02em तूविद्युम्नंवीभासा\raisebox{0.25ex}{\accentmark{15}{\char"1CD2}}ह\raisebox{0.25ex}{\accentmark{12}{\char"0951}}म् \nolinebreak\hskip 0.08em \mbox{॥७॥}}} \par}

\vspace{0.8em}
                % --- SUBSECTIONS (Mantras) ---
                
                % Conditional rendering based on output_mode
                \par\filbreak

\phantomsection

\addcontentsline{toc}{subsection}{ऋक् ८}

{\centering \footnotesize \textcolor{AccentPurple}{॥ रैभसूनू काश्यपौ पवमानः सोमः अनुष्टुप् ॥} \par}

\vspace{0.12em}

{\centering \textbf{\textcolor{AccentBlue}{आभी\raisebox{0.25ex}{\accentmark{15}{\char"1CD2}}ना\raisebox{0.25ex}{\accentmark{12}{\char"0951}}वन्तेअद्रू\raisebox{0.25ex}{\accentmark{15}{\char"1CD2}}हाः\kern 0.05em ।\hskip 0.12em plus 0.04em minus 0.02em प्रीय\raisebox{0.15ex}{\accentmark{12}{\char"1CF9}}मिन्द्रा\raisebox{0.25ex}{\accentmark{12}{\char"0951}}स्यका\raisebox{0.25ex}{\accentmark{15}{\char"1CD2}}म्यम्\raisebox{0.25ex}{\accentmark{12}{\char"0951}}\kern 0.05em ।\hskip 0.12em plus 0.04em minus 0.02em वत्सन्नपू\raisebox{0.25ex}{\accentmark{15}{\char"1CD2}}र्वाआ\raisebox{0.25ex}{\accentmark{15}{\char"1CD2}}यूनि\kern 0.05em ।\hskip 0.12em plus 0.04em minus 0.02em जातं\raisebox{0.25ex}{\accentmark{15}{\char"1CD2}}रिहन्तीमा\raisebox{0.25ex}{\accentmark{12}{\char"0951}}ताराः \nolinebreak\hskip 0.08em \mbox{॥८॥}}} \par}

\vspace{0.8em}
                % --- SUBSECTIONS (Mantras) ---
                
                % Conditional rendering based on output_mode
                \par\filbreak

\phantomsection

\addcontentsline{toc}{subsection}{ऋक् ९}

{\centering \footnotesize \textcolor{AccentPurple}{॥ प्रजापतिर्वैश्वामित्रो वाच्यो वा पवमानः सोमः अनुष्टुप् ॥} \par}

\vspace{0.12em}

{\centering \textbf{\textcolor{AccentBlue}{प्रासुन्वानाय\raisebox{0.25ex}{\accentmark{15}{\char"1CD2}}आ\raisebox{0.25ex}{\accentmark{15}{\char"1CD2}}न्धा\raisebox{0.25ex}{\accentmark{12}{\char"0951}}सः\kern 0.05em ।\hskip 0.12em plus 0.04em minus 0.02em मर्क्तोना\raisebox{0.15ex}{\accentmark{12}{\char"1CF9}}वा\raisebox{0.25ex}{\accentmark{12}{\char"0951}}ष्टात\raisebox{0.25ex}{\accentmark{15}{\char"1CD2}}द्वा\raisebox{0.25ex}{\accentmark{15}{\char"1CD2}}चा\raisebox{0.25ex}{\accentmark{12}{\char"0951}}ः\kern 0.05em ।\hskip 0.12em plus 0.04em minus 0.02em आपश्वा\raisebox{0.25ex}{\accentmark{15}{\char"1CD2}}ना\raisebox{0.25ex}{\accentmark{12}{\char"0951}}माराधा\raisebox{0.25ex}{\accentmark{15}{\char"1CD2}}सम्\raisebox{0.25ex}{\accentmark{12}{\char"0951}}\kern 0.05em ।\hskip 0.12em plus 0.04em minus 0.02em हाता\raisebox{0.25ex}{\accentmark{15}{\char"1CD2}}माख\raisebox{0.25ex}{\accentmark{15}{\char"1CD2}}न्नाभ्र\raisebox{0.25ex}{\accentmark{15}{\char"1CD2}}ङ्गव\raisebox{0.25ex}{\accentmark{12}{\char"0951}}ः \nolinebreak\hskip 0.08em \mbox{॥९॥}}} \par}

\vspace{0.8em}

            % --- End of Kandah Closing Floral Line & Counts ---
            \par\nopagebreak\vspace{0.9em plus 0.3em minus 0.1em}
            \centerline{\textcolor{black!35}{\rule{1.8cm}{0.35pt}}\;\raisebox{-0.22ex}{\textcolor{vedicmaroon}{\large\SixFlowerPetalDotted}}\;\textcolor{black!35}{\rule{1.8cm}{0.35pt}}}
            \nopagebreak\vspace{0.45em}
            \centerline{\devafont\textbf{\textcolor{AccentBlue}{॥ ऋक् सङ्ख्या : ९ ॥}}}
        
            \clearpage

            % --- SECTION ---
            
            \phantomsection  

            \section*{\centering नवम खण्डः}
            \addcontentsline{toc}{section}{नवम खण्डः} 
            \markright{नवम खण्डः}
            \nopagebreak % Keep section title with first mantra
            
            % --- SECTION-LEVEL PROCEDURE LINK ---

                % --- SUBSECTIONS (Mantras) ---
                
                % Conditional rendering based on output_mode
                \par\filbreak

\phantomsection

\addcontentsline{toc}{subsection}{ऋक् १}

{\centering \footnotesize \textcolor{AccentPurple}{॥ कविर्भार्गवः पवमानः सोमः जगती ॥} \par}

\vspace{0.12em}

{\centering \textbf{\textcolor{AccentBlue}{आ\raisebox{0.25ex}{\accentmark{15}{\char"1CD2}}भी\raisebox{0.25ex}{\accentmark{15}{\char"1CD2}}प्रीया\raisebox{0.25ex}{\accentmark{15}{\char"1CD2}}णी\raisebox{0.25ex}{\accentmark{12}{\char"0951}}पवातेचानो\raisebox{0.25ex}{\accentmark{12}{\char"0951}}हितः\kern 0.05em ।\hskip 0.12em plus 0.04em minus 0.02em ना\raisebox{0.25ex}{\accentmark{15}{\char"1CD2}}मा\raisebox{0.25ex}{\accentmark{12}{\char"0951}}नीयह्वो\raisebox{0.25ex}{\accentmark{15}{\char"1CD2}}आ\raisebox{0.25ex}{\accentmark{15}{\char"1CD2}}धीयेषूवा\raisebox{0.25ex}{\accentmark{15}{\char"1CD2}}र्धा\raisebox{0.25ex}{\accentmark{12}{\char"0951}}ते\kern 0.05em ।\hskip 0.12em plus 0.04em minus 0.02em आ\raisebox{0.25ex}{\accentmark{15}{\char"1CD2}}सू\raisebox{0.25ex}{\accentmark{15}{\char"1CD2}}र्यास्यबृहा\raisebox{0.25ex}{\accentmark{15}{\char"1CD2}}तो\raisebox{0.25ex}{\accentmark{15}{\char"1CD2}}बृहन्ना\raisebox{0.25ex}{\accentmark{15}{\char"1CD2}}धी\kern 0.05em ।\hskip 0.12em plus 0.04em minus 0.02em रा\raisebox{0.25ex}{\accentmark{15}{\char"1CD2}}थं\raisebox{0.25ex}{\accentmark{15}{\char"1CD2}}विष्व\raisebox{0.25ex}{\accentmark{15}{\char"1CD2}}ञ्चमऋहद्वीचाक्षाण\raisebox{0.25ex}{\accentmark{15}{\char"1CD2}}ः \nolinebreak\hskip 0.08em \mbox{॥१॥}}} \par}

\vspace{0.8em}
                % --- SUBSECTIONS (Mantras) ---
                
                % Conditional rendering based on output_mode
                \par\filbreak

\phantomsection

\addcontentsline{toc}{subsection}{ऋक् २}

{\centering \footnotesize \textcolor{AccentPurple}{॥ कविर्भार्गवः पवमानः सोमः जगती ॥} \par}

\vspace{0.12em}

{\centering \textbf{\textcolor{AccentBlue}{आचोद\raisebox{0.25ex}{\accentmark{15}{\char"1CD2}}सो\raisebox{0.25ex}{\accentmark{15}{\char"1CD2}}नोधन्वन्त्विन्दा\raisebox{0.25ex}{\accentmark{12}{\char"0951}}वः\kern 0.05em ।\hskip 0.12em plus 0.04em minus 0.02em प्रस्वाना\raisebox{0.15ex}{\accentmark{12}{\char"1CF9}}सो\raisebox{0.25ex}{\accentmark{12}{\char"0951}}बृहा\raisebox{0.25ex}{\accentmark{15}{\char"1CD2}}द्देवे\raisebox{0.25ex}{\accentmark{15}{\char"1CD2}}षूहा\raisebox{0.25ex}{\accentmark{15}{\char"1CD2}}रा\raisebox{0.25ex}{\accentmark{12}{\char"0951}}यः\kern 0.05em ।\hskip 0.12em plus 0.04em minus 0.02em वी\raisebox{0.25ex}{\accentmark{15}{\char"1CD2}}चीदश्नाना\raisebox{0.25ex}{\accentmark{15}{\char"1CD2}}ईषा\raisebox{0.25ex}{\accentmark{15}{\char"1CD2}}योआ\raisebox{0.25ex}{\accentmark{15}{\char"1CD2}}रा\raisebox{0.25ex}{\accentmark{12}{\char"0951}}तयः\kern 0.05em ।\hskip 0.12em plus 0.04em minus 0.02em अर्योना\raisebox{0.25ex}{\accentmark{12}{\char"0951}}स्स\raisebox{0.25ex}{\accentmark{15}{\char"1CD2}}न्तुसानीषन्तुनोधी\raisebox{0.25ex}{\accentmark{15}{\char"1CD2}}याः \nolinebreak\hskip 0.08em \mbox{॥२॥}}} \par}

\vspace{0.8em}
                % --- SUBSECTIONS (Mantras) ---
                
                % Conditional rendering based on output_mode
                \par\filbreak

\phantomsection

\addcontentsline{toc}{subsection}{ऋक् ३}

{\centering \footnotesize \textcolor{AccentPurple}{॥ कविर्भार्गवः पवमानः सोमः जगती ॥} \par}

\vspace{0.12em}

{\centering \textbf{\textcolor{AccentBlue}{ए\raisebox{0.25ex}{\accentmark{15}{\char"1CD2}}ष\raisebox{0.25ex}{\accentmark{15}{\char"1CD2}}प्र\raisebox{0.25ex}{\accentmark{15}{\char"1CD2}}कोशेमाधू\raisebox{0.25ex}{\accentmark{12}{\char"0951}}मंअचिक्रदत्\kern 0.05em ।\hskip 0.12em plus 0.04em minus 0.02em इन्द्रस्याव\raisebox{0.15ex}{\accentmark{12}{\char"1CF9}}ज्रो\raisebox{0.25ex}{\accentmark{12}{\char"0951}}वापू\raisebox{0.15ex}{\accentmark{12}{\char"1CF9}}षो\raisebox{0.25ex}{\accentmark{12}{\char"0951}}वापू\raisebox{0.25ex}{\accentmark{15}{\char"1CD2}}ष्टा\raisebox{0.25ex}{\accentmark{12}{\char"0951}}मः\kern 0.05em ।\hskip 0.12em plus 0.04em minus 0.02em आ\raisebox{0.25ex}{\accentmark{15}{\char"1CD2}}भीमृत\raisebox{0.15ex}{\accentmark{12}{\char"1CF9}}स्या\raisebox{0.25ex}{\accentmark{12}{\char"0951}}सूदूघा\raisebox{0.25ex}{\accentmark{12}{\char"0951}}घृतश्रूता\raisebox{0.25ex}{\accentmark{12}{\char"0951}}ः\kern 0.05em ।\hskip 0.12em plus 0.04em minus 0.02em वाश्रा\raisebox{0.25ex}{\accentmark{15}{\char"1CD2}}अ\raisebox{0.25ex}{\accentmark{15}{\char"1CD2}}ऋषन्तीपाया\raisebox{0.25ex}{\accentmark{12}{\char"0951}}साचाधे\raisebox{0.25ex}{\accentmark{15}{\char"1CD2}}ना\raisebox{0.25ex}{\accentmark{15}{\char"1CD2}}वा\raisebox{0.25ex}{\accentmark{12}{\char"0951}}ः \nolinebreak\hskip 0.08em \mbox{॥३॥}}} \par}

\vspace{0.8em}
                % --- SUBSECTIONS (Mantras) ---
                
                % Conditional rendering based on output_mode
                \par\filbreak

\phantomsection

\addcontentsline{toc}{subsection}{ऋक् ४}

{\centering \footnotesize \textcolor{AccentPurple}{॥ सिकता निवावरी पवमानः सोमः जगती ॥} \par}

\vspace{0.12em}

{\centering \textbf{\textcolor{AccentBlue}{प्रो\raisebox{0.25ex}{\accentmark{15}{\char"1CD2}}आ\raisebox{0.25ex}{\accentmark{12}{\char"0951}}यासिदिन्दूरिन्द्रा\raisebox{0.25ex}{\accentmark{12}{\char"0951}}स्यानिष्कृ\raisebox{0.25ex}{\accentmark{12}{\char"0951}}तम्\kern 0.05em ।\hskip 0.12em plus 0.04em minus 0.02em साखासख्यु\raisebox{0.25ex}{\accentmark{12}{\char"0951}}र्नाप्रा\raisebox{0.25ex}{\accentmark{15}{\char"1CD2}}मीनातीसङ्गीर\raisebox{0.25ex}{\accentmark{12}{\char"0951}}म्\kern 0.05em ।\hskip 0.12em plus 0.04em minus 0.02em म\raisebox{0.25ex}{\accentmark{15}{\char"1CD2}}र्या\raisebox{0.25ex}{\accentmark{12}{\char"0951}}यिवायुवातीभि\raisebox{0.25ex}{\accentmark{12}{\char"0951}}स्सामऋष\raisebox{0.25ex}{\accentmark{15}{\char"1CD2}}ति\kern 0.05em ।\hskip 0.12em plus 0.04em minus 0.02em सो\raisebox{0.15ex}{\accentmark{12}{\char"1CF8}}म\raisebox{0.25ex}{\accentmark{12}{\char"0951}}ःकाला\raisebox{0.15ex}{\accentmark{12}{\char"1CF9}}शेशाता\raisebox{0.25ex}{\accentmark{15}{\char"1CD2}}या\raisebox{0.25ex}{\accentmark{12}{\char"0951}}मानापाथा \nolinebreak\hskip 0.08em \mbox{॥४॥}}} \par}

\vspace{0.8em}
                % --- SUBSECTIONS (Mantras) ---
                
                % Conditional rendering based on output_mode
                \par\filbreak

\phantomsection

\addcontentsline{toc}{subsection}{ऋक् ५}

{\centering \footnotesize \textcolor{AccentPurple}{॥ कविर्भार्गवः पवमानः सोमः जगती ॥} \par}

\vspace{0.12em}

{\centering \textbf{\textcolor{AccentBlue}{धर्ता\raisebox{0.25ex}{\accentmark{15}{\char"1CD2}}दीवा\raisebox{0.25ex}{\accentmark{15}{\char"1CD2}}ःपा\raisebox{0.25ex}{\accentmark{12}{\char"0951}}वातेकृत्वि\raisebox{0.25ex}{\accentmark{15}{\char"1CD2}}योरासा\raisebox{0.25ex}{\accentmark{12}{\char"0951}}ः\kern 0.05em ।\hskip 0.12em plus 0.04em minus 0.02em द\raisebox{0.15ex}{\accentmark{12}{\char"1CF8}}क्षोदेवा\raisebox{0.25ex}{\accentmark{15}{\char"1CD2}}नामानुमा\raisebox{0.25ex}{\accentmark{15}{\char"1CD2}}द्यो\raisebox{0.25ex}{\accentmark{12}{\char"0951}}नृभिः\kern 0.05em ।\hskip 0.12em plus 0.04em minus 0.02em हा\raisebox{0.25ex}{\accentmark{15}{\char"1CD2}}री\raisebox{0.25ex}{\accentmark{12}{\char"0951}}सृजानोअ\raisebox{0.25ex}{\accentmark{15}{\char"1CD2}}त्योनसत्वा\raisebox{0.25ex}{\accentmark{12}{\char"0951}}भीः\kern 0.05em ।\hskip 0.12em plus 0.04em minus 0.02em वृ\raisebox{0.25ex}{\accentmark{12}{\char"0951}}थापाजांसि\raisebox{0.25ex}{\accentmark{12}{\char"0951}}कृणूषे\raisebox{0.15ex}{\accentmark{12}{\char"1CF9}}नादीष्वा \nolinebreak\hskip 0.08em \mbox{॥५॥}}} \par}

\vspace{0.8em}
                % --- SUBSECTIONS (Mantras) ---
                
                % Conditional rendering based on output_mode
                \par\filbreak

\phantomsection

\addcontentsline{toc}{subsection}{ऋक् ६}

{\centering \footnotesize \textcolor{AccentPurple}{॥ सिकता निवावरी पवमानः सोमः जगती ॥} \par}

\vspace{0.12em}

{\centering \textbf{\textcolor{AccentBlue}{वृ\raisebox{0.25ex}{\accentmark{15}{\char"1CD2}}षा\raisebox{0.25ex}{\accentmark{12}{\char"0951}}मातीनां\raisebox{0.25ex}{\accentmark{15}{\char"1CD2}}पा\raisebox{0.25ex}{\accentmark{12}{\char"0951}}वतेवीचाक्षाण\raisebox{0.25ex}{\accentmark{15}{\char"1CD2}}ः\kern 0.05em ।\hskip 0.12em plus 0.04em minus 0.02em सोमोअह्रां\raisebox{0.25ex}{\accentmark{12}{\char"0951}}प्रा\raisebox{0.25ex}{\accentmark{15}{\char"1CD2}}तारीतो\raisebox{0.15ex}{\accentmark{12}{\char"1CF9}}षा\raisebox{0.15ex}{\accentmark{12}{\char"1CF9}}सान्दीवा\raisebox{0.25ex}{\accentmark{15}{\char"1CD2}}ः\kern 0.05em ।\hskip 0.12em plus 0.04em minus 0.02em प्राणा\raisebox{0.25ex}{\accentmark{15}{\char"1CD2}}सिन्धू\raisebox{0.25ex}{\accentmark{12}{\char"0951}}नांकाला\raisebox{0.25ex}{\accentmark{15}{\char"1CD2}}शंअचि\raisebox{0.25ex}{\accentmark{15}{\char"1CD2}}क्रद\raisebox{0.25ex}{\accentmark{12}{\char"0951}}त्\kern 0.05em ।\hskip 0.12em plus 0.04em minus 0.02em इन्द्र\raisebox{0.25ex}{\accentmark{12}{\char"0951}}स्या\raisebox{0.25ex}{\accentmark{15}{\char"1CD2}}हा\raisebox{0.15ex}{\accentmark{12}{\char"1CF9}}र्द्या\raisebox{0.25ex}{\accentmark{12}{\char"0951}}वीशा\raisebox{0.15ex}{\accentmark{12}{\char"1CF9}}न्मा\raisebox{0.25ex}{\accentmark{12}{\char"0951}}नीषीभिः \nolinebreak\hskip 0.08em \mbox{॥६॥}}} \par}

\vspace{0.8em}
                % --- SUBSECTIONS (Mantras) ---
                
                % Conditional rendering based on output_mode
                \par\filbreak

\phantomsection

\addcontentsline{toc}{subsection}{ऋक् ७}

{\centering \footnotesize \textcolor{AccentPurple}{॥ वसुर्भारद्वाजः पवमानः सोमः जगती ॥} \par}

\vspace{0.12em}

{\centering \textbf{\textcolor{AccentBlue}{आ\raisebox{0.15ex}{\accentmark{12}{\char"1CF8}}सा\raisebox{0.25ex}{\accentmark{12}{\char"0951}}वीसो\raisebox{0.25ex}{\accentmark{15}{\char"1CD2}}मो\raisebox{0.25ex}{\accentmark{12}{\char"0951}}आरूषो\raisebox{0.25ex}{\accentmark{15}{\char"1CD2}}वृषाहा\raisebox{0.25ex}{\accentmark{15}{\char"1CD2}}रि\raisebox{0.25ex}{\accentmark{12}{\char"0951}}ः\kern 0.05em ।\hskip 0.12em plus 0.04em minus 0.02em रा\raisebox{0.25ex}{\accentmark{15}{\char"1CD2}}जे\raisebox{0.25ex}{\accentmark{12}{\char"0951}}वादस्मो\raisebox{0.25ex}{\accentmark{15}{\char"1CD2}}आभी\raisebox{0.25ex}{\accentmark{15}{\char"1CD2}}गा\raisebox{0.25ex}{\accentmark{15}{\char"1CD2}}आ\raisebox{0.25ex}{\accentmark{12}{\char"0951}}चिक्रदत्\kern 0.05em ।\hskip 0.12em plus 0.04em minus 0.02em पुनानो\raisebox{0.25ex}{\accentmark{15}{\char"1CD2}}वारामा\raisebox{0.15ex}{\accentmark{12}{\char"1CF9}}त्ये\raisebox{0.25ex}{\accentmark{12}{\char"0951}}ष्वव्या\raisebox{0.25ex}{\accentmark{12}{\char"0951}}म्\kern 0.05em ।\hskip 0.12em plus 0.04em minus 0.02em श्येनो\raisebox{0.15ex}{\accentmark{12}{\char"1CF9}}नयो\raisebox{0.15ex}{\accentmark{12}{\char"1CF9}}निङ्घृता\raisebox{0.15ex}{\accentmark{12}{\char"1CF9}}व\raisebox{0.25ex}{\accentmark{12}{\char"0951}}न्तामा\raisebox{0.25ex}{\accentmark{15}{\char"1CD2}}सा\raisebox{0.25ex}{\accentmark{12}{\char"0951}}दत् \nolinebreak\hskip 0.08em \mbox{॥७॥}}} \par}

\vspace{0.8em}
                % --- SUBSECTIONS (Mantras) ---
                
                % Conditional rendering based on output_mode
                \par\filbreak

\phantomsection

\addcontentsline{toc}{subsection}{ऋक् ८}

{\centering \textbf{\textcolor{AccentBlue}{॥ वत्सप्रिर्भालन्दःपवमानःसोमःजगती ॥ ॥ रेणुर्वैश्वामित्रःपवमानःसोमःजगती ॥ प्रा\raisebox{0.25ex}{\accentmark{15}{\char"1CD2}}देव\raisebox{0.25ex}{\accentmark{15}{\char"1CD2}}मा\raisebox{0.25ex}{\accentmark{15}{\char"1CD2}}च्छामाधूमन्ताइन्दा\raisebox{0.25ex}{\accentmark{12}{\char"0951}}वः\kern 0.05em ।\hskip 0.12em plus 0.04em minus 0.02em आसी\raisebox{0.25ex}{\accentmark{12}{\char"0951}}ष्यदन्तागावाआ\raisebox{0.25ex}{\accentmark{15}{\char"1CD2}}ना\raisebox{0.25ex}{\accentmark{15}{\char"1CD2}}धेनावाः\kern 0.05em ।\hskip 0.12em plus 0.04em minus 0.02em बर्ही\raisebox{0.25ex}{\accentmark{12}{\char"0951}}षादो\raisebox{0.25ex}{\accentmark{12}{\char"0951}}वाचाना\raisebox{0.15ex}{\accentmark{12}{\char"1CF9}}व\raisebox{0.25ex}{\accentmark{12}{\char"0951}}न्ताऊ\raisebox{0.25ex}{\accentmark{15}{\char"1CD2}}धा\raisebox{0.15ex}{\accentmark{12}{\char"1CF9}}भिः\kern 0.05em ।\hskip 0.12em plus 0.04em minus 0.02em परिस्रू\raisebox{0.15ex}{\accentmark{12}{\char"1CF9}}ता\raisebox{0.25ex}{\accentmark{12}{\char"0951}}मुस्रीया\raisebox{0.25ex}{\accentmark{12}{\char"0951}}निर्णीज\raisebox{0.25ex}{\accentmark{12}{\char"0951}}न्धिरे \nolinebreak\hskip 0.08em \mbox{॥८॥}}} \par}

\vspace{0.8em}
                % --- SUBSECTIONS (Mantras) ---
                
                % Conditional rendering based on output_mode
                \par\filbreak

\phantomsection

\addcontentsline{toc}{subsection}{ऋक् ९}

{\centering \textbf{\textcolor{AccentBlue}{त्रिरा\raisebox{0.25ex}{\accentmark{12}{\char"0951}}स्मैसप्ता\raisebox{0.25ex}{\accentmark{15}{\char"1CD2}}धेना\raisebox{0.25ex}{\accentmark{15}{\char"1CD2}}वो\raisebox{0.25ex}{\accentmark{12}{\char"0951}}दुदुह्रिरे\kern 0.05em ।\hskip 0.12em plus 0.04em minus 0.02em सत्या\raisebox{0.25ex}{\accentmark{15}{\char"1CD2}}माशीरंपारामे\raisebox{0.25ex}{\accentmark{15}{\char"1CD2}}व्यो\raisebox{0.25ex}{\accentmark{12}{\char"0951}}मनि\kern 0.05em ।\hskip 0.12em plus 0.04em minus 0.02em च\raisebox{0.25ex}{\accentmark{15}{\char"1CD2}}त्वारिन्यो\raisebox{0.25ex}{\accentmark{15}{\char"1CD2}}भुवा\raisebox{0.25ex}{\accentmark{12}{\char"0951}}नानीनिर्णीजे\kern 0.05em ।\hskip 0.12em plus 0.04em minus 0.02em चा\raisebox{0.25ex}{\accentmark{15}{\char"1CD2}}रूणिचक्रेया\raisebox{0.25ex}{\accentmark{15}{\char"1CD2}}दृतैरा\raisebox{0.25ex}{\accentmark{12}{\char"0951}}वार्धत \nolinebreak\hskip 0.08em \mbox{॥९॥}}} \par}

\vspace{0.8em}
                % --- SUBSECTIONS (Mantras) ---
                
                % Conditional rendering based on output_mode
                \par\filbreak

\phantomsection

\addcontentsline{toc}{subsection}{ऋक् १०}

{\centering \footnotesize \textcolor{AccentPurple}{॥ वेनो भार्गवः पवमानः सोमः जगती ॥} \par}

\vspace{0.12em}

{\centering \textbf{\textcolor{AccentBlue}{इन्द्रा\raisebox{0.25ex}{\accentmark{12}{\char"0951}}यसोमासू\raisebox{0.15ex}{\accentmark{12}{\char"1CF9}}षूतःपा\raisebox{0.25ex}{\accentmark{15}{\char"1CD2}}री\raisebox{0.25ex}{\accentmark{12}{\char"0951}}स्रवा\kern 0.05em ।\hskip 0.12em plus 0.04em minus 0.02em अपा\raisebox{0.25ex}{\accentmark{15}{\char"1CD2}}मी\raisebox{0.25ex}{\accentmark{12}{\char"0951}}वाभवातुराक्षा\raisebox{0.25ex}{\accentmark{12}{\char"0951}}सासाहा\kern 0.05em ।\hskip 0.12em plus 0.04em minus 0.02em माते\raisebox{0.25ex}{\accentmark{15}{\char"1CD2}}रासा\raisebox{0.15ex}{\accentmark{12}{\char"1CF8}}स्यमत्सतद्वा\raisebox{0.25ex}{\accentmark{12}{\char"0951}}यावीना\raisebox{0.25ex}{\accentmark{12}{\char"0951}}ः\kern 0.05em ।\hskip 0.12em plus 0.04em minus 0.02em द्रावी\raisebox{0.25ex}{\accentmark{12}{\char"0951}}णास्वन्ता\raisebox{0.25ex}{\accentmark{15}{\char"1CD2}}ईहा\raisebox{0.25ex}{\accentmark{15}{\char"1CD2}}सन्त्वि\raisebox{0.25ex}{\accentmark{15}{\char"1CD2}}न्दा\raisebox{0.25ex}{\accentmark{12}{\char"0951}}वः \nolinebreak\hskip 0.08em \mbox{॥१०॥}}} \par}

\vspace{0.8em}
                % --- SUBSECTIONS (Mantras) ---
                
                % Conditional rendering based on output_mode
                \par\filbreak

\phantomsection

\addcontentsline{toc}{subsection}{ऋक् ११}

{\centering \footnotesize \textcolor{AccentPurple}{॥ गृत्समदः शौनकः पवमानः सोमः जगती ॥} \par}

\vspace{0.12em}

{\centering \textbf{\textcolor{AccentBlue}{अञ्ज\raisebox{0.25ex}{\accentmark{15}{\char"1CD2}}तेव्यठ\raisebox{0.25ex}{\accentmark{12}{\char"0951}}ञ्जतेसा\raisebox{0.25ex}{\accentmark{15}{\char"1CD2}}म\raisebox{0.25ex}{\accentmark{12}{\char"0951}}ञ्जते\kern 0.05em ।\hskip 0.12em plus 0.04em minus 0.02em क्रातुंरिहन्ती\raisebox{0.25ex}{\accentmark{12}{\char"0951}}मध्वाभ्या\raisebox{0.25ex}{\accentmark{12}{\char"0951}}ठ\raisebox{0.25ex}{\accentmark{12}{\char"0951}}ञ्जते\kern 0.05em ।\hskip 0.12em plus 0.04em minus 0.02em सि\raisebox{0.25ex}{\accentmark{15}{\char"1CD2}}न्धोरुच्छ्वासे\raisebox{0.25ex}{\accentmark{15}{\char"1CD2}}पाताय\raisebox{0.25ex}{\accentmark{12}{\char"0951}}न्तामुक्षाणम्\raisebox{0.25ex}{\accentmark{12}{\char"0951}}\kern 0.05em ।\hskip 0.12em plus 0.04em minus 0.02em हिरण्यपावाःपाशू\raisebox{0.25ex}{\accentmark{15}{\char"1CD2}}मप्सू\raisebox{0.25ex}{\accentmark{15}{\char"1CD2}}गृ\raisebox{0.15ex}{\accentmark{12}{\char"1CF8}}भ्णते \nolinebreak\hskip 0.08em \mbox{॥११॥}}} \par}

\vspace{0.8em}
                % --- SUBSECTIONS (Mantras) ---
                
                % Conditional rendering based on output_mode
                \par\filbreak

\phantomsection

\addcontentsline{toc}{subsection}{ऋक् १२}

{\centering \footnotesize \textcolor{AccentPurple}{॥ पवित्र आङ्गिरसः पवमानः सोमः जगती ॥} \par}

\vspace{0.12em}

{\centering \textbf{\textcolor{AccentBlue}{पावि\raisebox{0.15ex}{\accentmark{12}{\char"1CF9}}त्र\raisebox{0.15ex}{\accentmark{12}{\char"1CF9}}न्ते\raisebox{0.25ex}{\accentmark{12}{\char"0951}}वी\raisebox{0.25ex}{\accentmark{15}{\char"1CD2}}ता\raisebox{0.25ex}{\accentmark{12}{\char"0951}}तंब्रह्मणस्पते\kern 0.05em ।\hskip 0.12em plus 0.04em minus 0.02em प्राभु\raisebox{0.15ex}{\accentmark{12}{\char"1CF9}}र्गा\raisebox{0.15ex}{\accentmark{12}{\char"1CF9}}त्राणीप\raisebox{0.25ex}{\accentmark{15}{\char"1CD2}}र्येषिविश्वा\raisebox{0.25ex}{\accentmark{15}{\char"1CD2}}ताः\kern 0.05em ।\hskip 0.12em plus 0.04em minus 0.02em आ\raisebox{0.25ex}{\accentmark{15}{\char"1CD2}}ता\raisebox{0.25ex}{\accentmark{12}{\char"0951}}प्ततानुर्नाता\raisebox{0.25ex}{\accentmark{15}{\char"1CD2}}दामो\raisebox{0.25ex}{\accentmark{15}{\char"1CD2}}आश्नुते\kern 0.05em ।\hskip 0.12em plus 0.04em minus 0.02em श्रूता\raisebox{0.25ex}{\accentmark{15}{\char"1CD2}}साइ\raisebox{0.15ex}{\accentmark{12}{\char"1CF9}}द्व\raisebox{0.15ex}{\accentmark{12}{\char"1CF9}}हन्तासन्तादा\raisebox{0.25ex}{\accentmark{12}{\char"0951}}शता \nolinebreak\hskip 0.08em \mbox{॥१२॥}}} \par}

\vspace{0.8em}

            % --- End of Kandah Closing Floral Line & Counts ---
            \par\nopagebreak\vspace{0.9em plus 0.3em minus 0.1em}
            \centerline{\textcolor{black!35}{\rule{1.8cm}{0.35pt}}\;\raisebox{-0.22ex}{\textcolor{vedicmaroon}{\large\SixFlowerPetalDotted}}\;\textcolor{black!35}{\rule{1.8cm}{0.35pt}}}
            \nopagebreak\vspace{0.45em}
            \centerline{\devafont\textbf{\textcolor{AccentBlue}{॥ ऋक् सङ्ख्या : १२ ॥}}}
        
            \clearpage

            % --- SECTION ---
            
            \phantomsection  

            \section*{\centering दशम खण्डः}
            \addcontentsline{toc}{section}{दशम खण्डः} 
            \markright{दशम खण्डः}
            \nopagebreak % Keep section title with first mantra
            
            % --- SECTION-LEVEL PROCEDURE LINK ---

                % --- SUBSECTIONS (Mantras) ---
                
                % Conditional rendering based on output_mode
                \par\filbreak

\phantomsection

\addcontentsline{toc}{subsection}{ऋक् १}

{\centering \footnotesize \textcolor{AccentPurple}{॥ अग्निश्चाक्षुषः पवमानः सोमः उष्णिक् ॥} \par}

\vspace{0.12em}

{\centering \textbf{\textcolor{AccentBlue}{इन्द्राम\raisebox{0.15ex}{\accentmark{12}{\char"1CF9}}च्छा\raisebox{0.25ex}{\accentmark{12}{\char"0951}}सूता\raisebox{0.25ex}{\accentmark{15}{\char"1CD2}}ईमे\kern 0.05em ।\hskip 0.12em plus 0.04em minus 0.02em वृ\raisebox{0.25ex}{\accentmark{15}{\char"1CD2}}षा\raisebox{0.25ex}{\accentmark{12}{\char"0951}}णंयन्तूहारा\raisebox{0.25ex}{\accentmark{12}{\char"0951}}याः\kern 0.05em ।\hskip 0.12em plus 0.04em minus 0.02em श्रूष्ठेजातासाइन्दावाःस्वर्वीदः \nolinebreak\hskip 0.08em \mbox{॥१॥}}} \par}

\vspace{0.8em}
                % --- SUBSECTIONS (Mantras) ---
                
                % Conditional rendering based on output_mode
                \par\filbreak

\phantomsection

\addcontentsline{toc}{subsection}{ऋक् २}

{\centering \footnotesize \textcolor{AccentPurple}{॥ चक्षुर्मानवः पवमानः सोमः उष्णिक् ॥} \par}

\vspace{0.12em}

{\centering \textbf{\textcolor{AccentBlue}{प्रा\raisebox{0.25ex}{\accentmark{15}{\char"1CD2}}धा\raisebox{0.25ex}{\accentmark{12}{\char"0951}}न्वासोमा\raisebox{0.25ex}{\accentmark{15}{\char"1CD2}}जा\raisebox{0.25ex}{\accentmark{15}{\char"1CD2}}ग्र\raisebox{0.25ex}{\accentmark{12}{\char"0951}}विः\kern 0.05em ।\hskip 0.12em plus 0.04em minus 0.02em इन्द्रा\raisebox{0.25ex}{\accentmark{15}{\char"1CD2}}ये\raisebox{0.25ex}{\accentmark{15}{\char"1CD2}}न्दोपा\raisebox{0.25ex}{\accentmark{15}{\char"1CD2}}री\raisebox{0.25ex}{\accentmark{12}{\char"0951}}स्रव\kern 0.05em ।\hskip 0.12em plus 0.04em minus 0.02em धुम\raisebox{0.25ex}{\accentmark{12}{\char"0951}}न्तंशुष्मामाभा\raisebox{0.25ex}{\accentmark{12}{\char"0951}}रास्वर्वी\raisebox{0.25ex}{\accentmark{15}{\char"1CD2}}दम् \nolinebreak\hskip 0.08em \mbox{॥२॥}}} \par}

\vspace{0.8em}
                % --- SUBSECTIONS (Mantras) ---
                
                % Conditional rendering based on output_mode
                \par\filbreak

\phantomsection

\addcontentsline{toc}{subsection}{ऋक् ३}

{\centering \footnotesize \textcolor{AccentPurple}{॥ पर्वतनारदौ काण्वौ (शिखण्डिम्यावस्परसौ काश्यपौ वा) पवमानः सोमः उष्णिक् ॥} \par}

\vspace{0.12em}

{\centering \textbf{\textcolor{AccentBlue}{सा\raisebox{0.15ex}{\accentmark{12}{\char"1CF8}}खा\raisebox{0.25ex}{\accentmark{12}{\char"0951}}यआ\raisebox{0.25ex}{\accentmark{15}{\char"1CD2}}नी\raisebox{0.25ex}{\accentmark{12}{\char"0951}}षीदत\kern 0.05em ।\hskip 0.12em plus 0.04em minus 0.02em पु\raisebox{0.25ex}{\accentmark{15}{\char"1CD2}}ना\raisebox{0.25ex}{\accentmark{15}{\char"1CD2}}नायप्रा\raisebox{0.25ex}{\accentmark{15}{\char"1CD2}}गा\raisebox{0.25ex}{\accentmark{12}{\char"0951}}यत\kern 0.05em ।\hskip 0.12em plus 0.04em minus 0.02em शी\raisebox{0.25ex}{\accentmark{15}{\char"1CD2}}शून्नाय\raisebox{0.25ex}{\accentmark{15}{\char"1CD2}}ज्ञैःपा\raisebox{0.25ex}{\accentmark{15}{\char"1CD2}}री\raisebox{0.25ex}{\accentmark{12}{\char"0951}}भूषताश्रिये\raisebox{0.25ex}{\accentmark{15}{\char"1CD2}} \nolinebreak\hskip 0.08em \mbox{॥३॥}}} \par}

\vspace{0.8em}
                % --- SUBSECTIONS (Mantras) ---
                
                % Conditional rendering based on output_mode
                \par\filbreak

\phantomsection

\addcontentsline{toc}{subsection}{ऋक् ४}

{\centering \footnotesize \textcolor{AccentPurple}{॥ पर्वतनारदौ काण्वौ पवमानः सोमः उष्णिक् ॥} \par}

\vspace{0.12em}

{\centering \textbf{\textcolor{AccentBlue}{तंवास्सखा\raisebox{0.25ex}{\accentmark{12}{\char"0951}}योमादा\raisebox{0.25ex}{\accentmark{12}{\char"0951}}य\kern 0.05em ।\hskip 0.12em plus 0.04em minus 0.02em पु\raisebox{0.25ex}{\accentmark{15}{\char"1CD2}}ना\raisebox{0.25ex}{\accentmark{15}{\char"1CD2}}नामाभी\raisebox{0.25ex}{\accentmark{15}{\char"1CD2}}गा\raisebox{0.25ex}{\accentmark{12}{\char"0951}}यता\kern 0.05em ।\hskip 0.12em plus 0.04em minus 0.02em शीशून्ना\raisebox{0.25ex}{\accentmark{15}{\char"1CD2}}हव्यै\raisebox{0.25ex}{\accentmark{15}{\char"1CD2}}स्वादयन्तागू\raisebox{0.25ex}{\accentmark{12}{\char"0951}}र्तीभिः \nolinebreak\hskip 0.08em \mbox{॥४॥}}} \par}

\vspace{0.8em}
                % --- SUBSECTIONS (Mantras) ---
                
                % Conditional rendering based on output_mode
                \par\filbreak

\phantomsection

\addcontentsline{toc}{subsection}{ऋक् ५}

{\centering \footnotesize \textcolor{AccentPurple}{॥ त्रित आप्त्यः पवमानः सोमः उष्णिक् ॥} \par}

\vspace{0.12em}

{\centering \textbf{\textcolor{AccentBlue}{प्राणा\raisebox{0.15ex}{\accentmark{12}{\char"1CF9}}शी\raisebox{0.15ex}{\accentmark{12}{\char"1CF9}}शूर्माही\raisebox{0.25ex}{\accentmark{15}{\char"1CD2}}नाम्\raisebox{0.25ex}{\accentmark{12}{\char"0951}}\kern 0.05em ।\hskip 0.12em plus 0.04em minus 0.02em हिन्व\raisebox{0.25ex}{\accentmark{15}{\char"1CD2}}न्नृता\raisebox{0.25ex}{\accentmark{15}{\char"1CD2}}स्यादी\raisebox{0.25ex}{\accentmark{15}{\char"1CD2}}धी\raisebox{0.25ex}{\accentmark{12}{\char"0951}}तिम्\kern 0.05em ।\hskip 0.12em plus 0.04em minus 0.02em वि\raisebox{0.25ex}{\accentmark{15}{\char"1CD2}}श्वापा\raisebox{0.15ex}{\accentmark{12}{\char"1CF9}}रीप्रीया\raisebox{0.15ex}{\accentmark{12}{\char"1CF9}}भु\raisebox{0.25ex}{\accentmark{12}{\char"0951}}वाद\raisebox{0.15ex}{\accentmark{12}{\char"1CF9}}धाद्वीता \nolinebreak\hskip 0.08em \mbox{॥५॥}}} \par}

\vspace{0.8em}
                % --- SUBSECTIONS (Mantras) ---
                
                % Conditional rendering based on output_mode
                \par\filbreak

\phantomsection

\addcontentsline{toc}{subsection}{ऋक् ६}

{\centering \footnotesize \textcolor{AccentPurple}{॥ मनुराप्सवः पवमानः सोमः उष्णिक् ॥} \par}

\vspace{0.12em}

{\centering \textbf{\textcolor{AccentBlue}{पावा\raisebox{0.25ex}{\accentmark{12}{\char"0951}}स्वादेवा\raisebox{0.25ex}{\accentmark{15}{\char"1CD2}}वी\raisebox{0.25ex}{\accentmark{12}{\char"0951}}तये\kern 0.05em ।\hskip 0.12em plus 0.04em minus 0.02em इन्द्रोधाराभीरो\raisebox{0.25ex}{\accentmark{15}{\char"1CD2}}जसा\kern 0.05em ।\hskip 0.12em plus 0.04em minus 0.02em आ\raisebox{0.25ex}{\accentmark{15}{\char"1CD2}}काला\raisebox{0.25ex}{\accentmark{15}{\char"1CD2}}शम्मा\raisebox{0.25ex}{\accentmark{15}{\char"1CD2}}धू\raisebox{0.25ex}{\accentmark{12}{\char"0951}}मान्सोमनस्सद\raisebox{0.25ex}{\accentmark{15}{\char"1CD2}}ः \nolinebreak\hskip 0.08em \mbox{॥६॥}}} \par}

\vspace{0.8em}
                % --- SUBSECTIONS (Mantras) ---
                
                % Conditional rendering based on output_mode
                \par\filbreak

\phantomsection

\addcontentsline{toc}{subsection}{ऋक् ७}

{\centering \footnotesize \textcolor{AccentPurple}{॥ अग्निश्चाक्षुषः पवमानः सोमः उष्णिक् ॥} \par}

\vspace{0.12em}

{\centering \textbf{\textcolor{AccentBlue}{सो\raisebox{0.25ex}{\accentmark{15}{\char"1CD2}}माःपूनाना\raisebox{0.25ex}{\accentmark{15}{\char"1CD2}}ऊमींणा\raisebox{0.25ex}{\accentmark{12}{\char"0951}}\kern 0.05em ।\hskip 0.12em plus 0.04em minus 0.02em अ\raisebox{0.25ex}{\accentmark{15}{\char"1CD2}}व्यंवा\raisebox{0.25ex}{\accentmark{15}{\char"1CD2}}रंवी\raisebox{0.25ex}{\accentmark{12}{\char"0951}}धा\raisebox{0.25ex}{\accentmark{12}{\char"0951}}वती\kern 0.05em ।\hskip 0.12em plus 0.04em minus 0.02em अ\raisebox{0.15ex}{\accentmark{12}{\char"1CF8}}ग्रे\raisebox{0.25ex}{\accentmark{12}{\char"0951}}वाच\raisebox{0.25ex}{\accentmark{15}{\char"1CD2}}ःपावमानःकानी\raisebox{0.25ex}{\accentmark{12}{\char"0951}}क्रदत् \nolinebreak\hskip 0.08em \mbox{॥७॥}}} \par}

\vspace{0.8em}
                % --- SUBSECTIONS (Mantras) ---
                
                % Conditional rendering based on output_mode
                \par\filbreak

\phantomsection

\addcontentsline{toc}{subsection}{ऋक् ८}

{\centering \textbf{\textcolor{AccentBlue}{॥ द्वितआप्त्यःपवमानःसोमःउष्णिक् ॥ ॥ पर्वतनारदौकाण्वौपवमानःसोमःउष्णिक् ॥ ॥ पर्वतनारदौकाण्वौ(शिखण्डिम्यावस्परसौकाश्यपौवा)पवमानःसोमःउष्णिक् ॥ प्रापू\raisebox{0.25ex}{\accentmark{12}{\char"0951}}नाना\raisebox{0.15ex}{\accentmark{12}{\char"1CF9}}या\raisebox{0.25ex}{\accentmark{12}{\char"0951}}वेधा\raisebox{0.25ex}{\accentmark{15}{\char"1CD2}}से\raisebox{0.25ex}{\accentmark{12}{\char"0951}}\kern 0.05em ।\hskip 0.12em plus 0.04em minus 0.02em सो\raisebox{0.15ex}{\accentmark{12}{\char"1CF8}}मा\raisebox{0.25ex}{\accentmark{12}{\char"0951}}यावा\raisebox{0.25ex}{\accentmark{15}{\char"1CD2}}चा\raisebox{0.25ex}{\accentmark{12}{\char"0951}}उच्यते\kern 0.05em ।\hskip 0.12em plus 0.04em minus 0.02em भृ\raisebox{0.25ex}{\accentmark{12}{\char"0951}}ति\raisebox{0.25ex}{\accentmark{12}{\char"0951}}र्न्ना\raisebox{0.25ex}{\accentmark{15}{\char"1CD2}}भारा\raisebox{0.25ex}{\accentmark{15}{\char"1CD2}}मा\raisebox{0.15ex}{\accentmark{12}{\char"1CF9}}तीभि\raisebox{0.25ex}{\accentmark{12}{\char"0951}}र्जुजो\raisebox{0.25ex}{\accentmark{15}{\char"1CD2}}षा\raisebox{0.25ex}{\accentmark{12}{\char"0951}}ते \nolinebreak\hskip 0.08em \mbox{॥८॥}}} \par}

\vspace{0.8em}
                % --- SUBSECTIONS (Mantras) ---
                
                % Conditional rendering based on output_mode
                \par\filbreak

\phantomsection

\addcontentsline{toc}{subsection}{ऋक् ९}

{\centering \textbf{\textcolor{AccentBlue}{गोम\raisebox{0.25ex}{\accentmark{12}{\char"0951}}न्नइन्दोअ\raisebox{0.25ex}{\accentmark{15}{\char"1CD2}}श्वा\raisebox{0.25ex}{\accentmark{12}{\char"0951}}वात्\kern 0.05em ।\hskip 0.12em plus 0.04em minus 0.02em सूता\raisebox{0.25ex}{\accentmark{15}{\char"1CD2}}स्सू\raisebox{0.25ex}{\accentmark{12}{\char"0951}}दक्षधनिवः\kern 0.05em ।\hskip 0.12em plus 0.04em minus 0.02em शूचि\raisebox{0.25ex}{\accentmark{12}{\char"0951}}ञ्चाव\raisebox{0.25ex}{\accentmark{15}{\char"1CD2}}र्ण\raisebox{0.25ex}{\accentmark{12}{\char"0951}}माधीगोषू\raisebox{0.25ex}{\accentmark{12}{\char"0951}}धाराया \nolinebreak\hskip 0.08em \mbox{॥९॥}}} \par}

\vspace{0.8em}
                % --- SUBSECTIONS (Mantras) ---
                
                % Conditional rendering based on output_mode
                \par\filbreak

\phantomsection

\addcontentsline{toc}{subsection}{ऋक् १०}

{\centering \textbf{\textcolor{AccentBlue}{अस्मा\raisebox{0.25ex}{\accentmark{15}{\char"1CD2}}भ्यन्त्वावासूवी\raisebox{0.25ex}{\accentmark{15}{\char"1CD2}}दम्\kern 0.05em ।\hskip 0.12em plus 0.04em minus 0.02em आभीवाणी\raisebox{0.25ex}{\accentmark{12}{\char"0951}}रनूषत\kern 0.05em ।\hskip 0.12em plus 0.04em minus 0.02em गोभिष्टेव\raisebox{0.15ex}{\accentmark{12}{\char"1CF9}}र्णा\raisebox{0.25ex}{\accentmark{12}{\char"0951}}माभी\raisebox{0.25ex}{\accentmark{15}{\char"1CD2}}वा\raisebox{0.25ex}{\accentmark{12}{\char"0951}}सयामसि \nolinebreak\hskip 0.08em \mbox{॥१०॥}}} \par}

\vspace{0.8em}
                % --- SUBSECTIONS (Mantras) ---
                
                % Conditional rendering based on output_mode
                \par\filbreak

\phantomsection

\addcontentsline{toc}{subsection}{ऋक् ११}

{\centering \footnotesize \textcolor{AccentPurple}{॥ अग्निश्चाक्षुषः पवमानः सोमः उष्णिक् ॥} \par}

\vspace{0.12em}

{\centering \textbf{\textcolor{AccentBlue}{पावा\raisebox{0.25ex}{\accentmark{12}{\char"0951}}तेहरियातोहारिः\kern 0.05em ।\hskip 0.12em plus 0.04em minus 0.02em आतिह्वा\raisebox{0.15ex}{\accentmark{12}{\char"1CF9}}रां\raisebox{0.25ex}{\accentmark{12}{\char"0951}}सीरं\raisebox{0.25ex}{\accentmark{15}{\char"1CD2}}ह्या\kern 0.05em ।\hskip 0.12em plus 0.04em minus 0.02em अभ्यर्ष\raisebox{0.25ex}{\accentmark{15}{\char"1CD2}}स्तोतृ\raisebox{0.15ex}{\accentmark{12}{\char"1CF9}}भ्यो\raisebox{0.25ex}{\accentmark{12}{\char"0951}}वीरा\raisebox{0.25ex}{\accentmark{15}{\char"1CD2}}वद्या\raisebox{0.25ex}{\accentmark{15}{\char"1CD2}}शा\raisebox{0.25ex}{\accentmark{12}{\char"0951}}ः \nolinebreak\hskip 0.08em \mbox{॥११॥}}} \par}

\vspace{0.8em}
                % --- SUBSECTIONS (Mantras) ---
                
                % Conditional rendering based on output_mode
                \par\filbreak

\phantomsection

\addcontentsline{toc}{subsection}{ऋक् १२}

{\centering \footnotesize \textcolor{AccentPurple}{॥ द्वित आप्त्यः पवमानः सोमः उष्णिक् ॥} \par}

\vspace{0.12em}

{\centering \textbf{\textcolor{AccentBlue}{पारीको\raisebox{0.25ex}{\accentmark{15}{\char"1CD2}}श\raisebox{0.25ex}{\accentmark{12}{\char"0951}}म्माधू\raisebox{0.15ex}{\accentmark{12}{\char"1CF9}}श्चूतम्\kern 0.05em ।\hskip 0.12em plus 0.04em minus 0.02em सो\raisebox{0.25ex}{\accentmark{15}{\char"1CD2}}म\raisebox{0.25ex}{\accentmark{12}{\char"0951}}ःपूनानो\raisebox{0.25ex}{\accentmark{15}{\char"1CD2}}अ\raisebox{0.25ex}{\accentmark{12}{\char"0951}}र्षति\kern 0.05em ।\hskip 0.12em plus 0.04em minus 0.02em आ\raisebox{0.25ex}{\accentmark{15}{\char"1CD2}}भी\raisebox{0.25ex}{\accentmark{12}{\char"0951}}वाणीऋ\raisebox{0.25ex}{\accentmark{12}{\char"0951}}षीणां\raisebox{0.15ex}{\accentmark{12}{\char"1CF9}}सप्ता\raisebox{0.25ex}{\accentmark{15}{\char"1CD2}}नु\raisebox{0.25ex}{\accentmark{12}{\char"0951}}षत \nolinebreak\hskip 0.08em \mbox{॥१२॥}}} \par}

\vspace{0.8em}

            % --- End of Kandah Closing Floral Line & Counts ---
            \par\nopagebreak\vspace{0.9em plus 0.3em minus 0.1em}
            \centerline{\textcolor{black!35}{\rule{1.8cm}{0.35pt}}\;\raisebox{-0.22ex}{\textcolor{vedicmaroon}{\large\SixFlowerPetalDotted}}\;\textcolor{black!35}{\rule{1.8cm}{0.35pt}}}
            \nopagebreak\vspace{0.45em}
            \centerline{\devafont\textbf{\textcolor{AccentBlue}{॥ ऋक् सङ्ख्या : १२ ॥}}}
        
            \clearpage

            % --- SECTION ---
            
            \phantomsection  

            \section*{\centering एकादश खण्डः}
            \addcontentsline{toc}{section}{एकादश खण्डः} 
            \markright{एकादश खण्डः}
            \nopagebreak % Keep section title with first mantra
            
            % --- SECTION-LEVEL PROCEDURE LINK ---

                % --- SUBSECTIONS (Mantras) ---
                
                % Conditional rendering based on output_mode
                \par\filbreak

\phantomsection

\addcontentsline{toc}{subsection}{ऋक् १}

{\centering \footnotesize \textcolor{AccentPurple}{॥ गौरवीतिः शाक्स्यः पवमानः सोमः ककुप ॥} \par}

\vspace{0.12em}

{\centering \textbf{\textcolor{AccentBlue}{पा\raisebox{0.15ex}{\accentmark{12}{\char"1CF8}}वा\raisebox{0.25ex}{\accentmark{12}{\char"0951}}स्वामा\raisebox{0.25ex}{\accentmark{15}{\char"1CD2}}धू\raisebox{0.25ex}{\accentmark{12}{\char"0951}}मक्तमः\kern 0.05em ।\hskip 0.12em plus 0.04em minus 0.02em इन्द्रा\raisebox{0.25ex}{\accentmark{12}{\char"0951}}यसोमक्रातूवि\raisebox{0.15ex}{\accentmark{12}{\char"1CF9}}क्ता\raisebox{0.25ex}{\accentmark{12}{\char"0951}}मोमा\raisebox{0.25ex}{\accentmark{15}{\char"1CD2}}दा\raisebox{0.25ex}{\accentmark{12}{\char"0951}}ः\kern 0.05em ।\hskip 0.12em plus 0.04em minus 0.02em मा\raisebox{0.25ex}{\accentmark{15}{\char"1CD2}}हिद्युक्षा\raisebox{0.25ex}{\accentmark{15}{\char"1CD2}}तामोमा\raisebox{0.25ex}{\accentmark{15}{\char"1CD2}}दा\raisebox{0.25ex}{\accentmark{12}{\char"0951}}ः \nolinebreak\hskip 0.08em \mbox{॥१॥}}} \par}

\vspace{0.8em}
                % --- SUBSECTIONS (Mantras) ---
                
                % Conditional rendering based on output_mode
                \par\filbreak

\phantomsection

\addcontentsline{toc}{subsection}{ऋक् २}

{\centering \footnotesize \textcolor{AccentPurple}{॥ ऊर्ध्वसद्मा आङ्गिरसः पवमानः सोमः ककुप ॥} \par}

\vspace{0.12em}

{\centering \textbf{\textcolor{AccentBlue}{आभि\raisebox{0.25ex}{\accentmark{15}{\char"1CD2}}द्युम्नं\raisebox{0.25ex}{\accentmark{15}{\char"1CD2}}बृहद्याशाः\kern 0.05em ।\hskip 0.12em plus 0.04em minus 0.02em ई\raisebox{0.25ex}{\accentmark{15}{\char"1CD2}}ळा\raisebox{0.25ex}{\accentmark{12}{\char"0951}}स्पतेदीदीही\raisebox{0.25ex}{\accentmark{15}{\char"1CD2}}देवादेवायूम्\kern 0.05em ।\hskip 0.12em plus 0.04em minus 0.02em वीको\raisebox{0.25ex}{\accentmark{15}{\char"1CD2}}शं\raisebox{0.25ex}{\accentmark{12}{\char"0951}}माध्यामंयू\raisebox{0.25ex}{\accentmark{12}{\char"0951}}वा \nolinebreak\hskip 0.08em \mbox{॥२॥}}} \par}

\vspace{0.8em}
                % --- SUBSECTIONS (Mantras) ---
                
                % Conditional rendering based on output_mode
                \par\filbreak

\phantomsection

\addcontentsline{toc}{subsection}{ऋक् ३}

{\centering \footnotesize \textcolor{AccentPurple}{॥ ऋजिश्वा भारद्वाजः पवमानः सोमः ककुप ॥} \par}

\vspace{0.12em}

{\centering \textbf{\textcolor{AccentBlue}{आ\raisebox{0.25ex}{\accentmark{15}{\char"1CD2}}सोतापा\raisebox{0.25ex}{\accentmark{15}{\char"1CD2}}री\raisebox{0.25ex}{\accentmark{12}{\char"0951}}षिञ्चता\kern 0.05em ।\hskip 0.12em plus 0.04em minus 0.02em अश्व\raisebox{0.25ex}{\accentmark{15}{\char"1CD2}}न्न\raisebox{0.15ex}{\accentmark{12}{\char"1CF9}}स्तो\raisebox{0.25ex}{\accentmark{12}{\char"0951}}मामप्तूरं\raisebox{0.25ex}{\accentmark{12}{\char"0951}}राजस्तू\raisebox{0.25ex}{\accentmark{12}{\char"0951}}रम्\kern 0.05em ।\hskip 0.12em plus 0.04em minus 0.02em वनप्राक्षा\raisebox{0.25ex}{\accentmark{12}{\char"0951}}मूतप्रूतम् \nolinebreak\hskip 0.08em \mbox{॥३॥}}} \par}

\vspace{0.8em}
                % --- SUBSECTIONS (Mantras) ---
                
                % Conditional rendering based on output_mode
                \par\filbreak

\phantomsection

\addcontentsline{toc}{subsection}{ऋक् ४}

{\centering \footnotesize \textcolor{AccentPurple}{॥ कृतयशा आङ्गिरसः पवमानः सोमः ककुप ॥} \par}

\vspace{0.12em}

{\centering \textbf{\textcolor{AccentBlue}{एता\raisebox{0.25ex}{\accentmark{15}{\char"1CD2}}मुत्य\raisebox{0.15ex}{\accentmark{12}{\char"1CF9}}म्मा\raisebox{0.25ex}{\accentmark{12}{\char"0951}}दच्यू\raisebox{0.25ex}{\accentmark{15}{\char"1CD2}}तम्\raisebox{0.25ex}{\accentmark{12}{\char"0951}}\kern 0.05em ।\hskip 0.12em plus 0.04em minus 0.02em साहा\raisebox{0.25ex}{\accentmark{15}{\char"1CD2}}स्रा\raisebox{0.25ex}{\accentmark{12}{\char"0951}}धारंवृषाभ\raisebox{0.15ex}{\accentmark{12}{\char"1CF9}}न्दी\raisebox{0.25ex}{\accentmark{12}{\char"0951}}वोदूहम्\raisebox{0.25ex}{\accentmark{12}{\char"0951}}\kern 0.05em ।\hskip 0.12em plus 0.04em minus 0.02em वि\raisebox{0.25ex}{\accentmark{15}{\char"1CD2}}श्वा\raisebox{0.25ex}{\accentmark{12}{\char"0951}}वा\raisebox{0.15ex}{\accentmark{12}{\char"1CF9}}सू\raisebox{0.25ex}{\accentmark{12}{\char"0951}}निबिभ्रा\raisebox{0.25ex}{\accentmark{12}{\char"0951}}तम् \nolinebreak\hskip 0.08em \mbox{॥४॥}}} \par}

\vspace{0.8em}
                % --- SUBSECTIONS (Mantras) ---
                
                % Conditional rendering based on output_mode
                \par\filbreak

\phantomsection

\addcontentsline{toc}{subsection}{ऋक् ५}

{\centering \footnotesize \textcolor{AccentPurple}{॥ ऋणञ्चयो राजर्षिः पवमानः सोमः यवमध्या गायत्री ॥} \par}

\vspace{0.12em}

{\centering \textbf{\textcolor{AccentBlue}{सा\raisebox{0.15ex}{\accentmark{12}{\char"1CF8}}सुन्वोयोवा\raisebox{0.25ex}{\accentmark{15}{\char"1CD2}}सू\raisebox{0.25ex}{\accentmark{12}{\char"0951}}नाम्\kern 0.05em ।\hskip 0.12em plus 0.04em minus 0.02em यो\raisebox{0.25ex}{\accentmark{15}{\char"1CD2}}राया\raisebox{0.15ex}{\accentmark{12}{\char"1CF9}}मानेता\raisebox{0.25ex}{\accentmark{15}{\char"1CD2}}याइळा\raisebox{0.25ex}{\accentmark{12}{\char"0951}}नाम्\kern 0.05em ।\hskip 0.12em plus 0.04em minus 0.02em सो\raisebox{0.25ex}{\accentmark{15}{\char"1CD2}}मोय\raisebox{0.25ex}{\accentmark{15}{\char"1CD2}}स्सू\raisebox{0.25ex}{\accentmark{12}{\char"0951}}क्षीतीना\raisebox{0.25ex}{\accentmark{15}{\char"1CD2}}म् \nolinebreak\hskip 0.08em \mbox{॥५॥}}} \par}

\vspace{0.8em}
                % --- SUBSECTIONS (Mantras) ---
                
                % Conditional rendering based on output_mode
                \par\filbreak

\phantomsection

\addcontentsline{toc}{subsection}{ऋक् ६}

{\centering \footnotesize \textcolor{AccentPurple}{॥ शक्तिर्वासिष्ठः पवमानः सोमः ककुप ॥} \par}

\vspace{0.12em}

{\centering \textbf{\textcolor{AccentBlue}{त्वंह्यांट्ट्यगदै\raisebox{0.25ex}{\accentmark{12}{\char"0951}}व्या\kern 0.05em ।\hskip 0.12em plus 0.04em minus 0.02em पावा\raisebox{0.25ex}{\accentmark{12}{\char"0951}}मानाजनी\raisebox{0.25ex}{\accentmark{12}{\char"0951}}मानीधूमा\raisebox{0.25ex}{\accentmark{15}{\char"1CD2}}क्तामः\kern 0.05em ।\hskip 0.12em plus 0.04em minus 0.02em अ\raisebox{0.25ex}{\accentmark{15}{\char"1CD2}}मृता\raisebox{0.25ex}{\accentmark{12}{\char"0951}}त्वा\raisebox{0.15ex}{\accentmark{12}{\char"1CF9}}या\raisebox{0.25ex}{\accentmark{12}{\char"0951}}घो\raisebox{0.25ex}{\accentmark{15}{\char"1CD2}}षा\raisebox{0.25ex}{\accentmark{15}{\char"1CD2}}यान् \nolinebreak\hskip 0.08em \mbox{॥६॥}}} \par}

\vspace{0.8em}
                % --- SUBSECTIONS (Mantras) ---
                
                % Conditional rendering based on output_mode
                \par\filbreak

\phantomsection

\addcontentsline{toc}{subsection}{ऋक् ७}

{\centering \footnotesize \textcolor{AccentPurple}{॥ ऊरुराङ्गिरसः पवमानः सोमः ककुप ॥} \par}

\vspace{0.12em}

{\centering \textbf{\textcolor{AccentBlue}{एषस्या\raisebox{0.25ex}{\accentmark{15}{\char"1CD2}}धा\raisebox{0.25ex}{\accentmark{15}{\char"1CD2}}रा\raisebox{0.25ex}{\accentmark{12}{\char"0951}}यासूत\raisebox{0.25ex}{\accentmark{15}{\char"1CD2}}ः\kern 0.05em ।\hskip 0.12em plus 0.04em minus 0.02em अव्यावा\raisebox{0.25ex}{\accentmark{15}{\char"1CD2}}रे\raisebox{0.25ex}{\accentmark{12}{\char"0951}}भिःपवतेमादि\raisebox{0.25ex}{\accentmark{12}{\char"0951}}न्तामः\kern 0.05em ।\hskip 0.12em plus 0.04em minus 0.02em क्रीळ\raisebox{0.25ex}{\accentmark{12}{\char"0951}}न्नूर्मिरापा\raisebox{0.25ex}{\accentmark{15}{\char"1CD2}}मी\raisebox{0.25ex}{\accentmark{12}{\char"0951}}वा \nolinebreak\hskip 0.08em \mbox{॥७॥}}} \par}

\vspace{0.8em}
                % --- SUBSECTIONS (Mantras) ---
                
                % Conditional rendering based on output_mode
                
                % --- SUBSECTIONS (Mantras) ---
                
                % Conditional rendering based on output_mode
                

            % --- End of Kandah Closing Floral Line & Counts ---
            \par\nopagebreak\vspace{0.9em plus 0.3em minus 0.1em}
            \centerline{\textcolor{black!35}{\rule{1.8cm}{0.35pt}}\;\raisebox{-0.22ex}{\textcolor{vedicmaroon}{\large\SixFlowerPetalDotted}}\;\textcolor{black!35}{\rule{1.8cm}{0.35pt}}}
            \nopagebreak\vspace{0.45em}
            \centerline{\devafont\textbf{\textcolor{AccentBlue}{॥ ऋक् सङ्ख्या : ९ ॥}}}


    % --- SUPERSECTION ---
    \clearpage           
    \phantomsection              
    \chapter*{\centering अथ आरण पाठः}
    \addcontentsline{toc}{chapter}{आरण पाठः}
    \markboth{आरण पाठः}{आरण पाठः}
    
    % GLUE FIX: Prevent page break between Chapter and First Section
    \nopagebreak

        

            % --- SECTION ---
            
            \phantomsection  

            \section*{\centering अथ प्रथमः खण्डः}
            \addcontentsline{toc}{section}{प्रथमः खण्डः} 
            \markright{प्रथमः खण्डः}
            \nopagebreak % Keep section title with first mantra
            
            % --- SECTION-LEVEL PROCEDURE LINK ---

                % --- SUBSECTIONS (Mantras) ---
                
                % Conditional rendering based on output_mode
                \par\filbreak

\phantomsection

\addcontentsline{toc}{subsection}{ऋक् १}

{\centering \textbf{\textcolor{AccentBlue}{इन्द्रोरा\raisebox{0.25ex}{\accentmark{15}{\char"1CD2}}जाजग\raisebox{0.25ex}{\accentmark{12}{\char"0951}}तश्वर्षाणीनां\kern 0.05em ।\hskip 0.12em plus 0.04em minus 0.02em आ\raisebox{0.15ex}{\accentmark{12}{\char"1CF8}}धि\raisebox{0.25ex}{\accentmark{12}{\char"0951}}क्षामा\raisebox{0.25ex}{\accentmark{15}{\char"1CD2}}विश्वा\raisebox{0.15ex}{\accentmark{12}{\char"1CF9}}रू\raisebox{0.25ex}{\accentmark{12}{\char"0951}}पंया\raisebox{0.25ex}{\accentmark{15}{\char"1CD2}}द\raisebox{0.25ex}{\accentmark{12}{\char"0951}}स्या\kern 0.05em ।\hskip 0.12em plus 0.04em minus 0.02em ता\raisebox{0.25ex}{\accentmark{15}{\char"1CD2}}तो\raisebox{0.25ex}{\accentmark{12}{\char"0951}}ददातीदाशू\raisebox{0.25ex}{\accentmark{15}{\char"1CD2}}षेवा\raisebox{0.25ex}{\accentmark{15}{\char"1CD2}}सू\raisebox{0.25ex}{\accentmark{12}{\char"0951}}नी\kern 0.05em ।\hskip 0.12em plus 0.04em minus 0.02em चोदद्राआधाऊ\raisebox{0.25ex}{\accentmark{15}{\char"1CD2}}पास्तुतञ्चीदर्वा\raisebox{0.25ex}{\accentmark{15}{\char"1CD2}}क् \nolinebreak\hskip 0.08em \mbox{॥१॥}}} \par}

\vspace{0.8em}
                % --- SUBSECTIONS (Mantras) ---
                
                % Conditional rendering based on output_mode
                \par\filbreak

\phantomsection

\addcontentsline{toc}{subsection}{ऋक् २}

{\centering \textbf{\textcolor{AccentBlue}{स\raisebox{0.25ex}{\accentmark{15}{\char"1CD2}}न्तेपा\raisebox{0.15ex}{\accentmark{12}{\char"1CF9}}यां\raisebox{0.25ex}{\accentmark{12}{\char"0951}}सीसा\raisebox{0.25ex}{\accentmark{15}{\char"1CD2}}मूयन्तूवा\raisebox{0.25ex}{\accentmark{15}{\char"1CD2}}जाः\raisebox{0.25ex}{\accentmark{12}{\char"0951}}\kern 0.05em ।\hskip 0.12em plus 0.04em minus 0.02em सं\raisebox{0.25ex}{\accentmark{15}{\char"1CD2}}वृष्ण्या\raisebox{0.25ex}{\accentmark{12}{\char"0951}}न्यभीमातीषा\raisebox{0.25ex}{\accentmark{15}{\char"1CD2}}हाः\raisebox{0.25ex}{\accentmark{12}{\char"0951}}\kern 0.05em ।\hskip 0.12em plus 0.04em minus 0.02em आप्यायमानोआमृ\raisebox{0.25ex}{\accentmark{15}{\char"1CD2}}ता\raisebox{0.25ex}{\accentmark{12}{\char"0951}}यसोमादीवि\raisebox{0.25ex}{\accentmark{15}{\char"1CD2}}श्रा\raisebox{0.25ex}{\accentmark{15}{\char"1CD2}}वां\raisebox{0.25ex}{\accentmark{12}{\char"0951}}स्यूक्तामा\raisebox{0.25ex}{\accentmark{15}{\char"1CD2}}नीधत्स्वा \nolinebreak\hskip 0.08em \mbox{॥२॥}}} \par}

\vspace{0.8em}
                % --- SUBSECTIONS (Mantras) ---
                
                % Conditional rendering based on output_mode
                \par\filbreak

\phantomsection

\addcontentsline{toc}{subsection}{ऋक् ३}

{\centering \textbf{\textcolor{AccentBlue}{त्वा\raisebox{0.25ex}{\accentmark{15}{\char"1CD2}}मीमाओषा\raisebox{0.25ex}{\accentmark{12}{\char"0951}}धीस्सोमा\raisebox{0.25ex}{\accentmark{15}{\char"1CD2}}वि\raisebox{0.25ex}{\accentmark{15}{\char"1CD2}}श्वाः\raisebox{0.25ex}{\accentmark{12}{\char"0951}}\kern 0.05em ।\hskip 0.12em plus 0.04em minus 0.02em त्वा\raisebox{0.25ex}{\accentmark{15}{\char"1CD2}}मापोआजनायस्त्वङगाः\kern 0.05em ।\hskip 0.12em plus 0.04em minus 0.02em त्वामा\raisebox{0.25ex}{\accentmark{12}{\char"0951}}तनोरूर्वा३्न्ता\raisebox{0.25ex}{\accentmark{15}{\char"1CD2}}री\raisebox{0.25ex}{\accentmark{12}{\char"0951}}क्षाम्\kern 0.05em ।\hskip 0.12em plus 0.04em minus 0.02em त्व\raisebox{0.15ex}{\accentmark{12}{\char"1CF8}}ञ्ज्यो\raisebox{0.15ex}{\accentmark{12}{\char"1CF8}}तीषावीतामोववर्क्था\kern 0.05em ।\hskip 0.12em plus 0.04em minus 0.02em मायीवर्चीआथो\raisebox{0.25ex}{\accentmark{12}{\char"0951}}भा\raisebox{0.25ex}{\accentmark{15}{\char"1CD2}}गम्\kern 0.05em ।\hskip 0.12em plus 0.04em minus 0.02em आ\raisebox{0.15ex}{\accentmark{12}{\char"1CF8}}थो\raisebox{0.25ex}{\accentmark{12}{\char"0951}}यज्ञा\raisebox{0.25ex}{\accentmark{15}{\char"1CD2}}स्यायत्पा\raisebox{0.25ex}{\accentmark{15}{\char"1CD2}}याः\kern 0.05em ।\hskip 0.12em plus 0.04em minus 0.02em प\raisebox{0.25ex}{\accentmark{15}{\char"1CD2}}रमेष्ठी\raisebox{0.25ex}{\accentmark{15}{\char"1CD2}}प्रा\raisebox{0.25ex}{\accentmark{15}{\char"1CD2}}जा\raisebox{0.25ex}{\accentmark{15}{\char"1CD2}}पा\raisebox{0.25ex}{\accentmark{12}{\char"0951}}तिः\kern 0.05em ।\hskip 0.12em plus 0.04em minus 0.02em दीवि\raisebox{0.25ex}{\accentmark{15}{\char"1CD2}}द्या\raisebox{0.25ex}{\accentmark{12}{\char"0951}}मी\raisebox{0.25ex}{\accentmark{12}{\char"0951}}वद्रंहतु \nolinebreak\hskip 0.08em \mbox{॥३॥}}} \par}

\vspace{0.8em}
                % --- SUBSECTIONS (Mantras) ---
                
                % Conditional rendering based on output_mode
                \par\filbreak

\phantomsection

\addcontentsline{toc}{subsection}{ऋक् ४}

{\centering \textbf{\textcolor{AccentBlue}{इ\raisebox{0.15ex}{\accentmark{12}{\char"1CF9}}दञ्छ\raisebox{0.15ex}{\accentmark{12}{\char"1CF9}}न्दो\raisebox{0.25ex}{\accentmark{12}{\char"0951}}जमादग्नेऋता\raisebox{0.25ex}{\accentmark{15}{\char"1CD2}}वृधाः\raisebox{0.25ex}{\accentmark{12}{\char"0951}}\kern 0.05em ।\hskip 0.12em plus 0.04em minus 0.02em ये\raisebox{0.15ex}{\accentmark{12}{\char"1CF8}}ना\raisebox{0.25ex}{\accentmark{12}{\char"0951}}देवा\raisebox{0.25ex}{\accentmark{15}{\char"1CD2}}सो\raisebox{0.25ex}{\accentmark{12}{\char"0951}}आमृतत्वामा\raisebox{0.25ex}{\accentmark{15}{\char"1CD2}}यान्\kern 0.05em ।\hskip 0.12em plus 0.04em minus 0.02em त\raisebox{0.25ex}{\accentmark{15}{\char"1CD2}}त्रद्या\raisebox{0.25ex}{\accentmark{15}{\char"1CD2}}वा\raisebox{0.25ex}{\accentmark{12}{\char"0951}}पृथिवी\raisebox{0.25ex}{\accentmark{15}{\char"1CD2}}धक्ता\raisebox{0.25ex}{\accentmark{15}{\char"1CD2}}मस्मा\raisebox{0.25ex}{\accentmark{15}{\char"1CD2}}न्\kern 0.05em ।\hskip 0.12em plus 0.04em minus 0.02em य\raisebox{0.15ex}{\accentmark{12}{\char"1CF8}}त्रा\raisebox{0.25ex}{\accentmark{15}{\char"1CD2}}देवानांगूह्यन्नीधाम \nolinebreak\hskip 0.08em \mbox{॥४॥}}} \par}

\vspace{0.8em}
                % --- SUBSECTIONS (Mantras) ---
                
                % Conditional rendering based on output_mode
                \par\filbreak

\phantomsection

\addcontentsline{toc}{subsection}{ऋक् ५}

{\centering \textbf{\textcolor{AccentBlue}{सा\raisebox{0.25ex}{\accentmark{15}{\char"1CD2}}मान्याय\raisebox{0.25ex}{\accentmark{15}{\char"1CD2}}न्त्यूपा\raisebox{0.25ex}{\accentmark{12}{\char"0951}}यन्त्यन्याः\raisebox{0.25ex}{\accentmark{15}{\char"1CD2}}\kern 0.05em ।\hskip 0.12em plus 0.04em minus 0.02em समा\raisebox{0.25ex}{\accentmark{15}{\char"1CD2}}नामुर्व\raisebox{0.25ex}{\accentmark{15}{\char"1CD2}}न्नद्याः१पृणन्ति\kern 0.05em ।\hskip 0.12em plus 0.04em minus 0.02em ता\raisebox{0.25ex}{\accentmark{15}{\char"1CD2}}मू\raisebox{0.25ex}{\accentmark{15}{\char"1CD2}}शू\raisebox{0.25ex}{\accentmark{15}{\char"1CD2}}चिंशूचायोदीदीवांसम्\raisebox{0.25ex}{\accentmark{12}{\char"0951}}\kern 0.05em ।\hskip 0.12em plus 0.04em minus 0.02em आपान्नापा\raisebox{0.15ex}{\accentmark{12}{\char"1CF8}}तामू\raisebox{0.15ex}{\accentmark{12}{\char"1CF9}}पा\raisebox{0.25ex}{\accentmark{12}{\char"0951}}यन्त्यापाः\raisebox{0.25ex}{\accentmark{12}{\char"0951}} \nolinebreak\hskip 0.08em \mbox{॥५॥}}} \par}

\vspace{0.8em}
                % --- SUBSECTIONS (Mantras) ---
                
                % Conditional rendering based on output_mode
                \par\filbreak

\phantomsection

\addcontentsline{toc}{subsection}{ऋक् ६}

{\centering \textbf{\textcolor{AccentBlue}{तेम\raisebox{0.25ex}{\accentmark{12}{\char"0951}}न्वतःप्रथामन्ना\raisebox{0.25ex}{\accentmark{15}{\char"1CD2}}मागोनां\raisebox{0.25ex}{\accentmark{12}{\char"0951}}\kern 0.05em ।\hskip 0.12em plus 0.04em minus 0.02em त्रिस्सप्ताः\raisebox{0.15ex}{\accentmark{12}{\char"1CF9}}पा\raisebox{0.15ex}{\accentmark{12}{\char"1CF9}}रा\raisebox{0.25ex}{\accentmark{12}{\char"0951}}मन्ना\raisebox{0.25ex}{\accentmark{15}{\char"1CD2}}माजाना\raisebox{0.25ex}{\accentmark{12}{\char"0951}}म्\kern 0.05em ।\hskip 0.12em plus 0.04em minus 0.02em ताजाना\raisebox{0.25ex}{\accentmark{12}{\char"0951}}तिरभ्या१नूषातक्षाः\raisebox{0.25ex}{\accentmark{15}{\char"1CD2}}\kern 0.05em ।\hskip 0.12em plus 0.04em minus 0.02em आविर्भू\raisebox{0.25ex}{\accentmark{12}{\char"0951}}वन्नारू\raisebox{0.15ex}{\accentmark{12}{\char"1CF9}}णीर्याशा\raisebox{0.25ex}{\accentmark{12}{\char"0951}}सा\raisebox{0.25ex}{\accentmark{15}{\char"1CD2}}गा\raisebox{0.25ex}{\accentmark{15}{\char"1CD2}}वाः\raisebox{0.25ex}{\accentmark{12}{\char"0951}} \nolinebreak\hskip 0.08em \mbox{॥६॥}}} \par}

\vspace{0.8em}
                % --- SUBSECTIONS (Mantras) ---
                
                % Conditional rendering based on output_mode
                \par\filbreak

\phantomsection

\addcontentsline{toc}{subsection}{ऋक् ७}

{\centering \textbf{\textcolor{AccentBlue}{साहार्षाभासाहावात्साऊ\raisebox{0.25ex}{\accentmark{15}{\char"1CD2}}देता\raisebox{0.25ex}{\accentmark{12}{\char"0951}}\kern 0.05em ।\hskip 0.12em plus 0.04em minus 0.02em विश्वारूपा\raisebox{0.25ex}{\accentmark{15}{\char"1CD2}}णीबि\raisebox{0.15ex}{\accentmark{12}{\char"1CF8}}भ्रा\raisebox{0.25ex}{\accentmark{12}{\char"0951}}तीद्यूर्ध्नीः\kern 0.05em ।\hskip 0.12em plus 0.04em minus 0.02em ऊरूः\raisebox{0.25ex}{\accentmark{15}{\char"1CD2}}पृथू\raisebox{0.25ex}{\accentmark{15}{\char"1CD2}}रायं\raisebox{0.25ex}{\accentmark{15}{\char"1CD2}}वोअस्तुलोकः\raisebox{0.25ex}{\accentmark{15}{\char"1CD2}}\kern 0.05em ।\hskip 0.12em plus 0.04em minus 0.02em ईम\raisebox{0.25ex}{\accentmark{15}{\char"1CD2}}आ\raisebox{0.25ex}{\accentmark{15}{\char"1CD2}}पा\raisebox{0.25ex}{\accentmark{12}{\char"0951}}सुप्रापाणा\raisebox{0.25ex}{\accentmark{15}{\char"1CD2}}ईह\raisebox{0.25ex}{\accentmark{15}{\char"1CD2}}स्ता \nolinebreak\hskip 0.08em \mbox{॥७॥}}} \par}

\vspace{0.8em}
                % --- SUBSECTIONS (Mantras) ---
                
                % Conditional rendering based on output_mode
                \par\filbreak

\phantomsection

\addcontentsline{toc}{subsection}{ऋक् ८}

{\centering \textbf{\textcolor{AccentBlue}{आरत्सू\raisebox{0.25ex}{\accentmark{12}{\char"0951}}पर्णाःपा\raisebox{0.25ex}{\accentmark{12}{\char"0951}}ततु\kern 0.05em ।\hskip 0.12em plus 0.04em minus 0.02em माश्वानामा\raisebox{0.25ex}{\accentmark{15}{\char"1CD2}}शी\raisebox{0.25ex}{\accentmark{12}{\char"0951}}तङ्करात्\kern 0.05em ।\hskip 0.12em plus 0.04em minus 0.02em तीर्था\raisebox{0.25ex}{\accentmark{15}{\char"1CD2}}नी\raisebox{0.25ex}{\accentmark{12}{\char"0951}}सुप्रा\raisebox{0.25ex}{\accentmark{15}{\char"1CD2}}पा\raisebox{0.25ex}{\accentmark{15}{\char"1CD2}}णानी\raisebox{0.25ex}{\accentmark{12}{\char"0951}}\kern 0.05em ।\hskip 0.12em plus 0.04em minus 0.02em सूग\raisebox{0.25ex}{\accentmark{15}{\char"1CD2}}ङगोषा\raisebox{0.25ex}{\accentmark{15}{\char"1CD2}}दनंकरात् \nolinebreak\hskip 0.08em \mbox{॥८॥}}} \par}

\vspace{0.8em}
                % --- SUBSECTIONS (Mantras) ---
                
                % Conditional rendering based on output_mode
                \par\filbreak

\phantomsection

\addcontentsline{toc}{subsection}{ऋक् ९}

{\centering \textbf{\textcolor{AccentBlue}{अग्नी\raisebox{0.25ex}{\accentmark{15}{\char"1CD2}}मीळेपूरोहि\raisebox{0.25ex}{\accentmark{12}{\char"0951}}तम्\kern 0.05em ।\hskip 0.12em plus 0.04em minus 0.02em यज्ञस्या\raisebox{0.25ex}{\accentmark{12}{\char"0951}}देवा\raisebox{0.25ex}{\accentmark{12}{\char"0951}}मृत्वी\raisebox{0.25ex}{\accentmark{15}{\char"1CD2}}जं\raisebox{0.25ex}{\accentmark{12}{\char"0951}}\kern 0.05em ।\hskip 0.12em plus 0.04em minus 0.02em हो\raisebox{0.25ex}{\accentmark{15}{\char"1CD2}}तारंरानाधा\raisebox{0.25ex}{\accentmark{15}{\char"1CD2}}ता\raisebox{0.25ex}{\accentmark{12}{\char"0951}}मम् \nolinebreak\hskip 0.08em \mbox{॥९॥}}} \par}

\vspace{0.8em}

            % --- End of Kandah Closing Floral Line & Counts ---
            \par\nopagebreak\vspace{0.9em plus 0.3em minus 0.1em}
            \centerline{\textcolor{black!35}{\rule{1.8cm}{0.35pt}}\;\raisebox{-0.22ex}{\textcolor{vedicmaroon}{\large\SixFlowerPetalDotted}}\;\textcolor{black!35}{\rule{1.8cm}{0.35pt}}}
            \nopagebreak\vspace{0.45em}
            \centerline{\devafont\textbf{\textcolor{AccentBlue}{॥ ऋक् सङ्ख्या : ९ ॥}}}
        
            \clearpage

            % --- SECTION ---
            
            \phantomsection  

            \section*{\centering अथ द्वितीयः खण्डः}
            \addcontentsline{toc}{section}{द्वितीयः खण्डः} 
            \markright{द्वितीयः खण्डः}
            \nopagebreak % Keep section title with first mantra
            
            % --- SECTION-LEVEL PROCEDURE LINK ---

                % --- SUBSECTIONS (Mantras) ---
                
                % Conditional rendering based on output_mode
                \par\filbreak

\phantomsection

\addcontentsline{toc}{subsection}{ऋक् १}

{\centering \textbf{\textcolor{AccentBlue}{म\raisebox{0.25ex}{\accentmark{15}{\char"1CD2}}न्ये\raisebox{0.25ex}{\accentmark{12}{\char"0951}}वान्द्यावापृथीवीसूभोज\raisebox{0.25ex}{\accentmark{12}{\char"0951}}सा\kern 0.05em ।\hskip 0.12em plus 0.04em minus 0.02em येआ\raisebox{0.25ex}{\accentmark{15}{\char"1CD2}}प्रा\raisebox{0.25ex}{\accentmark{12}{\char"0951}}थेथामामी\raisebox{0.25ex}{\accentmark{15}{\char"1CD2}}तामाभीयो\raisebox{0.25ex}{\accentmark{15}{\char"1CD2}}ज\raisebox{0.25ex}{\accentmark{12}{\char"0951}}नम्\kern 0.05em ।\hskip 0.12em plus 0.04em minus 0.02em द्या\raisebox{0.25ex}{\accentmark{15}{\char"1CD2}}वा\raisebox{0.25ex}{\accentmark{12}{\char"0951}}पृथीवीभा\raisebox{0.25ex}{\accentmark{15}{\char"1CD2}}वा\raisebox{0.25ex}{\accentmark{12}{\char"0951}}तंस्योने\kern 0.05em ।\hskip 0.12em plus 0.04em minus 0.02em ते\raisebox{0.25ex}{\accentmark{15}{\char"1CD2}}नोमुञ्चातामंहा\raisebox{0.25ex}{\accentmark{12}{\char"0951}}सः \nolinebreak\hskip 0.08em \mbox{॥१॥}}} \par}

\vspace{0.8em}
                % --- SUBSECTIONS (Mantras) ---
                
                % Conditional rendering based on output_mode
                \par\filbreak

\phantomsection

\addcontentsline{toc}{subsection}{ऋक् २}

{\centering \textbf{\textcolor{AccentBlue}{या\raisebox{0.25ex}{\accentmark{15}{\char"1CD2}}शासामाद्या\raisebox{0.25ex}{\accentmark{15}{\char"1CD2}}वा\raisebox{0.25ex}{\accentmark{12}{\char"0951}}पृथिवी\kern 0.05em ।\hskip 0.12em plus 0.04em minus 0.02em भा\raisebox{0.25ex}{\accentmark{15}{\char"1CD2}}गे\raisebox{0.25ex}{\accentmark{12}{\char"0951}}नइन्द्रबृहास्पा\raisebox{0.25ex}{\accentmark{15}{\char"1CD2}}ती\kern 0.05em ।\hskip 0.12em plus 0.04em minus 0.02em या\raisebox{0.25ex}{\accentmark{15}{\char"1CD2}}शोभा\raisebox{0.25ex}{\accentmark{15}{\char"1CD2}}ग\raisebox{0.25ex}{\accentmark{12}{\char"0951}}स्यविन्दतु\kern 0.05em ।\hskip 0.12em plus 0.04em minus 0.02em या\raisebox{0.15ex}{\accentmark{12}{\char"1CF8}}शो\raisebox{0.25ex}{\accentmark{12}{\char"0951}}माप्राती\raisebox{0.25ex}{\accentmark{12}{\char"0951}}मुच्यतां\kern 0.05em ।\hskip 0.12em plus 0.04em minus 0.02em यशस्व्याद्व्यस्यास्संसा\raisebox{0.25ex}{\accentmark{15}{\char"1CD2}}दाः\raisebox{0.25ex}{\accentmark{12}{\char"0951}}\kern 0.05em ।\hskip 0.12em plus 0.04em minus 0.02em आहं\raisebox{0.25ex}{\accentmark{15}{\char"1CD2}}प्रा\raisebox{0.25ex}{\accentmark{15}{\char"1CD2}}वादीता\raisebox{0.25ex}{\accentmark{15}{\char"1CD2}}स्यां\raisebox{0.25ex}{\accentmark{12}{\char"0951}}\kern 0.05em ।\hskip 0.12em plus 0.04em minus 0.02em आहं\raisebox{0.25ex}{\accentmark{15}{\char"1CD2}}भू\raisebox{0.25ex}{\accentmark{12}{\char"0951}}यासमूक्तामः\raisebox{0.25ex}{\accentmark{15}{\char"1CD2}} \nolinebreak\hskip 0.08em \mbox{॥२॥}}} \par}

\vspace{0.8em}
                % --- SUBSECTIONS (Mantras) ---
                
                % Conditional rendering based on output_mode
                \par\filbreak

\phantomsection

\addcontentsline{toc}{subsection}{ऋक् ३}

{\centering \textbf{\textcolor{AccentBlue}{प्र\raisebox{0.15ex}{\accentmark{12}{\char"1CF8}}क्षा\raisebox{0.25ex}{\accentmark{12}{\char"0951}}स्यावृ\raisebox{0.25ex}{\accentmark{15}{\char"1CD2}}ष्णो\raisebox{0.25ex}{\accentmark{12}{\char"0951}}आरूषस्यानू\raisebox{0.25ex}{\accentmark{15}{\char"1CD2}}मा\raisebox{0.25ex}{\accentmark{15}{\char"1CD2}}हाः\raisebox{0.25ex}{\accentmark{12}{\char"0951}}\kern 0.05em ।\hskip 0.12em plus 0.04em minus 0.02em प्रा\raisebox{0.25ex}{\accentmark{15}{\char"1CD2}}नुवो\raisebox{0.25ex}{\accentmark{12}{\char"0951}}चोवीद\raisebox{0.15ex}{\accentmark{12}{\char"1CF9}}था\raisebox{0.25ex}{\accentmark{12}{\char"0951}}जाता\raisebox{0.25ex}{\accentmark{15}{\char"1CD2}}वेदसे\kern 0.05em ।\hskip 0.12em plus 0.04em minus 0.02em वै\raisebox{0.25ex}{\accentmark{15}{\char"1CD2}}श्वानाराया\raisebox{0.25ex}{\accentmark{12}{\char"0951}}मा\raisebox{0.15ex}{\accentmark{12}{\char"1CF9}}तिर्ना\raisebox{0.15ex}{\accentmark{12}{\char"1CF9}}व्या\raisebox{0.25ex}{\accentmark{12}{\char"0951}}सेशू\raisebox{0.25ex}{\accentmark{15}{\char"1CD2}}चीः\raisebox{0.25ex}{\accentmark{12}{\char"0951}}\kern 0.05em ।\hskip 0.12em plus 0.04em minus 0.02em सो\raisebox{0.25ex}{\accentmark{15}{\char"1CD2}}मा\raisebox{0.25ex}{\accentmark{12}{\char"0951}}इवपवातेचा\raisebox{0.15ex}{\accentmark{12}{\char"1CF9}}रू\raisebox{0.25ex}{\accentmark{12}{\char"0951}}रग्ना\raisebox{0.25ex}{\accentmark{12}{\char"0951}}ये \nolinebreak\hskip 0.08em \mbox{॥३॥}}} \par}

\vspace{0.8em}
                % --- SUBSECTIONS (Mantras) ---
                
                % Conditional rendering based on output_mode
                \par\filbreak

\phantomsection

\addcontentsline{toc}{subsection}{ऋक् ४}

{\centering \textbf{\textcolor{AccentBlue}{वि\raisebox{0.15ex}{\accentmark{12}{\char"1CF8}}श्वे\raisebox{0.25ex}{\accentmark{12}{\char"0951}}देवा\raisebox{0.25ex}{\accentmark{15}{\char"1CD2}}मामा\raisebox{0.25ex}{\accentmark{12}{\char"0951}}शृण्वन्तु\raisebox{0.25ex}{\accentmark{15}{\char"1CD2}}यज्ञं\raisebox{0.25ex}{\accentmark{15}{\char"1CD2}}\kern 0.05em ।\hskip 0.12em plus 0.04em minus 0.02em ऊ\raisebox{0.15ex}{\accentmark{12}{\char"1CF9}}भेरो\raisebox{0.25ex}{\accentmark{15}{\char"1CD2}}द\raisebox{0.25ex}{\accentmark{12}{\char"0951}}सीआपान्ना\raisebox{0.15ex}{\accentmark{12}{\char"1CF9}}पा\raisebox{0.25ex}{\accentmark{12}{\char"0951}}च्चामन्मा\raisebox{0.25ex}{\accentmark{12}{\char"0951}}\kern 0.05em ।\hskip 0.12em plus 0.04em minus 0.02em मा\raisebox{0.25ex}{\accentmark{12}{\char"0951}}वोवा\raisebox{0.25ex}{\accentmark{12}{\char"0951}}चांसिपारीचार्क्षाणि\raisebox{0.25ex}{\accentmark{12}{\char"0951}}वोचं\kern 0.05em ।\hskip 0.12em plus 0.04em minus 0.02em सुम्नेष्विद्वोअन्ता\raisebox{0.25ex}{\accentmark{15}{\char"1CD2}}मा\raisebox{0.25ex}{\accentmark{12}{\char"0951}}मदेम \nolinebreak\hskip 0.08em \mbox{॥४॥}}} \par}

\vspace{0.8em}
                % --- SUBSECTIONS (Mantras) ---
                
                % Conditional rendering based on output_mode
                \par\filbreak

\phantomsection

\addcontentsline{toc}{subsection}{ऋक् ५}

{\centering \textbf{\textcolor{AccentBlue}{एवा\raisebox{0.25ex}{\accentmark{15}{\char"1CD2}}ह्या\raisebox{0.25ex}{\accentmark{15}{\char"1CD2}}सी\raisebox{0.25ex}{\accentmark{12}{\char"0951}}वीरायुः\kern 0.05em ।\hskip 0.12em plus 0.04em minus 0.02em एवा\raisebox{0.25ex}{\accentmark{15}{\char"1CD2}}शू\raisebox{0.25ex}{\accentmark{15}{\char"1CD2}}रा\raisebox{0.25ex}{\accentmark{12}{\char"0951}}ऊ\raisebox{0.25ex}{\accentmark{15}{\char"1CD2}}ता\raisebox{0.25ex}{\accentmark{15}{\char"1CD2}}स्थीरः\raisebox{0.25ex}{\accentmark{15}{\char"1CD2}}\kern 0.05em ।\hskip 0.12em plus 0.04em minus 0.02em एवा\raisebox{0.25ex}{\accentmark{15}{\char"1CD2}}तेरा\raisebox{0.25ex}{\accentmark{15}{\char"1CD2}}ध्याम्मा\raisebox{0.25ex}{\accentmark{15}{\char"1CD2}}नाः\raisebox{0.25ex}{\accentmark{12}{\char"0951}} \nolinebreak\hskip 0.08em \mbox{॥५॥}}} \par}

\vspace{0.8em}
                % --- SUBSECTIONS (Mantras) ---
                
                % Conditional rendering based on output_mode
                \par\filbreak

\phantomsection

\addcontentsline{toc}{subsection}{ऋक् ६}

{\centering \textbf{\textcolor{AccentBlue}{आ\raisebox{0.25ex}{\accentmark{15}{\char"1CD2}}प्रागा\raisebox{0.25ex}{\accentmark{12}{\char"0951}}त्भद्रा\raisebox{0.15ex}{\accentmark{12}{\char"1CF9}}यू\raisebox{0.25ex}{\accentmark{12}{\char"0951}}वातिः\raisebox{0.25ex}{\accentmark{12}{\char"0951}}\kern 0.05em ।\hskip 0.12em plus 0.04em minus 0.02em अ\raisebox{0.25ex}{\accentmark{15}{\char"1CD2}}ह्नः\raisebox{0.25ex}{\accentmark{12}{\char"0951}}केतु\raisebox{0.25ex}{\accentmark{15}{\char"1CD2}}न्सा\raisebox{0.25ex}{\accentmark{15}{\char"1CD2}}मी\raisebox{0.25ex}{\accentmark{12}{\char"0951}}प्स्यती\kern 0.05em ।\hskip 0.12em plus 0.04em minus 0.02em आ\raisebox{0.15ex}{\accentmark{12}{\char"1CF8}}भु\raisebox{0.25ex}{\accentmark{12}{\char"0951}}त्भाद्रा\raisebox{0.25ex}{\accentmark{15}{\char"1CD2}}नीवे\raisebox{0.25ex}{\accentmark{15}{\char"1CD2}}शा\raisebox{0.25ex}{\accentmark{12}{\char"0951}}नी\kern 0.05em ।\hskip 0.12em plus 0.04em minus 0.02em विश्वा\raisebox{0.25ex}{\accentmark{12}{\char"0951}}स्याज\raisebox{0.15ex}{\accentmark{12}{\char"1CF9}}ग\raisebox{0.25ex}{\accentmark{12}{\char"0951}}तोरा\raisebox{0.25ex}{\accentmark{15}{\char"1CD2}}त्रिः\raisebox{0.25ex}{\accentmark{12}{\char"0951}}\kern 0.05em ।\hskip 0.12em plus 0.04em minus 0.02em उग्रं\raisebox{0.25ex}{\accentmark{15}{\char"1CD2}}वा\raisebox{0.25ex}{\accentmark{15}{\char"1CD2}}जी\raisebox{0.25ex}{\accentmark{15}{\char"1CD2}}न्याक्र\raisebox{0.25ex}{\accentmark{15}{\char"1CD2}}मीत् \nolinebreak\hskip 0.08em \mbox{॥६॥}}} \par}

\vspace{0.8em}
                % --- SUBSECTIONS (Mantras) ---
                
                % Conditional rendering based on output_mode
                \par\filbreak

\phantomsection

\addcontentsline{toc}{subsection}{ऋक् ७}

{\centering \textbf{\textcolor{AccentBlue}{अग्नी\raisebox{0.25ex}{\accentmark{15}{\char"1CD2}}राश्मीज\raisebox{0.25ex}{\accentmark{15}{\char"1CD2}}न्मा\raisebox{0.25ex}{\accentmark{12}{\char"0951}}नाजाता\raisebox{0.25ex}{\accentmark{15}{\char"1CD2}}वे\raisebox{0.25ex}{\accentmark{12}{\char"0951}}दाः\kern 0.05em ।\hskip 0.12em plus 0.04em minus 0.02em घृतम्मे\raisebox{0.25ex}{\accentmark{15}{\char"1CD2}}चाक्षु\raisebox{0.25ex}{\accentmark{12}{\char"0951}}रामृत\raisebox{0.25ex}{\accentmark{12}{\char"0951}}म्माआसन्\kern 0.05em ।\hskip 0.12em plus 0.04em minus 0.02em त्रिधा\raisebox{0.25ex}{\accentmark{15}{\char"1CD2}}स्त्वा\raisebox{0.15ex}{\accentmark{12}{\char"1CF9}}र्कोरा\raisebox{0.25ex}{\accentmark{15}{\char"1CD2}}ज\raisebox{0.25ex}{\accentmark{12}{\char"0951}}सोविमा\raisebox{0.25ex}{\accentmark{15}{\char"1CD2}}नः\raisebox{0.25ex}{\accentmark{12}{\char"0951}}\kern 0.05em ।\hskip 0.12em plus 0.04em minus 0.02em आज\raisebox{0.25ex}{\accentmark{12}{\char"0951}}सृञ्ज्यो\raisebox{0.15ex}{\accentmark{12}{\char"1CF9}}तिर्हा\raisebox{0.25ex}{\accentmark{12}{\char"0951}}वीरास्मिसा\raisebox{0.25ex}{\accentmark{15}{\char"1CD2}}र्व\raisebox{0.25ex}{\accentmark{12}{\char"0951}}म् \nolinebreak\hskip 0.08em \mbox{॥७॥}}} \par}

\vspace{0.8em}
                % --- SUBSECTIONS (Mantras) ---
                
                % Conditional rendering based on output_mode
                \par\filbreak

\phantomsection

\addcontentsline{toc}{subsection}{ऋक् ८}

{\centering \textbf{\textcolor{AccentBlue}{पात्य\raisebox{0.15ex}{\accentmark{12}{\char"1CF9}}ग्नीर्वीपो\raisebox{0.15ex}{\accentmark{12}{\char"1CF9}}अ\raisebox{0.15ex}{\accentmark{12}{\char"1CF9}}ग्रं\raisebox{0.25ex}{\accentmark{12}{\char"0951}}पादं\raisebox{0.25ex}{\accentmark{15}{\char"1CD2}}वोः\kern 0.05em ।\hskip 0.12em plus 0.04em minus 0.02em पातीयह्न\raisebox{0.15ex}{\accentmark{12}{\char"1CF9}}श्चा\raisebox{0.15ex}{\accentmark{12}{\char"1CF9}}राणंसू\raisebox{0.25ex}{\accentmark{15}{\char"1CD2}}र्या\raisebox{0.25ex}{\accentmark{12}{\char"0951}}स्या\kern 0.05em ।\hskip 0.12em plus 0.04em minus 0.02em पा\raisebox{0.25ex}{\accentmark{15}{\char"1CD2}}तीना\raisebox{0.15ex}{\accentmark{12}{\char"1CF9}}भा\raisebox{0.25ex}{\accentmark{12}{\char"0951}}सप्ता\raisebox{0.25ex}{\accentmark{15}{\char"1CD2}}शी\raisebox{0.25ex}{\accentmark{12}{\char"0951}}र्षाणामग्नीः\raisebox{0.25ex}{\accentmark{15}{\char"1CD2}}\kern 0.05em ।\hskip 0.12em plus 0.04em minus 0.02em पा\raisebox{0.25ex}{\accentmark{15}{\char"1CD2}}तीर्देवा\raisebox{0.25ex}{\accentmark{15}{\char"1CD2}}ना\raisebox{0.25ex}{\accentmark{12}{\char"0951}}मूपामा\raisebox{0.15ex}{\accentmark{12}{\char"1CF8}}द\raisebox{0.25ex}{\accentmark{12}{\char"0951}}मृष्वा\raisebox{0.25ex}{\accentmark{15}{\char"1CD2}} \nolinebreak\hskip 0.08em \mbox{॥८॥}}} \par}

\vspace{0.8em}
                % --- SUBSECTIONS (Mantras) ---
                
                % Conditional rendering based on output_mode
                \par\filbreak

\phantomsection

\addcontentsline{toc}{subsection}{ऋक् ९}

{\centering \textbf{\textcolor{AccentBlue}{भ्रा\raisebox{0.25ex}{\accentmark{15}{\char"1CD2}}ज\raisebox{0.25ex}{\accentmark{12}{\char"0951}}न्त्यग्नेसमीधानादी\raisebox{0.25ex}{\accentmark{15}{\char"1CD2}}दीवः\raisebox{0.25ex}{\accentmark{15}{\char"1CD2}}\kern 0.05em ।\hskip 0.12em plus 0.04em minus 0.02em जिह्वा\raisebox{0.25ex}{\accentmark{15}{\char"1CD2}}चारत्यन्ता\raisebox{0.25ex}{\accentmark{15}{\char"1CD2}}रासा\raisebox{0.25ex}{\accentmark{15}{\char"1CD2}}नी\kern 0.05em ।\hskip 0.12em plus 0.04em minus 0.02em सत्व\raisebox{0.25ex}{\accentmark{15}{\char"1CD2}}न्नो\raisebox{0.25ex}{\accentmark{12}{\char"0951}}अग्नेपा\raisebox{0.25ex}{\accentmark{15}{\char"1CD2}}यासावासूवी\raisebox{0.25ex}{\accentmark{15}{\char"1CD2}}त्\kern 0.05em ।\hskip 0.12em plus 0.04em minus 0.02em रायिंवर्चोंदृशे\raisebox{0.25ex}{\accentmark{15}{\char"1CD2}}दाः\raisebox{0.25ex}{\accentmark{12}{\char"0951}} \nolinebreak\hskip 0.08em \mbox{॥९॥}}} \par}

\vspace{0.8em}

            % --- End of Kandah Closing Floral Line & Counts ---
            \par\nopagebreak\vspace{0.9em plus 0.3em minus 0.1em}
            \centerline{\textcolor{black!35}{\rule{1.8cm}{0.35pt}}\;\raisebox{-0.22ex}{\textcolor{vedicmaroon}{\large\SixFlowerPetalDotted}}\;\textcolor{black!35}{\rule{1.8cm}{0.35pt}}}
            \nopagebreak\vspace{0.45em}
            \centerline{\devafont\textbf{\textcolor{AccentBlue}{॥ ऋक् सङ्ख्या : ९ ॥}}}
        
            \clearpage

            % --- SECTION ---
            
            \phantomsection  

            \section*{\centering अथ तृतीयः खण्डः}
            \addcontentsline{toc}{section}{तृतीयः खण्डः} 
            \markright{तृतीयः खण्डः}
            \nopagebreak % Keep section title with first mantra
            
            % --- SECTION-LEVEL PROCEDURE LINK ---

                % --- SUBSECTIONS (Mantras) ---
                
                % Conditional rendering based on output_mode
                \par\filbreak

\phantomsection

\addcontentsline{toc}{subsection}{ऋक् १}

{\centering \textbf{\textcolor{AccentBlue}{इ\raisebox{0.15ex}{\accentmark{12}{\char"1CF8}}न्द्र\raisebox{0.25ex}{\accentmark{12}{\char"0951}}स्यानूवीर्या\raisebox{0.25ex}{\accentmark{12}{\char"0951}}णिप्रावो\raisebox{0.25ex}{\accentmark{12}{\char"0951}}चं\kern 0.05em ।\hskip 0.12em plus 0.04em minus 0.02em या\raisebox{0.15ex}{\accentmark{12}{\char"1CF8}}नी\raisebox{0.25ex}{\accentmark{12}{\char"0951}}चाका\raisebox{0.25ex}{\accentmark{15}{\char"1CD2}}राप्राथामा\raisebox{0.25ex}{\accentmark{12}{\char"0951}}नीवज्रीः\raisebox{0.25ex}{\accentmark{15}{\char"1CD2}}\kern 0.05em ।\hskip 0.12em plus 0.04em minus 0.02em आ\raisebox{0.25ex}{\accentmark{15}{\char"1CD2}}हन्ना\raisebox{0.25ex}{\accentmark{15}{\char"1CD2}}हीमन्वापा\raisebox{0.25ex}{\accentmark{15}{\char"1CD2}}स्ता\raisebox{0.25ex}{\accentmark{12}{\char"0951}}तर्दा\kern 0.05em ।\hskip 0.12em plus 0.04em minus 0.02em प्रा\raisebox{0.25ex}{\accentmark{15}{\char"1CD2}}वक्षा\raisebox{0.25ex}{\accentmark{12}{\char"0951}}णाभीनत्पार्वा\raisebox{0.25ex}{\accentmark{12}{\char"0951}}तानाम् \nolinebreak\hskip 0.08em \mbox{॥१॥}}} \par}

\vspace{0.8em}
                % --- SUBSECTIONS (Mantras) ---
                
                % Conditional rendering based on output_mode
                \par\filbreak

\phantomsection

\addcontentsline{toc}{subsection}{ऋक् २}

{\centering \textbf{\textcolor{AccentBlue}{बिभ्राट्बृ\raisebox{0.25ex}{\accentmark{12}{\char"0951}}हा\raisebox{0.25ex}{\accentmark{15}{\char"1CD2}}त्पीबतूसो\raisebox{0.25ex}{\accentmark{15}{\char"1CD2}}म्यम्माधू\kern 0.05em ।\hskip 0.12em plus 0.04em minus 0.02em आयूर्दध\raisebox{0.25ex}{\accentmark{12}{\char"0951}}द्यज्ञा\raisebox{0.15ex}{\accentmark{12}{\char"1CF9}}पा\raisebox{0.25ex}{\accentmark{12}{\char"0951}}तावाविहृतम्\kern 0.05em ।\hskip 0.12em plus 0.04em minus 0.02em वा\raisebox{0.25ex}{\accentmark{15}{\char"1CD2}}ताजूतोयोआ\raisebox{0.25ex}{\accentmark{12}{\char"0951}}भीराक्षा\raisebox{0.25ex}{\accentmark{12}{\char"0951}}तीत्मा\raisebox{0.25ex}{\accentmark{15}{\char"1CD2}}ना\kern 0.05em ।\hskip 0.12em plus 0.04em minus 0.02em प्राजाः\raisebox{0.25ex}{\accentmark{15}{\char"1CD2}}पिबर्तीबहुधा\raisebox{0.25ex}{\accentmark{15}{\char"1CD2}}वी\raisebox{0.25ex}{\accentmark{15}{\char"1CD2}}रा\raisebox{0.25ex}{\accentmark{12}{\char"0951}}ज\raisebox{0.15ex}{\accentmark{12}{\char"1CF9}}ति \nolinebreak\hskip 0.08em \mbox{॥२॥}}} \par}

\vspace{0.8em}
                % --- SUBSECTIONS (Mantras) ---
                
                % Conditional rendering based on output_mode
                \par\filbreak

\phantomsection

\addcontentsline{toc}{subsection}{ऋक् ३}

{\centering \textbf{\textcolor{AccentBlue}{वा\raisebox{0.25ex}{\accentmark{15}{\char"1CD2}}स\raisebox{0.25ex}{\accentmark{12}{\char"0951}}न्तइ\raisebox{0.15ex}{\accentmark{12}{\char"1CF8}}न्नुरन्याः\raisebox{0.15ex}{\accentmark{12}{\char"1CF9}}\kern 0.05em ।\hskip 0.12em plus 0.04em minus 0.02em ग्रीष्मइन्नु\raisebox{0.25ex}{\accentmark{15}{\char"1CD2}}र\raisebox{0.25ex}{\accentmark{15}{\char"1CD2}}त्न्याः\raisebox{0.25ex}{\accentmark{12}{\char"0951}}\kern 0.05em ।\hskip 0.12em plus 0.04em minus 0.02em वा\raisebox{0.25ex}{\accentmark{15}{\char"1CD2}}र्षाणइन्नुशारा\raisebox{0.25ex}{\accentmark{15}{\char"1CD2}}दाः\raisebox{0.25ex}{\accentmark{12}{\char"0951}}\kern 0.05em ।\hskip 0.12em plus 0.04em minus 0.02em हेमन्ता\raisebox{0.15ex}{\accentmark{12}{\char"1CF9}}शी\raisebox{0.15ex}{\accentmark{12}{\char"1CF9}}शीरइन्नुरन्त्याः\raisebox{0.25ex}{\accentmark{12}{\char"0951}} \nolinebreak\hskip 0.08em \mbox{॥३॥}}} \par}

\vspace{0.8em}
                % --- SUBSECTIONS (Mantras) ---
                
                % Conditional rendering based on output_mode
                \par\filbreak

\phantomsection

\addcontentsline{toc}{subsection}{ऋक् ४}

{\centering \textbf{\textcolor{AccentBlue}{ता\raisebox{0.25ex}{\accentmark{15}{\char"1CD2}}वत्य\raisebox{0.15ex}{\accentmark{12}{\char"1CF9}}पीतोददतः\kern 0.05em ।\hskip 0.12em plus 0.04em minus 0.02em तावास्वादिष्ठाते\raisebox{0.25ex}{\accentmark{15}{\char"1CD2}}पीतो\kern 0.05em ।\hskip 0.12em plus 0.04em minus 0.02em प्रस्वद्वा\raisebox{0.25ex}{\accentmark{15}{\char"1CD2}}नोरासा\raisebox{0.25ex}{\accentmark{12}{\char"0951}}नां\kern 0.05em ।\hskip 0.12em plus 0.04em minus 0.02em तूविग्री\raisebox{0.25ex}{\accentmark{15}{\char"1CD2}}वा\raisebox{0.25ex}{\accentmark{12}{\char"0951}}इवेरते \nolinebreak\hskip 0.08em \mbox{॥४॥}}} \par}

\vspace{0.8em}
                % --- SUBSECTIONS (Mantras) ---
                
                % Conditional rendering based on output_mode
                \par\filbreak

\phantomsection

\addcontentsline{toc}{subsection}{ऋक् ५}

{\centering \textbf{\textcolor{AccentBlue}{इन्द्र\raisebox{0.15ex}{\accentmark{12}{\char"1CF9}}स्फाढूतावृ\raisebox{0.25ex}{\accentmark{12}{\char"0951}}त्राहा\raisebox{0.25ex}{\accentmark{15}{\char"1CD2}}\kern 0.05em ।\hskip 0.12em plus 0.04em minus 0.02em परस्पा\raisebox{0.25ex}{\accentmark{15}{\char"1CD2}}नोवा\raisebox{0.25ex}{\accentmark{15}{\char"1CD2}}रेण्यः\kern 0.05em ।\hskip 0.12em plus 0.04em minus 0.02em सा\raisebox{0.25ex}{\accentmark{15}{\char"1CD2}}नो\raisebox{0.25ex}{\accentmark{12}{\char"0951}}र\raisebox{0.15ex}{\accentmark{12}{\char"1CF9}}क्षिषच्चार\raisebox{0.15ex}{\accentmark{12}{\char"1CF9}}मंसा\raisebox{0.15ex}{\accentmark{12}{\char"1CF9}}मा\raisebox{0.25ex}{\accentmark{15}{\char"1CD2}}ध्यामं\kern 0.05em ।\hskip 0.12em plus 0.04em minus 0.02em सा\raisebox{0.25ex}{\accentmark{15}{\char"1CD2}}पश्चात्पातूसा\raisebox{0.25ex}{\accentmark{15}{\char"1CD2}}पू\raisebox{0.25ex}{\accentmark{12}{\char"0951}}राः \nolinebreak\hskip 0.08em \mbox{॥५॥}}} \par}

\vspace{0.8em}
                % --- SUBSECTIONS (Mantras) ---
                
                % Conditional rendering based on output_mode
                \par\filbreak

\phantomsection

\addcontentsline{toc}{subsection}{ऋक् ६}

{\centering \textbf{\textcolor{AccentBlue}{सह\raisebox{0.25ex}{\accentmark{15}{\char"1CD2}}स्राशीर्षापूरू\raisebox{0.25ex}{\accentmark{12}{\char"0951}}षः\kern 0.05em ।\hskip 0.12em plus 0.04em minus 0.02em सहस्राक्षा\raisebox{0.25ex}{\accentmark{15}{\char"1CD2}}स्साहस्रा\raisebox{0.25ex}{\accentmark{15}{\char"1CD2}}पात्\kern 0.05em ।\hskip 0.12em plus 0.04em minus 0.02em सा\raisebox{0.15ex}{\accentmark{12}{\char"1CF8}}भू\raisebox{0.25ex}{\accentmark{12}{\char"0951}}मिंस\raisebox{0.25ex}{\accentmark{15}{\char"1CD2}}र्वातो\raisebox{0.25ex}{\accentmark{12}{\char"0951}}वृत्वा\kern 0.05em ।\hskip 0.12em plus 0.04em minus 0.02em अत्यातिष्ठद्दशांगूलम् \nolinebreak\hskip 0.08em \mbox{॥६॥}}} \par}

\vspace{0.8em}
                % --- SUBSECTIONS (Mantras) ---
                
                % Conditional rendering based on output_mode
                \par\filbreak

\phantomsection

\addcontentsline{toc}{subsection}{ऋक् ७}

{\centering \textbf{\textcolor{AccentBlue}{त्रीपा\raisebox{0.25ex}{\accentmark{15}{\char"1CD2}}दूर्ध्वा\raisebox{0.25ex}{\accentmark{12}{\char"0951}}ऊदू\raisebox{0.25ex}{\accentmark{15}{\char"1CD2}}यत्पूरू\raisebox{0.25ex}{\accentmark{15}{\char"1CD2}}षः\kern 0.05em ।\hskip 0.12em plus 0.04em minus 0.02em पादो\raisebox{0.25ex}{\accentmark{12}{\char"0951}}स्येहा\raisebox{0.15ex}{\accentmark{12}{\char"1CF9}}भा\raisebox{0.25ex}{\accentmark{12}{\char"0951}}वत्पूनाः\kern 0.05em ।\hskip 0.12em plus 0.04em minus 0.02em ता\raisebox{0.25ex}{\accentmark{15}{\char"1CD2}}थाविष्वंव्या१क्रमात्\kern 0.05em ।\hskip 0.12em plus 0.04em minus 0.02em अ\raisebox{0.15ex}{\accentmark{12}{\char"1CF9}}श\raisebox{0.25ex}{\accentmark{12}{\char"0951}}ना\raisebox{0.25ex}{\accentmark{12}{\char"0951}}नशने\raisebox{0.25ex}{\accentmark{15}{\char"1CD2}}आभि \nolinebreak\hskip 0.08em \mbox{॥७॥}}} \par}

\vspace{0.8em}
                % --- SUBSECTIONS (Mantras) ---
                
                % Conditional rendering based on output_mode
                \par\filbreak

\phantomsection

\addcontentsline{toc}{subsection}{ऋक् ८}

{\centering \textbf{\textcolor{AccentBlue}{पूरू\raisebox{0.25ex}{\accentmark{12}{\char"0951}}षाए\raisebox{0.25ex}{\accentmark{15}{\char"1CD2}}वेदंस\raisebox{0.25ex}{\accentmark{15}{\char"1CD2}}र्वं\raisebox{0.25ex}{\accentmark{12}{\char"0951}}\kern 0.05em ।\hskip 0.12em plus 0.04em minus 0.02em य\raisebox{0.25ex}{\accentmark{15}{\char"1CD2}}त्भूतंया\raisebox{0.25ex}{\accentmark{15}{\char"1CD2}}च्चाभा\raisebox{0.25ex}{\accentmark{15}{\char"1CD2}}व्यं\raisebox{0.25ex}{\accentmark{12}{\char"0951}}\kern 0.05em ।\hskip 0.12em plus 0.04em minus 0.02em पादो\raisebox{0.25ex}{\accentmark{12}{\char"0951}}स्यास\raisebox{0.15ex}{\accentmark{12}{\char"1CF9}}र्वा\raisebox{0.25ex}{\accentmark{12}{\char"0951}}भूतानी\kern 0.05em ।\hskip 0.12em plus 0.04em minus 0.02em त्रीपाद\raisebox{0.25ex}{\accentmark{12}{\char"0951}}स्यामृ\raisebox{0.15ex}{\accentmark{12}{\char"1CF9}}त\raisebox{0.25ex}{\accentmark{12}{\char"0951}}न्दीवी\raisebox{0.25ex}{\accentmark{15}{\char"1CD2}} \nolinebreak\hskip 0.08em \mbox{॥८॥}}} \par}

\vspace{0.8em}
                % --- SUBSECTIONS (Mantras) ---
                
                % Conditional rendering based on output_mode
                \par\filbreak

\phantomsection

\addcontentsline{toc}{subsection}{ऋक् ९}

{\centering \textbf{\textcolor{AccentBlue}{ता\raisebox{0.25ex}{\accentmark{15}{\char"1CD2}}वा\raisebox{0.25ex}{\accentmark{12}{\char"0951}}नस्यामाहीमा\raisebox{0.25ex}{\accentmark{15}{\char"1CD2}}\kern 0.05em ।\hskip 0.12em plus 0.04em minus 0.02em ता\raisebox{0.25ex}{\accentmark{15}{\char"1CD2}}तो\raisebox{0.15ex}{\accentmark{12}{\char"1CF8}}ज्यायांश्चापू\raisebox{0.25ex}{\accentmark{15}{\char"1CD2}}रु\raisebox{0.25ex}{\accentmark{12}{\char"0951}}षः\kern 0.05em ।\hskip 0.12em plus 0.04em minus 0.02em ऊता\raisebox{0.15ex}{\accentmark{12}{\char"1CF9}}मृ\raisebox{0.25ex}{\accentmark{12}{\char"0951}}तत्वस्येशा\raisebox{0.25ex}{\accentmark{12}{\char"0951}}नः\kern 0.05em ।\hskip 0.12em plus 0.04em minus 0.02em याद\raisebox{0.25ex}{\accentmark{12}{\char"0951}}न्नेनातीरो\raisebox{0.15ex}{\accentmark{12}{\char"1CF9}}हा\raisebox{0.25ex}{\accentmark{12}{\char"0951}}ति \nolinebreak\hskip 0.08em \mbox{॥९॥}}} \par}

\vspace{0.8em}
                % --- SUBSECTIONS (Mantras) ---
                
                % Conditional rendering based on output_mode
                \par\filbreak

\phantomsection

\addcontentsline{toc}{subsection}{ऋक् १०}

{\centering \textbf{\textcolor{AccentBlue}{ता\raisebox{0.15ex}{\accentmark{12}{\char"1CF8}}तो\raisebox{0.25ex}{\accentmark{12}{\char"0951}}वीराड\raisebox{0.25ex}{\accentmark{15}{\char"1CD2}}जा\raisebox{0.25ex}{\accentmark{15}{\char"1CD2}}यत\kern 0.05em ।\hskip 0.12em plus 0.04em minus 0.02em वीरा\raisebox{0.25ex}{\accentmark{15}{\char"1CD2}}जोआ\raisebox{0.25ex}{\accentmark{15}{\char"1CD2}}धीपू\raisebox{0.25ex}{\accentmark{15}{\char"1CD2}}रु\raisebox{0.25ex}{\accentmark{12}{\char"0951}}षः\kern 0.05em ।\hskip 0.12em plus 0.04em minus 0.02em सा\raisebox{0.25ex}{\accentmark{15}{\char"1CD2}}जातो\raisebox{0.25ex}{\accentmark{15}{\char"1CD2}}अत्या\raisebox{0.25ex}{\accentmark{12}{\char"0951}}रिच्यत\kern 0.05em ।\hskip 0.12em plus 0.04em minus 0.02em पश्चात्भू\raisebox{0.25ex}{\accentmark{15}{\char"1CD2}}मि\raisebox{0.15ex}{\accentmark{12}{\char"1CF8}}मा\raisebox{0.15ex}{\accentmark{12}{\char"1CF8}}थो\raisebox{0.25ex}{\accentmark{12}{\char"0951}}पूरः\raisebox{0.25ex}{\accentmark{15}{\char"1CD2}} \nolinebreak\hskip 0.08em \mbox{॥१०॥}}} \par}

\vspace{0.8em}

            % --- End of Kandah Closing Floral Line & Counts ---
            \par\nopagebreak\vspace{0.9em plus 0.3em minus 0.1em}
            \centerline{\textcolor{black!35}{\rule{1.8cm}{0.35pt}}\;\raisebox{-0.22ex}{\textcolor{vedicmaroon}{\large\SixFlowerPetalDotted}}\;\textcolor{black!35}{\rule{1.8cm}{0.35pt}}}
            \nopagebreak\vspace{0.45em}
            \centerline{\devafont\textbf{\textcolor{AccentBlue}{॥ ऋक् सङ्ख्या : १० ॥}}}
        
            \clearpage

            % --- SECTION ---
            
            \phantomsection  

            \section*{\centering अथ चतुर्थः खण्डः}
            \addcontentsline{toc}{section}{चतुर्थः खण्डः} 
            \markright{चतुर्थः खण्डः}
            \nopagebreak % Keep section title with first mantra
            
            % --- SECTION-LEVEL PROCEDURE LINK ---

                % --- SUBSECTIONS (Mantras) ---
                
                % Conditional rendering based on output_mode
                \par\filbreak

\phantomsection

\addcontentsline{toc}{subsection}{ऋक् १}

{\centering \textbf{\textcolor{AccentBlue}{य\raisebox{0.25ex}{\accentmark{15}{\char"1CD2}}त्पू\raisebox{0.25ex}{\accentmark{15}{\char"1CD2}}रूषेणाहावीषाः\raisebox{0.25ex}{\accentmark{15}{\char"1CD2}}\kern 0.05em ।\hskip 0.12em plus 0.04em minus 0.02em देवा\raisebox{0.25ex}{\accentmark{15}{\char"1CD2}}यज्ञामा\raisebox{0.25ex}{\accentmark{15}{\char"1CD2}}त\raisebox{0.25ex}{\accentmark{12}{\char"0951}}न्वाताः\kern 0.05em ।\hskip 0.12em plus 0.04em minus 0.02em वस\raisebox{0.25ex}{\accentmark{12}{\char"0951}}न्ताये\raisebox{0.25ex}{\accentmark{12}{\char"0951}}षामासीदा\raisebox{0.25ex}{\accentmark{15}{\char"1CD2}}ज्यं\raisebox{0.25ex}{\accentmark{12}{\char"0951}}\kern 0.05em ।\hskip 0.12em plus 0.04em minus 0.02em ग्रीष्मा\raisebox{0.25ex}{\accentmark{15}{\char"1CD2}}इध्माः\raisebox{0.15ex}{\accentmark{12}{\char"1CF9}}शा\raisebox{0.15ex}{\accentmark{12}{\char"1CF8}}रा\raisebox{0.25ex}{\accentmark{12}{\char"0951}}द्धाविः \nolinebreak\hskip 0.08em \mbox{॥१॥}}} \par}

\vspace{0.8em}
                % --- SUBSECTIONS (Mantras) ---
                
                % Conditional rendering based on output_mode
                \par\filbreak

\phantomsection

\addcontentsline{toc}{subsection}{ऋक् २}

{\centering \textbf{\textcolor{AccentBlue}{सप्ता\raisebox{0.25ex}{\accentmark{15}{\char"1CD2}}स्या\raisebox{0.25ex}{\accentmark{12}{\char"0951}}सन्पारीधा\raisebox{0.25ex}{\accentmark{15}{\char"1CD2}}याः\kern 0.05em ।\hskip 0.12em plus 0.04em minus 0.02em त्रीस्सप्त\raisebox{0.25ex}{\accentmark{15}{\char"1CD2}}स्समीधाः\raisebox{0.25ex}{\accentmark{15}{\char"1CD2}}कृ\raisebox{0.25ex}{\accentmark{12}{\char"0951}}ताः\kern 0.05em ।\hskip 0.12em plus 0.04em minus 0.02em देवायद्य\raisebox{0.25ex}{\accentmark{15}{\char"1CD2}}ज्ञ\raisebox{0.25ex}{\accentmark{15}{\char"1CD2}}न्तन्वानः\raisebox{0.25ex}{\accentmark{15}{\char"1CD2}}\kern 0.05em ।\hskip 0.12em plus 0.04em minus 0.02em आब\raisebox{0.25ex}{\accentmark{12}{\char"0951}}ध्नन्पू\raisebox{0.25ex}{\accentmark{15}{\char"1CD2}}रू\raisebox{0.25ex}{\accentmark{12}{\char"0951}}षंहाविः\raisebox{0.25ex}{\accentmark{15}{\char"1CD2}} \nolinebreak\hskip 0.08em \mbox{॥२॥}}} \par}

\vspace{0.8em}
                % --- SUBSECTIONS (Mantras) ---
                
                % Conditional rendering based on output_mode
                \par\filbreak

\phantomsection

\addcontentsline{toc}{subsection}{ऋक् ३}

{\centering \textbf{\textcolor{AccentBlue}{हा\raisebox{0.25ex}{\accentmark{15}{\char"1CD2}}री\raisebox{0.25ex}{\accentmark{12}{\char"0951}}तइन्द्रश्माश्रू\raisebox{0.25ex}{\accentmark{15}{\char"1CD2}}णीः\kern 0.05em ।\hskip 0.12em plus 0.04em minus 0.02em ऊतो\raisebox{0.15ex}{\accentmark{12}{\char"1CF9}}ते\raisebox{0.25ex}{\accentmark{12}{\char"0951}}हा\raisebox{0.25ex}{\accentmark{12}{\char"0951}}रीताहा\raisebox{0.25ex}{\accentmark{15}{\char"1CD2}}रीः\kern 0.05em ।\hskip 0.12em plus 0.04em minus 0.02em तन्त्वा\raisebox{0.25ex}{\accentmark{12}{\char"0951}}स्तुवन्तीकारा\raisebox{0.25ex}{\accentmark{12}{\char"0951}}वाः\kern 0.05em ।\hskip 0.12em plus 0.04em minus 0.02em वार्षा\raisebox{0.15ex}{\accentmark{12}{\char"1CF9}}सोवा\raisebox{0.25ex}{\accentmark{12}{\char"0951}}नार्गा\raisebox{0.25ex}{\accentmark{12}{\char"0951}}वः \nolinebreak\hskip 0.08em \mbox{॥३॥}}} \par}

\vspace{0.8em}
                % --- SUBSECTIONS (Mantras) ---
                
                % Conditional rendering based on output_mode
                \par\filbreak

\phantomsection

\addcontentsline{toc}{subsection}{ऋक् ४}

{\centering \textbf{\textcolor{AccentBlue}{ऋ\raisebox{0.25ex}{\accentmark{15}{\char"1CD2}}ष्या\raisebox{0.25ex}{\accentmark{12}{\char"0951}}सइन्द्रभूं\raisebox{0.25ex}{\accentmark{15}{\char"1CD2}}ईति\kern 0.05em ।\hskip 0.12em plus 0.04em minus 0.02em मा\raisebox{0.25ex}{\accentmark{15}{\char"1CD2}}घा\raisebox{0.25ex}{\accentmark{12}{\char"0951}}वन्विन्द्राभूंई\raisebox{0.25ex}{\accentmark{15}{\char"1CD2}}ति\kern 0.05em ।\hskip 0.12em plus 0.04em minus 0.02em इ\raisebox{0.25ex}{\accentmark{15}{\char"1CD2}}न्द्र\raisebox{0.25ex}{\accentmark{12}{\char"0951}}स्ता\raisebox{0.25ex}{\accentmark{15}{\char"1CD2}}ता\raisebox{0.25ex}{\accentmark{12}{\char"0951}}रपूतः \nolinebreak\hskip 0.08em \mbox{॥४॥}}} \par}

\vspace{0.8em}
                % --- SUBSECTIONS (Mantras) ---
                
                % Conditional rendering based on output_mode
                \par\filbreak

\phantomsection

\addcontentsline{toc}{subsection}{ऋक् ५}

{\centering \textbf{\textcolor{AccentBlue}{य\raisebox{0.25ex}{\accentmark{15}{\char"1CD2}}द्वार्चो\raisebox{0.25ex}{\accentmark{12}{\char"0951}}हीरण्य\raisebox{0.25ex}{\accentmark{15}{\char"1CD2}}स्य\raisebox{0.25ex}{\accentmark{12}{\char"0951}}\kern 0.05em ।\hskip 0.12em plus 0.04em minus 0.02em यद्वा\raisebox{0.25ex}{\accentmark{15}{\char"1CD2}}व\raisebox{0.25ex}{\accentmark{15}{\char"1CD2}}र्चो\raisebox{0.25ex}{\accentmark{12}{\char"0951}}ग\raisebox{0.15ex}{\accentmark{12}{\char"1CF9}}वा\raisebox{0.25ex}{\accentmark{12}{\char"0951}}मूता\raisebox{0.25ex}{\accentmark{12}{\char"0951}}\kern 0.05em ।\hskip 0.12em plus 0.04em minus 0.02em सत्या\raisebox{0.25ex}{\accentmark{15}{\char"1CD2}}स्या\raisebox{0.25ex}{\accentmark{12}{\char"0951}}ब्र\raisebox{0.15ex}{\accentmark{12}{\char"1CF9}}ह्माणोव\raisebox{0.25ex}{\accentmark{15}{\char"1CD2}}र्चाः\raisebox{0.25ex}{\accentmark{12}{\char"0951}}\kern 0.05em ।\hskip 0.12em plus 0.04em minus 0.02em ते\raisebox{0.25ex}{\accentmark{15}{\char"1CD2}}ना\raisebox{0.25ex}{\accentmark{12}{\char"0951}}मसंसृजाम\raisebox{0.25ex}{\accentmark{12}{\char"0951}}सी \nolinebreak\hskip 0.08em \mbox{॥५॥}}} \par}

\vspace{0.8em}
                % --- SUBSECTIONS (Mantras) ---
                
                % Conditional rendering based on output_mode
                \par\filbreak

\phantomsection

\addcontentsline{toc}{subsection}{ऋक् ६}

{\centering \textbf{\textcolor{AccentBlue}{चित्रन्देवा\raisebox{0.25ex}{\accentmark{15}{\char"1CD2}}नामूद\raisebox{0.25ex}{\accentmark{12}{\char"0951}}कादनी\raisebox{0.25ex}{\accentmark{12}{\char"0951}}कं\kern 0.05em ।\hskip 0.12em plus 0.04em minus 0.02em च\raisebox{0.15ex}{\accentmark{12}{\char"1CF8}}क्षुर्मित्र\raisebox{0.25ex}{\accentmark{15}{\char"1CD2}}स्यावा\raisebox{0.25ex}{\accentmark{15}{\char"1CD2}}रूणास्याग्नेः\raisebox{0.25ex}{\accentmark{15}{\char"1CD2}}\kern 0.05em ।\hskip 0.12em plus 0.04em minus 0.02em आप्राद्या\raisebox{0.25ex}{\accentmark{15}{\char"1CD2}}वापृथीवी\raisebox{0.25ex}{\accentmark{12}{\char"0951}}अन्तारी\raisebox{0.25ex}{\accentmark{12}{\char"0951}}क्षं\kern 0.05em ।\hskip 0.12em plus 0.04em minus 0.02em सू\raisebox{0.15ex}{\accentmark{12}{\char"1CF8}}र्यात्माजग\raisebox{0.25ex}{\accentmark{12}{\char"0951}}तास्तस्थूषा\raisebox{0.25ex}{\accentmark{12}{\char"0951}}श्च \nolinebreak\hskip 0.08em \mbox{॥६॥}}} \par}

\vspace{0.8em}
                % --- SUBSECTIONS (Mantras) ---
                
                % Conditional rendering based on output_mode
                \par\filbreak

\phantomsection

\addcontentsline{toc}{subsection}{ऋक् ७}

{\centering \textbf{\textcolor{AccentBlue}{या\raisebox{0.25ex}{\accentmark{15}{\char"1CD2}}दीत\raisebox{0.15ex}{\accentmark{12}{\char"1CF8}}स्त\raisebox{0.25ex}{\accentmark{15}{\char"1CD2}}न्वोमामा\raisebox{0.25ex}{\accentmark{12}{\char"0951}}\kern 0.05em ।\hskip 0.12em plus 0.04em minus 0.02em दोषा\raisebox{0.25ex}{\accentmark{12}{\char"0951}}रासास्यभेजिरे\raisebox{0.25ex}{\accentmark{12}{\char"0951}}\kern 0.05em ।\hskip 0.12em plus 0.04em minus 0.02em ना\raisebox{0.25ex}{\accentmark{15}{\char"1CD2}}राशं\raisebox{0.15ex}{\accentmark{12}{\char"1CF9}}सेनासोमे\raisebox{0.25ex}{\accentmark{12}{\char"0951}}ना\kern 0.05em ।\hskip 0.12em plus 0.04em minus 0.02em आ\raisebox{0.25ex}{\accentmark{15}{\char"1CD2}}हन्त\raisebox{0.25ex}{\accentmark{15}{\char"1CD2}}त्पू\raisebox{0.25ex}{\accentmark{15}{\char"1CD2}}नारा\raisebox{0.15ex}{\accentmark{12}{\char"1CF8}}दे\raisebox{0.25ex}{\accentmark{12}{\char"0951}}दे \nolinebreak\hskip 0.08em \mbox{॥७॥}}} \par}

\vspace{0.8em}
                % --- SUBSECTIONS (Mantras) ---
                
                % Conditional rendering based on output_mode
                \par\filbreak

\phantomsection

\addcontentsline{toc}{subsection}{ऋक् ८}

{\centering \textbf{\textcolor{AccentBlue}{यूष्मा\raisebox{0.25ex}{\accentmark{15}{\char"1CD2}}दस्पारासस्पारी\kern 0.05em ।\hskip 0.12em plus 0.04em minus 0.02em यस्ये\raisebox{0.25ex}{\accentmark{12}{\char"0951}}दमारा\raisebox{0.25ex}{\accentmark{12}{\char"0951}}जोयू\raisebox{0.25ex}{\accentmark{15}{\char"1CD2}}जाः\raisebox{0.25ex}{\accentmark{12}{\char"0951}}\kern 0.05em ।\hskip 0.12em plus 0.04em minus 0.02em तुजे\raisebox{0.25ex}{\accentmark{12}{\char"0951}}ज\raisebox{0.25ex}{\accentmark{15}{\char"1CD2}}ने\raisebox{0.25ex}{\accentmark{12}{\char"0951}}वानंस्वाः१\kern 0.05em ।\hskip 0.12em plus 0.04em minus 0.02em इन्द्र\raisebox{0.25ex}{\accentmark{12}{\char"0951}}स्यार\raisebox{0.15ex}{\accentmark{12}{\char"1CF9}}न्त्यं\raisebox{0.25ex}{\accentmark{12}{\char"0951}}बृहा\raisebox{0.25ex}{\accentmark{15}{\char"1CD2}}त् \nolinebreak\hskip 0.08em \mbox{॥८॥}}} \par}

\vspace{0.8em}
                % --- SUBSECTIONS (Mantras) ---
                
                % Conditional rendering based on output_mode
                \par\filbreak

\phantomsection

\addcontentsline{toc}{subsection}{ऋक् ९}

{\centering \textbf{\textcolor{AccentBlue}{सा\raisebox{0.25ex}{\accentmark{15}{\char"1CD2}}हास्तन्न\raisebox{0.25ex}{\accentmark{12}{\char"0951}}इन्द्रादध्यो\raisebox{0.25ex}{\accentmark{15}{\char"1CD2}}जाः\kern 0.05em ।\hskip 0.12em plus 0.04em minus 0.02em ईशेह्या३स्यामा\raisebox{0.25ex}{\accentmark{12}{\char"0951}}हातोवी\raisebox{0.25ex}{\accentmark{12}{\char"0951}}रप्सी\raisebox{0.25ex}{\accentmark{15}{\char"1CD2}}न्\kern 0.05em ।\hskip 0.12em plus 0.04em minus 0.02em क्रा\raisebox{0.25ex}{\accentmark{15}{\char"1CD2}}तून्नानिर्मणं\raisebox{0.25ex}{\accentmark{15}{\char"1CD2}}स्था\raisebox{0.25ex}{\accentmark{12}{\char"0951}}वीरं\raisebox{0.25ex}{\accentmark{15}{\char"1CD2}}चावा\raisebox{0.25ex}{\accentmark{15}{\char"1CD2}}जं\raisebox{0.25ex}{\accentmark{12}{\char"0951}}\kern 0.05em ।\hskip 0.12em plus 0.04em minus 0.02em वृ\raisebox{0.15ex}{\accentmark{12}{\char"1CF8}}त्रे\raisebox{0.25ex}{\accentmark{15}{\char"1CD2}}षुश\raisebox{0.15ex}{\accentmark{12}{\char"1CF9}}त्रू\raisebox{0.25ex}{\accentmark{12}{\char"0951}}न्सू\raisebox{0.15ex}{\accentmark{12}{\char"1CF9}}ह\raisebox{0.25ex}{\accentmark{12}{\char"0951}}ना\raisebox{0.25ex}{\accentmark{12}{\char"0951}}कृधी\raisebox{0.25ex}{\accentmark{12}{\char"0951}}नः \nolinebreak\hskip 0.08em \mbox{॥९॥}}} \par}

\vspace{0.8em}
                % --- SUBSECTIONS (Mantras) ---
                
                % Conditional rendering based on output_mode
                \par\filbreak

\phantomsection

\addcontentsline{toc}{subsection}{ऋक् १०}

{\centering \textbf{\textcolor{AccentBlue}{इन्द्राज्ये\raisebox{0.15ex}{\accentmark{12}{\char"1CF9}}ष्१न्नआभा\raisebox{0.25ex}{\accentmark{15}{\char"1CD2}}रा\kern 0.05em ।\hskip 0.12em plus 0.04em minus 0.02em ओ\raisebox{0.25ex}{\accentmark{12}{\char"0951}}जीष्ठंपूपूरीश्रा\raisebox{0.25ex}{\accentmark{15}{\char"1CD2}}वाः\raisebox{0.25ex}{\accentmark{12}{\char"0951}}\kern 0.05em ।\hskip 0.12em plus 0.04em minus 0.02em य\raisebox{0.25ex}{\accentmark{15}{\char"1CD2}}द्या\raisebox{0.25ex}{\accentmark{15}{\char"1CD2}}धृ\raisebox{0.25ex}{\accentmark{12}{\char"0951}}क्षेमवज्रहा\raisebox{0.25ex}{\accentmark{15}{\char"1CD2}}स्तारोद\raisebox{0.25ex}{\accentmark{12}{\char"0951}}सि\kern 0.05em ।\hskip 0.12em plus 0.04em minus 0.02em ओभे\raisebox{0.25ex}{\accentmark{15}{\char"1CD2}}सू\raisebox{0.25ex}{\accentmark{12}{\char"0951}}शिप्रवप्राः \nolinebreak\hskip 0.08em \mbox{॥१०॥}}} \par}

\vspace{0.8em}

            % --- End of Kandah Closing Floral Line & Counts ---
            \par\nopagebreak\vspace{0.9em plus 0.3em minus 0.1em}
            \centerline{\textcolor{black!35}{\rule{1.8cm}{0.35pt}}\;\raisebox{-0.22ex}{\textcolor{vedicmaroon}{\large\SixFlowerPetalDotted}}\;\textcolor{black!35}{\rule{1.8cm}{0.35pt}}}
            \nopagebreak\vspace{0.45em}
            \centerline{\devafont\textbf{\textcolor{AccentBlue}{॥ ऋक् सङ्ख्या : १० ॥}}}
        
            \clearpage

            % --- SECTION ---
            
            \phantomsection  

            \section*{\centering अथ पञ्चमः खण्डः}
            \addcontentsline{toc}{section}{पञ्चमः खण्डः} 
            \markright{पञ्चमः खण्डः}
            \nopagebreak % Keep section title with first mantra
            
            % --- SECTION-LEVEL PROCEDURE LINK ---

                % --- SUBSECTIONS (Mantras) ---
                
                % Conditional rendering based on output_mode
                \par\filbreak

\phantomsection

\addcontentsline{toc}{subsection}{ऋक् १}

{\centering \textbf{\textcolor{AccentBlue}{ऊ\raisebox{0.25ex}{\accentmark{15}{\char"1CD2}}दु\raisebox{0.25ex}{\accentmark{12}{\char"0951}}क्तामंवारूणापा\raisebox{0.15ex}{\accentmark{12}{\char"1CF9}}शा\raisebox{0.25ex}{\accentmark{12}{\char"0951}}मस्मात्\kern 0.05em ।\hskip 0.12em plus 0.04em minus 0.02em आबा\raisebox{0.25ex}{\accentmark{12}{\char"0951}}धामं\raisebox{0.15ex}{\accentmark{12}{\char"1CF9}}वीमा\raisebox{0.25ex}{\accentmark{12}{\char"0951}}ध्यामं\raisebox{0.25ex}{\accentmark{15}{\char"1CD2}}श्ना\raisebox{0.25ex}{\accentmark{12}{\char"0951}}थाय\kern 0.05em ।\hskip 0.12em plus 0.04em minus 0.02em आथा\raisebox{0.25ex}{\accentmark{12}{\char"0951}}दित्यव्रातेवा\raisebox{0.25ex}{\accentmark{15}{\char"1CD2}}यन्ता\raisebox{0.25ex}{\accentmark{15}{\char"1CD2}}वा\raisebox{0.25ex}{\accentmark{12}{\char"0951}}\kern 0.05em ।\hskip 0.12em plus 0.04em minus 0.02em आ\raisebox{0.25ex}{\accentmark{15}{\char"1CD2}}नागसोआ\raisebox{0.25ex}{\accentmark{15}{\char"1CD2}}दीतयेस्याम \nolinebreak\hskip 0.08em \mbox{॥१॥}}} \par}

\vspace{0.8em}
                % --- SUBSECTIONS (Mantras) ---
                
                % Conditional rendering based on output_mode
                \par\filbreak

\phantomsection

\addcontentsline{toc}{subsection}{ऋक् २}

{\centering \textbf{\textcolor{AccentBlue}{त्वा\raisebox{0.15ex}{\accentmark{12}{\char"1CF8}}या\raisebox{0.25ex}{\accentmark{12}{\char"0951}}वायंपावा\raisebox{0.25ex}{\accentmark{12}{\char"0951}}मानेनसोम\kern 0.05em ।\hskip 0.12em plus 0.04em minus 0.02em भा\raisebox{0.15ex}{\accentmark{12}{\char"1CF8}}रे\raisebox{0.25ex}{\accentmark{12}{\char"0951}}कृतं\raisebox{0.25ex}{\accentmark{15}{\char"1CD2}}वी\raisebox{0.25ex}{\accentmark{15}{\char"1CD2}}ची\raisebox{0.25ex}{\accentmark{12}{\char"0951}}नुयामाशश्वात्\kern 0.05em ।\hskip 0.12em plus 0.04em minus 0.02em त\raisebox{0.15ex}{\accentmark{12}{\char"1CF8}}न्नो\raisebox{0.25ex}{\accentmark{15}{\char"1CD2}}मित्रोवा\raisebox{0.25ex}{\accentmark{15}{\char"1CD2}}रूणोमामहन्तां\kern 0.05em ।\hskip 0.12em plus 0.04em minus 0.02em आ\raisebox{0.15ex}{\accentmark{12}{\char"1CF8}}दि\raisebox{0.25ex}{\accentmark{12}{\char"0951}}तिस्सिन्धुःपृथीवीमूतद्यौः \nolinebreak\hskip 0.08em \mbox{॥२॥}}} \par}

\vspace{0.8em}
                % --- SUBSECTIONS (Mantras) ---
                
                % Conditional rendering based on output_mode
                \par\filbreak

\phantomsection

\addcontentsline{toc}{subsection}{ऋक् ३}

{\centering \textbf{\textcolor{AccentBlue}{ईमां\raisebox{0.25ex}{\accentmark{15}{\char"1CD2}}वृषा\raisebox{0.25ex}{\accentmark{12}{\char"0951}}णंकृणूतैकम्माम् \nolinebreak\hskip 0.08em \mbox{॥३॥}}} \par}

\vspace{0.8em}
                % --- SUBSECTIONS (Mantras) ---
                
                % Conditional rendering based on output_mode
                \par\filbreak

\phantomsection

\addcontentsline{toc}{subsection}{ऋक् ४}

{\centering \textbf{\textcolor{AccentBlue}{साना\raisebox{0.25ex}{\accentmark{12}{\char"0951}}इ\raisebox{0.15ex}{\accentmark{12}{\char"1CF9}}न्द्रा\raisebox{0.25ex}{\accentmark{12}{\char"0951}}यायाज्या\raisebox{0.25ex}{\accentmark{12}{\char"0951}}वे\kern 0.05em ।\hskip 0.12em plus 0.04em minus 0.02em वा\raisebox{0.25ex}{\accentmark{15}{\char"1CD2}}रू\raisebox{0.25ex}{\accentmark{12}{\char"0951}}णा\raisebox{0.25ex}{\accentmark{15}{\char"1CD2}}यामारू\raisebox{0.25ex}{\accentmark{15}{\char"1CD2}}त्भ्यः\kern 0.05em ।\hskip 0.12em plus 0.04em minus 0.02em वरिवोवि\raisebox{0.25ex}{\accentmark{15}{\char"1CD2}}त्पारि\raisebox{0.25ex}{\accentmark{12}{\char"0951}}स्रवा \nolinebreak\hskip 0.08em \mbox{॥४॥}}} \par}

\vspace{0.8em}
                % --- SUBSECTIONS (Mantras) ---
                
                % Conditional rendering based on output_mode
                \par\filbreak

\phantomsection

\addcontentsline{toc}{subsection}{ऋक् ५}

{\centering \textbf{\textcolor{AccentBlue}{एना\raisebox{0.15ex}{\accentmark{12}{\char"1CF9}}वि\raisebox{0.15ex}{\accentmark{12}{\char"1CF9}}श्वा\raisebox{0.25ex}{\accentmark{12}{\char"0951}}नर्या\raisebox{0.25ex}{\accentmark{12}{\char"0951}}आ\raisebox{0.25ex}{\accentmark{15}{\char"1CD2}}\kern 0.05em ।\hskip 0.12em plus 0.04em minus 0.02em द्युम्नानीमा\raisebox{0.25ex}{\accentmark{15}{\char"1CD2}}नूषाणां\kern 0.05em ।\hskip 0.12em plus 0.04em minus 0.02em सी\raisebox{0.25ex}{\accentmark{15}{\char"1CD2}}षा\raisebox{0.25ex}{\accentmark{12}{\char"0951}}सन्तोवनामहे\kern 0.05em ।\hskip 0.12em plus 0.04em minus 0.02em त्वामेताद\raisebox{0.25ex}{\accentmark{12}{\char"0951}}धारया\kern 0.05em ।\hskip 0.12em plus 0.04em minus 0.02em कृष्णासूरो\raisebox{0.25ex}{\accentmark{15}{\char"1CD2}}ही\raisebox{0.25ex}{\accentmark{15}{\char"1CD2}}णीषुच\kern 0.05em ।\hskip 0.12em plus 0.04em minus 0.02em पा\raisebox{0.25ex}{\accentmark{15}{\char"1CD2}}रू\raisebox{0.25ex}{\accentmark{12}{\char"0951}}ष्णीषूरू\raisebox{0.25ex}{\accentmark{15}{\char"1CD2}}शत्पा\raisebox{0.25ex}{\accentmark{15}{\char"1CD2}}याः \nolinebreak\hskip 0.08em \mbox{॥५॥}}} \par}

\vspace{0.8em}
                % --- SUBSECTIONS (Mantras) ---
                
                % Conditional rendering based on output_mode
                \par\filbreak

\phantomsection

\addcontentsline{toc}{subsection}{ऋक् ६}

{\centering \textbf{\textcolor{AccentBlue}{प्राथश्चायास्या\raisebox{0.25ex}{\accentmark{12}{\char"0951}}सा\raisebox{0.25ex}{\accentmark{12}{\char"0951}}प्रा\raisebox{0.15ex}{\accentmark{12}{\char"1CF9}}थश्चानां\kern 0.05em ।\hskip 0.12em plus 0.04em minus 0.02em आनु\raisebox{0.25ex}{\accentmark{12}{\char"0951}}ष्टभस्याहावी\raisebox{0.15ex}{\accentmark{12}{\char"1CF9}}षो\raisebox{0.25ex}{\accentmark{12}{\char"0951}}हावीर्यात्\kern 0.05em ।\hskip 0.12em plus 0.04em minus 0.02em धातू\raisebox{0.25ex}{\accentmark{15}{\char"1CD2}}र्द्यू\raisebox{0.25ex}{\accentmark{12}{\char"0951}}ता\raisebox{0.25ex}{\accentmark{12}{\char"0951}}नास्सवी\raisebox{0.25ex}{\accentmark{12}{\char"0951}}तुश्चा\raisebox{0.15ex}{\accentmark{12}{\char"1CF9}}वि\raisebox{0.25ex}{\accentmark{12}{\char"0951}}ष्णोः\kern 0.05em ।\hskip 0.12em plus 0.04em minus 0.02em रा\raisebox{0.25ex}{\accentmark{15}{\char"1CD2}}थन्तारामा\raisebox{0.25ex}{\accentmark{15}{\char"1CD2}}जभारावासीष्ठम् \nolinebreak\hskip 0.08em \mbox{॥६॥}}} \par}

\vspace{0.8em}
                % --- SUBSECTIONS (Mantras) ---
                
                % Conditional rendering based on output_mode
                \par\filbreak

\phantomsection

\addcontentsline{toc}{subsection}{ऋक् ७}

{\centering \textbf{\textcolor{AccentBlue}{इ\raisebox{0.25ex}{\accentmark{15}{\char"1CD2}}न्द्राइ\raisebox{0.25ex}{\accentmark{15}{\char"1CD2}}द्धरीयो\raisebox{0.15ex}{\accentmark{12}{\char"1CF8}}स्सा\raisebox{0.25ex}{\accentmark{15}{\char"1CD2}}चा\raisebox{0.25ex}{\accentmark{12}{\char"0951}}\kern 0.05em ।\hskip 0.12em plus 0.04em minus 0.02em सम्मिश्लोआ\raisebox{0.15ex}{\accentmark{12}{\char"1CF9}}वा\raisebox{0.25ex}{\accentmark{12}{\char"0951}}चोयू\raisebox{0.25ex}{\accentmark{15}{\char"1CD2}}जा\raisebox{0.25ex}{\accentmark{12}{\char"0951}}\kern 0.05em ।\hskip 0.12em plus 0.04em minus 0.02em इन्द्रो\raisebox{0.25ex}{\accentmark{12}{\char"0951}}वज्री\raisebox{0.15ex}{\accentmark{12}{\char"1CF9}}ही\raisebox{0.25ex}{\accentmark{12}{\char"0951}}रण्यायाः \nolinebreak\hskip 0.08em \mbox{॥७॥}}} \par}

\vspace{0.8em}
                % --- SUBSECTIONS (Mantras) ---
                
                % Conditional rendering based on output_mode
                \par\filbreak

\phantomsection

\addcontentsline{toc}{subsection}{ऋक् ८}

{\centering \textbf{\textcolor{AccentBlue}{इ\raisebox{0.25ex}{\accentmark{15}{\char"1CD2}}न्द्रावाजेषुनोअव\kern 0.05em ।\hskip 0.12em plus 0.04em minus 0.02em साहा\raisebox{0.25ex}{\accentmark{15}{\char"1CD2}}स्राप्रथनेषुच\kern 0.05em ।\hskip 0.12em plus 0.04em minus 0.02em उग्रा\raisebox{0.25ex}{\accentmark{15}{\char"1CD2}}उग्रा\raisebox{0.15ex}{\accentmark{12}{\char"1CF9}}भीरू\raisebox{0.25ex}{\accentmark{12}{\char"0951}}तीभीः\raisebox{0.25ex}{\accentmark{12}{\char"0951}} \nolinebreak\hskip 0.08em \mbox{॥८॥}}} \par}

\vspace{0.8em}
                % --- SUBSECTIONS (Mantras) ---
                
                % Conditional rendering based on output_mode
                \par\filbreak

\phantomsection

\addcontentsline{toc}{subsection}{ऋक् १०}

{\centering \textbf{\textcolor{AccentBlue}{आ\raisebox{0.25ex}{\accentmark{15}{\char"1CD2}}रूरूचादूषा\raisebox{0.25ex}{\accentmark{15}{\char"1CD2}}साः\raisebox{0.25ex}{\accentmark{12}{\char"0951}}\kern 0.05em ।\hskip 0.12em plus 0.04em minus 0.02em प्रश्नी\raisebox{0.25ex}{\accentmark{12}{\char"0951}}राग्रा\raisebox{0.25ex}{\accentmark{15}{\char"1CD2}}यूः\kern 0.05em ।\hskip 0.12em plus 0.04em minus 0.02em उक्षा\raisebox{0.25ex}{\accentmark{15}{\char"1CD2}}मीमा\raisebox{0.25ex}{\accentmark{15}{\char"1CD2}}तीभूवा\raisebox{0.25ex}{\accentmark{15}{\char"1CD2}}नेषुवाजयुः\kern 0.05em ।\hskip 0.12em plus 0.04em minus 0.02em मायाविनो\raisebox{0.25ex}{\accentmark{12}{\char"0951}}ममिरे\kern 0.05em ।\hskip 0.12em plus 0.04em minus 0.02em अ\raisebox{0.25ex}{\accentmark{15}{\char"1CD2}}स्यामाया\raisebox{0.25ex}{\accentmark{15}{\char"1CD2}}या\raisebox{0.25ex}{\accentmark{12}{\char"0951}}नृचाक्षा\raisebox{0.25ex}{\accentmark{12}{\char"0951}}सः\kern 0.05em ।\hskip 0.12em plus 0.04em minus 0.02em पीतारोग\raisebox{0.25ex}{\accentmark{15}{\char"1CD2}}र्भामाद\raisebox{0.25ex}{\accentmark{12}{\char"0951}}धुः \nolinebreak\hskip 0.08em \mbox{॥१०॥}}} \par}

\vspace{0.8em}

            % --- End of Kandah Closing Floral Line & Counts ---
            \par\nopagebreak\vspace{0.9em plus 0.3em minus 0.1em}
            \centerline{\textcolor{black!35}{\rule{1.8cm}{0.35pt}}\;\raisebox{-0.22ex}{\textcolor{vedicmaroon}{\large\SixFlowerPetalDotted}}\;\textcolor{black!35}{\rule{1.8cm}{0.35pt}}}
            \nopagebreak\vspace{0.45em}
            \centerline{\devafont\textbf{\textcolor{AccentBlue}{॥ ऋक् सङ्ख्या : ९ ॥}}}
        
            \clearpage

            % --- SECTION ---
            
            \phantomsection  

            \section*{\centering अथ षष्ठः खण्डः}
            \addcontentsline{toc}{section}{षष्ठः खण्डः} 
            \markright{षष्ठः खण्डः}
            \nopagebreak % Keep section title with first mantra
            
            % --- SECTION-LEVEL PROCEDURE LINK ---

                % --- SUBSECTIONS (Mantras) ---
                
                % Conditional rendering based on output_mode
                \par\filbreak

\phantomsection

\addcontentsline{toc}{subsection}{ऋक् १}

{\centering \textbf{\textcolor{AccentBlue}{यज्जा\raisebox{0.25ex}{\accentmark{15}{\char"1CD2}}या\raisebox{0.25ex}{\accentmark{12}{\char"0951}}था\raisebox{0.25ex}{\accentmark{12}{\char"0951}}पूर्व्यामाघा\raisebox{0.25ex}{\accentmark{12}{\char"0951}}वन्वृत्राहा\raisebox{0.25ex}{\accentmark{15}{\char"1CD2}}था\raisebox{0.25ex}{\accentmark{12}{\char"0951}}या\kern 0.05em ।\hskip 0.12em plus 0.04em minus 0.02em त\raisebox{0.15ex}{\accentmark{12}{\char"1CF8}}त्पृथिवीमाप्रा\raisebox{0.25ex}{\accentmark{12}{\char"0951}}था\raisebox{0.25ex}{\accentmark{12}{\char"0951}}या\kern 0.05em ।\hskip 0.12em plus 0.04em minus 0.02em ताद\raisebox{0.25ex}{\accentmark{12}{\char"0951}}स्तभ्नाऊतादीव\raisebox{0.25ex}{\accentmark{12}{\char"0951}}म् \nolinebreak\hskip 0.08em \mbox{॥१॥}}} \par}

\vspace{0.8em}
                % --- SUBSECTIONS (Mantras) ---
                
                % Conditional rendering based on output_mode
                \par\filbreak

\phantomsection

\addcontentsline{toc}{subsection}{ऋक् २}

{\centering \textbf{\textcolor{AccentBlue}{अग्नआयूं\raisebox{0.25ex}{\accentmark{12}{\char"0951}}षिपावसे\kern 0.05em ।\hskip 0.12em plus 0.04em minus 0.02em आसूर्वोर्जामी\raisebox{0.25ex}{\accentmark{15}{\char"1CD2}}षञ्चनः\kern 0.05em ।\hskip 0.12em plus 0.04em minus 0.02em आरे\raisebox{0.25ex}{\accentmark{15}{\char"1CD2}}बा\raisebox{0.25ex}{\accentmark{15}{\char"1CD2}}धास्वदु\raisebox{0.25ex}{\accentmark{12}{\char"0951}}च्छूना\raisebox{0.25ex}{\accentmark{12}{\char"0951}}म् \nolinebreak\hskip 0.08em \mbox{॥२॥}}} \par}

\vspace{0.8em}
                % --- SUBSECTIONS (Mantras) ---
                
                % Conditional rendering based on output_mode
                \par\filbreak

\phantomsection

\addcontentsline{toc}{subsection}{ऋक् ३}

{\centering \textbf{\textcolor{AccentBlue}{आयं\raisebox{0.25ex}{\accentmark{15}{\char"1CD2}}गौःप्रश्नी\raisebox{0.25ex}{\accentmark{12}{\char"0951}}रक्रामीत्\kern 0.05em ।\hskip 0.12em plus 0.04em minus 0.02em आ\raisebox{0.25ex}{\accentmark{15}{\char"1CD2}}सा\raisebox{0.25ex}{\accentmark{15}{\char"1CD2}}दन्माता\raisebox{0.15ex}{\accentmark{12}{\char"1CF9}}रं\raisebox{0.25ex}{\accentmark{12}{\char"0951}}पूनः\kern 0.05em ।\hskip 0.12em plus 0.04em minus 0.02em पी\raisebox{0.25ex}{\accentmark{15}{\char"1CD2}}तार\raisebox{0.25ex}{\accentmark{12}{\char"0951}}ञ्चाप्राया\raisebox{0.25ex}{\accentmark{15}{\char"1CD2}}न्स्वाः \nolinebreak\hskip 0.08em \mbox{॥३॥}}} \par}

\vspace{0.8em}
                % --- SUBSECTIONS (Mantras) ---
                
                % Conditional rendering based on output_mode
                \par\filbreak

\phantomsection

\addcontentsline{toc}{subsection}{ऋक् ४}

{\centering \textbf{\textcolor{AccentBlue}{अ\raisebox{0.25ex}{\accentmark{15}{\char"1CD2}}न्ता\raisebox{0.25ex}{\accentmark{15}{\char"1CD2}}श्चा\raisebox{0.25ex}{\accentmark{12}{\char"0951}}रतिरोचन\raisebox{0.25ex}{\accentmark{15}{\char"1CD2}}\kern 0.05em ।\hskip 0.12em plus 0.04em minus 0.02em अस्य\raisebox{0.25ex}{\accentmark{15}{\char"1CD2}}प्राणा\raisebox{0.25ex}{\accentmark{15}{\char"1CD2}}दपानातीव्य\raisebox{0.25ex}{\accentmark{15}{\char"1CD2}}ख्यन्मा\raisebox{0.25ex}{\accentmark{15}{\char"1CD2}}हीषोदीवम् \nolinebreak\hskip 0.08em \mbox{॥४॥}}} \par}

\vspace{0.8em}
                % --- SUBSECTIONS (Mantras) ---
                
                % Conditional rendering based on output_mode
                \par\filbreak

\phantomsection

\addcontentsline{toc}{subsection}{ऋक् ५}

{\centering \textbf{\textcolor{AccentBlue}{त्रिंश\raisebox{0.25ex}{\accentmark{15}{\char"1CD2}}द्धामावी\raisebox{0.25ex}{\accentmark{15}{\char"1CD2}}राजति\kern 0.05em ।\hskip 0.12em plus 0.04em minus 0.02em वा\raisebox{0.15ex}{\accentmark{12}{\char"1CF8}}क्पा\raisebox{0.25ex}{\accentmark{12}{\char"0951}}तंगाया\raisebox{0.25ex}{\accentmark{12}{\char"0951}}धीयते\kern 0.05em ।\hskip 0.12em plus 0.04em minus 0.02em प्रा\raisebox{0.25ex}{\accentmark{15}{\char"1CD2}}तीव\raisebox{0.25ex}{\accentmark{15}{\char"1CD2}}स्तोरा\raisebox{0.25ex}{\accentmark{15}{\char"1CD2}}हद्यू\raisebox{0.25ex}{\accentmark{15}{\char"1CD2}}भीः \nolinebreak\hskip 0.08em \mbox{॥५॥}}} \par}

\vspace{0.8em}

            % --- End of Kandah Closing Floral Line & Counts ---
            \par\nopagebreak\vspace{0.9em plus 0.3em minus 0.1em}
            \centerline{\textcolor{black!35}{\rule{1.8cm}{0.35pt}}\;\raisebox{-0.22ex}{\textcolor{vedicmaroon}{\large\SixFlowerPetalDotted}}\;\textcolor{black!35}{\rule{1.8cm}{0.35pt}}}
            \nopagebreak\vspace{0.45em}
            \centerline{\devafont\textbf{\textcolor{AccentBlue}{॥ ऋक् सङ्ख्या : ५ ॥}}}


    % --- SUPERSECTION ---
    \clearpage           
    \phantomsection              
    \chapter*{\centering अथ शाक्वर पर्वा}
    \addcontentsline{toc}{chapter}{शाक्वर पर्वा}
    \markboth{शाक्वर पर्वा}{शाक्वर पर्वा}
    
    % GLUE FIX: Prevent page break between Chapter and First Section
    \nopagebreak

        

            % --- SECTION ---
            
            \phantomsection  

            \section*{\centering अथ शाक्वर पर्वा}
            \addcontentsline{toc}{section}{शाक्वर पर्वा} 
            \markright{शाक्वर पर्वा}
            \nopagebreak % Keep section title with first mantra
            
            % --- SECTION-LEVEL PROCEDURE LINK ---

                % --- SUBSECTIONS (Mantras) ---
                
                % Conditional rendering based on output_mode
                \par\filbreak

\phantomsection

\addcontentsline{toc}{subsection}{ऋक् १}

{\centering \textbf{\textcolor{AccentBlue}{वीदा\raisebox{0.25ex}{\accentmark{15}{\char"1CD2}}मा\raisebox{0.25ex}{\accentmark{12}{\char"0951}}घवन्वीदा\raisebox{0.25ex}{\accentmark{15}{\char"1CD2}}गातूमा\raisebox{0.25ex}{\accentmark{15}{\char"1CD2}}नूशं\raisebox{0.25ex}{\accentmark{12}{\char"0951}}सीषोदीशाः\raisebox{0.25ex}{\accentmark{12}{\char"0951}}\kern 0.05em ।\hskip 0.12em plus 0.04em minus 0.02em शि\raisebox{0.25ex}{\accentmark{15}{\char"1CD2}}क्षा\raisebox{0.25ex}{\accentmark{12}{\char"0951}}शचीनांपतेपूर्वीणां\raisebox{0.25ex}{\accentmark{15}{\char"1CD2}}पू\raisebox{0.25ex}{\accentmark{12}{\char"0951}}रो\raisebox{0.15ex}{\accentmark{12}{\char"1CF8}}वसो\kern 0.05em ।\hskip 0.12em plus 0.04em minus 0.02em आ\raisebox{0.25ex}{\accentmark{15}{\char"1CD2}}भीष्ट्वा\raisebox{0.25ex}{\accentmark{15}{\char"1CD2}}मा\raisebox{0.15ex}{\accentmark{12}{\char"1CF9}}भीष्टी\raisebox{0.25ex}{\accentmark{12}{\char"0951}}भिःप्रा\raisebox{0.25ex}{\accentmark{15}{\char"1CD2}}चे\raisebox{0.25ex}{\accentmark{12}{\char"0951}}तानःप्रा\raisebox{0.25ex}{\accentmark{15}{\char"1CD2}}चेतायेन्द्रा\raisebox{0.15ex}{\accentmark{12}{\char"1CF9}}द्यु\raisebox{0.25ex}{\accentmark{12}{\char"0951}}म्ना\raisebox{0.25ex}{\accentmark{15}{\char"1CD2}}या\raisebox{0.25ex}{\accentmark{12}{\char"0951}}नाईषे\kern 0.05em ।\hskip 0.12em plus 0.04em minus 0.02em एवा\raisebox{0.25ex}{\accentmark{15}{\char"1CD2}}हीशक्रो\raisebox{0.25ex}{\accentmark{15}{\char"1CD2}}रायेवा\raisebox{0.25ex}{\accentmark{15}{\char"1CD2}}जा\raisebox{0.25ex}{\accentmark{12}{\char"0951}}यवज्रीवः\kern 0.05em ।\hskip 0.12em plus 0.04em minus 0.02em शा\raisebox{0.25ex}{\accentmark{15}{\char"1CD2}}वीष्ठवज्रिन्वृञ्ज\raisebox{0.25ex}{\accentmark{15}{\char"1CD2}}सेमंहिष्ठ\raisebox{0.25ex}{\accentmark{15}{\char"1CD2}}वज्रिन्वृन्जसाआ\raisebox{0.15ex}{\accentmark{12}{\char"1CF9}}या\raisebox{0.25ex}{\accentmark{12}{\char"0951}}हीपीबामात्स्वा \nolinebreak\hskip 0.08em \mbox{॥१॥}}} \par}

\vspace{0.8em}
                % --- SUBSECTIONS (Mantras) ---
                
                % Conditional rendering based on output_mode
                \par\filbreak

\phantomsection

\addcontentsline{toc}{subsection}{ऋक् २}

{\centering \textbf{\textcolor{AccentBlue}{वीदाराये\raisebox{0.25ex}{\accentmark{15}{\char"1CD2}}सूवी\raisebox{0.25ex}{\accentmark{15}{\char"1CD2}}र्यं\raisebox{0.25ex}{\accentmark{12}{\char"0951}}\kern 0.05em ।\hskip 0.12em plus 0.04em minus 0.02em भूवोवा\raisebox{0.15ex}{\accentmark{12}{\char"1CF9}}जा\raisebox{0.25ex}{\accentmark{12}{\char"0951}}नांपातिर्वाशंआ\raisebox{0.25ex}{\accentmark{15}{\char"1CD2}}नू\kern 0.05em ।\hskip 0.12em plus 0.04em minus 0.02em मंहिष्ठवज्रिन्वृञ्ज\raisebox{0.25ex}{\accentmark{15}{\char"1CD2}}सेयश्शावी\raisebox{0.25ex}{\accentmark{12}{\char"0951}}ष्ठश्शू\raisebox{0.15ex}{\accentmark{12}{\char"1CF9}}रा\raisebox{0.15ex}{\accentmark{12}{\char"1CF9}}णाम्\kern 0.05em ।\hskip 0.12em plus 0.04em minus 0.02em योमं\raisebox{0.25ex}{\accentmark{15}{\char"1CD2}}ही\raisebox{0.25ex}{\accentmark{12}{\char"0951}}ष्ठोमाघो\raisebox{0.25ex}{\accentmark{15}{\char"1CD2}}नां\raisebox{0.25ex}{\accentmark{12}{\char"0951}}\kern 0.05em ।\hskip 0.12em plus 0.04em minus 0.02em चिकित्वोआभी\raisebox{0.15ex}{\accentmark{12}{\char"1CF9}}नो\raisebox{0.25ex}{\accentmark{12}{\char"0951}}नाये\raisebox{0.25ex}{\accentmark{15}{\char"1CD2}}न्द्रोवीदेता\raisebox{0.25ex}{\accentmark{15}{\char"1CD2}}मूस्तुही\kern 0.05em ।\hskip 0.12em plus 0.04em minus 0.02em ई\raisebox{0.25ex}{\accentmark{15}{\char"1CD2}}शे\raisebox{0.15ex}{\accentmark{12}{\char"1CF9}}ही\raisebox{0.25ex}{\accentmark{12}{\char"0951}}शक्र\raisebox{0.15ex}{\accentmark{12}{\char"1CF9}}स्ता\raisebox{0.25ex}{\accentmark{12}{\char"0951}}मूताये\raisebox{0.25ex}{\accentmark{12}{\char"0951}}हावा\raisebox{0.25ex}{\accentmark{12}{\char"0951}}माहे\raisebox{0.15ex}{\accentmark{12}{\char"1CF9}}जेता\raisebox{0.25ex}{\accentmark{12}{\char"0951}}रामापा\raisebox{0.25ex}{\accentmark{12}{\char"0951}}राजितं\kern 0.05em ।\hskip 0.12em plus 0.04em minus 0.02em सा\raisebox{0.25ex}{\accentmark{15}{\char"1CD2}}नाः\raisebox{0.25ex}{\accentmark{12}{\char"0951}}परीषा\raisebox{0.25ex}{\accentmark{15}{\char"1CD2}}द\raisebox{0.25ex}{\accentmark{15}{\char"1CD2}}तिद्वीषः\raisebox{0.25ex}{\accentmark{12}{\char"0951}}क्रातुच्छन्दोऋतंबृहात् \nolinebreak\hskip 0.08em \mbox{॥२॥}}} \par}

\vspace{0.8em}
                % --- SUBSECTIONS (Mantras) ---
                
                % Conditional rendering based on output_mode
                \par\filbreak

\phantomsection

\addcontentsline{toc}{subsection}{ऋक् ३}

{\centering \textbf{\textcolor{AccentBlue}{इन्द्रंधानास्यासाताये\kern 0.05em ।\hskip 0.12em plus 0.04em minus 0.02em हावा\raisebox{0.25ex}{\accentmark{12}{\char"0951}}महेजेतारामापाराजित\raisebox{0.25ex}{\accentmark{15}{\char"1CD2}}म्\kern 0.05em ।\hskip 0.12em plus 0.04em minus 0.02em सा\raisebox{0.25ex}{\accentmark{15}{\char"1CD2}}नाः\raisebox{0.25ex}{\accentmark{12}{\char"0951}}पारीषद\raisebox{0.25ex}{\accentmark{12}{\char"0951}}तिद्वीषस्सा\raisebox{0.25ex}{\accentmark{15}{\char"1CD2}}नाः\raisebox{0.25ex}{\accentmark{12}{\char"0951}}पा\raisebox{0.25ex}{\accentmark{15}{\char"1CD2}}री\raisebox{0.25ex}{\accentmark{12}{\char"0951}}षाद\raisebox{0.25ex}{\accentmark{15}{\char"1CD2}}तिद्वीषाः\raisebox{0.25ex}{\accentmark{12}{\char"0951}}\kern 0.05em ।\hskip 0.12em plus 0.04em minus 0.02em पू\raisebox{0.15ex}{\accentmark{12}{\char"1CF8}}र्वास्याया\raisebox{0.25ex}{\accentmark{12}{\char"0951}}क्ते\raisebox{0.25ex}{\accentmark{12}{\char"0951}}अद्रिवः\kern 0.05em ।\hskip 0.12em plus 0.04em minus 0.02em सुम्नाधे\raisebox{0.25ex}{\accentmark{12}{\char"0951}}हिनो\raisebox{0.25ex}{\accentmark{12}{\char"0951}}वसो\kern 0.05em ।\hskip 0.12em plus 0.04em minus 0.02em पूर्तीश्शा\raisebox{0.25ex}{\accentmark{15}{\char"1CD2}}वीष्ठशस्य\raisebox{0.25ex}{\accentmark{15}{\char"1CD2}}ते\kern 0.05em ।\hskip 0.12em plus 0.04em minus 0.02em व\raisebox{0.25ex}{\accentmark{15}{\char"1CD2}}शीहीशक्रोनून\raisebox{0.15ex}{\accentmark{12}{\char"1CF9}}न्त\raisebox{0.15ex}{\accentmark{12}{\char"1CF9}}न्नव्यं\raisebox{0.25ex}{\accentmark{12}{\char"0951}}सन्या\raisebox{0.25ex}{\accentmark{15}{\char"1CD2}}से\kern 0.05em ।\hskip 0.12em plus 0.04em minus 0.02em प्रा\raisebox{0.25ex}{\accentmark{15}{\char"1CD2}}भो\raisebox{0.25ex}{\accentmark{15}{\char"1CD2}}र्जनस्यवृत्रा\raisebox{0.25ex}{\accentmark{12}{\char"0951}}ह\raisebox{0.25ex}{\accentmark{15}{\char"1CD2}}न्सामर्ये\raisebox{0.25ex}{\accentmark{15}{\char"1CD2}}षु\raisebox{0.25ex}{\accentmark{12}{\char"0951}}स्तवामहे\kern 0.05em ।\hskip 0.12em plus 0.04em minus 0.02em शू\raisebox{0.25ex}{\accentmark{15}{\char"1CD2}}रोयो\raisebox{0.25ex}{\accentmark{15}{\char"1CD2}}गोषूगच्छा\raisebox{0.25ex}{\accentmark{12}{\char"0951}}ति\kern 0.05em ।\hskip 0.12em plus 0.04em minus 0.02em सा\raisebox{0.15ex}{\accentmark{12}{\char"1CF8}}खा\raisebox{0.25ex}{\accentmark{12}{\char"0951}}सूशे\raisebox{0.25ex}{\accentmark{15}{\char"1CD2}}वोआद्रा\raisebox{0.25ex}{\accentmark{15}{\char"1CD2}}युः\raisebox{0.25ex}{\accentmark{15}{\char"1CD2}} \nolinebreak\hskip 0.08em \mbox{॥३॥}}} \par}

\vspace{0.8em}
                % --- SUBSECTIONS (Mantras) ---
                
                % Conditional rendering based on output_mode
                \par\filbreak

\phantomsection

\addcontentsline{toc}{subsection}{ऋक् ४}

{\centering \textbf{\textcolor{AccentBlue}{एवा\raisebox{0.25ex}{\accentmark{15}{\char"1CD2}}ह्ये\raisebox{0.25ex}{\accentmark{12}{\char"0951}}वाएवाह्याग्ने\raisebox{0.25ex}{\accentmark{12}{\char"0951}}एवाही\raisebox{0.25ex}{\accentmark{12}{\char"0951}}न्द्राएवा\raisebox{0.25ex}{\accentmark{15}{\char"1CD2}}ही\raisebox{0.25ex}{\accentmark{15}{\char"1CD2}}पूष\raisebox{0.25ex}{\accentmark{12}{\char"0951}}न्एवाही\raisebox{0.25ex}{\accentmark{15}{\char"1CD2}}दे\raisebox{0.25ex}{\accentmark{12}{\char"0951}}वाः \nolinebreak\hskip 0.08em \mbox{॥४॥}}} \par}

\vspace{0.8em}

            % --- End of Kandah Closing Floral Line & Counts ---
            \par\nopagebreak\vspace{0.9em plus 0.3em minus 0.1em}
            \centerline{\textcolor{black!35}{\rule{1.8cm}{0.35pt}}\;\raisebox{-0.22ex}{\textcolor{vedicmaroon}{\large\SixFlowerPetalDotted}}\;\textcolor{black!35}{\rule{1.8cm}{0.35pt}}}
            \nopagebreak\vspace{0.45em}
            \centerline{\devafont\textbf{\textcolor{AccentBlue}{॥ ऋक् सङ्ख्या : ४ ॥}}}


% --- Summary Table ---
\clearpage
\phantomsection
\markboth{पूर्वार्चिकम् सङ्ख्या}{पूर्वार्चिकम् सङ्ख्या}
\fancyhead[LO]{पूर्वार्चिकम् सङ्ख्या}
\fancyhead[RE]{पूर्वार्चिकम् सङ्ख्या}
\chapter*{\centering ॥ पूर्वार्चिकम् सङ्ख्या ॥}
\addcontentsline{toc}{chapter}{पूर्वार्चिकम् सङ्ख्या}

\begin{center}
{\small % Reduce font size for the table
\renewcommand{\arraystretch}{1.3} % Increase vertical spacing between rows
\begin{longtable}{|p{6cm}|c|}
\hline
\textbf{खण्डः}  & \textbf{ऋक् सङ्ख्या}   \\
\hline\hline
\endfirsthead

\hline
\textbf{खण्डः}  & \textbf{ऋक् सङ्ख्या}   \\
\hline\hline
\endhead

\hline
\endfoot

\textbf{\textcolor{AccentBlue}{आग्नेयपाठः}}  & \textbf{\textcolor{AccentBlue}{११६}}   \\
\hline
प्रथम खण्डः  & १०   \\
द्वितीय खण्डः  & १०   \\
तृतीय खण्डः  & १४   \\
चतुर्थ खण्डः  & १०   \\
पञ्चम खण्डः  & १०   \\
षष्ठ खण्डः  & ८   \\
सप्तम खण्डः  & १०   \\
अष्टम खण्डः  & ८   \\
नवम खण्डः  & १०   \\
दशम खण्डः  & ६   \\
एकादश खण्डः  & १०   \\
द्वादश खण्डः  & १०   \\
\hline
\textbf{\textcolor{AccentBlue}{तद्वपाठः}}  & \textbf{\textcolor{AccentBlue}{११८}}   \\
\hline
प्रथम खण्डः  & १०   \\
द्वितीय खण्डः  & १०   \\
तृतीय खण्डः  & १०   \\
चतुर्थ खण्डः  & १०   \\
पञ्चम खण्डः  & १०   \\
षष्ठ खण्डः  & १०   \\
सप्तम खण्डः  & १०   \\
अष्टम खण्डः  & ९   \\
नवम खण्डः  & १०   \\
दशम खण्डः  & १०   \\
एकादश खण्डः  & ९   \\
द्वादश खण्डः  & १०   \\
\hline
\textbf{\textcolor{AccentBlue}{बृहतिपाठः}}  & \textbf{\textcolor{AccentBlue}{८०}}   \\
\hline
प्रथम खण्डः  & १०   \\
द्वितीय खण्डः  & १०   \\
तृतीय खण्डः  & १०   \\
चतुर्थ खण्डः  & १०   \\
पञ्चम खण्डः  & १०   \\
षष्ठ खण्डः  & १०   \\
सप्तम खण्डः  & १०   \\
अष्टम खण्डः  & १०   \\
\hline
\textbf{\textcolor{AccentBlue}{असाविपाठः}}  & \textbf{\textcolor{AccentBlue}{५७}}   \\
\hline
प्रथम खण्डः  & १०   \\
द्वितीय खण्डः  & ९   \\
तृतीय खण्डः  & १०   \\
चतुर्थ खण्डः  & १०   \\
पञ्चम खण्डः  & ८   \\
षष्ठ खण्डः  & १०   \\
\hline
\textbf{\textcolor{AccentBlue}{ऐन्द्रपाठः}}  & \textbf{\textcolor{AccentBlue}{९७}}   \\
\hline
प्रथम खण्डः  & ११   \\
द्वितीय खण्डः  & १०   \\
तृतीय खण्डः  & ८   \\
चतुर्थ खण्डः  & १०   \\
पञ्चम खण्डः  & १०   \\
षष्ठ खण्डः  & ८   \\
सप्तम खण्डः  & १०   \\
अष्टम खण्डः  & १०   \\
नवम खण्डः  & १०   \\
दशम खण्डः  & १०   \\
\hline
\textbf{\textcolor{AccentBlue}{पवमानपाठः}}  & \textbf{\textcolor{AccentBlue}{१२०}}   \\
\hline
प्रथम खण्डः  & १०   \\
द्वितीय खण्डः  & १०   \\
तृतीय खण्डः  & १०   \\
चतुर्थ खण्डः  & १४   \\
पञ्चम खण्डः  & १२   \\
षष्ठ खण्डः  & १०   \\
सप्तम खण्डः  & १२   \\
अष्टम खण्डः  & ९   \\
नवम खण्डः  & १२   \\
दशम खण्डः  & १२   \\
एकादश खण्डः  & ९   \\
\hline
\textbf{\textcolor{AccentBlue}{आरण पाठः}}  & \textbf{\textcolor{AccentBlue}{५२}}   \\
\hline
प्रथमः खण्डः  & ९   \\
द्वितीयः खण्डः  & ९   \\
तृतीयः खण्डः  & १०   \\
चतुर्थः खण्डः  & १०   \\
पञ्चमः खण्डः  & ९   \\
षष्ठः खण्डः  & ५   \\
\hline
\textbf{\textcolor{AccentBlue}{शाक्वर पर्वा}}  & \textbf{\textcolor{AccentBlue}{४}}   \\
\hline
शाक्वर पर्वा  & ४   \\
\hline
\textbf{आहत्य}  & \textbf{६४४}   \\
\hline
\end{longtable}
}
\end{center}
\clearpage
\fancyhead[LO]{\rightmark}
\fancyhead[RE]{\leftmark}

\clearpage

  \printindex

\end{document}